% English translation, made from our own transcription (original.tex), not from any existing
% translation. All mathematics, equation numbers, footnote markers, \origpage markers and figure
% paths are preserved unchanged; only the prose (and prose set inside \text{...}, HOUSESTYLE R21)
% is translated. Legendre's own period terms are rendered literally (see glossary.yaml); the
% transcriber's \ednote notes on printed misprints are carried over unchanged.
\documentclass{article}
\usepackage{readmasters}
\begin{document}

\origpage{i}

\section*{EXERCISES IN INTEGRAL CALCULUS ON VARIOUS ORDERS OF TRANSCENDENTS AND ON QUADRATURES;}

By A. M. LEGENDRE, Knight of the Empire, Member of the Institute and of the Legion of Honour, Titular Councillor of the Imperial University, etc.

PARIS,

Madame the Widow COURCIER, PRINTER-BOOKSELLER FOR MATHEMATICS.

1811.

\origpage{1}

\section*{EXERCISES IN INTEGRAL CALCULUS.}

\section*{FIRST PART.}

\section*{ON ELLIPTIC FUNCTIONS.}

Having exhausted the differential formulas that can be integrated algebraically as well as by arcs of a circle or by logarithms, the geometers turned to seeking out all those that are integrable by arcs of an ellipse or by arcs of a hyperbola${}^{(1)}$. There was reason to believe that these transcendents held the first rank after the circular and logarithmic functions, and it mattered for the progress of the new calculus to bring back to a well-known degree of difficulty all the integrals that were susceptible of this reduction.

The formulas that can be integrated in this way proved very numerous; but there was no connection between the results, and they were far from forming a theory.

An Italian geometer of great sagacity opened the way to deeper speculations${}^{(2)}$. He proved that on any ellipse or

\textbf{(1)} \emph{Maclaurin.} Traité des Fluxions.~--- \emph{D'Alembert.} Mém.\ de Berlin. 1746.

\textbf{(2)} \emph{Fagnani.} Produzioni matematiche, Tom.\ II.

\origpage{2}
on any given hyperbola one can assign, in an infinity of ways, two arcs whose difference is equal to an algebraic quantity. He demonstrated at the same time that the \emph{lemniscate} enjoys this singular property, that its arcs can be multiplied or divided algebraically, like the arcs of a circle, although each of them is a transcendent of a higher order.

Euler, by a combination that may be regarded as most fortunate, although such strokes of chance never befall any but those who know how to bring them about, found the complete algebraic integral of a differential equation composed of two separate but similar terms, each of which is integrable only by arcs of conic sections${}^{(1)}$.

This important discovery gave its author occasion to compare, in a more general manner than had been done before him, not only the arcs of one and the same ellipse or one and the same hyperbola, but in general all the transcendents comprised in the formula $\displaystyle\int\frac{P\,dx}{R}$, where $P$ is a rational function of $x$, and $R$ a radical of the form $\sqrt{(\alpha+\beta x+\gamma x^{2}+\delta x^{3}+\varepsilon x^{4})}$, $\alpha$, $\beta$, $\gamma$, $\delta$, $\varepsilon$, being constants.

The integral found by Euler was too remarkable not to fix the particular attention of the geometers. Lagrange wished to bring this integration back within the ordinary procedures of analysis; he succeeded in it by a very ingenious method${}^{(2)}$, whose application rises gradually from the lower transcendents to the Eulerian transcendents; but he tried in vain to arrive at a result more general than Euler's.

Shortly afterwards, Landen, an English geometer, demonstrated that any arc of a hyperbola can be measured by two arcs of an ellipse${}^{(3)}$; a memorable discovery, which reduces to arcs of an ellipse alone all the integrals that until then had been expressible only by the rectification of the two curves.

Finally Lagrange distinguished himself anew in the same career,

\textbf{(1)} \emph{Euler.} Novi Com.\ Petrop.\ Tom.\ VI et VII.

\textbf{(2)} Mém.\ de Turin, Tom.\ IV.

\textbf{(3)} Mathematical Memoirs, by John Landen, 1780.

\origpage{3}
by giving a general method${}^{(1)}$ for bringing back, by successive transformations, the integral $\displaystyle\int\frac{P\,dx}{R}$ to the integral of a similar formula which, by the arrangement of its coefficients, is easy to evaluate by approximation. These transformations have the twofold aim of serving for the comparison of a sequence of transcendents formed according to the same law, and of leading to the most rapid approximations of which these functions are susceptible.

Such were the principal discoveries of the geometers in the theory of the transcendents designated by $\displaystyle\int\frac{P\,dx}{R}$, when I published my Researches on integration by arcs of an ellipse${}^{(2)}$. The first part had been composed before I had knowledge of Landen's theorem; it contained new views on the use of arcs of an ellipse, and in particular a means of avoiding the use of arcs of a hyperbola in the integral calculus, by supplying their place with a suitably constructed table of arcs of an ellipse. I then gave a new demonstration of Landen's theorem, and I proved by the same method that any given ellipse forms part of an infinite sequence of ellipses so connected among themselves that by the rectification of two of these ellipses, taken at will, one obtains the rectification of all the others. These ellipses having a common semi-major axis equal to unity, and their eccentricities varying according to a known law from zero to unity, one can by this theorem reduce the rectification of a given ellipse to that of two other ellipses differing as little from the circle as one pleases. This was one step further in a difficult career.

But this subject, and in general the theory of the transcendents designated by $\displaystyle\int\frac{P\,dx}{R}$, called for treatment in a more methodical and more thorough manner. This is what I attempted to do in my \emph{Memoir on Elliptic Transcendents}, published in 1793. I proposed in that work to compare with one another all the functions comprised under this denomination, to classify them into different species, to reduce each to the simplest form

\textbf{(1)} Nouveaux Mémoires de Turin, ann.\ 1784 et 1785, Tom.\ II.

\textbf{(2)} Mém.\ de l'Acad.\ des Sciences de Paris, ann.\ 1786.

\origpage{4}
of which it is susceptible, to evaluate them by the promptest and easiest approximations; finally, to form from the whole of this theory a sort of algorithm that might contribute to extending the domain of analysis.

Having resumed the sequel of these researches after a long interruption, I have succeeded in perfecting this theory in some parts, chiefly in that which concerns the elliptic functions of the third species. These improvements were of a nature to bring some change to my first work; I therefore thought I should treat this subject in a new order, giving it greater developments and illustrating it by chosen examples. It is the result of this latest work that I present at this moment to the geometers; I hope they will be good enough to receive it as a new branch of analysis that may offer fine and useful applications.

\section*{General idea of the different kinds of transcendents contained in the integral formula $\displaystyle\int\frac{P\,dx}{R}$.}

\subsection*{(1).}

We represent by $P$ any rational function whatever of $x$, and by $R$ the radical $\sqrt{(\alpha+\beta x+\gamma x^{2}+\delta x^{3}+\varepsilon x^{4})}$: this radical remaining the same, an infinity of values can be given to $P$; but there does not thereby result an infinity of transcendents of different nature. One can always, by partial integrations, reduce the integral $\displaystyle\int\frac{P\,dx}{R}$ to an algebraic part, plus a certain number of transcendents, which are always of the same form and of the same nature. This is what we have to develop.

Suppose first that $P$ is an integral function of $x$, so that we have $P=A+Bx+Cx^{2}+\ldots+Kx^{k}$. If we represent, for brevity, the integral $\displaystyle\int\frac{x^{m}\,dx}{R}$ by $\Pi^{m}$, it is clear that we shall have
\[ \int\frac{P\,dx}{R}=A\Pi^{0}+B\Pi^{1}+C\Pi^{2}+\ldots+K\Pi^{k}: \]
now, by differentiating the quantity $x^{m-3}R$, and returning from the differential to the integral, we find this formula
\[ \begin{aligned} x^{m-3}R={}&(m-3)\alpha\Pi^{m-4}+(m-\tfrac{5}{2})\beta\Pi^{m-3}+(m-2)\gamma\Pi^{m-2}\\ &+(m-\tfrac{3}{2})\delta\Pi^{m-1}+(m-1)\varepsilon\Pi^{m}; \end{aligned} \]

\origpage{5}
whence it follows that $\Pi^{m}$, $\Pi^{m-1}$, etc.\ can be successively lowered to inferior degrees, and that the reduction can be continued until we have only quantities below $\Pi^{3}$; for the reduction of $\Pi^{3}$ is still possible, because in the case of $m=3$, the term that would seem to have to contain $\Pi^{-1}$ vanishes. Therefore, $P$ being an integral function of $x$, the integral $\displaystyle\int\frac{P\,dx}{R}$ can always be reduced to an algebraic quantity, plus a transcendent which will constantly be of the form
\[ \int(A+Bx+Cx^{2})\frac{dx}{R}. \]

\subsection*{(2).}

Let us now consider the formula in all its generality, and let $P$ be any rational function of $x$: we shall proceed as if it were a question of integrating the rational fraction $P\,dx$; we shall first extract the integral part that may be contained in it, and this part will be treated as has just been said. We shall then decompose the remainder into several partial fractions, according to the number of factors of the denominator; from this there will result as many terms in the total integral, and each of these terms can be represented by the general expression $\displaystyle N\int\frac{dx}{(1+nx)^{k}R}$. Now, it is easy to reduce all terms of this species to similar terms, where $k=1$; this we shall prove in a moment. Let us observe beforehand that if the denominator, instead of being complex, were simply $x^{k}$, the integral $\displaystyle\int\frac{dx}{x^{k}R}$, and all similar ones where $k>1$, could be determined by means of the integral $\displaystyle\int\left(\frac{A}{x}+B+Cx+Dx^{2}\right)\frac{dx}{R}$, it being sufficient for this to give negative values to $m-4$ in the formula of article (1).

Let us return to the general term $\displaystyle\int\frac{dx}{(1+nx)^{k}R}$, which we shall call $\Gamma^{k}$. If we put $1+nx=\omega$, and take the coefficients $\alpha'$, $\beta'$, etc., in such a way that we have the identical equation
\[ \begin{aligned} &\alpha+\beta\left(\frac{\omega-1}{n}\right)+\gamma\left(\frac{\omega-1}{n}\right)^{2}+\delta\left(\frac{\omega-1}{n}\right)^{3}+\varepsilon\left(\frac{\omega-1}{n}\right)^{4}\\ &=\alpha'+\beta'\omega+\gamma'\omega^{2}+\delta'\omega^{3}+\varepsilon'\omega^{4}, \end{aligned} \]

\origpage{6}
we shall find, by differentiating the quantity $\omega^{-k+1}R$, this general formula:
\[ \begin{aligned} \frac{-R}{n(1+nx)^{k-1}}={}&(k-1)\alpha'\Gamma^{k}+(k-\tfrac{3}{2})\beta'\Gamma^{k-1}+(k-2)\gamma'\Gamma^{k-2}\\ &+(k-\tfrac{5}{2})\delta'\Gamma^{k-3}+(k-3)\varepsilon'\Gamma^{k-4}. \end{aligned} \]
Whence it follows that the quantities $\Gamma^{2}$, $\Gamma^{3}$, etc., can be determined by means of $\Gamma^{1}$, $\Gamma^{0}$, $\Gamma^{-1}$, $\Gamma^{-2}$; now, $\displaystyle\Gamma^{0}=\int\frac{dx}{R}$, $\displaystyle\Gamma^{-1}=\int(1+nx)\frac{dx}{R}$, $\displaystyle\Gamma^{-2}=\int(1+nx)^{2}\frac{dx}{R}$. Therefore in general the formula $\displaystyle\int\frac{dx}{(1+nx)^{k}R}$, and all similar formulas in which $k$ is greater than unity, will always reduce to an algebraic part, plus an integral of the form
\[ \int\left(\frac{A}{1+nx}+B+Cx+Dx^{2}\right)\frac{dx}{R}. \]

\subsection*{(3).}

We conclude from this that, whatever the rational function of $x$ represented by $P$, the integral $\displaystyle\int\frac{P\,dx}{R}$ will always be decomposable into three principal parts, the first algebraic, the second of the form $\displaystyle\int(A+Bx+Cx^{2})\frac{dx}{R}$, and the third containing one or more terms of the form $\displaystyle N\int\frac{dx}{(1+nx)R}$, where the coefficient $n$ may be real or imaginary.

It is seen by this analysis that the number of transcendents comprised in the formula $\displaystyle\int\frac{P\,dx}{R}$ is very limited. There exist only two principal species, one of the form $\displaystyle\int(A+Bx+Cx^{2})\frac{dx}{R}$, the other of the form $\displaystyle\int\frac{dx}{(1+nx)R}$. I observe even that as long as $n$ is real, this second species is comprised in the first, and reduces to it immediately, by putting $1+nx=\dfrac{1}{z}$.

These general reductions once perceived, we shall follow another route in order to arrive at a more precise knowledge of the same transcendents.

\origpage{7}

\section*{Manner of making the odd powers of the variable disappear under the radical.}

\subsection*{(4).}

The radical being always $\sqrt{(\alpha+\beta x+\gamma x^{2}+\delta x^{3}+\varepsilon x^{4})}$, let us suppose that the quantity under the sign is decomposed into two real factors $\zeta+2\eta x+\theta x^{2}$, $\lambda+2\mu x+\nu x^{2}$. These factors must be of the same sign for the radical to be real; thus we may suppose
\[ \zeta+2\eta x+\theta x^{2}=(\lambda+2\mu x+\nu x^{2})y^{2}. \]
From this we draw,
\[ x=\frac{\mu y^{2}-\eta+\sqrt{[(\mu y^{2}-\eta)^{2}+(\lambda y^{2}-\zeta)(\theta-\nu y^{2})]}}{\theta-\nu y^{2}}; \]
and if we call, for brevity, $Y$ the radical of this expression, we shall have $\dfrac{dx}{R}=\dfrac{dy}{Y}$. The value of $x$ being substituted in $P$, we see that the differential $\dfrac{P\,dx}{R}$ or $\dfrac{P\,dy}{Y}$ will always decompose into two parts, one rational and integrable by the ordinary rules, the other affected by the radical $Y$, but such that there will be only even powers of $y$ under the radical and outside the radical. Thus the integration of the formula $\dfrac{P\,dx}{R}$ always reduces to that of a formula of the same nature, in which $P$ and $R$ contain only even powers of $x$.

\subsection*{(5).}

This method is general; however, since the value of $x$ which is to be substituted in $P$ is somewhat complicated, it will not be useless to show that one can reach the same end by a much simpler substitution, which consists in putting $x=\dfrac{p+qy}{1+y}$, $p$ and $q$ being two indeterminate constants.

Let us take up again the factors $\zeta+2\eta x+\theta x^{2}$, $\lambda+2\mu x+\nu x^{2}$, and substitute in each of them the value of $x$; we may disregard the common denominator which comes out of the radical, and in order that the

\origpage{8}
odd powers of $y$ may disappear in the numerators, we must have
\[ \begin{aligned} &\zeta+\eta(p+q)+\theta pq=0,\\ &\lambda+\mu(p+q)+\nu pq=0. \end{aligned} \]
These two equations will give rational values of $p+q$ and $pq$; but in order that $p$ and $q$ be real, it is further necessary that $(p+q)^{2}-4pq$, or $(p-q)^{2}$, be positive. Now, two cases may be distinguished:

1°. If the quantity $\alpha+\beta x+\gamma x^{2}+\text{etc.}$ does not have all its factors real, we may suppose that $\lambda+2\mu x+\nu x^{2}$ contains two imaginary ones, and that we have in consequence $\lambda\nu>\mu^{2}$. But the second of the equations in $p$ and $q$ gives
\[ (p-q)^{2}=\left(\frac{\lambda+\nu pq}{\mu}-\frac{2\mu}{\nu}\right)^{2}+\frac{4\lambda\nu-4\mu^{2}}{\nu^{2}}; \]
hence in this first case the second member is positive, and the values of $p$ and $q$ will be real.

2°. If the quantity $\alpha+\beta x+\text{etc.}$ has all its factors real, and if, in decomposing it into two factors of the second degree, as we have just done, there do not result real values for $p$ and $q$; then it will be observed that there are two other ways of decomposing the quantity of the fourth degree into two factors of the second, and one may be assured that these two new combinations will give real values for $p$ and $q$. Here is the demonstration.

Let $a$, $b$, $c$, $d$ be the four values of $x$ that are obtained by putting $R=0$; the equations which determine $p$ and $q$ can be written thus:
\[ \begin{aligned} &ab-\tfrac{1}{2}(a+b)(p+q)+pq=0.\\ &cd-\tfrac{1}{2}(c+d)(p+q)+pq=0, \end{aligned} \]
\ednote{The print ends the first of these two equations with a period, where a comma is expected.}
and from them it results that
\[ \begin{aligned} &\frac{p+q}{2}=\frac{ab-cd}{a+b-c-d},\\ &\left(\frac{p-q}{2}\right)^{2}=\frac{a-c.a-d.b-c.b-d}{(a+b-c-d)^{2}}. \end{aligned} \]

\origpage{9}
Now, one is master of making the numerator of this last quantity positive; for suppose that $a$, $b$, $c$, $d$ are written in the order of their magnitude, the negative roots coming after the positive ones, then the differences $a-b$, $a-c$, etc.\ will all be positive, and $p-q$ will be real. We shall again have $p-q$ real, if we take for $a$ and $b$ the extremes in magnitude of the roots $a$, $b$, $c$, $d$. Finally, the third combination would give $p-q$ imaginary.

Therefore by the substitution of $x=\dfrac{p+qy}{1+y}$, it is always possible to make the odd powers of the variable disappear under the radical, and at the same time one obtains this advantage, that the two factors, under the new radical, are real and of the form $f+gy^{2}$, which is an essential point, and one that would not always be obtained by the first method.

Let us remark that there are two particular cases to examine.

1°. When one of the roots $a$ and $b$ is equal to one of the roots $c$ and $d$, the radical $R$ simplifies, and the integral depends only on arcs of a circle and logarithms.

2°. If we have $a+b=c+d$, a simple permutation suffices to cease having $a+b=c+d$; but one may also profit from this circumstance to facilitate the transformation. Indeed, the two factors of the quantity under the radical being then of the form $\lambda+\nu(x^{2}+2mx)$, $\zeta+\theta(x^{2}+2mx)$, it is clear that if we put $x+m=y$, the odd powers of the variable will disappear.

\section*{Reduction of the differential $\dfrac{P\,dx}{R}$ to the form $\dfrac{Q\,d\varphi}{\sqrt{(1-c^{2}\sin.^{2}\varphi)}}$.}

\subsection*{(6).}

Since by a first preparation we can arrange that there be no odd powers of the variable under the radical, we shall henceforth put $R=\sqrt{(\alpha+\beta x^{2}+\gamma x^{4})}$. We shall also suppose that $P$ is an even function of $x$; for whatever the rational function $P$, we can always put $P=M+Nx$, $M$ and $N$ being even functions of $x$: now the part $\dfrac{Nx\,dx}{R}$ is brought back to the ordinary rules, by putting $x^{2}=y$; thus the whole difficulty reduces to integrating the formula $\dfrac{M\,dx}{R}$, in which $M$ is an even function of $x$.

\origpage{10}

This being laid down, we shall prove, by the enumeration of the different cases, that the differential $\dfrac{dx}{R}$ can always be brought back to the form $\dfrac{m\,d\varphi}{\sqrt{(1-c^{2}\sin^{2}.\varphi)}}$, where $c$ is smaller than unity.

\emph{First case.} Suppose the factors of $\alpha+\beta x^{2}+\gamma x^{4}$ imaginary, and represent this quantity by $\lambda^{2}+2\lambda\mu x^{2}\cos\theta+\mu^{2}x^{4}$, $\lambda$ and $\mu$ being positive, and $\cos\theta$ being able to be positive or negative. We shall put $x=\sqrt{\left(\dfrac{\lambda}{\mu}\right)}\cdot\operatorname{tang}\tfrac{1}{2}\varphi$, and taking $c=\sin\tfrac{1}{2}\theta$, we shall have the transformed expression
\[ \frac{dx}{R}=\frac{1}{2\sqrt{\lambda\mu}}\cdot\frac{d\varphi}{\sqrt{(1-c^{2}\sin^{2}\varphi)}}. \]

\emph{Second case.} Let $\alpha+\beta x^{2}+\gamma x^{4}=m^{2}(1+p^{2}x^{2})(1-q^{2}x^{2})$; the limit of $x$ being $\dfrac{1}{q}$, we shall put $x=\dfrac{1}{q}\cos\varphi$ and $\dfrac{p^{2}}{p^{2}+q^{2}}=c^{2}$, which will give
\[ \frac{dx}{R}=\frac{c}{mp}\cdot\frac{-d\varphi}{\sqrt{(1-c^{2}\sin^{2}\varphi)}}. \]
If we wished the transformed expression to be positive, we should have to put $\cot\varphi=\sqrt{(1-c^{2})}\operatorname{tang}\psi$, which directly gives
\[ x^{2}=\frac{\sin^{2}\psi}{q^{2}+p^{2}\cos^{2}\psi}, \]
and the transformed expression would be
\[ \frac{dx}{R}=\frac{c}{mp}\cdot\frac{d\psi}{\sqrt{(1-c^{2}\sin^{2}\psi)}}. \]

\emph{Third case.} Let $\alpha+\beta x^{2}+\gamma x^{4}=m^{2}(1+p^{2}x^{2})(x^{2}-q^{2})$; we shall put $x=\dfrac{q}{\cos\varphi}$, $\dfrac{1}{1+p^{2}q^{2}}=c^{2}$, and we shall have
\[ \frac{dx}{R}=\frac{c}{m}\cdot\frac{d\varphi}{\sqrt{(1-c^{2}\sin^{2}\varphi)}}. \]

\emph{Fourth case.} Let $\alpha+\beta x^{2}+\gamma x^{4}=m^{2}(1+p^{2}x^{2})(1+q^{2}x^{2})$; we shall suppose $p>q$, and putting $x=\dfrac{\operatorname{tang}\varphi}{p}$, $\dfrac{p^{2}-q^{2}}{p^{2}}=c^{2}$, we shall have
\[ \frac{dx}{R}=\frac{1}{mp}\cdot\frac{d\varphi}{\sqrt{(1-c^{2}\sin^{2}\varphi)}}. \]

\emph{Fifth case.} Let $\alpha+\beta x^{2}+\gamma x^{4}=m^{2}(1-p^{2}x^{2})(1-q^{2}x^{2})$,

\origpage{11}
we shall suppose $p>q$, and putting $px=\sin\varphi$, $\dfrac{q}{p}=c$, the transformed expression will be
\[ \frac{dx}{R}=\frac{1}{mp}\cdot\frac{d\varphi}{\sqrt{(1-c^{2}\sin^{2}\varphi)}}. \]

This formula will serve from $x=0$ to $x=\dfrac{1}{p}$; the radical $R$ would be imaginary from $x=\dfrac{1}{p}$ to $x=\dfrac{1}{q}$; but it becomes real again from $x=\dfrac{1}{q}$ to $x=\infty$. In this last case, we must write $R^{2}=m^{2}(p^{2}x^{2}-1)(q^{2}x^{2}-1)$, and putting $qx=\dfrac{1}{\sin\varphi}$, $\dfrac{q}{p}=c$, the transformed expression will be $\dfrac{-d\varphi}{mp\sqrt{(1-c^{2}\sin^{2}\varphi)}}$. If we wish it to be positive, we shall have to put, as above, $\cot\varphi=\sqrt{(1-c^{2})}\cdot\operatorname{tang}\psi$, which will directly give $x^{2}=\dfrac{1-c^{2}\sin^{2}\psi}{q^{2}\cos^{2}\psi}$, and the transformed expression will be
\[ \frac{dx}{R}=\frac{1}{mp}\cdot\frac{d\psi}{\sqrt{(1-c^{2}\sin^{2}\psi)}}, \]
a formula absolutely similar to the first; whence it follows that the integral of $\dfrac{dx}{R}$ for the case of $x>\dfrac{1}{q}$ will always be deduced from the same integral, taken on the supposition $x<\dfrac{1}{p}$.

\emph{Sixth case.} Finally, let $\alpha+\beta x^{2}+\gamma x^{4}=m^{2}(x^{2}-q^{2})(p^{2}-x^{2})$; then $x$ must lie between $p$ and $q$. Let $p>q$, and let $x=\dfrac{q}{\cos\psi}$ be put; we shall have first $\dfrac{dx}{R}=\dfrac{d\psi}{m\sqrt{(p^{2}-q^{2}-p^{2}\sin^{2}\psi)}}$. In this formula, the angle $\psi$ has a limit; to introduce an unlimited one, let $\sin\psi=c\sin\varphi$ and $c^{2}=\dfrac{p^{2}-q^{2}}{p^{2}}$, and we shall have the transformed expression
\[ \frac{dx}{R}=\frac{1}{mp}\cdot\frac{d\varphi}{\sqrt{(1-c^{2}\sin^{2}\varphi)}}, \]
which we should have reached directly by putting
\[ x^{2}=\frac{q^{2}}{1-c^{2}\sin^{2}\varphi}. \]

\subsection*{(7).}

It may now be remarked that, setting aside the first case, the values of $x^{2}$ which effect the sought reduction are always

\origpage{12}
of the form $\dfrac{A+B\sin^{2}\varphi}{C+D\sin^{2}\varphi}$, where the coefficients are constant. It remains only to substitute this value of $x^{2}$ in $P$, and the formula $\displaystyle\int\frac{P\,dx}{R}$ will be transformed into another, $\displaystyle\int\frac{Q\,d\varphi}{\sqrt{(1-c^{2}\sin^{2}\varphi)}}$, in which $c$ will be smaller than unity, and $Q$ will be an even rational function of $\sin\varphi$, which will contain $\sin\varphi$ to the same degree as $P$ contains $x$.

We set aside the first case, because the form of the value of $x$ is a little different, and because moreover, by following the method of article 5, this case is avoided, since one always falls upon real factors. But we shall return in another article to the case of imaginary factors.

\section*{Development of the formula $\displaystyle\int\frac{Q\,d\varphi}{\sqrt{(1-c^{2}\sin^{2}\varphi)}}$.}

\subsection*{(8).}

Henceforth we shall represent the radical $\sqrt{(1-c^{2}\sin^{2}\varphi)}$ by $\Delta$, and also by $\Delta(\varphi)$ and $\Delta(c,\varphi)$, when it is regarded as a function of $\varphi$, or as a function of $c$ and $\varphi$.

Let us consider first the case where $Q$ is an integral function, and let $Q=A+B\sin^{2}\varphi+C\sin^{4}\varphi+\text{etc.}$; if we represent for brevity the integral $\displaystyle\int\frac{\sin^{2k}\varphi\cdot d\varphi}{\Delta}$ by $Z^{2k}$, we shall have
\[ \int\frac{Q\,dx}{\Delta}=AZ^{0}+BZ^{2}+CZ^{4}+\text{etc.} \]
\ednote{Printed $Q\,dx$ in the numerator; the integrand is in $\varphi$, so $Q\,d\varphi$ is expected. An apparent misprint.}
But by differentiating the quantity $\Delta\cos\varphi\sin^{2k-3}\varphi$, we easily find the formula
\[ \Delta\cos\varphi\sin^{2k-3}\varphi=(2k-3)Z^{2k-4}-(1+c^{2})(2k-2)Z^{2k-2}+c^{2}(2k-1)Z^{2k}; \]
whence it follows that $Z^{4}$, $Z^{6}$, etc.\ can be expressed by means of $Z^{0}$ and $Z^{2}$, and that thus, in the case where $Q$ is an integral function, the integral $\displaystyle\int\frac{Q\,d\varphi}{\Delta}$ is equal to an algebraic quantity, plus an integral of the form $\displaystyle\int(A+B\sin^{2}\varphi)\frac{d\varphi}{\Delta}$.

\subsection*{(9).}

Let $Q$ now be any rational function of $\sin^{2}\varphi$; after extracting the integral part, and treating it as

\origpage{13}
has just been said, the development of the fractional part can always be made in such a way that each partial fraction is of the form $\dfrac{N}{(1+n\sin^{2}\varphi)^{k}}$, $n$ and $N$ being constant coefficients, real or imaginary. It will therefore be a matter of reducing the formula $\displaystyle\int\frac{d\varphi}{(1+n\sin^{2}\varphi)^{k}\Delta}$, and all similar ones, to the simplest formulas of the same species. Now, if we put for a moment $\sin\varphi=x$, this formula will become
\[ \int\frac{dx}{(1+nx^{2})^{k}\sqrt{[1-(1+c^{2})x^{2}+c^{2}x^{4}]}}. \]
Let us consider, for greater generality, the formula
\[ \Pi^{k}=\int\frac{dx}{(1+nx^{2})^{k}R}, \]
in which $R$ represents the radical $\sqrt{(\alpha+\beta x^{2}+\gamma x^{4})}$; we shall find by the means already indicated,
\[ \begin{aligned} \frac{xR}{(1+nx^{2})^{k-1}}={}&(2k-2)\left(\alpha-\frac{\beta}{n}+\frac{\gamma}{n^{2}}\right)\Pi^{k}-(2k-3)\left(\alpha-\frac{2\beta}{n}+\frac{3\gamma}{n^{2}}\right)\Pi^{k-1}\\ &+(2k-4)\left(-\frac{\beta}{n}+\frac{3\gamma}{n^{2}}\right)\Pi^{k-2}-(2k-5)\frac{\gamma}{n^{2}}\Pi^{k-3}. \end{aligned} \]
From this it results that the term $\Pi^{k}$, and all similar ones where $k$ is greater than unity, can be expressed partly algebraically, partly by the quantities $\Pi^{1}$, $\Pi^{0}$, $\Pi^{-1}$, which are the simplest of their species. It may even be observed that if $1+nx^{2}$ were a divisor of the quantity under the radical, in which case we should have $\alpha-\dfrac{\beta}{n}+\dfrac{\gamma}{n^{2}}=0$, then the preceding formula would have one term fewer; thus the quantity $\Pi^{k}$ would no longer depend on anything but the two $\Pi^{0}$ and $\Pi^{-1}$.

Therefore the determination of $\Pi^{k}$ will in general depend on the formula
\[ \int\left(\frac{A}{1+nx^{2}}+B+Cx^{2}\right)\frac{dx}{R}, \]
and in particular, when $1+nx^{2}$ is a divisor of $R^{2}$, this denominator can be made to disappear, and the determination of $\Pi^{k}$ will depend only on the formula $\displaystyle\int(B+Cx^{2})\frac{dx}{R}$.

\origpage{14}

\subsection*{(10).}

The application of this result to the integral $\displaystyle\int\frac{d\varphi}{(1+n\sin^{2}\varphi)^{k}\Delta}$ shows that in general this integral will depend on the following,
\[ \int\left(\frac{A}{1+n\sin^{2}\varphi}+B+C\sin^{2}\varphi\right)\frac{d\varphi}{\Delta}. \]
And in the particular case where $1+n\sin^{2}\varphi$ is a factor of $\cos^{2}\varphi\cdot\Delta^{2}$, the integral can be freed of the denominator, and will depend only on the integral formula
\[ \int(B+C\sin^{2}\varphi)\frac{d\varphi}{\Delta}. \]
This particular case will occur if $n=-1$, and if $n=-c^{2}$, in which cases the formulas are $\displaystyle\int\frac{d\varphi}{\cos^{2k}\varphi\cdot\Delta}$, and $\displaystyle\int\frac{d\varphi}{\Delta^{2k+1}}$; it is indeed easy to reduce both to an algebraic part, plus an integral of the form $\displaystyle\int(A+B\sin^{2}\varphi)\frac{d\varphi}{\Delta}$. To these two cases may be added that of $\displaystyle\int\frac{d\varphi}{\sin^{2k}\varphi\cdot\Delta}$, which is susceptible of a similar reduction, and which will be effected by the formula of article 8.

It results from this analysis that the general formula $\displaystyle\int\frac{Q\,d\varphi}{\Delta}$, suitably reduced, will contain, 1°.\ an algebraic part; 2°.\ an integral of the form $\displaystyle\int(A+B\sin^{2}\varphi)\frac{d\varphi}{\Delta}$; 3° one or more parts of the form $\displaystyle\int\frac{N}{1+n\sin^{2}\varphi}\cdot\frac{d\varphi}{\Delta}$, in each of which the coefficients $N$ and $n$ will have any values whatever, real or imaginary.

\section*{Definition of the elliptic functions, and their division into three species.}

\subsection*{(11).}

Since the transcendents which we have considered always reduce to these two forms $\displaystyle\int(A+B\sin^{2}\varphi)\frac{d\varphi}{\Delta}$, $\displaystyle\int\frac{d\varphi}{(1+n\sin^{2}\varphi)\Delta}$, it is clear that they are comprised in the general formula
\[ H=\int\frac{A+B\sin^{2}\varphi}{1+n\sin^{2}\varphi}\cdot\frac{d\varphi}{\Delta}. \]

\origpage{15}
We shall henceforth call \emph{elliptic functions} or \emph{transcendents} the integrals comprised in this formula. The transcendent $H$ will be supposed to vanish or begin when $\varphi=0$; its extent will be determined by the variable $\varphi$, which will be called \emph{the amplitude}; the constant $c$, always smaller than unity, will be called \emph{the modulus}; finally $\sqrt{(1-c^{2})}$, which we shall constantly denote by $b$, may be called \emph{the complement of the modulus}.

The quantity $H$, when one considers only the variability of $\varphi$, can be denoted by $H(\varphi)$; if one considers at once the variation of the modulus and of the amplitude, it can be represented by $H(c,\varphi)$, and so on, if one wished to take account of the inequality of the coefficients.

\subsection*{(12).}

It suffices to know all the values of $H$, from $\varphi=0$ to $\varphi=90^{\circ}=\dfrac{\pi}{2}$, and the value of this transcendent for any amplitude whatever will be easily known. Indeed, let $\varphi=i\pi\pm\alpha$, $i$ being an integer, and $\alpha$ an arc smaller than $\dfrac{\pi}{2}$, then we shall have
\[ H(\varphi)=2iH\left(\frac{\pi}{2}\right)\pm H(\alpha). \]
In this respect the function $H$ is like the arcs of an ellipse, and in general the arcs of all oval curves composed of four equal and similar parts: an arc, however large it may be, and containing, if one likes, several circumferences, is always expressed without difficulty by the quarter of the curve and a portion of that quarter. The determinate function $H\left(\dfrac{\pi}{2}\right)$ is in a way the unit of the functions $H$, or the function become complete: we shall denote it by $H^{1}$.

It does not appear that the function $H$, taken in all its generality, can be reduced to arcs of an ellipse; this holds only when $n=0$, or when some substitution can make the denominator $1+n\sin^{2}\varphi$ disappear, which we do not believe possible in general. Thus the denomination of elliptic function is improper in some respects; we nevertheless adopt it because of the great analogy that will be found between the properties of this function and those of arcs of an ellipse.

\origpage{16}

\subsection*{(13).}

Since every integral denoted by $\displaystyle\int\frac{P\,dx}{R}$ reduces to an algebraic part, plus a certain number of terms each of which can be assimilated to the function $H$, it follows that for the integrals in question one might admit only a single species of transcendents, represented by the function $H$, and in which the coefficients $A$, $B$, $n$ would be at will real or imaginary. But in order to penetrate the nature of these integrals thoroughly, and to be able to establish between them the comparisons and reductions of which they are capable, it is necessary to divide the function $H$ into several distinct species, whose properties will become more apparent when each is considered in isolation.

I observe first that arcs of an ellipse are contained in the formula $H$. Indeed, the equation of the ellipse being $y^{2}=\dfrac{b^{2}}{a^{2}}(a^{2}-x^{2})$, if we put $x=a\sin\varphi$, we shall have
\[ y=b\cos\varphi,\quad\text{and}\quad\sqrt{(dx^{2}+dy^{2})}=d\varphi\sqrt{(a^{2}\cos^{2}\varphi+b^{2}\sin^{2}\varphi)}, \]
a formula which reduces to $d\varphi\sqrt{(1-c^{2}\sin^{2}\varphi)}$, on putting $a=1$ and $c^{2}=1-b^{2}$.

\rmfigure{figures/fig-1.png}{Fig.~1.}{A quadrant of the ellipse BMKNA with centre C and semi-axes CA and CB, inside the dotted circular quadrant DM'K'N'A, with the ordinates PMM', IKK' and QNN' standing on CA, the lines CO and CZ from the centre, the segment OM and the segment ZN.}

Let then the semi-major axis be $CA=1$ and the semi-minor axis $CB=b$; if on the circumscribed circle $DM'A$ we take the arc $DM'=\varphi$, and let fall from the point $M'$ the perpendicular $M'P$ on the major axis, we shall determine on the ellipse an arc $BM=E$, whose value will be
\[ E=\int\Delta\,d\varphi. \]
This function or transcendent $E$ constitutes one of the species into which we shall divide the function $H$; it is deduced from the general formula by putting $n=0$, $A=1$, $B=-c^{2}$.

\rmfigure{figures/fig-2.png}{Fig.~2.}{A hyperbola branch rising from A through M, N and K towards O, with the axis CAPF, the perpendicular CB, the straight line CV from C, the ordinate PM, the tangent at M extended down to Z, and the segment CZ.}

The equation of the hyperbola being $y^{2}=\dfrac{\beta^{2}}{\alpha^{2}}(x^{2}-\alpha^{2})$, if we likewise put $x=\dfrac{\alpha}{\cos\varphi}$, we shall have
\[ y=\frac{\beta\sin\varphi}{\cos\varphi},\quad\text{and}\quad\sqrt{(dx^{2}+dy^{2})}=\frac{d\varphi}{\cos^{2}\varphi}\sqrt{(\beta^{2}+\alpha^{2}\sin^{2}\varphi)}; \]
but in order to have a radical entirely similar to that of the arc of an ellipse, other denominations must be taken. Let then the transverse semi-axis be $CA=c$, its conjugate $CB=b$, the distance from the

\origpage{17}
centre to the focus $CF=1$, and the ordinate $y=b^{2}\operatorname{tang}\varphi$; we shall have the abscissa $x=\dfrac{c}{\cos\varphi}\sqrt{(1-c^{2}\sin^{2}\varphi)}$, whence results $\sqrt{(dx^{2}+dy^{2})}=\displaystyle\int\frac{b^{2}\,d\varphi}{\Delta\cos^{2}\varphi}$. But on differentiating the quantity $\Delta\operatorname{tang}\varphi$, we have
\[ d(\Delta\operatorname{tang}\varphi)=\frac{b^{2}\,d\varphi}{\Delta\cos^{2}\varphi}-\frac{b^{2}\,d\varphi}{\Delta}+\Delta\,d\varphi. \]
Hence if we call $\Upsilon$ the arc of a hyperbola $AM$ which corresponds to the amplitude $\varphi$, or whose extreme ordinate is $PM=b^{2}\operatorname{tang}\varphi$, we shall have
\[ \Upsilon=\Delta\operatorname{tang}\varphi-\int\Delta\,d\varphi+b^{2}\int\frac{d\varphi}{\Delta}: \]
and one may remark that the algebraic part $\Delta\operatorname{tang}\varphi$ is nothing other than the tangent $MZ$ terminated by the perpendicular $CZ$ let fall from the centre onto this tangent.

The function $\Upsilon$ is comprised in the formula $H$, since it is deduced from it by putting $A=b^{2}$, $B=-b^{2}c^{2}$, $n=-1$; if we take this function for the second species of the transcendents contained in the formula $H$, the integral $\displaystyle\int\frac{d\varphi}{\Delta}$ can be expressed by the arcs $E$ and $\Upsilon$ by means of the equation
\[ b^{2}\int\frac{d\varphi}{\Delta}=E+\Upsilon-\Delta\operatorname{tang}\varphi. \]
Finally, as $H$ can be put in the form
\[ H=A'\int\frac{d\varphi}{\Delta}+B'\int\Delta\,d\varphi+C'\int\frac{d\varphi}{(1+n\sin^{2}\varphi)\Delta}; \]
there remains only a third species of transcendents to be considered, represented by the formula
\[ \Pi=\int\frac{d\varphi}{(1+n\sin^{2}\varphi)\Delta}. \]

Such, then, would be the division of elliptic functions into three species, if one admitted the arcs $E$ and $\Upsilon$ as constituting the first two species, and this is indeed the first idea that presents itself in this kind of investigation.

\subsection*{(14).}

But, by a more thorough examination of the properties of these

\origpage{18}
functions, one finds that the function $F=\displaystyle\int\frac{d\varphi}{\Delta}$ ought to be preferred to the arc of a hyperbola $\Upsilon$, in making it one of the three species of elliptic functions, and that this function $\displaystyle\int\frac{d\varphi}{\Delta}$ ought even to be considered as simpler than the arc of an ellipse $\displaystyle\int\Delta\,d\varphi$.

Indeed, we shall prove below that the function $F$ can be expressed by two arcs of an ellipse; but the converse does not hold, and an arc of an ellipse cannot be expressed by two functions such as $F$, which already indicates that the function $E$ is of a more composite nature than the function $F$. But this consequence shows itself still more clearly by examining the respective properties of these functions.

The most remarkable property of the functions $F$ is that one can determine, by purely algebraic operations, a function equal to the sum or to the difference of two others; whence it follows that one can determine algebraically a function that is a multiple, a submultiple, or in general in a rational ratio to a given function; a property which the functions $F$ share with arcs of a circle and logarithms, and which holds even when these functions, considered as integrals or as arcs of a curve, do not have the common origin $\varphi=0$, and begin at any points whatever.

Arcs of an ellipse and arcs of a hyperbola are, of all the other transcendents, those which come nearest to enjoying the same property, but they do not enjoy it in an absolute manner. Thus, two arcs being given on one of these curves, reckoned from the same point where $\varphi=0$, one can find algebraically, not an arc equal to their sum, but an arc equal to this sum, plus or minus an algebraic quantity, which proves that the arcs $\Upsilon$ and $E$ are of a more composite nature than the functions $F$. These latter must therefore be regarded as holding the first rank after arcs of a circle and logarithms~${}^{(1)}$.

\textbf{(1)} These functions unite so great a number of properties that, when they shall be more generally known, it will doubtless be judged necessary to impose a particular name upon them, and to designate the function of $c$ and $\varphi$ equal to $\displaystyle\int\frac{d\varphi}{\Delta}$ as one designates the arc whose sine is $x$, or the number whose logarithm is $y$. It seems that the function $F$ would be fairly well characterized by giving it the name \emph{Nome}, because this function has the property of regulating everything that concerns the comparison of elliptic functions. Perhaps it would be fitting at the same time to give the names \emph{Epinome} and \emph{Paranome} to the functions $E$ and $\Pi$ which constitute the other two species.

\origpage{19}

Finally, a motive which alone would suffice to make the functions $F$ preferred to the functions $\Upsilon$ in the classification of elliptic functions is that one cannot suppose $\varphi>\tfrac{1}{2}\pi$ in the function $\Upsilon$, whereas one can suppose $\varphi$ of any magnitude whatever in the functions $F$ and $E$, which, in general, increase almost proportionally to the angle $\varphi$. Now the applications of the integral calculus, principally those which concern mechanics, often give occasion to attribute to the principal variable $\varphi$ any values whatever, composed, if one wishes, of several circumferences. The use of the functions $\Upsilon$ would therefore give rise to embarrassments or even to errors, which is never to be feared in that of the functions $F$.

\subsection*{(15).}

This being laid down, the elliptic functions or transcendents comprised in the formula $H$ will be divided into three species:

The first and the simplest is represented by the formula $F=\displaystyle\int\frac{d\varphi}{\Delta}$;

The second is the arc of an ellipse, reckoned from the minor axis, and whose expression is $E=\displaystyle\int\Delta\,d\varphi$;

Finally, the third species is represented by the formula $\Pi=\displaystyle\int\frac{d\varphi}{(1+n\sin^{2}\varphi)\Delta}$; it contains, besides the modulus $c$ common to the other two species, a parameter $n$ which may be at will positive or negative, real or imaginary.

One might believe that the case where $n$ is imaginary differs essentially from the case where $n$ is real, and that it requires the formation of a fourth species of elliptic functions; but by reductions and transformations which we shall set out below, it will be seen that this fourth species is needless to consider, and that the third species may thus be restricted to the cases where $n$ is real alone. We shall regard only as necessary and common conditions

\origpage{20}
to the three species of functions, that the amplitude $\varphi$ and the modulus $c$ be real, and that at the same time $c$ be smaller than unity.

\section*{Comparison of elliptic functions of the first species.}

\subsection*{(16).}

All geometers know the complete algebraic integral which Euler has given of the differential equation
\[ \frac{dx}{\sqrt{(\alpha+\beta x+\gamma x^{2}+\delta x^{3}+\varepsilon x^{4})}}+\frac{dy}{\sqrt{(\alpha+\beta y+\gamma y^{2}+\delta y^{3}+\varepsilon y^{4})}}=0. \]
According to the reductions indicated in article 6, this equation can, without losing any of its generality, be put in the form
\[ \frac{d\varphi}{\sqrt{(1-c^{2}\sin^{2}\varphi)}}+\frac{d\psi}{\sqrt{(1-c^{2}\sin^{2}\psi)}}=0; \]
and then its integral is $F(\varphi)+F(\psi)=F(\mu)$, $\mu$ being an arbitrary constant. But the same integral, found by Euler's method, is expressed thus:
\[ \cos\varphi\cos\psi-\sin\varphi\sin\psi\sqrt{(1-c^{2}\sin^{2}\mu)}=\cos\mu\ldots(a'), \]
and here is how this result may be verified \emph{a posteriori}.

I observe first that the equation $(a')$ can be put in one or the other of the two following forms:
\[ \left.\begin{aligned} \cos\mu\cos\varphi+\sin\mu\sin\varphi\,\Delta(\psi)&=\cos\psi\\ \cos\mu\cos\psi+\sin\mu\sin\psi\,\Delta(\varphi)&=\cos\varphi \end{aligned}\right\}\ldots(b'); \]
for these equations, each freed from the radical it contains, lead to the same result as the equation $(a')$ freed from its radical.

This being so, if we differentiate the equation $(a')$ after having divided each member by $\sin\varphi\sin\psi$, in order to make the radical disappear, we shall have
\[ \frac{d\varphi}{\sin\varphi}(\cos\psi-\cos\mu\cos\varphi)+\frac{d\psi}{\sin\psi}(\cos\varphi-\cos\mu\cos\psi)=0. \]
Substituting in this the values of $\cos\psi-\cos\mu\cos\varphi$ and $\cos\varphi-\cos\mu\cos\psi$ given by the equations $(b')$, we shall have
\[ \frac{d\varphi}{\Delta(\varphi)}+\frac{d\psi}{\Delta(\psi)}=0. \]

\origpage{21}
From this one sees that the transcendental equation $F(\varphi)+F(\psi)=F(\mu)$ is satisfied, whether by the algebraic equation $(a')$ or by one of the equations $(b')$; for these three equations may be employed indifferently one for another.

If we had $c=0$, the function $F(\varphi)$ would reduce to the arc $\varphi$, and then, having $\varphi+\psi=\mu$, we should conclude
\[ \begin{gathered} \cos\varphi\cos\psi-\sin\varphi\sin\psi=\cos\mu\\ \cos\mu\cos\varphi+\sin\mu\sin\varphi=\cos\psi\\ \cos\mu\cos\psi+\sin\mu\sin\psi=\cos\varphi. \end{gathered} \]
This is indeed what the equations $(a')$ and $(b')$ give in the case $c=0$, which reduces the radicals $\Delta(\mu)$, $\Delta(\varphi)$, $\Delta(\psi)$ to unity.

\subsection*{(17).}

\rmfigure{figures/fig-3.png}{Fig.~3.}{A spherical triangle ABC with the side AC prolonged beyond C and a second arc drawn from the point D on BC up to meet that prolongation, together with a dashed arc from a to beta beside AB and short dotted marks lettered alpha, a at A and beta, b at B.}

Here it is proper to observe, following Lagrange~${}^{(1)}$, that if we construct a spherical triangle whose sides $AB$, $AC$, $BC$ are respectively equal to the amplitudes $\mu$, $\varphi$, $\psi$, then the opposite angles $C$, $B$, $A$ will be such that we shall have
\[ \begin{aligned} \cos C&=\frac{\cos\mu-\cos\varphi\cos\psi}{\sin\varphi\sin\psi}=-\sqrt{(1-c^{2}\sin^{2}\mu)}\\ \cos B&=\frac{\cos\varphi-\cos\mu\cos\psi}{\sin\mu\sin\psi}=\sqrt{(1-c^{2}\sin^{2}\varphi)}\\ \cos A&=\frac{\cos\psi-\cos\mu\cos\varphi}{\sin\mu\sin\varphi}=\sqrt{(1-c^{2}\sin^{2}\psi)}; \end{aligned} \]
whence we deduce
\[ \begin{gathered} \sin^{2}C=c^{2}\sin^{2}\mu,\quad\sin^{2}B=c^{2}\sin^{2}\varphi,\quad\sin^{2}A=c^{2}\sin^{2}\psi,\\ \text{or}\quad c=\frac{\sin C}{\sin\mu}=\frac{\sin B}{\sin\varphi}=\frac{\sin A}{\sin\psi}. \end{gathered} \]
These equations agree with the known properties of spherical triangles; they prove at the same time that in every spherical triangle formed by three sides $\varphi$, $\psi$, $\mu$ which satisfy the transcendental equation $F(\varphi)+F(\psi)=F(\mu)$, the ratio of the sine of each

\textbf{(1)} \emph{Théorie des Fonct.\ analyt.}, pag.\ 85.

\origpage{22}
angle to the sine of the opposite side will be equal to the modulus~${}^{(1)}$. One sees moreover that in the triangle $ABC$ the two angles $A$ and $B$ are always acute, and the angle $C$ always obtuse.

If we then put $F(\mu)+F(\theta)=F(\nu)$, it will be necessary to observe that in the new spherical triangle which will have for sides $\mu$, $\theta$, $\nu$, the angle opposite the side $\mu$ must be acute and have the same sine as the angle opposite the side $\mu$ in the first triangle, so that its ratio to $\sin\mu$ be equal to $c$ in both triangles; thus the angle opposite the side $\mu$ in the second triangle must be the supplement of the angle opposite the side $\mu$ in the first; whence arises this construction.

Prolong the side $AC$ towards $E$, make $CD=\theta$; from the point $D$ and with an interval $DE=\mu$, describe an arc which will meet $CE$ in $E$, and will determine the second triangle $CDE$, in which we shall have $CE=\nu$, and the side $\nu$ will satisfy the equation $F(\nu)=F(\varphi)+F(\psi)+F(\theta)$.

This construction can be continued as far as the limits of the sides of spherical triangles permit, that is, so long as the greatest side is not greater than a semicircumference. And it is evident that the same construction could be used to find successively the angles $\varphi_{2}$, $\varphi_{3}$, $\varphi_{4}$, etc., such that we should have $F(\varphi_{2})=2F(\varphi)$, $F(\varphi_{3})=3F(\varphi)$, etc., which would serve for the multiplication of the function $F$. But, as we repeat, these constructions can extend only to the limits of spherical triangles, whereas analysis knows no bound.

It is worth remarking that the consideration of the spherical triangle $ABC$ offers a very simple means of verifying the integral $(a')$. Indeed, if, regarding $\mu$ and $C$ as constants, we vary the sides $\varphi$ and $\psi$ infinitely little, we shall have $A\alpha=\beta b$, or $-d\varphi\cos A=d\psi\cos B$; but we have found $\cos A=\Delta(\psi)$ and $\cos B=\Delta(\varphi)$; therefore $\dfrac{d\varphi}{\Delta(\varphi)}+\dfrac{d\psi}{\Delta(\psi)}=0$.

Thus a very simple application of spherical trigonometry

\textbf{(1)} Another spherical triangle formed by the three angles $\varphi$, $\psi$, $180^{\circ}-\mu$, and by the three sides $B$, $A$, $180^{\circ}-C$, would equally satisfy the equation $F(\varphi)+F(\psi)=F(\mu)$. But the ratio which was $c$ in the first triangle would become $\dfrac{1}{c}$ in this one.

\origpage{23}
would have sufficed to find the complete algebraic integral of the transcendental equation $\dfrac{d\varphi}{\Delta(\varphi)}+\dfrac{d\psi}{\Delta(\psi)}=0$.

\subsection*{(18).}

Two elliptic functions of the first species $F(\varphi)$, $F(\psi)$ being given, if we wish to find a third function $F(\mu)$ equal to their sum, we must determine $\mu$ by the equation $(a')$, which will give the formulas
\[ \begin{aligned} \sin\mu &= \frac{\sin\varphi\cos\psi\,\Delta(\psi)+\sin\psi\cos\varphi\,\Delta(\varphi)}{1-c^{2}\sin^{2}\varphi\sin^{2}\psi} \\ \cos\mu &= \frac{\cos\varphi\cos\psi-\sin\varphi\sin\psi\,\Delta(\varphi)\,\Delta(\psi)}{1-c^{2}\sin^{2}\varphi\sin^{2}\psi} \\ \Delta(\mu) &= \frac{\Delta(\varphi)\,\Delta(\psi)-c^{2}\sin\varphi\sin\psi\cos\varphi\cos\psi}{1-c^{2}\sin^{2}\varphi\sin^{2}\psi} \\ \operatorname{tang}\mu &= \frac{\operatorname{tang}\varphi\,\Delta(\psi)+\operatorname{tang}\psi\,\Delta(\varphi)}{1-\operatorname{tang}\varphi\operatorname{tang}\psi\,\Delta(\varphi)\,\Delta(\psi)}. \end{aligned} \]
According to this last formula, if we took two auxiliary angles $\varphi'$, $\psi'$, such that
\[ \operatorname{tang}\varphi'=\operatorname{tang}\varphi\,\Delta(\psi) \quad\text{and}\quad \operatorname{tang}\psi'=\operatorname{tang}\psi\,\Delta(\varphi), \]
it would follow that
\[ \mu=\varphi'+\psi', \]
which is a means of calculating the angle $\mu$ easily by the tables of sines.

If we put $\mu=\frac{1}{2}\pi$, that is, if we have $F(\varphi)+F(\psi)=F^{1}$, the two functions $F(\varphi)$, $F(\psi)$ will be in some way complements of one another, since their sum is equal to the complete function $F^{1}$; then we shall have immediately by the equation $(a')$,
\[ b\operatorname{tang}\varphi\operatorname{tang}\psi=1: \]
this is the necessary relation between the angles $\varphi$ and $\psi$, in order that we may have $F(\varphi)+F(\psi)=F(\frac{1}{2}\pi)=F^{1}$. From this we deduce for the expression of $\psi$ in terms of $\varphi$,
\[ \sin\psi=\frac{\cos\varphi}{\Delta(\varphi)}, \quad \cos\psi=\frac{b\sin\varphi}{\Delta(\varphi)}, \quad \Delta(\psi)=\frac{b}{\Delta(\varphi)}. \]

\subsection*{(19).}

Two functions $F(\varphi)$, $F(\psi)$ being given, if we wish to know

\origpage{24}
a third function $F(\mu)$ which shall be equal to their difference, so that we have
\[ F(\varphi)-F(\psi)=F(\mu), \]
the solution of this problem will be easily deduced from that of the preceding problem, and we shall have as result the formulas
\[ \begin{aligned} \sin\mu &= \frac{\sin\varphi\cos\psi\,\Delta(\psi)-\sin\psi\cos\varphi\,\Delta(\varphi)}{1-c^{2}\sin^{2}\varphi\sin^{2}\psi} \\ \cos\mu &= \frac{\cos\varphi\cos\psi+\sin\varphi\sin\psi\,\Delta(\varphi)\,\Delta(\psi)}{1-c^{2}\sin^{2}\varphi\sin^{2}\psi} \\ \Delta(\mu) &= \frac{\Delta(\varphi)\,\Delta(\psi)+c^{2}\sin\varphi\sin\psi\cos\varphi\cos\psi}{1-c^{2}\sin^{2}\varphi\sin^{2}\psi} \\ \operatorname{tang}\mu &= \frac{\operatorname{tang}\varphi\,\Delta(\psi)-\operatorname{tang}\psi\,\Delta(\varphi)}{1+\operatorname{tang}\varphi\operatorname{tang}\psi\,\Delta(\varphi)\,\Delta(\psi)}. \end{aligned} \]
If in this last formula we put $\operatorname{tang}\varphi'=\operatorname{tang}\varphi\,\Delta(\psi)$ and $\operatorname{tang}\psi'=\operatorname{tang}\psi\,\Delta(\varphi)$, we shall have $\mu=\varphi'-\psi'$.

These formulas for the difference could be deduced from the formulas for the sum, by simply changing in the latter the sign of $\psi$, and keeping $\Delta(\psi)$ positive. It is needless, moreover, to point out the analogy which prevails between the values of $\sin\mu$, $\cos\mu$, and those of $\sin(\varphi\pm\psi)$, $\cos(\varphi\pm\psi)$; they would coincide entirely if we had $c=0$.

\subsection*{(20).}

Since we know algebraically the amplitude of the function equal to the sum or to the difference of two given functions, it is clear that we can find algebraically a multiple of a given function, and that in general we can resolve, on the multiplication and division of elliptic functions of the first species, the same problems that are resolved on the multiplication and division of arcs of a circle. Let us denote by $\varphi_{n}$ the amplitude of the function which contains $n$ times the function whose amplitude is $\varphi$, so that we have $F(\varphi_{n})=nF(\varphi)$; the question would be to have the general expression of $\sin$ and $\cos\varphi_{n}$, by means of $\sin$ and $\cos\varphi$.

The extreme cases present no difficulty. If we have $c=0$, then $\varphi_{n}=n\varphi$, and the expression of $\sin\varphi_{n}$, as well as that of $\cos\varphi_{n}$, are deduced from the known formula
\[ \cos\varphi_{n}+\sqrt{-1}\cdot\sin\varphi_{n}=(\cos\varphi+\sqrt{-1}\sin\varphi)^{n}. \]

\origpage{25}
If we have $c=1$, then the function $F(\varphi)$ becomes $\displaystyle\int\frac{d\varphi}{\cos\varphi}$, so that we have $F(\varphi)=\frac{1}{2}\log\left(\dfrac{1+\sin\varphi}{1-\sin\varphi}\right)$, and the formula for the multiplication of functions is
\[ \frac{1+\sin\varphi_{n}}{1-\sin\varphi_{n}}=\left(\frac{1+\sin\varphi}{1-\sin\varphi}\right)^{n}. \]
The difficulty is to find the general expression of $\sin\varphi_{n}$, when $c$ has any value whatever between $0$ and $1$.

\subsection*{(21).}

Let us consider first the simplest case, which is that of duplication; it will suffice to put $\psi=\varphi$ in the formulas of article 18, which will give
\[ \begin{aligned} \sin\varphi_{2} &= \frac{2\sin\varphi\cos\varphi\,\Delta(\varphi)}{1-c^{2}\sin^{4}\varphi} \\ \cos\varphi_{2} &= \frac{1-2\sin^{2}\varphi+c^{2}\sin^{4}\varphi}{1-c^{2}\sin^{4}\varphi} \\ \Delta(\varphi_{2}) &= \frac{1-2c^{2}\sin^{2}\varphi+c^{2}\sin^{4}\varphi}{1-c^{2}\sin^{4}\varphi} \\ \operatorname{tang}\tfrac{1}{2}\varphi_{2} &= \operatorname{tang}\varphi\,\Delta(\varphi). \end{aligned} \]
The last of these formulas would be deduced immediately from the consideration of the spherical triangle $ABC$ which, being isosceles in the case in question, gives $\operatorname{tang}\frac{1}{2}BA=\cos B\operatorname{tang}BC=\operatorname{tang}\varphi\,\Delta(\varphi)$.

Thus, taking the auxiliary $B$ such that $\sin B=c\sin\varphi$, we shall have $\varphi_{2}$ by the equation $\operatorname{tang}\frac{1}{2}\varphi_{2}=\cos B\operatorname{tang}\varphi$; we shall find in like manner $\varphi_{4}$ by means of $\varphi_{2}$, $\varphi_{8}$ by means of $\varphi_{4}$, etc., so that the function $F$ will be multiplied by 2, 4, 8, 16, etc.

Conversely, $\varphi_{2}$ being given, we shall find $\varphi$ by the equation
\[ \sin\varphi=\frac{\sin\frac{1}{2}\varphi_{2}}{\sqrt{[\frac{1}{2}+\frac{1}{2}\Delta(\varphi_{2})]}}, \]
or else, calling by analogy $\varphi_{\frac{1}{2}}$ the amplitude of the function which is equal to the half of $F(\varphi)$, we shall have
\[ \sin\varphi_{\frac{1}{2}}=\frac{\sin\frac{1}{2}\varphi}{\sqrt{(\frac{1}{2}+\frac{1}{2}\Delta)}}. \]
If we put $c\sin\varphi=\sin\beta$, which gives $\Delta=\cos\beta$, we shall have more

\origpage{26}
simply
\[ \sin\varphi_{\frac{1}{2}}=\frac{\sin\frac{1}{2}\varphi}{\cot\frac{1}{2}\beta}, \]
\ednote{Printed $\cot\frac{1}{2}\beta$ in the denominator; since $\sqrt{(\frac{1}{2}+\frac{1}{2}\cos\beta)}=\cos\frac{1}{2}\beta$, the expected reading is $\cos\frac{1}{2}\beta$. An apparent misprint.} and we shall thus satisfy the equation $F(\varphi_{\frac{1}{2}})=\frac{1}{2}F(\varphi)$.

We shall find in like manner the amplitudes $\varphi_{\frac{1}{4}}$, $\varphi_{\frac{1}{8}}$, etc., by which the bisection of the function $F(\varphi)$ can be continued indefinitely.

\subsection*{(22).}

Let us come to multiplication by any number. For this, let us consider three consecutive amplitudes $\varphi_{n-1}$, $\varphi_{n}$, $\varphi_{n+1}$, which, as the notation indicates, correspond to the multiple functions $(n-1)F$, $nF$, $(n+1)F$. The formulas for the sum and the difference of two functions will apply to the equations $F(\varphi_{n+1})=F(\varphi_{n})+F(\varphi)$, $F(\varphi_{n-1})=F(\varphi_{n})-F(\varphi)$, and there will result
\[ \begin{aligned} \sin\varphi_{n+1}+\sin\varphi_{n-1} &= \frac{2\Delta\cos\varphi\sin\varphi_{n}}{1-c^{2}\sin^{2}\varphi\sin^{2}\varphi_{n}} \\ \cos\varphi_{n+1}+\cos\varphi_{n-1} &= \frac{2\cos\varphi\sin\varphi_{n}}{1-c^{2}\sin^{2}\varphi\sin^{2}\varphi_{n}}. \end{aligned} \]
\ednote{In the second formula the numerator is printed $2\cos\varphi\sin\varphi_{n}$; the sum of the cosines requires $2\cos\varphi\cos\varphi_{n}$, as the division on p. 27 confirms. An apparent misprint.} These formulas, in which $\Delta$ and $\varphi$ remain constantly the same while $n$ varies, appear as convenient as possible for drawing from them the successive values of $\sin\varphi_{2}$, $\sin\varphi_{3}$, $\cos\varphi_{2}$, $\cos\varphi_{3}$, etc.

Let, for brevity, $2\Delta\cos\varphi=p$, $1-c^{2}\sin^{4}\varphi=q$, $c^{2}\sin^{4}\varphi=r$; we shall find for the sequence of sines,
\[ \begin{aligned} \sin\varphi_{2} &= \frac{p}{q}\sin\varphi \\ \sin\varphi_{3} &= \frac{p^{2}-q^{2}}{q^{2}-rp^{2}}\sin\varphi \\ \sin\varphi_{4} &= \frac{(1+r)p^{2}-2q^{2}}{q^{4}-rp^{4}}\,pq\sin\varphi \\ \sin\varphi_{5} &= \frac{q^{6}-3p^{2}q^{4}+(1+2r)p^{4}q^{2}-r^{2}p^{6}}{q^{6}-3rp^{2}q^{4}+r(r+2)p^{4}q^{2}-rp^{6}}\sin\varphi \\ &\text{etc.}, \end{aligned} \]
and for that of the cosines, we have, putting moreover $s=1-2\sin^{2}\varphi+r=2\cos^{2}\varphi-q$,

\origpage{27}
\[ \begin{aligned} \cos\varphi_{2} &= \frac{1-2\sin^{2}\varphi+r}{q}=\frac{s}{q} \\ \cos\varphi_{3} &= \frac{2sq-q^{2}+p^{2}r}{q^{2}-rp^{2}}\cos\varphi \\ \cos\varphi_{4} &= \frac{q^{4}-2q^{2}p^{2}\sin^{2}\varphi+rp^{4}}{q^{4}-rp^{4}} \\ \cos\varphi_{5} &= \left(\frac{2sq(q^{2}-rp^{2})(2sq-q^{2}+rp^{2})}{q^{6}-3rp^{2}q^{4}+(r^{2}+2r)p^{4}q^{2}-rp^{6}}-1\right)\cos\varphi \\ &\text{etc.} \end{aligned} \]
But these expressions, when fully developed, become very prolix, and their law is very difficult to perceive.

When it is a question of calculating $\varphi_{n}$ trigonometrically by means of $\varphi$, one will arrive at it easily in this manner. The two general formulas, divided one by the other, will give
\[ \frac{\sin\varphi_{n+1}+\sin\varphi_{n-1}}{\cos\varphi_{n+1}+\cos\varphi_{n-1}}=\Delta\operatorname{tang}\varphi_{n}, \]
which reduces to this very simple formula,
\[ \operatorname{tang}(\tfrac{1}{2}\varphi_{n+1}+\tfrac{1}{2}\varphi_{n-1})=\Delta\operatorname{tang}\varphi_{n}; \]
whence we draw successively
\[ \begin{aligned} &\operatorname{tang}\tfrac{1}{2}\varphi_{2}=\Delta\operatorname{tang}\varphi \\ &\operatorname{tang}(\tfrac{1}{2}\varphi_{3}+\tfrac{1}{2}\varphi)=\Delta\operatorname{tang}\varphi_{2} \\ &\operatorname{tang}(\tfrac{1}{2}\varphi_{4}+\tfrac{1}{2}\varphi_{2})=\Delta\operatorname{tang}\varphi_{3} \\ &\qquad\qquad\text{etc.} \end{aligned} \]
We shall thus know successively, by a very simple calculation, the values of $\varphi_{2}$, $\varphi_{3}$, etc. One might also proceed by larger intervals in order to arrive sooner at a distant term; for we have in like manner
\[ \operatorname{tang}(\tfrac{1}{2}\varphi_{n+i}+\tfrac{1}{2}\varphi_{n-i})=\Delta i\operatorname{tang}\varphi_{n}, \]
$\Delta i$ being the $\Delta$ which corresponds to the amplitude $\varphi i$.

\subsection*{(23).}

The division of a given elliptic function of the first species into a certain number of equal parts is an

\origpage{28}
algebraic problem which will be resolved by the development of the formulas that serve for multiplication. We have already seen the formulas for dividing by 2 or by a power of 2. Suppose we wish to divide the function $F$ into three equal parts; we shall call $\varphi_{3}$ the amplitude of the given function, and $\varphi$ that of the function which is the third of it. Putting then $\sin\varphi_{3}=a$ and $\sin\varphi=x$, we shall have to determine $x$ the equation
\[ a=\frac{3x-4(1+c^{2})x^{3}+6c^{2}x^{5}-c^{4}x^{9}}{1-6c^{2}x^{4}+4c^{2}(1+c^{2})x^{6}-3c^{4}x^{8}}, \]
an equation which is of the ninth degree. It would be of the twenty-fifth degree for quintisection, and so on.

The equations are lower by half when it is a question of dividing the complete function $F^{1}$ into an odd number of parts; in this respect the division of the functions $F$ is like that of arcs of a circle; whereas the division of any arc whatever into $n$ equal parts requires the resolution of an equation of the $n^{\text{th}}$ degree, that of the quarter of the circumference requires only the resolution of an equation of degree $\dfrac{n-1}{2}$.

Suppose in general $\varphi_{n}=\frac{1}{2}\pi$, so that we have $F(\varphi)=\dfrac{1}{n}F^{1}$; we shall then have also $F(\varphi_{n-i})+F(\varphi_{i})=nF(\varphi)=F^{1}$; which gives the formula $\operatorname{tang}(\varphi_{n-i})=\dfrac{1}{b}\cot\varphi_{i}$, whence we deduce
\[ \begin{gathered} \operatorname{tang}\varphi_{n-1}=\frac{1}{b}\cot\varphi \\ \operatorname{tang}\varphi_{n-2}=\frac{1}{b}\cot\varphi_{2} \\ \operatorname{tang}\varphi_{n-3}=\frac{1}{b}\cot\varphi_{3} \\ \text{etc.} \end{gathered} \]
But because of $\varphi_{n}=\frac{1}{2}\pi$, we have
\[ \begin{gathered} \operatorname{tang}(\tfrac{1}{4}\pi+\tfrac{1}{2}\varphi_{n-2})=\Delta\operatorname{tang}\varphi_{n-1}=\frac{\Delta}{b}\cot\varphi \\ \operatorname{tang}(\tfrac{1}{2}\varphi_{n-1}+\tfrac{1}{2}\varphi_{n-3})=\Delta\operatorname{tang}\varphi_{n-2} \\ \text{etc.} \end{gathered} \]
From these result formulas fairly simple for determining $\varphi_{n-1}$, $\varphi_{n-2}$, etc.; developing those which give $\varphi_{2}$, $\varphi_{3}$, etc., we shall have

\origpage{29}
by the meeting of these two sequences the equation which is to determine $\varphi$, and the calculation will be less complicated than by the development of $\sin\varphi_{n}$ or of $\cos\varphi_{n}$.

When $n=3$, we shall have immediately $\operatorname{tang}(\frac{1}{4}\pi+\frac{1}{2}\varphi)=\dfrac{\Delta}{b}\cot\varphi$, or $b\sin\varphi=\Delta(1-\sin\varphi)$. Let $\sin\varphi=x$, and the equation for determining $x$ will be
\[ 0=1-2x+2c^{2}x^{3}-c^{2}x^{4}: \]
this is the equation which gives the trisection of the function $F^{1}$, and one sees that its degree $=\dfrac{9-1}{2}$.

When $n=5$, we shall have to eliminate $\varphi_{3}$ from the two equations
\[ \begin{gathered} \operatorname{tang}(\tfrac{1}{4}\pi+\tfrac{1}{2}\varphi_{3})=\frac{\Delta}{b}\cot\varphi \\ \operatorname{tang}(\tfrac{1}{2}\varphi_{3}+\tfrac{1}{2}\varphi)=\Delta\operatorname{tang}\varphi_{2}, \end{gathered} \]
and then put in place of $\operatorname{tang}\varphi_{2}$ its value $\dfrac{2\Delta\sin\varphi\cos\varphi}{1-2\sin^{2}\varphi+c^{2}\sin^{4}\varphi}$; we shall thus obtain, putting $\sin\varphi=x$,
\[ \frac{1+x}{bx}\sqrt{(1-c^{2}x^{2})}=\frac{1+2x-2c^{2}x^{3}-c^{2}x^{4}}{1-2x+2c^{2}x^{3}-c^{2}x^{4}}, \]
an equation for the quintisection of the function $F^{1}$, which, when fully developed, will rise to the degree $12=\dfrac{25-1}{2}$.

\subsection*{(24).}

Let us return to the equation of the trisection; it offers in its resolution some peculiarities which deserve to be remarked.

Let first $x=y+\frac{1}{2}$, and this equation will become
\[ y^{4}-\tfrac{3}{2}y^{2}+\left(\frac{2}{c^{2}}-1\right)y-\tfrac{3}{16}=0. \]
Suppose that the first member is the product of the two factors $y^{2}-py+q$, $y^{2}+py-r$; we shall have to determine $p$, $q$, $r$ the equations $q-r=p^{2}-\frac{3}{2}$, $qr=\frac{3}{16}$, $p(q+r)=\dfrac{2}{c^{2}}-1$. From the first two we draw $(q+r)^{2}=p^{4}-3p^{2}+3$, and this value being

\origpage{30}
substituted in the third, there results $p^{6}-3p^{4}+3p^{2}=\left(\dfrac{2}{c^{2}}-1\right)^{2}$, or $(p^{2}-1)^{3}=\dfrac{4b^{2}}{c^{4}}$. Let therefore
\[ k=\sqrt[3]{\left(\frac{4b^{2}}{c^{4}}\right)}, \]
and we shall have $p^{2}-1=k$, or $p=\sqrt{(1+k)}$; whence results
\[ \begin{aligned} q &= \tfrac{1}{2}k-\tfrac{1}{4}+\tfrac{1}{2}\sqrt{(k^{2}-k+1)} \\ r &= \tfrac{1}{2}k-\tfrac{1}{4}-\tfrac{1}{2}\sqrt{(k^{2}-k+1)}. \end{aligned} \]
Thus the four values of $y$ will be
\[ \begin{aligned} y &= \tfrac{1}{2}\sqrt{(1+k)}\pm\tfrac{1}{2}\sqrt{[2-k-2\sqrt{(1-k+k^{2})}]} \\ y &= -\tfrac{1}{2}\sqrt{(1+k)}\pm\tfrac{1}{2}\sqrt{[2-k+2\sqrt{(1-k+k^{2})}]}. \end{aligned} \]
The first two are imaginary, and of the other two only the positive root suits the question, so that we shall have $y+\frac{1}{2}$, or
\[ \sin\varphi=\tfrac{1}{2}-\tfrac{1}{2}\sqrt{(1+k)}+\tfrac{1}{2}\sqrt{[2-k+2\sqrt{(1-k+k^{2})}]}; \]
so that this value can be constructed geometrically, when $k$ is given. Now, whatever $k$ may be, we have
\[ c^{2}=\frac{2}{1+\sqrt{(1+k^{3})}}, \]
a value which is always comprised between the limits 1 and 0.

Let $k=1$, we shall have
\[ \begin{gathered} c^{2}=\frac{2}{1+\sqrt{2}}=2(\sqrt{2}-1) \quad\text{and} \\ \sin\varphi=\tfrac{1}{2}-\tfrac{1}{2}\sqrt{2}+\tfrac{1}{2}\sqrt{3}. \end{gathered} \]

Let $k=2$, we shall have
\[ \begin{gathered} c^{2}=\tfrac{1}{2}=b^{2} \quad\text{and} \\ \sin\varphi=\tfrac{1}{2}-\tfrac{1}{2}\sqrt{3}+\sqrt{(\tfrac{1}{2}\sqrt{3})} \\ \cos^{2}\varphi=(-1+\sqrt{3})\sqrt{(\tfrac{1}{2}\sqrt{3})}=\sqrt{(2\sqrt{3}-3)}. \end{gathered} \]
\ednote{The sign before the 1 in $(-1+\sqrt{3})$ is faint and partly obscured by a reader's pencil stroke through the following $+$; it is read as $-1$, which the value of $\cos^{2}\varphi$ requires.} Here are two cases in which the trisection of the function $F^{1}$ is effected by simple extractions of square roots.

The applications that we shall have occasion to make in what follows,

\origpage{31}
require us to consider further the case of $c=\frac{1}{2}\sqrt{(2+\sqrt{3})}=\cos\dfrac{\pi}{12}$, and that of $c=\frac{1}{2}\sqrt{(2-\sqrt{3})}=\sin\dfrac{\pi}{12}$.

Let first $c=\frac{1}{2}\sqrt{(2+\sqrt{3})}$, we shall have $k^{3}=\dfrac{4b^{2}}{c^{4}}=\dfrac{4b^{2}c^{2}}{c^{6}}=\dfrac{1}{4c^{6}}$; hence $k=\dfrac{1}{2c^{2}}\sqrt[3]{2}$; from this value we shall find
\[ \sin\varphi=\sqrt{3}-1. \]
But as the reductions needed to arrive at this result would be very laborious, it is simpler to return to the original equation
\[ 0=1-2x+\frac{2+\sqrt{3}}{4}(2x^{3}-x^{4}), \]
and one will easily verify that this equation is satisfied by putting $x=\sqrt{3}-1$; then we have $\sin^{2}\varphi=4-2\sqrt{3}$, $\cos^{2}\varphi=2\sqrt{3}-3$, and $\operatorname{tang}^{2}\varphi=\dfrac{2}{\sqrt{3}}$, or $\operatorname{tang}\varphi=\sqrt{\left(\dfrac{2}{\sqrt{3}}\right)}$. This is the value of $\varphi$ which gives $F(\varphi)=\frac{1}{3}F^{1}$, when $c=\cos\dfrac{\pi}{12}$.

Let now $c=\frac{1}{2}\sqrt{(2-\sqrt{3})}$, we shall again have $k^{3}=\dfrac{1}{4c^{6}}$, and $k=\dfrac{1}{2c^{2}}\sqrt[3]{2}=(4+2\sqrt{3})\sqrt[3]{2}$.
The substitution of this value in the general formula will make $\sin\varphi$ known; but to arrive at the simplest result, here is the route that must be followed.

I take up again the general equation $0=1-2x+c^{2}(2x^{3}-x^{4})$; I put $x^{2}=1-y$, which amounts to supposing $y=\cos^{2}\varphi$, the transformed equation will be
\[ 0=c^{4}y^{4}+6b^{2}c^{2}y^{2}+4b^{2}(b^{2}-c^{2})y-3b^{4}. \]
Let $y=\dfrac{b}{c}z$, we shall have anew
\[ 0=z^{4}+6z^{2}+\left(\frac{4b}{c}-\frac{4c}{b}\right)z-3. \]
In the present case, we have $\dfrac{b}{c}-\dfrac{c}{b}=\dfrac{b^{2}-c^{2}}{bc}=4(b^{2}-c^{2})=2\sqrt{3}$; so that the equation becomes
\[ 0=z^{4}+6z^{2}+8\sqrt{3}\cdot z-3. \]

\origpage{32}
This equation being divided by $z+\sqrt{3}$, a factor useless for our purpose, the quotient is
\[ 0=z^{3}-z^{2}\sqrt{3}+9z-\sqrt{3}. \]
Let finally $z=\dfrac{1+2v}{\sqrt{3}}$, and the transformed equation will be
\[ 0=v^{3}+6v+2, \]
whence one draws $v=\sqrt[3]{2}-\sqrt[3]{4}$. Let $\sqrt[3]{4}=m$, and we shall have $v=\frac{1}{2}m^{2}-m$; hence $z=\dfrac{m^{2}-2m+1}{\sqrt{3}}$, $y=\dfrac{b}{c}z=(2+\sqrt{3})z=\dfrac{m^{2}-2m+1}{2\sqrt{3}-3}=\cos^{2}\varphi$; hence finally $\cos\varphi=\dfrac{m-1}{\sqrt{(2\sqrt{3}-3)}}=(m-1)\sqrt{\left(\dfrac{2+\sqrt{3}}{\sqrt{3}}\right)}$. This is the simplest expression by which one can determine $\varphi$ so as to satisfy the equation $F(\varphi)=\frac{1}{3}F^{1}$, in the case of $c=\frac{1}{2}\sqrt{(2-\sqrt{3})}=\sin\dfrac{\pi}{12}$.

To effect the trisection of the function $F^{1}$ completely, we must still have the value of $\varphi_{2}$ which gives $F(\varphi_{2})=\frac{2}{3}F^{1}$: now $\varphi_{2}$ is easy to deduce from $\varphi$, either by the formulas of duplication, or by the particular formulas
\[ \begin{gathered} \operatorname{tang}\varphi_{2}=\frac{\cot\varphi}{b} \\ \cos\varphi_{2}=1-\sin\varphi. \end{gathered} \]
The first results from the equation $b\operatorname{tang}\varphi\operatorname{tang}\psi=1$ which holds (art.~18) when $F(\varphi)+F(\psi)=F^{1}$; the second results from the combination of the equations $\operatorname{tang}\varphi_{2}=\dfrac{\cot\varphi}{b}$, $\operatorname{tang}\frac{1}{2}\varphi_{2}=\Delta(\varphi)\operatorname{tang}\varphi=\dfrac{1-\cos\varphi_{2}}{\sin\varphi_{2}}$, and $b\sin\varphi=\Delta(1-\sin\varphi)$. The formula $\cos\varphi_{2}=1-\sin\varphi$ is the more remarkable in that it does not contain the modulus $c$, and so holds whatever the modulus. We shall therefore also have rationally $\sin\varphi_{2}=\dfrac{\cot\varphi}{b}(1-\sin\varphi)$.

The value of $\sin\varphi_{2}$ can be determined directly by the trisection of the function $F(\pi)$, for it is evident that we have $F(\varphi_{2})=\frac{2}{3}F^{1}=\frac{1}{3}F(\pi)$. Let therefore $\sin\varphi_{2}=y$, and putting $a=\sin\pi=0$ in the formula of triplication (art.~23), we shall have, to determine $y$, the equation
\[ 0=3-4(1+c^{2})y^{2}+6c^{2}y^{4}-c^{4}y^{8}, \]
which has only the difficulty of the fourth degree.

\origpage{33}
\subsection*{(25).}

Let us now occupy ourselves with the equation for the quintisection, which, when fully developed, becomes
\[ \begin{aligned} 0={}&1-4x^{2}+5c^{2}x^{4}+4b^{2}c^{2}x^{6}-5c^{4}x^{8}+4c^{6}x^{10}-c^{6}x^{12} \\ &-2x[1-5c^{2}x^{2}+(6c^{2}+4c^{4})x^{4}-(4c^{2}+6c^{4})x^{6}+5c^{4}x^{8}-c^{6}x^{10}]. \end{aligned} \]

To give at least a few particular examples of the solution of this equation, one could assign a value to $x$ (taking care that it be greater than the root of the equation $0=1-2x-4x^{2}$); and from the supposed value of $x$, we see that $c^{2}$ would be determined by an equation of the third degree.

But the search for particular cases of solution is a question of indeterminate analysis which can be resolved more easily in the following manner.

Let $p=1-c^{2}x^{4}$, $q=2x(1-c^{2}x^{2})$, and the equation of article 23 will be
\[ \frac{p+q}{p-q}=\frac{(1+x)\sqrt{(1-c^{2}x^{2})}}{bx}. \]
Raising each member to the square, and subtracting unity from both sides, we shall have
\[ \frac{4pq}{(p-q)^{2}}=\frac{p+q}{b^{2}x^{2}}. \]
Let $p=mq$, this equation will give $b^{2}x^{2}=\dfrac{(m+1)(m-1)^{2}}{4m}q$; let therefore again
\[ n=\frac{(m+1)(m-1)^{2}}{2m}, \]
and we shall have, restoring the value of $q$, $b^{2}x=n(1-c^{2}x^{2})$, which gives
\[ c^{2}=\frac{x-n}{x-nx^{2}}. \]
But from the equation $p=mq$, or $1-c^{2}x^{4}=2mx(1-c^{2}x^{2})$, we draw
\[ c^{2}=\frac{2mx-1}{2mx^{3}-x^{4}}. \]
Equating these two values of $c^{2}$, there will come $0=1-(2m+n)x+x^{2}$.

From this is deduced a very simple general solution of the problem we have proposed to ourselves. Having taken for $m$ any value whatever

\origpage{34}
comprised between 1 and $\dfrac{\sqrt{5}+1}{2}=1.618$ (for these values correspond to the limits $c=1$ and $c=0$), one deduces $n=\dfrac{(m+1)(m-1)^{2}}{2m}$; then
\[ \begin{aligned} x &= m+\tfrac{1}{2}n-\sqrt{[(m+\tfrac{1}{2}n)^{2}-1]}, \\ c^{2} &= \frac{1-n(m+\frac{1}{2}n)-n\sqrt{[(m+\frac{1}{2}n)^{2}-1]}}{1-n(m+\frac{1}{2}n)+n\sqrt{[(m+\frac{1}{2}n)^{2}-1]}}. \end{aligned} \]

Let, for example, $m=\dfrac{3}{2}$, we shall have $n=\dfrac{5}{24}$, whence
\[ \begin{aligned} x &= \frac{77-5\sqrt{145}}{48} \\ c^{2} &= \frac{767-25\sqrt{145}}{767+25\sqrt{145}}. \end{aligned} \]
From this value of the modulus $c$ and that of $x=\sin\varphi$, we shall satisfy the equation $F(\varphi)=\frac{1}{5}F^{1}$; then it will be easy to have the values of $\varphi_{2}$, $\varphi_{3}$, $\varphi_{4}$, by which we shall complete the division of the complete function $F^{1}$ into five equal parts.

\subsection*{(26).}

We have shown how one finds the relation between $\varphi_{n}$ and $\varphi$, so that $F(\varphi_{n})=nF(\varphi)$, $n$ being a whole number. If it were necessary to find the relation between $\varphi$ and $\psi$ so that $F(\psi)=\dfrac{m}{n}F(\varphi)$, $m$ and $n$ being integers, one would take an auxiliary angle $\omega$, such that one had at once $nF(\psi)=F(\omega)$ and $mF(\varphi)=F(\omega)$. The first condition would give an algebraic equation between the sines of the angles $\psi$ and $\omega$, the second would give one between those of the angles $\varphi$ and $\omega$; whence, eliminating $\omega$, we shall have the sought relation between $\varphi$ and $\psi$.

In general we see that it will always be possible to satisfy the transcendent equation
\[ 0=mF(\varphi)+nF(\psi)+pF(\omega)+\text{etc.}, \]
by an algebraic equation, the coefficients $m$, $n$, $p$, etc. being whole numbers at will, which do not all have the same sign. This is a quite simple consequence of the fact that the transcendent equation $F(\varphi)+F(\psi)=F(\mu)$ is represented by an algebraic equation.

\subsection*{(27).}

Hence it results that one can always find the complete algebraic integral

\origpage{35}
of the differential equation
\[ \frac{m\,d\varphi}{\sqrt{(1-c^{2}\sin^{2}\varphi)}}\pm\frac{n\,d\psi}{\sqrt{(1-c^{2}\sin^{2}\psi)}}=0, \]
$m$ and $n$ being whole numbers; for the integral is at first $mF(\varphi)\pm nF(\psi)=\text{const.}$ And if we suppose that when $\psi=0$, we have $\varphi=\mu$, the constant will be $mF(\mu)$, and we shall have the integral
\[ mF(\varphi)\pm nF(\psi)=mF(\mu). \]
Let $\omega$ be an auxiliary such that $F(\mu)=F(\varphi)\pm F(\omega)$, we shall have at the same time $mF(\omega)=nF(\psi)$; if we then express these two equations in algebraic terms, and eliminate $\omega$, we shall have the sought algebraic integral in which $\mu$ will be the arbitrary constant.

In general;\ednote{Printed with a semicolon after ``général''; a comma is expected, an apparent misprint.} if we had the following equation, in which $m$, $n$, $p$, etc. are positive or negative integers,
\[ 0=\frac{m\,d\varphi}{\sqrt{(1-c^{2}\sin^{2}\varphi)}}+\frac{n\,d\psi}{\sqrt{(1-c^{2}\sin^{2}\psi)}}+\frac{p\,d\omega}{\sqrt{(1-c^{2}\sin^{2}\omega)}}+\text{etc.}, \]
the complete integral will be $F(\mu)=mF(\varphi)+nF(\psi)+pF(\omega)+\text{etc.}$, $\mu$ being the arbitrary constant, and this integral can always be represented by an algebraic equation, whatever the number of terms, provided it be not infinite.

If we denote by $R(x)$ the radical $\sqrt{(\alpha+\beta x+\gamma x^{2}+\delta x^{3}+\varepsilon x^{4})}$, and by $R(y)$, $R(z)$, etc. similar radicals in $y$, $z$, etc.; if moreover $m$, $n$, $p$, etc. denote positive or negative whole numbers, it is clear that the equation
\[ 0=\frac{m\,dx}{R(x)}+\frac{n\,dy}{R(y)}+\frac{p\,dz}{R(z)}+\text{etc.}, \]
can always be reduced to the preceding form, and that thus it will always have a complete algebraic integral.

Nothing would prevent us from supposing that $z$ and the following variables were given algebraic functions of $x$ and $y$; then the preceding equation would contain only two variables, and despite the infinity of forms of which it is susceptible, it would always admit a complete algebraic integral. This is perhaps the only manner of

\origpage{36}
generalizing the result of Euler, concerning the equation $\dfrac{dx}{R(x)}+\dfrac{dy}{R(y)}=0$. Lagrange proposed to find cases of integrability of the equation $\dfrac{dx}{\sqrt{X}}+\dfrac{dy}{\sqrt{Y}}=0$, without supposing that the two polynomials $X$ and $Y$ are entirely similar${}^{(1)}$; but it does not appear that the researches of this great Geometer led him beyond Euler's equation; for the equation which he gives on page 119, as being more general, reduces to it immediately by putting $v=ky$, and giving the coefficient $k$ a suitable value. Thus it is very doubtful whether, with two terms only, Euler's equation can be generalized; but with a greater number, we see that it admits a great extension.

\subsection*{(28).}

Since the functions $F$ can be multiplied or divided at will, this property furnishes a fairly simple means of evaluating them by approximation. First we shall suppose that $\varphi$ does not exceed $\frac{1}{2}\pi$, for according to article 12, all cases reduce to that one.

This being laid down, we shall determine $\varphi_{\frac{1}{2}}$ by article 21, in such a way that $F(\varphi_{\frac{1}{2}})=\frac{1}{2}F(\varphi)$, and the integral $\displaystyle\int\frac{d\varphi_{\frac{1}{2}}}{\Delta(\varphi_{\frac{1}{2}})}$ will have an extent smaller than $\displaystyle\int\frac{d\varphi}{\Delta(\varphi)}$, by nearly one half, if $c$ is not too near unity. By a second bisection, we can make it so that $F(\varphi_{\frac{1}{4}})=\frac{1}{4}F(\varphi)$, and the integral that $F(\varphi_{\frac{1}{4}})$ represents will again have an extent nearly twice smaller, and so on. But when the amplitude $\psi$ of the formula under consideration has become very small, the function $F(\psi)$ reduces sensibly to the arc $\psi$. Therefore whatever the first value of $\varphi$, the corresponding function $F(\varphi)$ will be equal to the last term of the sequence $\varphi$, $2\varphi_{\frac{1}{2}}$, $4\varphi_{\frac{1}{4}}$, $8\varphi_{\frac{1}{8}}$, etc., and in most cases, we shall obtain this limit by a fairly short calculation.

\textbf{(1)} Mémoires de Turin, vol.\ IV.

\origpage{37}
\section*{EXAMPLE.}

Let it be proposed to find the value of $F$ when $c=\sqrt{\left(\dfrac{2+\sqrt{3}}{4}\right)}=\sin 75^{\circ}$ and $\operatorname{tang}\varphi=\sqrt{\left(\dfrac{2}{\sqrt{3}}\right)}$, we shall find by the formulas of art.\ 21, what follows
\[ \begin{aligned} \varphi &= 47^{\circ}\;3'\;30''91 & \qquad \beta &= 45^{\circ}\;0'\;0''00 \\ \varphi_{\frac{1}{2}} &= 25\;\;36\;\;\;5.64 & \beta_{\frac{1}{2}} &= 24\;\;40\;\;10.94 \\ \varphi_{\frac{1}{4}} &= 13\;\;\;6\;\;30.98 & \beta_{\frac{1}{4}} &= 12\;\;39\;\;15.83 \\ \varphi_{\frac{1}{8}} &= \;\;6\;\;35\;\;40.74 & \beta_{\frac{1}{8}} &= \;\;6\;\;22\;\;\;8.40 \\ \varphi_{\frac{1}{16}} &= \;\;3\;\;18\;\;\;8.75 & &\text{etc.} \\ &\text{etc.} & & \end{aligned} \]
\ednote{The first row is printed $30''91$; the same angle is given as $30''94$ on p.~91, and the exact value is $30.95''$.}
The last two values give
\[ \begin{aligned} 8\varphi_{\frac{1}{8}} &= 52^{\circ}\;45'\;25''.92 \\ 16\varphi_{\frac{1}{16}} &= 52\;\;\;50\;\;\;20\;.00 \end{aligned} \]
Their difference is $4'\;54''.08$; and as, by the nature of these approximations, one result must approach the limit about four times more closely than the preceding, we shall add to the last result the third of the difference $4'\;54''.08$, which is $1'\;38''.03$, and we shall thus have for the value of $F$, the very close arc
\[ 52^{\circ}\;51'\;58''.03, \]
which in parts of the radius $=\frac{\pi}{2}\times 0.5874013=0.9226878$.

This method could become very long when $c$ is very near unity, and $\varphi$ little different from $90^{\circ}$: we shall give hereafter a more expeditious one.

\section*{Application to the Lemniscate.}

\subsection*{(29).}

\rmfigure{figures/fig-4.png}{Fig.~4.}{A lemniscate with its double point C on the axis AB, the points L, M, H, K, O, I on the upper left loop, M' on the lower left loop and N', N on the right loop, with the chord CM and the ordinate MP to the axis.}

The Lemniscate is, as is known, a curve of the fourth order which has for its equation $(x^{2}+y^{2})^{2}=a^{2}(x^{2}-y^{2})$. We suppose the semi-axis $CA=a$, the abscissa $CP=x$, and the ordinate $PM=y$; if moreover we put the chord $CM=z$, we shall have in functions of $z$, $x=\dfrac{z}{a}\sqrt{\left(\dfrac{a^{2}+z^{2}}{2}\right)}$, $y=\dfrac{z}{a}\sqrt{\left(\dfrac{a^{2}-z^{2}}{2}\right)}$, whence results the arc $AM$,

\origpage{38}
\[ s=\int\frac{-a^{2}\,dz}{\sqrt{(a^{4}-z^{4})}}. \]

Let again $z=a\cos\varphi$, and $c^{2}=\frac{1}{2}$, we shall have
\[ s=\frac{a}{\sqrt{2}}\int\frac{d\varphi}{\sqrt{(1-c^{2}\sin^{2}\varphi)}}=\frac{a}{\sqrt{2}}F(\varphi), \]
a result at which we should have arrived directly by putting
\[ \begin{aligned} x &= a\Delta\cos\varphi \\ y &= \frac{a}{\sqrt{2}}\sin\varphi\cos\varphi. \end{aligned} \]
With these values, in which $\Delta$ is always taken positively, according to our usage, one follows the curve over its whole circuit, in this manner:

In the first quarter $AMC$, the values of $\varphi$ extend from $\varphi=0$ to $\varphi=\frac{1}{2}\pi$; in the second quarter $CNB$, they extend from $\varphi=\frac{1}{2}\pi$ to $\varphi=\pi$; in the third quarter $BN'C$, from $\varphi=\pi$ to $\varphi=\frac{3}{2}\pi$, and finally in the fourth $CM'A$, from $\varphi=\frac{3}{2}\pi$ to $\varphi=2\pi$.

This being laid down, if we put $\dfrac{a}{\sqrt{2}}=1$, we shall have $s=F(\varphi)$; but whatever the line one takes for unity, we see that the arcs of the Lemniscate enjoy all the properties of elliptic functions of the first species, that is to say, that they can be added, subtracted or divided algebraically like arcs of a circle. Thus, any arc $MO$ being given, one can, reckoning from a given point $H$, determine algebraically another arc $HI$ or $HL$ which shall be in a given rational ratio to the arc $MO$. For this it suffices to determine $\varphi$ from an equation of the form $F(\beta)-F(\alpha)=\pm\dfrac{m}{n}[F(\varphi)-F(\gamma)]$, to which there always corresponds an algebraic equation.

A fortiori can one divide a given arc $MO$ into a certain number of equal parts. For example, if it is a question of the quarter of the curve $AMC$, one will determine its midpoint $K$ by means of the equation $\sin^{2}\varphi=\dfrac{1}{1+b}=\dfrac{1}{1+c}$, or by the chord $CK=\dfrac{1}{c}\cos\varphi=\sqrt{(2\sqrt{2}-2)}$; if one wishes to determine the arc $AM$ equal to the third of $AOC$, it will be necessary,

\origpage{39}
according to art.\ 24, to put $\cos\varphi=\sqrt[4]{(2\sqrt{3}-3)}$, or to take the chord $CM=\sqrt{[2\sqrt{(2\sqrt{3}-3)}]}$.

The arcs of the Lemniscate represent the elliptic functions of the first species in the case where the modulus $c=\sqrt{\frac{1}{2}}=b$. It would be rather curious to inquire whether there is some other algebraic curve whose arcs represent the function $F$ for another value of $c$; but this inquiry is none the less difficult.

It would present no difficulty if one admitted arcs of a circle and logarithms in the expression of the coordinates. Indeed it is a question of satisfying the equation $dx^{2}+dy^{2}=\dfrac{d\varphi^{2}}{1-c^{2}\sin^{2}\varphi}$, which can be put under this form
\[ dx^{2}+dy^{2}=\frac{d\varphi^{2}(\cos^{2}\varphi+b^{2}\sin^{2}\varphi)(\cos^{2}\psi+\sin^{2}\psi)}{(1-c^{2}\sin^{2}\varphi^{2})^{2}}; \]
\ednote{The denominator is printed with $\sin^{2}\varphi^{2}$; $(1-c^{2}\sin^{2}\varphi)^{2}$ is meant, an apparent misprint.}
this is what one obtains generally by the values
\[ \begin{aligned} dx &= \frac{d\varphi(\cos\varphi\cos\psi-b\sin\varphi\sin\psi)}{1-c^{2}\sin^{2}\varphi} \\ dy &= -\frac{d\varphi(\cos\varphi\sin\psi+b\sin\varphi\cos\psi)}{1-c^{2}\sin^{2}\varphi}, \end{aligned} \]
in which one may take at will $\sin\psi$ as a function of $\sin$ and $\cos\varphi$. To confine ourselves to a particular case, let $\psi=0$, we shall have
\[ \begin{aligned} dx &= \frac{d\varphi\cos\varphi}{1-c^{2}\sin^{2}\varphi} \\ dy &= -\frac{b\,d\varphi\sin\varphi}{1-c^{2}\sin^{2}\varphi}, \end{aligned} \]
whence one deduces by integrating
\[ \begin{aligned} x &= \frac{1}{2c}\log\left(\frac{1+c\sin\varphi}{1-c\sin\varphi}\right) \\ y &= \frac{1}{c}\operatorname{arc}\operatorname{tang}\left(\frac{c\cos\varphi}{b}\right). \end{aligned} \]
The curve described from these two equations will therefore be such that any arc $s$, reckoned from $\varphi=0$, will have for its value $F(\varphi)$, and will generally represent an elliptic function of the first species, whatever its modulus.

\origpage{40}
\section*{Application to the motion of the simple pendulum.}

\subsection*{(30).}

The elliptic functions of the first species receive an immediate application in the determination of the motion of the simple pendulum. Let $1$ be gravity, $L$ the length of the pendulum, $H$ the height due to the velocity at the lowest point, $\psi$ the angle by which the pendulum is displaced from the vertical at the end of the time $t$, we shall have
\[ t=\int\frac{\tfrac{1}{2}d\psi\sqrt{L}}{\sqrt{\left(\dfrac{H}{2L}-\sin^{2}\tfrac{1}{2}\psi\right)}}. \]
This general formula offers two cases to consider, according as $\dfrac{H}{2L}$ shall be greater or less than unity.

In the first case, it is clear that the body will turn ceaselessly in the same direction, and will have in its successive revolutions the same velocities at the same points of the circumference. Let $\dfrac{2L}{H}=c^{2}$, and we shall have
\[ t=c\sqrt{L}\int\frac{\tfrac{1}{2}d\psi}{\sqrt{(1-c^{2}\sin^{2}\tfrac{1}{2}\psi)}}=c\sqrt{L}\cdot F(\tfrac{1}{2}\psi); \]
whence results the time of an entire revolution
\[ T=2c\sqrt{L}\cdot F^{1}. \]

In the second case, which is properly that of oscillations, we shall put $\dfrac{H}{2L}=c^{2}$, $\sin\frac{1}{2}\psi=c\sin\varphi$, and we shall have
\[ t=\sqrt{L}\int\frac{d\varphi}{\sqrt{(1-c^{2}\sin^{2}\varphi)}}=\sqrt{L}\cdot F(\varphi); \]
hence the time of a half-oscillation $=\sqrt{L}\cdot F^{1}$, and the time of the entire oscillation $T=2\sqrt{L}\cdot F^{1}$.

When the oscillations are infinitely small, we have $F^{1}=\frac{1}{2}\pi$, and $T=\pi\sqrt{L}$, which agrees with the known formulas.

Since the time employed in traversing any arc, reckoning from the vertical, is represented by an elliptic function of the first species, it follows that, an arc traversed being given

\origpage{41}
in the time $t$, one can find algebraically another arc traversed in a time that is a multiple of $t$, or in general commensurable with $t$. One can also find an arc such that the time through this arc shall be equal to the sum or to the difference of the times through two or more given arcs, and this, whether these arcs end at the vertical or do not end there.

\rmfigure{figures/fig-5.png}{Fig.~5.}{A circular arc through its lowest point A rising to B, with the vertical line AC, the points I and O on the arc between A and B, and the chords AI and AO.}

If one wishes to divide into two equal parts the time of the half-oscillation, it will be necessary, according to the known formulas, to put $\sin^{2}\varphi=\dfrac{1}{1+b}$, which will give $\sin\frac{1}{2}\psi=c\sin\varphi=\sqrt{(1-b)}$. Therefore if $AB$ is the arc of the half-oscillation, and if, after having divided the arc $AB$ into two equal parts at the point $I$, one takes the arc $AO$ such that the chord $AO$ is to the chord $AI$ as $\sqrt{2}$ is to $1$, the time of the half-oscillation will be divided into two equal parts at the point $O$.

\section*{Comparison of the elliptic functions of the second species.}

\subsection*{(31).}

Suppose that the amplitudes $\varphi$, $\psi$, $\mu$ are such that we have $F(\varphi)+F(\psi)-F(\mu)=0$, I say that we shall have at the same time $E(\varphi)+E(\psi)-E(\mu)=P$, $P$ being an algebraic quantity. Indeed, if we differentiate this equation, regarding $\mu$ as constant, we shall have
\[ dP=d\varphi\,\Delta(\varphi)+d\psi\,\Delta(\psi). \]
Putting in place of $\Delta(\varphi)$ and $\Delta(\psi)$ their values drawn from the equations $(b')$, there will come
\[ dP=d\varphi\left(\frac{\cos\varphi-\cos\psi\cos\mu}{\sin\psi\sin\mu}\right)+d\psi\left(\frac{\cos\psi-\cos\varphi\cos\mu}{\sin\varphi\sin\mu}\right), \]
or
\[ dP=\frac{\tfrac{1}{2}d(\sin^{2}\varphi+\sin^{2}\psi+2\cos\mu\cos\varphi\cos\psi)}{\sin\mu\sin\varphi\sin\psi}; \]
But from the equation $(a')$, we deduce
\[ \sin^{2}\varphi+\sin^{2}\psi+2\cos\mu\cos\varphi\cos\psi=1+\cos^{2}\mu+c^{2}\sin^{2}\mu\sin^{2}\varphi\sin^{2}\psi; \]
hence
\[ dP=\frac{\tfrac{1}{2}d(c^{2}\sin^{2}\mu\sin^{2}\varphi\sin^{2}\psi)}{\sin\mu\sin\varphi\sin\psi}=c^{2}d(\sin\mu\sin\varphi\sin\psi); \]

\origpage{42}
hence $P=c^{2}\sin\mu\sin\varphi\sin\psi$ without a constant, because $P$ must vanish when $\varphi=0$.

Thus while the elliptic functions of the first species have between them the relation $F(\varphi)+F(\psi)-F(\mu)=0$, the corresponding functions of the second species satisfy the equation
\[ E(\varphi)+E(\psi)-E(\mu)=c^{2}\sin\mu\sin\varphi\sin\psi\ldots(c'): \]
this is the general formula which will serve to compare with one another the functions of the second species, as we have compared those of the first.

If we considered more generally the function $G(\varphi)$ composed of the first and the second species, namely,
\[ G(\varphi)=E(\varphi)+kF(\varphi), \]
$k$ being any constant coefficient, it is visible that we should similarly have
\[ G(\varphi)+G(\psi)-G(\mu)=c^{2}\sin\mu\sin\varphi\sin\psi; \]
so that all the consequences that we shall draw from the equation $(c')$ for the functions $E$ will apply generally to the functions $G$.

\subsection*{(32).}

Let us denote as above by $\varphi_{2}$, $\varphi_{3}$, $\varphi_{4}$, etc. the amplitudes which give $F(\varphi_{2})=2F(\varphi)$, $F(\varphi_{3})=3F(\varphi)$, etc., we shall have, by virtue of the equation $(c')$,
\[ \begin{aligned} 2E(\varphi)-E(\varphi_{2}) &= c^{2}\sin\varphi\sin\varphi\sin\varphi_{2} \\ 3E(\varphi)-E(\varphi_{3}) &= c^{2}\sin\varphi(\sin\varphi\sin\varphi_{2}+\sin\varphi_{2}\sin\varphi_{3}) \\ 4E(\varphi)-E(\varphi_{4}) &= c^{2}\sin\varphi(\sin\varphi\sin\varphi_{2}+\sin\varphi_{2}\sin\varphi_{3}+\sin\varphi_{3}\sin\varphi_{4}) \\ &\text{etc.} \end{aligned} \]
Hence the same relation between $\varphi_{n}$ and $\varphi$, which gives $F(\varphi_{n})=nF(\varphi)$, will give $nE(\varphi)-E(\varphi_{n})$ equal to an algebraic quantity.

In general, if $m$, $n$, $p$, etc. are positive or negative whole numbers, one can arrange that
\[ mE(\varphi)+nE(\psi)+pE(\omega)+\text{etc.} \]
be equal to an algebraic quantity: for this one must establish between

\origpage{43}
the angles $\varphi$, $\psi$, $\omega$, etc., the relation which gives
\[ 0=mF(\varphi)+nF(\psi)+pF(\omega)+\text{etc.} \]
We shall observe that the functions $F(\varphi)$, $E(\varphi)$ are in general of the same sign as $\varphi$; when $\varphi$ changes sign, they also change it, and keep the same value. This being laid down, it would seem that one can satisfy the equation
\[ 0=mF(\varphi)+nF(\psi)+pF(\omega)+\text{etc.} \]
in many different ways; for one is master to change the sign of each coefficient, provided one changes at the same time the sign of the corresponding amplitude. But one can confine oneself to considering the functions $F(\varphi)$, $F(\psi)$, $F(\omega)$, etc. as always positive; and in that case, there will never be more than one relation between the angles $\varphi$, $\psi$, $\omega$, etc., which will give $0=mF(\varphi)+nF(\psi)+pF(\omega)+\text{etc.}$; then we see that the coefficients $m$, $n$, $p$, etc. cannot all be of the same sign.

\subsection*{(33).}

Such in general are the relations which the elliptic functions of the first and second species have between them: we must now enter into some details on the numerous corollaries that one can draw from the equation of arcs
\[ E(\varphi)+E(\psi)-E(\mu)=c^{2}\sin\varphi\sin\psi\sin\mu, \]
combined with the algebraic equation which it supposes and which can be put under one of these three forms,
\[ \begin{aligned} \cos\mu &= \cos\varphi\cos\psi-\sin\varphi\sin\psi\Delta(\mu) \\ \cos\psi &= \cos\mu\cos\varphi+\sin\mu\sin\varphi\Delta(\psi) \\ \cos\varphi &= \cos\mu\cos\psi+\sin\mu\sin\psi\Delta(\varphi). \end{aligned} \]
From these equations are deduced the following values, which will serve to express the second member $c^{2}\sin\varphi\sin\psi\sin\mu$ as a function of two only of the amplitudes $\varphi$, $\psi$, $\mu$.
\[ \begin{aligned} \sin\mu &= \frac{\sin\varphi\cos\psi\Delta(\psi)+\sin\psi\cos\varphi\Delta(\varphi)}{1-c^{2}\sin^{2}\varphi\sin^{2}\psi} \\ \sin\psi &= \frac{\sin\mu\cos\varphi\Delta(\varphi)-\sin\varphi\cos\mu\Delta(\mu)}{1-c^{2}\sin^{2}\mu\sin^{2}\varphi} \\ \sin\varphi &= \frac{\sin\mu\cos\psi\Delta(\psi)-\sin\psi\cos\mu\Delta(\mu)}{1-c^{2}\sin^{2}\mu\sin^{2}\psi}, \end{aligned} \]

\origpage{44}

This being laid down, here are the principal corollaries which result from the equations mentioned, for the comparison of arcs of an ellipse.

I. If we put $\mu=\frac{1}{2}\pi$, the equation of arcs will become
\[ E(\varphi)+E(\psi)-E^{1}=c^{2}\sin\varphi\sin\psi, \]
and the corresponding algebraic equation will be
\[ b\operatorname{tang}\varphi\operatorname{tang}\psi=1. \]
From it we draw $\operatorname{tang}\psi=\dfrac{1}{b}\cot\varphi$, $\sin\psi=\dfrac{\cos\varphi}{\Delta(\varphi)}$, $\cos\psi=\dfrac{b\sin\varphi}{\Delta(\varphi)}$, and $c^{2}\sin\varphi\sin\psi=\dfrac{c^{2}\sin\varphi\cos\varphi}{\Delta(\varphi)}$. We should similarly have $\operatorname{tang}\varphi=\dfrac{1}{b}\cot\psi$, $\sin\varphi=\dfrac{\cos\psi}{\Delta(\psi)}$, $\cos\varphi=\dfrac{b\sin\psi}{\Delta(\psi)}$, and $c^{2}\sin\varphi\sin\psi=\dfrac{c^{2}\sin\psi\cos\psi}{\Delta(\psi)}$.

If, from this relation between the amplitudes $\varphi$ and $\psi$, we determine on the ellipse the arcs $BM=E(\varphi)$, $BN=E(\psi)$, we shall have
\[ BM-AN=c^{2}\sin\varphi\sin\psi. \]
In this equation consists the theorem of Fagnani. It follows that, the arc $BM$ terminated at the minor axis being given, one can find a second arc $AN$ terminated at the major axis, such that the difference of these arcs $BM-AN$ is equal to an algebraic quantity $c^{2}\sin\varphi\sin\psi$; and this quantity is represented by the part of the tangent $OM$ or $ZN$ terminated at the perpendicular $CO$ or $CZ$, dropped from the centre of the ellipse.

When $\varphi$ and $\psi$ are equal to one and the same angle $\theta$, the two points $M$ and $N$ coincide in one and the same point $K$. Then we have $\operatorname{tang}^{2}\theta=\dfrac{1}{b}$, $\sin^{2}\theta=\dfrac{1}{1+b}$, $\cos^{2}\theta=\dfrac{b}{1+b}$, which gives
\[ BK-AK=1-b. \]
Therefore then the difference of the arcs $BK$, $AK$ is equal to the difference of the semi-axes $CA$, $CB$; at the same time we have
\[ \begin{aligned} BK &= \tfrac{1}{2}E^{1}+\tfrac{1}{2}(1-b) \\ AK &= \tfrac{1}{2}E^{1}-\tfrac{1}{2}(1-b). \end{aligned} \]
There is therefore at the point $K$ a sort of bisection of the quadrant of the ellipse,

\origpage{45}
since each of the parts $BK$, $AK$ can be measured by the half of this quadrant of the ellipse.

II. Let in the general equation of arcs $\varphi=\psi=\theta$, we shall have
\[ 2E(\theta)-E(\mu)=c^{2}\sin^{2}\theta\sin\mu, \]
and the corresponding algebraic equation will be $\cos^{2}\theta-\sin^{2}\theta\Delta(\mu)=\cos\mu$. This is the equation which serves for the duplication of elliptic arcs or for their bisection.

From it we draw for the duplication
\[ \sin\mu=\frac{2\sin\theta\cos\theta\Delta(\theta)}{1-c^{2}\sin^{4}\theta},\qquad \operatorname{tang}\tfrac{1}{2}\mu=\Delta(\theta)\operatorname{tang}\theta, \]
and for the bisection,
\[ \sin^{2}\theta=\frac{1-\cos\mu}{1+\Delta(\mu)}. \]
From this we see that, the arc $BM=E(\theta)$ being given, one can find another arc $BN=E(\mu)$ which is measured by the double of $BM$, so that we have $2BM-BN=c^{2}\sin^{2}\theta\sin\mu$.

\rmfigure{figures/fig-6.png}{Fig.~6.}{A quadrant of an ellipse with centre C and semi-axes CA and CB, with the points M, I, N, K, P, Q marked in turn along the arc from B to A.}

And conversely, the arc $BN=E(\mu)$ being given, one can find a second arc $BM=E(\theta)$ which is measured by the half of the arc $BN$; we shall have in effect $BM=\frac{1}{2}BN+\frac{1}{2}c^{2}\sin^{2}\theta\sin\mu=\frac{1}{2}BN+\dfrac{1-\Delta(\mu)}{2}\operatorname{tang}\frac{1}{2}\mu$.

When we put $\mu=\frac{1}{2}\pi$, the point $M$ falls at $K$, and we have as above $BK=\frac{1}{2}E^{1}+\frac{1}{2}(1-b)$.

To have anew the point of bisection of the arc $BK$, we must put $\operatorname{tang}^{2}\mu=\dfrac{1}{b}$, $\sin^{2}\mu=\dfrac{1}{1+b}$, $\Delta(\mu)=\sqrt{b}$, and we shall have
\[ \sin^{2}\theta=\frac{1-\sqrt{\left(\dfrac{1}{1+b}\right)}}{1+\sqrt{b}}. \]
\ednote{Printed $\sqrt{\left(\frac{1}{1+b}\right)}$; since $\sin^{2}\mu=\frac{1}{1+b}$, the numerator should be $1-\cos\mu=1-\sqrt{\left(\frac{b}{1+b}\right)}$; an apparent misprint.}
This value determines the arc $BI=E(\theta)$ which is measured by the half of $BK$ or by the quarter of $E^{1}$; we have in effect
\[ BI=\tfrac{1}{2}BK+\tfrac{1}{2}\cdot\frac{1-\sqrt{b}}{\sqrt{(1+b)}+\sqrt{b}}, \]

\origpage{46}
or
\[ BI=\tfrac{1}{4}E^{1}+\tfrac{1}{4}(1-b)+\tfrac{1}{2}\cdot\frac{1-\sqrt{b}}{\sqrt{(1+b)}+\sqrt{b}}, \]
and one may continue this sort of bisection to infinity.

\rmfigure{figures/fig-7.png}{Fig.~7.}{An elliptic arc starting at B, with the points M, O, N, D marked in turn along it.}

III. Having taken at will the arc $BD$ reckoned from the minor axis, with any point $M$ on this arc, there will always be another point $N$ on the same arc, such that the difference of the arcs $BM$, $DN$ will be equal to an algebraic quantity.

For this it suffices to put $BD=E(\varphi)$, $BM=E(\psi)$, $BN=E(\mu)$, and to determine $\mu$ as a function of $\varphi$ and $\psi$, as if one wished to satisfy the equation $F(\mu)=F(\varphi)-F(\psi)$, which will be done by the formula of article 19. By this means we shall have $E(\mu)+E(\psi)-E(\varphi)=c^{2}\sin\mu\sin\varphi\sin\psi$, or $BM-DN=c^{2}\sin\mu\sin\varphi\sin\psi$.

When $\mu=\psi$, the two arcs $BM$, $BN$ coincide in a single one $BO$, and each of the arcs $BO$, $OD$ is measured by half of the arc $BD$. This is a case which was examined in Corollary II.

\rmfigure{figures/fig-8.png}{Fig.~8.}{An elliptic arc starting at B, with the points N, D, M, P marked in turn along it.}

IV. Given any arc $BD$, terminating at the minor axis, with a point $M$ to serve as origin of a second arc, one can determine this second arc $MP$ or $NM$, in whichever direction one wishes, in such a way that its difference from the arc $BD$ is equal to an algebraic quantity.

In the first case, if one puts $BM=E(\varphi)$, $BD=E(\psi)$, $BP=E(\mu)$, and determines $\mu$ from the equation $F(\varphi)+F(\psi)=F(\mu)$, or from the formulas of article 18, we shall have $BD-MP=c^{2}\sin\varphi\sin\psi\sin\mu$.

In the second case, if one puts $BM=E(\varphi)$, $BD=E(\psi)$, $BN=E(\mu)$, and determines $\mu$ from the equation $F(\varphi)-F(\psi)=F(\mu)$, or from the formulas of article 19, we shall have $BD-MN=c^{2}\sin\mu\sin\varphi\sin\psi$.

\rmfigure{figures/fig-9.png}{Fig.~9.}{An elliptic arc starting at B, with the points N, O, H, M, I, D, P marked in turn along it.}

V. Given any arc $OD$ whose origin is no longer at the extremity of the minor axis, with a point $M$ to serve as origin of a second arc, one will be able to determine this second arc $MN$ or $MP$ in such a way that its difference from the arc $OD$ is equal to an algebraic quantity.

For, 1°. by Corollary III, one can find $BH-OD=$ an algebraic quantity, and then $BH-MN=$ an algebraic quantity; hence $MN-OD$ will again be an algebraic quantity; 2°. having found the point $H$ by Corol.\ III, one can find the

\origpage{47}
point $P$ by Corol.\ IV, so that $BH-MP$ is equal to an algebraic quantity, and then $OD-MP$ will also be equal to an algebraic quantity.

Thus one can find on the ellipse an infinity of arcs equal to a given arc, plus or minus a difference assignable geometrically, so that the arc will take its origin at will at any point of the ellipse, and will be directed in whichever sense one wishes.

VI. Whatever the arc $OP$ and the point $M$ taken on this arc, there will always be on the same arc another point $D$ such that the difference of the arcs $OM$, $DP$ will be equal to an algebraic quantity.

This follows immediately from Corol.\ V.

When the points $M$ and $D$ coincide in a single point $I$, each of the arcs $OI$, $IP$ is measured by half of $OP$, and we have a first bisection of this arc. One will be able likewise to find a second, a third, etc.\ to infinity.

VII. Given the arc $BM$ whose origin is at the minor axis, one can find an arc $BN$ which is equal to any multiple of the arc $BM$, plus or minus an algebraic quantity.

For by putting $BM=E(\varphi)$, $BN=E(\psi)$, if one satisfies the equation $F(\psi)=nF(\varphi)$, one will have at the same time $nE(\varphi)-E(\psi)=$ an algebraic quantity. In this case, $\psi$ would be what we have denoted above by $\varphi_{n}$, and we have seen the manner of determining $\varphi_{n}$ by means of $\varphi$.

VIII. Conversely, given any arc $BN=E(\psi)$, one will be able, by the solution of an algebraic equation, to determine the arc $BM=E(\varphi)$ which is equal to a submultiple of the arc $BN$, plus an algebraic quantity.

For example, for the trisection of the arc $BN$, one must determine $\sin\varphi$ by the equation
\[ \sin\psi=\frac{3\sin\varphi-4(1+c^{2})\sin^{3}\varphi+6c^{2}\sin^{5}\varphi-c^{4}\sin^{9}\varphi}{1-6c^{2}\sin^{4}\varphi+4c^{2}(1+c^{2})\sin^{6}\varphi-3c^{4}\sin^{8}\varphi}, \]
an equation which is in general of the ninth degree, but which reduces to the fourth when $\psi=\frac{1}{2}\pi$.

IX. In general one will be able to find, by the solution of an algebraic equation, an arc $E(\varphi)$ which is equal to a rational part

\origpage{48}
$\frac{m}{n}$ of the given arc $E(\psi)$, plus or minus an algebraic quantity. The equation will be the same as that which would give $F(\varphi)=\frac{m}{n}F(\psi)$.

It will be the same if the given arc is not terminated at the minor axis. For if $OP$ is this arc, one will seek by Corol.\ III the arc $BM$ equal to $OP+$ an algebraic quantity, and it will then only be a matter of finding an arc $BN=\frac{m}{n}BM\pm$ an algebraic quantity. One will be able afterwards to give the arc $BM$ another origin at will.

X. Two arcs being given anywhere one wishes on the ellipse, one will be able to find an arc equal to their sum or to their difference, plus or minus a straight line assignable geometrically. This is a consequence of Corol.\ III and V.

Thus all the comparisons that are ordinarily made of arcs of a circle by way of analysis hold equally for arcs of an ellipse, except for the difference which affects all the results, but which can be made to disappear in many cases, when the origin of the sought arc is arbitrary. The disappearance of this difference, when it can take place, adds one more degree of interest to the problems; it then brings the properties of arcs of an ellipse closer to those of elliptic functions of the first species. This is why we believe we should give some examples of it here.

\subsection*{(34).}

Problem I. \emph{To determine on the quadrant of an ellipse $BKA$ an arc $MP$ which is precisely equal to half of the arc $BKA$.}

Let $\varphi$ be the amplitude of the point $M$, $\psi$ that of the point $P$, $\theta$ the amplitude of the point $K$, first point of bisection of the arc $BKA$, for which we have $\sin^{2}\theta=\frac{1}{1+b}$ and $E(\theta)=\frac{1}{2}E^{1}+\frac{1}{2}(1-b)$. Suppose we have $F(\varphi)+F(\theta)-F(\psi)=0$, it will follow that $E(\varphi)+E(\theta)-E(\psi)=c^{2}\sin\varphi\sin\psi\sin\theta$; hence
\[ E(\psi)-E(\varphi)=\tfrac{1}{2}E^{1}+\tfrac{1}{2}(1-b)-c^{2}\sin\varphi\sin\psi\sin\theta. \]
Hence if we wish to have $MP=\frac{1}{2}E^{1}$, we must put
\[ c^{2}\sin\varphi\sin\psi\sin\theta=\tfrac{1}{2}(1-b), \]
which will give $\sin\varphi\sin\psi=\frac{1}{2}\sin\theta$. But on the other hand the equation

\origpage{49}
$F(\varphi)+F(\theta)-F(\psi)=0$ gives $\cos\varphi\cos\psi+\sin\varphi\sin\psi\Delta(\theta)=\cos\theta$, and since $\Delta(\theta)=\sqrt{b}$, we shall have at once $\sin\varphi\sin\psi=\frac{1}{2}\sin\theta$, and $\cos\varphi\cos\psi=\frac{1}{2}\cos\theta$. From these equations we draw
\[ \begin{aligned} \cos(\psi-\varphi)&=\tfrac{1}{2}(\cos\theta+\sin\theta)=\cos\frac{\pi}{4}\cos\left(\theta-\frac{\pi}{4}\right)\\ \cos(\psi+\varphi)&=\tfrac{1}{2}(\cos\theta-\sin\theta)=\cos\frac{\pi}{4}\cos\left(\theta+\frac{\pi}{4}\right). \end{aligned} \]
Thus the angles $\psi-\varphi$ and $\psi+\varphi$ will be known. One could also determine directly the values of $\sin\varphi$ and $\sin\psi$ by means of the formulas
\[ \begin{aligned} \sin\varphi&=\tfrac{1}{4}\sqrt{(3+4\sin\theta+2\sin^{2}\theta)}-\tfrac{1}{4}\sqrt{(3-4\sin\theta+2\sin^{2}\theta)}\\ \sin\psi&=\tfrac{1}{4}\sqrt{(3+4\sin\theta+2\sin^{2}\theta)}+\tfrac{1}{4}\sqrt{(3-4\sin\theta+2\sin^{2}\theta)}, \end{aligned} \]
and these sines will be the abscissas of the sought points $M$ and $P$.

\subsection*{(35).}

Problem II. \emph{To determine on the quadrant of an ellipse $BMA$ an arc $NQ$ which is equal to the third of $BMA$.}

If we suppose $F(\varphi)=\frac{1}{3}F^{1}$ and $F(\varphi_{2})=2F(\varphi)$, we shall have $\sin\varphi$ by the solution of the equation
\[ 0=1-2\sin\varphi+2c^{2}\sin^{3}\varphi-c^{2}\sin^{4}\varphi, \]
and $\varphi_{2}$ will be deduced from $\varphi$, either by the equation $\cos\varphi_{2}=1-\sin\varphi$, or by the equation $\sin\varphi_{2}=\frac{\cos\varphi}{\Delta}$.

This being posed, we shall have $3E(\varphi)-E^{1}=c^{2}\sin\varphi\sin\varphi_{2}(1+\sin\varphi)$, or
\[ E(\varphi)=\tfrac{1}{3}E^{1}+\tfrac{1}{3}c^{2}\sin\varphi\sin\varphi_{2}(1+\sin\varphi). \]
Suppose anew $F(\varphi)+F(\psi)-F(\omega)=0$, we shall have
\[ E(\varphi)+E(\psi)-E(\omega)=c^{2}\sin\varphi\sin\psi\sin\omega. \]
Hence if we wish to have $E(\omega)-E(\psi)=\frac{1}{3}E^{1}$, we must put $\sin\psi\sin\omega=\frac{1}{3}\sin\varphi_{2}(1+\sin\varphi)$. But besides, by virtue of the supposition made, we have the equation $\cos\psi\cos\omega+\sin\psi\sin\omega\Delta(\varphi)=\cos\varphi$; hence the unknowns $\psi$ and $\omega$ must be determined by the equations
\[ \begin{aligned} \sin\psi\sin\omega&=\tfrac{1}{3}\sin\varphi_{2}(1+\sin\varphi)=\frac{\cos^{3}\varphi}{3b\sin\varphi}\\ \cos\psi\cos\omega&=\tfrac{1}{3}\cos\varphi(2-\sin\varphi) \end{aligned} \]

\origpage{50}
They will immediately give the values of $\cos(\omega-\psi)$ and $\cos(\omega+\psi)$, from which those of $\psi$ and $\omega$ will be concluded. Knowing these values, we shall put $BN=E(\psi)$, $BQ=E(\omega)$, and the arc $NQ$ will be equal to the third of the quadrant of the ellipse $E^{1}$.

One can demonstrate that the problem is always possible by the following consideration. Let $BI=E(\varphi)$ be taken, we shall have $BI>\frac{1}{3}E^{1}$. Let there then be taken the arc $AP$ corresponding to $BI$, in such a way that the difference $BI-AP$ is equal to an algebraic quantity, that is to say, let the amplitude $\sigma$ of the point $P$ be taken so that we have $b\operatorname{tang}\varphi\operatorname{tang}\sigma=1$, we shall have $AP<\frac{1}{3}E^{1}$. But if one imagines that the arc $\frac{1}{3}E^{1}$ is carried along the quadrant of the ellipse, in such a way that its origin traverses successively the whole interval from $B$ to $P$, this arc being represented at $B$ by $BI>\frac{1}{3}E^{1}$, and at $P$ by $AP<\frac{1}{3}E^{1}$, it must, by virtue of the law of continuity, be represented exactly at some intermediate point $N$ by the arc $NQ=\frac{1}{3}E^{1}$.

To apply the general solution to a particular case, let $c^{2}=\frac{1}{2}=b^{2}$, we have found above (art.~24) $\cos\varphi=\sqrt[4]{(2\sqrt{3}-3)}=\sqrt[4]{(8\cos 30^{\circ}\sin^{2}15^{\circ})}$; we shall therefore have by trigonometric calculation,
\[ \begin{aligned} L\cos\varphi&=9.9166532 & &\\ L\sin\varphi&=9.7517253\ldots\ldots & \sin\varphi&=0.5645797\\ & & 2-\sin\varphi&=1.4354203\\ L\sin\psi\sin\omega&=9.6716281 & \sin\psi\sin\omega&=0.4694920\\ L\cos\psi\cos\omega&=9.5965110 & \cos\psi\cos\omega&=0.3949217\\ & & \cos(\omega+\psi)&=-0.0745703\\ & & \cos(\omega-\psi)&=0.8644137\\ \omega+\psi&=94^{\circ}16'35''5 & \omega&=62^{\circ}13'49''3\\ \omega-\psi&=30.11.3.1 & \psi&=32^{\circ}2'46''2 \end{aligned} \]
These values are only approximate, but the formulas found give the rigorous solution, which can even be constructed geometrically.

\rmfigure{figures/fig-10.png}{Fig.~10.}{An elliptic arc starting at B, with the points M, N, D, P, Q, R marked in turn along it.}

\subsection*{(36).}

Problem III. \emph{Given the two arcs $MN$ and $PQ$, situated as one wishes on the circumference of the ellipse, to find a third arc $DR$ equal to their sum $MN+PQ$.}

Let $\alpha$, $\beta$, $\delta$, $\varepsilon$ be the amplitudes of the given points $M$, $N$, $P$, $Q$,

\origpage{51}
so that we have $MN=E(\beta)-E(\alpha)$, $PQ=E(\varepsilon)-E(\delta)$; let $\psi$ and $\omega$ be the amplitudes of the sought points $D$, $R$, so that we have $DR=E(\omega)-E(\psi)$.

Let further $\lambda$ and $\mu$ be two amplitudes such that we have $F(\lambda)=F(\beta)-F(\alpha)$, $F(\mu)=F(\varepsilon)-F(\delta)$; we shall have at the same time
\[ \begin{aligned} E(\lambda)+E(\alpha)-E(\beta)&=c^{2}\sin\alpha\sin\beta\sin\lambda\\ E(\mu)+E(\delta)-E(\varepsilon)&=c^{2}\sin\delta\sin\varepsilon\sin\mu; \end{aligned} \]
hence
\[ MN+PQ=E(\lambda)+E(\mu)-c^{2}\sin\alpha\sin\beta\sin\lambda-c^{2}\sin\delta\sin\varepsilon\sin\mu. \]
Let finally $\varphi$ be a new amplitude such that $F(\varphi)=F(\lambda)+F(\mu)$, we shall have
\[ E(\lambda)+E(\mu)-E(\varphi)=c^{2}\sin\varphi\sin\lambda\sin\mu; \]
hence
\[ MN+PQ-E(\varphi)=c^{2}\sin\varphi\sin\lambda\sin\mu-c^{2}\sin\alpha\sin\beta\sin\lambda-c^{2}\sin\delta\sin\varepsilon\sin\mu. \]
Let us now put $F(\varphi)=F(\omega)-F(\psi)$, there will result
\[ E(\varphi)+E(\psi)-E(\omega)=c^{2}\sin\varphi\sin\psi\sin\omega; \]
but by hypothesis, we have $E(\omega)-E(\psi)=MN+PQ$; hence
\[ \sin\varphi\sin\psi\sin\omega=\sin\alpha\sin\beta\sin\lambda+\sin\delta\sin\varepsilon\sin\mu-\sin\lambda\sin\mu\sin\varphi. \]
Let, for brevity, $\sin\alpha\sin\beta\sin\lambda+\sin\delta\sin\varepsilon\sin\mu=M$, and we shall have
\[ \sin\omega\sin\psi=\frac{M}{\sin\varphi}-\sin\lambda\sin\mu. \]
Besides, the equation $F(\varphi)=F(\omega)-F(\psi)$ gives $\cos\omega\cos\psi+\sin\omega\sin\psi\Delta(\varphi)=\cos\varphi$; hence we shall have, to determine $\psi$ and $\omega$, the two equations
\[ \begin{aligned} \sin\omega\sin\psi&=\frac{M}{\sin\varphi}-\sin\lambda\sin\mu\\ \cos\omega\cos\psi&=\cos\varphi-\frac{\Delta(\varphi)}{\sin\varphi}M+\sin\lambda\sin\mu\Delta(\varphi). \end{aligned} \]
If in addition one observes that the equation $F(\varphi)=F(\lambda)+F(\mu)$ gives $\cos\varphi=\cos\lambda\cos\mu-\sin\lambda\sin\mu\Delta(\varphi)$, one will deduce from the equations

\origpage{52}
above,
\[ \begin{aligned} \cos(\omega+\psi)&=\cos(\lambda-\mu)-\frac{M}{\sin\varphi}[1+\Delta(\varphi)]\\ \cos(\omega-\psi)&=\cos(\lambda+\mu)+\frac{Mc^{2}\sin\varphi}{1+\Delta(\varphi)}. \end{aligned} \]
The immediate data being $\alpha$, $\beta$, $\delta$, $\varepsilon$, one deduces from them $\lambda$ and $\mu$ by the equations $F(\lambda)=F(\beta)-F(\alpha)$, $F(\mu)=F(\varepsilon)-F(\delta)$, which give
\[ \begin{aligned} \sin\lambda&=\frac{\sin\beta\cos\alpha\Delta(\alpha)-\sin\alpha\cos\beta\Delta(\beta)}{1-c^{2}\sin\alpha\sin^{2}\beta}\\ \sin\mu&=\frac{\sin\varepsilon\cos\delta\Delta(\delta)-\sin\delta\cos\varepsilon\Delta(\varepsilon)}{1-c^{2}\sin\delta\sin^{2}\varepsilon}. \end{aligned} \]
\ednote{In both denominators the print sets the first sine with no exponent ($\sin\alpha$, $\sin\delta$); by the addition formula $\sin^{2}\alpha$ and $\sin^{2}\delta$ are expected (cf.\ $\sin^{2}\mu\sin^{2}\lambda$ below). Apparent misprints or dropped superscripts.}
The value of $M$ is therefore known; next one will deduce $\sin\varphi$ and $\Delta(\varphi)$ from the equation $F(\varphi)=F(\lambda)+F(\mu)$, which gives
\[ \begin{aligned} \sin\varphi&=\frac{\sin\lambda\cos\mu\Delta(\mu)+\sin\mu\cos\lambda\Delta(\lambda)}{1-c^{2}\sin^{2}\mu\sin^{2}\lambda}\\ \Delta(\varphi)&=\frac{\Delta(\mu)\Delta(\lambda)-c^{2}\sin\lambda\sin\mu\cos\lambda\cos\mu}{1-c^{2}\sin^{2}\mu\sin^{2}\lambda}. \end{aligned} \]
Now the values of $\sin\lambda$, $\cos\lambda$, $\Delta(\lambda)$ are given as functions of $\beta$ and $\alpha$ by the formulas of art.~19; the same holds for the values of $\sin\mu$, $\cos\mu$, $\Delta(\mu)$ expressed as functions of $\varepsilon$ and $\delta$. We shall therefore know all the quantities which compose the values of $\cos(\omega-\psi)$ and $\cos(\omega+\psi)$.

This problem can serve to solve many others, and notably to find an arc which is exactly in a rational ratio with a given arc; but for this it is necessary that in each application the values found for $\cos(\omega+\psi)$ and $\cos(\omega-\psi)$ be each contained between the limits $+1$ and $-1$, without which the problem would become impossible.

\section*{Comparison of the arcs of a hyperbola.}

\subsection*{(37).}

We have already found (art.~13) that the arc $AM$ denoted by $\Upsilon$ has for its expression
\[ \Upsilon=\Delta\operatorname{tang}\varphi-E(\varphi)+b^{2}F(\varphi). \]

\origpage{53}
The point $M$, extremity of the arc $AM$, is supposed to have the amplitude $\varphi$. Let us consider two other points $N$ and $O$ whose amplitudes are $\psi$ and $\omega$, and suppose that we have the equation $F(\varphi)+F(\psi)-F(\omega)=0$; then the arcs $AM$, $AN$, $AO$, denoted respectively by $\Upsilon(\varphi)$, $\Upsilon(\psi)$, $\Upsilon(\omega)$, will have between them this relation
\[ \begin{aligned} \Upsilon(\varphi)+\Upsilon(\psi)-\Upsilon(\omega)={}&\Delta(\varphi)\operatorname{tang}\varphi+\Delta(\psi)\operatorname{tang}\psi-\Delta(\omega)\operatorname{tang}\omega\\ &-E(\varphi)-E(\psi)+E(\omega), \end{aligned} \]
or, putting in the known value of $E(\varphi)+E(\psi)-E(\omega)$,
\[ \begin{aligned} \Upsilon(\varphi)+\Upsilon(\psi)-\Upsilon(\omega)={}&\Delta(\varphi)\operatorname{tang}\varphi+\Delta(\psi)\operatorname{tang}\psi-\Delta(\omega)\operatorname{tang}\omega\\ &-c^{2}\sin\varphi\sin\psi\sin\omega. \end{aligned} \]
This is the fundamental equation from which one can make on the arcs of a hyperbola the same comparisons that we have made on the arcs of an ellipse, but observing that in the hyperbola one cannot give the amplitudes a value greater than $\frac{1}{2}\pi$, and that when $\varphi=\frac{1}{2}\pi$, the arc $AM$ becomes infinite.

The quantity $\Delta\operatorname{tang}\varphi$ is nothing other than the tangent $MZ$, terminated by the perpendicular $CZ$ let fall from the centre onto this tangent. Thus $\Delta\operatorname{tang}\varphi-\Upsilon(\varphi)$ is the excess of the tangent $MZ$ over the arc $AZ$\ednote{Printed $AZ$; the arc meant is $AM$, an apparent misprint.}. If we call $G(\varphi)$ this function, we shall have for each point $M$,
\[ G(\varphi)=E(\varphi)-b^{2}F(\varphi), \]
and when the three points $M$, $N$, $O$, determined by the amplitudes $\varphi$, $\psi$, $\omega$, are linked together by the relation $F(\varphi)+F(\psi)-F(\omega)=0$, the corresponding functions $G(\varphi)$, $G(\psi)$, $G(\omega)$ relating to the hyperbola will satisfy the equation
\[ G(\varphi)+G(\psi)-G(\omega)=c^{2}\sin\varphi\sin\psi\sin\omega, \]
an equation entirely similar to that which holds in the ellipse, and from which similar corollaries will be deduced, save the restrictions peculiar to the hyperbola and of which we have already spoken.

We have found that when one puts $\sin^{2}\theta=\frac{1}{1+b}$, we have $F(\theta)=\frac{1}{2}F^{1}$, and $E(\theta)=\frac{1}{2}E^{1}+\frac{1}{2}(1-b)$, we shall therefore have similarly for the hyperbola,
\[ G(\theta)=\tfrac{1}{2}G^{1}+\tfrac{1}{2}(1-b), \]

\origpage{54}
$G^{1}$ is the difference between the infinite branch of the hyperbola $AMO$ and its asymptote $CV$, which is supposed to meet it at an infinitely distant point. This difference, estimated by means of the functions $E^{1}$, $F^{1}$, has for value $G^{1}=E^{1}-b^{2}F^{1}$; but without recourse to these functions, it can be deduced from the preceding equation, which gives
\[ G^{1}=2G(\theta)-(1-b). \]
Thus the quantity $G^{1}$, difference of two infinites, will be determined by the quantity $G(\theta)$, relating to the point $K$ whose amplitude is $\theta$, and which has for coordinates $y=b^{2}\operatorname{tang}\theta=b\sqrt{b}$, $x=c\sqrt{(1+b)}$.

The amplitude $\theta$ is that which gives $F(\theta)=\frac{1}{2}F^{1}$; if one seeks successively by the formulas of bisection the amplitudes $\theta'$, $\theta''$, etc.\ such that we have $F(\theta')=\frac{1}{2}F(\theta)=\frac{1}{4}F^{1}$, $F(\theta'')=\frac{1}{2}F(\theta')=\frac{1}{8}F^{1}$, etc., we shall have at the same time
\[ \begin{gathered} G(\theta)=2G(\theta')-c^{2}\sin^{2}\theta'\sin\theta\\ G(\theta')=2G(\theta'')-c^{2}\sin^{2}\theta''\sin\theta'\\ \text{etc.} \end{gathered} \]
From which it follows that the quantity $G^{1}$ will be determined by the last term of the sequence $G(\theta)$, $G(\theta')$, $G(\theta'')$, etc., prolonged as far as one wishes. Now when $\varphi$ has become very small, the quantity $G(\varphi)$ has for its very close approximate value $c^{2}\sin\varphi$; thus the repeated bisection of the function $G(\theta)$ furnishes a means of determining by approximation the value of the transcendent $G^{1}$, and the same method applies to any function $G(\varphi)$ of which one would wish to have an approximate value. But we shall give hereafter for this object more expeditious methods.

\section*{Particular development of the formula $Z=\displaystyle\int\frac{(f+gx^{2})\,dx}{\sqrt{(\alpha^{2}+2\alpha\beta x^{2}\cos\theta+\beta^{2}x^{4})}}$,}

\subsection*{(38).}

This formula is met with rather often in applications, and moreover it is necessary to examine particularly the case of the imaginary factors of which we spoke (art.~7).

\origpage{55}

The variable $x$ is susceptible of all values, from $x=0$ to $x=\infty$; but as, on putting $x=\infty$, we have $Z=\displaystyle\int\frac{g}{\beta}\,dx$, in order to rid ourselves of the term which may become infinite, we shall consider simply the formula
\[ X=\int\left(\frac{(f+gx^{2})\,dx}{\sqrt{(\alpha^{2}+2\alpha\beta x^{2}\cos\theta+\beta^{2}x^{4})}}-\frac{g}{\beta}\,dx\right), \]
which by this means will always have a finite value.

It is now a question of transforming this expression in such a way that the binomial factors of the quantity under the radical become real. For this one may indifferently make one of the four following suppositions:
\[ \begin{gathered} \sqrt{(\alpha^{2}+2\alpha\beta x^{2}\cos\theta+\beta^{2}x^{4})}=2xy\sqrt{\alpha\beta},\\ \alpha+\beta x^{2}=2xy\sqrt{\alpha\beta},\\ \beta x^{2}+\alpha\cos\theta+\sqrt{(\alpha^{2}+2\alpha\beta x^{2}\cos\theta+\beta^{2}x^{4})}=2\alpha y^{2},\\ x=\frac{1-y}{1+y}\sqrt{\frac{\alpha}{\beta}}, \end{gathered} \]
and the transformed in $y$ will have the required conditions. Let us confine ourselves to presenting the result of the third supposition: it gives $x^{2}=\frac{\alpha}{\beta}\left(y^{2}-\cos\theta-\frac{\sin^{2}\theta}{4y^{2}}\right)$, and
\[ X=\int\frac{\frac{dy}{\sqrt{\alpha\beta}}\left(f-\frac{g\alpha}{\beta}\cos\theta-\frac{g\alpha\sin^{2}\theta}{2\beta y^{2}}\right)}{\sqrt{[(y^{2}-\cos^{2}\tfrac{1}{2}\theta)(y^{2}+\sin^{2}\tfrac{1}{2}\theta)]}}, \]
where one sees that indeed the two factors of the quantity under the radical are real. One sees also that the least value of $y$ being $\cos\frac{1}{2}\theta$, one may put $y=\frac{\cos\frac{1}{2}\theta}{\cos\varphi}$, which will give the following transformed, in which we have put $c=\sin\frac{1}{2}\theta$,
\[ X=\int\frac{f\beta-g\alpha+2g\alpha c^{2}\sin^{2}\varphi}{\beta\sqrt{\alpha\beta}}\cdot\frac{d\varphi}{\sqrt{(1-c^{2}\sin^{2}\varphi)}}, \]
and it must be remarked that in this formula we suppose $\sqrt{\alpha\beta}$ real; for if the term $2\alpha\beta x^{2}\cos\theta$ were negative, one would throw the sign $-$ upon $\cos\theta$.

\origpage{56}

This being posed, we shall have, according to the accustomed expressions,
\[ X=\frac{f\beta+g\alpha}{\beta\sqrt{\alpha\beta}}F-\frac{2g\alpha}{\beta\sqrt{\alpha\beta}}E; \]
this is the required integral taken from $x=0$, for we have
\[ \cos^{2}\varphi=\frac{-\beta x^{2}-\alpha\cos\theta+\sqrt{(\alpha^{2}+2\alpha\beta x^{2}\cos\theta+\beta^{2}x^{4})}}{2\alpha\sin^{2}\tfrac{1}{2}\theta}, \]
and conversely $x\sqrt{\frac{\beta}{\alpha}}=\frac{\sin\varphi}{\cos\varphi}\sqrt{(1-c^{2}\sin^{2}\varphi)}$; whence one sees that on putting $x=0$, we have $\varphi=0$, and that on putting $x=\infty$, we have $\varphi=\frac{1}{2}\pi$.

The total value of the integral, taken from $x=0$ to $x=\infty$, will therefore be
\[ X^{1}=\frac{f\beta+g\alpha}{\beta\sqrt{\alpha\beta}}F^{1}-\frac{2g\alpha}{\beta\sqrt{\alpha\beta}}E^{1}. \]
Here are some applications of these formulas which will lead to rather remarkable results.

\section*{EXAMPLE I.}

\subsection*{(39).}

Let it be proposed to evaluate the two integrals
\[ M=\int\frac{z^{-\frac{2}{3}}\,dz}{\sqrt{(1-z^{2})}},\quad N=\int\frac{z^{\frac{1}{3}}\,dz}{\sqrt{(1-z^{2})}} \]
between the limits $z=0$, $z=1$. We know besides that the product of these integrals $=\frac{3}{2}\pi$ (\emph{Calc.\ intég.\ d'Euler}, vol.~1, p.~244).

If in the first we put $z=(1-x^{2})^{-\frac{3}{2}}$\ednote{Printed $1-x^{2}$; the transformed integral and the limits $x=0$ to $x=\infty$ require $z=(1+x^{2})^{-\frac{3}{2}}$, an apparent misprint.}, we shall have the transformed $M=\displaystyle\int\frac{3\,dx}{\sqrt{(3+3x^{2}+x^{4})}}$, which must be integrated from $x=0$ to $x=\infty$. This formula being compared with the formula $X$, we shall have $f=3$, $g=0$, $\alpha=\sqrt{3}$, $\beta=1$, $\cos\theta=\frac{1}{2}\sqrt{3}$, $c=\sin\frac{1}{2}\theta=\frac{1}{2}\sqrt{(2-\sqrt{3})}=\sin 15^{\circ}$. Thus the indefinite integral is $M=\frac{3}{\sqrt[4]{3}}F$, and the complete integral, on putting $x=\infty$ or $\varphi=\frac{1}{2}\pi$, is
\[ M^{1}=\frac{3}{\sqrt[4]{3}}F^{1}. \]

\origpage{57}

If in the second formula we put $z=(1-x^{2})^{\frac{3}{2}}$, the transformed will be $N=\displaystyle\int\frac{3\,dx\,(1-x^{2})}{\sqrt{(3-3x^{2}+x^{4})}}$, and the integral must be taken from $x=0$ to $x=1$. Comparing this formula with the formula $X$, we shall have $f=3$, $g=-3$, $\alpha=\sqrt{3}$, $\beta=1$, $\cos\theta=-\frac{1}{2}\sqrt{3}$, $c'=\cos 15^{\circ}=\frac{1}{2}\sqrt{(2+\sqrt{3})}$; this modulus being different from that of the other formula, I distinguish it by an accent; I have therefore by substitution
\[ N=-3x+\frac{3-3\sqrt{3}}{\sqrt[4]{3}}F(c',\varphi)+\frac{6\sqrt{3}}{\sqrt[4]{3}}E(c',\varphi). \]
Now the relation between $\varphi$ and $x$ being
\[ \frac{3+2\sqrt{3}}{2}\cos^{2}\varphi=\tfrac{3}{2}-x^{2}+\sqrt{(3-3x^{2}+x^{4})}, \]
if we put $x=1$, we shall have $\cos^{2}\varphi=2\sqrt{3}-3$, or else $\operatorname{tang}\varphi=\sqrt{\left(\frac{2}{\sqrt{3}}\right)}$. But we have already seen (art 24) that in the case of $c=\frac{1}{2}\sqrt{(2+\sqrt{3})}$ and $\operatorname{tang}\varphi=\sqrt{\frac{2}{\sqrt{3}}}$, we have exactly $F(\varphi)=\frac{1}{3}F^{1}$; it follows by the formulas of art.~32 that we have at the same time $E(\varphi)=\frac{1}{3}E^{1}+\frac{1}{2\sqrt[4]{3}}$. Hence the integral $N$ taken from $x=0$ to $x=1$ will have this value
\[ N^{1}=\frac{1-\sqrt{3}}{\sqrt[4]{3}}F^{1}(c')+\frac{2\sqrt{3}}{\sqrt[4]{3}}E^{1}(c'). \]
Since then we have already found $M^{1}=\frac{3}{\sqrt[4]{3}}F^{1}(c)$, and the product $M^{1}N^{1}=\frac{3}{2}\pi$, we shall have between the three functions $F^{1}(c)$, $F^{1}(c')$, $E^{1}(c')$ this very remarkable relation:
\[ \tfrac{1}{2}\pi=F^{1}(c)\cdot\left[\frac{1-\sqrt{3}}{\sqrt{3}}F^{1}(c')+2E^{1}(c')\right], \]
whence it follows that in this particular case the arc of an ellipse $E^{1}(c')$, which is a function of the second species, can be expressed by the two functions of the first species $F^{1}(c)$, $F^{1}(c')$. Besides, one may observe that $c'=\sqrt{(1-c^{2})}=b$, that is to say, the moduli $c$ and $c'$ are complements of one another.

\origpage{58}

\section*{EXAMPLE II.}

\subsection*{(40).}

It is proposed to evaluate by elliptic functions the two integrals $P=\displaystyle\int\frac{z^{-\frac{1}{3}}\,dz}{\sqrt{(1-z^{2})}}$, $Q=\displaystyle\int\frac{z^{\frac{2}{3}}\,dz}{\sqrt{(1-z^{2})}}$, from $z=0$ to $z=1$: we know besides that their product $=\frac{3}{4}\pi$.

If in the first we put $z=(1-x^{2})^{\frac{3}{2}}$, we shall have the transformed $P=\displaystyle\int\frac{3\,dx}{\sqrt{(3-3x^{2}+x^{4})}}$, which must be integrated from $x=0$ to $x=1$. This being treated like the integral $N$ of the preceding example, we shall have as result
\[ P^{1}=\frac{1}{\sqrt[4]{3}}F^{1}(c'). \]

Let us come to the formula $Q$, and put first $z^{2}=y^{-3}$, we shall have the transformed $Q=\frac{3}{2}\displaystyle\int\frac{dy}{y^{2}\sqrt{(y^{3}-1)}}$ which must be integrated from $y=1$ to $y=\infty$. We have first, by integrating by parts,
\[ Q=\frac{3}{2}\cdot\frac{\sqrt{(y^{3}-1)}}{y}-\frac{3}{4}\int\frac{y\,dy}{\sqrt{(y^{3}-1)}}. \]
Let then $y=1+x^{2}$, and we shall have $\displaystyle\int\frac{y\,dy}{\sqrt{(y^{3}-1)}}=\displaystyle\int\frac{2(1+x^{2})\,dx}{\sqrt{(3+3x^{2}+x^{4})}}$. This being posed, the value of $Q$ may be put in the form
\[ Q=\frac{3}{2}\cdot\frac{x\sqrt{(3+3x^{2}+x^{4})}}{1+x^{2}}-\frac{3}{2}x+\frac{3}{2}\int\left[dx-\frac{(1+x^{2})\,dx}{\sqrt{(3+3x^{2}+x^{4})}}\right]. \]
This integral must be taken from $x=0$ to $x=\infty$, and as the part outside the sign vanishes at these two limits, we shall have simply
\[ Q=\frac{3}{2}\int\left[dx-\frac{(1+x^{2})\,dx}{\sqrt{(3+3x^{2}+x^{4})}}\right]. \]
Comparing with the formula $X$, we shall have $f=-\frac{3}{2}$, $g=-\frac{3}{2}$, $\alpha=\sqrt{3}$, $\beta=1$, $\cos\theta=\frac{1}{2}\sqrt{3}$, $c=\frac{1}{2}\sqrt{(2-\sqrt{3})}$; hence the required integral is
\[ Q^{1}=-\left(\frac{3+2\sqrt{3}}{2\sqrt[4]{3}}\right)F^{1}(c)+\frac{3\sqrt{3}}{\sqrt[4]{3}}E^{1}(c). \]
\ednote{The numerator is printed $3+2\sqrt{3}$; with $f=g=-\frac{3}{2}$, $\alpha=\sqrt{3}$, $\beta=1$ the formula for $X^{1}$ gives $3+3\sqrt{3}$. An apparent misprint.}

\origpage{59}
Now since we have $P^{1}Q^{1}=\frac{3}{4}\pi$, we shall have a relation between three elliptic functions, which, with that which we have found, offers these two results:
\[ \begin{gathered} \frac{\pi}{4}=F^{1}(b)\cdot\left[E^{1}(c)-\left(\frac{1+\sqrt{3}}{2\sqrt{3}}\right)F^{1}(c)\right]\\ \frac{\pi}{4}=F^{1}(c)\cdot\left[E^{1}(b)-\left(\frac{\sqrt{3}-1}{2\sqrt{3}}\right)F^{1}(b)\right]; \end{gathered} \]
whence one sees that the functions of the second species $E^{1}(c)$, $E^{1}(b)$ can, in this particular case, be expressed by the two functions of the first species $F^{1}(b)$, $F^{1}(c)$. To these two relations, which are already very remarkable, one must join a third, $F^{1}(b)=\sqrt{3}\,F^{1}(c)$, which will be demonstrated in the following example.

\section*{EXAMPLE III.}

\subsection*{(41).}

Let it be proposed to evaluate the integral $R=\displaystyle\int dz\,(1-z^{3})^{-\frac{2}{3}}$, taken from $z=0$ to $z=1$.

We shall first put $1-z^{3}=\left(\frac{z}{y}\right)^{\frac{3}{2}}$, which will give the transformed $R=\displaystyle\int\frac{dy}{\sqrt{(4y^{3}+1)}}$ to be integrated from $y=0$ to $y=\infty$. Let then $m=\sqrt[3]{4}$ and $my=x^{2}-1$, we shall have the new transformed $R=\frac{2}{m}\displaystyle\int\frac{dx}{\sqrt{(x^{4}-3x^{2}+3)}}$, which must be integrated from $x=1$ to $x=\infty$. This integral is, apart from the coefficient, the same as the integral $P$ of the preceding example; thus having regard to the limits of $R$, and putting $\sqrt[4]{3}=n$, we shall have the required value
\[ R^{1}=\frac{4}{3mn}F^{1}(b); \]
but there is another way of finding the value of $R$.

Let $1-z^{3}=\left(1-\frac{z}{y}\right)^{3}$, we shall first find the transformed $R=\displaystyle\int\frac{dy\sqrt{3}}{\sqrt{(4y^{3}-1)}}$, which must be integrated from $y=1$ to $y=\infty$. Let then $my=1+x^{2}$, we shall have $R=\frac{2\sqrt{3}}{m}\displaystyle\int\frac{dx}{\sqrt{(x^{4}+3x^{2}+3)}}$, a new formula which must be integrated from $x=\sqrt{(m-1)}$ to

\origpage{60}
$x=\infty$. Now this integral is similar to the formula $M$ of Example I, and we shall obtain likewise $R=\frac{2\sqrt{3}}{mn}F(c,\varphi)+\text{const.}$, observing that this integral must be taken from the value of $\varphi$ which gives $x=\sqrt{(m-1)}$ to the value of $\varphi$ which gives $x=\infty$; the latter is $\varphi=\frac{1}{2}\pi$, the other being named $\theta$, we shall have $R^{1}=\frac{2\sqrt{3}}{mn}[F^{1}(c)-F(c,\theta)]$. But in general,
\[ \cos^{2}\varphi=\frac{-x^{2}-\frac{3}{2}+\sqrt{(x^{4}+3x^{2}+3)}}{2c^{2}\sqrt{3}}: \]
hence by putting $x^{2}=m-1$, it will come
\[ \cos^{2}\theta=\frac{-m-\frac{1}{2}+\sqrt{(m^{2}+m+1)}}{2c^{2}\sqrt{3}}=\frac{(m-1)^{2}}{2\sqrt{3}-3}. \]
Now, according to article 24, this value of $\theta$ is that which for the modulus $c=\frac{1}{2}\sqrt{(2-\sqrt{3})}$ gives $F(\theta)=\frac{1}{3}F^{1}$; hence
\[ R^{1}=\frac{4\sqrt{3}}{3mn}F^{1}(c). \]
Comparing this value with that which we found by the other method, there results this new relation
\[ F^{1}(b)=\sqrt{3}\cdot F^{1}(c), \]
which, being joined to the two already found, shows that a single one of the four transcendents $F^{1}(c)$, $F^{1}(b)$, $E^{1}(c)$, $E^{1}(b)$ suffices to determine the other three. We have, for example, the equations
\[ \begin{gathered} \frac{\pi}{4\sqrt{3}}=F^{1}(c)\left[E^{1}(c)-\left(\frac{\sqrt{3}+1}{2\sqrt{3}}\right)F^{1}(c)\right]\\ \frac{\pi\sqrt{3}}{4}=F^{1}(b)\left[E^{1}(b)-\left(\frac{\sqrt{3}-1}{2\sqrt{3}}\right)F^{1}(b)\right], \end{gathered} \]
which serve to determine the function $E^{1}(c)$ by means of $F^{1}(c)$, and the function $E^{1}(b)$ by means of $F^{1}(b)$.

\origpage{61}

\section*{Theorem on the functions of the first and second species whose moduli are complements of one another.}

\subsection*{(42).}

In the comparisons just established between the functions $F^{1}(c)$, $E^{1}(c)$, which relate to the modulus $c=\frac{1}{2}\sqrt{(2-\sqrt{3})}=\sin 15^{\circ}$, and the functions $F^{1}(b)$, $E^{1}(b)$, which relate to the complementary modulus $b=\frac{1}{2}\sqrt{(2+\sqrt{3})}=\cos 15^{\circ}$, the three equations found lead to this remarkable result
\[ \frac{\pi}{2}=F^{1}(c)E^{1}(b)+F^{1}(b)E^{1}(c)-F^{1}(b)F^{1}(c)\ldots(d') \]
in which one sees that the two quantities $b$, $c$ may be interchanged, and that this equation is therefore verified in two cases, that of $c=\sin 15^{\circ}$, and that of $c=\sin 75^{\circ}$. It would be easy to demonstrate directly that it is still true in two other cases, when $c$ is infinitely small, and when $c=\sqrt{\frac{1}{2}}=b$; but we shall prove in general that it holds whatever $c$ may be.

To abbreviate the notation, let us denote simply by $F$, $E$, the quantities $F^{1}(c)$, $E^{1}(c)$, and by $F'$, $E'$ the quantities $F^{1}(b)$, $E^{1}(b)$, and suppose
\[ P=FE'+F'E-FF', \]
$P$ being a function of $c$ as yet unknown.

I differentiate the two members with respect to $c$, which is the only variable they contain. Now, since $E(\varphi)=\displaystyle\int\Delta\,d\varphi$, $F(\varphi)=\displaystyle\int\frac{d\varphi}{\Delta}$, $\Delta^{2}=1-c^{2}\sin^{2}\varphi$, the differentiation gives
\[ \begin{aligned} \frac{dE}{dc}&=-\int\frac{c\,d\varphi\sin^{2}\varphi}{\Delta}=\frac{1}{c}(E-F)\\ \frac{dF}{dc}&=\int\frac{c\,d\varphi\sin^{2}\varphi}{\Delta^{3}}=\frac{1}{c}\int\frac{d\varphi}{\Delta^{3}}-\frac{1}{c}\int\frac{d\varphi}{\Delta}. \end{aligned} \]
But by the formulas of art.\ 9, we have $\displaystyle\int\frac{d\varphi}{\Delta^{3}}=\frac{1}{b^{2}}\displaystyle\int\Delta\,d\varphi-\frac{c^{2}\sin\varphi\cos\varphi}{b^{2}\Delta}$, and in the case at hand, $\varphi=\frac{1}{2}\pi$, the second term vanishes: thus we shall have
\[ \frac{dF}{dc}=\frac{1}{b^{2}c}(E-b^{2}F). \]

\origpage{62}
Similarly we shall have $\frac{dE'}{db}=\frac{1}{b}(E'-F')$, $\frac{dF'}{db}=\frac{1}{c^{2}b}(E'-c^{2}F')$, and because $b\,db+c\,dc=0$, we shall deduce
\[ \begin{aligned} \frac{dE'}{dc}&=-\frac{c}{b^{2}}(E'-F')\\ \frac{dF'}{dc}&=-\frac{1}{b^{2}c}(E'-c^{2}F'). \end{aligned} \]
Substituting these values in that of $dP$, we shall have $dP=0$; hence $P=\text{const.}$ But we found in a particular case $P=\frac{1}{2}\pi$; hence the equation $(d')$ holds generally, whatever $c$ may be.

When $c=\sqrt{\frac{1}{2}}=b$, the equation $(d')$ gives
\[ \frac{\pi}{2}=F^{1}(2E^{1}-F^{1}). \]
Thus in this particular case, $E^{1}$ is again determined by $F^{1}$.

There may be a few other particular cases where the function of the second species $E^{1}(c)$ is expressed by the function of the second\ednote{Printed ``seconde''; $F^{1}(c)$ is the function of the first kind, so ``première'' is meant. An apparent misprint.} species $F^{1}(c)$; we shall indeed show that each known particular case makes known an infinity of others; but the general determination appears impossible, according to the following researches.

\section*{Differential equations expressing the mutual connection of the functions $E$ and $F$.}

\subsection*{(43).}

If we differentiate with respect to $c$ the two formulas $E=\displaystyle\int\Delta\,d\varphi$, $F=\displaystyle\int\frac{d\varphi}{\Delta}$, we shall have, as above,
\[ \begin{aligned} \frac{dE}{dc}&=\int\frac{-c\,d\varphi\sin^{2}\varphi}{\Delta}=\frac{1}{c}(E-F)\\ \frac{dF}{dc}&=\frac{1}{b^{2}c}(E-b^{2}F)-\frac{c}{b^{2}}\cdot\frac{\sin\varphi\cos\varphi}{\Delta}. \end{aligned} \]
From this result the two formulas
\[ \begin{aligned} F&=E-c\,\frac{dE}{dc}\\ E&=b^{2}\left(E+c\,\frac{dF}{dc}\right)+\frac{c^{2}\sin\varphi\cos\varphi}{\Delta}, \end{aligned}\ldots(e') \]
\ednote{Inside the bracket the print reads $E$, though the glyph also carries a reader's pen mark and the reading is not certain; the derivation and the restatement on p. 64 require $F$. An apparent misprint.}
which contain the mutual relations of the functions $E$ and $F$.

\origpage{63}

These equations may be simplified still further by putting $F=\frac{M}{c}$ and $E=Nc$, which gives
\[ \begin{aligned} M&=-c^{3}\frac{dN}{dc}\\ N&=\frac{b^{2}}{c}\cdot\frac{dM}{dc}+\frac{c\sin\varphi\cos\varphi}{\Delta}. \end{aligned} \]
But however simple these may be, one cannot integrate them by ordinary means, even if one supposed $\varphi=\frac{1}{2}\pi$, which would change the functions $F$ and $E$ into $F^{1}$ and $E^{1}$. Thus there is hardly any hope of determining $E^{1}$ in general by $F^{1}$, and still less the indefinite function $E$ by the function $F$; which maintains and justifies the distinction we have made between the elliptic functions of the first and of the second species.

A rather important observation presents itself on this subject. It will be seen by the methods to be set out below that the function $F$ can be expressed in finite terms by the function $E$ and another function of the same species, whose modulus and amplitude are deduced, according to a known law, from the modulus and amplitude belonging to the function $E$. This expression of $F$ by two functions of the species of $E$ is a veritable integral which must satisfy the equations $(e')$, and which does not fall within the ordinary procedures of integration. The success obtained in this particular case may therefore become a useful example in other researches.

\subsection*{(44).}

Let us return to the equations $(e')$, and see what use they might have.

If we had a table of arcs of an ellipse drawn up for all values of $c$ and $\varphi$, at equal and sufficiently close intervals, this same table could offer, for each arc $E$, the differential coefficient $\frac{dE}{dc}$, since, if $\alpha$ is the difference between two consecutive values of $c$, the differential coefficient is deduced from the finite differences by the formula
\[ \alpha\frac{dE}{dc}=\Delta E-\tfrac{1}{2}\Delta^{2}E+\tfrac{1}{3}\Delta^{3}E-\text{etc.}, \]
$\Delta E$, $\Delta^{2}E$, etc. being the first, second, etc. differences relative

\origpage{64}
to $c$, given immediately by the numbers of the Table. A Table so arranged would make known the function $F$ for every value of $c$ and $\varphi$, since we have $F=E-c\frac{dE}{dc}$. This is the means we had proposed in the Memoirs cited, of 1786, to avoid the use of arcs of a hyperbola, or of the functions $F$.

A Table of the values of the function $F$, drawn up in the same manner with the differential coefficient $\frac{dF}{dc}$ corresponding to each term, would be equally suited to replace the arcs of an ellipse, since the second of the equations $(e')$ gives
\[ E=b^{2}\left(F+c\frac{dF}{dc}\right)+\frac{c^{2}\sin\varphi\cos\varphi}{\Delta}. \]

The labour of forming a somewhat extensive Table, such as we have just proposed, is too considerable for one to believe that it will ever be undertaken; if however the intervals by which one makes $c$ and $\varphi$ increase were not too close, the Table could still be very useful, without the labour of its construction exceeding the patience of a zealous calculator.

I would therefore propose, as a thing quite practicable, that one construct a Table of arcs of an ellipse, or of functions $E(c,\varphi)$, for all values of $\varphi$, from degree to degree, and for all values of the modulus $c$ put in the form $\sin\theta$, also from degree to degree. This Table would contain only 8100 numbers, calculated with seven significant figures; and by adding to it the first and second differences, one could extend its use to all values of $c$ and $\varphi$. It would be easy, in each particular case, to draw from this Table the value of the coefficient $\frac{dE}{dc}$, at least for the values of $c$ and of $\varphi$ included in the Table. Finally, one could in strictness avoid entirely the formation of the coefficients $\frac{dE}{dc}$, or that of a separate Table for the functions $F$, since it will be demonstrated below that every function $F$ or every coefficient $\frac{dE}{dc}$ can be expressed by two arcs of an ellipse.

\origpage{65}

\section*{Development of the functions $F^{1}$ and $E^{1}$ in series.}

\subsection*{(45).}

From the equations $(e')$ one can deduce two differential equations of the second order suited to determine the functions $F$ and $E$ separately; these equations are
\[ \begin{aligned} (1-c^{2})\frac{ddF}{dc^{2}}+\frac{1-3c^{2}}{c}\cdot\frac{dF}{dc}-F+\frac{\sin\varphi\cos\varphi}{\Delta^{3}}&=0\\ (1-c^{2})\frac{ddE}{dc^{2}}+\frac{1-c^{2}}{c}\cdot\frac{dE}{dc}+E-\frac{\sin\varphi\cos\varphi}{\Delta}&=0; \end{aligned} \]
and when one considers the complete functions, these equations become
\[ \begin{aligned} (1-c^{2})\frac{ddF^{1}}{dc^{2}}+\frac{1-3c^{2}}{c}\cdot\frac{dF^{1}}{dc}-F^{1}&=0\\ (1-c^{2})\frac{ddE^{1}}{dc^{2}}+\frac{1-c^{2}}{c}\cdot\frac{dE^{1}}{dc}+E^{1}&=0. \end{aligned} \]
These will be simplified still further by putting $c=\sin\theta$, which would give
\[ \begin{aligned} \frac{ddF^{1}}{d\theta^{2}}+\cot 2\theta\cdot\frac{dF^{1}}{d\theta}-F^{1}&=0\\ \frac{ddE^{1}}{d\theta^{2}}+\frac{1}{\sin 2\theta}\cdot\frac{dE^{1}}{d\theta}+E^{1}&=0. \end{aligned} \]
\ednote{Misprints, kept as printed: the substitution $c=\sin\theta$ gives $2\cot 2\theta$ in the first equation and $\frac{2}{\sin 2\theta}$ in the second. The print has no 2 before $\cot 2\theta$ and a 1 over $\sin 2\theta$; a reader has added a handwritten 2 in both places.}
Finally, one may regard $F^{1}$ and $E^{1}$ as functions of the complementary modulus $b$, which will give the differential equations
\[ \begin{aligned} (1-b^{2})\frac{ddF^{1}}{db^{2}}+\frac{1-3b^{2}}{b}\cdot\frac{dF^{1}}{db}-F^{1}&=0\\ (1-b^{2})\frac{ddE^{1}}{db^{2}}-\left(\frac{1+b^{2}}{b}\right)\cdot\frac{dE^{1}}{db}+E^{1}&=0; \end{aligned} \]
whence one sees that the equation in $F^{1}$ is of the same form, whether one considers $F^{1}$ as a function of $b$ or as a function of $c$.

These equations are useful for making known the law of the development of the functions $F^{1}$, $E^{1}$ in series. There is no difficulty in developing these functions according to the powers of $c$; for the expressions $F=\displaystyle\int d\varphi\,(1-c^{2}\sin^{2}\varphi)^{-\frac{1}{2}}$, $E=\displaystyle\int d\varphi\,(1-c^{2}\sin^{2}\varphi)^{\frac{1}{2}}$ being integrated

\origpage{66}
from $\varphi=0$ to $\varphi=\frac{1}{2}\pi$, one draws from them immediately
\[ \begin{aligned} F^{1}&=\frac{\pi}{2}\left(1+\frac{1^{2}}{2^{2}}c^{2}+\frac{1^{2}\cdot 3^{2}}{2^{2}\cdot 4^{2}}\cdot c^{4}+\frac{1^{2}\cdot 3^{2}\cdot 5^{2}}{2^{2}\cdot 4^{2}\cdot 6^{2}}c^{6}+\text{etc.}\right)\\ E^{1}&=\frac{\pi}{2}\left(1-\frac{1}{2^{2}}c^{2}-\frac{1^{2}}{2^{2}\cdot 4^{2}}\cdot 3c^{4}-\frac{1^{2}\cdot 3^{2}}{2^{2}\cdot 4^{2}\cdot 6^{2}}\cdot 5c^{6}-\text{etc.}\right) \end{aligned} \]
But these series are no longer sufficiently convergent when $c$ is very near unity, and then it is fitting to arrange them according to the powers of $b$, considering $b$ as very small.

Now from the equation $E^{1}=b^{2}\left(F^{1}+c\frac{dF^{1}}{dc}\right)=b^{2}\cdot\frac{d(cF^{1})}{dc}$, one may, as a first approximation, put $E^{1}=1$, which will give $d(cF^{1})=\frac{dc}{b^{2}}=\frac{dc}{1-c^{2}}$, and consequently $cF^{1}=\frac{1}{2}\log\left(\frac{1+c}{1-c}\right)$; but in the same case, we have $c=1-\frac{1}{2}b^{2}$, hence the first approximate value of $F^{1}$ is $F^{1}=\log\left(\frac{4}{b}\right)$.

Let now $F^{1}=P\log\left(\frac{4}{b}\right)+Q$, $P$ and $Q$ being sequences arranged according to the powers of $b$; if one substitutes this value in the differential equation
\[ (1-b^{2})\frac{ddF^{1}}{db^{2}}+\frac{1-3b^{2}}{b}\cdot\frac{dF^{1}}{db}-F^{1}=0, \]
one will find that the equation for determining $P$ is absolutely similar to that which determines $F^{1}$; the latter is of the same form, whether $F$ is considered as a function of $c$ or as a function of $b$; thus we shall have first
\[ P=1+\frac{1^{2}}{2^{2}}b^{2}+\frac{1^{2}\cdot 3^{2}}{2^{2}\cdot 4^{2}}b^{4}+\frac{1^{2}\cdot 3^{2}\cdot 5^{2}}{2^{2}\cdot 4^{2}\cdot 6^{2}}b^{6}+\text{etc.} \]
Let us denote the successive coefficients by $1$, $m'$, $m''$, etc., so that we have
\[ P=1+m'b^{2}+m''b^{4}+m'''b^{6}+\text{etc.}; \]
we shall then suppose
\[ \begin{aligned} F^{1}={}&(1+m'b^{2}+m''b^{4}+m'''b^{6}+\text{etc.})\log\frac{4}{b}\\ &-m'A'b^{2}-m''A''b^{4}-m'''A'''b^{6}-\text{etc.} \end{aligned} \]
Substituting this value in the differential equation above, we

\origpage{67}
shall have, for determining the new coefficients $A'$, $A''$, etc., the following equation, whose terms follow a very simple law:
\[ \begin{aligned} 0={}&2&&+6m'b^{2}&&+10m''b^{4}&&+14m'''b^{6}&&+18m^{\text{iv}}b^{8}&&+\text{etc.}\\ &-4m'&&-6m''b^{2}&&-12m'''b^{4}&&-16m^{\text{iv}}b^{6}&&-20m^{\text{v}}b^{8}&&-\text{etc.}\\ &-2^{2}m'A'&&-4^{2}m''A''b^{2}&&-6^{2}m'''A'''b^{4}&&-8^{2}m^{\text{iv}}A^{\text{iv}}b^{6}&&-10^{2}m^{\text{v}}A^{\text{v}}b^{8}&&-\text{etc.}\\ &&&+3^{2}m'A'b^{2}&&+5^{2}m''A''b^{4}&&+7^{2}m'''A'''b^{6}&&+9^{2}m^{\text{iv}}A^{\text{iv}}b^{8}&&+\text{etc.} \end{aligned} \]
\ednote{Second row: $-6m''b^{2}$ as printed; the row's law ($-4k\,m^{(k)}b^{2k-2}$: $-4m'$, $-12m'''b^{4}$, $-16m^{\text{iv}}b^{6}$) requires $-8m''b^{2}$.}
Observing next that we have $2^{2}m'=1$, $4^{2}m''=3^{2}m'$, $6^{2}m'''=5^{2}m''$, etc., we find successively
\[ \begin{aligned} A'&=1\\ A''&=1+\frac{2}{3\cdot 4}\\ A'''&=1+\frac{2}{3\cdot 4}+\frac{2}{5\cdot 6}\\ A^{\text{iv}}&=1+\frac{2}{3\cdot 4}+\frac{2}{5\cdot 6}+\frac{2}{7\cdot 8}\\ &\text{etc.} \end{aligned} \]
so that we have in general $A^{(n)}=A^{(n-1)}+\frac{2}{(2n-1)2n}$. From this law, it is easy to see what $A^{(n)}$ becomes when $n$ is very large; for this purpose consider the sequence
\[ y=\frac{2x^{2}}{1\cdot 2}+\frac{2x^{4}}{3\cdot 4}+\frac{2x^{6}}{5\cdot 6}+\frac{2x^{8}}{7\cdot 8}+\text{etc.}, \]
which may be put in this form
\[ \begin{aligned} y={}&2x\left(x+\frac{x^{3}}{3}+\frac{x^{5}}{5}+\frac{x^{7}}{7}+\text{etc.}\right)\\ &-\left(x^{2}+\frac{x^{4}}{2}+\frac{x^{6}}{3}+\frac{x^{8}}{4}+\text{etc.}\right); \end{aligned} \]
by summing these series we shall have $y=x\log\left(\frac{1+x}{1-x}\right)+\log(1-x^{2})=(1+x)\log(1+x)+(1-x)\log(1-x)$. Let $x=1$, and we shall have $y=2\log 2$; hence the limit of the quantities $A'$, $A''$, $A'''$, etc. is $2\log 2$ or 1.386, etc. This being so, the complete value of $F^{1}$ is developed thus:

\origpage{68}
\[ \begin{aligned} F^{1}={}&\log\frac{4}{b}+\frac{1^{2}}{2^{2}}b^{2}\left(\log\frac{4}{b}-1\right)\\ &+\frac{1^{2}\cdot 3^{2}}{2^{2}\cdot 4^{2}}b^{4}\left(\log\frac{4}{b}-1-\frac{2}{3\cdot 4}\right)\\ &+\frac{1^{2}\cdot 3^{2}\cdot 5^{2}}{2^{2}\cdot 4^{2}\cdot 6^{2}}b^{6}\left(\log\frac{4}{b}-1-\frac{2}{3\cdot 4}-\frac{2}{5\cdot 6}\right)\\ &+\frac{1^{2}\cdot 3^{2}\cdot 5^{2}\cdot 7^{2}}{2^{2}\cdot 4^{2}\cdot 6^{2}\cdot 8^{2}}b^{8}\left(\log\frac{4}{b}-1-\frac{2}{3\cdot 4}-\frac{2}{5\cdot 6}-\frac{2}{7\cdot 8}\right)\\ &+\text{etc.} \end{aligned} \]
Knowing $F^{1}$, we shall have $E^{1}$ immediately from the equation $E^{1}=b^{2}F^{1}+b^{2}c\frac{dF^{1}}{dc}=b^{2}F^{1}-b(1-b^{2})\frac{dF^{1}}{db}$, and there results
\[ \begin{aligned} E^{1}={}&1+\frac{1}{2}b^{2}\left(\log\frac{4}{b}-\frac{1}{1\cdot 2}\right)\\ &+\frac{1^{2}}{2^{2}}\cdot\frac{3}{4}b^{4}\left(\log\frac{4}{b}-A'-\frac{1}{3\cdot 4}\right)\\ &+\frac{1^{2}\cdot 3^{2}}{2^{2}\cdot 4^{2}}\cdot\frac{5}{6}b^{6}\left(\log\frac{4}{b}-A''-\frac{1}{5\cdot 6}\right)\\ &+\frac{1^{2}\cdot 3^{2}\cdot 5^{2}}{2^{2}\cdot 4^{2}\cdot 6^{2}}\cdot\frac{7}{8}b^{8}\left(\log\frac{4}{b}-A'''-\frac{1}{7\cdot 8}\right)\\ &+\text{etc.}, \end{aligned} \]
an expression whose law is manifest and which agrees with that which we have given in another form (\emph{Mém. de l'Acad.}, 1786, pag.~630).

\section*{On the changes that may be made to the parameter in the elliptic functions of the third species.}

\subsection*{(46).}

Let $p=\frac{\operatorname{tang}\varphi}{\Delta}$, and let $\alpha$ be an indeterminate constant; we find by differentiation,
\[ \frac{dp}{1+\alpha p^{2}}=\frac{d\varphi}{\Delta}\cdot\frac{1-c^{2}\sin^{4}\varphi}{(1-c^{2}\sin^{2}\varphi)\cos^{2}\varphi+\alpha\sin^{2}\varphi}. \]
Suppose that the denominator $(1-c^{2}\sin^{2}\varphi)\cos^{2}\varphi+\alpha\sin^{2}\varphi$ is equal to the product of the two factors $(1+n\sin^{2}\varphi)\left(1+\frac{c^{2}}{n}\sin^{2}\varphi\right)$; the

\origpage{69}
value of $\alpha$ must be
\[ \alpha=(1+n)\left(1+\frac{c^{2}}{n}\right) \]
and then, the differential equation being decomposed thus,
\[ \frac{dp}{1+\alpha p^{2}}=\frac{d\varphi}{\Delta}\left[\frac{1}{1+n\sin^{2}\varphi}+\frac{1}{1+\frac{c^{2}}{n}\sin^{2}\varphi}-1\right], \]
its integral is
\[ \Pi(n)+\Pi\left(\frac{c^{2}}{n}\right)=F+\frac{1}{\sqrt{\alpha}}\operatorname{arc}\operatorname{tang}\frac{\sqrt{\alpha}\operatorname{tang}\varphi}{\Delta}\ldots(f'), \]
a very remarkable formula by means of which every function $\Pi$ whose parameter is greater than the modulus $c$ can always be transformed into another whose parameter will be less than $c$; for of the two parameters $n$ and $\frac{c^{2}}{n}$, it is evident that one is always greater than $c$ and the other less. The functions $\Pi$ which are thus transformed into one another have, moreover, the same modulus $c$ and the same amplitude $\varphi$; this is why, instead of denoting them by $\Pi(n,c,\varphi)$, $\Pi\left(\frac{c^{2}}{n},c,\varphi\right)$, we denote them simply by $\Pi(n)$, $\Pi\left(\frac{c^{2}}{n}\right)$, indicating only the element by which the two functions differ from one another.

The general formula $(f')$ applied to the cases $n=c$ and $n=-c$ furnishes these two corollaries,
\[ \begin{aligned} \Pi(c)&=\tfrac{1}{2}F+\frac{\frac{1}{2}}{1+c}\operatorname{arc}\operatorname{tang}\frac{(1+c)\operatorname{tang}\varphi}{\Delta}\\ \Pi(-c)&=\tfrac{1}{2}F+\frac{\frac{1}{2}}{1-c}\operatorname{arc}\operatorname{tang}\frac{(1-c)\operatorname{tang}\varphi}{\Delta}. \end{aligned} \]
Thus in these two cases, the function of the third species is reduced immediately to a function of the first species.

\subsection*{(47).}

There are three cases to consider in the formula $(f')$, according as $\alpha$ is positive, zero, or negative.

The coefficient $\alpha$ will be positive whenever the parameter $n$ is positive, or whenever, $n$ being negative, it is included between

\origpage{70}
$-1$ and $-c^{2}$; in other words $\alpha$ will be positive if the parameter $n$ is of one of the two forms $n=\cot^{2}\theta$, $n=-1+b^{2}\sin^{2}\theta$. Then the equation $(f')$ will contain a real arc of a circle, and by this equation the two functions $\Pi(n)$, $\Pi\left(\frac{c^{2}}{n}\right)$ will be brought back one to the other.

In this first case, if we make $\varphi=\frac{1}{2}\pi$, we shall have between the complete functions this relation
\[ \Pi^{1}(n)+\Pi^{1}\left(\frac{c^{2}}{n}\right)=F^{1}+\frac{\pi}{2\sqrt{\alpha}}. \]

\subsection*{(48).}

If we make $\alpha=0$, we shall have $n=-1$ or $n=-c^{2}$, then the integral $\displaystyle\int\frac{dp}{1+\alpha p^{2}}$ reducing to $p$, the equation $(f')$ becomes
\[ \Pi(-1)+\Pi(-c^{2})=F+\frac{\operatorname{tang}\varphi}{\Delta}. \]
This result is easy to verify; we have indeed, by the reductions already known,
\[ \begin{aligned} \Pi(-1)&=\int\frac{d\varphi}{\Delta\cos^{2}\varphi}=F-\frac{1}{b^{2}}E+\frac{1}{b^{2}}\Delta\operatorname{tang}\varphi\\ \Pi(-c^{2})&=\int\frac{d\varphi}{\Delta^{3}}=\frac{1}{b^{2}}E-\frac{c^{2}\sin\varphi\cos\varphi}{b^{2}\Delta}; \end{aligned} \]
and the sum of these two quantities reduces to $F+\frac{\operatorname{tang}\varphi}{\Delta}$.

The function $\Pi(-1)$, which is brought back immediately to the functions of the first and of the second species, is that which gives the rectification of the hyperbola (art.\ 13). It has an infinite value when $\varphi=\frac{1}{2}\pi$, because then it represents the total length of the curve up to its infinite extremity; and this is also what the preceding formula gives; but this formula would fail if one wished to make $\varphi>\frac{1}{2}\pi$; it would seem to give a finite value for $\Pi(-1)$, whereas this value, composed of the part where $\varphi=\frac{1}{2}\pi$ and of another part, is necessarily infinite. This at least is what appears to result from the integral formula $\Pi(-1)$ considered in itself, and without relation to any curve; for we always suppose $\Delta$ positive.

\subsection*{(49).}

It remains to examine the case where $\alpha$ is negative; this case will occur

\origpage{71}
whenever the parameter is negative, but not included between $-1$ and $-c^{2}$; so that $n$ must be represented by one or other of the formulas $n=-\frac{1}{\sin^{2}\theta}$, $n=-c^{2}\sin^{2}\theta$.

Let then $n=-c^{2}\sin^{2}\theta$, and $\alpha=-\beta$; we shall have $\beta=\frac{\cos^{2}\theta}{\sin^{2}\theta}(1-c^{2}\sin^{2}\theta)$, and the integral $\displaystyle\int\frac{dp}{1+\alpha p^{2}}$ becomes $\displaystyle\int\frac{dp}{1-\beta p^{2}}$ or $\frac{1}{2\sqrt{\beta}}\log\left(\frac{1+p\sqrt{\beta}}{1-p\sqrt{\beta}}\right)$; hence the equation $(f')$ must then be replaced by the following:
\[ \Pi(n)+\Pi\left(\frac{c^{2}}{n}\right)=F+\frac{1}{2\sqrt{\beta}}\log\left(\frac{\Delta+\sqrt{\beta}\operatorname{tang}\varphi}{\Delta-\sqrt{\beta}\operatorname{tang}\varphi}\right). \]
Putting in this equation the values of $n$ and of $\beta$, and observing that $\sqrt{(1-c^{2}\sin^{2}\theta)}$ may be denoted by $\Delta(\theta)$, while $\sqrt{(1-c^{2}\sin^{2}\varphi)}$ is denoted by $\Delta(\varphi)$, we shall have this general formula for the case of $\alpha$ negative:
\[ \Pi(-c^{2}\sin^{2}\theta)+\Pi\left(-\frac{1}{\sin^{2}\theta}\right)=F+\frac{\operatorname{tang}\theta}{2\Delta(\theta)}\log\left(\frac{\Delta(\varphi)\operatorname{tang}\theta+\Delta(\theta)\operatorname{tang}\varphi}{\Delta(\varphi)\operatorname{tang}\theta-\Delta(\theta)\operatorname{tang}\varphi}\right). \]
It is to be remarked that the second member of this equation becomes infinite when $\varphi=\theta$: indeed, as we have
\[ \Pi\left(-\frac{1}{\sin^{2}\theta}\right)=\int\frac{d\varphi}{\left(1-\frac{\sin^{2}\varphi}{\sin^{2}\theta}\right)\Delta}, \]
one sees that the denominator of the differential is zero when $\varphi=\theta$, which makes the integral infinite in this case.

When afterwards one supposes $\varphi>\theta$, the denominator becomes negative, and the value of $\Pi\left(-\frac{1}{\sin^{2}\theta}\right)$ becomes finite again through the mutual destruction of the infinite parts of opposite signs. But then the formula needs to be rectified so as not to offer the logarithm of a negative quantity, and it will have to be written thus:
\[ \Pi(-c^{2}\sin^{2}\theta)+\Pi\left(-\frac{1}{\sin^{2}\theta}\right)=F+\frac{\operatorname{tang}\theta}{2\Delta(\theta)}\log\frac{\Delta(\theta)\operatorname{tang}\varphi+\Delta(\varphi)\operatorname{tang}\theta}{\Delta(\theta)\operatorname{tang}\varphi-\Delta(\varphi)\operatorname{tang}\theta}: \]
it will hold from $\varphi=\theta$ to $\varphi=\frac{1}{2}\pi$.

When $\varphi=\frac{1}{2}\pi$, this formula gives between the complete functions

\origpage{72}
the relation
\[ \Pi^{1}(-c^{2}\sin^{2}\theta)+\Pi^{1}\left(-\frac{1}{\sin^{2}\theta}\right)=F^{1}, \]
in which one sees that the logarithms have entirely disappeared.

I observe on this last equation that the function $\Pi(-c^{2}\sin^{2}\theta)$, which represents the integral $\displaystyle\int\frac{d\varphi}{(1-c^{2}\sin^{2}\theta\sin^{2}\varphi)\Delta}$, is necessarily greater than $\displaystyle\int\frac{d\varphi}{\Delta}$ or $F$, and that thus $\Pi^{1}(-c^{2}\sin^{2}\theta)$ is greater than $F^{1}$; the preceding equation cannot therefore subsist unless $\Pi^{1}\left(-\frac{1}{\sin^{2}\theta}\right)$ is negative. Now, this result is easily explained by the nature of the function $\Pi\left(-\frac{1}{\sin^{2}\theta}\right)$, which is positive from $\varphi=0$ to $\varphi=\theta$, and negative from $\varphi=\theta$ to $\varphi=\frac{1}{2}\pi$; it is consequently necessary that the negative part be greater than the positive part.

\subsection*{(50).}

\emph{Moreover, to avoid every anomaly foreign to the object with which we are concerned, and to consider among the functions $\Pi$ only those which are positive and finite, we shall set aside, in all that follows, the case where the parameter $n$ is at once negative and greater than unity.}

If this case should be met with, one could, by means of the formula of art.\ 49, reduce the proposed function $\Pi\left(-\frac{1}{\sin^{2}\theta}\right)$ to the function $\Pi(-c^{2}\sin^{2}\theta)$, which is subject to no difficulty; the anomaly would then fall on the logarithmic term joined to this function.

This being so, the functions $\Pi$, with regard to the various values of the parameter, relate to three principal cases which require particular developments.

\emph{First case.} If the parameter $n$ is positive, one can always give it the form $n=\cot^{2}\theta$, and one will then have $\sqrt{\alpha}=\frac{\sqrt{(1-b^{2}\sin^{2}\theta)}}{\sin\theta\cos\theta}=\frac{\Delta(b,\theta)}{\sin\theta\cos\theta}$.

\emph{Second case.} If the parameter $n$ is negative, but included between $-1$ and $-c^{2}$, one can denote it by the formula $n=-1+b^{2}\sin^{2}\theta$, and then one will have $\sqrt{\alpha}=\frac{b^{2}\sin\theta\cos\theta}{\sqrt{(1-b^{2}\sin^{2}\theta)}}=\frac{b^{2}\sin\theta\cos\theta}{\Delta(b,\theta)}$.

\origpage{73}

\emph{Third case.} If the parameter $n$ is negative and smaller than $c^{2}$, one can make $n=-c^{2}\sin^{2}\theta$, and one will have $\sqrt{(\alpha)}=\frac{\cos\theta}{\sin\theta}\Delta(\theta)\cdot\sqrt{-1}$.

It will be seen below that the first two cases could be deemed to make but one, inasmuch as a function belonging to one of these cases can be transformed into a function belonging to the other case. It is not the same with the third case, which differs essentially from the other two; for the functions that relate to the first two cases always bring in arcs of a circle in their comparison; whereas the functions that relate to the third admit only logarithms in their comparison.

It would seem natural to add to these three cases that in which one would suppose the parameter $n$ imaginary; but we shall prove below that this fourth case is useless to consider, and that one can always supply its place by suitable transformations.

\subsection*{(51).}

Let now $p=\frac{\sin\varphi\cos\varphi}{\Delta}$, and $k$ an indeterminate coefficient; we shall have by differentiation,
\[ \frac{dp}{1+kp^{2}}=\frac{1-2\sin^{2}\varphi+c^{2}\sin^{4}\varphi}{1+(k-c^{2})\sin^{2}\varphi-k\sin^{4}\varphi}\cdot\frac{d\varphi}{\Delta}; \]
if we then make the denominator
\[ 1+(k-c^{2})\sin^{2}\varphi-k\sin^{4}\varphi=(1+n\sin^{2}\varphi)(1-m\sin^{2}\varphi), \]
there will result $k=mn$, $(1+n)(1-m)=b^{2}$, and the differential equation will take the form
\[ \frac{dp}{1+kp^{2}}=\frac{d\varphi}{\Delta}\left[\frac{1+n}{n}\cdot\frac{1}{1+n\sin^{2}\varphi}-\frac{1-m}{m}\cdot\frac{1}{1-m\sin^{2}\varphi}-\frac{c^{2}}{mn}\right], \]
whence we draw, by integrating,
\[ \frac{1+n}{n}\Pi(n)-\left(\frac{1-m}{m}\right)\Pi(-m)=\frac{c^{2}}{mn}F+\frac{1}{\sqrt{mn}}\operatorname{arc}\operatorname{tang}\frac{\sqrt{mn}\sin\varphi\cos\varphi}{\Delta}\ldots(g'). \]

By this formula the two functions $\Pi(n)$, $\Pi(-m)$ can be reduced one to the other, provided that between the parameters $n$ and $-m$ we have the relation
\[ (1+n)(1-m)=b^{2}. \]

\origpage{74}
This relation is such that, if we make $n=\cot^{2}\theta$, we shall have $-m=-1+b^{2}\sin^{2}\theta$; whence one sees that the first two cases of no.\ 50 make, properly speaking, only one, since the reduction of one function to the other can be made immediately by means of the formula
\[ \begin{aligned} \Pi(\cot^{2}\theta)&=\frac{b^{2}\sin^{2}\theta\cos^{2}\theta}{1-b^{2}\sin^{2}\theta}\Pi(-1+b^{2}\sin^{2}\theta)\\ &\quad+\frac{c^{2}\sin^{2}\theta}{1-b^{2}\sin^{2}\theta}F+\frac{\sin\theta\cos\theta}{\Delta(b,\theta)}\operatorname{arc}\operatorname{tang}\frac{\sin\varphi\cos\varphi\Delta(b,\theta)}{\operatorname{tang}\theta\Delta(c,\varphi)}. \end{aligned} \]

\subsection*{(52).}

If in the formula $(g')$ we make $m=c^{2}\sin^{2}\theta$, the equation of condition $(1+n)(1-m)=b^{2}$ will give $n(\cos^{2}\theta+b^{2}\sin^{2}\theta)=-c^{2}\cos^{2}\theta$, whence it follows that $n$ can also be represented by the value $n=-c^{2}\sin^{2}\lambda$, and then we shall have between $\theta$ and $\lambda$ the relation $\sin^{2}\lambda(\cos^{2}\theta+b^{2}\sin^{2}\theta)=\cos^{2}\theta$, from which results
\[ 1=b\operatorname{tang}\theta\operatorname{tang}\lambda; \]
it is the same relation that gives $F(\theta)+F(\lambda)=F^{1}$.

We shall therefore have in this case the general formula
\[ \begin{aligned} &\frac{1-c^{2}\sin^{2}\theta}{c^{2}\sin^{2}\theta}\Pi(-c^{2}\sin^{2}\theta)+\frac{1-c^{2}\sin^{2}\lambda}{c^{2}\sin^{2}\lambda}\Pi(-c^{2}\sin^{2}\lambda)\\ &=\frac{F}{\sin^{2}\theta\sin^{2}\lambda}-\frac{1}{2\sin\theta\sin\lambda}\log\left(\frac{\Delta+c^{2}\sin\theta\sin\lambda\sin\varphi\cos\varphi}{\Delta-c^{2}\sin\theta\sin\lambda\sin\varphi\cos\varphi}\right), \end{aligned} \]
\ednote{The right-hand side is printed without $c^{2}$ in its denominators; formula $(g')$ with $m=c^{2}\sin^{2}\theta$, $n=-c^{2}\sin^{2}\lambda$ gives $\frac{F}{c^{2}\sin^{2}\theta\sin^{2}\lambda}$ and $\frac{1}{2c^{2}\sin\theta\sin\lambda}$, and only these yield the case $\theta=\lambda$ on p.~75. An apparent misprint.}
by means of which one can reduce one to the other the functions $\Pi(-c^{2}\sin^{2}\theta)$, $\Pi(-c^{2}\sin^{2}\lambda)$, provided that between the parameters we have the relation $1=b\operatorname{tang}\theta\operatorname{tang}\lambda$.

Thus, not only can the functions $\Pi$ whose parameter is negative and greater than unity be reduced to the functions whose parameter is less than $c^{2}$ (art.\ 49), but these can further be reduced to the case where the parameter, always of the form $-c^{2}\sin^{2}\theta$, does not exceed $1-b$; for in the case where we should have $\theta=\lambda$, the equation $1=b\operatorname{tang}\theta\operatorname{tang}\lambda$ gives $\sin^{2}\theta=\frac{1}{1+b}$, and consequently $c^{2}\sin^{2}\theta=1-b$; in every other case, one of the parameters $-c^{2}\sin^{2}\theta$, $-c^{2}\sin^{2}\lambda$, disregarding its sign, will necessarily be smaller than $1-b$.

The case of $\theta=\lambda$ deserves to be remarked; then the general formula

\origpage{75}
gives
\[ \Pi(-1+b)=\frac{1+b}{2b}F-\frac{1}{4b}\log\left(\frac{\Delta+(1-b)\sin\varphi\cos\varphi}{\Delta-(1-b)\sin\varphi\cos\varphi}\right). \]
Thus the function of the third species $\Pi(-1+b)$ is reduced indefinitely to the first species, and we have in particular for the expression of the complete function,
\[ \Pi^{1}(-1+b)=\frac{1+b}{2b}F^{1}. \]

By choosing other values of $p$ or other differentials than $\frac{dp}{1+kp^{2}}$, one could arrive at other formulas of reduction for the functions $\Pi$; but most of these formulas would fall back into the two that we have found, or would be only a combination of these two formulas and of those that we shall set out in the following chapter; this is why we shall not dwell on them further.

\section*{Comparison of the elliptic functions of the third species.}

\subsection*{(53).}

Consider the two similar functions
\[ \Pi(\varphi)=\int\frac{d\varphi}{(1+n\sin^{2}\varphi)\Delta(\varphi)},\quad \Pi(\psi)=\int\frac{d\psi}{(1+n\sin^{2}\psi)\Delta\psi}, \]
and suppose, as we did for the functions of the first and second species, that we have the equation $F(\varphi)+F(\psi)-F(\mu)=0$, $\mu$ being a constant; then if we make $\Pi(\varphi)+\Pi(\psi)-\Pi(\mu)=Q$, we shall have
\[ Q=\int\left[\frac{d\varphi}{(1+n\sin^{2}\varphi)\Delta(\varphi)}+\frac{d\psi}{(1+n\sin^{2}\psi)\Delta(\psi)}\right], \]
this integral being taken in such a way that it vanishes when $\varphi=0$ or when $\psi=\mu$.

But since $\mu$ is constant, we have $\frac{d\varphi}{\Delta(\varphi)}+\frac{d\psi}{\Delta(\psi)}=0$, from which results
\[ Q=\int\frac{d\varphi}{\Delta(\varphi)}\cdot\frac{n(\sin^{2}\psi-\sin^{2}\varphi)}{1+n(\sin^{2}\psi+\sin^{2}\varphi)+n^{2}\sin^{2}\psi\sin^{2}\varphi}. \]

\origpage{76}
Now we found (art.\ 31) $d\varphi\Delta(\varphi)+d\psi\Delta(\psi)=c^{2}d(\sin\mu\sin\varphi\sin\psi)$; the first member reduces to $\frac{d\varphi}{\Delta(\varphi)}[\Delta^{2}(\varphi)-\Delta^{2}(\psi)]=\frac{c^{2}d\varphi}{\Delta(\varphi)}(\sin^{2}\psi-\sin^{2}\varphi)$; therefore, putting $\sin^{2}\varphi+\sin^{2}\psi=p$, $\sin\varphi\sin\psi=q$, we have~$\ldots\ \frac{d\varphi}{\Delta(\varphi)}(\sin^{2}\psi-\sin^{2}\varphi)=\sin\mu\cdot dq$; therefore, finally,
\[ Q=\int\frac{n\sin\mu\cdot dq}{1+np+n^{2}q^{2}}. \]
But from the equation $\cos\varphi\cos\psi-\sin\varphi\sin\psi\Delta(\mu)=\cos\mu$, we deduce $1-p+q^{2}=[\cos\mu+q\Delta(\mu)]^{2}$, and consequently
\[ p=\sin^{2}\mu-2q\cos\mu\Delta(\mu)+c^{2}q^{2}\sin^{2}\mu; \]
putting this value of $p$ in the preceding formula, we have
\[ Q=\int\frac{n\sin\mu\cdot dq}{1+n\sin^{2}\mu-2nq\cos\mu\Delta(\mu)+q^{2}(n^{2}+nc^{2}\sin^{2}\mu)}. \]
Carrying out the integration, then, and putting as above $\alpha=(1+n)\left(1+\frac{c^{2}}{n}\right)$, we shall have $Q$, or
\[ \Pi(\varphi)+\Pi(\psi)-\Pi(\mu)=\frac{1}{\sqrt{\alpha}}\operatorname{arc}\operatorname{tang}\left(\frac{n\sqrt{\alpha}\sin\mu\sin\varphi\sin\psi}{1+n-n\cos\mu\cos\varphi\cos\psi}\right)\ldots(h'). \]
This is the general formula which, for functions of the third species, corresponds to the formula $F(\varphi)+F(\psi)-F(\mu)=0$ for functions of the first species, and to the formula $E(\varphi)+E(\psi)-E(\mu)=c^{2}\sin\mu\sin\varphi\sin\psi$ for functions of the second species.

Thus the difference, which is zero in functions of the first species and algebraic in those of the second species, is expressed by an arc of a circle or by a logarithm in functions of the third species. I say \emph{arc of a circle} or \emph{logarithm}, for one sees that the second member of equation $(h')$ will be an arc of a circle or a logarithm, according as $\alpha$ is positive or negative, that is, according as the function $\Pi$ relates to the first two cases or to the third case of article 50.

\subsection*{(54).}

From equation $(h')$ and from all those which can be formed similarly between three functions $\Pi$, we shall conclude that if $i$, $k$, $l$, etc.\ are positive integers, it will always be possible to arrange so

\origpage{77}
that we have
\[ i\Pi(\varphi)+k\Pi(\psi)+l\Pi(\omega)+\text{etc.}=W, \]
$W$ being a quantity determinable by arcs of a circle or by logarithms. For this one must establish between the amplitudes $\varphi$, $\psi$, $\omega$, etc., the relation which gives
\[ iF(\varphi)+kF(\psi)+lF(\omega)+\text{etc.}=0, \]
and this relation can always be expressed by an algebraic equation between the sines of the angles $\varphi$, $\psi$, $\omega$, etc. This result suffers no exception, and would hold even when the parameter $n$ were imaginary. But one can consider these properties from a still more general point of view.

\subsection*{(55).}

Let $P$ or $P(\varphi)$ be a rational even function of $\sin\varphi$, and let
\[ Z(\varphi)=\int\frac{P(\varphi)\cdot d\varphi}{\sqrt{(1-c^{2}\sin^{2}\varphi)}}. \]
Let $Z(\psi)$ be a similar function of $\psi$, these functions or integrals being taken in such a way that they vanish when the amplitudes $\varphi$ and $\psi$ are zero. Suppose that we always have the equation $F(\varphi)+F(\psi)-F(\mu)=0$, which, taking $\mu$ constant, gives $\frac{d\varphi}{\Delta(\varphi)}+\frac{d\psi}{\Delta(\psi)}=0$; we shall then have
\[ Z(\varphi)+Z(\psi)-Z(\mu)=\int\frac{d\varphi}{\Delta(\varphi)}[P(\varphi)-P(\psi)]. \]
Let us put as above $\sin^{2}\varphi+\sin^{2}\psi=p$, $\sin\varphi\sin\psi=q$; there will result
\[ \begin{aligned} \sin^{2}\varphi&=\tfrac{1}{2}p+\tfrac{1}{2}\sqrt{(p^{2}-4q^{2})}\\ \sin^{2}\psi&=\tfrac{1}{2}p-\tfrac{1}{2}\sqrt{(p^{2}-4q^{2})}. \end{aligned} \]
Let us now substitute the value of $\sin^{2}\varphi$ in $P(\varphi)$; the result will be of the form
\[ P(\varphi)=M+N\sqrt{(p^{2}-4q^{2})}, \]
$M$ and $N$ being rational functions of $p$ and $q$. We shall likewise have
\[ P(\psi)=M-N\sqrt{(p^{2}-4q^{2})}; \]
therefore $P(\varphi)-P(\psi)=2N\sqrt{(p^{2}-4q^{2})}=2N(\sin^{2}\varphi-\sin^{2}\psi)$, and

\origpage{78}
\[ \int\frac{d\varphi}{\Delta(\varphi)}[P(\varphi)-P(\psi)]=\int\frac{2N\,d\varphi(\sin^{2}\varphi-\sin^{2}\psi)}{\Delta(\varphi)}=-2\int N\,dq\sin\mu; \]
therefore, finally,
\[ Z(\varphi)+Z(\psi)-Z(\mu)=-2\sin\mu\int N\,dq. \]
In this formula $N$ is a rational function of $p$ and of $q$; if one substitutes in it the value of $p$ in terms of $q$, namely $p=\sin^{2}\mu-2q\cos\mu\Delta(\mu)+c^{2}q^{2}\sin^{2}\mu$, $N$ will be a rational function of $q$ alone, and thus the value of $Z(\varphi)+Z(\psi)-Z(\mu)$ can always be determined by arcs of a circle and by logarithms.

In general, if $i$, $k$, $l$, etc.\ denote positive or negative whole numbers, and if one establishes between the angles $\varphi$, $\psi$, $\omega$, etc.\ the relation which gives $iF(\varphi)+kF(\psi)+lF(\omega)+\text{etc.}=0$, we shall have at the same time
\[ iZ(\varphi)+kZ(\psi)+lZ(\omega)+\text{etc.}=W, \]
$W$ being a quantity determinable by arcs of a circle and by logarithms. The same property will hold even if $P$ should contain odd powers of $\sin\varphi$; for the part of $dZ$ affected by the odd powers would be integrated by arcs of a circle and by logarithms.

It is absolutely the same for the function
\[ Z(x)=\int\frac{P\,dx}{\sqrt{(\alpha+\beta x+\gamma x^{2}+\delta x^{3}+\varepsilon x^{4})}}, \]
$P$ being a rational function of $x$, and one can always find an algebraic equation between $x$, $y$, $z$, etc., such that the quantity
\[ iZ(x)+kZ(y)+lZ(z)+\text{etc.} \]
is determinable by arcs of a circle and logarithms.

\subsection*{(56).}

Let us return to formula $(h')$, from which must be deduced everything that concerns the comparison of functions of the third species. We have already observed that in this formula the arc of a circle must be replaced by a logarithm when the parameter is of the form $n=-c^{2}\sin^{2}\theta$; then we have $\sqrt{\alpha}=\frac{\cos\theta}{\sin\theta}\Delta(\theta)\cdot\sqrt{-1}$, and since in general $\frac{1}{\sqrt{-1}}\operatorname{arc}\operatorname{tang}z\sqrt{-1}=\frac{1}{2}\log\frac{1+z}{1-z}$, if we put

\origpage{79}
for brevity,
\[ \begin{aligned} \frac{\sin\theta\cos\mu\Delta(\mu)-\cos\theta\sin\mu\Delta(\theta)}{1-c^{2}\sin^{2}\theta\sin^{2}\mu}&=\sin\mu',\\ \frac{\sin\theta\cos\mu\Delta(\mu)+\cos\theta\sin\mu\Delta(\theta)}{1-c^{2}\sin^{2}\theta\sin^{2}\mu}&=\sin\mu'', \end{aligned} \]
formula $(h')$ will change into this one:
\[ \Pi(\varphi)+\Pi(\psi)-\Pi(\mu)=\frac{\operatorname{tang}\theta}{2\Delta(\theta)}\log\left(\frac{1+c^{2}\sin\theta\sin\mu'\sin\varphi\sin\psi}{1+c^{2}\sin\theta\sin\mu''\sin\varphi\sin\psi}\right), \]
and one may remark that the auxiliary angles $\mu'$, $\mu''$ are those which would give
\[ F(\theta)-F(\mu)=F(\mu'),\quad F(\theta)+F(\mu)=F(\mu''). \]
For the rest, since one can immediately convert into logarithms the arcs whose tangents are of the form $t\sqrt{-1}$, one may confine oneself to employing only equation $(h')$ in all the comparisons that will have to be made of the functions $\Pi$, according to the various values of $n$; but to facilitate applications, we believe we should set down here the formulas which hold in some of the simplest cases, adding to them those which concern the functions of the first and second species.

\subsection*{(57).}

Let, 1°. $\mu=\frac{1}{2}\pi$; we shall have successively for the three species of elliptic functions:
\[ \begin{aligned} F(\varphi)+F(\psi)-F^{1}&=0\\ E(\varphi)+E(\psi)-E^{1}&=c^{2}\sin\varphi\sin\psi\\ \Pi(\varphi)+\Pi(\psi)-\Pi^{1}&=\frac{1}{\sqrt{\alpha}}\operatorname{arc}\operatorname{tang}\frac{n\sqrt{\alpha}\sin\varphi\sin\psi}{1+n}. \end{aligned} \]
The relation between $\varphi$ and $\psi$ is given by the equation $b\operatorname{tang}\varphi\operatorname{tang}\psi=1$, from which one draws $\sin\psi=\frac{\cos\varphi}{\Delta(\varphi)}$ and $\sin\varphi=\frac{\cos\psi}{\Delta(\psi)}$. As for the value of $\alpha$, it is expressed according to the different forms of $n$, as has been seen (art.\ 50).

If in particular we have $\varphi=\psi$, which gives $\operatorname{tang}\varphi=\frac{1}{\sqrt{b}}$, $\sin\varphi=\frac{1}{\sqrt{(1+b)}}$, the preceding formulas become

\origpage{80}
\[ \begin{aligned} F(\varphi)&=\tfrac{1}{2}F^{1}\\ E(\varphi)&=\tfrac{1}{2}E^{1}+\tfrac{1}{2}(1-b)\\ \Pi(\varphi)&=\tfrac{1}{2}\Pi^{1}+\frac{1}{2\sqrt{\alpha}}\operatorname{arc}\operatorname{tang}\frac{n\sqrt{\alpha}}{(1+n)(1+b)}; \end{aligned} \]
they thus serve for the bisection of the complete function.

Let, 2°. $\psi=\varphi$ and $\mu=\varphi_{2}$; the amplitude $\varphi_{2}$ will be deduced from $\varphi$ by the formulas of article 21, and we shall have for the duplication of the functions the formulas
\[ \begin{aligned} 2F(\varphi)-F(\varphi_{2})&=0\\ 2E(\varphi)-E(\varphi_{2})&=c^{2}\sin^{2}\varphi\sin\varphi_{2}\\ 2\Pi(\varphi)-\Pi(\varphi_{2})&=\frac{1}{\sqrt{\alpha}}\operatorname{arc}\operatorname{tang}\left(\frac{n\sqrt{\alpha}\sin^{2}\varphi\sin\varphi_{2}}{1+n-n\cos^{2}\varphi\cos\varphi_{2}}\right). \end{aligned} \]
The same formulas will serve for bisection, by determining $\varphi$ by means of $\varphi_{2}$.

Let, 3°. $\psi=\varphi_{2}$ and $\mu=\varphi_{3}$, $\varphi_{3}$ being the amplitude which gives $F(\varphi_{3})=3F(\varphi)$, and which is determined by the formulas of article 22; we shall have for the triplication of the functions the formulas
\[ \begin{aligned} 3F(\varphi)-F(\varphi_{3})&=0\\ 3E(\varphi)-E(\varphi_{3})&=c^{2}\sin\varphi\sin\varphi_{2}(\sin\varphi+\sin\varphi_{3})\\ 3\Pi(\varphi)-\Pi(\varphi_{3})&=\frac{1}{\sqrt{\alpha}}\operatorname{arc}\operatorname{tang}\left(\frac{n\sqrt{\alpha}\sin\varphi\sin\varphi_{2}\sin\varphi_{3}}{1+n-n\cos\varphi\cos\varphi_{2}\cos\varphi_{3}}\right)\\ &\quad+\frac{1}{\sqrt{\alpha}}\operatorname{arc}\operatorname{tang}\left(\frac{n\sqrt{\alpha}\sin^{2}\varphi\sin\varphi_{2}}{1+n-n\cos^{2}\varphi\cos\varphi_{2}}\right). \end{aligned} \]
In the case where $\varphi_{3}=\frac{1}{2}\pi$, we have for the trisection of the complete functions the formulas
\[ \begin{aligned} 3F(\varphi)&=F^{1}\\ 3E(\varphi)&=E^{1}+c^{2}\sin\varphi\sin\varphi_{2}(1+\sin\varphi)\\ 3\Pi(\varphi)&=\Pi^{1}+\frac{1}{\sqrt{\alpha}}\operatorname{arc}\operatorname{tang}\left(\frac{n\sqrt{\alpha}\sin\varphi\sin\varphi_{2}}{1+n}\right)\\ &\quad+\frac{1}{\sqrt{\alpha}}\operatorname{arc}\operatorname{tang}\left(\frac{n\sqrt{\alpha}\sin\varphi\sin\varphi_{2}}{1+n-n\cos^{2}\varphi\cos\varphi_{2}}\right). \end{aligned} \]
\ednote{In the last line the numerator is printed $n\sqrt{\alpha}\sin\varphi\sin\varphi_{2}$; the second term of the triplication formula above, of which this is the case $\varphi_{3}=\frac{1}{2}\pi$, gives $n\sqrt{\alpha}\sin^{2}\varphi\sin\varphi_{2}$. An apparent misprint.}
As for the values of $\varphi$ and $\varphi_{2}$, they will be found by articles 23 and 24.

\origpage{81}
\section*{Formation of an infinite series of elliptic functions of the first species, linked together by constant ratios.}

\subsection*{(58).}

Up to now we have compared elliptic functions of the first species with one another only insofar as they had the same modulus, or could be considered as representing different arcs of one and the same curve; these comparisons were then extended, on the same principle, to functions of the second and third species; and the theorems contained in the formulas $(f')$ and $(g')$ still suppose that the modulus is the same in the two functions compared.

We shall now show that one can, by a very simple law, form an infinity of elliptic functions of the first species, which differ from one another both in modulus and in amplitude, but which have the very remarkable property of being in constant ratios to one another.

Let us consider the two functions $F(c,\varphi)=\displaystyle\int\frac{d\varphi}{\sqrt{(1-c^{2}\sin^{2}\varphi)}}$, $F(c',\varphi')=\displaystyle\int\frac{d\varphi'}{\sqrt{(1-c'^{2}\sin^{2}\varphi')}}$; I say that if we put $c'=\frac{2\sqrt{c}}{1+c}$, and determine $\varphi'$ from the equation $\sin(2\varphi'-\varphi)=c\sin\varphi$, we shall have generally
\[ F(c',\varphi')=\frac{1+c}{2}F(c,\varphi). \]
Indeed the assumed equation gives first $\cos(2\varphi'-\varphi)=\Delta$; thus we shall have successively:
\[ \begin{aligned} \cos2\varphi'&=\Delta\cos\varphi-c\sin^{2}\varphi\\ 2\cos^{2}\varphi'&=1-c\sin^{2}\varphi+\Delta\cos\varphi\\ 2\sin^{2}\varphi'&=1+c\sin^{2}\varphi-\Delta\cos\varphi\\ 2\sin\varphi'\cos\varphi'&=\sin\varphi(c\cos\varphi+\Delta)\\ 2d\varphi'&=\frac{d\varphi}{\Delta}(c\cos\varphi+\Delta)\\ \sqrt{(1-c'^{2}\sin^{2}\varphi')}&=\frac{c\cos\varphi+\Delta}{1+c}. \end{aligned} \]

\origpage{82}
From the last two we draw $\frac{d\varphi'}{\sqrt{(1-c'^{2}\sin^{2}\varphi')}}=\frac{1+c}{2}\cdot\frac{d\varphi}{\Delta}$, which gives on integrating, $F(c',\varphi')=\frac{1+c}{2}F(c,\varphi)$. We add no constant, because the amplitudes $\varphi$ and $\varphi'$ vanish at the same time. One sees therefore from this result that the functions $F(c',\varphi')$, $F(c,\varphi)$ will be in a constant ratio to one another, whatever the amplitudes $\varphi'$ and $\varphi$, provided they are linked together by the equation $\sin(2\varphi'-\varphi)=c\sin\varphi$.

Let us observe that, since the arc $\varphi$ increases indefinitely, and may comprise as many circumferences as one pleases, the arc $\varphi'$ also increases indefinitely, but in such a way that $2\varphi'-\varphi$ is always enclosed between the limits $+\theta$ and $-\theta$, $\theta$ being the arc whose sine is $c$. Indeed, since we have $\sin(2\varphi'-\varphi)=c\sin\varphi$ and $\cos(2\varphi'-\varphi)=\Delta$, $\Delta$ being always positive, one sees that $2\varphi'-\varphi$ is always equal to the smallest arc $\delta$, positive or negative, determined by the equation $\sin\delta=c\sin\varphi=\sin\theta\sin\varphi$; thus we always have $\varphi'=\frac{1}{2}\varphi+\frac{1}{2}\delta$ or $\varphi=2\varphi'-\delta$. In accordance with this observation, no ambiguity need ever be feared in determining the respective values of $\varphi'$ and $\varphi$.

If we put $\varphi'=\frac{1}{2}\pi$, we shall have $\varphi=\pi$ and $F(c,\varphi)=2F^{1}(c)$; thus the complete functions $F^{1}(c')$, $F^{1}(c)$ have between them this relation
\[ F^{1}(c')=(1+c)F^{1}(c). \]

\subsection*{(59).}

Let us now conceive that, starting from the given term $c$, one forms an infinite series of moduli $c$, $c'$, $c''$, $c'''$, etc., according to the law
\[ c'=\frac{2\sqrt{c}}{1+c},\quad c''=\frac{2\sqrt{c'}}{1+c'},\quad c'''=\frac{2\sqrt{c''}}{1+c''},\quad\text{etc.} \]
This series of moduli, which is continually increasing, will have unity as its limit, and will sensibly attain this limit after a fairly small number of terms.

If by analogy we call $b'$, $b''$, $b'''$, etc.\ the complements of the moduli $c'$, $c''$, $c'''$, etc., the series $b'$, $b''$, $b'''$, etc.\ will be continually decreasing, and each term will be deduced from the preceding modulus according to this law:
\[ b'=\frac{1-c}{1+c},\quad b''=\frac{1-c'}{1+c'},\quad b'''=\frac{1-c''}{1+c''},\quad\text{etc.} \]

\origpage{83}

Let next $\varphi$, $\varphi'$, $\varphi''$, $\varphi'''$, etc.\ be the series of amplitudes each of which is deduced from the preceding by the formulas
\[ \begin{gathered} \sin(2\varphi'-\varphi)=c\sin\varphi\\ \sin(2\varphi''-\varphi')=c'\sin\varphi'\\ \sin(2\varphi'''-\varphi'')=c''\sin\varphi''\\ \text{etc.} \end{gathered} \]
we shall form in this way an infinite series of functions of the first species $F(c,\varphi)$, $F(c',\varphi')$, $F(c'',\varphi'')$, etc., between which we shall have the equations
\[ \begin{aligned} &F(c',\varphi')=\frac{1+c}{2}F(c,\varphi)\\ &F(c'',\varphi'')=\frac{1+c'}{2}F(c',\varphi')=\frac{1+c}{2}\cdot\frac{1+c'}{2}F(c,\varphi)\\ &F(c''',\varphi''')=\frac{1+c''}{2}F(c'',\varphi'')=\frac{1+c}{2}\cdot\frac{1+c'}{2}\cdot\frac{1+c''}{2}F(c,\varphi),\\ &\qquad\text{etc.} \end{aligned} \]
whence it follows that any two of these functions are always in a constant ratio to one another for all values of the corresponding amplitudes.

As for the complete functions, their ratios will likewise be constant, and the equation $F^{1}(c')=(1+c)F^{1}(c)$, already found, will give successively
\[ \begin{aligned} &F^{1}(c')=(1+c)F^{1}(c)\\ &F^{1}(c'')=(1+c')F^{1}(c')=(1+c)(1+c')F^{1}(c)\\ &F^{1}(c''')=(1+c'')F^{1}(c'')=(1+c)(1+c')(1+c'')F^{1}(c)\\ &\qquad\text{etc.} \end{aligned} \]

\subsection*{(60).}

These results can be given a still greater extension. Indeed, the infinite series of moduli $c$, $c'$, $c''$, etc., which is increasing in one direction and has unity as its limit, can be prolonged to infinity in the opposite direction, where it will be decreasing and will have zero as its limit. Let us denote by $c$, $c^{0}$, $c^{00}$, $c^{000}$, etc.\ this decreasing series; the law which links two consecutive terms will likewise be
\[ c=\frac{2\sqrt{c^{0}}}{1+c^{0}},\quad c^{0}=\frac{2\sqrt{c^{00}}}{1+c^{00}},\quad c^{00}=\frac{2\sqrt{c^{000}}}{1+c^{000}},\quad\text{etc.} \]

\origpage{84}
But to form each term by means of the preceding one, it will be simpler to use the following values, where $b^{0}$, $b^{00}$, $b^{000}$, etc.\ denote the complements of the moduli $c^{0}$, $c^{00}$, $c^{000}$, etc.,
\[ c^{0}=\frac{1-b}{1+b},\quad c^{00}=\frac{1-b^{0}}{1+b^{0}},\quad c^{000}=\frac{1-b^{00}}{1+b^{00}},\quad\text{etc.} \]
This being laid down, if we denote by $\varphi$, $\varphi^{0}$, $\varphi^{00}$, $\varphi^{000}$, etc.\ the series of amplitudes corresponding to the moduli $c$, $c^{0}$, $c^{00}$, $c^{000}$, etc., we shall have first $\sin(2\varphi-\varphi^{0})=c^{0}\sin\varphi^{0}$, an equation which must be solved in order to deduce $\varphi^{0}$ from $\varphi$; from it we draw successively
\[ \begin{aligned} \sin\varphi^{0}&=(1+b)\frac{\sin\varphi\cos\varphi}{\Delta}\\ \cos\varphi^{0}&=\frac{1-(1+b)\sin^{2}\varphi}{\Delta}\\ \Delta(c^{0},\varphi^{0})&=\frac{1-(1-b)\sin^{2}\varphi}{\Delta}\\ \operatorname{tang}\varphi^{0}&=\frac{(1+b)\operatorname{tang}\varphi}{1-b\operatorname{tang}^{2}\varphi}. \end{aligned} \]
This last can be put in the very simple form $\operatorname{tang}(\varphi^{0}-\varphi)=b\operatorname{tang}\varphi$, and from it one would deduce without ambiguity,
\[ \varphi^{0}=2\varphi-c^{0}\sin2\varphi+\tfrac{1}{2}c^{02}\sin4\varphi-\tfrac{1}{3}c^{03}\sin6\varphi+\text{etc.}, \]
which agrees with the remark we made on the always limited value of $2\varphi'-\varphi$, which applies to $2\varphi-\varphi^{0}$, $2\varphi^{0}-\varphi^{00}$, etc.

Let us suppose, then, that after having determined the decreasing moduli $c^{0}$, $c^{00}$, $c^{000}$, etc., as well as their complements $b^{0}$, $b^{00}$, $b^{000}$, etc.\ as has just been said, one calculates successively the amplitudes $\varphi^{0}$, $\varphi^{00}$, $\varphi^{000}$, etc.\ by the formulas
\[ \begin{gathered} \operatorname{tang}(\varphi^{0}-\varphi)=b\operatorname{tang}\varphi\\ \operatorname{tang}(\varphi^{00}-\varphi^{0})=b^{0}\operatorname{tang}\varphi^{0}\\ \operatorname{tang}(\varphi^{000}-\varphi^{00})=b^{00}\operatorname{tang}\varphi^{00}\\ \text{etc.} \end{gathered} \]
Then we shall have an infinite series of functions $F(c,\varphi)$, $F(c^{0},\varphi^{0})$, $F(c^{00},\varphi^{000})$\ednote{Printed $\varphi^{000}$; the series and the next page give $F(c^{00},\varphi^{00})$, so probably a misprint for $\varphi^{00}$.}, etc.\ linked together by constant ratios, in this manner:

\origpage{85}
\[ \begin{gathered} F(c,\varphi)=\frac{1+c^{0}}{2}F(c^{0},\varphi^{0})\\ F(c^{0},\varphi^{0})=\frac{1+c^{00}}{2}F(c^{00},\varphi^{00})\\ F(c^{00},\varphi^{00})=\frac{1+c^{000}}{2}F(c^{000},\varphi^{000})\\ \text{etc.} \end{gathered} \]
When we put $\varphi=\frac{1}{2}\pi$, we have $\varphi^{0}=\pi$, $\varphi^{00}=2\pi$, $\varphi^{000}=4\pi$, etc.; whence it follows that the complete functions have between them these relations
\[ \begin{gathered} F^{1}(c)=(1+c^{0})F^{1}(c^{0})\\ F^{1}(c^{0})=(1+c^{00})F^{1}(c^{00})\\ F^{1}(c^{00})=(1+c^{000})F^{1}(c^{000})\\ \text{etc.} \end{gathered} \]

\section*{Application of the same law to elliptic functions of the second species.}

\subsection*{(61).}

Let us now consider the two functions $E(c,\varphi)=\displaystyle\int d\varphi\,\Delta(c,\varphi)$, $E(c',\varphi')=\displaystyle\int d\varphi'\,\Delta(c',\varphi')$; if in the second formula we substitute for $d\varphi'$ and $\Delta(c',\varphi')$ the values found in art.~58, we shall have
\[ \begin{aligned} (2+c)\,E(c',\varphi')&=\int\frac{d\varphi}{\Delta}(c\cos\varphi+\Delta)^{2}\\ &=2E(c,\varphi)-b^{2}F(c,\varphi)+2c\sin\varphi; \end{aligned} \]
\ednote{The factor reads $(2+c)$, its right part covered by a reader's ink stroke; the display below and the computation require $(2+2c)$. Whether a second $2$ is hidden under the mark or the print omits it cannot be settled from the scan.}
whence one sees that the function of the first species $F(c,\varphi)$ can be expressed by the two arcs of an ellipse $E(c,\varphi)$, $E(c',\varphi')$, since we have
\[ b^{2}F(c,\varphi)=2E(c,\varphi)-(2+2c)E(c',\varphi')+2c\sin\varphi. \]

\subsection*{(62).}

It was found (art.~13) that the expression for the arc of a hyperbola is
\[ \Upsilon=\Delta\operatorname{tang}\varphi-E(c,\varphi)+b^{2}F(c,\varphi); \]
if one substitutes in it the value of $b^{2}F(c,\varphi)$, this expression will become
\[ \Upsilon=\Delta\operatorname{tang}\varphi+E(c,\varphi)-2(1+c)E(c',\varphi')+2c\sin\varphi; \]
whence it follows that an arc of a hyperbola can always be expressed by two

\origpage{86}
arcs of an ellipse, which is the beautiful theorem with which Landen has enriched Geometry${}^{(1)}$.

We have already called $G(\varphi)$ the difference between the arc of a hyperbola $\Upsilon$ and its tangent $\Delta\operatorname{tang}\varphi$, terminated at the perpendicular dropped from the centre; this difference is therefore expressed in general by two arcs of an ellipse, so that we have
\[ G(\varphi)=2(1+c)E(c',\varphi')-E(c,\varphi)-2c\sin\varphi; \]
and in particular when we put $\varphi=\frac{1}{2}\pi$, we have for the expression of the difference between the asymptote and the curve,
\[ G^{1}=2(1+c)E(c',\varphi')-E^{1}(c)-2c. \]
But since $\varphi=\frac{1}{2}\pi$, we have $\sin(2\varphi'-\frac{1}{2}\pi)=c\sin\varphi=c$, or $\cos2\varphi'=-c$, which gives $\sin^{2}\varphi'=\dfrac{1+c}{2}=\dfrac{1}{1+b'}$; therefore the arc $E(c',\varphi')$ is the one which is measured by half of $E^{1}(c')$, and we have $E'(c',\varphi')$\ednote{Printed $E'$ with a prime, where the same arc is $E(c',\varphi')$ in the line above; probably a misprint or a damaged glyph for $E$.}${}=\frac{1}{2}E^{1}(c')+\frac{1}{2}(1-b')$. Substituting this value in that of $G^{1}$, there will result
\[ G^{1}=(1+c)E^{1}(c')-E^{1}(c); \]
therefore the transcendent $G^{1}$ is equal to the difference of the two quadrants of an ellipse which have for semi-axes, the one $1+c$ and $1-c$, the other $1$ and $\sqrt{(1-c^{2})}$.

\subsection*{(63).}

The value of $F$ found in no.~61 is the integral of particular form which satisfies the differential equations $(e')$ of article~43; we shall deduce from it new properties of the functions $E$.

Let us consider an infinite series of functions $E(c,\varphi)$, $E(c',\varphi')$, $E(c'',\varphi'')$, etc.\ formed according to the same law as the functions of the first species $F(c,\varphi)$, $F(c',\varphi')$, $F(c'',\varphi'')$, etc.; we shall have first the two equations
\[ \begin{aligned} &\tfrac{1}{2}b^{2}F(c,\varphi)=E(c,\varphi)-(1+c)E(c',\varphi')+c\sin\varphi\\ &\tfrac{1}{2}b'^{2}F(c',\varphi')=E(c',\varphi')-(1+c')E(c'',\varphi'')+c'\sin\varphi'; \end{aligned} \]

\textbf{(1)} One would easily deduce from the same formulas that every arc of an ellipse can be expressed by two arcs of a hyperbola.

\origpage{87}
but we have moreover $F(c',\varphi')=\dfrac{1+c}{2}F(c,\varphi)$; eliminating therefore from these three equations the functions of the first species, we shall have between the three consecutive functions of the second species $E(c,\varphi)$, $E(c',\varphi')$, $E(c'',\varphi'')$, this equation
\[ \begin{aligned} 0={}&\tfrac{1}{2}b'(1+b')E(c,\varphi)-(2+b')E(c',\varphi')+2(1+c')E(c'',\varphi'')\\ &+\tfrac{1}{2}b'(1-b')\sin\varphi-2c'\sin\varphi', \end{aligned} \]
an equation which can be considered as a sort of particular integral of the differential equation of the second order $0=(1+c^{2})\dfrac{ddE}{dc^{2}}+\text{etc.}$ given in art.~45.

The same equation can be applied to three consecutive ellipses, taken not only in the infinite series $E(c,\varphi)$, $E(c',\varphi')$, $E(c'',\varphi'')$, $E(c''',\varphi''')$, etc., but in general in the doubly infinite series
\[ \ldots E(c'',\varphi''),\ E(c',\varphi'),\ E(c,\varphi),\ E(c^{0},\varphi^{0}),\ E(c^{00},\varphi^{00}),\ldots \]
whose extremes are, on the one hand, the ellipse which has $1$ for its eccentricity and which reduces to its major axis; on the other hand, the ellipse which has $0$ for its eccentricity and which merges into the circle: it follows therefore that by the indefinite rectification of two ellipses of this series, one obtains the indefinite rectification of all the others.

The general formula is simplified when it is a question of the definite rectification of these ellipses. Indeed if we put $\varphi''=\frac{1}{2}\pi$, we shall have $\varphi'=\pi$ and $\varphi=2\pi$, which will give $E(c'',\varphi'')=E^{1}(c'')$, $E(c',\varphi')=2E^{1}(c')$, $E(c,\varphi)=4E^{1}(c)$; therefore we shall have between the three quadrants of an ellipse $E^{1}(c)$, $E^{1}(c')$, $E^{1}(c'')$, this equation
\[ 0=b'(1+b')E^{1}(c)-(2+b')E^{1}(c')+(1+c')E^{1}(c''). \]
We shall have a similar equation between any three consecutive terms taken in the general series of ellipses. Thus the circumference of a proposed ellipse can always be determined exactly by the circumferences of two ellipses, differing as little from the straight line as one pleases, by prolonging the series in one direction, or differing as little from the circle as one pleases, by prolonging the series in the other direction.

\origpage{88}

In this latter case, one would have to make use of the successive equations
\[ \begin{aligned} &0=(1+c^{0})E^{1}(c)-(2+b^{0})E^{1}(c^{0})+b^{0}(1+b^{0})E^{1}(c^{00})\\ &0=(1+c^{00})E^{1}(c^{0})-(2+b^{00})E^{1}(c^{00})+b^{00}(1+b^{00})E^{1}(c^{000})\\ &\text{etc.}, \end{aligned} \]
which one would prolong as far as one had thought fit to prolong the series of moduli $c$, $c^{0}$, $c^{00}$, etc.

The exact and absolute determination of which we have just spoken may be regarded as a very remarkable theorem in the theory of transcendents; but if one's only aim is to obtain approximations, one will arrive at them more easily by the method which we shall give hereafter.

\subsection*{(64).}

We found (art.~41) that the two functions $E^{1}(c)$, $F^{1}(c)$, both for the modulus $c=\sin15^{\circ}$ and for its complement $c=\sin75^{\circ}$, can be determined by any one of the four functions supposed known. Therefore the two series of ellipses formed, the one according to the modulus $c=\sin15^{\circ}=\sqrt{\left(\dfrac{2-\sqrt{3}}{4}\right)}$, the other according to the modulus $c=\sin75^{\circ}=\sqrt{\left(\dfrac{2+\sqrt{3}}{4}\right)}$, are such that, the circumference of a single one of these ellipses being known, one can find the circumference of all the others. It is the same with the two sequences of functions of the first species $F^{1}(c)$ formed according to the same moduli, and a single known term in these four series will suffice to make all the others known.

We likewise found (art.~42) that when $c=\sin45^{\circ}=\sqrt{\tfrac{1}{2}}$, the functions $F^{1}(c)$, $E^{1}(c)$ can be determined one by the other; but $F^{1}(c)$ is determined generally by the two quantities $E^{1}(c)$, $E^{1}(c')$, or by the two $E^{1}(c)$, $E^{1}(c^{0})$, for from our formulas one easily finds
\[ bF^{1}(c)=-E^{1}(c)+(1+b)E^{1}(c^{0}). \]
Therefore in general all the ellipses which make up the series formed according to the modulus $c=\sin45^{\circ}$ are such that, the circumference of one of them being known, one can determine that of all the others.

\origpage{89}

To these three general cases one can join a fourth. Suppose that we have $b=c^{0}=\dfrac{1-b}{1+b}$, which gives $b=-1+\sqrt{2}$; and $c^{2}=-2+2\sqrt{2}$, we shall have $F^{1}(b)=F^{1}(c^{0})=\dfrac{1}{1+c^{0}}F^{1}(c)=\dfrac{1}{\sqrt{2}}F^{1}(c)$, $E^{1}(b)=E^{1}(c^{0})$. Substituting these values in formula $(d')$, we shall have
\[ \frac{\pi}{\sqrt{2}}=F^{1}(c)\,[E^{1}(c)+\sqrt{2}E^{1}(c^{0})-F^{1}(c)]. \]
But we have besides $bF^{1}(c)=(1+b)E^{1}(c^{0})-E^{1}(c)$; thus one can determine $E^{1}(c)$ by $E^{1}(c^{0})$, and in general in the series of ellipses formed according to the modulus $c=\sqrt{(2\sqrt{2}-2)}$, the circumference of a single ellipse being known, one can determine that of all the others.

\section*{Method of approximation applied to elliptic functions of the first species.}

\subsection*{(65).}

The function $F(c,\varphi)$ being proposed, of which one wishes to have an approximate value, one will calculate the decreasing moduli $c^{0}$, $c^{00}$, $c^{000}$, etc.\ and the increasing amplitudes $\varphi^{0}$, $\varphi^{00}$, $\varphi^{000}$, etc.\ by the formulas of art.~60; one will thus have successively
\[ \begin{aligned} F(c,\varphi)&=\frac{1+c^{0}}{2}F(c^{0},\varphi^{0})\\ &=\frac{1+c^{0}}{2}\cdot\frac{1+c^{00}}{2}F(c^{00},\varphi^{00})\\ &=\frac{1+c^{0}}{2}\cdot\frac{1+c^{00}}{2}\cdot\frac{1+c^{000}}{2}F(c^{000},\varphi^{000})\\ &\quad\text{etc.} \end{aligned} \]
But when $c$ has become very small, we have $\Delta=1$ and $\displaystyle\int\frac{d\varphi}{\Delta}=\varphi$. Let then $\Phi$ be the limit of the angles $\frac{1}{2}\varphi^{0}$, $\frac{1}{4}\varphi^{00}$, $\frac{1}{8}\varphi^{000}$, etc., a limit which they will always sensibly attain after a certain number of terms, and we shall have the function sought
\[ F(c,\varphi)=\Phi(1+c^{0})(1+c^{00})(1+c^{000}),\ \text{etc.} \]
When $\varphi=\frac{1}{2}\pi$, the limit $\Phi$ will likewise be $\frac{1}{2}\pi$, so

\origpage{90}
that we shall have
\[ F^{1}(c)=\frac{\pi}{2}(1+c^{0})(1+c^{00})(1+c^{000}),\ \text{etc.} \]

The constant product $(1+c^{0})(1+c^{00})(1+c^{000})$, etc.\ which we shall represent by $K$, can also be expressed in this manner:
\[ K=\frac{2\sqrt{c^{0}}}{c}\cdot\frac{2\sqrt{c^{00}}}{c^{0}}\cdot\frac{2\sqrt{c^{000}}}{c^{00}},\ \text{etc.}; \]
and in this form it is easy to calculate by logarithms. $K$ being known, we shall have $F(c,\varphi)=K\Phi$ and $F^{1}(c)=K\cdot\dfrac{\pi}{2}$.

To facilitate the calculation of the decreasing moduli, one will determine an auxiliary angle $\mu$ by the equation $\sin\mu=c$, which will give $b=\cos\mu$ and $c^{0}=\operatorname{tang}^{2}\frac{1}{2}\mu$. Putting likewise $c^{0}=\operatorname{tang}^{2}\frac{1}{2}\mu=\sin\mu^{0}$, which will determine a new angle $\mu^{0}$, we shall have $c^{00}=\operatorname{tang}^{2}\frac{1}{2}\mu^{00}$\ednote{Printed $\mu^{00}$; by the rule just stated the angle is $\mu^{0}$, so probably a misprint.}, and so on.

When one has arrived at a very small $c$, one may, to avoid angles that are too small, calculate the following term $c^{0}$ by the formula
\[ c^{0}=\frac{1}{4}c^{2}+\frac{1\cdot3}{4\cdot6}c^{4}+\frac{1\cdot3\cdot5}{4\cdot6\cdot8}c^{6}+\text{etc.}, \]
of which the first term, or at most the first two terms, will suffice.

As for the calculation of the angles $\varphi^{0}$, $\varphi^{00}$, etc., we have nothing to add to the simplicity of the formula $\operatorname{tang}(\varphi^{0}-\varphi)=b\operatorname{tang}\varphi$, which is very well suited to trigonometric calculation. We shall only recall what was said in art.~60, that the angle $\varphi^{0}-\varphi$ is almost equal to $\varphi$ when $c$ is very small, and that in general its difference from $\varphi$ is less than the angle whose sine is $c$. One must therefore take for the angle $\varphi^{0}-\varphi$, not always the smallest angle that the tables of sines give, but the one which comes close to $\varphi$, and which may comprise several circumferences.

\section*{EXAMPLE.}

\subsection*{(66).}

The value of the function $F(c,\varphi)$ is required when $c=\frac{1}{2}\sqrt{(2+\sqrt{3})}=\sin75^{\circ}$ and $\operatorname{tang}\varphi=\sqrt{\left(\dfrac{2}{\sqrt{3}}\right)}$.

\origpage{91}

Here first is a table showing the calculation of the moduli and their complements:
\[ \begin{aligned} &\textit{Values of $c$ and $b$.} && \textit{Their logarithms.}\\ c&=\sin\ 75^{\circ}\ 0'\ 0''00\ldots\ldots && 9.9849438\\ b&=\cos\ 75.\ 0.\ 0.00\ldots\ldots && 9.4129962\\ c^{0}&=\left\{\begin{array}{l}\operatorname{tang}^{2}37.30.\ 0.00\\ \sin\ 36.\ 4.16.47\end{array}\right\}\ldots\ldots && 9.7699610\\ b^{0}&=\cos\ 36.\ 4.16.47\ldots\ldots && 9.9075648\\ c^{00}&=\left\{\begin{array}{l}\operatorname{tang}^{2}18.\ 2.\ 8.235\\ \sin\ 6.\ 5.\ 9.38\end{array}\right\}\ldots\ldots && 9.0253880\\ b^{00}&=\cos\ 6.\ 5.\ 9.38\ldots\ldots && 9.9975452\\ c^{000}&=\left\{\begin{array}{l}\operatorname{tang}^{2}\ 3.\ 2.34.69\\ \sin\ 0.\ 9.42.90\end{array}\right\}\ldots\ldots && 7.4511672\\ b^{000}&=\cos\ 0.\ 9.42.90\ldots\ldots && 9.9999982\\ c^{0000}&=\tfrac{1}{4}(c^{000})^{2}+\text{etc.}\ldots\ldots && 4.3002761\\ b^{0000}&\ldots\ldots && 0.0000000 \end{aligned} \]
These logarithms will readily give the value of $K$, for which it will suffice to use four factors. We shall thus have $\log K=0.2460561$ and $K=1.7622037$. From this there results first $F^{1}=\dfrac{\pi}{2}K=2.7680631$; we shall then find, by the calculation of the amplitudes,
\[ \begin{aligned} \varphi&=47^{\circ}\ 3'\ 30''94\\ \varphi^{0}&=62.36.\ 3.10\\ \varphi^{00}&=119.55.47.67\\ \varphi^{000}&=240.\ 0.\ 0.19\\ \varphi^{0000}&=480.\ 0.\ 0.00. \end{aligned} \]
The other values of $\varphi$ would increase in double ratio; whence it follows that the limit of the angles $\varphi$, $\dfrac{\varphi^{0}}{2}$, $\dfrac{\varphi^{00}}{4}$, etc. is $30^{\circ}$ or $\dfrac{\pi}{6}$; therefore we have $F(c,\varphi)=K\dfrac{\pi}{6}=0.9226877$, which agrees with the value found in art.~28.

Let us remark that, since the limit $\Phi=30^{\circ}=\dfrac{1}{3}\cdot\dfrac{\pi}{2}$, we have $F=\frac{1}{3}F^{1}$; and indeed this equality holds rigorously by virtue of the value we have taken for $\operatorname{tang}\varphi$. See article~24.

\origpage{92}
\subsection*{(67).}

The value of the amplitude $\varphi$ being given, one sees that it is easy to find the function $F$ with all the necessary exactness; conversely, it may be useful to determine the amplitude when the value of the function $F$ is supposed known. For this purpose, one must calculate the terms of the series $c^{0}$, $c^{00}$, etc. as far as a term small enough to be neglected. Let this term be, for example, $c^{00000}$, or for brevity $c^{05}$; since we have $\varphi^{05}=2\varphi^{04}-c^{05}\sin^{2}\varphi^{04}+\text{etc.}$\ednote{translator's note: Printed with a raised $2$ after sin; the series of art.\ 60 gives $\sin 2\varphi^{04}$, so the raised figure is probably a misplaced $2$.}, we shall have, with sufficient exactness, $\varphi^{05}=2\varphi^{04}$, and all the more $\varphi^{06}=2\varphi^{05}$, etc. Therefore the limit of the angles $\varphi$, $\frac{1}{2}\varphi^{0}$, $\frac{1}{4}\varphi^{00}$, etc. will be $\frac{1}{16}\varphi^{04}$: this limit has been designated by $\Phi$, so that we shall have $\varphi^{04}=16\Phi$. Moreover, the value of $F$ being given, we know $\Phi$ by the equation $\Phi=\dfrac{F}{K}$; therefore we shall also know $\varphi^{04}=16\Phi$. This done, we shall calculate successively the values of $\varphi^{000}$, $\varphi^{00}$, $\varphi^{0}$, $\varphi$, by means of the equations
\[ \begin{aligned} \sin(2\varphi^{03}-\varphi^{04})&=c^{04}\sin\varphi^{04}\\ \sin(2\varphi^{02}-\varphi^{03})&=c^{03}\sin\varphi^{03}\\ \sin(2\varphi^{0}-\varphi^{02})&=c^{02}\sin\varphi^{02}\\ \sin(2\varphi-\varphi^{0})&=c^{0}\sin\varphi^{0}, \end{aligned} \]
and we shall have the required amplitude $\varphi$.

This method applies in particular to the solution of the equation $F(\psi)=nF(\varphi)$, whatever $n$ may be. $\varphi$ being given, we shall know $F(\varphi)$ and $nF(\varphi)$, or $F(\psi)$; then from $F(\psi)$ we shall deduce the amplitude $\psi$, as has just been explained. This means will never require more than a small number of operations, unless $1-c$ is excessively small; whereas if $n$ were somewhat large or merely fractional, the algebraic methods we have given for determining $\varphi_{n}$ from the equation $F(\varphi_{n})=nF(\varphi)$ would become very long or even impracticable.

\subsection*{(68).}

The method we have just set out is in general very expeditious; it will require only that a few more terms be calculated, both in the sequence of moduli and in that of amplitudes, when $c$ is extremely close to unity; but one can be assured that this number of terms will never be very considerable, for in the hypothesis where one would have to continue the sequence $c$, $c^{0}$, $c^{00}$, etc. as far as the tenth term in order that this term should be of a decimal unit of the tenth order, the primitive value of $b$ would have to be smaller than $10^{-21}$, and consequently that of $1-c$ smaller than $10^{-42}$.

\origpage{93}
The universality of the method is sufficiently established by this observation; nevertheless, when $1-c$ is extremely small, one may take advantage of this circumstance to simplify the calculations, and proceed to the approximations in another manner.

In the case in question, the quantity $b$ is very small, and since we have $F=\displaystyle\int\frac{d\varphi}{\sqrt{(\cos^{2}\varphi+b^{2}\sin^{2}\varphi)}}$, if the amplitude $\varphi$ is such that $\operatorname{tang}\varphi$ is much smaller than $\dfrac{1}{b}$, we shall have, with sufficient approximation,
\[ F=\int\frac{d\varphi}{\cos\varphi}=\log\operatorname{tang}(45^{\circ}+\tfrac{1}{2}\varphi). \]

If $\operatorname{tang}\varphi$ is greater than $\dfrac{1}{b}$, it will be necessary to transform the proposed formula $F(c,\varphi)$ into another in which $b$ is much smaller, so that in the transformed formula $b\operatorname{tang}\varphi$ becomes a negligible quantity. This is easily obtained by calculating, as far as the suitable term, the increasing moduli $c'$, $c''$, etc. and the amplitudes $\varphi'$, $\varphi''$, etc., by the formulas of article~59.

For this purpose let $b=\sin\lambda$; we shall have $b'=\operatorname{tang}^{2}\frac{1}{2}\lambda$; let again $b'=\operatorname{tang}^{2}\frac{1}{2}\lambda=\sin\lambda'$; we shall have $b''=\operatorname{tang}^{2}\frac{1}{2}\lambda'$, and so on as far as a term $b^{\mu}$ having the required degree of smallness; we shall then calculate the amplitudes $\varphi$, $\varphi'$, $\varphi''$, $\varphi'''$, etc., as far as a term $\varphi^{\mu}$ corresponding to the last value of $b$; and this term being named $\Phi'$, if we put
\[ K'=\frac{2}{1+c}\cdot\frac{2}{1+c'}\cdot\frac{2}{1+c''}\,\text{etc.}=\sqrt{\left(\frac{c'c''c'''\ldots\text{etc.}}{c}\right)}, \]
we shall have
\[ F(c,\varphi)=K'\log\operatorname{tang}(45^{\circ}+\tfrac{1}{2}\Phi'). \]
This value will be exact if the factors of which $K'$ is composed, and the sequence of which $\Phi''$\ednote{translator's note: Printed $\Phi''$; the term was named $\Phi'$ just above, so probably a misprint for $\Phi'$.} is the last term, are prolonged to infinity; but after a rather small number of terms one will in general obtain every degree of approximation one may desire. One would have at the same time

\origpage{94}
for the complete function $F^{1}(c)$, either the value given by the preceding formula by calculating the angle $\Phi'$ by means of $\varphi=90^{\circ}$, or the value $F^{1}(c)=K'\cdot\dfrac{1}{2^{\mu}}\log\dfrac{4}{b^{\mu}}$.

\section*{EXAMPLE.}

\subsection*{(69).}

Let, as above, $c=\sin75^{\circ}$, $\operatorname{tang}\varphi=\sqrt{\left(\dfrac{2}{\sqrt{3}}\right)}$; this case is little favorable to the application of the preceding method, because $b$ is not a very small quantity.

We shall first calculate the increasing moduli $c'$, $c''$, $c'''$, etc., and their complements $b'$, $b''$, $b'''$, etc., as follows:
\[ \begin{aligned} &\textit{Values of $b$ and $c$.} && \textit{Their logarithms.}\\ b&=\sin\ 15^{\circ}\ 0'\ 0''00\ldots\ldots && 9.4129962\\ c&=\cos\ 15.\ 0.\ 0.00\ldots\ldots && 9.9849438\\ b'&=\left\{\begin{array}{l}\operatorname{tang}^{2}\ 7.30.\ 0.00\\ \sin\ 0.59.35.24\end{array}\right\}\ldots\ldots && 8.2388582\\ c'&\ldots\ldots && 9.9999348\\ b''&=\left\{\begin{array}{l}\operatorname{tang}^{2}\ 0.29.47.62\\ \sin\ 0.\ 0.\text{etc.}\end{array}\right\}\ldots\ldots && 3.8757219\\ c''&\ldots\ldots && 0.0000000\\ b'''&=(\tfrac{1}{2}b'')^{2}\ldots\ldots && 1.1493838\\ c'''&\ldots\ldots && 0.0000000. \end{aligned} \]
\ednote{translator's note: The logarithm of $b''$ is printed $3.8757219$; from $\log b'=8.2388582$ it should be $5.8757219$ (as the $b'''$ row, $1.1493838$, presupposes), so the first figure is a misprint.}
Then the calculation of the amplitudes $\varphi'$, $\varphi''$, etc. gives the following results:
\[ \begin{aligned} \varphi&=47^{\circ}\ 3'\ 30''95,\quad 2\varphi'-\varphi=45^{\circ}\\ \varphi'&=46.\ 1.45.475\\ \varphi''&=46.\ 1.29.41\\ \varphi'''&=46.\ 1.29.41. \end{aligned} \]
The value of the factor $K'$ reduces in this example to $\sqrt{\dfrac{c'}{c}}$, so we have $\log K'=0.0074955$; and since we have found $\Phi'=46^{\circ}1'29''41$, we shall have
\[ F(\varphi)=K'\log\operatorname{tang}68^{\circ}\ 0'\ 44''705. \]

\origpage{95}
This log-tangent, taken from the tables and then multiplied by the modulus in order to make of it a hyperbolic logarithm, will give
\[ \log F=9{,}9650547, \]
or $F=0.9226877$, as has already been found in no.~66.

If in the same example one wishes to have the value of the complete function, one must put $\varphi=90^{\circ}$, and calculate successively the values of $\varphi'$, $\varphi''$, etc., which will give
\[ \begin{aligned} \varphi'&=82^{\circ}\ 30'\ 0''\ 00\\ \varphi''&=82.28.\ 2.84\\ \varphi'''&=\varphi''. \end{aligned} \]
Therefore the limit $\Phi'=82^{\circ}28'2''84$, and thus the complete function
\[ F^{1}=K'\log\operatorname{tang}86^{\circ}\ 14'\ 1''42, \]
or $\log F^{1}=0.4421761$, which agrees with the value found in no.~66.

We shall remark that in this example it was necessary to have recourse to a particular means in order to determine, with the suitable precision, the amplitude $\varphi''$, which is sensibly indistinguishable from the limit $\Phi'$. The equation $\sin(2\varphi''-\varphi')=c'\sin\varphi'$, the most direct for determining $\varphi''$, is not suited to giving quite exactly the fractions of a second contained in $\varphi''$, because these fractions have very little influence on the sine of an angle of $82^{\circ}$, too close to the right angle. In this case, and in all similar ones that may be met with, one will determine $\varphi''$ much more exactly by means of the equation $\operatorname{tang}(\varphi'-\varphi'')=b''\operatorname{tang}\varphi''$, or $\varphi'-\varphi''=Rb''\operatorname{tang}\varphi''$, $R$ being the number of seconds contained in the radius: for this, one must put in the second member $\varphi'$ in place of $\varphi''$, which will give a first approximate value of $\varphi'-\varphi''$, and consequently one of $\varphi''$. Substituting this value anew in place of $\varphi''$ in the second member, we shall have a second value of $\varphi'-\varphi''$, which will be approximate at least to hundredths of a second.

In the example in question, we shall therefore first put $\varphi'-\varphi''=Rb''\operatorname{tang}82^{\circ}30'$, which will give $\varphi'-\varphi''=1'57''68$ and $\varphi''=82^{\circ}28'2''32$. Substituting this value anew in place of $\varphi''$, we shall have more exactly $\varphi'-\varphi''=Rb''\operatorname{tang}82^{\circ}\,28'\,2''.32=1'\,57''\,167$, whence $\varphi''=82^{\circ}\,28'\,2''\,833$.

\origpage{96}
\subsection*{(70).}

One now sees that there are two methods for determining the approximate value of a function of the first species $F(c,\varphi)$ whose modulus and amplitude are known. If $c$ is smaller than $\sqrt{\frac{1}{2}}$ or $\sin45^{\circ}$, it will be proper to use the method of no.~65, according to which, calculating the decreasing moduli $c^{0}$, $c^{00}$, $c^{000}$, etc. and the increasing amplitudes $\varphi^{0}$, $\varphi^{00}$, $\varphi^{000}$, etc., one obtains the formula $F(c,\varphi)=K\Phi$.

If the modulus $c$ is greater than $\sin45^{\circ}$, it will be proper to use the method of no.~68, which consists in calculating the increasing moduli $c'$, $c''$, $c'''$, etc., and the decreasing amplitudes $\varphi'$, $\varphi''$, $\varphi'''$, etc., from which one concludes $F(c,\varphi)=K'\log\operatorname{tang}(45^{\circ}+\frac{1}{2}\Phi')$.

These two methods can each be extended at will beyond the limit we have fixed. But in order to diminish as much as possible the number of transformed functions, it is well to keep to this limit, and in this way one will never need to calculate more than three terms, both of the series of moduli and of that of amplitudes, in order to obtain a result approximate to the seventh place of decimals.

It is very remarkable that the function $F$ can always be expressed at will by an arc of a circle or by a logarithm; but one sees that it comes nearer to arcs of a circle if $c^{2}<\frac{1}{2}$, and that it has more affinity with logarithms if $c^{2}>\frac{1}{2}$.

\subsection*{(71).}

We have shown in article 67 how one determines the amplitude $\varphi$ corresponding to a given value of the function $F(c,\varphi)$. The method set out in that article has all the perfection one could desire when $c^{2}$ is $<\frac{1}{2}$, and it can be applied with success when $c^{2}$ comes much nearer to unity. Nevertheless, if one wishes the number of transformed functions to be as small as possible, one must, when $c^{2}$ is $>\frac{1}{2}$, apply the method of approximation of article 68.

For this one will begin by calculating the sequence $c'$, $c''$, $c'''$, etc. as far as a term that does not differ sensibly from unity. This limit is $c'''$ or even $c''$, when one does not wish to push the exactness beyond six or seven decimals. From this one will derive $K'=\sqrt{\left(\dfrac{c'c''c'''\,\text{etc.}}{c}\right)}$; $K'$ being known, we shall know $\log\operatorname{tang}(45^{\circ}+\frac{1}{2}\Phi')=\dfrac{F}{K'}$. From this

\origpage{97}
hyperbolic logarithm, or from the common logarithm corresponding to it, one will derive the value of the angle $\Phi'$, the limit of the angles $\varphi'$, $\varphi''$, $\varphi'''$, etc. If one has stopped at $c'''$, neglecting the difference $1-c'''$, then $\varphi'''$ may be taken for the limit $\Phi'$. Knowing $\varphi'''$, one will go back successively to the values of $\varphi''$, $\varphi'$, $\varphi$ by the equations
\[ \begin{aligned} \operatorname{tang}(\varphi''-\varphi''')&=b'''\operatorname{tang}\varphi'''\\ \operatorname{tang}(\varphi'-\varphi'')&=b''\operatorname{tang}\varphi''\\ \operatorname{tang}(\varphi-\varphi')&=b'\operatorname{tang}\varphi'; \end{aligned} \]
we shall therefore know the required amplitude $\varphi$.

By these methods one can determine the function $F$ when its amplitude is known, or conversely, determine the amplitude when the function is known; so that it will suffice to calculate four to five terms at most of the series of moduli, and an equal number of terms of the series of amplitudes, in order to have ten exact decimals, if the tables one uses are to ten decimals. If they had twenty, it would suffice to calculate one more term in the series mentioned to have results exact to the twentieth decimal, and so on.

\section*{Particular properties of the functions $F(c)$, $F(b)$, whose moduli are complements of one another.}

\subsection*{(72).}

We have seen that, by prolonging the indefinite sequence of moduli derived from a given modulus $c$, both in the sense in which they converge to the limit 0 and in that in which they converge to the limit 1, one has, for the indefinite functions and for the definite functions, these two series of equations
\[ \begin{aligned} F(c,\varphi)&=\frac{1+c^{0}}{2}F(c^{0},\varphi^{0}) &\qquad& F^{1}(c)=(1+c^{0})F^{1}(c^{0})\\ F(c,\varphi)&=\frac{1+c^{0}}{2}\cdot\frac{1+c^{00}}{2}F(c^{00},\varphi^{00}) && F^{1}(c)=(1+c^{0})(1+c^{00})F^{1}(c^{00})\\ F(c,\varphi)&=\frac{1+c^{0}}{2}\cdot\frac{1+c^{00}}{2}\cdot\frac{1+c^{000}}{2}F(c^{000},\varphi^{000}) && F^{1}(c)=(1+c^{0})(1+c^{00})(1+c^{000})F^{1}(c^{000})\\ &\text{etc.} && \text{etc.} \end{aligned} \]

\origpage{98}
\[ \begin{aligned} F(c,\varphi)&=\frac{2}{1+c}F(c',\varphi') &\qquad& F^{1}(c)=\frac{1}{1+c}F^{1}(c')\\ F(c,\varphi)&=\frac{2}{1+c}\cdot\frac{2}{1+c'}F(c'',\varphi'') && F^{1}(c)=\frac{1}{1+c}\cdot\frac{1}{1+c'}F^{1}(c'')\\ F(c,\varphi)&=\frac{2}{1+c}\cdot\frac{2}{1+c'}\cdot\frac{2}{1+c''}F(c''',\varphi''') && F^{1}(c)=\frac{1}{1+c}\cdot\frac{1}{1+c'}\cdot\frac{1}{1+c''}F^{1}(c''')\\ &\text{etc.} && \text{etc.} \end{aligned} \]
Let us put in these last formulas $C$, $B$, $C'$, $B'$, etc. in place of $c$, $b$, $c'$, $b'$; we shall first have
\[ F(C,\varphi)=\frac{2}{1+C}F(C',\varphi'),\quad F^{1}(C)=\frac{1}{1+C}F^{1}(C'). \]
Let us then suppose $C=b$, and consequently $B=c$; we shall have $C'=\dfrac{2\sqrt{C}}{1+C}=\dfrac{2\sqrt{b}}{1+b}=b^{0}$, consequently $B'=c^{0}$; likewise we should have in the following transformed functions $C''=b^{00}$, $B''=c^{00}$, $C'''=b^{000}$, $B'''=c^{000}$, etc. But we have at the same time the equation $c^{0}=\dfrac{1-b}{1+b}$, which gives $\dfrac{1}{1+C}=\dfrac{1}{1+b}=\dfrac{1+c^{0}}{2}$, $\dfrac{1}{1+C'}=\dfrac{1}{1+b^{0}}=\dfrac{1+c^{00}}{2}$, etc; therefore the second series of equations transferred to the complementary modulus will give
\[ \begin{aligned} F(b,\varphi)&=(1+c^{0})F(b^{0},\varphi') &\qquad& F^{1}(b)=\frac{1+c^{0}}{2}F^{1}(b^{0})\\ F(b,\varphi)&=(1+c^{0})(1+c^{00})F(b^{00},\varphi'') && F^{1}(b)=\frac{1+c^{0}}{2}\cdot\frac{1+c^{00}}{2}F^{1}(b^{00})\\ F(b,\varphi)&=(1+c^{0})(1+c^{00})(1+c^{000})F(b^{000},\varphi''') && F^{1}(b)=\frac{1+c^{0}}{2}\cdot\frac{1+c^{00}}{2}\cdot\frac{1+c^{000}}{2}F^{1}(b^{000})\\ &\text{etc.} && \text{etc.} \end{aligned} \]
Now as $b^{i}$ becomes greater, we have more and more exactly $F(b^{i},\varphi^{i})=\log\operatorname{tang}(45^{\circ}+\frac{1}{2}\varphi^{i})$ and $F^{1}(b^{i})=\log\left(\dfrac{4}{c^{i}}\right)$; therefore we shall have, by an ever-increasing approximation,
\[ \begin{aligned} F(b,\varphi)&=(1+c^{0})\log\operatorname{tang}(45^{\circ}+\tfrac{1}{2}\varphi') &\qquad& F^{1}(b)=\frac{1+c^{0}}{2}\log\frac{4}{c^{0}}\\ F(b,\varphi)&=(1+c^{0})(1+c^{00})\log\operatorname{tang}(45^{\circ}+\tfrac{1}{2}\varphi'') && F^{1}(b)=\frac{1+c^{0}}{2}\cdot\frac{1+c^{00}}{2}\log\frac{4}{c^{00}}\\ F(b,\varphi)&=(1+c^{0})(1+c^{00})(1+c^{000})\log\operatorname{tang}(45^{\circ}+\tfrac{1}{2}\varphi''') && F^{1}(b)=\frac{1+c^{0}}{2}\cdot\frac{1+c^{00}}{2}\cdot\frac{1+c^{000}}{2}\log\frac{4}{c^{000}}\\ &\text{etc.} && \text{etc.} \end{aligned} \]

\origpage{99}
The first series of equations similarly gives, for the modulus $c$, the formulas of approximation
\[ \begin{aligned} F(c,\varphi)&=\frac{1+c^{0}}{2}\varphi^{0} &\qquad& F^{1}(c)=(1+c^{0})\frac{\pi}{2}\\ F(c,\varphi)&=\frac{1+c^{0}}{2}\cdot\frac{1+c^{00}}{2}\varphi^{00} && F^{1}(c)=(1+c^{0})(1+c^{00})\frac{\pi}{2}\\ F(c,\varphi)&=\frac{1+c^{0}}{2}\cdot\frac{1+c^{00}}{2}\cdot\frac{1+c^{000}}{2}\varphi^{000} && F^{1}(c)=(1+c^{0})(1+c^{00})(1+c^{000})\frac{\pi}{2}\\ &\text{etc.} && \text{etc.} \end{aligned} \]
Various consequences are to be drawn from these formulas.

\subsection*{(73).}

In the case of $b=c=\sqrt{\frac{1}{2}}$, where we have $c'=b^{0}$, $b'=c^{0}$, $c''=b^{00}$, $b''=c^{00}$, etc., these formulas offer for the same function $F(c,\varphi)$ two series of values which, being compared term by term, give
\[ \begin{aligned} F(c^{0},\varphi^{0})&=2F(c',\varphi') &\qquad& F^{1}(c^{0})=\tfrac{1}{2}F^{1}(c')\\ F(c^{00},\varphi^{00})&=4F(c'',\varphi'') && F^{1}(c^{00})=\tfrac{1}{4}F^{1}(c'')\\ F(c^{000},\varphi^{000})&=8F(c''',\varphi''') && F^{1}(c^{000})=\tfrac{1}{8}F^{1}(c''')\\ &\text{etc.} && \text{etc.} \end{aligned} \]
These formulas hold rigorously: the following hold only approximately, but the approximation is ever increasing:
\[ \begin{aligned} \tfrac{1}{2}\varphi^{0}&=\log\operatorname{tang}(45^{\circ}+\tfrac{1}{2}\varphi') &\qquad& \frac{\pi}{2}=\tfrac{1}{2}\log\frac{4}{c^{0}}\\ \tfrac{1}{4}\varphi^{00}&=\log\operatorname{tang}(45^{\circ}+\tfrac{1}{2}\varphi'') && \frac{\pi}{2}=\tfrac{1}{4}\log\frac{4}{c^{00}}\\ \tfrac{1}{8}\varphi^{000}&=\log\operatorname{tang}(45^{\circ}+\tfrac{1}{2}\varphi''') && \frac{\pi}{2}=\tfrac{1}{8}\log\frac{4}{c^{000}}\\ &\text{etc.} && \text{etc.} \end{aligned} \]
It results from this that if the series $\varphi$, $\dfrac{\varphi^{0}}{2}$, $\dfrac{\varphi^{00}}{4}$, $\dfrac{\varphi^{000}}{8}$, etc. has the limit $\Phi$, and if in the contrary sense the series $\varphi$, $\varphi'$, $\varphi''$, $\varphi'''$, etc. has the limit $\Phi'$, we shall have exactly (in the hypothesis of the modulus $c=\sin45^{\circ}$),
\[ \Phi=\log\operatorname{tang}(45^{\circ}+\tfrac{1}{2}\Phi'). \]

\origpage{100}
If we make in particular $\varphi=\frac{1}{2}\pi$, we shall find successively
\[ \begin{aligned} \varphi'&=67^{\circ}30'0''00\\ \varphi''&=66.30.54.36\\ \varphi'''&=66.30.47.74\\ \varphi^{\text{iv}}&=\varphi'''. \end{aligned} \]
Thus the limit $\Phi'$, as far as it can be determined with a table to seven decimals, is $66^{\circ}30'47''74$, and there results from it
\[ \frac{\pi}{2}=\log\operatorname{tang}78^{\circ}15'23''87; \]
which is indeed a very close value of $\dfrac{\pi}{2}$. One could increase the degree of approximation indefinitely by calculating the limit $\Phi'$ with more extensive tables.

But from the other series of equations there results a still easier means of having the logarithmic value of the number $\pi$, since one sees that this number is equal to the limit toward which the successive terms
\[ \log\frac{4}{c^{0}},\quad \tfrac{1}{2}\log\frac{4}{c^{00}},\quad \tfrac{1}{4}\log\frac{4}{c^{000}},\quad \text{etc.} \]
rapidly converge. To learn to what point each term of this sequence approaches the true value of $\pi$, I designate by $c^{n}$ the $n^{\text{th}}$ term of the sequence $c^{0}$, $c^{00}$, etc., and I call $x$ the approximate value of $\pi$ given by this term, so that we shall have $x=\dfrac{1}{2^{n-1}}\log\dfrac{4}{c^{n}}$. I then suppose that for the term $c^{n+1}$ the quantity $x$ becomes $x-\omega$; we shall have $x-\omega=\dfrac{1}{2^{n}}\log\dfrac{4}{c^{n+1}}$. But according to article 65, we have $c^{n+1}=\frac{1}{4}(c^{n})^{2}+\dfrac{1\cdot3}{4\cdot6}(c^{n})^{4}$, therefore $\dfrac{4}{c^{n+1}}=\left(\dfrac{4}{c^{n}}\right)^{2}[1-\frac{1}{2}(c^{n})^{2}]$, and consequently $\log\dfrac{4}{c^{n+1}}=2\log\left(\dfrac{4}{c^{n}}\right)-2c^{n+1}$; therefore
\[ x-\omega=\frac{1}{2^{n-1}}\log\frac{4}{c^{n}}-\frac{1}{2^{n-1}}c^{n+1}; \]
and thus $\omega=\dfrac{1}{2^{n-1}}c^{n+1}$. But since we have in a very approximate manner

\origpage{101}
$\pi=\dfrac{1}{2^{n}}\log\dfrac{4}{c^{n+1}}$, it follows that $c^{n+1}=4e^{-2^{n}\pi}$; therefore $\omega=2^{3-n}e^{-2^{n}\pi}$, and $\log\omega=-(n-3)\log2-2^{n}\pi$, or in common logarithms,
\[ \log\omega=-(n-3)\log2-2^{n-3}(10.915). \]
When $n=3$, this formula gives $\log\omega=-10.915$; thus the third term $\frac{1}{4}\log\dfrac{4}{c^{000}}$ already gives the value of $\pi$ exact to ten decimals.

When $n=4$, we have $\log\omega=-22.131$; thus the fourth term gives 22 exact decimals, the fifth gives 44, the sixth 88, and the seventh 175. Whence one sees that $\dfrac{1}{64}\log\dfrac{4}{c^{07}}$ is a much closer value of $\pi$ than that which Vega gave with 140 decimals. As to the value of $c^{07}$, it is deduced successively from that of $c$ by simple extractions of square roots, since from a modulus $c$ one passes to the following modulus $c^{0}$ by means of the formula $c^{0}=\dfrac{1-\sqrt{(1-c^{2})}}{1+\sqrt{(1-c^{2})}}$.

\subsection*{(74).}

Let us now suppose that the two complete functions $F^{1}(b)$, $F^{1}(c)$ are to each other in a known ratio, so that we have $F^{1}(b)=mF^{1}(c)$; it follows from the formulas of no.~72 that we shall have, 1°. $\dfrac{m\pi}{2}=\log\operatorname{tang}(45^{\circ}+\frac{1}{2}\Phi')$, $\Phi'$ being the limit of the angles $\varphi'$, $\varphi''$, $\varphi'''$, etc. calculated on the supposition $\varphi=\frac{1}{2}\pi$; that is to say, $\dfrac{m\pi}{2}$ will be the limit of the sequence
\[ \log\operatorname{tang}(45^{\circ}+\tfrac{1}{2}\varphi'),\quad \log\operatorname{tang}(45^{\circ}+\tfrac{1}{2}\varphi''),\quad \log\operatorname{tang}(45^{\circ}+\tfrac{1}{2}\varphi'''),\quad \text{etc.}; \]
2°. that $\dfrac{m\pi}{2}$ will equally be the limit of the sequence
\[ \tfrac{1}{2}\log\frac{4}{c^{0}},\quad \tfrac{1}{4}\log\frac{4}{c^{00}},\quad \tfrac{1}{8}\log\frac{4}{c^{000}},\quad \text{etc.} \]
The case of $b=c$ is included in these formulas by supposing $m=1$; but there are other cases in which one can make the application.

Thus we have found that when $c=\sin15^{\circ}$, we have $F^{1}(b)=\sqrt{3}\,F^{1}(c)$. Therefore, calculating the sequence $c^{0}$, $c^{00}$, etc. from the modulus $c=\sin15^{\circ}=\frac{1}{2}\sqrt{(2-\sqrt{3})}$, one sees that $\dfrac{\pi\sqrt{3}}{2}$ will be equal to the limit

\origpage{102}
of the sequence
\[ \tfrac{1}{2}\log\frac{4}{c^{0}},\quad \tfrac{1}{4}\log\frac{4}{c^{00}},\quad \tfrac{1}{8}\log\frac{4}{c^{000}},\quad \text{etc.}; \]
the first term would give $\pi=\dfrac{2}{\sqrt{3}}\log\dfrac{2}{\operatorname{tang}7^{\circ}\frac{1}{2}}=3.141636$, a value which differs from the true one only in the fifth decimal.

We shall see hereafter that, by making $c=\sin\alpha$, $\alpha$ being determined by the equation $\sin2\alpha=\operatorname{tang}^{2}15^{\circ}$, we have $F^{1}(b)=3F^{1}(c)$; thus, calculating the sequence of moduli $c^{0}$, $c^{00}$, etc. from the value $c=\sin\alpha$, we shall have $\dfrac{3\pi}{2}=\frac{1}{2}\log\dfrac{4}{c^{0}}$, $\frac{1}{4}\log\dfrac{4}{c^{00}}$, etc.; and by the first term of this sequence we obtain $\pi=3.1415926627$, the error being only one decimal unit of the eighth order; whence one may conclude that at the fifth term the approximation would be equivalent to about 128 decimals.

Finally, we have seen that, taking $b=\sqrt{2}-1$, $c=\sqrt{(2\sqrt{2}-2)}$, the functions $F^{1}(c)$, $F^{1}(b)$ are such that we have $F^{1}(c)=\sqrt{2}\,F^{1}(b)$; in this case one would make $m=\dfrac{1}{\sqrt{2}}$, and the preceding formulas will offer a new application by giving the approximate value of $\dfrac{\pi}{2\sqrt{2}}$.

\section*{Method of approximation applied to elliptic functions of the second species.}

\subsection*{(75).}

Let us consider generally a function $G$ composed of the first and the second species, so that we have
\[ G=\int(A+B\sin^{2}\varphi)\frac{d\varphi}{\Delta}. \]
If we substitute in it the values $\sin^{2}\varphi=\frac{1}{2}(1+c^{0}\sin^{2}\varphi^{0}-\Delta^{0}\cos\varphi^{0})$, $\dfrac{d\varphi}{\Delta}=\dfrac{1+c^{0}}{2}\cdot\dfrac{d\varphi^{0}}{\Delta^{0}}$, we shall have the transformed function
\[ G=\frac{1+c^{0}}{2}(G^{0}-\tfrac{1}{2}B\sin\varphi^{0}), \]
where we suppose $G^{0}=\displaystyle\int(A^{0}+B^{0}\sin^{2}\varphi^{0})\frac{d\varphi^{0}}{\Delta^{0}}$, $A^{0}=A+\frac{1}{2}B$, $B^{0}=\frac{1}{2}Bc^{0}$.

\origpage{103}
We shall then find, by similar transformations,
\[ \begin{aligned} G^{0}&=\frac{1+c^{00}}{2}(G^{00}-\tfrac{1}{2}B^{0}\sin\varphi^{00})\\ G^{00}&=\frac{1+c^{000}}{2}(G^{000}-\tfrac{1}{2}B^{00}\sin\varphi^{000})\\ &\text{etc.} \end{aligned} \]
Thus the integral formula $G$ is reduced successively to the similar integrals $G^{0}$, $G^{00}$, $G^{000}$, etc., and a single one of these integrals suffices to determine all the others.

Let us suppose that this sequence is prolonged as far as a term $G^{\mu}$ remote enough that the corresponding modulus $c^{\mu}$ is of the order of the quantities one wishes to neglect; then we shall have $\Delta^{\mu}=1$, and the value of $G^{\mu}$ reduces at first to $\displaystyle\int(A^{\mu}+B^{\mu}\sin^{2}\varphi^{\mu})\,d\varphi^{\mu}$; let us observe next that we have $B^{0}=\frac{1}{2}Bc^{0}$, $B^{00}=\frac{1}{2}B^{0}c^{00}=\frac{1}{4}Bc^{0}c^{00}$, $B^{000}=\frac{1}{8}Bc^{0}c^{00}c^{000}$, etc.; therefore the sequence $B^{0}$, $B^{00}$, $B^{000}$, etc. decreases more rapidly than the sequence $c^{0}$, $c^{00}$, $c^{000}$, etc. Therefore one may put, after a certain number of terms, $B^{\mu}=0$, which will give $G^{\mu}=A^{\mu}\varphi^{\mu}=A^{\mu}2^{\mu}\Phi$, $\Phi$ being the limit of the angles $\varphi$, $\frac{1}{2}\varphi^{0}$, $\frac{1}{4}\varphi^{00}$, etc.

As for $A^{\mu}$, since we have $A^{0}=A+\frac{1}{2}B$, $A^{00}=A^{0}+\frac{1}{2}B^{0}$, etc., the general expression of $A^{\mu}$ will be
\[ A^{\mu}=A+\tfrac{1}{2}B\left(1+\frac{c^{0}}{2}+\frac{c^{0}c^{00}}{4}+\frac{c^{0}c^{00}c^{000}}{8}+\text{etc.}\right). \]
Now it is easy to see that the value of $G$, determined by that of $G^{\mu}$, is composed of two parts; one $KA^{\mu}\Phi$, $K$ denoting the product $(1+c^{0})(1+c^{00})(1+c^{000})$, etc., the other algebraic or periodic, namely:
\[ -\left(\frac{1+c^{0}}{2}\right)\frac{B}{2}\sin\varphi^{0}-\left(\frac{1+c^{0}}{2}\right)\left(\frac{1+c^{00}}{2}\right)\frac{Bc^{0}}{4}\sin\varphi^{00}-\text{etc.} \]
But since we have $1+c^{0}=\dfrac{2\sqrt{c^{0}}}{c}$, $1+c^{00}=\dfrac{2\sqrt{c^{00}}}{c^{0}}$, etc., this second part can be put in the form
\[ -\frac{B}{c}\left(\frac{\sqrt{c^{0}}}{2}\sin\varphi^{0}+\frac{\sqrt{c^{0}c^{00}}}{4}\sin\varphi^{00}+\text{etc.}\right). \]

\origpage{104}
Therefore the total value of $G$ will be
\[ G=KA^{\mu}\Phi-\frac{B}{c}\left(\frac{\sqrt{c^{0}}}{2}\sin\varphi^{0}+\frac{\sqrt{c^{0}c^{00}}}{4}\sin\varphi^{00}+\frac{\sqrt{c^{0}c^{00}c^{000}}}{8}\sin\varphi^{000}+\text{etc.}\right). \]
In the case where $\varphi=\frac{1}{2}\pi$, the algebraic part vanishes, and we have simply $G^{1}=KA^{\mu}\cdot\frac{1}{2}\pi$.

The series which compose the value of $G$ will be very convergent if we have $c^{2}<\frac{1}{2}$, and then three terms suffice to give the value of $G$ approximate to about seven decimals. The same formulas may be employed with advantage even if the modulus $c$ were much nearer to unity; for if we had, for example, $c=0.985$, there would result approximately $c^{0}=\sqrt{\frac{1}{2}}$; whence one sees that there would be only one more term to calculate, that is to say four terms in all, to have the value of $G$ with the same degree of approximation.

But if $c$ is much nearer to unity, or if, having $c^{2}>\frac{1}{2}$, one wishes to obtain the same degree of approximation with the least possible number of transformed functions, then it will be proper to use the second method of which we made use in such a case for the functions of the first species.

\subsection*{(76).}

This method consists in prolonging the series of transformed functions in an inverse order. Thus the equation between $G$ and $G^{0}$, which gives
\[ G^{0}=\frac{2}{1+c^{0}}G+\tfrac{1}{2}B\sin\varphi^{0}, \]
applies similarly to the two functions $G$ and $G^{1}$\ednote{translator's note: Printed $G^{1}$ (superscript one); the function meant is $G'$, the first of the inverse transforms, as the next display shows; an apparent misprint.}, and from it we deduce
\[ G=\frac{2}{1+c}G'+\tfrac{1}{2}B'\sin\varphi, \]
an equation in which we must suppose $G'=\displaystyle\int(A'+B'\sin^{2}\varphi')\frac{d\varphi'}{\Delta'}$, $A'=A-\frac{1}{2}B'$, $B'=\dfrac{2B}{c}$. We shall similarly have, by further transformations,
\[ \begin{aligned} G'&=\frac{2}{1+c'}G''+\tfrac{1}{2}B''\sin\varphi''\\ G''&=\frac{2}{1+c''}G'''+\tfrac{1}{2}B'''\sin\varphi''\\ &\text{etc.} \end{aligned} \]
\ednote{translator's note: The first line prints $\sin\varphi''$ (a reader's pencil stroke crosses it); by the preceding formula $G=\frac{2}{1+c}G'+\frac{1}{2}B'\sin\varphi$ the term should be $\sin\varphi'$, so an apparent misprint.}

\origpage{105}
But when the sequence $G'$, $G''$, etc.\ has been continued to a sufficiently remote term $G^{\mu}$, we may put $c^{\mu}=1$, $\Delta^{\mu}=\cos\varphi^{\mu}$, and then the function $G^{\mu}$ becomes $\displaystyle\int(A^{\mu}+B^{\mu}\sin^{2}\varphi^{\mu})\frac{d\varphi^{\mu}}{\cos\varphi^{\mu}}$, which gives on integrating,
\[ G^{\mu}=(A^{\mu}+B^{\mu})\log\operatorname{tang}(45^{\circ}+\tfrac{1}{2}\varphi^{\mu})-B^{\mu}\sin\varphi^{\mu}, \]
and $\varphi^{\mu}$ will then be equal to the limit $\Phi'$.

We shall therefore know the value of $G$ with all the exactness one could wish, by continuing the sequence of moduli $c'$, $c''$, $c'''$, etc., and that of their complements $b'$, $b''$, $b'''$, etc., until the last term of the latter, $b^{\mu}$, is small enough to be neglected. Then we may put $c^{\mu}=1$, and, calculating an equal number of terms of the series of amplitudes $\varphi'$, $\varphi''$, $\varphi'''$, etc., we shall have the last term $\varphi^{\mu}$, which may be taken for the limit $\Phi'$.

Making, then, the substitutions necessary to obtain the value of $G$ by means of $G^{\mu}$, and observing that $\dfrac{2}{1+c}=\dfrac{c'}{\sqrt{c}}$, $\dfrac{2}{1+c'}=\dfrac{c''}{\sqrt{c'}}$, etc., we shall find this general formula
\[ \begin{aligned} G=K'(A^{\mu}+B^{\mu})\log\operatorname{tang}(45^{\circ}+\tfrac{1}{2}\varphi^{\mu})&-\frac{B}{c}\cdot\frac{2^{\mu}c^{\mu}\sin\varphi^{\mu}}{\sqrt{(cc'c''\ldots c^{\mu-1})}}\\ &+\frac{B}{c}\left(\sin\varphi+\frac{2}{\sqrt{c}}\sin\varphi'+\frac{4}{\sqrt{cc'}}\sin\varphi''\ldots+\frac{2^{\mu-1}\sin\varphi^{\mu-1}}{\sqrt{(cc'c''\ldots c^{\mu-2})}}\right). \end{aligned} \]
We have put, as above, $K'=c^{\mu}\sqrt{\left(\dfrac{c'c''\ldots c^{\mu-1}}{c}\right)}$, or simply $K'=\sqrt{\left(\dfrac{c'c''\ldots c^{\mu-1}}{c}\right)}$, because $c^{\mu}$ may be taken for unity. As to the value of $A^{\mu}+B^{\mu}$, if we call it $L^{\mu}$, we shall have
\[ L^{\mu}=A+\frac{2^{\mu}B}{cc'c''\ldots c^{\mu-1}}-\frac{B}{c}-\frac{2B}{cc'}-\frac{4B}{cc'c''}\ldots-\frac{2^{\mu-1}B}{cc'c''\ldots c^{\mu-1}}. \]
But this quantity, expressed thus by a divergent sequence, would be little convenient for calculations of approximation, and it is fitting to put it in another form. Now, on expressing $L^{\mu}$ by means of $L$ and the successive differences $L'-L$, $L''-L'$, etc., we find

\origpage{106}
easily
\[ L^{\mu}=A+\frac{B}{c}+\frac{2B}{cc'}(1-c')+\frac{4B}{cc'c''}(1-c'')+\text{etc.}; \]
or, what comes to the same thing,
\[ L^{\mu}=A+B+\frac{Bb}{c}\left[\sqrt{b'}+\sqrt{\left(\frac{b'b''}{c}\right)}+\sqrt{\left(\frac{b'b''b'''}{cc'}\right)}\ldots+\sqrt{\left(\frac{b'b''\ldots b^{\mu}}{cc'\ldots c^{\mu-2}}\right)}\right], \]
and in this form the convergence of the sequence is manifest.

We can likewise make the sequence of algebraic quantities contained in the value of $G$ convergent. For this, let
\[ \psi(\mu)=\frac{2^{\mu}c^{\mu}\sin\varphi^{\mu}}{\sqrt{(cc'c''\ldots c^{\mu-1})}}-\sin\varphi-\frac{2}{\sqrt{c}}\sin\varphi'-\frac{4}{\sqrt{cc'}}\sin\varphi''\ldots-\frac{2^{\mu-1}\sin\varphi^{\mu-1}}{\sqrt{(cc'c''\ldots c^{\mu-2})}}, \]
we deduce from this
\[ \begin{aligned} \psi(\mu+1)-\psi(\mu)&=\frac{2^{\mu+1}c^{\mu+1}\sin\varphi^{\mu+1}}{\sqrt{(cc'\ldots c^{\mu})}}-\frac{2^{\mu}c^{\mu}\sin\varphi^{\mu}}{\sqrt{(cc'\ldots c^{\mu-1})}}-\frac{2^{\mu}\sin\varphi^{\mu}}{\sqrt{(cc'\ldots c^{\mu-1})}}\\ &=\frac{2^{\mu+1}}{\sqrt{(cc'c''\ldots c^{\mu})}}\left(c^{\mu+1}\sin\varphi^{\mu+1}-\frac{c^{\mu}}{c^{\mu+1}}\sin\varphi^{\mu}\right). \end{aligned} \]
Now when $b^{\mu}$ has become very small, the equation between $\varphi^{\mu}$ and $\varphi^{\mu+1}$, which is $\operatorname{tang}(\varphi^{\mu}-\varphi^{\mu+1})=b^{\mu+1}\operatorname{tang}\varphi^{\mu+1}$, gives very nearly $\varphi^{\mu}=\varphi^{\mu+1}+b^{\mu+1}\operatorname{tang}\varphi^{\mu+1}$, or $\sin\varphi^{\mu}=(1+b^{\mu+1})\sin\varphi^{\mu+1}$. Hence $c^{\mu+1}\sin\varphi^{\mu+1}-\dfrac{c^{\mu}}{c^{\mu+1}}\sin\varphi^{\mu}=\dfrac{b^{\mu+1}}{c^{\mu+1}}(1-b^{\mu+1})\sin\varphi^{\mu+1}$. This quantity is of the order $b^{\mu+1}$; therefore, to express $\psi(\mu)$, we shall have this very convergent sequence:
\[ \begin{aligned} \psi(\mu)=c\sin\varphi&+\frac{2}{\sqrt{c}}\left(c'\sin\varphi'-\frac{c}{c'}\sin\varphi\right)\\ &+\frac{4}{\sqrt{cc'}}\left(c''\sin\varphi''-\frac{c'}{c''}\sin\varphi'\right)\\ &+\frac{8}{\sqrt{cc'c''}}\left(c'''\sin\varphi'''-\frac{c''}{c'''}\sin\varphi''\right)\\ &\ \ \vdots\\ &+\frac{2^{\mu}}{\sqrt{(cc'\ldots c^{\mu-1})}}\left(c^{\mu}\sin\varphi^{\mu}-\frac{c^{\mu-1}}{c^{\mu}}\sin\varphi^{\mu-1}\right). \end{aligned} \]

\origpage{107}
This being laid down, the value of $G$, expressed wholly in convergent sequences, will be
\[ G=K'L^{\mu}\log\operatorname{tang}(45^{\circ}+\tfrac{1}{2}\varphi^{\mu})-\frac{B}{c}\psi(\mu). \]
It must not be lost sight of that this approximation is founded on having substituted, in the value of $G^{\mu}$, $\cos\varphi^{\mu}$ for the radical $\Delta^{\mu}$, whose true value is $\sqrt{[\cos^{2}\varphi^{\mu}+(b^{\mu}\sin\varphi^{\mu})^{2}]}$. This substitution is permitted so long as $b^{\mu}\operatorname{tang}\varphi^{\mu}$ is a negligible quantity; but it could no longer take place if $\varphi^{\mu}$ were so near $90^{\circ}$ that $b^{\mu}\operatorname{tang}\varphi^{\mu}$ ceased to be very small. Suppose, which is the most unfavourable case, that $\varphi^{\mu}=90^{\circ}$; then the sequences will have to be continued one term further, or to the $(\mu+1)^{\text{th}}$ term, which will give $\operatorname{tang}\varphi^{\mu+1}=\dfrac{1}{\sqrt{(b^{\mu+1})}}$, or $b^{\mu+1}\operatorname{tang}\varphi^{\mu+1}=\sqrt{(b^{\mu+1})}$, a quantity whose square $b^{\mu+1}$ is entirely negligible with respect to unity, since $b^{\mu}$ was already supposed negligible. The formula will therefore apply without any difficulty; thus the means of preventing every exception is to take $\mu$ large enough that $\varphi^{\mu-1}$ does not exceed $90^{\circ}$, a condition which will always be easy to fulfil.

To obtain the value of the complete function $G^{1}$, one could put $\varphi=90^{\circ}$ in the general formula; but this formula would not be appreciably simplified, and it would still contain an indefinite number of terms. It will be simpler, conforming to the preceding observation, to put $\varphi^{\mu-1}=90^{\circ}$, which gives $\operatorname{tang}\varphi^{\mu}=\dfrac{1}{\sqrt{b^{\mu}}}$, $\sin\varphi^{\mu}=\dfrac{1}{\sqrt{(1+b^{\mu})}}$, and, going backwards, $\varphi^{\mu-2}=\pi$, $\varphi^{\mu-3}=2\pi\ldots\varphi^{\mu-i}=2^{i-2}\pi$, $\varphi'=2^{\mu-3}\pi$, $\varphi=2^{\mu-2}\pi=2^{\mu-1}\dfrac{\pi}{2}$. We shall therefore have $G(\varphi)=2^{\mu-1}G^{1}$; but in the hypothesis that $b^{\mu}$ is of the order of negligible quantities, we have $\sin\varphi^{\mu}=1-\frac{1}{2}b^{\mu}$, and $\log\operatorname{tang}(45^{\circ}+\frac{1}{2}\varphi^{\mu})=\frac{1}{2}\log\left(\dfrac{1+\sin\varphi^{\mu}}{1-\sin\varphi^{\mu}}\right)=\frac{1}{2}\log\left(\dfrac{4}{b^{\mu}}\right)$; we shall have at the same time the quantity $\psi(\mu)=\dfrac{2^{\mu}c^{\mu}}{cK'}\sin\varphi^{\mu}-\dfrac{2^{\mu-1}}{cK'}=\dfrac{2^{\mu-1}}{cK'}$; hence
\[ 2^{\mu-1}G^{1}=\frac{K'L^{\mu}}{2}\log\left(\frac{4}{b^{\mu}}\right)-\frac{2^{\mu-1}B}{c^{2}K'}, \]

\origpage{108}
and finally
\[ G^{1}=\frac{K'L^{\mu}}{2^{\mu}}\log\left(\frac{4}{b^{\mu}}\right)-\frac{B}{c^{2}K'}. \]
These formulas for determining both the indefinite function $G$ and the complete function $G^{1}$ allow the approximation to be pushed as far as one wishes. It will be necessary to continue the calculation of the increasing moduli $c'$, $c''$, $c'''$, etc., and that of their complements $b'$, $b''\ldots b^{\mu}$, until $b^{\mu}$ belongs to the order of decimals that one wishes to neglect, or is less than this order. It will be necessary at the same time, if it concerns the non-complete function $G$, not to give $\varphi$ a value greater than $2^{\mu-2}\pi$. We shall then calculate, up to the same limit, the values of $K'$, $L^{\mu}$ and $\psi(\mu)$.

\subsection*{(77).}

To apply these formulas to arcs of an ellipse, or to functions properly so called of the second species, let $G=E=\displaystyle\int\Delta\,d\varphi$; we shall have $A=1$, $B=-c^{2}$.

By the first method, the expression of the indefinite arc will be
\[ E(c,\varphi)=KL\Phi+\frac{c\sqrt{c^{0}}}{2}\sin\varphi^{0}+\frac{c\sqrt{c^{0}c^{00}}}{4}\sin\varphi^{00}+\frac{c\sqrt{c^{0}c^{00}c^{000}}}{8}\sin\varphi^{000}+\text{etc.} \]
and that of the quarter of the ellipse
\[ E^{1}(c)=KL\frac{\pi}{2}. \]
In these formulas, $K$ represents the product $(1+c^{0})(1+c^{00})(1+c^{000})$, etc., continued up to a factor $1+c^{\mu}$, where $c^{\mu}$ is of the order of negligible quantities. This same quantity may be put in the form $K=\dfrac{2\sqrt{c^{0}}}{c}\cdot\dfrac{2\sqrt{c^{00}}}{c^{0}}\cdot\dfrac{2\sqrt{c^{000}}}{c^{00}}\cdot\text{etc.}$, more convenient for logarithmic calculation. Finally the quantity $L$ is given by the convergent series
\[ L=1-\frac{c^{0}}{2}-\frac{c^{2}c^{0}}{4}-\frac{c^{2}c^{0}c^{00}}{4}-\text{etc.}, \]
\ednote{Printed $\frac{c^{0}}{2}$ in the first term and $4$ in the third denominator; the series that agrees with the second form below (and numerically) is $1-\frac{c^{2}}{2}-\frac{c^{2}c^{0}}{4}-\frac{c^{2}c^{0}c^{00}}{8}-\text{etc.}$; apparent misprints.}
to which one may give this other form
\[ L=\frac{c^{2}}{4c^{0}}\left(1-\frac{c^{02}c^{00}}{2}-\frac{c^{02}c^{00}c^{000}}{4}-\frac{c^{02}c^{00}c^{000}c^{0000}}{8}-\text{etc.}\right); \]

\origpage{109}
so that, on putting $KL=M$, we should have
\[ M=\left(\frac{c^{2}}{2\sqrt{c^{0}}}\right)\cdot\frac{2\sqrt{c^{00}}}{c^{0}}\cdot\frac{2\sqrt{c^{000}}}{c^{00}}\,\text{etc.}\left(1-\frac{c^{02}c^{00}}{2}-\frac{c^{02}c^{00}c^{000}}{4}-\text{etc.}\right) \]
These formulas offer the most convergent series that one could wish for the rectification of the ellipse. They apply, as we have already said, not only to values of $c$ smaller than $\sqrt{\frac{1}{2}}$, but to much larger values very close to unity.

\section*{EXAMPLE.}

\subsection*{(78).}

Let it be proposed to find the arc $E(c,\varphi)$ when $c=\sin 75^{\circ}$ and $\operatorname{tang}\varphi=\sqrt{\dfrac{2}{\sqrt{3}}}$.

By the values already calculated (art.~66), we have $\log K=0.2460561$, $\Phi=30^{\circ}$, $\varphi^{0}=62^{\circ}36'3''.10$, $\varphi^{00}=119^{\circ}55'47''67$, $\varphi^{000}=240^{\circ}0'0''19$, then we shall find
\[ \begin{aligned} L&=0.3888658\\ \frac{c\sqrt{c^{0}}}{2}\sin\varphi^{0}&=0.3290186\\ \frac{c\sqrt{c^{0}c^{00}}}{4}\sin\varphi^{00}&=0.0522872\\ \frac{c\sqrt{c^{0}c^{00}c^{000}}}{8}\sin\varphi^{000}&=-0.0013888\\ \frac{c\sqrt{c^{0}c^{00}c^{000}c^{0000}}}{16}\sin\varphi^{0000}&=0.0000010 \end{aligned} \]
The sum of the algebraic parts is 0.3799180, so we have $E(c,\varphi)=KL\cdot\frac{\pi}{6}+0.3799180$; we shall have at the same time the complete function $E^{1}(c)=KL\cdot\frac{\pi}{2}$, whence would result $E(c,\varphi)=\frac{1}{3}E^{1}(c)+0.3799180$. But since the arc $E(c,\varphi)$ is measured by the third of $E^{1}(c)$, we shall find by the formulas of trisection (art.~57) $E(c,\varphi)=\frac{1}{3}E^{1}(c)+\dfrac{c^{2}\cos^{3}\varphi}{3b}$; the preceding result agrees perfectly with this formula, for we have in fact $\dfrac{c^{2}\cos^{3}\varphi}{3b}=0.3799179$.

The same calculation gives moreover $\log E^{1}(c)=0.0319757$ and $E^{1}(c)=1.0764049$; hence in the case proposed, the arc $E(c,\varphi)=0.7387196$; and we see that the part determined algebraically forms more than half the arc. For the rest, the value found for $E^{1}(c)$ agrees with

\origpage{110}
that which would be deduced from the formula of article 41, by putting in it for $F^{1}(c)$ its value found in art.~66.

\subsection*{(79).}

If the modulus $c$ differs little enough from unity that it is judged necessary to apply the second method, it will be necessary, in the formulas of art.~76, to put $A=1$, $B=-c^{2}$, which will give
\[ L^{\mu}=b^{2}-bc\left[\sqrt{b'}+\sqrt{\left(\frac{b'b''}{c}\right)}+\sqrt{\left(\frac{b'b''b'''}{cc'}\right)}\ldots+\sqrt{\left(\frac{b'b''\ldots b^{\mu}}{cc'\ldots c^{\mu-2}}\right)}\right]. \]
As to the value of $K'$ and that of the function $\psi(\mu)$, they do not depend at all on the values of $A$ and $B$; so they remain the same as in the article cited. We shall therefore have, from these values, the indefinite function
\[ E(c,\varphi)=K'L^{\mu}\operatorname{tang}(45^{\circ}+\tfrac{1}{2}\Phi')+c\psi(\mu), \]
\ednote{Printed without $\log$ and with $\Phi'$; by the general formula of art.\ 76 and the case $\mu=2$ below, one expects $\log\operatorname{tang}(45^{\circ}+\frac{1}{2}\varphi^{\mu})$; an apparent misprint.}
and the complete function
\[ E^{1}(c)=\frac{K'L^{\mu}}{2^{\mu}}\log\frac{4}{b^{\mu}}+\frac{1}{K'}. \]
These formulas are the simplest that analysis can offer for the calculation of arcs of an ellipse when the modulus is very near unity; they suppose that for $\mu$ we take a number of units large enough that $b^{\mu}$ belongs to the order of decimals that one wishes to neglect.

If we wish to have the values of $E(c,\varphi)$ and of $E^{1}(c)$ approximated only to the seventh order of decimals, it will suffice to put $\mu=2$ in the preceding formulas; we shall have first $L^{\mu}=b^{2}-bc\sqrt{b'}-bc\sqrt{\left(\dfrac{b'b''}{c}\right)}$, and because we have in general $b\sqrt{b'}=1-c$, $b'\sqrt{b''}=1-c'$, $c'=\dfrac{2\sqrt{c}}{1+c}$, we find more simply $L^{\mu}=\dfrac{b^{2}}{1+\sqrt{c}}$: but in the same case we have $K'=\sqrt{\dfrac{c'}{c}}$. Hence the two general formulas reduce to these
\[ \begin{aligned} E(c,\varphi)&=\frac{b^{2}}{1+\sqrt{c}}\sqrt{\frac{c'}{c}}\log\operatorname{tang}(45^{\circ}+\tfrac{1}{2}\varphi'')+c^{2}\sin\varphi\\ &\quad+2\sqrt{c}\left(c'\sin\varphi'-\frac{c}{c'}\sin\varphi\right)+4\sqrt{\frac{c}{c'}}(\sin\varphi''-c'\sin\varphi'),\\ E^{1}(c)&=\frac{b^{2}}{1+\sqrt{c}}\cdot\frac{1}{4}\sqrt{\frac{c'}{c}}\cdot\log\frac{4}{b''}+\sqrt{\frac{c}{c'}}. \end{aligned} \]

\origpage{111}
The last is the same as $E^{1}(c)=\dfrac{b^{2}}{1+\sqrt{c}}F^{1}(c)+\sqrt{\dfrac{c}{c'}}$, and in this form it offers a very simple relation between the functions $F^{1}(c)$, $E^{1}(c)$, all the closer the smaller $b$ is.

\subsection*{(80).}

Let us apply these formulas to the example of article 78, which is a little favourable case, since the value of $c$ differs appreciably from unity. We already know from article 69 the values of $c'$, $b'$, $c''$, $b''$, as well as those of the amplitudes $\varphi$, $\varphi'$, $\varphi''$. By means of these values, the calculation of the complete function $E^{1}(c)$ will be done thus:
\[ \begin{aligned} &1+\sqrt{c}= && 1.9828152 &\qquad& \log(1+\sqrt{c})\ldots && 0.2972822\\ &\log 4\ldots\ldots && 0.6020600 && \log b^{2}\ldots\ldots && 8.8259924\\ &\log b''\ldots\ldots && 5.8757219 && \text{diff.}\ldots\ldots && 8.5287102\\ &\log\frac{4}{b''}\ldots && 4.7263381 && \log K'\ldots\ldots && 0.0074955\\ &\tfrac{1}{4}\log\frac{4}{b''}= && 1.1815845 && \log 1.1815\ \text{etc.}\ldots && 0.0724648\\ &&&&& \log 2.3025\ \text{etc.}\ldots && 0.3622157\\ &&&&& \log x\ldots\ldots && 8.9708862 \end{aligned} \]
\[ \begin{aligned} x&=0.0935161\\ \frac{1}{K'}&=0.9828891\\ \text{Sum}\ldots\ E^{1}(c)&=1.0764052 \end{aligned} \]
This value agrees with that which we have already found, as far as the small neglected parts allow.

As to the calculation of the arc $E(c,\varphi)$, we have already found the value of $E(c,\varphi)$\ednote{Printed $E(c,\varphi)$; the logarithm taken next is that of $F(c,\varphi)$, so $F(c,\varphi)$ is expected here; an apparent misprint.} $=K'\log\operatorname{tang}(45^{\circ}+\frac{1}{2}\varphi'')$, there results
\[ \begin{aligned} \log F(c,\varphi)&=9.9650547\\ \log\frac{b^{2}}{1+\sqrt{c}}&=8.5287102\\ \log y&=8.4937649,\quad y=0.0311720. \end{aligned} \]
We then find the algebraic parts

\origpage{112}
\[ \begin{aligned} &c^{2}\sin\varphi && =0.6830127\\ &2\sqrt{c}\left(c'\sin\varphi'-\frac{c}{c'}\sin\varphi\right) && =0.0243225\\ &4\sqrt{\frac{c}{c'}}(\sin\varphi''-c'\sin\varphi') && =0.0002123\\ &\text{sum}\ldots\ldots && 0.7075475\\ &\text{adding}\ldots\ y && =0.0311720\\ &\text{we have}\ E(c,\varphi) && =0.7387195 \end{aligned} \]
which again agrees with the value found by the first method.

\subsection*{(81).}

To apply the preceding formulas to the rectification of the hyperbola, it is only a question of evaluating the integral $G=\displaystyle\int\frac{c^{2}\,d\varphi\cos^{2}\varphi}{\Delta}$, which represents the quantity $\Delta\operatorname{tang}\varphi-\Upsilon$, the difference between the arc of a hyperbola and its tangent. Now, putting in the general formula of art.~75 $A=c^{2}$, $B=-c^{2}$, we have the indefinite function
\[ \begin{aligned} G(c,\varphi)=\tfrac{1}{2}Kc^{2}\Phi&\left(1-\frac{c^{0}}{2}-\frac{c^{0}c^{00}}{4}-\frac{c^{0}c^{00}c^{000}}{8}-\text{etc.}\right)\\ &+\frac{c\sqrt{c^{0}}}{2}\sin\varphi+\frac{c\sqrt{c^{0}c^{00}}}{4}\sin\varphi^{00}+\frac{c\sqrt{c^{0}c^{00}c^{000}}}{8}\sin\varphi^{000}+\text{etc.} \end{aligned} \]
\ednote{The second line prints $\sin\varphi$ in its first term; the corresponding series on p. 108 has $\sin\varphi^{0}$, so either a misprint or a superscript lost in printing.}
and the definite function
\[ G^{1}(c)=Kc^{2}\cdot\frac{\pi}{4}\left(1-\frac{c^{0}}{2}-\frac{c^{0}c^{00}}{4}-\frac{c^{0}c^{00}c^{000}}{8}-\text{etc.}\right): \]
this is the difference between the asymptote and the curve, which is thus expressed very simply by a very convergent series.

When the hyperbola is equilateral, the quantity $G^{1}(c)$, which in general has for its value $E^{1}(c)-b^{2}F^{1}(c)$, reduces to $E^{1}(c)-\frac{1}{2}F^{1}(c)$; but as in the particular case where $c^{2}=\frac{1}{2}=b^{2}$ we have $\dfrac{\pi}{2}=(2E^{1}-F^{1})F^{1}$, the value of $G^{1}$ will reduce further to $\dfrac{\frac{1}{4}\pi}{F^{1}}$, and because $F^{1}=\frac{1}{2}\pi K$, we shall have still more simply $G^{1}=\dfrac{1}{2K}$. We found above $\log K=0.0720074$, so we shall have $\log G^{1}=9.6269626$.

The formula

\origpage{113}
$\Upsilon=\Delta\operatorname{tang}\varphi-G(c,\varphi)$ may be regarded as applicable to all cases; nevertheless if $c$ were extremely near unity, one could, to determine the function $G$, make use of the formulas of article 76, after putting $A=c^{2}$, $B=-c^{2}$.

\section*{Remarkable formulas for determining the complete functions $F^{1}(c)$, $E^{1}(c)$, when $c$ is little removed from one of its limits.}

\subsection*{(82).}

When the modulus $c$ differs very little from unity, we have found (art.~76) the two formulas
\[ \begin{aligned} F^{1}(c)&=\tfrac{1}{4}\sqrt{\left(\frac{c'}{c}\right)}\log\frac{4}{b''}\\ E^{1}(c)&=\frac{b^{2}}{1+\sqrt{c}}F^{1}(c)+\sqrt{\left(\frac{c}{c'}\right)}. \end{aligned} \]
These formulas will give about ten exact decimals if $c=\sin 80^{\circ}$; they would give a greater number if $c$ were larger; they will give only seven when one makes $c=\sin 64^{\circ}$, and so in proportion in the other cases. For the rest their use is very convenient since it requires only that one calculate with precision the values of $c'$, $b'$ and $b''$. For this purpose one must recall that, on putting $c=\cos\mu$, which gives $b=\sin\mu$, we shall have $c'=\dfrac{\sqrt{\cos\mu}}{\cos^{2}\frac{1}{2}\mu}$, $b'=\operatorname{tang}^{2}\frac{1}{2}\mu$. If $\mu$ is very small, the tables will give $\cos\mu$ and $\cos\frac{1}{2}\mu$ with precision; knowing $\cos\frac{1}{2}\mu$ and $\sin\mu$, we shall deduce from them $\sin\frac{1}{2}\mu=\dfrac{\sin\mu}{2\cos\frac{1}{2}\mu}$, and $\operatorname{tang}\frac{1}{2}\mu=\dfrac{\sin\frac{1}{2}\mu}{\cos\frac{1}{2}\mu}$, which will make $b'$ known; finally, to avoid angles that are too small, we shall observe that $b''$ is deduced from $b'$ as $c^{0}$ from $c$, and that thus we have $b''=\frac{1}{4}b'^{2}+\dfrac{1\cdot 3}{4\cdot 6}b'^{4}+\text{etc.}$; whence results $\dfrac{4}{b''}=\left(\dfrac{4}{b'}\right)^{2}(1-\frac{1}{2}b'^{2})$, and consequently $\frac{1}{2}\log\dfrac{4}{b''}=\log\dfrac{4}{b'}-\frac{1}{4}b'^{2}m$, $m$ being put for the number $0.43429$ etc.; whence we see that with a slight correction easy to calculate, we shall deduce $\log\dfrac{4}{b''}$ from $\log\dfrac{4}{b'}$. The rest presents no difficulty, remembering that the logarithms in the formula are hyperbolic logarithms which will have to be reduced to common logarithms.

\origpage{114}
\subsection*{(83).}

If we substitute the value of $E^{1}(c)$ given by the second of the preceding formulas in the equation of article 42, we shall deduce
\[ F^{1}(c)=\frac{\tfrac{1}{2}\pi-\sqrt{\left(\frac{1+c}{2}\sqrt{c}\right)}\cdot F^{1}(b)}{E^{1}(b)-\left(\frac{c^{2}+\sqrt{c}}{1+\sqrt{c}}\right)F^{1}(b)}, \]
and if we observe that the value of $F^{1}(c)$ contains a logarithm which the second member could not offer through the development of the quantities $F^{1}(b)$, $E^{1}(b)$, in series ordered according to the powers of $b$ supposed very small, we shall conclude that the preceding equation cannot subsist unless the numerator and the denominator of the second member are appreciably equal to zero. Thus we must have the two equations
\[ \begin{aligned} F^{1}(b)&=\frac{\tfrac{1}{2}\pi}{\sqrt{\left(\frac{1+c}{2}\sqrt{c}\right)}}=\tfrac{1}{2}\pi\sqrt{\left(\frac{c'}{c}\right)}\\ E^{1}(b)&=\frac{1+c\sqrt{c}}{1+\sqrt{c}}\cdot\tfrac{1}{2}\pi\sqrt{c'}. \end{aligned} \]
Let us change $b$ into $c$, and reciprocally, in order to refer these new formulas to the modulus $c$; we shall have
\[ \begin{aligned} F^{1}(c)&=\frac{\tfrac{1}{2}\pi}{\sqrt{\left(\frac{1+b}{2}\sqrt{b}\right)}}=\tfrac{1}{2}\pi\sqrt{\left(\frac{b^{0}}{b}\right)}\\ E^{1}(c)&=\frac{1+b\sqrt{b}}{1+\sqrt{b}}\cdot\tfrac{1}{2}\pi\sqrt{b^{0}}=(2-b^{0})\frac{\pi}{2}\sqrt{\left(\frac{b}{b^{0}}\right)}. \end{aligned} \]
These formulas are still simpler than those of the preceding article; they suppose no other preliminary calculation than that of $b$ and $b^{0}$; now on putting $c=\sin\mu$, we have $b=\cos\mu$, and $b^{0}=\dfrac{\sqrt{\cos\mu}}{\cos^{2}\frac{1}{2}\mu}$: besides, $\log(2-b^{0})$ is very easy to deduce from $\log b^{0}$; for as $1-b^{0}$ is very small, if one puts $\log b^{0}=-\delta$, we shall have $\log(2-b^{0})=\delta-\delta^{2}(2.3025,\text{etc.})$.

These formulas are no less exact for small values of $c$ than those of the preceding article for values of $c$ little different from unity; the one set and the other agree with the formulas of article 45, and one may satisfy oneself that they are equivalent to the development of a fairly large number of terms of the latter.

\origpage{115}
\subsection*{(84).}

Let us apply these formulas to the calculation of the functions $F^{1}(c)$, $F^{1}(b)$, supposing $c=\dfrac{\sqrt[4]{4}-\sqrt[4]{3}}{1+\sqrt{3}}$ and $b=\dfrac{\sqrt[4]{4}+\sqrt[4]{3}}{1+\sqrt{3}}$. Let $c=\sin\alpha$, $b=\cos\alpha$, we shall have $\sin 2\alpha=\operatorname{tang}^{2}15^{\circ}$ and $\operatorname{tang}(45^{\circ}+\alpha)=\sqrt{\left(\dfrac{2}{\sqrt{3}}\right)}$; from one or the other of these equations we shall draw
\[ \alpha=2^{\circ}\,3'\,30''.94970. \]
This angle being rather small, some precautions will be needed in calculating the values of the quantities sought to ten decimals. We find first directly from the tables, $\log\sin 2\alpha=8.8561049048$; then the known value of $\alpha$ gives $\log\cos\alpha=9.9997196211$; from this we deduce
\[ \log\sin\alpha=8.55535\ 52880. \]
For the calculation of the formula $F^{1}(c)=\frac{1}{2}\pi\sqrt{\dfrac{b^{0}}{b}}$, I put $k=\sqrt{\dfrac{b^{0}}{b}}=\dfrac{1}{\cos\frac{1}{2}\alpha\sqrt[4]{\cos\alpha}}$, and I find
\[ \begin{aligned} &\cos\tfrac{1}{2}\alpha\ldots\ldots && 9.99992\ 99167\\ &\sqrt[4]{\cos\alpha}\ldots\ldots && 9.99992\ 99053\\ &\text{sum}\ldots\ldots && 9.99985\ 98220\\ &k\ldots\ldots && 0.00014\ 01780\\ &\tfrac{1}{2}\pi\ldots\ldots && 0.19611\ 98770\\ &\text{Hence}\ \log F^{1}(c) && =0.19626\ 00550 \end{aligned} \]
The calculation of $F^{1}(b)$ will be done by the first formula, which gives $F^{1}(b)=\frac{1}{4}k\log\dfrac{4}{c^{00}}$. Now we have $c^{0}=\operatorname{tang}^{2}\frac{1}{2}\alpha$; thus $\log c^{0}$ is deduced from the logarithms already found in this way:

\origpage{116}
\[ \begin{aligned} &\cos\tfrac{1}{2}\alpha\ldots\ldots && 9.99992\ 99167\\ &2\ldots\ldots && 0.30102\ 99957\\ &2\cos\tfrac{1}{2}\alpha\ldots\ldots && 0.30095\ 99124\\ &\sin\alpha\ldots\ldots && 8.55535\ 52880\\ &\sin\tfrac{1}{2}\alpha\ldots\ldots && 8.25439\ 53756\\ &\cos\tfrac{1}{2}\alpha\ldots\ldots && 9.99992\ 99167\\ &\operatorname{tang}\tfrac{1}{2}\alpha\ldots\ldots && 8.25446\ 54589\\ &c^{0}\ldots\ldots && 6.50893\ 09178\\ &2\ldots\ldots && 0.30102\ 99957\\ &\tfrac{1}{2}c^{0}\ldots\ldots && 6.20790\ 09221\\ &\tfrac{1}{4}c^{02}\ldots\ldots && 2.41580\ 18442. \end{aligned} \]
It is further necessary to seek the value of $c^{00}$; but, as has already been observed, one may dispense with the calculation of $c^{00}$ by the formula $\frac{1}{2}\log\dfrac{4}{c^{00}}=\log\dfrac{4}{c^{0}}-\frac{1}{4}c^{02}(0.434\ \text{etc.})$. In this way we obtain in common logarithms $\frac{1}{4}\log\dfrac{4}{c^{00}}=2.0465645311$. Taking the third of this number, and putting for brevity $\beta=0.6821881770$, we shall have $\frac{1}{3}F^{1}(b)=\beta Mk$, $M$ being put for the factor 2.3025 etc.
\[ \begin{aligned} &\beta\ldots\ldots && 9.83390\ 41883\\ &M\ldots\ldots && 0.36221\ 56887\\ &k\ldots\ldots && 0.00014\ 01780\\ &\tfrac{1}{3}F^{1}(b)\ldots\ldots && 0.19626\ 00550. \end{aligned} \]
This value is exactly the same as that found for the logarithm of $F^{1}(c)$; thus we have $F^{1}(b)=3F^{1}(c)$. This equality is at least very closely approximate, since the formulas from which it was drawn could, with more extensive tables, give the values of $F^{1}(b)$ and $F^{1}(c)$ approximated at least to 15 decimals; but it is very probable that it is rigorous, and this we shall have occasion to confirm by a particular demonstration.

The equality $F^{1}(b)=3F^{1}(c)$ being supposed exact on the strength of this result, we find by other considerations that the functions $F^{1}(c)$ and $F^{1}(\sin 45^{\circ})$ have between them this relation $F^{1}(c)=\dfrac{2\sqrt[4]{3}}{3}\cos 15^{\circ}\cdot F^{1}(\sin 45^{\circ})$;

\origpage{117}
this will be easily verified by the value found for $F^{1}(c)$, and by that of $F^{1}(\sin 45^{\circ})$ taken from the little Table which follows, calculated to ten decimals.
\[ \begin{aligned} c && \log F^{1}(c) &\qquad& c && \log F^{1}(c)\\ \sin 0^{\circ} && 0.19611\ 98770 && \sin 90^{\circ} && \ldots\text{infinite}\ldots\\ \sin 15^{\circ} && 0.20361\ 53664 && \sin 75^{\circ} && 0.44217\ 59938\\ \sin 30^{\circ} && 0.22679\ 32597 && \sin 60^{\circ} && 0.33375\ 26135\\ \sin 45^{\circ} && 0.26812\ 72224 && \sin 45^{\circ} && 0.26812\ 72224 \end{aligned} \]

We append here a more extensive Table of the logarithms of the complete functions $F^{1}(c)$, $E^{1}(c)$, calculated to seven decimals, for all the angles of the modulus, degree by degree, the modulus being put in the form $c=\sin\theta$. This Table will be very useful in the theory of elliptic functions; it has been calculated by the abridged formulas that we have just reported, from $\theta=0^{\circ}$ to $\theta=20^{\circ}$, and from $\theta=90^{\circ}$ to $\theta=70^{\circ}$; for the other values of $\theta$, the formulas of articles 65 and 77 have been followed.

\origpage{118}
\[ \begin{aligned} \text{Angle of the modulus.} && \text{Log. }F^{1} && \text{Log. }E^{1} &\qquad& \text{Angle of the modulus.} && \text{Log. }F^{1} && \text{Log. }E^{1}\\ 0^{\circ} && 0.1961199 && 0.1961199 && 90^{\circ} && \text{Infinite.} && 0.0000000\\ 1 && 0.1961530 && 0.1960868 && 89 && 0.7351923 && 0.0003264\\ 2 && 0.1962522 && 0.1959876 && 88 && 0.6760272 && 0.0012533\\ 3 && 0.1964176 && 0.1958222 && 87 && 0.6373550 && 0.0022773\\ 4 && 0.1966493 && 0.1955908 && 86 && 0.6077507 && 0.0037396\\ 5 && 0.1969474 && 0.1952932 && 85 && 0.5834041 && 0.0054653\\ 6 && 0.1973119 && 0.1949295 && 84 && 0.5625137 && 0.0074221\\ 7 && 0.1977430 && 0.1944998 && 83 && 0.5441205 && 0.0095837\\ 8 && 0.1982409 && 0.1940041 && 82 && 0.5276129 && 0.0119271\\ 9 && 0.1988057 && 0.1934423 && 81 && 0.5125914 && 0.0144321\\ 10 && 0.1994378 && 0.1928147 && 80 && 0.4987770 && 0.0170811\\ 11 && 0.2001373 && 0.1921212 && 79 && 0.4859667 && 0.0198581\\ 12 && 0.2009044 && 0.1913618 && 78 && 0.4740076 && 0.0227487\\ 13 && 0.2017396 && 0.1905367 && 77 && 0.4627820 && 0.0257397\\ 14 && 0.2026431 && 0.1896460 && 76 && 0.4521964 && 0.0288191\\ 15 && 0.2036154 && 0.1886896 && 75 && 0.4421760 && 0.0319757\\ 16 && 0.2046567 && 0.1876678 && 74 && 0.4326595 && 0.0351996\\ 17 && 0.2057675 && 0.1865805 && 73 && 0.4235961 && 0.0384811\\ 18 && 0.2069484 && 0.1854280 && 72 && 0.4149432 && 0.0418119\\ 19 && 0.2081997 && 0.1842104 && 71 && 0.4066648 && 0.0451835\\ 20 && 0.2095220 && 0.1829278 && 70 && 0.3987298 && 0.0485886\\ 21 && 0.2109158 && 0.1815802 && 69 && 0.3911076 && 0.0520197\\ 22 && 0.2123818 && 0.1801680 && 68 && 0.3837870 && 0.0554720\\ 23 && 0.2139206 && 0.1786913 && 67 && 0.3767358 && 0.0589377\\ 24 && 0.2155328 && 0.1771502 && 66 && 0.3699400 && 0.0624119\\ 25 && 0.2172193 && 0.1755451 && 65 && 0.3633838 && 0.0658895\\ 26 && 0.2189809 && 0.1738761 && 64 && 0.3570533 && 0.0693649\\ 27 && 0.2208180 && 0.1721435 && 63 && 0.3509356 && 0.0728337\\ 28 && 0.2227319 && 0.1703476 && 62 && 0.3450196 && 0.0762928\\ 29 && 0.2247233 && 0.1684886 && 61 && 0.3392950 && 0.0797376\\ 30 && 0.2267933 && 0.1665669 && 60 && 0.3337526 && 0.0831643\\ 31 && 0.2289427 && 0.1645829 && 59 && 0.3283839 && 0.0865695\\ 32 && 0.2311728 && 0.1625368 && 58 && 0.3231814 && 0.0899500\\ 33 && 0.2334846 && 0.1604292 && 57 && 0.3181381 && 0.0933029\\ 34 && 0.2358794 && 0.1582605 && 56 && 0.3132474 && 0.0966256\\ 35 && 0.2383586 && 0.1560313 && 55 && 0.3085036 && 0.0999152\\ 36 && 0.2409233 && 0.1537418 && 54 && 0.3039013 && 0.1031694\\ 37 && 0.2435751 && 0.1513928 && 53 && 0.2994353 && 0.1063861\\ 38 && 0.2463154 && 0.1489848 && 52 && 0.2951012 && 0.1095630\\ 39 && 0.2491459 && 0.1465186 && 51 && 0.2908945 && 0.1126982\\ 40 && 0.2520685 && 0.1439949 && 50 && 0.2868114 && 0.1157899\\ 41 && 0.2550848 && 0.1414143 && 49 && 0.2828480 && 0.1188364\\ 42 && 0.2581965 && 0.1387776 && 48 && 0.2790011 && 0.1218362\\ 43 && 0.2614061 && 0.1360858 && 47 && 0.2752673 && 0.1247878\\ 44 && 0.2647155 && 0.1333399 && 46 && 0.2716436 && 0.1276898\\ 45 && 0.2681272 && 0.1305409 && 45 && 0.2681272 && 0.1305409 \end{aligned} \]
\ednote{Three entries are printed as given but differ from the exact values by more than rounding: Log.\ $F^{1}$ at $85^{\circ}$ is printed 0.5834041 (exact 0.5833963), Log.\ $F^{1}$ at $69^{\circ}$ is printed 0.3911076 (exact 0.3911115), and Log.\ $E^{1}$ at $88^{\circ}$ is printed 0.0012533 (exact 0.0011208). The other entries agree with the exact values to within a few units of the seventh decimal.}

\origpage{119}
\section*{Method of approximation applied to elliptic functions of the third species.}

\subsection*{(85).}

\ednote{The article number is printed (85), and the next article is also numbered (85); art.\ 88 cites this equation as n° 84, so (84) is meant. An apparent misprint.} Instead of considering the simple formula $\Pi=\displaystyle\int\frac{d\varphi}{(1+n\sin^{2}\varphi)\Delta}$ which belongs to the third species, we shall consider more generally the formula
\[ H=\int\left(A+\frac{B\sin^{2}\varphi}{1+n\sin^{2}\varphi}\right)\frac{d\varphi}{\Delta}, \]
which is composed of a function of the third species and one of the first, affected by constant coefficients. If the calculation were applied to the simple function $\Pi$, the first transform would contain a part affected by the function of the first species $F(c^{0},\varphi^{0})$; this is why it is fitting, for uniformity of results, to admit both kinds of function in the primitive formula, as we have already done for the functions $G$ (art.\ 76).

This being laid down, if we put $\sin^{2}\varphi=\frac{1}{2}(1+c^{0}\sin^{2}\varphi^{0}-\Delta^{0}\cos\varphi^{0})$, and $\frac{d\varphi}{\Delta}=\frac{1+c^{0}}{2}\cdot\frac{d\varphi^{0}}{\Delta^{0}}$, we shall have for the first transform
\[ H=\frac{1+c^{0}}{2}\left[H^{0}-\frac{\tfrac{1}{2}B}{1+n}\int\frac{d\varphi^{0}\cos\varphi^{0}}{1+n^{0}\sin^{2}\varphi^{0}}\right], \]
the values of $H^{0}$ and of its coefficients being
\[ \begin{aligned} H^{0}&=\int\left(A^{0}+\frac{B^{0}\sin^{2}\varphi^{0}}{1+n^{0}\sin^{2}\varphi^{0}}\right)\frac{d\varphi^{0}}{\Delta^{0}}.\\ n^{0}&=\frac{n(n+c^{2})}{(1+b)^{2}(1+n)}\\ A^{0}&=A+\frac{\tfrac{1}{2}B}{1+n}\\ B^{0}&=B\cdot\frac{c^{2}+2n+n^{2}}{2(1+b)^{2}(1+n)^{2}}. \end{aligned} \]
Thus the function $H(n,c,\varphi)$ is brought back to a similar function $H(n^{0},c^{0},\varphi^{0})$, in which the three elements $n$, $c$, $\varphi$ have undergone changes suited to facilitate approximations.

The same law will hold in the subsequent transforms; so that

\origpage{120}
by making
\[ \begin{aligned} H^{00}&=\int\left(A^{00}+\frac{B^{00}\sin^{2}\varphi^{00}}{1+n^{00}\sin^{2}\varphi^{00}}\right)\frac{d\varphi^{00}}{\Delta^{00}}\\ n^{00}&=\frac{n^{0}(n^{0}+c^{02})}{(1+b^{0})^{2}(1+n^{0})}\\ A^{00}&=A^{0}+\frac{\tfrac{1}{2}B^{0}}{1+n^{0}}\\ B^{00}&=B^{0}\cdot\frac{c^{02}+2n^{0}+n^{02}}{2(1+b^{0})^{2}(1+n^{0})^{2}}; \end{aligned} \]
we shall have the second transform
\[ H^{0}=\frac{1+c^{00}}{2}\left[H^{00}-\frac{\tfrac{1}{2}B^{0}}{1+n^{0}}\int\frac{d\varphi^{00}\cos\varphi^{00}}{1+n^{00}\sin^{2}\varphi^{00}}\right]. \]
One may continue these transformations indefinitely, whence will result an infinite sequence of functions $H(n,c,\varphi)$, $H(n^{0},c^{0},\varphi^{0})$, $H(n^{00},c^{00},\varphi^{00})$, etc., such that, any one of them being known, all the others can be determined.

\subsection*{(85).}

Let us first examine in some detail the law of progression of the parameters $n^{0}$, $n^{00}$, etc. If we put $n=mc$, and likewise $n^{0}=m^{0}c^{0}$, the equation which determines $n^{0}$ will give $m^{0}=\frac{m(m+c)}{1+mc}$. But one may always suppose the parameter $n<c$, which gives $m<1$; we shall therefore also have $m^{0}<1$, and even $m^{0}<m$; for from the preceding value one deduces $1+\frac{m^{0}}{m}=\frac{(1+c)(1+m)}{1+cm}$, $1-\frac{m^{0}}{m}=\frac{(1-c)(1-m)}{1+cm}$, which proves that $\frac{m^{0}}{m}$, positive or negative, is always smaller than unity. Hence the terms $m$, $m^{0}$, $m^{00}$, etc. are all smaller than unity and form a decreasing sequence; and consequently the sequence of parameters $n$, $n^{0}$, $n^{00}$, etc. decreases more rapidly than the sequence of moduli $c$, $c^{0}$, $c^{00}$, etc., each term of the first being smaller than the corresponding term of the second.

With regard to the signs of these parameters, there are two cases to be considered. 1\textsuperscript{o}. If $n$ is positive, or if $n$ is negative and greater than $c^{2}$, so that it is included in the form $-1+b^{2}\sin^{2}\theta$, then $n^{0}$ will be positive, and so will all the following parameters $n^{00}$, $n^{000}$, etc. In this first case is included that of $n=\pm c$

\origpage{121}
which will give $n^{0}=c^{0}$, $n^{00}=c^{00}$, etc.; but this case is useless to consider, because, according to article 46, the function $H$ is then reduced to the first species.

2\textsuperscript{o}. If $n$ is negative and smaller than $c^{2}$, so that we have $n=-c^{2}\sin^{2}\theta$, we may likewise put $n^{0}=-c^{02}\sin^{2}\theta^{0}$, and the angle $\theta^{0}$ will be determined by the equation
\[ \operatorname{tang}(\theta^{0}-\theta)=b\operatorname{tang}\theta, \]
which is the same law according to which the amplitude $\varphi^{0}$ is deduced from $\varphi$.

Calculating therefore successively $\theta^{0}$, $\theta^{00}$, $\theta^{000}$, etc. according to this law, we shall form the sequence of parameters $n^{0}=-c^{02}\sin^{2}\theta^{0}$, $n^{00}=-c^{002}\sin^{2}\theta^{00}$, etc., which decreases, as one sees, with still greater rapidity than in the first case.

Knowing thus the law of progression of the parameters $n$, let us come to that of the coefficients $A$ and $B$ in the successive transforms. Let, for brevity,
\[ k=\frac{c^{2}+2n+n^{2}}{(1+b)^{2}(1+n)^{2}}=\frac{(1+n)^{2}-b^{2}}{(1+b)^{2}(1+n)^{2}}; \]
and let $k^{0}$ denote a quantity composed of $c^{0}$, $b^{0}$, $n^{0}$, as $k$ is of $c$, $b$, $n$, and so on, we shall have
\[ B^{0}=\tfrac{1}{2}kB,\quad B^{00}=\tfrac{1}{4}kk^{0}B,\quad B^{000}=\tfrac{1}{8}kk^{0}k^{00}B,\quad\text{etc.} \]
But since, after a certain number $\mu$ of terms, the $c^{\mu}$ and $n^{\mu}$ taken from the sequences $c^{0}$, $c^{00}$, etc., $n^{0}$, $n^{00}$, etc., are small enough to be regarded as nil, it is clear that the same will hold of the corresponding term $k^{\mu}$ of the sequence $k$, $k^{0}$, $k^{00}$, etc., and that thus one may with perfect safety put $B^{\mu}=0$. As to the value of $A^{\mu}$, it will be given by the series
\[ A^{\mu}=A+\frac{\tfrac{1}{2}B}{1+n}+\frac{\tfrac{1}{4}kB}{1+n^{0}}+\frac{\tfrac{1}{8}kk^{0}B}{1+n^{00}}+\text{etc.} \]
which must be continued up to the term that would contain $k^{\mu}$, exclusively.

Let us always call $\Phi$ the limit of the angles $\varphi$, $\frac{\varphi^{0}}{2}$, $\frac{\varphi^{00}}{4}$, etc., this limit being deemed attained after the number of terms $\mu$; we shall have $\varphi^{\mu}=2^{\mu}\Phi$, and because $c^{\mu}$ is of the order of negligible quantities, we shall have $H^{\mu}=2^{\mu}A^{\mu}\Phi$.

\origpage{122}
\subsection*{(86).}

It is now easy to have the result of the successive transforms up to the $\mu^{\text{th}}$ inclusively. Let always
\[ K^{\mu}=(1+c^{0})(1+c^{00})\ldots(1+c^{\mu}),\quad\text{or}\quad K^{\mu}=\frac{2\sqrt{c^{0}}}{c}\cdot\frac{2\sqrt{c^{00}}}{c^{0}}\cdot\frac{2\sqrt{c^{000}}}{c^{00}}\cdot\text{etc.}, \]
the value of the function $H$ will be given in general by the formula
\[ \begin{aligned} H=K^{\mu}A^{\mu}\Phi&-\frac{\sqrt{c^{0}}}{c}\cdot\frac{\tfrac{1}{2}B}{1+n}\cdot\frac{\operatorname{arc}\operatorname{tang}(\sqrt{n^{0}}\sin\varphi^{0})}{\sqrt{n^{0}}}\\ &-\frac{\sqrt{c^{0}}}{c}\cdot\frac{\sqrt{c^{00}}}{c^{0}}\cdot\frac{\tfrac{1}{4}kB}{1+n^{0}}\cdot\frac{\operatorname{arc}\operatorname{tang}(\sqrt{n^{00}}\sin\varphi^{00})}{\sqrt{n^{00}}}\\ &-\frac{\sqrt{c^{0}}}{c}\cdot\frac{\sqrt{c^{00}}}{c^{0}}\cdot\frac{\sqrt{c^{000}}}{c^{00}}\cdot\frac{\tfrac{1}{8}kk^{0}B}{1+n^{00}}\cdot\frac{\operatorname{arc}\operatorname{tang}(\sqrt{n^{000}}\sin\varphi^{000})}{\sqrt{n^{000}}}\\ &-\text{etc.} \end{aligned} \]
This series being continued up to the term whose rank is $\mu$ inclusively, the error of the formula, measured by the following terms, can only be of the order $c^{\mu+1}$; so that if one wishes to push the precision only to quantities of the order $c^{\mu}$, there will be one term fewer to calculate, both in the preceding series and in the value of $K$ and in that of $A^{\mu}$.

Moreover, the arcs of a circle which enter into the general formula would be changed into logarithms, if $n$ were of the form $-c^{2}\sin^{2}\theta$; but this change is subject to no difficulty.

When $\varphi=\frac{1}{2}\pi$ or $=$ a multiple of $\frac{1}{2}\pi$, all the arcs of a circle disappear, and we have simply $H=KA^{\mu}\Phi$. Hence the value of the complete function is
\[ H^{1}=KA^{\mu}\cdot\tfrac{1}{2}\pi. \]
Note that the value of $H$ expressed by the successive transforms would terminate of itself, if one had one of the equalities $n=-1+b$, $n^{0}=-1+b^{0}$, $n^{00}=-1+b^{00}$, etc.: the function $H$ would then be reduced indefinitely to the first species.

In order that we may have in general $-1+b^{\mu}=n^{\mu}=-(c^{\mu}\sin\theta^{\mu})^{2}$, it is necessary that $\cot\theta^{\mu}=\sqrt{b^{\mu}}$, or that $F(c^{\mu},\theta^{\mu})=\frac{1}{2}F^{1}(c^{\mu})$. But the successive parameters $n=-c^{2}\sin^{2}\theta$, $n^{0}=-c^{02}\sin^{2}\theta^{0}$, $n^{00}=-c^{002}\sin^{2}\theta^{00}$, etc. being formed according to the same law that connects the amplitudes $\theta$, $\theta^{0}$, $\theta^{00}$, etc., we have $F(c,\theta)=\frac{1+c^{0}}{2}F(c^{0},\theta^{0})$, $F(c^{0},\theta^{0})=\frac{1+c^{00}}{2}F(c^{00},\theta^{00})$, etc.;

\origpage{123}
and in general $F(c,\theta)=\frac{(1+c^{0})(1+c^{00})\ldots(1+c^{\mu})}{2^{\mu}}F(c^{\mu},\theta^{\mu})=\frac{1}{2^{\mu+1}}F^{1}(c)$. The cases which give rise to the reduction of the function $H$ are therefore all those in which the parameter $n$, of the form $-c^{2}\sin^{2}\theta$, is such that we have $F(c,\theta)=\frac{1}{2^{\mu+1}}F^{1}(c)$, $\mu$ being an integer. The same reduction would still take place if we had $F(c,\theta)=\frac{2\nu+1}{2^{\mu+1}}F^{1}(c)$, $2\nu+1$ being any odd number whatever.

The approximation formulas which we have just given for determining the value of the functions $H$ and $H^{1}$ can be applied to all cases, and most often only very few terms will have to be calculated in order to obtain a high degree of precision. However, if the difference $1-c$ were so small that one had to continue the sequence $c^{0}$, $c^{00}$, $c^{000}$, etc. quite far in order to arrive at a negligible term $c^{\mu}$, it might be preferable to follow the second method, that is, to make the transformations in an inverse order, as we have done in such a case with respect to the functions of the first and second species.

\subsection*{(87).}

And first, if the difference $1-c$ is already small enough to be neglected, we may integrate the formula $H=\displaystyle\int\left(A+\frac{B\sin^{2}\varphi}{1+n\sin^{2}\varphi}\right)\frac{d\varphi}{\Delta}$ immediately, by putting $\cos\varphi$ in place of $\Delta$, which will give
\[ H=\left(A+\frac{B}{1+n}\right)\log\operatorname{tang}(45^{\circ}+\tfrac{1}{2}\varphi)-\frac{B}{1+n}\cdot\frac{\operatorname{arc}\operatorname{tang}(\sqrt{n}\sin\varphi)}{\sqrt{n}}. \]
This integral is rigorous when $c=1$; but if there is a difference, however small, between 1 and $c$, it cannot be applied without error to values of $\varphi$ greater than a certain limit. Indeed, the quantity $\Delta$, which in general is $\sqrt{(\cos^{2}\varphi+b^{2}\sin^{2}\varphi)}$, reduces to $\cos\varphi$ only when $\cot\varphi$ is deemed much greater than $b$. For this reason it is fitting to seek a formula which gives the value of $H$ for every value of $\varphi$, on the sole hypothesis that one may neglect the terms affected by the factor $b$ raised to the square or to a higher power.

For this purpose we shall suppose that there has been made in the function $H$

\origpage{124}
a preparation relative to the parameter. If $n$ is negative, we always suppose it either of the form $-c^{2}\sin^{2}\theta$, or of the form $-1+b^{2}\sin^{2}\theta$; if the latter case occurs, the function $H$ will be changed into another in which the parameter is positive, which will be done by the formula of no.~51.

This being so, to have the value of the function $H$, I put it under the form
\[ H=\left(A+\frac{B}{1+n}\right)F(c,\varphi)-\frac{B}{1+n}P, \]
and it remains to determine $P$ by the formula
\[ P=\int\frac{d\varphi\cos^{2}\varphi}{(1+n\sin^{2}\varphi)\Delta}. \]
I observe next that we have $\frac{\cos\varphi}{\Delta}=1-\frac{b^{2}\sin^{2}\varphi}{\Delta(\cos\varphi+\Delta)}$; therefore if we put
\[ Q=\int\frac{2d\varphi\cos\varphi\sin^{2}\varphi}{(1+n\sin^{2}\varphi)\Delta(\cos\varphi+\Delta)}; \]
we shall have
\[ P=\frac{\operatorname{arc}\operatorname{tang}(\sqrt{n}\sin\varphi)}{\sqrt{n}}-\tfrac{1}{2}b^{2}Q. \]
Everything is thus reduced to finding an approximate value of the integral $Q$. For this purpose, I observe that we have $\cos\varphi+\Delta<2\Delta$, and $\cos\varphi+\Delta>2\cos\varphi$; therefore if we put
\[ \begin{aligned} Q'&=\int\frac{d\varphi\sin^{2}\varphi}{(1+n\sin^{2}\varphi)\Delta}\\ Q''&=\int\frac{d\varphi\cos\varphi\sin^{2}\varphi}{(1+n\sin^{2}\varphi)(1-c^{2}\sin^{2}\varphi)}, \end{aligned} \]
we shall have $Q<Q'$ and $Q>Q''$. Now in the first place, the combination of the integrals $P$ and $Q'$ gives $(1+n)Q'+P=\displaystyle\int\frac{d\varphi}{\Delta}$, whence results $Q'=\frac{F(c,\varphi)-P}{1+n}$, therefore we have
\[ P>\frac{\operatorname{arc}\operatorname{tang}(\sqrt{n}\sin\varphi)}{\sqrt{n}}-\frac{\tfrac{1}{2}b^{2}}{1+n}F(c,\varphi)+\frac{\tfrac{1}{2}b^{2}}{1+n}P, \]
or
\[ P>\left(1+\frac{\tfrac{1}{2}b^{2}}{1+n}\right)\frac{\operatorname{arc}\operatorname{tang}(\sqrt{n}\sin\varphi)}{\sqrt{n}}-\frac{\tfrac{1}{2}b^{2}}{1+n}F(c,\varphi) \]
We find next by direct integration,
\[ (c^{2}+n)Q''=\frac{1}{2c}\log\left(\frac{1+c\sin\varphi}{1-c\sin\varphi}\right)-\frac{1}{\sqrt{n}}\operatorname{arc}\operatorname{tang}(\sqrt{n}\sin\varphi); \]

\origpage{125}
therefore we have
\[ P<\left(1+\frac{\tfrac{1}{2}b^{2}}{1+n}\right)\frac{\operatorname{arc}\operatorname{tang}(\sqrt{n}\sin\varphi)}{\sqrt{n}}-\frac{\tfrac{1}{4}b^{2}}{1+n}\log\left(\frac{1+c\sin\varphi}{1-c\sin\varphi}\right). \]
The greatest difference between these two limits of $P$ occurs when $\varphi=90^{\circ}$, and as we have in this case $F^{1}(c)=\log\frac{4}{b}$, $\log\left(\frac{1+c}{1-c}\right)=2\log\frac{2}{b}$, this difference $=\frac{b^{2}}{1+n}\cdot\frac{1}{2}\log 2$. It is, as one sees, of the order of the quantities one wishes to neglect. Thus we have with sufficient exactness,
\[ P=\frac{1}{\sqrt{n}}\operatorname{arc}\operatorname{tang}(\sqrt{n}\sin\varphi)-\frac{\tfrac{1}{2}b^{2}}{1+n}F(c,\varphi), \]
which gives the function sought
\[ H=\left(A+\frac{B}{c+n}\right)F(c,\varphi)-\frac{B}{1+n}\cdot\frac{\operatorname{arc}\operatorname{tang}(\sqrt{n}\sin\varphi)}{\sqrt{n}}; \]
and there results for the complete function,
\[ H^{1}=\left(A+\frac{B}{1+n}\right)(1+\tfrac{1}{4}b^{2})\log\frac{4}{b}-\frac{B}{1+n}\cdot\frac{\operatorname{arc}\operatorname{tang}(\sqrt{n})}{\sqrt{n}}. \]
These values of $H$ and of $H^{1}$ are those which it is fitting to employ if quantities of the order $b^{2}$ are negligible, and then there is no need to have recourse to the transformations; but it is fitting to examine particularly the case where $n$ is negative.

The case of $n=-1+b^{2}\sin^{2}\theta$, where $n+1$ would be of the order $b^{2}$, does not allow the preceding formulas to be applied to it with success, which is why we have avoided it by a transformation which makes $n$ positive.

The case of $n=-c^{2}\sin^{2}\theta$ suffers no difficulty, if $\theta$ is not too near $90^{\circ}$, because then $\frac{b^{2}}{1+n}$ remains a very small quantity; only in this case, the expression $\frac{\operatorname{arc}\operatorname{tang}(\sqrt{n}\sin\varphi)}{\sqrt{n}}$ must be changed into this one, $\frac{1}{2c\sin\theta}\log\left(\frac{1+c\sin\theta\sin\varphi}{1-c\sin\theta\sin\varphi}\right)$. But if at the same time $\theta$ is very near $90^{\circ}$, the denominator $1+n$, equivalent to $b^{2}+c^{2}\cos^{2}\theta$, becomes very small; yet as long as $\cos\theta$ is much greater than $b$, the fraction $\frac{b^{2}}{1+n}$ or $\frac{b^{2}}{b^{2}+c^{2}\cos^{2}\theta}$ will remain very small, and the

\origpage{126}
preceding method will be applicable; but if $\cos\theta$ is a very small quantity of the order $b$ or of a higher order, the fraction $\frac{b^{2}}{1+n}$ ceases to be negligible. To obviate this inconvenience, one must have recourse to the formula of no.~52, by means of which every function of the third species whose parameter is $n=-c^{2}\sin^{2}\theta$ can be transformed into another function whose parameter is $n'=-c^{2}\sin^{2}\lambda$, provided that between the angles $\theta$ and $\lambda$ we have the relation $b\operatorname{tang}\theta\operatorname{tang}\lambda=1$. Then we have $\frac{b^{2}}{1+n}\cdot\frac{b^{2}}{1+n'}=b^{2}$, so that of the two factors $\frac{b^{2}}{1+n}$, $\frac{b^{2}}{1+n'}$ there will always be one smaller than $b$; whence it follows that, with the aid of a transformation, if it is necessary, the error of the formula will not exceed quantities of the order $b$; and because in the most unfavourable case, which is that of $1+n=1+n'=b$, the value of $H^{1}$ becomes very large and is of the order $b^{-1}$, or more exactly of the order $b^{-1}\log\left(\frac{4}{b}\right)$ which differs very little from it, it follows that, the absolute error being of the order $b$, the relative error is in effect of the order $b^{2}$, as in the cases where $n$ is positive.

We have entered into these details to prove that when $1-c$ is small enough to be negligible, the formulas found for $H$ and $H^{1}$ can give a sufficient approximation, without requiring any transformations other than those indicated by articles 51 and 52. But the method which we are going to expound has the advantage of giving an indefinite approximation, which applies generally to all cases.

\subsection*{(88).}

The equation found in no.~84 between $H$ and $H^{0}$ applies to the two consecutive functions $H'$ and $H$, and gives the transform
\[ H=\frac{2}{1+c}H'+\frac{\tfrac{1}{2}B'}{1+n'}\cdot\int\frac{d\varphi\cos\varphi}{1+n\sin^{2}\varphi}; \]
we shall have at the same time
\[ H'=\int\left(A'+\frac{B'\sin^{2}\varphi'}{1+n'\sin^{2}\varphi'}\right)\frac{d\varphi'}{\Delta'}, \]
and the coefficients $n'$, $A'$, $B'$ must be deduced from the equations
\[ n=\frac{n'(n'+c'^{2})}{(1+b')^{2}(1+n')},\quad B=\frac{B'}{2}\cdot\frac{(1+n')^{2}-b'^{2}}{(1+b')^{2}(1+n')^{2}},\quad A=A'+\frac{\tfrac{1}{2}B'}{1+n'}. \]

\origpage{127}
We know how $c'$, $b'$, $\varphi'$ are calculated by means of $c$ and $\varphi$; it is therefore only a matter of examining the law of formation of the quantities $n'$, $A'$, $B'$. And first, to deduce $n'$ from $n$, we have the formula
\[ n'=\frac{2}{(1+c)^{2}}[n-c+\sqrt{(1+n)(c^{2}+n)}]; \]
whence one sees that we must set aside the case in which we should have $n=-1+b^{2}\sin^{2}\theta$; for if this case occurred, the parameter $n'$ would become imaginary, and we should thus fall into a greater difficulty than the one we wish to resolve. Fortunately there is no inconvenience in making this exception, since by the formula of no.~51, every proposed function of the third species whose parameter is of the form $-1+b^{2}\sin^{2}\theta$ can be transformed into another whose parameter is positive. We can therefore limit ourselves to considering two principal cases, that where $n$ is positive and that where $n$ is of the form $-c^{2}\sin^{2}\theta$. In these two cases, one can assure oneself that the value of $n'$ will always be of the same form as that of $n$, so that nothing can stop the calculation of the successive transforms.

\subsection*{(89).}

Let then in the first place $n=-c^{2}\sin^{2}\theta$, and likewise $n'=-c'^{2}\sin^{2}\theta'$; we shall have by the preceding formula,
\[ \sin^{2}\theta'=\tfrac{1}{2}+\tfrac{1}{2}c\sin^{2}\theta-\tfrac{1}{2}\cos\theta\Delta(\theta). \]
This is the same law according to which the amplitude $\varphi'$ is deduced from $\varphi$; it may be represented more simply by the equation $\sin(2\theta'-\theta)=c\sin\theta$. By virtue of this law, the angles $\theta$, $\theta'$, $\theta''$, etc. form a decreasing sequence which approaches without end a limit $\Theta$, and which sensibly attains it after a small number of terms. The sequence of parameters $n$, $n'$, $n''$, etc. therefore likewise converges toward the limit $-\sin^{2}\Theta$, since moreover the sequence $c$, $c'$, $c''$, etc. has unity for its limit.

In the second place, if $n$ is positive, let $1+n=\frac{b^{2}}{\sin^{2}\lambda}$, $\lambda$ being a very small angle determined by the value $\sin\lambda=\frac{b}{\sqrt{(1+n)}}$, so that if we put $n=\cot^{2}\theta$, we should have $\sin\lambda=b\sin\theta$. This being so, the value of $n'$ will be given by the equation
\[ 1+n'=b'\cot^{2}\tfrac{1}{2}\lambda; \]

\origpage{128}
and if we likewise put $1+n'=\frac{b'^{2}}{\sin^{2}\lambda'}$, the angle $\lambda'$ will be deduced from $\lambda$ by the very simple formula
\[ \sin\lambda'=\sqrt{b'}\cdot\operatorname{tang}\tfrac{1}{2}\lambda. \]
We shall deduce likewise $\lambda''$ from $\lambda'$, $\lambda'''$ from $\lambda''$, etc., so that we shall form the sequence of parameters $n$, $n'$, $n''$, etc. according to the equations $1+n'=\frac{b'^{2}}{\sin^{2}\lambda'}$, $1+n''=\frac{b''^{2}}{\sin^{2}\lambda''}$, etc.

Let us remark that, starting from the angle $\lambda$, which is very small, the terms of the series $\lambda$, $\lambda'$, $\lambda''$, etc. diminish continually, but in such a way that the ratios $\frac{\sin\lambda}{b}$, $\frac{\sin\lambda'}{b'}$, $\frac{\sin\lambda''}{b''}$, etc. converge toward a constant quantity which will be their limit. Indeed, from the preceding values one deduces $\frac{\sin\lambda'}{b'}=\frac{\sin\lambda}{b}\cdot\frac{1+c}{1+\cos\lambda}$; we should have in the same way $\frac{\sin\lambda''}{b''}=\frac{\sin\lambda'}{b'}\cdot\frac{1+c'}{1+\cos\lambda'}$, etc. Hence after a small number of terms $\frac{\sin\lambda^{\mu}}{b^{\mu}}$ will be constant, and consequently also $1+n^{\mu}$; whence one sees that the sequence of parameters $n$, $n'$, $n''$, etc. converges, as in the first case, toward a limit which it sensibly attains at the end of a few terms.

The very simple means which we have just indicated for calculating the sequence of parameters $n$, $n'$, $n''$, etc., when the first, $n$, is positive, can also be applied without difficulty to the case where $n$ is of the form $-c^{2}\sin^{2}\theta$; so that we may suppose it employed generally in all cases; and on this supposition, we are going to continue the calculations necessary to arrive at the expression of the value of $H$.

\subsection*{(90).}

We must first have the law of the coefficients $B^{\mu}$, $A^{\mu}$; for this, if we put
\[ h'=\frac{(1+b')^{2}(1+n')^{2}}{(1+n')^{2}-b'^{2}},\quad h''=\frac{(1+b'')^{2}(1+n'')^{2}}{(1+n'')^{2}-b''^{2}},\quad\text{etc.}, \]
we shall have successively
\[ B'=2h'B,\quad B''=2h''B'=4h'h''B,\quad B'''=8h'h''h'''B,\quad\text{etc.} \]

\origpage{129}
But from the equations $1+n'=b'\cot^{2}\frac{1}{2}\lambda$, $1+n=\frac{b^{2}}{\sin^{2}\lambda}$, $b^{2}=\frac{4b'}{(1+b')^{2}}$, we draw
\[ h'=\frac{1+n'}{1+n}\cdot\frac{1}{\cos\lambda}; \]
we shall have likewise $h''=\frac{1+n''}{1+n'}\cdot\frac{1}{\cos\lambda'}$, $h'''=\frac{1+n'''}{1+n''}\cdot\frac{1}{\cos\lambda''}$, etc.; therefore
\[ B'=\frac{2B(1+n')}{(1+n)\cos\lambda},\quad B''=\frac{4B(1+n'')}{(1+n)\cos\lambda\cos\lambda'},\quad B'''=\frac{8B(1+n''')}{(1+n)\cos\lambda\cos\lambda'\cos\lambda''},\quad\text{etc.} \]
and in general $B^{\mu}=\frac{2^{\mu}B(1+n^{\mu})}{(1+n)\cos\lambda\cos\lambda'\ldots\cos\lambda^{\mu-1}}$.

Finally from the equations $A'=A-\frac{\frac{1}{2}B'}{1+n'}$, $A''=A'-\frac{\frac{1}{2}B''}{1+n''}$, etc., we deduce generally
\[ A^{\mu}=A-\frac{\tfrac{1}{2}B'}{1+n'}-\frac{\tfrac{1}{2}B''}{1+n''}-\ldots-\frac{\tfrac{1}{2}B^{\mu}}{1+n^{\mu}}; \]
or, making the substitutions,
\[ A^{\mu}=A-\frac{B}{1+n}\left(\frac{1}{\cos\lambda}+\frac{2}{\cos\lambda\cos\lambda'}+\frac{4}{\cos\lambda\cos\lambda'\cos\lambda''}\ldots+\frac{2^{\mu-1}}{\cos\lambda\cos\lambda'\ldots\cos\lambda^{\mu-1}}\right). \]
These values will make known that of the coefficient $A^{\mu}+\frac{B^{\mu}}{1+n^{\mu}}$ which enters into the expression of $H^{\mu}$. Let this coefficient $=L^{\mu}$, we shall have
\[ \begin{aligned} L^{\mu}=A&+\frac{2^{\mu}B}{(1+n)\cos\lambda\cos\lambda'\ldots\cos\lambda^{\mu-1}}\\ &-\frac{B}{1+n}\left(\frac{1}{\cos\lambda}+\frac{2}{\cos\lambda\cos\lambda'}+\frac{4}{\cos\lambda\cos\lambda'\cos\lambda''}\ldots+\frac{2^{\mu-1}}{\cos\lambda\cos\lambda'\ldots\cos\lambda^{\mu-1}}\right); \end{aligned} \]
But as the series contained in this formula is divergent, it will be fitting to give it another form. For this purpose I observe that the preceding formula gives
\[ L^{\mu+1}-L^{\mu}=\frac{2^{\mu}B(1-\cos\lambda^{\mu})}{(1+n)\cos\lambda\cos\lambda'\ldots\cos\lambda^{\mu}}. \]
Now we have seen that as $\mu$ increases, the quantity $\frac{\sin\lambda^{\mu}}{b^{\mu}}$ approaches

\origpage{130}
more and more a certain limit, and that thus $\sin\lambda^{\mu}$ is in general of the order $b^{\mu}$; therefore $1-\cos\lambda^{\mu}$ is of the order $(b^{\mu})^{2}$, or $b^{\mu+1}$; therefore the difference $L^{\mu+1}-L^{\mu}$ is a very small quantity of the order $2^{\mu}b^{\mu+1}$, intermediate between $b^{\mu}$ and $b^{\mu+1}$, but nearer to the latter. We shall therefore have for the expression of $L^{\mu}$ this highly convergent series
\[ \begin{aligned} L^{\mu}=A+\frac{B}{1+n}&+\frac{B(1-\cos\lambda)}{(1+n)\cos\lambda}+\frac{2B(1-\cos\lambda')}{(1+n)\cos\lambda\cos\lambda'}\\ &+\frac{4B(1-\cos\lambda'')}{(1+n)\cos\lambda\cos\lambda'\cos\lambda''}+\text{etc.}, \end{aligned} \]
which must be continued up to the term \ldots\ $\frac{2^{\mu-1}B(1-\cos\lambda^{\mu-1})}{(1+n)\cos\lambda\cos\lambda'\cos\lambda''\ldots\cos\lambda^{\mu-1}}$ inclusively, if one wishes to neglect only the quantities of the order $2^{\mu}b^{\mu+1}$.

\subsection*{(91).}

This being so, the successive transforms being
\[ \begin{aligned} H&=(1+b')H'+\frac{\tfrac{1}{2}B'}{1+n'}\cdot\frac{\operatorname{arc}\operatorname{tang}(\sqrt{n}\sin\varphi)}{\sqrt{n}}\\ H'&=(1+b'')H''+\frac{\tfrac{1}{2}B''}{1+n''}\cdot\frac{\operatorname{arc}\operatorname{tang}(\sqrt{n'}\sin\varphi')}{\sqrt{n'}}\\ &\text{etc.}, \end{aligned} \]
we may express the function $H$ generally by means of $H^{\mu}$; and as the series are continued to a number of terms $\mu$ sufficient for the value of $H^{\mu}$ to be given with all the necessary exactness by the formula
\[ H^{\mu}=\left(A^{\mu}+\frac{B^{\mu}}{1+n^{\mu}}\right)\log\operatorname{tang}(45^{\circ}+\tfrac{1}{2}\varphi^{\mu})-\frac{B^{\mu}}{1+n^{\mu}}\cdot\frac{\operatorname{arc}\operatorname{tang}(\sqrt{n^{\mu}}\sin\varphi^{\mu})}{\sqrt{n^{\mu}}}, \]
there will result from the substitutions this general value of $H$, where we have put for brevity, $I^{\mu}=(1+b')(1+b'')(1+b''')\ldots(1+b^{\mu})$,

\origpage{131}
\[ \begin{aligned} H=I^{\mu}L^{\mu}\log\operatorname{tang}(45^{\circ}+\tfrac{1}{2}\varphi^{\mu})&-I^{\mu}\cdot\frac{B^{\mu}}{1+n^{\mu}}\cdot\frac{\operatorname{arc}\operatorname{tang}(\sqrt{n^{\mu}}\sin\varphi^{\mu})}{\sqrt{n^{\mu}}}\\ &+\frac{\tfrac{1}{2}B'}{1+n'}\cdot\frac{\operatorname{arc}\operatorname{tang}(\sqrt{n}\sin\varphi)}{\sqrt{n}}\\ &+(1+b')\cdot\frac{\tfrac{1}{2}B''}{1+n''}\cdot\frac{\operatorname{arc}\operatorname{tang}(\sqrt{n'}\sin\varphi')}{\sqrt{n'}}\\ &+(1+b')(1+b'')\cdot\frac{\tfrac{1}{2}B'''}{1+n'''}\cdot\frac{\operatorname{arc}\operatorname{tang}(\sqrt{n''}\sin\varphi'')}{\sqrt{n''}}\\ &\ \vdots\\ &+I^{\mu-1}\cdot\frac{\tfrac{1}{2}B^{\mu}}{1+n^{\mu}}\cdot\frac{\operatorname{arc}\operatorname{tang}(\sqrt{n^{\mu-1}}\sin\varphi^{\mu-1})}{\sqrt{n^{\mu-1}}}. \end{aligned} \]

\subsection*{(92).}

This formula has the inconvenience of offering a divergent series, and it is important to give it another form. For this purpose let us separate the first term $I^{\mu}L^{\mu}\log\operatorname{tang}(45^{\circ}+\tfrac{1}{2}\varphi^{\mu})$, and call the remainder $Z^{\mu}$; let us further denote by $T^{i}$ the quantity $\frac{\operatorname{arc}\operatorname{tang}(\sqrt{n^{i}}\sin\varphi^{i})}{\sqrt{n^{i}}}$; we shall have
\[ Z^{\mu+1}-Z^{\mu}=\frac{B^{\mu}}{1+n^{\mu}}\cdot I^{\mu}T^{\mu}-\frac{B^{\mu+1}}{1+n^{\mu+1}}I^{\mu+1}T^{\mu+1}+\frac{\tfrac{1}{2}B^{\mu+1}}{1+n^{\mu+1}}I^{\mu}T^{\mu}. \]
This equation can be put under the form
\[ Z^{\mu+1}-Z^{\mu}=I^{\mu}T^{\mu}\left(\frac{B^{\mu}}{1+n^{\mu}}+\frac{\tfrac{1}{2}B^{\mu+1}}{1+n^{\mu+1}}\right)-I^{\mu+1}T^{\mu+1}\frac{B^{\mu+1}}{1+n^{\mu+1}}; \]
but if for a moment we put $\Lambda^{i}=\frac{1}{\cos\lambda}\cdot\frac{1}{\cos\lambda'}\ldots\frac{1}{\cos\lambda^{i}}$, we shall have $\frac{B^{\mu}}{1+n^{\mu}}=\frac{2^{\mu}B}{1+n}\Lambda^{\mu-1}$ and $\frac{B^{\mu+1}}{1+n^{\mu+1}}=\frac{2^{\mu+1}B}{1+n}\Lambda^{\mu}$; therefore
\[ Z^{\mu+1}-Z^{\mu}=\frac{2^{\mu+1}B}{1+n}I^{\mu}\Lambda^{\mu}\left[\frac{1+\cos\lambda^{\mu}}{2}T^{\mu}-(1+b^{\mu+1})T^{\mu+1}\right]. \]
Now on the one hand the quantity $I^{\mu}\Lambda^{\mu}$ cannot pass a certain limit not far above unity; on the other hand the differences $1-\cos\lambda^{\mu}$, $T^{\mu}-T^{\mu+1}$ are of the order $b^{\mu+1}$; therefore $Z^{\mu+1}-Z^{\mu}$ is a very small quantity of the order $2^{\mu+1}b^{\mu+1}$ which may be regarded as intermediate between $b^{\mu}$ and $b^{\mu+1}$, but much nearer to the latter.

\origpage{132}
We may therefore express $Z^{\mu}$ by the highly convergent series $Z+(Z'-Z)+(Z''-Z')+\text{etc.}$, and we shall have for the general expression of the function $H$ this formula
\[ \begin{aligned} H=I^{\mu}L^{\mu}\log\operatorname{tang}(45^{\circ}+\tfrac{1}{2}\varphi^{\mu})&-\frac{B}{1+n}\cdot\frac{\operatorname{arc}\operatorname{tang}(\sqrt{n}\sin\varphi)}{\sqrt{n}}\\ &-\frac{2B}{1+n}\cdot\frac{1}{\cos\lambda}\left[(1+b')\frac{\operatorname{arc}\operatorname{tang}(\sqrt{n'}\sin\varphi')}{\sqrt{n'}}-\left(\frac{1+\cos\lambda}{2}\right)\cdot\frac{\operatorname{arc}\operatorname{tang}(\sqrt{n}\sin\varphi)}{\sqrt{n}}\right]\\ &-\frac{4B}{1+n}\cdot\frac{1+b'}{\cos\lambda\cos\lambda'}\left[(1+b'')\frac{\operatorname{arc}\operatorname{tang}(\sqrt{n''}\sin\varphi'')}{\sqrt{n''}}-\left(\frac{1+\cos\lambda'}{2}\right)\cdot\frac{\operatorname{arc}\operatorname{tang}(\sqrt{n'}\sin\varphi')}{\sqrt{n'}}\right]\\ &-\frac{8B}{1+n}\cdot\frac{1+b'\cdot 1+b''}{\cos\lambda\cos\lambda'\cos\lambda''}\left[(1+b''')\frac{\operatorname{arc}\operatorname{tang}(\sqrt{n'''}\sin\varphi''')}{\sqrt{n'''}}-\left(\frac{1+\cos\lambda''}{2}\right)\frac{\operatorname{arc}\operatorname{tang}(\sqrt{n''}\sin\varphi'')}{\sqrt{n''}}\right]\\ &-\text{etc.} \end{aligned} \]

This formula, founded on the one which we have taken for the value of $H^{\mu}$, supposes that the angle $\varphi^{\mu}$ is always less than the limit given by the equation $\cot\varphi^{\mu}=\sqrt{b^{\mu}}$. When $\varphi^{\mu}$ is equal to this limit, or is merely close to it, the supposed value of $H^{\mu}$ is exact up to quantities of the order $b^{\mu}$, and then the series which gives the value of $H$ must be continued up to the $\mu^{\text{th}}$ term $-\frac{2^{\mu}B}{1+n}$, etc. inclusively. In the case where $\varphi^{\mu}$ is remote from the limit in question, which will happen when $\varphi$ does not exceed $90^{\circ}$, as one may always suppose, the series may be calculated only up to the $(\mu-1)^{\text{th}}$ term inclusively, and the result will still be exact up to quantities of the order $b^{\mu}$.

To deduce from the general formula the value of the complete function $H^{1}$, one could put $\varphi=90^{\circ}$, and continue the series, as has just been said, up to the $(\mu-1)^{\text{th}}$ term inclusively. But it is simpler to put $\varphi=2^{\mu-1}\frac{\pi}{2}=2^{\mu-2}\pi$, which will give successively $\varphi'=2^{\mu-3}\pi$, $\varphi''=2^{\mu-4}\pi\ldots\varphi^{\mu-1}=\tfrac{1}{2}\pi$, $\cot\varphi^{\mu}=\sqrt{b^{\mu}}$. We thus arrive at the limit of the values of $\varphi^{\mu}$, and neglecting quantities of the order $b^{\mu}$, we shall have $\sin\varphi^{\mu}=1$, $\log\operatorname{tang}(45^{\circ}+\tfrac{1}{2}\varphi^{\mu})=\tfrac{1}{2}\log\left(\frac{4}{b^{\mu}}\right)$, and by the substitution of all these values, we find
\[ \begin{aligned} 2^{\mu-1}H^{1}=I^{\mu}L^{\mu}\tfrac{1}{2}\log\frac{4}{b^{\mu}}&-\frac{2^{\mu-1}B}{1+n}\cdot\frac{1+b'\cdot 1+b''\ldots 1+b^{\mu-1}}{\cos\lambda\cos\lambda'\ldots\cos\lambda^{\mu-2}}\cdot\frac{\operatorname{arc}\operatorname{tang}\sqrt{n^{\mu-1}}}{\sqrt{n^{\mu-1}}}\\ &-\frac{2^{\mu}B}{1+n}\left(\frac{1+b'\cdot 1+b''\ldots 1+b^{\mu-1}}{\cos\lambda\cos\lambda'\ldots\ldots\cos\lambda^{\mu-1}}\left[(1+b^{\mu})\frac{\operatorname{arc}\operatorname{tang}\sqrt{n^{\mu}}}{\sqrt{n^{\mu}}}-\left(\frac{1+\cos\lambda^{\mu-1}}{2}\right)\frac{\operatorname{arc}\operatorname{tang}\sqrt{n^{\mu-1}}}{\sqrt{n^{\mu-1}}}\right]. \right. \end{aligned} \]
\ednote{The parenthesis opened before the second fraction of the last line is not closed in the print.}

\origpage{133}
But the quantities $1-\cos\lambda^{\mu-1}$, $n^{\mu}-n^{\mu-1}$ are of the order $b^{\mu}$, and consequently negligible; we have therefore more simply
\[ H^{1}=I^{\mu}\left[\frac{L^{\mu}}{2^{\mu}}\log\frac{4}{b^{\mu}}-\frac{B}{1+n}\cdot\frac{1}{\cos\lambda}\cdot\frac{1}{\cos\lambda'}\ldots\frac{1}{\cos\lambda^{\mu-2}}\cdot\frac{\operatorname{arc}\operatorname{tang}\sqrt{n^{\mu-1}}}{\sqrt{n^{\mu-1}}}\right], \]
a formula in which one must substitute the known value of $L^{\mu}$ and that of $I^{\mu}=(1+b')(1+b'')\ldots(1+b^{\mu-1})=\frac{2\sqrt{b'}}{b}\cdot\frac{2\sqrt{b''}}{b'}\cdot\frac{2\sqrt{b'''}}{b''}$, etc., this product being continued until the number of its factors is $\mu-1$.

Such are the simplest formulas by which one can calculate the values of the functions $H$ and $H^{1}$, when the modulus $c$ is extremely near unity; they will give whatever degree of approximation one wishes to obtain, by sufficiently prolonging the series of which they are composed.

\subsection*{(93).}

If we now compare this second method with the first, we shall not fail to remark that the latter is subject to no exception, and that it requires neither considerations of limits nor transformations; it is therefore preferable to employ the first method, even in cases where $c$ would be very near unity, since this method then has no other inconvenience than that of employing a few more terms, and the number of these terms can never be very great.

We shall remark further that the formulas of the first method may be simplified in some respects, by employing subsidiary angles $\lambda^{0}$, $\lambda^{00}$, $\lambda^{000}$, etc. to calculate both the series of parameters $n^{0}$, $n^{00}$, etc., and that of the quantities $k$, $k^{0}$, $k^{00}$, etc., after the example of what we have done in the second method. For this purpose we shall have the successive formulas
\[ \begin{aligned} 1+n&=b\cot^{2}\tfrac{1}{2}\lambda^{0} &&\ldots\ldots &\operatorname{tang}\tfrac{1}{2}\lambda^{0}&=\sqrt{\left(\frac{b}{1+n}\right)}\\ 1+n^{0}&=b^{0}\cot^{2}\tfrac{1}{2}\lambda^{00}=\frac{b^{02}}{\sin^{2}\lambda^{0}} &&\ldots\ldots &\operatorname{tang}\tfrac{1}{2}\lambda^{00}&=\frac{\sin\lambda^{0}}{\sqrt{b^{0}}}\\ 1+n^{00}&=b^{00}\cot^{2}\tfrac{1}{2}\lambda^{000}=\frac{b^{002}}{\sin^{2}\lambda^{00}} &&\ldots\ldots &\operatorname{tang}\tfrac{1}{2}\lambda^{000}&=\frac{\sin\lambda^{00}}{\sqrt{b^{00}}}\\ 1+n^{000}&=b^{000}\cot^{2}\tfrac{1}{2}\lambda^{0000}=\frac{b^{0002}}{\sin^{2}\lambda^{000}} &&\ldots\ldots &\operatorname{tang}\tfrac{1}{2}\lambda^{0000}&=\frac{\sin\lambda^{000}}{\sqrt{b^{000}}}\\ &\text{etc.} &&& &\text{etc.} \end{aligned} \]

\origpage{134}
from which we deduce
\[ \begin{aligned} k&=\frac{1+n^{0}}{1+n}\cos\lambda^{0} &&\ldots\ldots & \frac{k}{1+n^{0}}&=\frac{1}{1+n}\cos\lambda^{0}\\ k^{0}&=\frac{1+n^{00}}{1+n^{0}}\cos\lambda^{00} &&\ldots\ldots & \frac{kk^{0}}{1+n^{00}}&=\frac{1}{1+n}\cos\lambda^{0}\cos\lambda^{00}\\ k^{00}&=\frac{1+n^{000}}{1+n^{00}}\cos\lambda^{000} &&\ldots\ldots & \frac{kk^{0}k^{00}}{1+n^{000}}&=\frac{1}{1+n}\cos\lambda^{0}\cos\lambda^{00}\cos\lambda^{000}\\ &\text{etc.} &&& &\text{etc.} \end{aligned} \]
Finally, if we put
\[ \begin{aligned} A^{\mu}=A+\frac{\tfrac{1}{2}B}{1+n}(1&+\tfrac{1}{2}\cos\lambda^{0}+\tfrac{1}{4}\cos\lambda^{0}\cos\lambda^{00}\ldots\\ &\ldots+\frac{1}{2^{\mu}}\cos\lambda^{0}\cos\lambda^{00}\ldots\cos\lambda^{\mu}), \end{aligned} \]
and always take $K^{\mu}=(1+c^{0})(1+c^{00})\ldots(1+c^{\mu})$, the value of $H$ will be
\[ \begin{aligned} H=K^{\mu}A^{\mu}\Phi&-\frac{\sqrt{c^{0}}}{2c}\cdot\frac{B}{1+n}\cdot\frac{\operatorname{arc}\operatorname{tang}(\sqrt{n^{0}}\sin\varphi^{0})}{\sqrt{n^{0}}}\\ &-\frac{\sqrt{c^{0}}}{2c}\cdot\frac{\sqrt{c^{00}}}{2c^{0}}\cdot\frac{B}{1+n}\cdot\cos\lambda^{0}\cdot\frac{\operatorname{arc}\operatorname{tang}(\sqrt{n^{00}}\sin\varphi^{00})}{\sqrt{n^{00}}}\\ &-\frac{\sqrt{c^{0}}}{2c}\cdot\frac{\sqrt{c^{00}}}{2c^{0}}\cdot\frac{\sqrt{c^{000}}}{2c^{00}}\cdot\frac{B}{1+n}\cdot\cos\lambda^{0}\cos\lambda^{00}\cdot\frac{\operatorname{arc}\operatorname{tang}(\sqrt{n^{000}}\sin\varphi^{000})}{\sqrt{n^{000}}}\\ &-\text{etc.} \end{aligned} \]
The limit of $n^{\mu}$ being zero in all cases, that of $\sin\lambda^{\mu}$ will be $b^{\mu}$, and consequently that of $\cos\lambda^{\mu}$ will be $c^{\mu}$. From this one sees how convergent this series is.

\section*{On the most general cases in which the reduction of elliptic functions of the third species can be carried out.}

\subsection*{(94).}

If we differentiate with respect to $n$ each member of the equation $\Pi=\displaystyle\int\frac{d\varphi}{(1+n\sin^{2}\varphi)\Delta}$, we shall have
\[ n\frac{d\Pi}{dn}=\int\frac{d\varphi}{(1+n\sin^{2}\varphi)^{2}\Delta}-\int\frac{d\varphi}{(1+n\sin^{2}\varphi)\Delta}: \]

\origpage{135}
now by the formula of no.\ 9 we find, putting $\alpha=(1+n)\left(1+\frac{c^{2}}{n}\right)$,
\[ \begin{aligned} \frac{2\alpha}{n}\int\frac{d\varphi}{(1+n\sin^{2}\varphi)^{2}\Delta}=\frac{\Delta\sin\varphi\cos\varphi}{1+n\sin^{2}\varphi}&-\frac{c^{2}}{n^{2}}\int(1+n\sin^{2}\varphi)\frac{d\varphi}{\Delta}\\ &+\left(1+\frac{2+2c^{2}}{n}+\frac{3c^{2}}{n^{2}}\right)\int\frac{d\varphi}{(1+n\sin^{2}\varphi)\Delta}. \end{aligned} \]
Combining these two equations with each other, and observing that we have
\[ c^{2}\int(1+n\sin^{2}\varphi)\frac{d\varphi}{\Delta}=(c^{2}+n)F-nE, \]
we shall arrive at the following equation, in which only $n$ is to be regarded as variable in $\Pi$ and in $\alpha$,
\[ \tfrac{1}{2}\Pi\,d\alpha+\alpha\,d\Pi=-\tfrac{1}{2}c^{2}F\cdot\frac{dn}{n^{2}}-\tfrac{1}{2}(F-E)\frac{dn}{n}+\tfrac{1}{2}\Delta\sin\varphi\cos\varphi\cdot\frac{dn}{1+n\sin^{2}\varphi}. \]
Multiplying both sides by $\alpha^{-\frac{1}{2}}$ and integrating, we shall have
\[ \sqrt{\alpha}\cdot\Pi=-\tfrac{1}{2}c^{2}F\int\frac{dn}{n^{2}\sqrt{\alpha}}-\tfrac{1}{2}(F-E)\int\frac{dn}{n\sqrt{\alpha}}+\tfrac{1}{2}\Delta\sin\varphi\cos\varphi\int\frac{dn}{(1+n\sin^{2}\varphi)\sqrt{\alpha}}. \]
Now, in order to carry out the integrations, we must consider successively the three cases of art.\ 50.

\section*{First case $n=\cot^{2}\theta$.}

\subsection*{(95).}

Then we have $dn=-\frac{2\cos\theta}{\sin^{3}\theta}\,d\theta$, $\sqrt{\alpha}=\frac{\Delta(b,\theta)}{\sin\theta\cos\theta}$, and the general equation becomes
\[ \begin{aligned} \frac{\Delta(b,\theta)}{\sin\theta\cos\theta}\Pi=c^{2}F\int\frac{d\theta\sin^{2}\theta}{\cos^{2}\theta\Delta(b,\theta)}&+(F-E)\int\frac{d\theta}{\Delta(b,\theta)}\\ &-\Delta\sin\varphi\cos\varphi\int\frac{d\theta\cos^{2}\theta}{(\sin^{2}\theta+\sin^{2}\varphi\cos^{2}\theta)\Delta(b,\theta)}. \end{aligned} \]
The integral $\displaystyle\int\frac{d\theta}{\Delta(b,\theta)}$ is represented as usual by $F(b,\theta)$; and according to the known reductions, we have
\[ \int\frac{c^{2}\,d\theta\sin^{2}\theta}{\cos^{2}\theta\Delta(b,\theta)}=\frac{\sin\theta}{\cos\theta}\Delta(b,\theta)-E(b,\theta); \]
finally, if for brevity we put $\cot^{2}\varphi=n'$, the integral $\Delta\sin\varphi\cos\varphi\displaystyle\int\frac{d\theta\cos^{2}\theta}{(\sin^{2}\theta+\sin^{2}\varphi\cos^{2}\theta)\Delta(b,\theta)}$ can be decomposed thus:
\[ -\frac{\Delta\sin\varphi}{\cos\varphi}\int\frac{d\theta}{\Delta(b,\theta)}+\frac{\Delta}{\sin\varphi\cos\varphi}\int\frac{d\theta}{(1+n'\sin^{2}\theta)\Delta(b,\theta)}, \]

\origpage{136}
and in this form it can be represented by
\[ -\frac{\Delta\sin\varphi}{\cos\varphi}F(b,\theta)+\frac{\Delta}{\sin\varphi\cos\varphi}\Pi(n',b,\theta). \]
Substituting all these values, and denoting, for greater clarity, by $\Pi(n,c,\varphi)$, $\Delta(c,\varphi)$, $F(c,\varphi)$, $E(c,\varphi)$ the functions $\Pi$, $\Delta$, $F$, $E$, we shall have, for the first case, the following general formula:
\[ \begin{aligned} &\frac{\Delta(b,\theta)}{\sin\theta\cos\theta}\Pi(n,c,\varphi)+\frac{\Delta(c,\varphi)}{\sin\varphi\cos\varphi}\Pi(n',b,\theta)\\ &=\left\{\begin{aligned} &C+\frac{\sin\theta}{\cos\theta}\Delta(b,\theta)F(c,\varphi)+\frac{\sin\varphi}{\cos\varphi}\Delta(c,\varphi)F(b,\theta)\\ &+F(c,\varphi)F(b,\theta)-F(c,\varphi)E(b,\theta)-E(c,\varphi)F(b,\theta) \end{aligned}\right.\qquad(i'). \end{aligned} \]
The constant $C$, added to this formula, might be a function of $\varphi$, since it results from an integration in which $\varphi$ was regarded as constant; but as the form of the integral is such that the quantities $\varphi$ and $\theta$ can be interchanged, provided that $c$ and $b$ are interchanged likewise, it follows that $C$ contains neither $\theta$ nor $\varphi$, and that $C$ is thus an absolute constant.

To determine its value, let us take a particular case, and suppose $\varphi$ and $\theta$ both infinitely small. This supposition gives, rejecting the infinitely small quantities of the second order, $F(c,\varphi)=E(c,\varphi)=\varphi$, $F(b,\theta)=E(b,\theta)=\theta$, $\Delta(c,\varphi)=\Delta(b,\theta)=1$, $\Pi(n,c,\varphi)=\displaystyle\int\frac{d\varphi}{1+n\sin^{2}\varphi}=\frac{\operatorname{arc}\operatorname{tang}\varphi\sqrt{(1+n)}}{\sqrt{(1+n)}}=\theta\operatorname{arc}\operatorname{tang}\frac{\varphi}{\theta}$, and similarly $\Pi(n',b,\theta)=\varphi\operatorname{arc}\operatorname{tang}\frac{\theta}{\varphi}$. Substituting all these values, there will result
\[ C=\operatorname{arc}\operatorname{tang}\frac{\varphi}{\theta}+\operatorname{arc}\operatorname{tang}\frac{\theta}{\varphi}=\tfrac{1}{2}\pi. \]

\subsection*{(96).}

If we wish to apply the formula of the preceding art.\ to the case where $\varphi=\tfrac{1}{2}\pi$, we must, in order to avoid the infinite quantities which occur in both members, put $\varphi=\tfrac{1}{2}\pi-\omega$, and suppose $\omega$ infinitely small. We shall thus have $n'=\cot^{2}\varphi=\operatorname{tang}^{2}\omega=\omega^{2}$; and since $n'$ is infinitely small, the function $\Pi(n',b,\theta)$ will be expressed in this manner
\[ \Pi(n',b,\theta)=\int\frac{d\theta}{(1+n'\sin^{2}\theta)\Delta(b,\theta)}=\int\frac{d\theta}{\Delta(b,\theta)}-n'\int\frac{d\theta\sin^{2}\theta}{\Delta(b,\theta)}: \]

\origpage{137}
we shall therefore have
\[ \frac{\Delta}{\sin\varphi\cos\varphi}\Pi(n',b,\theta)-\frac{\Delta\sin\varphi}{\cos\varphi}F(b,\theta)=\omega\Delta F(b,\theta)-\omega\Delta\int\frac{d\theta\sin^{2}\theta}{\Delta(b,\theta)}, \]
from which one sees that the first member reduces to zero when we put $\omega=0$, and that the supposition $\varphi=\tfrac{1}{2}\pi$ thus gives this formula
\[ \frac{\Delta(b,\theta)}{\sin\theta\cos\theta}\Pi^{1}(n,c)=\tfrac{1}{2}\pi+\left.\begin{aligned} &\frac{\sin\theta}{\cos\theta}\Delta(b,\theta)F^{1}(c)+F^{1}(c)F(b,\theta)\\ &-F^{1}(c)E(b,\theta)-E^{1}(c)F(b,\theta) \end{aligned}\right\}\;(k'). \]
Whenever, therefore, the parameter $n$ is positive, the complete function of the third species $\Pi^{1}(n,c)$ can always be determined by functions of the first and second species, namely, by the complete functions $F^{1}(c)$, $E^{1}(c)$, which relate to the modulus $c$, and by the incomplete functions $F(b,\theta)$, $E(b,\theta)$, which relate to the complementary modulus $b$, and whose amplitude $\theta$ is deduced from the parameter $n$ by the equation $\cot\theta=\sqrt{n}$.

This theorem is very important in the theory of elliptic functions, and its use will be all the more extensive in that, in many problems of geometry or mechanics which depend on elliptic functions, one often needs only the definite or complete integrals. We shall give various applications of it below.

\subsection*{(97).}

If in the equation $(k')$ we put $\frac{c^{2}}{n}$ in place of $n$, and put $\frac{c^{2}}{n}=\cot^{2}\lambda$, in order to put $\lambda$ at the same time in place of $\theta$, we shall have
\[ \begin{aligned} \frac{\Delta(b,\lambda)}{\sin\lambda\cos\lambda}\Pi^{1}\left(\frac{c^{2}}{n},c\right)=\tfrac{1}{2}\pi&+\frac{\sin\lambda}{\cos\lambda}\Delta(b,\lambda)F^{1}(c)+F^{1}(c)F(b,\lambda)\\ &-F^{1}(c)E(b,\lambda)-E^{1}(c)F(b,\lambda). \end{aligned} \]
But by virtue of the equation $c^{2}=n\cot^{2}\lambda$ or $c\operatorname{tang}\theta\operatorname{tang}\lambda=1$, which relates to the modulus $b$, we have (no.\ 18) $F(b,\theta)+F(b,\lambda)=F^{1}(b)$, $E(b,\theta)+E(b,\lambda)=E^{1}(b)+b^{2}\sin\lambda\sin\theta$; we have at the same time $\sin\lambda=\frac{\cos\theta}{\Delta(b,\theta)}$, $\cos\lambda=\frac{c\sin\theta}{\Delta(b,\lambda)}$\ednote{Printed $\Delta(b,\lambda)$ in the denominator of $\cos\lambda$; the value $\cos\lambda=c\sin\theta:\Delta(b,\theta)$ is meant, an apparent misprint.}, $\Delta(b,\lambda)=\frac{c}{\Delta(b,\theta)}$, $\frac{\Delta(b,\lambda)}{\sin\lambda\cos\lambda}$, $=\frac{\Delta(b,\theta)}{\sin\theta\cos\theta}=\sqrt{\alpha}$,

\origpage{138}
Adding, therefore, the preceding equation to the equation $(i')$, and observing that by the formula of no.\ 47 we have $\Pi^{1}(n,c)+\Pi^{1}\left(\frac{c^{2}}{n},c\right)=F^{1}(c)+\frac{\pi}{2\sqrt{\alpha}}$, we find that the angles $\theta$ and $\lambda$ disappear entirely from the calculation, and that there remains between the complete functions, for the moduli $c$ and $b$, this relation:
\[ \frac{\pi}{2}=F^{1}(c)E^{1}(b)+F^{1}(b)E^{1}(c)-F^{1}(c)F^{1}(b), \]
which is a new demonstration of the theorem of art.\ 42.

\subsection*{(98).}

Since the complete function $\Pi'(n,c)$\ednote{The print shows a prime here where the work writes $\Pi^{1}$ for the complete function throughout; probably a damaged or wrong superscript $1$.} can be expressed by functions of the first and second species, we shall be able to express in the same manner every incomplete function $\Pi(n,c,\varphi)$ whose amplitude $\varphi$ is such that we have $F(c,\varphi)=kF^{1}(c)$, $k$ being a rational number. Indeed, if this relation holds, it follows from the formulas demonstrated above that we have $\Pi(n,c,\varphi)=k\Pi^{1}(n,c)+W$, $W$ being a quantity determinable by arcs of a circle. Hence there is an infinity of cases in which the function $\Pi$, whose parameter is arbitrary provided it be positive, can be reduced to functions of the first and second species, and these cases are determined by a general symptom.

The amplitude $\varphi$ being given, one can find by a fairly simple calculation whether the condition just mentioned can be fulfilled, that is, whether the function $F(\varphi)$ is in a rational ratio with the complete function $F^{1}$. For this one must calculate the approximate value of the function $F(\varphi)$ by the method of article 65. When the angle $\Phi$, the limit of the angles $\varphi$, $\frac{\varphi^{0}}{2}$, $\frac{\varphi^{00}}{4}$, etc., has been determined, one must examine whether this limit $\Phi$ stands to the right angle in a ratio that is at once simple and very close. If the ratio were somewhat composite, there would be no occasion for any reduction, given that the reduction requires the ratio to be exact, and that, were it at all composite, the equation which determines $\sin\varphi$ would be of a very high degree, which would make it improbable that the given value of $\varphi$ satisfied this equation. But if the ratio in question is expressed in simple numbers and appears at least very close, there will be reason to believe that it is exact. Then one will make sure, by the rigorous formulas, whether the given value of $\sin\varphi$ corresponds to the

\origpage{139}
ratio $k$ found between $F(c,\varphi)$ and $F^{1}(c)$, or between $\Phi$ and $\tfrac{1}{2}\pi$. When this is so, one will likewise find, by the rigorous formulas, the value of $W$ which gives $\Pi(n,c,\varphi)=k\Pi^{1}(n,c)+W$; thus we shall be able to determine $\Pi(n,c,\varphi)$ by functions of the first and second species.

The method just indicated is nothing other than the method of approximation by which the function of the first species $F(c,\varphi)$ is evaluated. But for the same object one may employ another method which seems to lead more directly to the goal.

The amplitude $\varphi$ being given, one will calculate successively the amplitudes $\varphi_{2}$, $\varphi_{3}$, $\varphi_{4}$, etc., which give $F(\varphi_{2})=2F(\varphi)$, $F(\varphi_{3})=3F(\varphi)$, $F(\varphi_{4})=4F(\varphi)$, etc.; this will be done by the formulas
\[ \begin{aligned} \operatorname{tang}\tfrac{1}{2}\varphi_{2}&=\Delta\operatorname{tang}\varphi\\ \operatorname{tang}(\tfrac{1}{2}\varphi_{3}+\tfrac{1}{2}\varphi)&=\Delta\operatorname{tang}\varphi_{2}\\ \operatorname{tang}(\tfrac{1}{2}\varphi_{4}+\tfrac{1}{2}\varphi_{2})&=\Delta\operatorname{tang}\varphi_{3}\\ \operatorname{tang}(\tfrac{1}{2}\varphi_{5}+\tfrac{1}{2}\varphi_{3})&=\Delta\operatorname{tang}\varphi_{4}\\ &\text{etc.} \end{aligned} \]
It is evident that if the angle $\varphi$ is such that we have $F(c,\varphi)=kF^{1}(c)$, $k$ being a rational number $\frac{\mu}{\nu}$, there will necessarily be in the sequence $\varphi_{2}$, $\varphi_{3}$, $\varphi_{4}$, etc. a term $\varphi_{\nu}$ such that $\varphi_{\nu}=\mu\cdot\frac{\pi}{2}$, that is to say, we must meet in this sequence a term $\varphi_{\nu}$ which is a multiple of the right angle. And since, moreover, for the reasons we have set out, the ratio $\frac{\mu}{\nu}$ must be expressed in fairly simple numbers, we shall never have more than a small number of terms to calculate in the sequence $\varphi_{2}$, $\varphi_{3}$, $\varphi_{4}$, etc., in order to recognize whether the functions $F(\varphi)$ and $F^{1}$ are commensurable with each other, and consequently whether $\Pi(n,c,\varphi)$ can be measured by $\Pi^{1}(n,c)$.

\subsection*{(99).}

To return to the equation $(i')$, we shall observe that this equation will serve in general to determine one of the functions $\Pi(n,c,\varphi)$, $\Pi(n',b,\theta)$ by means of the other, supposed known. These two functions have moduli that are complements of each other, and their amplitudes are deduced reciprocally from their parameters by the equations $n=\cot^{2}\theta$, $n'=\cot^{2}\varphi$.

\origpage{140}
Suppose that the angle $\theta$, deduced from the parameter $n$, is such that we have $F(b,\theta)=kF^{1}(b)$, $k$ being a rational number; we shall have, by the known properties, $E(b,\theta)=kE^{1}(b)+V$, and $\Pi(n',b,\theta)=\Pi^{1}(n',b)+W$\ednote{Printed without the factor $k$; the corresponding statement in art.\ 102 reads $k\Pi^{1}(n',b)+W$, so the $k$ has apparently dropped out.}, the quantities $V$ and $W$ being determinable, the first algebraically, the second by arcs of a circle.

If we now substitute these values in the equation $(i')$, and then put for $\Pi^{1}(n',b)$ its value deduced from the formula $(k')$, we shall have, after the reductions, this very simple result
\[ \begin{aligned} \frac{\Delta(b,\theta)}{\sin\theta\cos\theta}\Pi(n,c,\varphi)=(1-k)\frac{\pi}{2}&+\frac{\sin\theta}{\cos\theta}\Delta(b,\theta)F(c,\varphi)\\ &-VF(c,\varphi)-\frac{W\Delta(c,\varphi)}{\sin\varphi\cos\varphi}. \end{aligned} \]
From which one sees that the function $\Pi$ can then be reduced indefinitely to functions of the first species, and the only condition necessary for this reduction is that the parameter $n$, represented by $\cot^{2}\theta$, be such that the ratio of $F(b,\theta)$ to $F^{1}(b)$ is rational. Thus we have here again, for these reductions, a general symptom which comprises an infinity of particular cases.

The case of $n=c$, which was solved in no.\ 46, is included in the general symptom, since, putting $\cot^{2}\theta=c$, the resulting value of $\theta$ is that which, applied to the modulus $b$, gives $F(b,\theta)=\tfrac{1}{2}F^{1}(b)$. Accordingly, putting $k=\tfrac{1}{2}$ and substituting the values of $V$ and $W$, which are $V=\tfrac{1}{2}(1-c)$, $W=\frac{\sin\varphi\cos\varphi}{2\Delta}\operatorname{arc}\operatorname{tang}.\frac{\Delta\cos\varphi}{(1+c)\sin\varphi}=\frac{\sin\varphi\cos\varphi}{2\Delta}\times\left(\frac{\pi}{2}-\operatorname{arc}\operatorname{tang}.\frac{(1+c)\operatorname{tang}\varphi}{\Delta}\right)$, we recover the formula
\[ \Pi(c)=\tfrac{1}{2}F+\frac{\tfrac{1}{2}F}{1+c}\operatorname{arc}\operatorname{tang}.\frac{(1+c)\operatorname{tang}\varphi}{\Delta}. \]
\ednote{Printed $\frac{\frac{1}{2}F}{1+c}$; the arc-tangent term comes from $W$ alone and carries no factor $F$ (compare the analogous formula for $\Pi(-c)$ in art.\ 102, numerator $\tfrac{1}{2}$), so the $F$ in the numerator is apparently a misprint.}

\section*{Second case, $n=-1+b^{2}\sin^{2}\theta$.}

\subsection*{(100).}

In this case we shall have $\frac{dn}{\sqrt{\alpha}}=2\,d\theta\Delta(b,\theta)$, $\sqrt{\alpha}=\frac{b^{2}\sin\theta\cos\theta}{\Delta(b,\theta)}$, and the general equation of no.\ 94 will become
\[ \begin{aligned} \frac{b^{2}\sin\theta\cos\theta}{\Delta(b,\theta)}\Pi=(F-E)\int\frac{d\theta}{\Delta(b,\theta)}&-c^{2}F\int\frac{d\theta}{(1-b^{2}\sin^{2}\theta)^{\frac{3}{2}}}\\ &+\Delta\sin\varphi\cos\varphi\int\frac{d\theta\Delta(b,\theta)}{\cos^{2}\varphi+b^{2}\sin^{2}\theta\sin^{2}\varphi}. \end{aligned} \]

\origpage{141}
Now by the known formulas we have $\displaystyle\int\frac{d\theta}{\Delta(b,\theta)}=F(b,\theta)$ and $\displaystyle\int\frac{d\theta}{(1-b^{2}\sin^{2}\theta)^{\frac{3}{2}}}=\frac{1}{c^{2}}E(b,\theta)-\frac{b^{2}}{c^{2}}\cdot\frac{\sin\theta\cos\theta}{\Delta(b,\theta)}$; moreover, putting $b^{2}\operatorname{tang}^{2}\varphi=n$\ednote{Printed $n$; the parameter so defined is $n'$, as in art.\ 102, so the prime has apparently dropped out.}, we have
\[ \Delta\sin\varphi\cos\varphi\int\frac{d\theta\Delta(b,\theta)}{\cos^{2}\varphi+b^{2}\sin^{2}\theta\sin^{2}\varphi}=\frac{\Delta}{\sin\varphi\cos\varphi}\Pi(n',b,\theta)-\frac{\Delta\cos\varphi}{\sin\varphi}F(b,\varphi); \]
\ednote{Printed $F(b,\varphi)$ at the end of this display; $F(b,\theta)$ is meant, as the formula $(l')$ below shows, an apparent misprint.}
hence, making all these substitutions, we shall have for the second case the general formula
\[ \frac{b^{2}\sin\theta\cos\theta}{\Delta(b,\theta)}[\Pi(n,c,\varphi)-F(c,\varphi)]=\left\{\begin{aligned} &\frac{\Delta(c,\varphi)}{\sin\varphi\cos\varphi}[\Pi(n',b,\theta)-\cos^{2}\varphi F(b,\theta)]\\ &+F(c,\varphi)F(b,\theta)-E(c,\varphi)F(b,\theta)-F(c,\varphi)E(b,\theta) \end{aligned}\right\}\;(l'). \]
No constant has been added, because, on putting $\theta=0$, the two members vanish.

\subsection*{(101).}

If we wish to know what this equation gives when $\varphi=\tfrac{1}{2}\pi$, we must find the value of $\frac{\Delta(c,\varphi)}{\sin\varphi\cos\varphi}\Pi(n',b,\theta)$ at this limit; and for this we must first suppose $\varphi=\tfrac{1}{2}\pi-\omega$, $\omega$ being infinitely small. Now since $n'=b^{2}\operatorname{tang}^{2}\varphi=b^{2}\cot^{2}\omega$, it follows that $\frac{b^{2}}{n'}=\operatorname{tang}^{2}\omega$, and the formula of no.\ 46 gives
\[ \Pi(n',b,\theta)+\Pi(\operatorname{tang}^{2}\omega,b,\theta)=F(b,\theta)+\frac{\sin\varphi\cos\varphi}{\Delta(c,\varphi)}\operatorname{arc}\operatorname{tang}.\frac{\Delta(c,\varphi)\operatorname{tang}\theta}{\sin\varphi\cos\varphi\Delta(b,\theta)}. \]
But $\omega$ being infinitely small, we have
\[ \Pi(\operatorname{tang}^{2}\omega,b,\theta)=F(b,\theta)-M\operatorname{tang}^{2}\omega, \]
$M$ being a finite quantity; therefore
\[ \frac{\Delta(c,\varphi)}{\sin\varphi\cos\varphi}\Pi(n',b,\theta)=\frac{\cos\varphi\Delta(c,\varphi)}{\sin^{3}\varphi}M+\operatorname{arc}\operatorname{tang}.\frac{\Delta(c,\varphi)\operatorname{tang}\theta}{\sin\varphi\cos\varphi\Delta(b,\theta)}. \]
The second member reduces to $\tfrac{1}{2}\pi$ when we put $\varphi=\tfrac{1}{2}\pi$; therefore we shall have, to determine $\Pi^{1}(n,c)$, the equation
\[ \frac{b^{2}\sin\theta\cos\theta}{\Delta(b,\theta)}[\Pi^{1}(n,c)-F^{1}(c)]=\left.\begin{aligned} &\frac{\pi}{2}+F^{1}(c)F(b,\theta)\\ &-E^{1}(c)F(b,\theta)-F^{1}(c)E(b,\theta) \end{aligned}\right\}\;(m'). \]

\origpage{142}
from which one sees that the complete function of the third species $\Pi^{1}(n,c)$ is expressed generally by functions of the first and second species.

We shall observe that the formulas $(l')$ and $(m')$ could have been deduced from those given in art.\ 51, 95 and 96; but it was not useless to show how one could arrive at them directly.

The incomplete function $\Pi(n,c,\varphi)$ will be expressed likewise by functions of the first and second species, if the amplitude $\varphi$ is such that we have $F(c,\varphi)=kF^{1}(c)$, $k$ being a rational number. For then we have $\Pi(n,c,\varphi)=k\Pi^{1}(n,c)+W$, $W$ being a quantity determinable by arcs of a circle.

\subsection*{(102).}

Suppose that the angle $\theta$, determined by the parameter $n=-1+b^{2}\sin^{2}\theta$, is such that for the modulus $b$ we have $F(b,\theta)=kF^{1}(b)$, $k$ being rational; then we shall have, by the known properties, $\Pi(n',b,\theta)=k\Pi^{1}(n',b)+W$, $W$ being a quantity determinable by arcs of a circle. We shall have at the same time $E(b,\theta)=kE^{1}(b)+V$, $V$ being a quantity determinable algebraically. It will therefore suffice to have the value of $\Pi^{1}(n',b)$, where we have $n'=b^{2}\operatorname{tang}^{2}\varphi$.

To this end, I observe that the formula of no.\ 47 gives
\[ \Pi^{1}(n',b)+\Pi^{1}(\cot^{2}\varphi,b)=F^{1}(b)+\frac{\pi\sin\varphi\cos\varphi}{2\Delta(c,\varphi)}; \]
then by the formula of art.\ 96, we have
\[ \Pi^{1}(\cot^{2}\varphi,b)=\frac{\sin\varphi\cos\varphi}{\Delta(c,\varphi)}\left\{\begin{aligned} &\frac{\pi}{2}+\frac{\sin\varphi}{\cos\varphi}\Delta(c,\varphi)F^{1}(b)+F^{1}(b)F(c,\varphi)\\ &-E^{1}(b)F(c,\varphi)-F^{1}(b)E(c,\varphi) \end{aligned}\right\}. \]
Making all these substitutions in the equation $(l')$ and reducing, we shall have
\[ \Pi(n,c,\varphi)=\left(1-\frac{\Delta(b,\theta)}{b^{2}\sin\theta\cos\theta}V\right)F(c,\varphi)+\frac{\Delta(c,\varphi)\Delta(b,\theta)}{b^{2}\sin\varphi\cos\varphi\sin\theta\cos\theta}W. \]
From which one sees that the function of the third species $\Pi(n,c,\varphi)$ will be reduced indefinitely to the function of the first species $F(c,\varphi)$; and the only condition necessary for this reduction is that the angle $\theta$

\origpage{143}
deduced from the parameter $n=-1+b^{2}\sin^{2}\theta$ be such that we have $F(b,\theta)=kF^{1}(b)$, $k$ being a rational number.

Thus in the case of $n=-c$, where we have $\sin^{2}\theta=\frac{1}{1+c}$, this value of $\theta$ is the one which gives $F(c,\theta)=\frac{1}{2}F^{1}(b)$\ednote{Printed $F(c,\theta)$; the angle $\theta$ belongs to the module $b$ (art.\ 102: $F(b,\theta)=kF^{1}(b)$), so $F(b,\theta)$ is meant.}; the reduction is therefore possible. We then find by the formulas of no.\ 57, $V=\frac{1}{2}(1-c)$, $W=\frac{\sin\varphi\cos\varphi}{2\Delta}\operatorname{arc}\operatorname{tang}\frac{(1-c)\operatorname{tang}\varphi}{\Delta}$; and substituting these values in the formula of the preceding article, we derive from it
\[ \Pi(-c)=\tfrac{1}{2}F+\frac{\frac{1}{2}}{1-c}\operatorname{arc}\operatorname{tang}.\frac{(1-c)\operatorname{tang}\varphi}{\Delta}, \]
which agrees with the formula of no.\ 46.

\subsection*{(103).}

If in the equation $(l')$ we put $n=-c^{2}$ or $\theta=\frac{1}{2}\pi$, we shall again find the same equation as in no.\ 48. But if we put $n=-1+\omega^{2}$, $\omega$ being supposed infinitely small, we shall have $\sin\theta=\frac{\omega}{b}$, or $\theta=\frac{\omega}{b}$, and the substitution made, neglecting the higher powers of $\omega$, will give
\[ \Pi^{1}(n,c)=\frac{\frac{1}{2}\pi}{b\omega}+F^{1}(c)-\frac{1}{b^{2}}E^{1}(c). \]
Such, then, must be the value of the integral $\displaystyle\int\frac{d\varphi}{(\cos^{2}\varphi+\omega^{2}\sin^{2}\varphi)\Delta}$, when we suppose $\omega$ infinitely small and $\varphi=\frac{1}{2}\pi$.

To verify this result by direct integration, I observe that $b$ being the value of $\Delta$ when $\varphi=\frac{1}{2}\pi$, we may put
\[ \Pi=\int\frac{d\varphi}{\cos^{2}\varphi+\omega^{2}\sin^{2}\varphi}\cdot\frac{1}{b}+\int\frac{d\varphi}{\cos^{2}\varphi+\omega^{2}\sin^{2}\varphi}\left(\frac{1}{\Delta}-\frac{1}{b}\right). \]
The first part has for its integral $\frac{1}{b\omega}\operatorname{arc}\operatorname{tang}(\omega\operatorname{tang}\varphi)$, and on putting $\varphi=\frac{1}{2}\pi$ it becomes $\frac{\frac{1}{2}\pi}{b\omega}$. The second part decomposes into these two others
\[ T=\int\frac{d\varphi}{\cos^{2}\varphi}\left(\frac{1}{\Delta}-\frac{1}{b}\right),\qquad V=\int\frac{\omega^{2}\sin^{2}\varphi\,d\varphi}{\cos^{2}\varphi(\cos^{2}\varphi+\omega^{2}\sin^{2}\varphi)}\left(\frac{1}{b}-\frac{1}{\Delta}\right). \]
Now by integrating by parts we have $T=\operatorname{tang}\varphi\left(\frac{1}{\Delta}-\frac{1}{b}\right)-c^{2}\displaystyle\int\frac{d\varphi\sin^{2}\varphi}{\Delta^{3}}$,

\origpage{144}
an expression whose first term $\operatorname{tang}\varphi\left(\frac{1}{\Delta}-\frac{1}{b}\right)=\frac{(b-\Delta)\operatorname{tang}\varphi}{b\Delta}=-\frac{c^{2}\sin\varphi\cos\varphi}{b\Delta(b+\Delta)}$ vanishes when $\varphi=\frac{1}{2}\pi$; therefore we have simply
\[ T=-\int\frac{c^{2}\sin^{2}\varphi\,d\varphi}{\Delta^{3}}=\int\frac{\Delta^{2}-1}{\Delta^{3}}\,d\varphi=F^{1}(c)-\frac{1}{b^{2}}E^{1}(c). \]

We already have the first two terms of the value of $\Pi^{1}(n,c)$; a third term would be given by the integral $V$, which can be put in the form
\[ V=\omega^{2}\int\frac{c^{2}\,d\varphi\sin^{2}\varphi}{(\cos^{2}\varphi+\omega^{2}\sin^{2}\varphi)(b+\Delta)b\Delta}; \]
and proceeding in the same manner, we find that the first term of $V$, developed according to ascending powers of $\omega$, is at the same time the first term of the simpler integral $\omega^{2}\displaystyle\int\frac{c^{2}\,d\varphi\sin^{2}\varphi}{2b^{3}(\cos^{2}\varphi+\omega^{2}\sin^{2}\varphi)}$, which has for its expression $\frac{c^{2}\omega}{2b^{3}}\cdot\frac{\pi}{2}$. Uniting all these parts, we shall have
\[ \Pi^{1}(n,c)=\frac{\pi}{2b\omega}+F^{1}(c)-\frac{1}{b^{2}}E^{1}(c)+\frac{\pi c^{2}}{4b^{3}}\omega+\text{etc.} \]

\section*{Third case, $n=-c^{2}\sin^{2}\theta$.}

\subsection*{(104).}

The substitution of this value of $n$ in the general equation of no.\ 94 gives
\[ \begin{aligned} \frac{\cos\theta}{\sin\theta}\Delta(c,\theta)\Pi&=(F-E)\int\frac{d\theta}{\Delta(c,\theta)}-F\int\frac{d\theta}{\sin^{2}\theta\Delta(c,\theta)}\\ &\quad+\Delta\sin\varphi\cos\varphi\int\frac{c^{2}\,d\theta\sin^{2}\theta}{(1-c^{2}\sin^{2}\theta\sin^{2}\varphi)\Delta(c,\theta)}, \end{aligned} \]
In this formula and in all those that will be deduced from it, the functions $\Delta$, $F$, $E$ being always relative to the modulus $c$, we shall add to them only the designation of the amplitude, that is, instead of $\Delta(c,\theta)$, $F(c,\theta)$, $E(c,\theta)$, we shall write simply $\Delta(\theta)$, $F(\theta)$, $E(\theta)$. This being laid down, we have by the known formulas,
\[ \int\frac{d\theta}{\sin^{2}\theta\Delta(\theta)}=F(\theta)-E(\theta)-\Delta(\theta)\cot\theta; \]
and if we put $n'=-c^{2}\sin^{2}\varphi$, we shall have
\[ \int\frac{c^{2}\sin^{2}\theta}{1-c^{2}\sin^{2}\varphi\sin^{2}\theta}\cdot\frac{d\theta}{\Delta(\theta)}=\frac{1}{\sin^{2}\varphi}[\Pi(n',\theta)-F(\theta)]. \]

\origpage{145}
Therefore the general formula for the third case is
\[ \begin{aligned} \cot\theta\Delta(\theta)[\Pi(n,\varphi)-F(\varphi)]&=\cot\varphi\Delta(\varphi)[\Pi(n',\theta)-F(\theta)]\\ &\quad+E(\theta)F(\varphi)-F(\theta)E(\varphi)\ldots(n'). \end{aligned} \]
No constant has been added, because the form of the equation shows that if we change $\varphi$ into $\theta$ and conversely, the constant would have to change sign, which proves that it is zero.

\subsection*{(105).}

If in the preceding formula we put $\varphi=\frac{1}{2}\pi$, we shall have immediately
\[ \Pi^{1}(n)=F^{1}+\frac{\operatorname{tang}\theta}{\Delta(\theta)}[F^{1}E(\theta)-E^{1}F(\theta)]\ldots(p'). \]
Thus in the third case, as in the first two, the complete function of the third species $\Pi^{1}(n)$ will always be expressed by functions of the first and second species.

The same holds of the incomplete function $\Pi(n,c,\varphi)$ or $\Pi(n,\varphi)$, if the amplitude $\varphi$ is such that we have $F(\varphi)=kF^{1}$, $k$ being a rational number; for then we should have $\Pi(n,\varphi)=k\Pi^{1}(n)+W$, $W$ being a quantity determinable by logarithms.

It is remarkable that the formulas $(n')$ and $(p')$ which we have just found for the third case are simpler than the analogous formulas for the first and second cases, since they do not contain the functions $F$ and $E$ relative to the complementary modulus $b$.

Let us remark further that the value of $\Pi^{1}(n)$ would be indeterminate if we had $\theta=\frac{1}{2}\pi$ or $n=-c^{2}$. But then we have generally, for every value of $\varphi$,
\[ \Pi(-c^{2},\varphi)=\frac{1}{b^{2}}E(\varphi)-\frac{c^{2}\sin\varphi\cos\varphi}{b^{2}\Delta(\varphi)}; \]
and consequently when $\varphi=\frac{1}{2}\pi$, we have
\[ \Pi^{1}(-c^{2})=\frac{1}{b^{2}}E^{1}. \]

\subsection*{(106).}

Let us return to the general formula $(n')$, and suppose that the angle $\theta$, deduced from the parameter $n=-c^{2}\sin^{2}\theta$, is such that we have $F(\theta)=kF^{1}$, $k$ being rational; we shall then have $E(\theta)=kE^{1}+V$, and

\origpage{146}
$\Pi(n',\theta)=k\Pi^{1}(n')+W$, $V$ being a quantity determinable algebraically, and $W$ being determinable by logarithms. These values, and that of the function $\Pi^{1}(n')$ drawn from the formula $(p')$, being substituted in the equation $(n')$, there will result
\[ \Pi(n,\varphi)=\left(1+\frac{V\operatorname{tang}\theta}{\Delta(\theta)}\right)F(\varphi)+\frac{W\cot\varphi\Delta(\varphi)}{\cot\theta\Delta(\theta)}. \]
From which it follows that the function $\Pi(n,\varphi)$ can be brought back indefinitely to functions of the first species, if the parameter satisfies the condition mentioned. We have already given (no.\ 86) a similar symptom of reduction, but the one just indicated is much more general.

Let, for example, $n=-1+b=-c^{2}\sin^{2}\theta$; we shall have $\sin^{2}\theta=\frac{1}{1+b}$; which gives $F(\theta)=\frac{1}{2}F^{1}$. We then have by the formulas of no.\ 57,
\[ E(\theta)=\tfrac{1}{2}E^{1}+\tfrac{1}{2}(1-b), \]
and
\[ \Pi(n',\theta)=\tfrac{1}{2}\Pi^{1}(n')+\frac{\sin\varphi}{4\Delta\cos\varphi}\log\left(\frac{\Delta-(1-b)\sin\varphi\cos\varphi}{\Delta+(1-b)\sin\varphi\cos\varphi}\right); \]
whence results
\[ \Pi(n,\varphi)=\frac{1+b}{2b}F(\varphi)+\frac{1}{4b}\log\left(\frac{\Delta-(1-b)\sin\varphi\cos\varphi}{\Delta+(1-b)\sin\varphi\cos\varphi}\right), \]
which agrees with the formula of no.\ 52.

We shall conclude these researches with a general observation: namely, that whenever the function of the third species $\Pi$ is indefinitely reducible to the lower species, it is always reducible to the first species, that is, it is expressed by the function $F(c,\varphi)$ alone. There is no exception to this general rule except when $n=-c^{2}$, and when $n=-1$, the only cases in which the function of the second species $E(\varphi)$ enters into the expression of $\Pi$, as one sees by the formulas of no.\ 48.

\origpage{147}
\section*{General reduction of elliptic functions whose parameter is imaginary.}

\subsection*{(107).}

Let $p=\frac{\sin\varphi\cos\varphi}{(1+\zeta\sin^{2}\varphi)\Delta}$, $\zeta$ being a constant coefficient; by differentiation, and introducing a second coefficient $k$, we shall have
\[ \frac{dp}{1+kp^{2}}=\frac{d\varphi}{\Delta}\cdot\frac{1-(2+\zeta)\sin^{2}\varphi+(1+2\zeta)c^{2}\sin^{4}\varphi-\zeta c^{2}\sin^{6}\varphi}{(1+\zeta\sin^{2}\varphi)^{2}(1-c^{2}\sin^{2}\varphi)+k\sin^{2}\varphi(1-\sin^{2}\varphi)}. \]
Suppose that the denominator of the second member is equal to the product of these three factors:
\[ (1+n\sin^{2}\varphi)(1+n'\sin^{2}\varphi)(1-m\sin^{2}\varphi); \]
the three equations
\[ \begin{gathered} nn'm=c^{2}\zeta^{2}\\ (1+n)(1+n')(1-m)=b^{2}(1+\zeta)^{2}\\ (n+c^{2})(n'+c^{2})(m-c^{2})=b^{2}c^{2}k. \end{gathered} \]
must then be satisfied. For this purpose, suppose that $n$ and $n'$ are known; by their means the three other quantities $\zeta$, $m$, and $k$ must be determined. And first, the equation which determines $\zeta$ being
\[ 1=\frac{c^{2}\zeta^{2}}{nn'}+\frac{b^{2}(1+\zeta)^{2}}{(1+n)(1+n')}; \]
if, for brevity, we put $M=b^{2}+c^{2}\left(\frac{1+n}{n}\right)\left(\frac{1+n'}{n'}\right)$, we shall have
\[ M\zeta=-b^{2}+\sqrt{\left[\frac{1+n}{n'}\cdot\frac{1+n'}{n'}(c^{2}+n)(c^{2}+n')\right]}. \]
\ednote{Printed $\frac{1+n}{n'}$ as the first factor under the radical instead of $\frac{1+n}{n}$; an apparent misprint, since solving the preceding equation for $\zeta$ with $M$ as defined gives the radicand $\frac{1+n}{n}\cdot\frac{1+n'}{n'}(c^{2}+n)(c^{2}+n')$.}
$\zeta$ being known, we shall have $m$ and $k$ by the equations
\[ \begin{aligned} m&=\frac{c^{2}\zeta^{2}}{nn'}\\ k&=(n+c^{2})(n'+c^{2})\cdot\frac{m-c^{2}}{b^{2}c^{2}}. \end{aligned} \]
This being laid down, the differential equation can be given the form
\[ \frac{dp}{1+kp^{2}}=\frac{d\varphi}{\Delta}\left[\frac{A}{1+n\sin^{2}\varphi}+\frac{A'}{1+n'\sin^{2}\varphi}+\frac{B}{1-m\sin^{2}\varphi}+\frac{1}{\zeta}\right], \]

\origpage{148}
and the coefficients $A$, $A'$, $B$ will be determined in this manner:
\[ \begin{aligned} A&=\frac{n^{3}+(2+\zeta)n^{2}+(1+2\zeta)c^{2}n+\zeta c^{2}}{n'(n-n')(n+m)}\\ A'&=\frac{n'^{3}+(2+\zeta)n'^{2}+(1+2\zeta)c^{2}n'+\zeta c^{2}}{n(n'-n)(n'+m)}\\ B&=\frac{m^{3}-(2+\zeta)m^{2}+(1+2\zeta)c^{2}m-\zeta c^{2}}{m(m+n)(m+n')}. \end{aligned} \]
\ednote{The denominators of $A$ and $A'$ are printed with $n'$ and $n$ as first factors; an apparent misprint for $n(n-n')(n+m)$ and $n'(n'-n)(n'+m)$: only these give $A+A'=1-\frac{1}{\zeta}-B$, as the value of $P$ below requires.}
We shall then have the integral
\[ A\Pi(n)+A'\Pi(n')+B\Pi(-m)+\frac{1}{\zeta}F=\int\frac{dp}{1+kp^{2}}. \]

\subsection*{(108).}

This result holds whatever $n$ and $n'$ may be. Let us therefore suppose that these parameters are imaginary and that we have
\[ \begin{aligned} n&=\nu(\cos\theta+\sqrt{-1}\sin\theta)\\ n'&=\nu(\cos\theta-\sqrt{-1}\sin\theta). \end{aligned} \]
From these values, if for brevity we put $h=\frac{b^{2}\nu}{\nu^{2}+2c^{2}\nu\cos\theta+c^{2}}$\ednote{The exponent of the last $c$ is partly covered by a reader's ink stroke; it is read as $c^{2}$, which is also what the equation for $\zeta$ on p.~147 requires (with $b^{2}+c^{2}=1$).}, we shall find
\[ \begin{aligned} m&=\frac{c^{2}\zeta^{2}}{\nu^{2}}\\ \zeta&=-\nu h+\nu\sqrt{(1+2h\cos\theta+h^{2})}\\ k&=(c^{4}+2c^{2}\nu\cos\theta+\nu^{2})\frac{m-c^{2}}{b^{2}c^{2}}; \end{aligned} \]
whence it is seen that the three quantities $\zeta$, $m$, $k$ are always real; the same will hold of the coefficient $B$, which can be put in the form
\[ B=\frac{m^{3}-(2+\zeta)m^{2}+(1+2\zeta)c^{2}m-\zeta c^{2}}{m^{3}+2m^{2}\nu\cos\theta+m\nu^{2}}. \]
As to the coefficients $A$ and $A'$, they are necessarily imaginary. To obtain their expression more easily, let $A=\frac{1}{2}(P+Q\sqrt{-1})$, $A'=\frac{1}{2}(P-Q\sqrt{-1})$; we shall find that the quantities $P$ and $Q$ are determined by the equations
\[ \begin{aligned} P&=1-\frac{1}{\zeta}-B\\ Q&=\frac{m(1-B)-2-\zeta}{\nu\sin\theta}-\frac{\cos\theta}{\sin\theta}\left(1+\frac{1}{\zeta}+\frac{1}{B}\right). \end{aligned} \]
\ednote{Printed $\frac{1}{B}$ as the last term in the bracket instead of $B$; an apparent misprint. The corresponding formula for $Q'$ on p. 149 has $+B'$, and only $+B$ reproduces the values $Q=-\frac{2}{c\sin\theta}(1+c\cos\theta)$ and $Q'=0$ given in art.\ 110.}

\origpage{149}
All the coefficients being thus known, we shall have, finally, the integral
\[ \begin{aligned} \frac{P+Q\sqrt{-1}}{2}\Pi(n)+\frac{P-Q\sqrt{-1}}{2}\Pi(n')&=-B\Pi(-m)-\frac{1}{\zeta}F\\ &\quad+\frac{1}{\sqrt{k}}\operatorname{arc}\operatorname{tang}.\frac{\sqrt{k}\sin\varphi\cos\varphi}{(1+\zeta\sin^{2}\varphi)\Delta}\ldots(q'). \end{aligned} \]

This integral does not suffice to determine the two functions $\Pi(n)$, $\Pi(n')$; but if one observes that the radical contained in the value of $\zeta$ may be taken with the sign $-$, and that there is thus a second value of $\zeta$, namely,
\[ \zeta'=-\nu h-\nu\sqrt{(1+2h\cos\theta+h^{2})}, \]
it will be seen that there must result from it the second integral that we need.

If we denote the new values of $m$, $k$, $B$, $P$, $Q$ by the same letters with accents, we shall therefore have
\[ \begin{aligned} m'&=\frac{c^{2}\zeta'^{2}}{\nu^{2}}\\ k'&=(c^{4}+2c^{2}\nu\cos\theta+\nu^{2})\cdot\frac{m'-c^{2}}{b^{2}c^{2}}\\ B'&=\frac{m'^{3}-(2+\zeta')m'^{2}+(1+2\zeta')c^{2}m'-c^{2}\zeta'}{m'^{3}+2m'^{2}\nu\cos\theta+m'\nu^{2}}\\ P'&=1-\frac{1}{\zeta'}-B'\\ Q'&=\frac{m'(1-B')-2-\zeta'}{\nu\sin\theta}-\frac{\cos\theta}{\sin\theta}\left(1+\frac{1}{\zeta'}+B'\right), \end{aligned} \]
and the second integral will be
\[ \begin{aligned} \frac{P'+Q'\sqrt{-1}}{2}\Pi(n)+\frac{P'-Q'\sqrt{-1}}{2}\Pi(n')&=-B'\Pi(-m')-\frac{1}{\zeta'}F\ldots(r').\\ &\quad+\frac{1}{\sqrt{k'}}\operatorname{arc}\operatorname{tang}\frac{\sqrt{k'}\sin\varphi\cos\varphi}{(1+\zeta'\sin^{2}\varphi)\Delta} \end{aligned} \]
It is now evident that these two integrals give the values of the functions $\Pi(n)$, $\Pi(n')$, each expressed by means of the functions $\Pi(-m)$ and $\Pi(-m')$. Hence in general \emph{every function of the third species whose parameter is imaginary can be expressed by two functions of the same species whose parameters are real.}

This theorem fully resolves the difficulty that might have been raised

\origpage{150}
over the division of elliptic functions into three species, had it not been established that functions whose parameter is imaginary can always be reduced to those whose parameter is real.

\subsection*{(109).}

We must now enter into some details on the results of the preceding solution.

I observe first that the parameters $-m$, $-m'$ are both less than unity, since by the equation which determines $\zeta$ we have
\[ 1-m=\frac{b^{2}(1+\zeta)^{2}}{1+2\nu\cos\theta+\nu^{2}}. \]
Secondly, if we substitute the value of $\zeta$ in the equation $m=\frac{c^{2}\zeta^{2}}{\nu^{2}}$, we shall have
\[ m-c^{2}=2c^{2}h[h+\cos\theta-\sqrt{(1+2h\cos\theta+h^{2})}]; \]
we shall similarly have
\[ m'-c^{2}=2c^{2}h[h+\cos\theta+\sqrt{(1+2h\cos\theta+h^{2})}]; \]
and the product of these equations gives
\[ (m-c^{2})(m'-c^{2})=-4c^{4}h^{2}\sin^{2}\theta: \]
hence we always have $m<c^{2}$ and $m'>c^{2}$. Thus of the two parameters $-m$, $-m'$, one belongs to the form $-c^{2}\sin^{2}\theta$, and the other to the form $-1+b^{2}\sin^{2}\theta$, so that we may put $m=c^{2}\sin^{2}\lambda$ and $m'=1-b^{2}\sin^{2}\mu$; one may also determine the angles $\lambda$ and $\mu$ directly by the formulas
\[ \begin{aligned} \sin\lambda&=-h+\sqrt{(1+2h\cos\theta+h^{2})}\\ \cos\mu&=\frac{2ch\sin\theta}{b\cos\lambda}. \end{aligned} \]
The corresponding values of $k$ and $k'$ being
\[ \begin{aligned} k&=(c^{4}+2c^{2}\nu\cos\theta+\nu^{2})\left(\frac{m-c^{2}}{b^{2}c^{2}}\right)\\ k'&=(c^{4}+2c^{2}\nu\cos\theta+\nu^{2})\left(\frac{m'-c^{2}}{b^{2}c^{2}}\right), \end{aligned} \]
we see that the first is negative and the second positive; whence it

\origpage{151}
follows that the reduced expression of every elliptic function whose parameter is imaginary will always contain an arc of a circle and a logarithm together with the two functions $\Pi(-m)$, $\Pi(-m')$ and the function of the first species $F$.

\subsection*{(110).}

The general formulas which we have just developed are subject to no exception, but it will not be useless to apply them to some cases which present remarkable reductions.

Let first $\nu=c$; we shall have in this case $\zeta'=-1$, $m'=1$, $B'=0$, $Q'=0$, $P'=2$, $k'=1+2c\cos\theta+c^{2}$, and the formula $(r')$ then becomes
\[ \Pi(n)+\Pi(n')=F+\frac{1}{\sqrt{k'}}\operatorname{arc}\operatorname{tang}.\frac{\sqrt{k'}\operatorname{tang}\varphi}{\Delta}. \]
This result would be deduced immediately from the formula of no.\ 46, since we have $nn'=c^{2}$.

The data for forming the other integral are found by our analysis as follows:
\[ \begin{aligned} \zeta&=\frac{c^{2}+c\cos\theta}{1+c\cos\theta}\\ m&=\left(\frac{c^{2}+c\cos\theta}{1+c\cos\theta}\right)^{2}\\ k&=-\frac{c^{2}\sin^{2}\theta(1+2c\cos\theta+c^{2})}{(1+c\cos\theta)^{2}}\\ B&=1-\frac{1}{\zeta}=-\frac{b^{2}}{c(c+\cos\theta)}\\ P&=0\\ Q&=-\frac{2}{c\sin\theta}(1+c\cos\theta). \end{aligned} \]
Substituting all these values in the equation $(q')$ and putting for brevity $\varepsilon=\frac{c\sin\theta\sqrt{(1+2c\cos\theta+c^{2})}}{1+c\cos\theta}$, we shall have
\[ \begin{aligned} \frac{1}{\sqrt{-1}}[\Pi(n)-\Pi(n')]&=\frac{b^{2}\sin\theta}{(1+c\cos\theta)(c+\cos\theta)}\Pi(-m)-\frac{\sin\theta}{c+\cos\theta}F\\ &\quad+\frac{1}{2\varepsilon}\log\left(\frac{(1+\zeta\sin^{2}\varphi)\Delta+\varepsilon\sin\varphi\cos\varphi}{(1+\zeta\sin^{2}\varphi)\Delta-\varepsilon\sin\varphi\cos\varphi}\right). \end{aligned} \]
These two equations give the value of $\Pi(n)$ and that of $\Pi(n')$

\origpage{152}
in the fairly general case where $n=c(\cos\theta+\sqrt{-1}\sin\theta)$ and $n'=c(\cos\theta-\sqrt{-1}\sin\theta)$.

If, however, we had $\cos\theta=-c$, a change would have to be made in the second equation; we should then have $m=0$, and the two terms $-B\Pi(-m)-\frac{1}{\zeta}F$ of the equation $(q')$ would reduce to the single term $-\left(B+\frac{1}{\zeta}\right)F$, or simply $-F$; so that the equation in question would become
\[ \Pi(n)-\Pi(n')=-\frac{c\sqrt{-1}}{b}F+\frac{\sqrt{-1}}{2b}\log\left(\frac{\Delta+c\sin\varphi\cos\varphi}{\Delta-c\sin\varphi\cos\varphi}\right); \]
we should have at the same time
\[ \Pi(n)+\Pi(n')=F+\frac{1}{b}\operatorname{arc}\operatorname{tang}\frac{b\operatorname{tang}\varphi}{\Delta}. \]
Thus in the case of $n=-c^{2}+bc\sqrt{-1}$, the function $\Pi(n)$ is reduced indefinitely to the first species. But although this case is perhaps the simplest of those where the parameter is imaginary\ednote{Printed ``imagiuaire'' with a turned letter instead of ``imaginaire''; an apparent misprint.}, we nevertheless see that the expression of the function $\Pi(n)$ contains at once an arc of a circle and a logarithm.

\subsection*{(111).}

Another case deserves to be examined with care; it is that in which we have
\[ n=-1+b(\cos\lambda+\sqrt{-1}\sin\lambda). \]
This value gives $\nu\cos\theta=-1+b\cos\lambda$, $\nu\sin\theta=b\sin\lambda$, whence results
\[ \begin{aligned} \nu^{2}&=1-2b\cos\lambda+b^{2}\\ h&=-\frac{1}{2\cos\theta}=\frac{\frac{1}{2}\nu}{1-b\cos\lambda}; \end{aligned} \]
We therefore have $\zeta=-\nu h\pm\nu h$, that is to say, the two values of $\zeta$ are $\zeta=0$, $\zeta'=-2\nu h$.

The value $\zeta=0$ gives $m=0$, which presents the same difficulty as in the last case of the preceding article. It will be resolved in the same manner; for, since we then have $\Pi(-m)=-F$\ednote{Printed clearly as $\Pi(-m)=-F$; with $m=0$ one has $\Pi(-m)=F$, which the reduction that follows requires. An apparent misprint.}, the two terms $-B\Pi(-m)-\frac{1}{\zeta}F$ reduce to $-\left(B+\frac{1}{\zeta}\right)F$; now, supposing first $\zeta$ infinitely small and then zero, we find

\origpage{153}
$B+\frac{1}{\zeta}=\frac{c^{2}}{\nu^{2}}$. Moreover, the substitutions give without difficulty
\[ \begin{aligned} P&=\frac{2b}{\nu^{2}}(b-\cos\lambda)\\ Q&=-\frac{2b\sin\lambda}{\nu^{2}}\\ k&=-\nu^{2}; \end{aligned} \]
we shall therefore have for the first equation
\[ \begin{aligned} &(b-\cos\lambda-\sqrt{-1}\sin\lambda)\Pi(n)+(b-\cos\lambda+\sqrt{-1}\sin\lambda)\Pi(n')\\ &\quad=-\frac{c^{2}}{b}F+\frac{\nu}{2b}\log\left(\frac{\Delta+\nu\sin\varphi\cos\varphi}{\Delta-\nu\sin\varphi\cos\varphi}\right); \end{aligned} \]
and one may remark that this equation would be deduced immediately from the formula of no.\ 51, by putting $n=-1+b\cos\lambda+b\sqrt{-1}\sin\lambda$; for then from the condition equation $(1+n)(1-m)=b^{2}$ we should deduce $1-m=b(\cos\lambda-\sqrt{-1}\sin\lambda)$, and consequently $-m=n'$.

To obtain the second equation, we must continue making the substitutions in the values of the coefficients: after many reductions they give the following results:
\[ \begin{aligned} \zeta'&=-\frac{\nu^{2}}{1-b\cos\lambda}\\ m'&=\frac{c^{2}\nu^{2}}{(1-b\cos\lambda)^{2}}\\ k'&=\frac{b^{2}\nu^{2}\sin^{2}\lambda}{(1-b\cos\lambda)^{2}}\\ B'&=\frac{b}{\nu^{2}}(b-\cos\lambda)\\ \tfrac{1}{2}P'&=\frac{1-b\cos\lambda}{\nu^{2}}\\ \tfrac{1}{2}Q'&=\frac{(b-\cos\lambda)\cdot(1-b\cos\lambda)}{\nu^{2}\sin\lambda}, \end{aligned} \]
and the second equation becomes
\[ \begin{aligned} &[\sin\lambda+\sqrt{-1}(b-\cos\lambda)]\Pi(n)+[\sin\lambda-\sqrt{-1}(b-\cos\lambda)]\Pi(n')\\ &\quad=-\frac{b\sin\lambda(b-\cos\lambda)}{1-b\cos\lambda}\Pi(-m')+\sin\lambda\cdot F+\frac{\nu}{b}\operatorname{arc}\operatorname{tang}\frac{b\nu\sin\lambda\sin\varphi\cos\varphi}{(1-b\cos\lambda-\nu^{2}\sin^{2}\varphi)\Delta}. \end{aligned} \]

\origpage{154}
Thus we shall know the functions $\Pi(n)$, $\Pi(n')$ by means of the single function of the third species $\Pi(-m')$.

If we had $\cos\lambda=b$, the formulas found would agree with those given in the preceding article for the case of $\cos\theta=-c$; and indeed in both cases we have $n=-c^{2}+bc\sqrt{-1}$.

The application of the preceding formula requires that two cases be carefully distinguished, according as $\cos\lambda$ is $>b$ or $<b$, in order to evaluate properly the arc of a circle $Z$ whose tangent is
\[ \frac{\nu b\sin\lambda\sin\varphi\cos\varphi}{(1-b\cos\lambda-\nu^{2}\sin^{2}\varphi)\Delta}. \]

Case 1°. $\cos\lambda>b$: the denominator $1-b\cos\lambda-\nu^{2}\sin^{2}\varphi$ will be positive for every value of $\varphi$, since, putting $\sin\varphi=1$, it reduces to $b(\cos\lambda-b)$, a positive quantity. Then the arc $Z$ increases only up to a certain limit which is its \emph{maximum}, after which it decreases, becomes zero, then negative. Thus a moving body which should describe the arc $Z$ while $\varphi$ increases uniformly would make oscillations about the starting point where $\varphi=0$, and would depart from it on either side only by a quantity determined by the value $\operatorname{tang}\varphi=\sqrt[4]{\left(\frac{1-b\cos\lambda}{b^{3}\cos\lambda-b^{4}}\right)}$. The arc $Z$ will therefore vanish when $\varphi=\frac{1}{2}\pi$, $\varphi=\pi$, etc.

Case 2°. $\cos\lambda<b$: then the denominator $1-b\cos\lambda-\nu^{2}\sin^{2}\varphi$ vanishes when $\sin\varphi=\frac{1-b\cos\lambda}{\nu^{2}}$, in which case we have $\cos\varphi=\frac{b(b-\cos\lambda)}{\nu^{2}}$\ednote{Printed $\sin\varphi$ and $\cos\varphi$ without the squares; the vanishing of $1-b\cos\lambda-\nu^{2}\sin^{2}\varphi$ gives $\sin^{2}\varphi=\frac{1-b\cos\lambda}{\nu^{2}}$ and $\cos^{2}\varphi=\frac{b(b-\cos\lambda)}{\nu^{2}}$. An apparent misprint.}, so that this value of $\varphi$ is real and $<\frac{1}{2}\pi$. At this point the arc $Z$, which has an infinite tangent, is equal to $\frac{1}{2}\pi$: it then increases continually with the angle $\varphi$, and when $\varphi=\frac{1}{2}\pi$, we have $Z=\pi$.

An error would therefore be committed in the application of the formula if, on putting $\varphi=\frac{1}{2}\pi$, one took $Z=0$. The value $Z=0$ holds only when $\cos\lambda>b$, but if we have $\cos\lambda<b$, we must take $Z=\pi$.

\subsection*{(112).}

It is very remarkable that in the case of an imaginary parameter of the form $n=-1+b(\cos\lambda+\sqrt{-1}\sin\lambda)$, the transformation which was used for the approximations (art.\ 84) leads immediately to the reduction of the function $\Pi(n)$.

\origpage{155}
Indeed, from this value of $n$, the formulas of the article cited give
\[ n^{0}=-\frac{\nu^{2}}{(1+b)^{2}}; \]
and since the new parameter $n^{0}$ is real, the solution is given immediately by the first transformed expression
\[ \Pi(n)=\frac{1+c^{0}}{2}\left[p\Pi^{0}+qF^{0}+\frac{\tfrac{1}{2}n}{1+n}\int\frac{d\varphi^{0}\cos\varphi^{0}}{1+n^{0}\sin^{2}\varphi^{0}}\right], \]
where we have
\[ \begin{aligned} p&=\frac{\sin\lambda}{\nu^{2}}[\sin\lambda+\sqrt{-1}(\cos\lambda-b)]\\ q&=\frac{c^{2}}{2b\nu^{2}}[\cos\lambda-b-\sqrt{-1}\sin\lambda]\\ \frac{\tfrac{1}{2}n}{1+n}&=\frac{1}{2b}[b-\cos\lambda+\sqrt{-1}\sin\lambda]. \end{aligned} \]
The value of $\Pi(n)$ expressed as a function of $n^{0}$, $c^{0}$, $\varphi^{0}$ is therefore
\[ \begin{aligned} \Pi(n)={}&\frac{1-b}{2b\nu^{2}}(\cos\lambda-b-\sqrt{-1}\sin\lambda)F(c^{0},\varphi^{0})\\ &+\frac{\sin\lambda}{(1+b)\nu^{2}}(\sin\lambda+(\cos\lambda-b)\sqrt{-1})\Pi(n^{0},c^{0},\varphi^{0})\\ &+\frac{b-\cos\lambda+\sqrt{-1}\sin\lambda}{4b\nu}\log\left(\frac{1+b+\nu\sin\varphi^{0}}{1+b-\nu\sin\varphi^{0}}\right). \end{aligned} \]
Whence it is seen that the function $\Pi(n)$, whose parameter is imaginary, is reduced immediately to the function $\Pi(n^{0},c^{0},\varphi^{0})$ whose parameter is real, and which belongs to the second case of article 50, since we have $n^{0}=-1+b^{02}\cos^{2}\frac{1}{2}\lambda$.

Knowing the value of $\Pi(n)$, we shall have at the same time that of $\Pi(n')$ by changing the sign of $\sqrt{-1}$ in the preceding formula. If one wishes to compare these two solutions, and for that purpose substitutes the values just found for $\Pi(n)$ and $\Pi(n')$ in the first formula of the preceding number, one will readily see, from the known relations between the amplitudes $\varphi$ and $\varphi^{0}$, that this equation becomes identical.

It will not be the same if the substitutions are made in the second equation. We then arrive at an equation between the functions $\Pi(n^{0},c^{0},\varphi^{0})$, $\Pi(-m',c,\varphi)$, which is doubtless true, but which

\origpage{156}
is not identical. This equation is
\[ \begin{aligned} \Pi(n^{0},c^{0},\varphi^{0})={}&-\frac{\tfrac{1}{2}b(1+b)(b-\cos\lambda)}{1-b\cos\lambda}\Pi(-m',c,\varphi)+\frac{1+b}{2}F(c,\varphi)\\ &+\frac{(1+b)\nu}{2b\sin\lambda}\operatorname{arc}\operatorname{tang}\frac{b\nu\sin\lambda\sin\varphi\cos\varphi}{(1-b\cos\lambda-\nu^{2}\sin^{2}\varphi)\Delta}. \end{aligned} \]
To verify it, at least in a particular case, let $\cos\lambda=b$; we shall have $\sin\lambda=c$, $\nu=c$, and this equation becomes
\[ \Pi(n^{0})=\frac{1+b}{2}F(c,\varphi)+\frac{1+b}{2b}\operatorname{arc}\operatorname{tang}\frac{b\operatorname{tang}\varphi}{\Delta}: \]
but in this same case we have $n^{0}=-\frac{c^{2}}{(1+b)^{2}}=-\left(\frac{1-b}{1+b}\right)=-c^{0}$; hence the formula of no.\ 46 applies to the function $\Pi(n^{0})$ and gives
\[ \Pi(n^{0})=\tfrac{1}{2}F^{0}+\frac{1}{2(1-c^{0})}\operatorname{arc}\operatorname{tang}\frac{(1-c^{0})\operatorname{tang}\varphi^{0}}{\Delta^{0}}, \]
we have moreover $F(c,\varphi)=\frac{1+c^{0}}{2}F^{0}=\frac{1}{1+b}F^{0}$; hence, in order that the two values of $\Pi(n^{0})$ may agree with each other, we must have
\[ \operatorname{arc}\operatorname{tang}\left(\frac{b\operatorname{tang}\varphi}{\Delta}\right)=\tfrac{1}{2}\operatorname{arc}\operatorname{tang}\frac{(1-c^{0})\operatorname{tang}\varphi^{0}}{\Delta^{0}}, \]
an equation which does indeed hold, by the known relations between $\varphi$ and $\varphi^{0}$ (art.\ 58).

\subsection*{(113).}

To compare our two solutions still better, let us seek, by the one and the other method, the value of the complete function $\Pi^{1}(n)$. In this case we have $\varphi=\frac{1}{2}\pi$, $\varphi^{0}=\pi$, and since in general $F(c,\varphi)=\frac{1+c^{0}}{2}F(c^{0},\varphi^{0})$, we must put $F(c^{0},\varphi^{0})=\frac{2}{1+c^{0}}F^{1}(c)=(1+b)F^{1}(c)$. Substituting these values in the formula of the preceding number, we shall have
\[ \begin{aligned} \Pi^{1}(n)={}&\frac{c^{2}}{2b\nu^{2}}(\cos\lambda-b-\sqrt{-1}\sin\lambda)F^{1}(c)\\ &+\frac{\sin\lambda}{(1+b)\nu^{2}}(\sin\lambda+(\cos\lambda-b)\sqrt{-1})\Pi(n^{0},c^{0},\pi \end{aligned} \]
\ednote{The closing parenthesis after $\pi$ is not printed at the end of the line; read $\Pi(n^{0},c^{0},\pi)$.}
It remains now to obtain the value of $\Pi(n^{0},c^{0},\pi)$ which is the same

\origpage{157}
as $2\Pi^{1}(n^{0},c^{0})$; this will be found by the formula of art.\ 101, because $n^{0}$ can be put in the form $n^{0}=-1+b^{02}\sin\theta^{2}$, by making $\theta=\frac{1}{2}(\pi-\lambda)$; we shall therefore have
\[ \begin{aligned} \Pi^{1}(n^{0},c^{0})=F^{1}(c^{0})+\frac{\Delta(b^{0},\theta)}{b^{02}\sin\theta\cos\theta}\Big[&\tfrac{1}{2}\pi+F^{1}(c^{0})F(b^{0},\theta)-E^{1}(c^{0})F(b^{0},\theta)\\ &-F^{1}(c^{0})E(b^{0},\theta)\Big]. \end{aligned} \]
Now by the known formulas we have $F^{1}(c^{0})=\frac{1+b}{2}F^{1}(c)$, $E^{1}(c^{0})=\frac{b}{1+b}F^{1}(c)+\frac{1}{1+b}E^{1}(c)$; on the other hand we also have $\Delta(b^{0},\theta)=\frac{\nu}{1+b}$, $b^{02}=\frac{4b}{(1+b)^{2}}$, $\frac{\Delta(b^{0},\theta)}{b^{02}\sin\theta\cos\theta}=\frac{\nu(1+b)}{2b\sin\lambda}$; hence
\[ \begin{aligned} \Pi^{1}(n^{0},c^{0})=\frac{1+b}{2}F^{1}(c)+\frac{(1+b)\nu}{2b\sin\lambda}\Big[&\tfrac{1}{2}\pi+\frac{1+b^{2}}{2+2b}F^{1}(c)F(b^{0},\theta)\\ &-\frac{1}{1+b}E^{1}(c)F(b^{0},\theta)-\left(\frac{1+b}{2}\right)F^{1}(c)E(b^{0},\theta)\Big]. \end{aligned} \]
We must now bring the functions $F(b^{0},\theta)$, $E(b^{0},\theta)$ back to a smaller modulus, which will be derived from $b^{0}$ according to the same law by which $c^{0}$ is reduced from $c$. But, if one did not take care, the notation might lead into error. Indeed, the quantity $b^{0}$ enters the preceding formula only as the complement of the modulus $c^{0}$, and not as a quantity derived from $b$ by the same law by which $c^{0}$ is derived from $c$. This law requires that $c^{0}$ be smaller than $c$, whereas $b^{0}$, the complement of $c^{0}$, is greater than $b$.

To avoid, then, every misunderstanding that might arise from this source, we shall put $b^{0}=C$, and we shall denote by $B$ the complement of $C$, so that we shall have $B=c^{0}$. Now we must express $F(C,\theta)$ and $E(C,\theta)$ by means of $F(C^{0},\theta^{0})$ and $E(C^{0},\theta^{0})$, which will be done by the formulas of arts.\ 60, 61, and, observing that the law of decreasing moduli gives $C^{0}=\frac{1-B}{1+B}=\frac{1-c^{0}}{1+c^{0}}=b$, we shall in this way have the values
\[ \begin{aligned} F(C,\theta)&=\frac{1+C^{0}}{2}F(C^{0},\theta^{0})=\frac{1+b}{2}F(b,\theta^{0})\\ E(C,\theta)&=-BF(C,\theta)+\frac{1+B}{2}E(C^{0},\theta^{0})+\frac{1-B}{2}\sin\theta^{0}\\ &=-\left(\frac{1-b}{2}\right)F(b,\theta^{0})+\frac{1}{1+b}E(b,\theta^{0})+\frac{b}{1+b}\sin\theta^{0}. \end{aligned} \]

\origpage{158}
As regards $\theta^{0}$, this angle is derived from $\theta$ by the formula $\operatorname{tang}(\theta^{0}-\theta)=B\operatorname{tang}\theta=\frac{1-b}{1+b}\operatorname{tang}\theta$, which gives
\[ \operatorname{tang}\theta^{0}=\frac{\sin 2\theta}{b+\cos 2\theta}=\frac{\sin\lambda}{b-\cos\lambda},\quad \sin\theta^{0}=\frac{\sin\lambda}{\nu}, \]
where it must be observed that one must take $\theta^{0}<\frac{1}{2}\pi$ if $\cos\lambda<b$, and $\theta^{0}>\frac{1}{2}\pi$ if $\cos\lambda>b$.

This being laid down, making all the substitutions, there will result
\[ \begin{aligned} \frac{\Pi^{1}(n^{0},c^{0})}{1+b}=\tfrac{1}{4}F^{1}(c)+\frac{\nu}{4b\sin\lambda}[&\pi+F^{1}(c)F(b,\theta^{0})-E^{1}(c)F(b,\theta^{0})\\ &-F^{1}(c)E(b,\theta^{0})]. \end{aligned} \]
Hence, finally, the complete function
\[ \begin{aligned} \Pi^{1}(n)={}&\tfrac{1}{2}F^{1}(c)-\left(\frac{1-b\cos\lambda}{2b\nu^{2}}\right)(b-\cos\lambda+\sqrt{-1}\sin\lambda)F^{1}(c)\\ &+\frac{\sin\lambda+(\cos\lambda-b)\sqrt{-1}}{2b\nu}[\pi+F^{1}(c)F(b,\theta^{0})-E^{1}(c)F(b,\theta^{0})\\ &\qquad-F^{1}(c)E(b,\theta^{0})]. \end{aligned} \]
To obtain the value of this same function by the formulas of art.\ 111, we must first put $m'$ in the form $1-b^{2}\sin^{2}\psi$, and then seek the value of $\Pi^{1}(-m')$ by the formula of art.\ 101; one will observe, moreover, that the angles $\psi$ and $\theta^{0}$ have between them the relation $1=c\operatorname{tang}\psi\operatorname{tang}\theta^{0}$, whence results $F(b,\psi)+F(b,\theta^{0})=F^{1}(b)$, $E(b,\psi)+E(b,\theta^{0})=E^{1}(b)+b^{2}\sin\psi\sin\theta^{0}$. The substitution being therefore made, one will find after the reductions the same value of $\Pi^{1}(n)$ as the preceding, which proves the agreement of the two methods.

\subsection*{(114).}

Since every elliptic function of the third species whose parameter is imaginary can be reduced to two functions of the same species whose parameter is real, it follows:

1°. That every complete function of the third species whose parameter is imaginary can always be reduced to the elliptic functions of the first and the second species.

2°. That every function of this kind, not complete, but whose amplitude $\varphi$ is such that we have $F(c,\varphi)=kF^{1}(c)$, $k$ being a rational number, can likewise be expressed by functions of the first and second species.

\origpage{159}

3°. That every function $\Pi(n)$ whose parameter is imaginary can be reduced indefinitely to the first species, if one can satisfy at the same time the two conditions necessary for the auxiliary functions $\Pi(-m)$, $\Pi(-m')$ to be susceptible of a like reduction. An example of these reductions has already been seen (art.\ 109) and it would be easy to produce many others.

\section*{On a general symptom for recognizing whether two given functions of the third species, which differ only in the parameters, can be reduced one to the other.}

\subsection*{(115).}

We have admitted three forms for the parameters, namely, $n=\cot^{2}\theta$, $n=-1+b^{2}\sin^{2}\theta$, $n=-c^{2}\sin^{2}\theta$; and we have shown that every function which belongs to the first form can be converted into a function which belongs to the second form, and conversely. This conversion is effected by the formula of no.\ 51, and without changing anything in the two other elements of each function, which are the modulus $c$ and the amplitude $\varphi$. But a function whose parameter belongs to the third form can never be reduced except to a function whose parameter is of the same form. For this reason we shall confine ourselves to comparing successively, for the three forms of the parameter, two functions which belong to one and the same form.

Case 1°. $n=\cot^{2}\theta$; the formula $(h')$ of no.\ 95 proves that the function $\Pi(\cot^{2}\theta,c,\varphi)$ can be reduced to the function $\Pi(b^{2}\operatorname{tang}^{2}\varphi,b,\theta)$ which differs from the first in its three elements. Similarly the function $\Pi(\cot^{2}\lambda,c,\varphi)$ can be reduced to the function $\Pi(b^{2}\operatorname{tang}^{2}\varphi,b,\lambda)$; but the functions $\Pi(b^{2}\operatorname{tang}^{2}\varphi,b,\theta)$, $\Pi(b^{2}\operatorname{tang}^{2}\varphi,b,\lambda)$, which differ only in the amplitude, can be compared with each other and with the complete function $\Pi^{1}(b^{2}\operatorname{tang}^{2}\varphi,b)$, in an infinity of ways, by the formulas of nos.\ 53 and 57. Suppose in general that we have the equation
\[ iF(b,\lambda)\pm kF(b,\theta)=lF^{1}(b), \]
$i$, $k$, $l$ being whole numbers; then it will follow, by known principles,
\[ i\Pi(b^{2}\operatorname{tang}^{2}\varphi,h,\lambda)\pm k\Pi(b^{2}\operatorname{tang}^{2}\varphi,b,\theta)=l\Pi^{1}(b^{2}\operatorname{tang}^{2}\varphi,b)+W, \]
\ednote{The second argument of the first $\Pi$ reads as $h$ in the scan, a misprint for $b$ (the other functions in this line and the text have $b$).}

\origpage{160}
$W$ being a quantity determinable by arcs of a circle. But the complete function of the third species $\Pi^{1}(b^{2}\operatorname{tang}^{2}\varphi,b)$ is reducible to the functions of the first and second species; hence the preceding equation gives the means of reducing the two functions $\Pi(b^{2}\operatorname{tang}^{2}\varphi,b,\lambda)$, $\Pi(b^{2}\operatorname{tang}^{2}\varphi,b,\theta)$ one to the other. Hence in general the two functions $\Pi(\cot^{2}\theta,c,\varphi)$, $\Pi(\cot^{2}\lambda,c,\varphi)$ can be reduced one to the other, if the angles $\theta$ and $\lambda$ of the parameters are such that we have $iF(b,\lambda)\pm kF(b,\theta)=lF^{1}(b)$, $i$, $k$, $l$ being whole numbers at will.

Case 2°. $n=-1+b^{2}\sin^{2}\theta$: the function $\Pi(n,c,\varphi)$ can be reduced to the function $\Pi(b^{2}\operatorname{tang}^{2}\varphi,b,\theta)$ by the formula of no.\ 101. Likewise if we have $n'=-1+b^{2}\sin^{2}\lambda$, the function $\Pi(n',c,\varphi)$ can be reduced to the function $\Pi(b^{2}\operatorname{tang}^{2}\varphi,b,\lambda)$. But it will be seen, as in the preceding case, that the two functions $\Pi(b^{2}\operatorname{tang}^{2}\varphi,b,\theta)$, $\Pi(b^{2}\operatorname{tang}^{2}\varphi,b,\lambda)$ can be reduced one to the other, if the angles $\theta$ and $\lambda$, derived from the parameters $n$ and $n'$, are such that we have $iF(b,\lambda)\pm kF(b,\theta)=lF^{1}(b)$, $k$, $i$, $l$ being whole numbers. Hence, this condition holding, it will always be possible to reduce the two functions $\Pi(n,c,\varphi)$, $\Pi(n',c,\varphi)$ one to the other.

Case 3°. $n=-c^{2}\sin^{2}\theta$: the function $\Pi(-c^{2}\sin^{2}\theta,c,\varphi)$ can be reduced to the function $\Pi(-c^{2}\sin^{2}\varphi,c,\theta)$ by the formula of no.\ 94; similarly the function $\Pi(-c^{2}\sin^{2}\lambda,c,\varphi)$ can be reduced to the function $\Pi(-c^{2}\sin^{2}\varphi,c,\lambda)$. Suppose that between the angles $\theta$ and $\lambda$, derived from the parameters, we have the relation $iF(c,\lambda)\pm kF(c,\theta)=lF^{1}(c)$, $i$, $k$, $l$ being whole numbers; then we shall have by known principles,
\[ i\Pi(-c^{2}\sin^{2}\varphi,c,\lambda)\pm k\Pi(-c^{2}\sin^{2}\varphi,c,\theta)=l\Pi^{1}(-c^{2}\sin^{2}\varphi,c)+W, \]
$W$ being a quantity determinable by logarithms. Hence the function $\Pi(-c^{2}\sin^{2}\varphi,c,\lambda)$ can be expressed by the function $\Pi(-c^{2}\sin^{2}\varphi,c,\theta)$, and conversely. Hence in general, if the stated condition holds, the two functions $\Pi(-c^{2}\sin^{2}\theta,c,\varphi)$, $\Pi(-c^{2}\sin^{2}\lambda,c,\varphi)$ will be reducible one to the other.

An example of this reduction is seen directly in the formula of no.\ 52. But the general symptom which we have just given furnishes the means of multiplying such comparisons to infinity, and proves that every given function of the third species can be transformed into an infinity of others which will differ from it only in the parameter.

\origpage{161}

The formula of no.\ 46, which is the simplest and the most remarkable of the formulas of comparison, is included in the general symptom given for the first two forms, since, putting $n=\cot^{2}\theta$ and $\frac{c^{2}}{n}=\cot^{2}\lambda$, we have between the angles $\theta$ and $\lambda$ the relation $c\operatorname{tang}\theta\operatorname{tang}\lambda=1$, whence results $F(b,\theta)+F(b,\lambda)=F^{1}(b)$.

The same formula would serve to compare the two functions $\Pi(-1+b^{2}\sin^{2}\theta)$, $\Pi(-1+b^{2}\sin^{2}\lambda)$; for, putting $n=-1+b^{2}\sin^{2}\theta$, $\frac{c^{2}}{n}=-1+b^{2}\sin^{2}\lambda$, there again results $c\operatorname{tang}\theta\operatorname{tang}\lambda=1$.

\section*{Example of a particular transformation of elliptic functions of the first and the second species.}

\subsection*{(116).}

Two elliptic functions of the first species can be transformed one into the other if their moduli belong to one and the same sequence or scale $\ldots c''$, $c'$, $c$, $c^{0}$, $c^{00}$, etc. formed according to the known law. This property extends to the functions of the second and third species, provided they are joined to a function of the first species, whose coefficient may change at will. But if the moduli of the two functions compared are not included in the sequence in question, there are very few cases where the reduction of one function to the other is possible; we know only a single example where the indefinite reduction takes place between two functions whose moduli are complements of each other, and it is this example that we are going to develop.

Consider the formula $R=\displaystyle\int\frac{dz}{(1-z^{3})^{\frac{2}{3}}}$ which is to be integrated from $z=0$; I put $1-z^{3}=\left(\frac{z}{u}\right)^{\frac{3}{2}}$, which gives $z^{3}=-\frac{1}{2}u^{-\frac{3}{2}}+\sqrt{(1+\frac{1}{4}u^{-3})}$\ednote{Printed $z^{3}$; the quadratic in $z^{\frac{3}{2}}$ requires $z^{\frac{3}{2}}$ here, an apparent misprint.}, and I have the transformed expression
\[ R=\int\frac{du}{\sqrt{(4u^{3}+1)}}. \]
Let $m=\sqrt[3]{4}$, and $mu=t^{2}-1$; we shall have for second transformed expression
\[ R=\frac{2}{m}\int\frac{dt}{\sqrt{(t^{4}-3t^{2}+3)}}. \]

\origpage{162}
Finally let $n=\sqrt[4]{3}$, $t=n\operatorname{tang}\frac{1}{2}\varphi$, $c^{2}=\frac{2+\sqrt{3}}{4}$; we shall have for the third transform
\[ R=\frac{1}{mn}\int\frac{d\varphi}{\sqrt{(1-c^{2}\sin^{2}\varphi)}}=\frac{1}{mn}(F(c,\varphi)+C). \]
To determine the constant, let $z=0$; we shall have $u=0$, $t=1$; and if we make $\operatorname{tang}\frac{1}{2}\lambda=\frac{1}{n}$, we shall have $\varphi=\lambda$, so that the integral sought will be
\[ R=\frac{1}{mn}[F(c,\varphi)-F(c,\lambda)]. \]
If we take a second angle $\mu$ such that $b\operatorname{tang}\mu\operatorname{tang}\lambda=1$, so that $F(c,\lambda)+F(c,\mu)=F^{1}(c)$, it will follow that $\operatorname{tang}\mu=\frac{1}{b}\cot\lambda=\frac{1}{b}\cdot\frac{n^{2}-1}{2n}=\frac{\sqrt{2}}{n}=\sqrt{\left(\frac{2}{\sqrt{3}}\right)}$. This value is the one which for the modulus $c$ gives $F(c,\mu)=\frac{1}{3}F^{1}(c)$; hence $F(c,\lambda)=\frac{2}{3}F^{1}(c)$, and finally
\[ R=\frac{1}{mn}[F(c,\varphi)-\tfrac{2}{3}F^{1}(c)]. \]
The entire integral taken up to $z=1$ will be found by making $\varphi=\pi$, and its value will be
\[ R'=\frac{1}{mn}\cdot\tfrac{4}{3}F^{1}(c). \]

\subsection*{(117).}

Let us now make use of another transformation to obtain the value of $R$. Let $\sqrt[3]{(1-z^{3})}=1-\frac{z}{y}$, which gives $(y^{3}-1)z^{2}+3yz=3y^{2}$, $z=\frac{y\sqrt{3}}{2}\cdot\frac{\sqrt{(4y^{3}-1)}-\sqrt{3}}{y^{3}-1}$; we shall first have $R=\displaystyle\int\frac{1}{zy}\left(\frac{dz}{z}-\frac{dy}{y}\right)$, and completing the substitution, we find
\[ R=\sqrt{3}\cdot\int\frac{-dy}{\sqrt{(4y^{3}-1)}}. \]
Now let $my=1+x^{2}$; we shall have
\[ R=\frac{2n^{2}}{m}\int\frac{-dx}{\sqrt{(x^{4}+3x^{2}+3)}}. \]
Finally let $x=n\cot\frac{1}{2}\omega$, and $b^{2}=\frac{2-\sqrt{3}}{4}$; we shall have for the last

\origpage{163}
transform
\[ R=\frac{n}{m}\int\frac{d\omega}{\sqrt{(1-b^{2}\sin^{2}\omega)}}=\frac{n}{m}F(b,\omega). \]
There is no constant to be added, because on making $z=0$ we have successively $y=\infty$, $z=\infty$\ednote{Printed $z=\infty$; $x=\infty$ is evidently meant, an apparent misprint.}, $\omega=0$.

If we make $z=1$ in order to have the complete or definite integral, we shall have $y=1$, $x^{2}=m-1$, $\cot\frac{1}{2}\omega=\frac{\sqrt{(m-1)}}{n}$. But the angle $\theta$ determined by the equation $\cot\frac{1}{2}\theta=\frac{\sqrt{(m-1)}}{n}$ is such that $F(b,\theta)=\frac{4}{3}F^{1}(b)$.

Indeed, we found by the formulas of the trisection (art.\ 24) that for the modulus $b=\frac{1}{2}\sqrt{(2-\sqrt{3})}$ the equation $F(b,\gamma)=\frac{1}{3}F^{1}(b)$ is satisfied by taking $\cos\gamma=(m-1)\sqrt{\left(\frac{2+n^{2}}{n^{2}}\right)}$; let $m-1=n^{2}r$, we shall have $\cos\gamma=nr\sqrt{(2+n^{2})}$, and $\sin\gamma=\sqrt{(1-2n^{2}r^{2}-n^{4}r^{2})}$. But the equation $m=1+n^{2}r$, when raised to the cube, gives $1=n^{2}r+n^{4}r^{2}+n^{2}r^{3}$; hence $1-2n^{2}r^{2}-n^{4}r^{2}=n^{2}(r-2r^{2}+2r^{3})$\ednote{Printed $+2r^{3}$; the preceding equation gives $+r^{3}$, an apparent misprint.}, and consequently $\sin\gamma=n(1-r)\sqrt{r}$, $\operatorname{tang}\gamma=\frac{1-r}{\sqrt{r(2+n^{2})}}$. Let $\delta$ be the angle which gives $F(b,\delta)+F(b,\gamma)=F^{1}(b)$, or $F(b,\delta)=\frac{2}{3}F^{1}(b)$; we shall have $c\operatorname{tang}\delta\operatorname{tang}\gamma=1$; hence $\operatorname{tang}\delta=\frac{\cot\gamma}{c}=\frac{2\cot\gamma}{\sqrt{(2+n^{2})}}=\frac{2\sqrt{r}}{1-r}$, and $\operatorname{tang}\frac{1}{2}\delta=\sqrt{r}=\frac{\sqrt{(m-1)}}{n}$. This value being that of $\cot\frac{1}{2}\theta$, it follows that $\frac{1}{2}\delta+\frac{1}{2}\theta=\frac{1}{2}\pi$, or $\theta=\pi-\delta$, hence $F(b,\theta)=2F^{1}(b)-\frac{2}{3}F^{1}(b)=\frac{4}{3}F^{1}(b)$. Hence, finally,
\[ R^{1}=\frac{n}{m}\cdot\tfrac{4}{3}F^{1}(b). \]
\ednote{Printed $R^{1}$ here, while the next sentence and p.~162 write $R'$; possibly a damaged prime.}
Comparing the two values found for $R'$, we draw from them
\[ F^{1}(c)=n^{2}F^{1}(b)=\sqrt{3}F^{1}(b), \]
which we have already found (no.\ 41).

And if we compare the two values of $R$, we shall have, whatever $\varphi$ may be, this general formula of reduction between two functions of different moduli
\[ F(c,\varphi)-\tfrac{2}{3}F^{1}(c)=\sqrt{3}F(b,\omega). \]

\origpage{164}

As for the relation between the angles $\omega$ and $\varphi$, it can be deduced from the preceding analysis, and there results
\[ \operatorname{tang}\tfrac{1}{2}\varphi=\frac{n\sin\tfrac{1}{2}\omega\Delta(b,\omega)+\cos^{3}\tfrac{1}{2}\omega}{n\cos\tfrac{1}{2}\omega\Delta(b,\omega)-\sin^{3}\tfrac{1}{2}\omega}. \]
Thus the simple substitution of this value in the integral $\displaystyle\int\frac{d\varphi}{\sqrt{(1-c^{2}\sin^{2}\varphi)}}$ must transform it into $n^{2}\displaystyle\int\frac{d\omega}{\sqrt{(1-b^{2}\sin^{3}\omega)}}$\ednote{Printed $\sin^{3}\omega$ instead of $\sin^{2}\omega$; an apparent misprint.}, and this one can easily verify.

The function $F(c,\varphi)$ is therefore expressed by means of the function $F(b,\omega)$ and the complete function $F^{1}(b)$; for we deduce from the preceding equations,
\[ F(c,\varphi)=\sqrt{3}[F(b,\omega)+\tfrac{2}{3}F^{1}(b)]. \]
The second member can then be reduced to a single function; for, having already $F(b,\delta)=\frac{2}{3}F^{1}(b)$, if we take an angle $\psi$ such that $F(b,\psi)=F(b,\omega)+F(b,\delta)$, which will be done by the formulas of art.\ 18, we shall have
\[ F(c,\varphi)=\sqrt{3}F(b,\psi); \]
so that the functions $F(c,\varphi)$, $F(b,\psi)$, whose moduli are complements of one another, will always be in a constant ratio to each other.

Moreover, one can immediately have the value of $\operatorname{tang}\frac{1}{2}\varphi$ expressed as a function of $\psi$; it suffices for this to substitute, in the value already found for $\operatorname{tang}\frac{1}{2}\varphi$, those of $\operatorname{tang}\frac{1}{2}\omega$ and $\Delta(b,\omega)$, which are
\[ \begin{aligned} \operatorname{tang}\tfrac{1}{2}\omega&=\frac{\sin\psi\Delta(b,\delta)-\sin\delta\Delta(b,\psi)}{\cos\psi+\cos\delta}\\ \Delta(b,\omega)&=\frac{\Delta(b,\psi)\Delta(b,\delta)+b^{2}\sin\psi\sin\delta\cos\psi\cos\delta}{1-c^{2}\sin^{2}\psi\sin^{2}\delta}, \end{aligned} \]
and this substitution presents no difficulty.

\subsection*{(118).}

Let us now seek the value of the function of the second species $E(c,\varphi)$ expressed by the function $E(b,\omega)$, or, if necessary, by the two functions $E(b,\omega)$, $F(b,\omega)$. The means which naturally presents itself is to substitute the known value of

\origpage{165}
$\operatorname{tang}\frac{1}{2}\varphi$ as a function of $\omega$ in the integral $\displaystyle\int\Delta\,d\varphi$; and as we have $\displaystyle\int\Delta\,d\varphi=\displaystyle\int\frac{d\varphi}{\Delta}(1-c^{2}\sin^{2}\varphi)=\displaystyle\int\frac{n^{2}\,d\omega}{\Delta(b,\omega)}(1-c^{2}\sin^{2}\varphi)$, the substitution would remain to be made only in the factor $1-c^{2}\sin^{2}\varphi$; now we find
\[ \sqrt{(1-c^{2}\sin^{2}\varphi)}=\frac{(1+\tfrac{1}{2}(n^{2}-1)\sin^{2}\omega)\Delta(b,\omega)-\tfrac{1}{2}n\sin\omega\cos\omega}{1+n^{2}-\tfrac{1}{2}n^{2}\sin^{2}\omega+n\sin\omega\cos\omega\Delta(b,\omega)}. \]
But we see that the result of the substitution will be very complicated, and that unless the integration offers unexpected reductions that are difficult to perceive, functions of the third species will necessarily enter into the expression of $E(c,\varphi)$, which would render this transformation illusory.

It therefore appears necessary to follow another route in order to reach the goal we have set ourselves. For this purpose, let us consider the integral formula
\[ T=\int\frac{x\,dx}{\sqrt[3]{(1-x^{3})}}: \]
if we make $1-x^{3}=\left(\frac{x}{u}\right)^{\frac{3}{2}}$, we shall have $x^{\frac{3}{2}}=-\frac{1}{2}u^{-\frac{3}{2}}+\sqrt{(1+\frac{1}{4}u^{-3})}$, $T=\displaystyle\int x^{\frac{1}{2}}u^{\frac{1}{2}}\,dx=-\frac{1}{2u}-\displaystyle\int\frac{du}{2u^{2}\sqrt{(4u^{3}+1)}}$. But integration by parts gives $\displaystyle\int\frac{du}{u^{2}\sqrt{(4u^{3}+1)}}=-\frac{1}{u}\sqrt{(4u^{3}+1)}+\displaystyle\int\frac{2u\,du}{\sqrt{(4u^{3}+1)}}$. Hence
\[ T=\frac{\sqrt{(4u^{3}+1)}-1}{2u}-\int\frac{u\,du}{\sqrt{(4u^{3}+1)}}. \]
Everything is therefore reduced to finding the integral $V=\displaystyle\int\frac{u\,du}{\sqrt{(4u^{3}+1)}}$. For this let $m^{3}=4$ and $mu=t^{2}-1$; we shall have
\[ V=\frac{2}{m^{2}}\int\frac{(t^{2}-1)\,dt}{\sqrt{(t^{4}-3t^{2}+3)}}. \]
Let again $n^{4}=3$, $t=n\operatorname{tang}\frac{1}{2}\varphi$, $c=\frac{1}{2}\sqrt{(2+\sqrt{3})}$, and the transform will be
\[ V=\frac{1}{m^{2}n}\int\frac{(n^{2}\operatorname{tang}^{2}\tfrac{1}{2}\varphi-1)\,d\varphi}{\sqrt{(1-c^{2}\sin^{2}\varphi)}}=\frac{n}{m^{2}}\int\frac{d\varphi}{\Delta\cos^{2}\tfrac{1}{2}\varphi}-\left(\frac{n^{2}+1}{m^{2}n}\right)F(c,\varphi). \]
But we have $\displaystyle\int\frac{d\varphi}{\Delta\cos^{2}\frac{1}{2}\varphi}=2\Delta\operatorname{tang}\frac{1}{2}\varphi+2F(c,\varphi)-2E(c,\varphi)$; hence

\origpage{166}
finally
\[ V=\frac{2n}{m^{2}}\left[\Delta\operatorname{tang}\tfrac{1}{2}\varphi+\frac{n^{2}-1}{2n^{2}}F(c,\varphi)-E(c,\varphi)\right]; \]
and the integral sought
\[ T=\frac{\sqrt{(4u^{3}+1)}-1}{2u}-\frac{2n}{m^{2}}\Delta\operatorname{tang}\tfrac{1}{2}\varphi+\frac{1-n^{2}}{m^{2}n}F+\frac{2n}{m^{2}}E+\text{const.} \]
The constant must be determined in such a way that the integral vanishes when $x=0$: then $u=0$, $t=1$, $\operatorname{tang}\frac{1}{2}\varphi=\frac{1}{n}$, $\varphi=\lambda$, $F(c,\varphi)=F(c,\lambda)=\frac{2}{3}F^{1}(c)$, $E(c,\lambda)=\frac{2}{3}E^{1}(c)+\frac{c^{2}\sin\mu\sin\lambda}{3}(2-\sin\mu)=\frac{2}{3}E^{1}(c)+\frac{n^{2}+1}{2n}$, $\Delta\operatorname{tang}\frac{1}{2}\varphi=\frac{1}{n(n^{2}+1)}$. Hence we have
\[ \begin{aligned} T={}&\frac{\sqrt{(4u^{3}+1)}-1}{2u}-\frac{2n}{m^{2}}\Delta\operatorname{tang}\tfrac{1}{2}\varphi-\left(\frac{n^{2}-1}{m^{2}n}\right)(F(c,\varphi)-\tfrac{2}{3}F^{1}(c))\\ &+\frac{m}{2(n^{2}+1)}+\frac{2n}{m^{2}}(E(c,\varphi)-\tfrac{2}{3}E^{1}(c)), \end{aligned} \]
or, to express everything as a function of $\varphi$:
\[ \begin{aligned} T={}&\tfrac{1}{2}m\cos\tfrac{1}{2}\varphi\cdot\frac{n(n^{2}+1)\Delta\sin\tfrac{1}{2}\varphi-\cos\tfrac{1}{2}\varphi}{n^{2}\sin^{2}\tfrac{1}{2}\varphi-\cos^{2}\tfrac{1}{2}\varphi}-\left(\frac{n^{2}-1}{m^{2}n}\right)[F(c,\varphi)-\tfrac{2}{3}F^{1}(c)]\\ &+\frac{2n}{m^{2}}[E(c,\varphi)-\tfrac{2}{3}E^{1}(c)]. \end{aligned} \]
When $x=1$, we have $\varphi=\pi$, and this value becomes
\[ T^{1}=\frac{2mn}{3}\left[E^{1}(c)-\left(\frac{n^{2}-1}{2n^{2}}\right)F^{1}(c)\right]. \]
Now by a property of the functions $F^{1}(c)$, $E^{1}(c)$, which was demonstrated in no.\ 39, we have $E^{1}(c)-\frac{n^{2}-1}{2n^{2}}F^{1}(c)=\frac{\pi\sqrt{3}}{4F^{1}(c)}$; hence
\[ T^{1}=\frac{m}{2n}\cdot\frac{\pi}{F^{1}(c)}. \]

\subsection*{(119).}

Let us now integrate by another procedure the proposed formula $T=\displaystyle\int\frac{x\,dx}{\sqrt[3]{(1-x^{3})}}$, and let, as above, $\sqrt[3]{(1-x^{3})}=1-\frac{x}{y}$.

\origpage{167}
The transform will be
\[ T=\frac{9\sqrt{3}}{2}\int\frac{y^{4}\,dy}{(y^{3}-1)^{2}\sqrt{(4y^{3}-1)}}-\int\frac{3y\,dy\,(1+\tfrac{1}{2}y^{3})}{(y^{3}-1)^{2}}. \]
This transform appears at first glance much more complicated than the one that resulted from the first method; and it seems that one cannot avoid introducing into the integral elliptic functions of the third species, even some whose parameter is imaginary. But these fears are dispelled on continuing the calculation, which offers very happy reductions.

I observe first that the rational part is integrated algebraically, and that we have
\[ \int\frac{3y\,dy\,(1+\tfrac{1}{2}y^{3})}{(y^{3}-1)^{2}}=-\tfrac{3}{2}\cdot\frac{y^{2}}{y^{3}-1}. \]
Next, if we make $Z=\dfrac{y^{2}\sqrt{(4y^{3}-1)}}{y^{3}-1}$, we shall have by differentiation,
\[ dZ=\frac{2y\,dy}{\sqrt{(4y^{3}-1)}}-\frac{9y^{4}\,dy}{(y^{3}-1)^{2}\sqrt{(4y^{3}-1)}}; \]
hence
\[ \int\frac{y^{4}\,dy}{(y^{3}-1)^{2}\sqrt{(4y^{3}-1)}}=\tfrac{2}{9}\int\frac{y\,dy}{\sqrt{(4y^{3}-1)}}-\tfrac{1}{9}Z. \]
By these various reductions we obtain
\[ T=\sqrt{3}\int\frac{y\,dy}{\sqrt{(4y^{3}-1)}}-\frac{\sqrt{3}}{2}\cdot\frac{y^{2}\sqrt{(4y^{3}-1)}}{y^{3}-1}+\tfrac{3}{2}\cdot\frac{y^{2}}{y^{3}-1}. \]
Everything is therefore reduced to finding the integral $U=\displaystyle\int\frac{y\,dy}{\sqrt{(4y^{3}-1)}}$; now if we make, as above, $my=1+\nu^{2}$, we shall have
\[ U=\frac{m}{2}\int\frac{d\nu\,(1+\nu^{2})}{\sqrt{(\nu^{4}+3\nu^{2}+3)}}. \]
If we then make $\nu=n\cot\frac{1}{2}\omega$, and $b=\frac{1}{2}\sqrt{(2-\sqrt{3})}$, we shall have
\[ U=-\frac{m}{4n}\int\frac{d\omega\,(1+n^{2}\cot^{2}\tfrac{1}{2}\omega)}{\Delta(b,\omega)}, \]
or
\[ U=-\frac{mn}{4}\int\frac{d\omega}{\sin^{2}\tfrac{1}{2}\omega\,\Delta(b,\omega)}+m\left(\frac{n^{2}-1}{4n}\right)\int\frac{d\omega}{\Delta(b,\omega)}, \]

\origpage{168}
and applying the known formulas, we have finally
\[ U=\frac{mn}{2}[E(b,\omega)+\cot\tfrac{1}{2}\omega\Delta(b,\omega)]-\frac{m(n^{2}+1)}{4n}F(b,\omega). \]
Now if in the equation $T=U\sqrt{3}-yx$ we substitute the value of $U$ and that of $yx$ as a function of $\omega$, we shall have the integral sought
\[ \begin{aligned} T={}&\frac{mn^{3}}{2}E(b,\omega)-\frac{mn}{4}(n^{2}+1)F(b,\omega)\\ &+\frac{m\sin\tfrac{1}{2}\omega}{2}\cdot\frac{(4-2n^{2})\cos^{2}\tfrac{1}{2}\omega-3(2-n^{2})\cos^{4}\tfrac{1}{2}\omega-1+n^{3}\sin\tfrac{1}{2}\omega\cos\tfrac{1}{2}\omega\Delta(b,\omega)}{n\cos\tfrac{1}{2}\omega\Delta(b,\omega)+\sin^{3}\tfrac{1}{2}\omega}, \end{aligned} \]
an integral to which no constant need be added, because it vanishes when $\omega=0$.

To have the value of the complete integral $T^{1}$, let $\omega=\theta$; we shall have as above $F(b,\omega)=\frac{4}{3}F^{1}(b)$, $E(b,\omega)=\frac{4}{3}E^{1}(b)+W$, and we find $W=-\frac{b^{2}}{3}\sin\delta\sin\gamma\,(2-\sin\gamma)=-\frac{b^{2}}{3}\sin^{3}\delta$; hence
\[ T^{1}=\frac{2m}{n}\left[E^{1}(b)-\left(\frac{n^{2}+1}{2n^{2}}\right)F^{1}(b)\right]+\frac{3m}{2n}[W+\cot\tfrac{1}{2}\theta\Delta(b,\theta)]-1. \]
But by the values already found, we have $\cot\frac{1}{2}\theta=\sqrt{r}$, $\Delta(b,\theta)=\frac{1}{n(1+r)\sqrt{r}}$, $\sin\delta=\frac{2\sqrt{r}}{1+r}$, $\sin\gamma=1-\cos\delta=\frac{2r}{1+r}$; moreover the relations between $m$, $n$, $r$ give $m=1+n^{2}r$, $n^{2}=\frac{1-3r^{2}}{r(1+r^{2})}$, $m=\frac{2(1-r^{2})}{1+r^{2}}$, $\sqrt{r}=\frac{n}{2}(1-r^{2})$. Substituting all these values as functions of $r$ in the algebraic part of $T^{1}$, and observing further that $0=1-9r^{2}+3r^{4}-3r^{6}$, we find that this part reduces to zero, so that we have simply
\[ T^{1}=\frac{2m}{n}\left[E^{1}(b)-\left(\frac{n^{2}+1}{2n^{2}}\right)F^{1}(b)\right]. \]
Comparing this value with the one found by the other method, there results the equation
\[ \frac{\pi}{4\sqrt{3}}=F^{1}(b)\left[E^{1}(b)-\left(\frac{1+\sqrt{3}}{2\sqrt{3}}\right)F^{1}(b)\right], \]
which agrees with the formula of art.\ 41.

\origpage{169}

We thus obtain, for the definite functions $F^{1}(c)$, $F^{1}(b)$, $E^{1}(c)$, $E^{1}(b)$, the same relations that we had already obtained by a simpler route. As for the indefinite functions $F(c,\varphi)$, $F(b,\omega)$, $E(c,\varphi)$, $E(b,\omega)$, they have among themselves the relations contained in the two equations
\[ \begin{gathered} F(c,\varphi)=n^{2}[F(b,\omega)+\tfrac{2}{3}F^{1}(b)],\\ \frac{mn}{2}[E(c,\varphi)-\tfrac{2}{3}E^{1}(c)]-\left(\frac{n^{2}-1}{m^{2}n}\right)[F(c,\varphi)-\tfrac{2}{3}F^{1}(c)]+\Phi\\ =\frac{mn^{3}}{2}E(b,\omega)-\frac{mn}{2}(n^{2}+1)F(b,\omega)+\Omega, \end{gathered} \]
$\Phi$ and $\Omega$ being known algebraic functions, the one of $\sin\frac{1}{2}\varphi$, the other of $\sin\frac{1}{2}\omega$; whence we see that the functions $F(c,\varphi)$, $E(c,\varphi)$ relating to the modulus $c$ can be expressed indefinitely by the functions $F(b,\omega)$, $E(b,\omega)$ relating to the complementary modulus $b$.

\section*{On the series which give, without transformations, the approximate values of the elliptic functions.}

\subsection*{(120).}

It may be necessary in several cases, especially in problems of mechanics, to express by the single variable $\varphi$ the elliptic functions on which these problems depend. The development is then made in the known manner; but the preceding methods will always serve to simplify greatly the determination of the coefficients.

Let us consider first the function $F=\displaystyle\int\frac{d\varphi}{\Delta}$, and suppose $\frac{1}{\Delta}=A-2B\cos2\varphi+4C\cos4\varphi-6D\cos6\varphi+\text{etc.}$, so that there results
\[ F=A\varphi-B\sin2\varphi+C\sin4\varphi-D\sin6\varphi+\text{etc.} \]
The question is to have, as simply as possible, the values of the coefficients $A$, $B$, $C$, etc. Now by a first development we have
\[ \frac{1}{\Delta}=1+\tfrac{1}{2}c^{2}\sin^{2}\varphi+\frac{1\cdot3}{2\cdot4}c^{4}\sin^{4}\varphi+\frac{1\cdot3\cdot5}{2\cdot4\cdot6}c^{6}\sin^{6}\varphi+\text{etc.} \]
Then putting, in place of the powers of the sines, their values in

\origpage{170}
linear cosines, we find a series of the supposed form, whose coefficients are easy to deduce from one another. We have first the first coefficient
\[ A=1+\frac{1^{2}}{2^{2}}c^{2}+\frac{1^{2}\cdot3^{2}}{2^{2}\cdot4^{2}}c^{4}+\frac{1^{2}\cdot3^{2}\cdot5^{2}}{2^{2}\cdot4^{2}\cdot6^{2}}c^{6}+\frac{1^{2}\cdot3^{2}\cdot5^{2}\cdot7^{2}}{2^{2}\cdot4^{2}\cdot6^{2}\cdot8^{2}}c^{8}+\text{etc.} \]
I represent it by
\[ A=1+m'c^{2}+m''c^{4}+m'''c^{6}+m^{\text{iv}}c^{8}+\text{etc.}; \]
the others will be successively
\[ \begin{aligned} B&=\tfrac{1}{2}m'c^{2}+\tfrac{2}{3}m''c^{4}+\tfrac{3}{4}m'''c^{6}+\tfrac{4}{5}m^{\text{iv}}c^{8}+\text{etc.}\\ C&=\tfrac{2}{3}\cdot\tfrac{1}{8}m''c^{4}+\tfrac{3}{4}\cdot\tfrac{2}{10}m'''c^{6}+\tfrac{4}{5}\cdot\tfrac{3}{12}m^{\text{iv}}c^{8}+\text{etc.}\\ D&=\tfrac{3}{4}\cdot\tfrac{2}{10}\cdot\tfrac{2}{18}m'''c^{6}+\tfrac{4}{5}\cdot\tfrac{3}{12}\cdot\tfrac{4}{21}m^{\text{iv}}c^{8}+\text{etc.}\\ E&=\tfrac{4}{5}\cdot\tfrac{3}{12}\cdot\tfrac{4}{21}\cdot\tfrac{3}{32}m^{\text{iv}}c^{8}+\text{etc.}\\ &\text{etc.} \end{aligned} \]
In general the $n^{\text{th}}$ of the quantities $C$, $D$, $E$, etc., will be deduced from the preceding one by omitting its first term, and multiplying the following terms by those of the series $\frac{n}{(n+1)(2n+2)}$, $\frac{2n}{(n+1)(2n+3)}$, $\frac{3n}{(n+1)(2n+4)}$, etc. By these series we can therefore find the values of $A$, $B$, $C$, etc., if the modulus $c$ is a sufficiently small quantity; but when $c$ differs little from unity, a great number of terms would have to be calculated to obtain only a mediocre approximation.

\subsection*{(121).}

The value of $\frac{1}{\Delta}$ being multiplied successively by $d\varphi$, $d\varphi\cos2\varphi$, $d\varphi\cos4\varphi$, etc., and then integrated from $\varphi=0$ to $\varphi=\frac{1}{2}\pi$, we shall have
\[ \begin{gathered} A\cdot\frac{\pi}{2}=\int\frac{d\varphi}{\Delta},\quad B\frac{\pi}{2}=\int-\frac{d\varphi\cos2\varphi}{\Delta},\\ 2C\cdot\frac{\pi}{2}=\int\frac{d\varphi\cos4\varphi}{\Delta},\quad 3D\cdot\frac{\pi}{2}=\int-\frac{d\varphi\cos6\varphi}{\Delta},\quad\text{etc.} \end{gathered} \]
Thus each coefficient can be determined separately by a

\origpage{171}
definite integral, and all these integrals depend only on the two functions $F^{1}$, $E^{1}$, which we know how to evaluate in all cases with all the necessary precision.

The first two, $A$ and $B$, are those which it is important to determine with the greatest precision; they are found by the formulas $\frac{\pi}{2}A=F^{1}$ and $\frac{\pi}{2}B=\frac{2}{c^{2}}(F^{1}-E^{1})-F^{1}$, and if we substitute for $F^{1}$ and $E^{1}$ their values given in nos.\ 65 and 77, we shall have
\[ \begin{gathered} A=\frac{2\sqrt{c^{0}}}{c}\cdot\frac{2\sqrt{c^{00}}}{c^{0}}\cdot\frac{2\sqrt{c^{000}}}{c^{00}}\cdot\text{etc.}\\ \frac{B}{A}=\frac{c^{0}}{2}+\frac{c^{0}c^{00}}{4}+\frac{c^{0}c^{00}c^{000}}{8}+\text{etc.} \end{gathered} \]
These two first coefficients being found, all the others can be deduced from them; for, differentiating the equation $\frac{1}{\Delta}=A-2B\cos2\varphi+4C\cos4\varphi-\text{etc.}$, we have
\[ \frac{c^{2}\sin\varphi\cos\varphi}{\Delta^{3}}=4B\sin2\varphi-16C\sin4\varphi+36D\sin6\varphi-\text{etc.} \]
Multiplying the first member by $\frac{2\Delta^{2}}{c^{2}}$, and the second by the equivalent quantity $\frac{2}{c^{2}}-1+\cos2\varphi$; multiplying likewise by $\sin2\varphi$ the two members of the equation $\frac{1}{\Delta}=A-2B\cos2\varphi+4C\cos4\varphi-\text{etc.}$, comparing the two products, and reducing all the terms to linear sines, we shall find
\[ \begin{aligned} 2\cdot3C&=4B\left(\frac{2}{c^{2}}-1\right)-A\\ 3\cdot5D&=16C\left(\frac{2}{c^{2}}-1\right)-1\cdot3B\\ 4\cdot7E&=36D\left(\frac{2}{c^{2}}-1\right)-2\cdot5C\\ 5\cdot9F&=64E\left(\frac{2}{c^{2}}-1\right)-3\cdot7D\\ &\text{etc.} \end{aligned} \]
In general, if $A^{n}$ denotes the $n^{\text{th}}$ term of the series $B$, $C$, $D$, $E$, etc., we shall have

\origpage{172}
\[ n(2n-1)A^{n}=(2n-2)^{2}A^{n-1}\left(\frac{2}{c^{2}}-1\right)-(n-2)(2n-3)A^{n-2}: \]
this is the law according to which each coefficient, beginning with $D$, can be deduced from the two preceding ones. The coefficient $C$, being excepted from the general law, is determined by the equation $6C=4B\left(\frac{2}{c^{2}}-1\right)-A$\ednote{The exponent of $c$ in this first occurrence is faintly printed and resembles $c^{4}$; it is read as $c^{2}$, as in the same equation later on this page and on p. 171.}, or more directly by the formula
\[ \begin{aligned} \frac{6C}{A}={}&c^{02}\left(1+\frac{c^{00}}{2}+\frac{c^{00}c^{000}}{4}+\frac{c^{00}c^{000}c^{0000}}{8}+\text{etc.}\right)\\ &+\frac{c^{00}}{2}\left(1+\frac{c^{000}}{2}+\frac{c^{000}c^{0000}}{4}+\text{etc.}\right). \end{aligned} \]
When $c$ is very near unity, we shall determine $E^{1}$ and $F^{1}$ by the formulas suited to this case, we shall deduce from them the coefficients $A$ and $B$ by the formulas $\frac{\pi}{2}A=F^{1}$, $\frac{\pi}{2}B=\frac{2}{c^{2}}(F^{1}-E^{1})-F^{1}$, we shall then calculate $C$ by the equation $6C=4B\left(\frac{2}{c^{2}}-1\right)-A$; the others will be deduced, each from the two preceding ones, by the general law which we have set forth; and this law will be of all the surer application in that the factor $\frac{2}{c^{2}}-1$ will be found, in the case in question, to differ little from unity.

\subsection*{(122).}

The elliptic functions of the second species can be developed in the same manner. Indeed, if we consider in general the function $G=\displaystyle\int(\alpha+\beta\sin^{2}\varphi)\frac{d\varphi}{\Delta}$, or, what comes to the same, $G=\displaystyle\int(\alpha'+\beta'\cos2\varphi)\frac{d\varphi}{\Delta}$, the preceding formulas will give
\[ \begin{aligned} G={}&\alpha'[A\varphi-B\sin2\varphi+C\sin4\varphi-D\sin6\varphi+\text{etc.}]\\ &-\beta'\left[B\varphi-\left(\frac{A+2C}{2}\right)\sin2\varphi+\frac{B+3D}{4}\sin4\varphi-\left(\frac{2C+4E}{6}\right)\sin6\varphi+\text{etc.}\right]. \end{aligned} \]
Thus this development presents no new difficulty, and is carried out by the same coefficients as that of the functions of the first species.

The functions of the third species can be treated in the same way.

\origpage{173}
Indeed, let $\Pi=\displaystyle\int\frac{d\varphi}{(1+n\sin^{2}\varphi)\Delta}$; if we suppose
\[ \Pi=M\varphi-N\sin2\varphi+P\sin4\varphi-Q\sin6\varphi+\text{etc.}, \]
and differentiate each member with respect to $\varphi$, there will result a value of $\frac{1}{\Delta}$ which, being compared with that of the preceding article, gives the equations
\[ \begin{aligned} N&=\frac{2}{n}(A-M)+M\\ 2P&=\frac{4}{n}(B-N)-2N-M\\ 3Q&=\frac{8}{n}(C-P)-4P-N\\ 4R&=\frac{12}{n}(D-Q)-6Q-2P\\ &\text{etc.} \end{aligned} \]
Whence it follows that, supposing $A$, $B$, $C$, etc. always known, it suffices to determine the first coefficient $M$ in order to know all the others. Now on making $\varphi=\frac{1}{2}\pi$, we have $\Pi^{1}=M\cdot\frac{1}{2}\pi$. Thus $M$ is determined by the complete function $\Pi^{1}$, which depends only on the functions of the first and second species.

One could also deduce the value of $M$ from those of the coefficients $A$, $B$, $C$, etc., by means of the known formula $\frac{\sqrt{(1+n)}}{1+n\sin^{2}\varphi}=1+2\alpha\cos2\varphi+2\alpha^{2}\cos4\varphi+\text{etc.}$, where $\alpha=\frac{\sqrt{(1+n)}-1}{\sqrt{(1+n)}+1}$; for on multiplying both members by $\frac{d\varphi}{\Delta}$ and integrating from $\varphi=0$ to $\varphi=\frac{1}{2}\pi$, we shall have
\[ M\sqrt{(1+n)}=A-2B\alpha+4C\alpha^{2}-6D\alpha^{3}+\text{etc.} \]

\section*{Surface of the oblique cone.}

\subsection*{(123).}

\rmfigure{figures/fig-11.png}{Fig.~11.}{A right triangle SBO with the vertical side SO, the line SA to the point A on the base BO, the half-circle on the diameter BA below the base with centre C, and the radius CM to a point M on it.}

Let us consider first the oblique cone with circular base. Let $S$ be the vertex of the cone, $C$ the centre of the base, $SO$ the perpendicular dropped from the vertex to the plane of the base, $SAB$ the section made in the cone by the plane $SCO$. If we make the radius $CA=1$, the height $SO=h$, the distance $CO=f$, and call $\omega$ the angle

\origpage{174}
$OCM$, the area $ASM$ corresponding to the angle $\omega$ will be expressed by the integral
\[ Z=\int\tfrac{1}{2}d\omega\sqrt{[h^{2}+(1-f\cos\omega)^{2}]}. \]
To reduce this integral to the ordinary form, let $\operatorname{tang}\frac{1}{2}\omega=m\operatorname{tang}\frac{1}{2}\varphi$, and let the indeterminate $m$ be taken so that
\[ m^{4}=\frac{h^{2}+(1-f)^{2}}{h^{2}+(1+f)^{2}}; \]
if we then make $\alpha=\frac{1-m^{2}}{1+m^{2}}$ and $c^{2}=\frac{1}{2}-\frac{1}{2}m^{2}\cdot\frac{h^{2}+1-f^{2}}{h^{2}+(1-f)^{2}}$, we shall have the transform
\[ Z=\frac{2m\sqrt{[h^{2}+(1-f)^{2}]}}{(1+m^{2})^{2}}\int\frac{d\varphi\sqrt{(1-c^{2}\sin^{2}\varphi)}}{(1+\alpha\cos\varphi)^{2}}. \]
To have the total surface of the cone, we must take this integral from $\varphi=0$ to $\varphi=\pi$, and double the result; or, what comes to the same, we must take the following integral from $\varphi=0$ to $\varphi=\frac{1}{2}\pi$,
\[ Z=\frac{2m\sqrt{[h^{2}+(1-f)^{2}]}}{(1+m^{2})^{2}}\int\left(\frac{\Delta\,d\varphi}{(1+\alpha\cos\varphi)^{2}}+\frac{\Delta\,d\varphi}{(1-\alpha\cos\varphi)^{2}}\right). \]

Let again $\alpha=\cos\theta$, or $m=\operatorname{tang}\frac{1}{2}\theta$. If we make for brevity $AS=k=\sqrt{[h^{2}+(1-f)^{2}]}$, we shall have on reducing,
\[ Z=2k\cot\tfrac{1}{2}\theta\int\frac{1+2\cot^{2}\theta-\cot^{2}\theta\sin^{2}\varphi}{(1+\cot^{2}\theta\sin^{2}\varphi)^{2}}\cdot\Delta\,d\varphi. \]
But by a first decomposition, we have
\[ \begin{aligned} \int\frac{1+2n-n\sin^{2}\varphi}{(1+n\sin^{2}\varphi)^{2}}\Delta\,d\varphi={}&\frac{c^{2}}{n}\int\frac{d\varphi}{\Delta}-\left(1+2c^{2}+\frac{3c^{2}}{n}\right)\int\frac{d\varphi}{(1+n\sin^{2}\varphi)\Delta}\\ &+(2+2n)\left(1+\frac{c^{2}}{n}\right)\int\frac{d\varphi}{(1+n\sin^{2}\varphi)^{2}\Delta}. \end{aligned} \]
Substituting then for the last integral the value which we found (art.\ 94), the second member will become
\[ \frac{n\Delta\sin\varphi\cos\varphi}{1+n\sin^{2}\varphi}-\int\frac{c^{2}\,d\varphi\sin^{2}\varphi}{\Delta}+(n+1)\int\frac{d\varphi}{(1+n\sin^{2}\varphi)\Delta}. \]
Hence, taking this integral from $\varphi=0$ to $\varphi=\frac{1}{2}\pi$, we

\origpage{175}
shall have the entire area of the cone
\[ Z^{1}=2k\cot\tfrac{1}{2}\theta[(n+1)\Pi^{1}(n,c)-F^{1}(c)+E^{1}(c)], \]
a formula in which $n=\cot^{2}\theta$.

Having already made the smallest apothem $AS=k$, if we make the greatest $BS=k'$, we shall have immediately, by virtue of the preceding values, $\operatorname{tang}\frac{1}{2}\theta=\sqrt{\left(\frac{k}{k'}\right)}$; moreover, if we call $\zeta$ the angle $\frac{1}{2}(ASO+BSO)$ which the perpendicular $AO$\ednote{Printed $AO$; the perpendicular from the vertex is $SO$, an apparent misprint.} makes with the line which divides the angle $ASB$ into two equal parts, we shall have $c=\sin\zeta$. Thus the quantities $c$, $k$, $\theta$, which enter as elements in the final result, are deduced immediately from the triangle $SAB$.

Moreover, since the complete function $\Pi^{1}(n,c)$ can be expressed by functions of the first and second species, we see that the surface of the oblique cone with circular base likewise depends only on the elliptic functions of the first and second species, so that it can be measured by arcs of an ellipse; a theorem which had not yet been demonstrated except in some particular cases.

The surface of the same cone being developed on a plane, there results a sector whose angle is determined by the formula
\[ V=\int\frac{d\omega\sqrt{[h^{2}+(1-f\cos\omega)^{2}]}}{h^{2}+f^{2}+1-2f\cos\omega}, \]
and by means of the same substitutions, we have the transform
\[ V=\frac{m}{k}\int\frac{\Delta\,d\varphi}{(1+n\sin^{2}\varphi)}=\frac{\operatorname{tang}\tfrac{1}{2}\theta}{k}\left[\left(1+\frac{c^{2}}{n}\right)\Pi(n,c,\varphi)-\frac{c^{2}}{n}F(c,\varphi)\right]. \]
This integral being taken from $\varphi=0$ to $\varphi=2\pi$, we shall have the total angle of the sector
\[ V^{1}=\frac{4\operatorname{tang}\tfrac{1}{2}\theta}{k}\left[\left(1+\frac{c^{2}}{n}\right)\Pi^{1}(n,c)-\frac{c^{2}}{n}F^{1}(c)\right]. \]

It may be remarked that if in the expression for $V$ we make $\varphi=\frac{1}{2}\pi$, the value of $V$ will become the quarter of the total angle $V^{1}$. We have at the same time $\omega=\theta$. Hence if we determine the point $M$ so that the angle $ACM=\theta$, the corresponding angle $V$ will be precisely the quarter of the angle resulting from the development of the entire surface of the cone.

\origpage{176}
\subsection*{(124).}

Suppose in the second place that the base of the cone is an ellipse. Let $h$ be the height of the cone, $f$ and $g$ the coordinates of the point where the perpendicular $h$ meets the plane of the base, $f$ being taken in the direction of the semi-major axis of the ellipse $1$, and $g$ in that of the conjugate semi-axis $b$. Let $\varphi$ be the amplitude of any point of the ellipse; the coordinates of this point will be $x=\sin\varphi$, $y=b\cos\varphi$, and the element of the curve $ds=\Delta\,d\varphi$; if from the vertex a perpendicular is drawn to the tangent at this point, the square of this perpendicular will have for its expression $h^{2}+\left(\frac{g\cos\varphi+bf\sin\varphi-b}{\Delta}\right)^{2}$, hence the element of the surface of the cone will be
\[ dZ=\tfrac{1}{2}d\varphi\sqrt{[h^{2}(1-c^{2}\sin^{2}\varphi)+(g\cos\varphi+bf\sin\varphi-b)^{2}]}. \]
This differential appears very complicated; yet if we make $\operatorname{tang}\frac{1}{2}\varphi=z$, which gives $\sin\varphi=\frac{2z}{1+z^{2}}$, $\cos\varphi=\frac{1-z^{2}}{1+z^{2}}$, $d\varphi=\frac{2\,dz}{1+z^{2}}$, we shall have the transform
\[ dZ=\frac{dz}{(1+z^{2})^{2}}\sqrt{[h^{2}(1+z^{2})^{2}-4c^{2}h^{2}z^{2}+(g-b+2bfz-(g+b)z^{2})^{2}]}; \]
and since the variable under the radical does not pass the fourth degree, it is clear that the integral $Z$ can be found by means of elliptic functions. Moreover, as it suffices, in order to have the total surface of the cone, to know the integral when $\varphi=2\pi$, it seems that the final result must depend only on the complete functions of the third species and of the lower species, which, by means of the known reductions, themselves depend only on the functions of the first and second species; so that the entire area of the cone can still be expressed by arcs of an ellipse. But it is necessary to enter into some details in order to justify this conclusion.

\subsection*{(125).}

To make the odd powers of the variable under the radical disappear in the value of $dZ$, we must first suppose $z=\frac{p+qx}{1+x}$; and because the quantity under the radical has, in the case in question, only simple imaginary factors, the result of the substitution, after $p$ and $q$ have been determined, will contain a new radical of the

\origpage{177}
form $\sqrt{[(1+\alpha^{2}x^{2})(1+\beta^{2}x^{2})]}$. To then bring the transformed expression back to elliptic functions, we shall have to put, supposing $\alpha>\beta$, $\alpha x=\operatorname{tang}\psi$, and $\psi$ will become the amplitude of these functions, while their modulus $c$ will be determined by the equation $c^{2}=1-\frac{\beta^{2}}{\alpha^{2}}$. The relation between $\varphi$ and $\psi$ will therefore be such that we have $\operatorname{tang}\frac{1}{2}\varphi=\frac{p\alpha+q\operatorname{tang}\psi}{\alpha+\operatorname{tang}\psi}$. Now I say that this equation can always be put into the form
\[ \operatorname{tang}(\psi-\nu)=A\operatorname{tang}\tfrac{1}{2}(\varphi-\mu), \]
$A$, $\mu$ and $\nu$ being constant quantities. Indeed, if we put $\operatorname{tang}\frac{1}{2}\mu=t$ and $\operatorname{tang}\nu=t'$, the preceding equation will give
\[ \frac{\operatorname{tang}\psi-t'}{1+t'\operatorname{tang}\psi}=\frac{A(\operatorname{tang}\tfrac{1}{2}\varphi-t)}{1+t\operatorname{tang}\tfrac{1}{2}\varphi}. \]
Substituting the value of $\operatorname{tang}\frac{1}{2}\varphi$ as a function of $\operatorname{tang}\psi$, we shall have
\[ \frac{\operatorname{tang}\psi-t'}{1+t'\operatorname{tang}\psi}=\frac{A(p\alpha+q\operatorname{tang}\psi-t\alpha-t\operatorname{tang}\psi)}{\alpha+\operatorname{tang}\psi+p\alpha t+qt\operatorname{tang}\psi}. \]
This equation having to hold whatever $\operatorname{tang}\psi$ may be, let 1°. $\operatorname{tang}\psi=t'$; we shall have
\[ t=\frac{p\alpha+qt'}{\alpha+t'}. \]
Let 2°. $\operatorname{tang}\psi=-\frac{1}{t'}$; we shall have
\[ t=\frac{1-\alpha t'}{p\alpha t'-q}. \]
Equating these two values of $t$, we shall have, to determine $t'$, the equation
\[ \frac{1-t'^{2}}{2t'}=\frac{\alpha^{2}(1+p^{2})-1-q^{2}}{2(1+pq)}=\cot2\nu. \]
\ednote{Printed denominator $2(1+pq)$; equating the two values of $t$ above gives $2\alpha(1+pq)$, an apparent misprint.}
Similarly we should have, to determine $t$, the equation
\[ \frac{1-t^{2}}{2t}=\frac{1-q^{2}+\alpha^{2}(1-p^{2})}{2(q-\alpha^{2}p)}=\cot\mu. \]
\ednote{Printed denominator $2(q-\alpha^{2}p)$; the two values of $t$ above give $2(q+\alpha^{2}p)$, an apparent misprint.}
3°. Finally, in order that the above equation may become entirely identical,

\origpage{178}
let $\operatorname{tang}\psi=\infty$; we shall have
\[ A=\frac{1}{t'}\cdot\frac{1+qt}{q-t}. \]
This being set down, the equation $\operatorname{tang}(\psi-\nu)=A\operatorname{tang}\frac{1}{2}(\varphi-\mu)$ holds for all corresponding values of $\varphi$ and $\psi$; it follows that to the three values $\varphi=\mu$, $\varphi=\pi+\mu$, $\varphi=2\pi+\mu$ there correspond these three values of the amplitude of the elliptic functions, $\psi=\nu$, $\psi=\frac{1}{2}\pi+\nu$, $\psi=\pi+\nu$. But to have the entire area of the cone, we must take the integral $Z$ from $\varphi=0$ to $\varphi=2\pi$; or, what comes to the same thing, from $\varphi=\mu$ to $\varphi=2\pi+\mu$; we must therefore take the elliptic functions that enter into the expression of $Z$ from $\psi=\nu$ to $\psi=\nu+\pi$; or, what comes to the same thing, from $\psi=0$ to $\psi=\pi$. Hence only complete functions will enter into the expression of the total area, and these functions always reduce to the functions of the first and of the second species; whence it follows that in the oblique cone with elliptic base, the entire area can still be expressed by arcs of an ellipse.

\section*{Construction of the shortest line on the surface of the spheroid.}

\subsection*{(126).}

We have given in the Mémoires de l'Institut, year 1806, the equations necessary for describing the shortest line on the surface of a spheroid, and we have shown that this curve, prolonged indefinitely, is composed of an infinity of equal and similar spires, contained between two parallels equally distant from the equator. We now propose to show how the different points of this curve are determined by means of elliptic functions, and for this purpose it will suffice to construct a single half of a spire, from the parallel to which it is tangent as far as the equator.

We must first recall what was shown in the cited memoir, that the shortest line drawn between two given points on the surface of a spheroid, a line called \emph{geodesic} when it is traced on the surface of the terrestrial spheroid, always forms part of a line of this nature perpendicular to the meridian of a

\origpage{179}
determined place, or tangent to the parallel of that place. We shall therefore suppose that through a point $A$ whose latitude is given, a perpendicular is made to pass to the meridian of the place $A$; the question is to construct this curve, that is to say, to determine, for each point $M$ whose latitude is given at will, the unknown elements of the triangle $APM$, namely, the distance $AM$ measured by the arc of the shortest line, the longitude of the point $M$ measured by the angle $P$ of the triangle $APM$, the angle $M$ of the same triangle which represents at this point the azimuthal direction of the curve; finally the arc $PM$ of the moving meridian, which will be known from its coordinates. Of these four elements, two can be determined algebraically, namely, the angle $M$ and the coordinates of the point $M$ in the plane $CMP$; the other two, namely, the arc $AM$ and the angle $APM$, can be assigned only by elliptic functions.

\rmfigure{figures/fig-12.png}{Fig.~12.}{The vertical line PrTC and the horizontal line CDE, with an outer curve running from P through A and Q round E to N and on past I, a dashed curve from P to D, a curve from P through M to N, a curve from A through M to I, the straight line TM, the arc MQ, and a short dashed line from r ending at m.}

Let the radius of the equator be $CA=A$, the semi-axis $CB=B$, the distance from the centre to the focus of the generating ellipse $\sqrt{(A^{2}-B^{2})}=A\delta$; let $l$ be the latitude of the point $A$, $\lambda$ that of the point $M$, $P$ the angle $APM$ which measures the longitude of the point $M$, $M$ the azimuthal angle $AMP$; let there be further the coordinates $CT=t$, $TM=u$, and the arc $AM=s$. The question is to determine all the quantities $t$, $u$, $M$, $s$, $P$ as functions of the variable $\lambda$ and of the given quantities.

It is useful for this purpose to substitute for the latitudes $l$ and $\lambda$ the reduced latitudes $l'$ and $\lambda'$, calculated by the formulas
\[ \operatorname{tang} l'=\frac{B}{A}\operatorname{tang} l,\qquad \operatorname{tang}\lambda'=\frac{B}{A}\operatorname{tang}\lambda. \]
This being set down, the elements that can be determined algebraically in the proposed question are $t=B\sin\lambda'$, $u=A\cos\lambda'$, $\sin M=\frac{\cos l'}{\cos\lambda'}$. It follows from the last equation that $l'$ is the \emph{maximum} of $\lambda'$, and $l$ that of $\lambda$. We also see from the values of the coordinates that $\lambda'$ is the amplitude of the arc $PM$; so that this arc, if it were needed, could be expressed by the formula $PM=A\cdot E(\delta,\lambda')$.

\subsection*{(127).}

To have the other two elements $s$ and $P$, let the angle $\varphi$ be taken such that $\sin\lambda'=\sin l'\cos\varphi$, and we shall have (cited Memoir, no.~4) the differential equations

\origpage{180}
\[ \begin{aligned} ds&=d\varphi\sqrt{(B^{2}+A^{2}\delta^{2}\sin^{2}l'\cos^{2}\varphi)}\\ dP&=\frac{\cos l'}{A}\cdot\frac{d\varphi\sqrt{(B^{2}+A^{2}\delta^{2}\sin^{2}l'\cos^{2}\varphi)}}{1-\sin^{2}l'\cos^{2}\varphi}. \end{aligned} \]
To bring these expressions back to the ordinary forms of elliptic functions, I put $c^{2}=\frac{A^{2}\delta^{2}\sin^{2}l'}{B^{2}+A^{2}\delta^{2}\sin^{2}l'}$, or more simply $c=\delta\sin l$, which gives $\sqrt{(B^{2}+A^{2}\delta^{2}\sin^{2}l')}=\frac{B}{b}$, and I have first $ds=\frac{B}{b}\,d\varphi\sqrt{(1-c^{2}\sin^{2}\varphi)}$; hence
\[ s=\frac{B}{b}E(c,\varphi). \]
This equation shows that every arc $s$ of the shortest line can be assimilated indefinitely to an arc of an ellipse whose semi-minor axis is $CP=B$ and whose semi-major axis is $CD=\frac{B}{b}$, a quantity greater than $B$ and less than $A$, since we have $\frac{B}{b}=A\frac{\sin l'}{\sin l}$. If on this quarter of an ellipse $PmD$ we take a point $m$ having $\varphi$ for its amplitude, or determined by the coordinates $Cr=B\cos\varphi$, $rm=\frac{B}{b}\sin\varphi$, the arc $Pm$ of the ellipse will be equal in length to the arc $AM$ of the shortest line.

These determinations hold throughout the whole extent of the curve we wish to construct; but it suffices to examine the course of this curve from $\varphi=0$ to $\varphi=90^{\circ}$, that is to say, from the point $A$ where we have $\lambda=l$, to the point $I$ where the curve meets the equator and where we have $\lambda=0$. At the point $A$ the curve is perpendicular to the meridian or tangent to the parallel; at the point $I$ it cuts the equator at an angle $I$ such that we have $\cos I=\sin M=\cos l'$, so that the angle $I$ is equal to the reduced latitude $l'$.

At this same point, the curve $AMI$ is equal to the quarter of an ellipse $BmD$\ednote{Printed ``BmD''; the quarter-ellipse named above is $PmD$, an apparent misprint.}, and we have $s=\frac{B}{b}E^{1}(c)$.

\subsection*{(128).}

Let us come now to the second formula, which gives the longitude $P$; we shall have first, putting $n=\operatorname{tang}^{2}l'$, $M=\frac{B}{bA\cos l'}=\frac{B}{A\cos l}$
\[ dP=M\cdot\frac{\Delta\,d\varphi}{1+n\sin^{2}\varphi}, \]

\origpage{181}
hence, integrating,
\[ P=\frac{M}{n}[(n+c^{2})\Pi(n,c,\varphi)-c^{2}F(c,\varphi)]. \]
By this formula we shall determine, for each point $M$ whose latitude is known, the longitude $P$ of that point; and the calculations can be carried out with all the degree of exactness that tables of logarithms allow, by means of the methods we have given for evaluating the functions $F$ and $\Pi$.

If we put $\varphi=90^{\circ}$, we shall have the position of the point $I$ where the curve cuts the equator by the formula
\[ P^{1}=\frac{M}{n}[(n+c^{2})\Pi^{1}(n,c)-c^{2}F^{1}(c)]; \]
and since the complete function $\Pi^{1}(n,c)$ can be expressed by functions of the first and of the second species, it follows that the point $I$ can be determined by these functions alone; or, what comes to the same thing, by arcs of an ellipse alone. We could determine in the same way an infinity of other points of the shortest line; and when the half-spire $AMI$ has been traced, we can likewise trace the continuation of this curve in the other half-spheroid, observing that the points situated on either side at equal latitudes have equal differences in longitude from the point $I$.

\subsection*{(129).}

It is remarkable that the formula which gives the value of the angle $P$ is absolutely of the same form as that which gives the value of the angle resulting from the development on a plane of a portion of the surface of an oblique cone with circular base (123).

This latter angle is easy to find to a certain degree of approximation, by simple geometrical constructions deduced from inscribed and circumscribed pyramids. Moreover, since the oblique cone has curvature in one sense only, this surface is held to be simpler than that of the spheroid, which is everywhere of double curvature; thus the problem of determining the shortest line on the surface of the spheroid must be regarded as reduced to a lesser difficulty by the following construction, which is easily drawn from our formulas.

\origpage{182}
\rmfigure{figures/fig-13.png}{Fig.~13.}{A triangle pfp' with the line pq from p to the point q on fp', the point c on pq, a dotted curve from p passing below the base pp' and up through n to q, and the dotted lines cn and fn.}

Make an isosceles triangle $pfp'$ such that its base $pp'$ representing\ednote{Printed ``représentant''; ``représente'' is expected, an apparent misprint.} the axis $2B$, the side $pf$ or $p'f$ represents the radius of the equator $A$; draw $pq$ in such a way that the angle $p'pq$ has $c$ for its sine, or that we have $\sin p'pq=\sin l\sin fpp'$. On $pq$ as diameter describe a circle $qnp$ whose plane is perpendicular to the plane of the triangle $fpq$. The cone which has this circle for its base, and the point $f$ for its vertex, is the one whose surface, when developed on a plane, will serve to describe the shortest line on the surface of the given spheroid. For this purpose, let $QM$ be a parallel whose latitude is $\lambda$; we shall determine successively the angles $\lambda'$, $\varphi$ and $\omega$, either by the equations $\operatorname{tang}\lambda'=\frac{B}{A}\operatorname{tang}\lambda$, $\cos\varphi=\frac{\sin\lambda'}{\sin l'}$, $\operatorname{tang}\frac{1}{2}\omega=\operatorname{tang}(45^{\circ}-\frac{1}{2}l')\cdot\operatorname{tang}\frac{1}{2}\varphi$, or by geometrical constructions which represent them. The angle $\omega$ being found, we shall draw in the plane of the base the radius $cn$, in such a way that the angle $qcn=\omega$; then drawing the straight line $fn$, and developing on a plane the surface $qfn$, the angle $f$ of the sector produced by the development will be equal to the angle $P$ of the spheroidal triangle $APM$, and will thus serve to determine the longitude of the point $M$.

When $\varphi=90^{\circ}$, we shall have $\omega=90^{\circ}-l'$; then the angle produced by the development of the sector $qfn$ will be half the angle produced by the development of the whole surface $fpnq$; and since this last angle is always less than two right angles, it follows that the angle $P$, measured then by the arc $EI$, is less than a right angle, which agrees with the known properties of the shortest line traced on a flattened spheroid.

\section*{Determination of the area of the ellipsoid.}

\subsection*{(130).}

Let the equation of the proposed surface be
\[ \frac{x^{2}}{a^{2}}+\frac{y^{2}}{b^{2}}+\frac{z^{2}}{c^{2}}=1; \]
the expression of the area, contained between given limits, will depend on the double integral
\[ S=\iint dx\,dy\sqrt{\left(\frac{1+\frac{c^{2}-a^{2}}{a^{2}}\cdot\frac{x^{2}}{a^{2}}+\frac{c^{2}-b^{2}}{b^{2}}\cdot\frac{y^{2}}{b^{2}}}{1-\frac{x^{2}}{a^{2}}-\frac{y^{2}}{b^{2}}}\right)}; \]

\origpage{183}
but either one of the two integrations cannot be carried out without introducing elliptic functions which would not allow the second integral to be obtained. To make the first integration possible, at least by arcs of a circle or by logarithms, we could apply the method that we gave in the Mémoires de l'Académie des Sciences, year 1788; but we shall arrive more directly at the same end in the following manner.

Let $z=c\cos\theta$; we shall have $\frac{x^{2}}{a^{2}}+\frac{y^{2}}{b^{2}}=\sin^{2}\theta$; let then $y=b\sin\theta\cos\varphi$; we shall have $x=a\sin\theta\sin\varphi$. From these equations, which are equivalent to the equation of the surface, we must express the element of the area as a function of the two new variables $\theta$ and $\varphi$. And first, to have the value of $dx\,dy$, I differentiate the equation $x=a\sin\theta\sin\varphi$, supposing $\theta$ constant; I have $dx=a\,d\varphi\cos\varphi\sin\theta$; then we must have the value of $dy$, supposing $x$ constant; this is what the equation $\frac{x^{2}}{a^{2}}+\frac{y^{2}}{b^{2}}=\sin^{2}\theta$ will give, whence we draw $y\,dy=b^{2}d\theta\sin\theta\cos\theta$; hence $dx\,dy=ab\,d\varphi\,d\theta\sin\theta\cos\theta$.

Substituting this value, as well as those of $x$ and $y$, in the expression of the area $S$; putting moreover, for brevity, $\frac{a^{2}-c^{2}}{a^{2}}=\delta$, $\frac{b^{2}-c^{2}}{b^{2}}=\varepsilon$, we shall have
\[ S=\iint ab\,d\varphi\,d\theta\sin\theta\sqrt{[1-(\delta\sin^{2}\varphi+\varepsilon\cos^{2}\varphi)\sin^{2}\theta]}. \]
The integral with respect to $\theta$ is easy to find by logarithms; but since the formula that would result for the second integration is not among those that we know how to reduce to known functions, we shall content ourselves with integrating by series.

Let then $\delta\sin^{2}\varphi+\varepsilon\cos^{2}\varphi=p$; we shall have, developing the radical,
\[ S=\iint ab\,d\varphi\,d\theta\sin\theta\left(1-\tfrac{1}{2}p\sin^{2}\theta-\tfrac{1\cdot1}{2\cdot4}p^{2}\sin^{4}\theta-\tfrac{1\cdot1\cdot3}{2\cdot4\cdot6}p^{3}\sin^{6}\theta-\text{etc.}\right). \]
Now, integrating from $\theta=0$ to $\theta=\frac{1}{2}\pi$, we have
\[ \int d\theta\sin\theta=1,\qquad \int d\theta\sin^{3}\theta=\frac{2}{3},\qquad \int d\theta\sin^{5}\theta=\frac{2\cdot4}{3\cdot5},\qquad \text{etc.} \]
Hence there results from the first integration,

\origpage{184}
\[ S=\int ab\,d\varphi\left(1-\tfrac{1}{3}p-\tfrac{1}{3\cdot5}p^{2}-\tfrac{1}{5\cdot7}p^{3}-\tfrac{1}{7\cdot9}p^{4}-\text{etc.}\right). \]
This formula must be integrated anew from $\varphi=0$ to $\varphi=\frac{1}{2}\pi$; let then, within these limits, $\displaystyle\int p\,d\varphi=\frac{\pi}{2}\cdot P'$, $\displaystyle\int p^{2}\,d\varphi=\frac{\pi}{2}\cdot P''$, $\displaystyle\int p^{3}\,d\varphi=\frac{\pi}{2}\cdot P'''$, etc.; and we shall have, after multiplying by 8, the total area of the ellipsoid
\[ 8S=4\pi ab\left(1-\tfrac{1}{3}P'-\tfrac{1}{3\cdot5}P''-\tfrac{1}{5\cdot7}P'''-\tfrac{1}{7\cdot9}P^{\text{iv}}-\text{etc.}\right), \]
a formula in which there remain to be substituted the following values:
\[ \begin{aligned} P'&=\delta\cdot\tfrac{1}{2}+\varepsilon\cdot\tfrac{1}{2},\\ P''&=\delta^{2}\cdot\tfrac{1\cdot3}{2\cdot4}+\delta\varepsilon\cdot\tfrac{1}{2}\cdot\tfrac{1}{2}+\varepsilon^{2}\cdot\tfrac{1\cdot3}{2\cdot4},\\ P'''&=\delta^{3}\cdot\tfrac{1\cdot3\cdot5}{2\cdot4\cdot6}+\delta^{2}\varepsilon\cdot\tfrac{1\cdot3}{2\cdot4}\cdot\tfrac{1}{2}+\delta\varepsilon^{2}\cdot\tfrac{1}{2}\cdot\tfrac{1\cdot3}{2\cdot4}+\varepsilon^{3}\cdot\tfrac{1\cdot3\cdot5}{2\cdot4\cdot6},\\ &\text{etc.} \end{aligned} \]
The law of these expressions is manifest, and we may observe that the general term $P^{(n)}$ is the coefficient of $z^{n}$ in the development of $(1-\delta z)^{-\frac{1}{2}}(1-\varepsilon z)^{-\frac{1}{2}}$, or in that of the product
\[ \left(1+\tfrac{1}{2}\delta z+\tfrac{1\cdot3}{2\cdot4}\delta^{2}z^{2}+\tfrac{1\cdot3\cdot5}{2\cdot4\cdot6}\delta^{3}z^{3}+\text{etc.}\right)\left(1+\tfrac{1}{2}\varepsilon z+\tfrac{1\cdot3}{2\cdot4}\varepsilon^{2}z^{2}+\tfrac{1\cdot3\cdot5}{2\cdot4\cdot6}\varepsilon^{3}z^{3}+\text{etc.}\right). \]

For the rest, the series $P'$, $P''$, $P'''$, etc. will necessarily be convergent if $\delta$ and $\varepsilon$ are both smaller than unity, which will always be the case by taking for $c$ the smallest of the three quantities $a$, $b$, $c$.

\subsection*{(131).}

The problem of determining the total area of the ellipsoid is solved by the series just found, as simply as it can be when one has no aim but to obtain an approximation. But we must have recourse to other means if we want a result depending only on elliptic functions or on algebraic quadratures.

\origpage{185}
In the preceding method we divided the ellipse which serves as base to the proposed surface into concentric elliptic zones, determined by the projections of the sections made in the ellipsoid parallel to its base. Any other way of dividing the base is equally admissible and must lead to the same result after the two integrations; but there is one in particular which deserves to be submitted to calculation, and which seems to promise elegant results. I mean the projections which arise from the lines of greatest and of least curvature traced on the surface of the ellipsoid.

Monge, in his Feuilles d'Analyse appliquée (edit.\ of the year 9, no.~19), has given the construction of these curves, as follows.

Having supposed $a>b$ and $b>c$, let the plane of $x$ and $y$ be taken as plane of projection, which is that of the principal section of which $a$ and $b$ are the semi-axes. Let there be described in this plane an auxiliary ellipse and hyperbola, having the common semi-axes $A$ and $B$ thus determined,
\[ \begin{aligned} A&=a\sqrt{\left(\frac{a^{2}-b^{2}}{a^{2}-c^{2}}\right)},\\ B&=b\sqrt{\left(\frac{a^{2}-b^{2}}{b^{2}-c^{2}}\right)}, \end{aligned} \]
the semi-axis $A$ being directed along the semi-axis $a$, and $B$ along $b$.

The equation of the hyperbola will be $y^{2}=\frac{B^{2}}{A^{2}}(x^{2}-A^{2})$, and the equation of the ellipse, $y^{2}=\frac{B^{2}}{A^{2}}(A^{2}-x^{2})$; they will touch at the vertex where $x=A$, and besides we may observe that by virtue of the suppositions made we have $A<a$.

This being set down, let $x=\alpha$, $y=\beta$ be the coordinates of a point taken at will on the hyperbola, so that we have $\beta^{2}=\frac{B^{2}}{A^{2}}(\alpha^{2}-A^{2})$; if, with the semi-axes $\alpha$ and $\beta$, taken in the same directions as $a$ and $b$, an ellipse is described, this ellipse, and all the ellipses described in the same manner according to the various values of $\alpha$ and $\beta$, with the sole condition that we have $\alpha<a$, will be the projections on the plane of $x$ and $y$ of all the lines of curvature of the same species, traced on the surface of the ellipsoid.

\origpage{186}
The projections of the lines of curvature of the other species, which will be hyperbolas, will be determined similarly by taking for the semi-axes of each hyperbola the coordinates $\alpha'$, $\beta'$ of one and the same point of the auxiliary ellipse, which must satisfy the equation $\beta'^{2}=\frac{B^{2}}{A^{2}}(A^{2}-\alpha'^{2})$.

The projections of the lines of curvature of the first species will therefore be the ellipses traced according to the equation
\[ y^{2}=\frac{\beta^{2}}{\alpha^{2}}(\alpha^{2}-x^{2})=\frac{B^{2}}{A^{2}}(\alpha^{2}-x^{2})\left(1-\frac{A^{2}}{\alpha^{2}}\right), \]
giving to $\alpha$ all the values from $\alpha=A$ to $\alpha=a$.

And the projections of the lines of curvature of the second species will be all the hyperbolas described according to the equation
\[ y^{2}=\frac{\beta'^{2}}{\alpha'^{2}}(x^{2}-\alpha'^{2})=\frac{B^{2}}{A^{2}}(x^{2}-\alpha'^{2})\left(\frac{A^{2}}{\alpha'^{2}}-1\right), \]
giving to $\alpha'$ all the values from $\alpha'=0$ to $\alpha'=A$.

Let us consider more particularly the ellipses which are the projections of the lines of curvature of the first species.

When we put $\alpha=A$, we have $\beta=0$, $y=0$, so that the ellipse is infinitely narrow, and reduces to its major axis. As the semi-axis $\alpha$ increases, the other semi-axis $\beta$ also increases; and finally when we put $\alpha=a$, the ellipse of projection coincides with the base of the ellipsoid, which is the last term of the projections.

\subsection*{(132).}

This being set down, let us seek the area of the ellipsoid which corresponds to the element of the base $dx\,dy$, contained between two infinitely near ellipses of projection. One of these ellipses having for its equation $y^{2}=\frac{\beta^{2}}{\alpha^{2}}(\alpha^{2}-x^{2})$, we may put $x=\alpha\sin\psi$, which will give $y=\beta\cos\psi$. Differentiating $x$ supposing $\alpha$ constant, we shall have $dx=\alpha\,d\psi\cos\psi$; if we then pass from this ellipse to the infinitely near ellipse, we must differentiate the equation
\[ y^{2}=\frac{B^{2}}{A^{2}}(\alpha^{2}-x^{2})\left(1-\frac{A^{2}}{\alpha^{2}}\right), \]
making $\alpha$ alone vary, which will give $y\,dy=\frac{B^{2}}{A^{2}}\alpha\,d\alpha\left(1-\frac{A^{2}x^{2}}{\alpha^{4}}\right)$,

\origpage{187}
hence the element of the projection
\[ dx\,dy=\frac{B^{2}}{A^{2}}\cdot\frac{d\alpha\,d\psi}{\beta}(\alpha^{2}-A^{2}\sin^{2}\psi). \]
When $\psi$ is made to vary supposing $\alpha$ constant, the point of projection varies only on one and the same ellipse of projection; likewise when $\alpha$ is made to vary supposing $\psi$ constant, the point of projection varies only on a hyperbola of projection; for, eliminating $\alpha$ and $\beta$ from the equations $\alpha=\frac{x}{\sin\psi}$, $\beta=\frac{y}{\cos\psi}$, $\beta^{2}=\frac{B^{2}}{A^{2}}(\alpha^{2}-A^{2})$, we have $y^{2}=\frac{B^{2}\cos^{2}\psi}{A^{2}\sin^{2}\psi}(x^{2}-A^{2}\sin^{2}\psi)$, which is the equation of a hyperbola of projection, whose semi-axes are $\alpha'=A\sin\psi$, $\beta'=B\cos\psi$.

From this we see that the element of projection whose value we have just given represents the quadrilateral contained between two ellipses and two infinitely near hyperbolas of projection.

If we now substitute the values of $x$, $y$ and $dx\,dy$ in the general expression of $S$, we shall have
\[ S=\iint\frac{B}{A}\cdot\frac{\alpha^{2}-A^{2}\sin^{2}\psi}{\sqrt{(\alpha^{2}-A^{2})}}\,d\alpha\,d\psi\sqrt{\left[\frac{1-\delta\frac{\alpha^{2}\sin^{2}\psi}{a^{2}}-\varepsilon\frac{\beta^{2}\cos^{2}\psi}{b^{2}}}{1-\frac{\alpha^{2}\sin^{2}\psi}{a^{2}}-\frac{\beta^{2}\cos^{2}\psi}{b^{2}}}\right]}. \]
This expression appears at first glance much more complicated than that to which we were led by the first method; but on examining it more closely, we find that it is susceptible of a very remarkable reduction.

Indeed, if we substitute the value $\beta^{2}=\frac{B^{2}}{A^{2}}(\alpha^{2}-A^{2})$, we shall recognize that the quantity under the radical decomposes into two distinct factors, one depending only on $\alpha$, the other depending only on $\psi$, so that we have
\[ S=\iint\frac{B}{A}\cdot\frac{\alpha^{2}-A^{2}\sin^{2}\psi}{\sqrt{(\alpha^{2}-A^{2})}}\,d\alpha\,d\psi\sqrt{\left[\frac{\sin^{2}\psi+\frac{a^{2}}{b^{2}}\cos^{2}\psi}{\sin^{2}\psi+\frac{B^{2}a^{2}}{A^{2}b^{2}}\cos^{2}\psi}\cdot\frac{1-\delta\frac{\alpha^{2}}{a^{2}}}{1-\frac{\alpha^{2}}{a^{2}}}\right]}. \]
If we now observe that for the proposed object we must take the integral with respect to $\psi$ from $\psi=0$ to $\psi=\frac{1}{2}\pi$, and

\origpage{188}
the integral with respect to $\alpha$ from $\alpha=A$ to $\alpha=a$, we shall see that these two integrals are entirely independent of one another, and that they can be sought separately, which will allow the entire area of the ellipsoid to be reduced to algebraic quadratures, and even, as will be proved afterwards, to elliptic functions alone.

\subsection*{(133).}

Let then
\[ \begin{aligned} M&=\int d\psi\sqrt{\left(\frac{\sin^{2}\psi+\frac{a^{2}}{b^{2}}\cos^{2}\psi}{\sin^{2}\psi+\frac{B^{2}a^{2}}{A^{2}b^{2}}\cos^{2}\psi}\right)}\\ N&=\int d\psi\sin^{2}\psi\sqrt{\left(\frac{\sin^{2}\psi+\frac{a^{2}}{b^{2}}\cos^{2}\psi}{\sin^{2}\psi+\frac{B^{2}a^{2}}{A^{2}b^{2}}\cos^{2}\psi}\right)}; \end{aligned} \]
these integrals having to be taken from $\psi=0$ to $\psi=\frac{1}{2}\pi$.

Let there be further
\[ \begin{aligned} P&=\int\frac{\alpha^{2}d\alpha}{\sqrt{(\alpha^{2}-A^{2})}}\sqrt{\left(\frac{a^{2}-\delta\alpha^{2}}{a^{2}-\alpha^{2}}\right)}\\ Q&=\int\frac{d\alpha}{\sqrt{(\alpha^{2}-A^{2})}}\sqrt{\left(\frac{a^{2}-\delta\alpha^{2}}{a^{2}-\alpha^{2}}\right)}, \end{aligned} \]
these integrals having to be taken from $\alpha=A$ to $\alpha=a$.

The required surface will be
\[ S=\frac{B}{A}PM-ABQN. \]
Thus everything reduces to finding the four integrals $P$, $Q$, $M$, $N$, within the designated limits.

To simplify the formulas $M$ and $N$, let $\operatorname{tang}\psi=\frac{a}{b}\operatorname{tang}\omega$, then $n=\frac{a^{2}-b^{2}}{b^{2}}$ and $c'^{2}=1-\frac{A^{2}}{B^{2}}=\frac{c^{2}}{b^{2}}\cdot\frac{a^{2}-b^{2}}{a^{2}-c^{2}}$; there will result from these substitutions
\[ \begin{aligned} M&=\frac{aA}{bB}\int\frac{d\omega}{(1+n\sin^{2}\omega)\sqrt{(1-c'^{2}\sin^{2}\omega)}}=\frac{aA}{bB}\Pi(n,c',\omega)\\ N&=\frac{a^{3}A}{b^{3}B}\int\frac{d\omega\sin^{2}\omega}{(1+n\sin^{2}\omega)^{2}\sqrt{(1-c'^{2}\sin^{2}\omega)}}. \end{aligned} \]

\origpage{189}
But by the known reductions we have
\[ \begin{aligned} \int\frac{d\omega\sin^{2}\omega}{(1+n\sin^{2}\omega)^{2}\Delta(c',\omega)}&=-\frac{n\sin\omega\cos\omega\,\Delta(c',\omega)}{2(n+1)(n+c'^{2})(1+n\sin^{2}\theta)}+\frac{F(c',\theta)}{2n(n+1)}\\ &\quad-\frac{E(c',\theta)}{2(n+1)(n+c'^{2})}+\frac{n^{2}-c'^{2}}{2n(n+1)(n+c'^{2})}\Pi(n,c',\theta). \end{aligned} \]
Hence, extending the integrals to $\theta=\frac{1}{2}\pi$, we shall have for the complete values of $M$ and $N$,
\[ \begin{aligned} M&=\frac{aA}{bB}\Pi^{1}(n,c')\\ N&=\frac{aA}{2bBn}\left[F^{1}(c')-\left(\frac{a^{2}-c^{2}}{a^{2}}\right)E^{1}(c')+\left(\frac{a^{2}-b^{2}-c^{2}}{b^{2}}\right)\Pi^{1}(n,c')\right]. \end{aligned} \]

\subsection*{(134).}

As regards the integrals $P$, $Q$, since $\alpha$ is always contained between $A$ and $a$, if one puts $A=a\sin\mu$, we may take
\[ \alpha^{2}=\frac{a^{2}\sin^{2}\mu}{1-\cos^{2}\mu\sin^{2}\zeta}, \]
and the substitution will give
\[ \begin{aligned} P&=ab\sin^{2}\mu\int\frac{d\zeta\sqrt{(1-b'^{2}\sin^{2}\zeta)}}{(1-\cos^{2}\mu\sin^{2}\zeta)^{2}}\\ Q&=\frac{b}{a}\int\frac{d\zeta\sqrt{(1-b'^{2}\sin^{2}\zeta)}}{1-\cos^{2}\mu\sin^{2}\zeta}. \end{aligned} \]
In these formulas we have put $b'^{2}=\frac{a^{2}}{b^{2}}\cdot\frac{b^{2}-c^{2}}{a^{2}-c^{2}}=\frac{A^{2}}{B^{2}}$, so that we have $b'^{2}+c'^{2}=1$, and thus the moduli $b'$ and $c'$ are complements of one another.

Let the parameter $n=\frac{a^{2}-b^{2}}{b^{2}}=\cot^{2}\lambda$; we shall have $b=a\sin\lambda$, and $b'^{2}=\frac{a^{2}}{b^{2}}\cos^{2}\mu=\frac{\cos^{2}\mu}{\sin^{2}\lambda}$, or $\cos^{2}\mu=b'^{2}\sin^{2}\lambda$; the integral $Q$ therefore relates to elliptic functions of the third species, whose parameter is $n'=-\cos^{2}\mu=-b'^{2}\sin^{2}\lambda$, and we shall find
\[ Q=\frac{1}{\sin\lambda}F(b',\zeta)-\frac{\cos^{2}\lambda}{\sin\lambda}\Pi(n',b',\zeta). \]
To have the value of $P$, we must first employ this formula

\origpage{190}
of reduction
\[ \begin{aligned} \int\frac{d\zeta\sqrt{(1-b'^{2}\sin^{2}\zeta)}}{(1+n'\sin^{2}\zeta)^{2}}&=\frac{n'\sin\zeta\cos\zeta\,\Delta(b',\zeta)}{2(1+n')(1+n'\sin^{2}\zeta)}-\frac{(b'^{2}+n')}{2n'(1+n')}F(b',\zeta)\\ &\quad+\frac{E(b',\zeta)}{2(1+n')}+\frac{n'^{2}+2n'+b'^{2}}{2n'(1+n')}\Pi(n',b',\zeta). \end{aligned} \]
Making\ednote{Printed ``Faisait''; ``Faisant'' is expected, an apparent misprint.} then $\zeta=\frac{1}{2}\pi$, we shall have the complete integrals
\[ \begin{aligned} P&=\frac{a^{2}}{2\sin\lambda}[\cos^{2}\lambda F^{1}(b')+\sin^{2}\lambda E^{1}(b')-(1-2\sin^{2}\lambda+b'^{2}\sin^{4}\lambda)\Pi^{1}(n',b')]\\ Q&=\frac{1}{\sin\lambda}F^{1}(b')-\frac{\cos^{2}\lambda}{\sin\lambda}\Pi^{1}(n',b'). \end{aligned} \]

\subsection*{(135).}

Now if in the expression of the area $S$ we substitute the values found for $M$, $N$, $P$, $Q$, there will result
\[ \begin{aligned} S={}&\frac{a^{2}}{2}\Pi^{1}(n,c')[\cot^{2}\lambda F^{1}(b')+E^{1}(b')+(\sin^{2}\mu-\cot^{2}\lambda)\Pi^{1}(n',b')]\\ &-\frac{a^{2}}{2}[F^{1}(b')-\cos^{2}\lambda\Pi^{1}(n',b')]\left[\frac{\sin^{2}\mu}{\cos^{2}\lambda}F^{1}(c')-E^{1}(c')+\frac{\cot^{2}\lambda-\sin^{2}\mu}{\cos^{2}\lambda}\Pi^{1}(n',c')\right]. \end{aligned} \]
By a first development which makes disappear the terms in which the two functions $\Pi^{1}$ are multiplied, we obtain
\[ \begin{aligned} S={}&\frac{a^{2}}{2}\Pi^{1}(n,c')\left[\left(\frac{\sin^{2}\mu}{\cos\lambda}-1\right)F^{1}(b')+E^{1}(b')\right]\\ &-\frac{a^{2}}{2}\left[\frac{\sin^{2}\mu}{\cos^{2}\lambda}F^{1}(c')-E^{1}(c')\right][F^{1}(b')-\cos^{2}\lambda\Pi(n',b')]. \end{aligned} \]
But by the known values of the functions $\Pi^{1}$, we have
\[ \begin{aligned} \Pi^{1}(n,c')={}&\sin^{2}\lambda F^{1}(c')+\frac{\sin\lambda\cos\lambda}{\Delta(b',\lambda)}[\tfrac{1}{2}\pi+F^{1}(c)F(b',\lambda)-F^{1}(c)E(b',\lambda)]\\ &-E^{1}(c)F(b',\lambda)\\ F^{1}(b)-\cos^{2}\lambda\Pi^{1}(n',b')={}&\sin^{2}\lambda F^{1}(b')+\frac{\sin\lambda\cos\lambda}{\Delta(b',\lambda)}[F^{1}(b')E(b',\lambda)\\ &-E^{1}(b')F(b',\lambda)]. \end{aligned} \]
Substituting and making the reductions to which the known formula $\frac{\pi}{2}=F^{1}(c')E^{1}(b')+F^{1}(b')E^{1}(c')-F^{1}(b')F^{1}(c')$ gives rise, we

\origpage{191}
we shall find
\[ \begin{aligned} S={}&\frac{\pi a^{2}}{4}\cdot\frac{\sin\lambda\cos\lambda}{\Delta(b',\lambda)}\left[\left(\frac{\sin^{2}\mu}{\cos^{2}\lambda}-1\right)[F^{1}(b')-F(b',\lambda)]\right.\\ &\left.+E^{1}(b')-E(b',\lambda)+\frac{\sin\lambda}{\cos\lambda}\Delta(b',\lambda)\right]. \end{aligned} \]
I observe that this result can be simplified still further by taking a new angle $\nu$ such that $F(b',\lambda)+F(b',\nu)=F^{1}(b')$; for this it is necessary that $c'\operatorname{tang}\lambda\operatorname{tang}\nu=1$, which gives $\cos\nu=\frac{c}{a}$; at the same time we shall have $E(b',\lambda)+E(b',\nu)=E^{1}(b')+\frac{\sin\nu\cos\nu}{\Delta(b',\nu)}$. Hence finally the entire area of the ellipsoid will have for its expression
\[ 8S=2\pi a^{2}\cdot\frac{\sin\nu\cos\nu}{\Delta(b',\nu)}\left[\frac{\cos^{2}\nu}{\sin^{2}\nu}F(b',\nu)+E(b',\nu)+\frac{\cos\nu}{\sin\nu}\Delta(b',\nu)\right]. \]
Moreover, since we have $b'^{2}=\frac{a^{2}}{b^{2}}\cdot\frac{b^{2}-c^{2}}{a^{2}-c^{2}}=\frac{b^{2}-c^{2}}{b^{2}\sin^{2}\nu}$, it follows that $\Delta(b',\nu)=\frac{c}{b}$, so that the value of $8S$ reduces to this very simple form,
\[ 8S=2\pi c^{2}+\frac{2\pi ab}{\sin\nu}\left[\frac{c^{2}}{a^{2}}F(b',\nu)+\frac{a^{2}-c^{2}}{a^{2}}E(b',\nu)\right], \]
a formula which can be applied immediately, remembering only that we have $\cos\nu=\frac{c}{a}$ and $b'^{2}=\frac{b^{2}-c^{2}}{b^{2}\sin^{2}\nu}$.

It is easy to verify this formula in the two cases where the ellipsoid becomes a solid of revolution; for if we have $a=b$, the formula gives for the area of the oblate spheroid,
\[ 8S=2\pi a^{2}+\frac{\pi b^{2}}{\sin\nu}\log\left(\frac{1+\sin\nu}{1-\sin\nu}\right); \]
and if we have $b=c$, the formula gives for the area of the prolate spheroid,
\[ 8S=\frac{2\pi ab}{\sin\nu}(\nu+\sin\nu\cos\nu), \]
which agrees with the known results.

We have therefore succeeded in expressing by very simple elliptic functions the entire area of the ellipsoid, which, according to the

\origpage{192}
other methods, seemed bound to be a much more complicated transcendent.

\subsection*{(136).}

It follows from this solution that the series found by the first method,
\[ 4\pi ab\left(1-\tfrac{1}{3}P'-\tfrac{1}{3\cdot5}P''-\tfrac{1}{5\cdot7}P'''-\text{etc.}\right), \]
is capable of being summed by means of elliptic functions; and this can be verified directly in the following manner.

Consider the function $\Gamma(y)$ or simply $\Gamma$, whose developed value would be
\[ \Gamma=y-\tfrac{1}{3}P'y^{3}-\tfrac{1}{3\cdot5}P''y^{5}-\tfrac{1}{5\cdot7}P'''y^{7}-\text{etc.}; \]
by successive differentiations we shall have,
\[ \begin{aligned} \frac{d\Gamma}{dy}&=1-P'y^{2}-\tfrac{1}{3}P''y^{4}-\tfrac{1}{5}P'''y^{6}-\text{etc.}\\ \frac{1}{y^{2}}\cdot\frac{d\Gamma}{dy}-\frac{1}{y}\cdot\frac{dd\Gamma}{dy^{2}}&=\frac{1}{y^{2}}+P'+P''y^{2}+P'''y^{4}+\text{etc.} \end{aligned} \]
But if we go back to the origin of the quantities $P'$, $P''$, $P'''$, etc., we shall find that the sequence $\frac{\pi}{2}(1+P'y^{2}+P''y^{4}+P'''y^{6}+\text{etc.})$ is nothing other than the integral, taken from $\varphi=0$ to $\varphi=\frac{1}{2}\pi$, of the differential $d\varphi(1+py^{2}+p^{2}y^{4}+p^{3}y^{6}+\text{etc.})=\frac{d\varphi}{1-py^{2}}=\frac{d\varphi}{1-(\delta\sin^{2}\varphi+\varepsilon\cos^{2}\varphi)y^{2}}$. This integral has for its value
\[ \frac{\tfrac{1}{2}\pi}{\sqrt{[(1-\delta y^{2})(1-\varepsilon y^{2})]}}. \]
Hence if, for brevity, we put $R=\sqrt{[1-(\delta+\varepsilon)y^{2}+\delta\varepsilon y^{4}]}$, we shall have
\[ \frac{1}{y^{2}}\cdot\frac{d\Gamma}{dy}-\frac{1}{y}\cdot\frac{dd\Gamma}{dy^{2}}=\frac{1}{y^{2}R}; \]
whence, multiplying by $dy$ and integrating,
\[ \frac{1}{y}\cdot\frac{d\Gamma}{dy}=\int-\frac{dy}{y^{2}R}=\frac{R}{y}-\delta\varepsilon\int\frac{y^{2}\,dy}{R}. \]

\origpage{193}
Multiplying by $y\,dy$ and integrating anew, we shall have
\[ \Gamma=\int R\,dy-\delta\varepsilon\cdot\frac{y^{2}}{2}\int\frac{y^{2}\,dy}{R}+\frac{\delta\varepsilon}{2}\int\frac{y^{4}\,dy}{R}, \]
or, what comes to the same thing,
\[ \Gamma=\frac{\delta\varepsilon}{2}(1-y^{2})\int\frac{y^{2}\,dy}{R}+\tfrac{1}{2}yR+\int\frac{dy}{2R}(1-\delta\varepsilon y^{2}). \]

Now let $y\sqrt{\delta}=\sin\psi$ and $k^{2}=\frac{\varepsilon}{\delta}=\frac{a^{2}}{b^{2}}\cdot\frac{b^{2}-c^{2}}{a^{2}-c^{2}}$; we shall have $\frac{dy}{R}=\frac{1}{\sqrt{\delta}}\cdot\frac{d\psi}{\sqrt{(1-k^{2}\sin^{2}\psi)}}$, and
\[ \int\frac{dy}{R}(1-\delta\varepsilon y^{2})=\frac{1-\delta}{2\sqrt{\delta}}F(k,\psi)+\frac{\sqrt{\delta}}{2}E(k,\psi). \]
It remains to make $y=1$ in these expressions, which will give $\sin\psi=\sqrt{\delta}$, $\cos\psi=\sqrt{(1-\delta)}=\frac{c}{a}$, $R=\sqrt{(1-\delta)}\cdot\sqrt{(1-\varepsilon)}=\frac{c^{2}}{ab}$; hence the value sought
\[ \Gamma=\tfrac{1}{2}\cdot\frac{c^{2}}{ab}+\frac{\cos^{2}\psi}{2\sin\psi}F(k,\psi)+\frac{\sin^{2}\psi}{2\sin\psi}E(k,\psi); \]
whence finally results the area of the ellipsoid
\[ 8S=2\pi c^{2}+\frac{2\pi ab}{\sin\psi}[\cos^{2}\psi F(k,\psi)+\sin^{2}\psi E(k,\psi)]; \]
and we see that this formula agrees entirely with the one we found by the second method, since the quantities $k$ and $\psi$ do not differ from those designated above by $b'$ and $\nu$.

We shall remark, finally, that one can determine in general the area of any quadrilateral whatever contained between two lines of curvature of one species and two lines of curvature of the other species. For this it suffices to substitute in the formula
\[ S=\frac{B}{A}\cdot PM-ABQN, \]
the values of the integrals $M$, $N$, $P$, $Q$, taken between the given limits. One must therefore take the two integrals $M$, $N$ between the two values of $\psi$ which correspond to the sides of the quadrilateral whose projections are hyperbolas, and the two integrals $P$ and $Q$ between

\origpage{194}
the two values of $x$ which correspond to the sides of the quadrilateral whose projections are ellipses. Thus all the quadrilaterals in question can be determined by elliptic functions, and the results will become simpler if they are applied to the entire zone contained between two lines of curvature of the same species.

Figure 14 may serve to make the constructions which this chapter supposes better understood.

\rmfigure{figures/fig-14.png}{Fig.~14.}{A quadrant with centre C and axes CA and CB, showing the ellipse BKA, the ellipse GNO, the hyperbola from O through M and K to I, the ellipses from b through n and m to a and from b' through n' and m' to a', the curves rising from e' through n, n' and N and from e through m and m', and the dotted lines BI, IA, fN, Ne, bM and Ma.}

$AKB$ is the principal ellipse whose semi-axes are $CA=a$, $CB=b$.

The auxiliary ellipse and hyperbola are represented by $ONG$ and $OMI$; they have for common axes $CO=A$, $CG=B$.

Having taken on the auxiliary hyperbola $OMI$ any point $M$ whose coordinates are $Ca=\alpha$, $aM=\beta$, one describes with the semi-axes $Ca=\alpha$, $Cb=\beta$, the ellipse $am\beta$\ednote{Printed $am\beta$ with the Greek letter; the ellipse through the points $a$, $m$, $b$ is presumably meant, $amb$, an apparent misprint, to be checked against the plate. Likewise $a'm'\beta'$ below.} which will be the projection of one of the lines of curvature of the first species.

The ellipse $a'm'\beta'$ represents the projection of another line of curvature of the same species.

Having taken on the auxiliary ellipse $ONG$ any point $N$ whose coordinates are $Ce=\alpha'$, $eN=\beta'$, one describes with the semi-axes $Ce=\alpha'$, $Cf=\beta'$, a hyperbola $emm'$ which will be the projection of one of the lines of curvature of the second species.

The hyperbola $e'nn'$ represents the projection of another line of curvature of the second species.

We have shown that one can determine in general by elliptic functions the area of any quadrilateral formed on the surface of the ellipsoid, which has for its projection $mnm'n'$.

\section*{On some general formulas which can be reduced to elliptic functions.}

\subsection*{(137).}

The nature of elliptic functions being known and thoroughly examined, it is important to reduce to these quantities as many integral formulas as possible. As their multitude is infinite and as varied as the substitutions that can be put to use, we shall content ourselves with indicating some cases in which this reduction takes place.

I. One can reduce to elliptic functions the integral\ldots\ $\displaystyle\int\frac{P\,dx}{\sqrt{(\alpha+\beta x^{2}+\gamma x^{4}+\delta x^{6})}}$, in which $P$ is a rational function of $x$.

\origpage{195}
For if we put $x^{2}=z$, we may suppose $P=M+N\sqrt{z}$, $M$ and $N$ being rational functions of $z$, and the proposed integral will become
\[ \int\frac{\tfrac{1}{2}M\,dz}{\sqrt{(\alpha z+\beta z^{2}+\gamma z^{3}+\delta z^{4})}}+\int\frac{\tfrac{1}{2}N\,dz}{\sqrt{(\alpha+\beta z+\gamma z^{2}+\delta z^{3})}}, \]
whose two parts are comprised in elliptic functions.

II. Every formula $\displaystyle\int\frac{P\,dx}{\sqrt[4]{(\alpha+\beta x^{2}+\gamma x^{4})}}$, in which $P$ is a rational function of $x$, can be reduced to elliptic functions.

For we can always put $P=M+Nx$, $M$ and $N$ being even rational functions of $x$. Consider the part
\[ \frac{Nx\,dx}{\sqrt[4]{(\alpha+\beta x^{2}+\gamma x^{4})}}. \]
If we put $\sqrt[4]{(\alpha+\beta x^{2}+\gamma x^{4})}=z$, which gives
\[ x^{2}=-\frac{\beta+\sqrt{(\beta^{2}-4\alpha\gamma+4\gamma z^{4})}}{2\gamma}, \]
it is clear that by the substitution of the value of $x^{2}$, $Nx\,dx$ will contain no other radical than $\sqrt{(\beta^{2}-4\alpha\gamma+4\gamma z^{4})}$; hence the whole difficulty reduces to integrating a quantity of the form
\[ \frac{Qz^{3}\,dz}{\sqrt{(\beta^{2}-4\alpha\gamma+4\gamma z^{4})}}, \]
in which $Q$ is a rational function of $z^{4}$.

As for the part $\frac{M\,dx}{\sqrt[4]{(\alpha+\beta x^{2}+\gamma x^{4})}}$, if we put $\sqrt[4]{(\alpha+\beta x^{2}+\gamma x^{4})}=xy$, we shall have
\[ \begin{aligned} x^{2}&=\frac{-\beta+\sqrt{(\beta^{2}-4\alpha\gamma+4\alpha y^{4})}}{2(\gamma-y^{4})}\\ \frac{dx}{xy}&=\frac{2x^{2}y^{2}\,dy}{\sqrt{(\beta^{2}-4\alpha\gamma+4\alpha y^{4})}}. \end{aligned} \]
From which we see that the transformed expression in $y$ will contain a rational part and one of the form
\[ \frac{Ry^{2}\,dy}{\sqrt{(\beta^{2}-4\alpha\gamma+4\alpha y^{4})}}, \]

\origpage{196}
$R$ being a rational function of $y^{4}$; hence the proposed formula is always reducible to elliptic functions.

III. One would resolve in absolutely the same manner the formula $\displaystyle\int P\,dx\sqrt[4]{(\alpha+\beta x^{2}+\gamma x^{4})}$, $P$ being a rational function of $x$.

These two cases comprise the formula $\displaystyle\int P\,dx(\alpha+\beta x^{2}+\gamma x^{4})^{\pm\frac{3}{4}}$, and also the formula $\displaystyle\int Q\,dy(\alpha+\beta y+\gamma y^{2})^{\pm\frac{1}{4}}$, $Q$ being a rational function of $y$; for by putting $y=x^{2}$, this last becomes a particular case of the other two.

IV. One can reduce to elliptic functions the formula $\displaystyle\int P\,dx(\alpha+\beta x+\gamma x^{2}+\delta x^{3})^{\pm\frac{1}{3}}$, $P$ being a rational function of $x$.

This reduction can be made in several ways; and first, if we make one or other of the suppositions
\[ \begin{aligned} \sqrt[3]{(\alpha+\beta x+\gamma x^{2}+\delta x^{3})}&=\sqrt[3]{\alpha}+z\\ \sqrt[3]{(\alpha+\beta x+\gamma x^{2}+\delta x^{3})}&=x\sqrt[3]{\delta}+z, \end{aligned} \]
\ednote{The first supposition is printed $\sqrt[3]{\alpha}+z$; the ``$\cdot x$'' after it in the scan is a reader's ink addition. The sense requires $\sqrt[3]{\alpha}+zx$.}
the transformed expression in $z$ will be comprised in elliptic functions. One can also make the coefficient $\alpha$ or the coefficient $\delta$ disappear at will from under the radical; for this one must put $x=m+y$, or $x=m+\frac{1}{y}$, and take for $m$ a real root of the equation $\alpha+\beta m+\gamma m^{2}+\delta m^{3}=0$. Suppose that $\delta$ has been made to disappear in this manner; then we shall put $\alpha+\beta x+\gamma x^{2}=z^{3}$, which gives
\[ x=-\frac{\beta+\sqrt{(\beta^{2}-4\gamma+4\gamma z^{3})}}{2\gamma}; \]
\ednote{The radicand is printed $\beta^{2}-4\gamma+4\gamma z^{3}$, with a blank space after the 4; $\beta^{2}-4\alpha\gamma+4\gamma z^{3}$ is expected, as in the formula that follows, an apparent dropped letter.}
and by substituting this value in the proposed formula, the whole difficulty will reduce to integrating a differential of the form $\frac{Q\,dz}{\sqrt{(\beta^{2}-4\alpha\gamma+4\gamma z^{3})}}$, $Q$ being a rational function of $z$.

V. One can reduce to elliptic functions the formula\ldots\ $\displaystyle\int\frac{P\,dx}{\sqrt{(\alpha+\beta x+\gamma x^{2}+\delta x^{3}+\gamma x^{4}+\beta x^{5}+\alpha x^{6})}}$, $P$ being a rational function of $x$.

\origpage{197}

For if we put $x^{2}+1=xz$, this formula will become at first
\[ \int\frac{Px^{-\frac{3}{2}}\,dx}{\sqrt{[\alpha(z^{3}-3z)+\beta(z^{2}-2)+\gamma z+\delta]}}; \]
then, as we have $x=\frac{1}{2}z\pm\frac{1}{2}\sqrt{(z^{2}-4)}$, this value being substituted in $P$, the result will be of the form $M\pm N\sqrt{(z^{2}-4)}$, $M$ and $N$ being rational functions of $z$. Moreover we have
\[ \begin{aligned} x+1&=x^{\frac{1}{2}}\sqrt{(z+2)}\\ x-1&=\pm x^{\frac{1}{2}}\sqrt{(z-2)}\\ 2x^{-\frac{1}{2}}&=\sqrt{(z+2)}\mp\sqrt{(z-2)}\\ x^{-\frac{3}{2}}\,dx&=\frac{\tfrac{1}{2}dz}{\sqrt{(z+2)}}\mp\frac{\tfrac{1}{2}dz}{\sqrt{(z-2)}}. \end{aligned} \]
Hence, all the substitutions having been made, the proposed formula will be composed of two parts, one of the form $\displaystyle\int\frac{Z\,dz}{\sqrt{[Q(z+2)]}}$, $Z$ being a rational function of $z$, and $Q$ denoting the polynomial $\alpha(z^{3}-3z)+\beta(z^{2}-2)+\gamma z+\delta$, the other being of the form $\displaystyle\int\frac{Z'\,dz}{\sqrt{[Q(z-2)]}}$, $Z'$ being likewise a rational function of $z$. Hence the transformed expression in $z$ will be integrable by elliptic functions.

VI. One will therefore be able to integrate in the same way the formula\ldots\ $\displaystyle\int\frac{P\,dy}{\sqrt{(\beta+\gamma y^{2}+\delta y^{4}+\gamma y^{6}+\beta y^{8})}}$, $P$ being a rational function of $y$.

For if we put $P=M+Ny$, $M$ and $N$ being even functions of $y$, and call $R$ the radical, the part $\frac{M\,dy}{R}$ falls under the preceding case, by making $\alpha=0$ and $y^{2}=x$. The other part $\frac{Ny\,dy}{R}$ is likewise reduced to an elliptic function by making $y^{2}=x$.

VII. One can bring to elliptic functions the formula $\displaystyle\int\frac{P\,dx}{\sqrt{(\alpha+\beta x^{4}+\gamma x^{8})}}$, $P$ being a rational function of $x$, provided that $\alpha$ and $\gamma$ are of the same sign.

For by making $\gamma=\alpha m^{8}$ and $mx=y$, this formula will be found comprised in the preceding case.

\origpage{198}

VIII. Let the formula $Z=\displaystyle\int\frac{Q\,d\varphi}{\sqrt{(1-m\sin^{2}\varphi)}}$ be proposed, $Q$ being an even rational function of $\sin\varphi$, and $m$ being greater than unity.

To bring this formula to elliptic functions, let first $m=\frac{1}{\sin^{2}\alpha}$; we see that $\alpha$ must be the greatest value of $\varphi$. Let then $\sin\varphi=\sin\alpha\sin\psi$, and $c=\sin\alpha$; the transformed expression will be $Z=\displaystyle\int\frac{cQ\,d\psi}{\sqrt{(1-c^{2}\sin^{2}\psi)}}$, and $Q$ will become by the substitution an even rational function of $\sin\psi$, so that the proposed formula will be brought to elliptic functions.

IX. If we had the formula $Z=\displaystyle\int\frac{Q\,d\varphi}{\sqrt{(1+m\sin^{2}\varphi)}}$, $m$ being positive and of any magnitude whatever, it would suffice, in order to bring this quantity to elliptic functions, to put $\varphi=\frac{1}{2}\pi-\psi$, and $\frac{m}{1+m}=c^{2}$.

X. Let the formula $Z=\displaystyle\int\frac{Q\,d\varphi}{\sqrt{(f+g\cos\varphi+h\sin\varphi)}}$ be proposed, in which $f$, $g$, $h$ are constants, and $Q$ a rational function of $\sin\varphi$ and $\cos\varphi$.

We shall first put $\varphi=2\psi+m$, and take the indeterminate $m$ so that $\operatorname{tang}m=\frac{h}{g}$; then the quantity $f+g\cos\varphi+h\sin\varphi$ will become of the form $f'\pm g'\sin^{2}\psi$, and we shall have $Q\,d\varphi=(M+N\sin\psi\cos\psi)\,d\psi$, $M$ and $N$ being even functions of $\sin\psi$. The part $\frac{N\,d\psi\sin\psi\cos\psi}{\sqrt{(f'\pm g'\sin^{2}\psi)}}$ can be made rational by putting $f'\pm g'\sin^{2}\psi=x^{2}$, so that there will remain to be integrated only the differential $\frac{M\,d\psi}{\sqrt{(f'\pm g'\sin^{2}\psi)}}$, and it is evident that this integral depends only on elliptic functions.

XI. If we have more generally
\[ Z=\int\frac{Q\,d\varphi}{\sqrt{(\alpha+\beta\cos\varphi+\gamma\sin\varphi+\delta\cos^{2}\varphi+\varepsilon\sin\varphi\cos\varphi+\zeta\sin^{2}\varphi)}}, \]
$Q$ being a rational function of $\sin\varphi$ and $\cos\varphi$; we shall put $\operatorname{tang}\frac{1}{2}\varphi=z$, which gives $\sin\varphi=\frac{2z}{1+z^{2}}$, $\cos\varphi=\frac{1-z^{2}}{1+z^{2}}$, $d\varphi=\frac{2\,dz}{1+z^{2}}$, and the transformed expression in $z$ will contain only one radical under which the variable does not

\origpage{199}
exceed the fourth degree; it will therefore be generally reducible to elliptic functions.

XII. Let finally $Z=\displaystyle\int\frac{Q\,d\varphi}{\sqrt{[(\alpha^{2}+\beta^{2}\sin^{2}\varphi)(\gamma+\delta\sin^{2}\varphi)]}}$, $Q$ being an even rational function of $\sin\varphi$; if we put $\operatorname{tang}\varphi=m\operatorname{tang}\psi$, which gives $\sin^{2}\varphi=\frac{m^{2}\sin^{2}\psi}{m^{2}\sin^{2}\psi+\cos^{2}\psi}$, $\cos^{2}\varphi=\frac{\cos^{2}\psi}{m^{2}\sin^{2}\psi+\cos^{2}\psi}$, $d\varphi=\frac{m\,d\psi}{m^{2}\sin^{2}\psi+\cos^{2}\psi}$, the transformed expression will be
\[ Z=\int\frac{Q\,d\psi}{\sqrt{(\alpha^{2}m^{2}\sin^{2}\psi+\alpha^{2}\cos^{2}\psi+\beta^{2}m^{2}\sin^{2}\psi)}\cdot\sqrt{(\gamma m^{2}\sin^{2}\psi+\gamma\cos^{2}\psi+\delta m^{2}\sin^{2}\psi)}}. \]
Let us take the indeterminate $m$ so that we have $\alpha^{2}m^{2}+\beta^{2}m^{2}=\alpha^{2}$, or $m^{2}=\frac{\alpha^{2}}{\alpha^{2}+\beta^{2}}$; the preceding integral will reduce to the form $\displaystyle\int\frac{mQ\,d\psi}{\sqrt{(A+B\sin^{2}\psi)}}$; and as by the same substitution $Q$ becomes an even rational function of $\sin\psi$, the formula $Z$ will again be brought to elliptic functions.

\subsection*{(138).}

We shall not multiply these examples of reduction further; they suffice to indicate the substitutions to be made in cases analogous to those we have developed. We shall conclude by gathering into one table some of the simplest formulas whose application presents itself frequently in the integrals which relate to elliptic functions.
\[ \left.\begin{aligned} \int\frac{d\varphi}{\Delta}&=F(c,\varphi)\\ \int\Delta\,d\varphi&=E(c,\varphi) \end{aligned}\right\}\ \text{Principal formulas.} \]
\[ \begin{aligned} \int\frac{d\varphi}{\Delta^{3}}&=\frac{1}{b^{2}}E(c,\varphi)-\frac{c^{2}\sin\varphi\cos\varphi}{b^{2}\Delta}\\ \int\frac{d\varphi\sin^{2}\varphi}{\Delta}&=\frac{1}{c^{2}}(F-E)\\ \int\frac{d\varphi\cos^{2}\varphi}{\Delta}&=\frac{1}{c^{2}}(E-b^{2}F)\\ \int\frac{d\varphi}{\Delta\cos^{2}\varphi}&=\frac{1}{b^{2}}(\Delta\operatorname{tang}\varphi+b^{2}F-E)\\ \int\frac{d\varphi\operatorname{tang}^{2}\varphi}{\Delta}&=\frac{1}{b^{2}}(\Delta\operatorname{tang}\varphi-E) \end{aligned} \]

\origpage{200}
\[ \begin{aligned} \int\frac{d\varphi\operatorname{tang}^{2}\frac{1}{2}\varphi}{\Delta}&=2\Delta\operatorname{tang}\tfrac{1}{2}\varphi+F-2E\\ \int\frac{d\varphi}{\Delta\cos^{2}\frac{1}{2}\varphi}&=2\Delta\operatorname{tang}\tfrac{1}{2}\varphi+2F-2E\\ \int\frac{\Delta\,d\varphi}{\cos^{2}\varphi}&=\Delta\operatorname{tang}\varphi+F-E\\ \int\Delta\,d\varphi\operatorname{tang}^{2}\varphi&=\Delta\operatorname{tang}\varphi+F-2E\\ \int\Delta^{3}\,d\varphi&=\frac{c^{2}}{3}\Delta\sin\varphi\cos\varphi+\frac{4-2c^{2}}{3}E-\frac{b^{2}}{3}F\\ \int\Delta\,d\varphi\sin^{2}\varphi&=-\tfrac{1}{3}\Delta\sin\varphi\cos\varphi+\frac{2c^{2}-1}{3c^{2}}E+\frac{b^{2}}{3c^{2}}F\\ \int\Delta\,d\varphi\cos^{2}\varphi&=\tfrac{1}{3}\Delta\sin\varphi\cos\varphi+\frac{1+c^{2}}{3c^{2}}E-\frac{b^{2}}{3c^{2}}F. \end{aligned} \]
All these integrals are taken counting from $\varphi=0$.

\section*{On the integral $V=\displaystyle\int\frac{(x^{2}+m)\,dx}{\sqrt{(\alpha+\beta x+\gamma x^{2}+\varepsilon x^{4})}}$.}

\subsection*{(139).}

This integral, in which the polynomial under the radical lacks a second term, is generally reducible to elliptic functions of the first and second species.

For if we decompose the polynomial $\alpha+\beta x+\gamma x^{2}+\varepsilon x^{4}$ into real factors of the second degree, in this manner:
\[ \alpha+\beta x+\gamma x^{2}+\varepsilon x^{4}=\varepsilon(x^{2}+\lambda x+\mu)(x^{2}-\lambda x+\nu), \]
and if then, to make the odd powers of the variable disappear from under the radical, we suppose $x=\frac{p+qy}{1+y}$, we shall have, to determine $p$ and $q$, the equations
\[ \begin{aligned} 2pq+\lambda(p+q)+2\mu&=0\\ 2pq-\lambda(p+q)+2\nu&=0, \end{aligned} \]
which will always give real values for $p$ and $q$, by taking suitably the values of $\lambda$, $\mu$, $\nu$ (art.~5).

This being so, if for brevity we put,
\[ \begin{aligned} M&=p^{2}-\lambda p+\nu=(p-\tfrac{1}{2}\lambda)(p-q)\\ N&=q^{2}-\lambda q+\nu=-(q-\tfrac{1}{2}\lambda)(p-q)\\ P&=p^{2}+\lambda p+\mu=(p+\tfrac{1}{2}\lambda)(p-q)\\ Q&=q^{2}+\lambda q+\mu=-(q+\tfrac{1}{2}\lambda)(p-q), \end{aligned} \]

\origpage{201}
we shall have
\[ \begin{aligned} x^{2}+\lambda x+\mu&=\frac{P+Qy^{2}}{(1+y)^{2}}\\ x^{2}-\lambda x+\nu&=\frac{M+Ny^{2}}{(1+y)^{2}}, \end{aligned} \]
and the transformed expression in $y$ will be
\[ V=(q-p)\int\frac{dy[(p+qy)^{2}+m(1+y)^{2}]}{(1+y)^{2}\sqrt{[\varepsilon(P+Qy^{2})(M+Ny^{2})]}}. \]
Let, for brevity, $Y=\sqrt{[\varepsilon(P+Qy^{2})(M+Ny^{2})]}$; we shall have by the development of the preceding quantity,
\[ V=(m+q^{2})(q-p)\int\frac{dy}{Y}-2q(q-p)^{2}\int\frac{dy}{(1+y)Y}+(q-p)^{3}\int\frac{dy}{(1+y)^{2}Y}. \]
But from the differential of $\frac{Y}{1+y}$ we find
\[ d\left(\frac{Y}{1+y}\right)=\frac{\varepsilon\,dy}{Y}\left[QN(y^{2}-1)+\frac{MQ+NP+2QN}{1+y}-\frac{(M+N)(P+Q)}{(1+y)^{2}}\right]. \]
Substituting in this quantity the values of $M$, $N$, $P$, $Q$, we shall have
\[ d\left(\frac{Y}{1+y}\right)=(q^{2}-\tfrac{1}{4}\lambda^{2})(q-p)^{2}\cdot\frac{(y^{2}-1)\varepsilon\,dy}{Y}+2q(q-p)^{3}\cdot\frac{\varepsilon\,dy}{(1+y)Y}-(q-p)^{4}\cdot\frac{\varepsilon\,dy}{(1+y)^{2}Y}. \]
From this we draw
\[ V=-\frac{1}{\varepsilon(q-p)}\cdot\frac{Y}{1+y}+(q^{2}-\tfrac{1}{4}\lambda^{2})(q-p)\int\frac{y^{2}\,dy}{Y}+(m+\tfrac{1}{4}\lambda^{2})(q-p)\int\frac{dy}{Y}. \]
Now the integrals $\displaystyle\int\frac{dy}{Y}$, $\displaystyle\int\frac{y^{2}\,dy}{Y}$ are brought by the known transformations to elliptic functions of the first and second species; hence the integral $V$ depends generally only on these two functions, and consequently can be determined by arcs of an ellipse.

\section*{On the integral $Z=\displaystyle\int\frac{dz}{(1+\mu z^{2})\sqrt[3]{(1+\nu z^{2})}}$.}

\subsection*{(140).}

This integral is readily brought to elliptic functions by the methods given in article 137; but the calculation deserves to be developed in some particular cases, on account of the

\origpage{202}
reductions which they present and which may cast a new light on the use of elliptic functions.

Suppose first $\nu$ positive, $\mu$ being at will positive or negative; if we put $\sqrt[3]{(1+\nu z^{2})}=y$, and $\alpha^{3}=\frac{\nu}{\mu}-1$, we shall have for first transformed expression
\[ Z=\frac{3\sqrt{\nu}}{2\mu}\int\frac{y\,dy}{(\alpha^{3}+y^{3})\sqrt{(y^{3}-1)}}. \]
This integral can be decomposed into two others; for this purpose let
\[ \begin{aligned} P&=\int\frac{y\,dy}{(\alpha^{2}-\alpha y+y^{2})\sqrt{(y^{3}-1)}}\\ Q&=\int\frac{y^{2}\,dy}{(\alpha^{3}+y^{3})\sqrt{(y^{3}-1)}}, \end{aligned} \]
and we shall have
\[ Z=\frac{3\sqrt{\nu}}{2\alpha\mu}(P-Q). \]
The integral $Q$ is found immediately; for if in place of $y$ we put its value as a function of $z$, we shall have
\[ Q=\frac{2\mu}{2\sqrt{\nu}}\int\frac{dz}{1+\mu z^{2}}=\frac{2\sqrt{\mu}}{3\sqrt{\nu}}\operatorname{arc}\operatorname{tang}(z\sqrt{\mu}). \]
\ednote{The first coefficient is printed $\frac{2\mu}{2\sqrt{\nu}}$; the substitution gives $\frac{2\mu}{3\sqrt{\nu}}$, which also agrees with the right-hand side, an apparent misprint.}
Thus the quantity $Q$ is given by an arc of a circle if $\mu$ is positive, and by a logarithm if $\mu$ is negative. Everything therefore reduces to finding the value of the other integral $P$.

For this, let $y=1+x^{2}$; we shall have
\[ P=\int\frac{2(1+x^{2})\,dx}{(x^{4}+(2-\alpha)x^{2}+\alpha^{2}-\alpha+1)\sqrt{(x^{4}+3x^{2}+3)}}. \]
This integral depends in general on elliptic functions whose parameter is imaginary, and we have given for this object the necessary formulas. But in the case of $\alpha=2$, the integral $P$ depends only on a function of the third species whose parameter is real, and it is this case whose calculation we are going to develop.

Having therefore made $\alpha=2$, or $\nu=9\mu$, it will be a matter of finding the integral
\[ P=\int\frac{2(1+x^{2})\,dx}{(x^{4}+3)\sqrt{(x^{4}+3x^{2}+3)}}. \]

\origpage{203}
For this, let $r=\sqrt[4]{3}$ and $x=r\operatorname{tang}\frac{1}{2}\varphi$; let further $c=\frac{1}{2}\sqrt{(2-\sqrt{3})}=\sin 15^{\circ}$; we shall have the transformed expression
\[ P=\int\frac{d\varphi}{\Delta}\left(\frac{1-r^{2}}{6r}+\frac{1}{6r}\cdot\frac{r^{2}+\cos\varphi}{1-\tfrac{1}{2}\sin^{2}\varphi}\right), \]
whence we draw
\[ P=\frac{1-r^{2}}{6r}F+\frac{r}{6}\Pi(-\tfrac{1}{2})+\frac{1}{6r}\int\frac{d\varphi\cos\varphi}{(1-\tfrac{1}{2}\sin^{2}\varphi)\Delta}. \]
To obtain this last integral, let $\sqrt{(1-c^{2}\sin^{2}\varphi)}=u\sin\varphi$; the differential $\frac{d\varphi\cos\varphi}{(1-\frac{1}{2}\sin^{2}\varphi)\Delta}$ will become $-\frac{du}{u^{2}-\frac{1}{4}r^{2}}$, and its integral will be $\frac{1}{r}\log\left(\frac{u+\frac{1}{2}r}{u-\frac{1}{2}r}\right)$; we have therefore
\[ P=\frac{1-r^{2}}{6r}F+\frac{r}{6}\Pi(-\tfrac{1}{2})+\frac{1}{6r^{2}}\log\left(\frac{\Delta+\frac{1}{2}r\sin\varphi}{\Delta-\frac{1}{2}r\sin\varphi}\right). \]
We add no constant, because this integral vanishes when $\varphi=0$, a value which corresponds to that of $z=0$.

The value of $P$ therefore depends on the elliptic function of the third species $\Pi(-\frac{1}{2})$, and we must examine whether, by reason of the parameter $-\frac{1}{2}$, this function might not be reducible indefinitely to a lower species.

Having the modulus $c=\frac{1}{2}\sqrt{(2-\sqrt{3})}$, and its complement $b=\frac{1}{2}\sqrt{(2+\sqrt{3})}$, we see immediately that the parameter $n=-\frac{1}{2}$ belongs to the second form of no.~50, and that by putting $n=-1+b^{2}\sin^{2}\theta$, we shall have $\sin^{2}\theta=\frac{n+1}{b^{2}}=4-2\sqrt{3}$, or $\sin\theta=\sqrt{3}-1$. Now this value of $\sin\theta$ is that which, according to article 24, satisfies the equation $F(b,\theta)=\frac{1}{3}F^{1}(b)$. Hence the reduction of the function $\Pi(-\frac{1}{2})$ is possible, and it will be carried out by the formula of no.~102, which gives
\[ \Pi(n)=\left(1-\frac{V\Delta(b,\theta)}{b^{2}\sin\theta\cos\theta}\right)F(c,\varphi)+\frac{W\Delta(b,\theta)\Delta(c,\varphi)}{b^{2}\sin\theta\cos\theta\sin\varphi\cos\varphi}. \]
It is therefore a matter of making the substitutions in this formula; now the known value of $\theta$ gives $\frac{\Delta(b,\theta)}{b^{2}\sin\theta\cos\theta}=\frac{2}{r}$; we have next by the formulas of no.~57, $V=\frac{1}{2}b^{2}\sin\theta\sin\theta_{2}(1+\sin\theta)$, $\cos\theta_{2}=1-\sin\theta$

\origpage{204}
$=2-\sqrt{3}$, $\sin\theta_{2}=r\sin\theta$; hence $V=\frac{1}{2r}$. We have in the second place by the formulas of the same article,
\[ W=\frac{1}{3\sqrt{\alpha'}}\operatorname{arc}\operatorname{tang}\frac{n'\sqrt{\alpha'}\sin\theta\sin\theta_{2}}{1+n'}+\frac{1}{3\sqrt{\alpha'}}\operatorname{arc}\operatorname{tang}\frac{n'\sqrt{\alpha'}\sin^{2}\theta\sin\theta_{2}}{1+n'-n'\cos^{2}\theta\cos\theta_{2}}. \]
Substituting the values of $\theta$ and $\theta_{2}$, and putting further $n'=b^{2}\operatorname{tang}^{2}\varphi$ and $\sqrt{\alpha'}=\frac{\Delta}{\sin\varphi\cos\varphi}$, we shall have
\[ W=\frac{\sin\varphi\cos\varphi}{3\Delta}\operatorname{arc}\operatorname{tang}\frac{r\operatorname{tang}\varphi}{2\Delta}+\frac{\sin\varphi\cos\varphi}{3\Delta}\operatorname{arc}\operatorname{tang}\frac{r\Delta\operatorname{tang}\varphi}{1+r^{2}\Delta^{2}}. \]
Hence the function $\Pi(-\frac{1}{2})$ is reduced to the first species by the formula
\[ \Pi(-\tfrac{1}{2})=\left(1-\frac{1}{r}\right)F+\frac{2}{3r}\operatorname{arc}\operatorname{tang}\frac{r\operatorname{tang}\varphi}{2\Delta}+\frac{2}{3r}\operatorname{arc}\operatorname{tang}\frac{r\Delta\operatorname{tang}\varphi}{1+r^{2}\Delta^{2}}. \]
\ednote{Printed $\left(1-\frac{1}{r}\right)F$; the formula of n° 102 with $\frac{V\Delta(b,\theta)}{b^{2}\sin\theta\cos\theta}=\frac{1}{2r}\cdot\frac{2}{r}$ gives $1-\frac{1}{r^{2}}$, which is also needed for $F$ to cancel in $P$ below, an apparent misprint.}
Substituting this value in that of the integral $P$, we shall have at last
\[ P=\frac{1}{6r}\log\left(\frac{\Delta+\frac{1}{2}r\sin\varphi}{\Delta-\frac{1}{2}r\sin\varphi}\right)+\tfrac{1}{9}\operatorname{arc}\operatorname{tang}\frac{r\operatorname{tang}\varphi}{2\Delta}+\tfrac{1}{9}\operatorname{arc}\operatorname{tang}\frac{r\Delta\operatorname{tang}\varphi}{1+r^{2}\Delta^{2}}. \]
\ednote{The logarithmic term is printed with $\frac{1}{6r}$; the expression of $P$ on p.~203 has $\frac{1}{6r^{2}}$, an apparent misprint.}
From which we see that the integral $P$ becomes finally independent of elliptic functions, and that its expression contains only arcs of a circle and logarithms. The same is consequently true of the proposed integral $Z$ in the case where we have $\alpha=2$, or $\frac{\nu}{\mu}=9$.

If we wish to have the value of $Z$ when $z=\infty$, we must put $\varphi=\pi$, which gives $P=\frac{2\pi}{9}$; moreover in the same case we have $Q=\frac{\pi}{9}$; hence $Z=\frac{\pi}{4\sqrt{\mu}}$.

\subsection*{(141).}

The numerous reductions which presented themselves in the preceding calculation must lead us to presume that it is possible to arrive at the same result without the aid of elliptic functions; this is what we shall now place beyond doubt in the following manner.

Let us remark first that in the formula $Z$, one can change the coefficient $\nu$ at will into another $\nu'$, provided it is of the same

\origpage{205}
sign, and that for this it suffices to put $z\sqrt{\nu}=z'\sqrt{\nu'}$. Let then $\nu=3$, which will give for the case we are considering $\mu=\frac{1}{3}$, and the proposed formula will be
\[ Z=\int\frac{3\,dz}{(3+z^{2})\sqrt[3]{(1+3z^{2})}}. \]
I put $z=\frac{1-x}{1+x}$, and I have the transformed expression
\[ \tfrac{1}{3}Z=-2^{\frac{5}{3}}\int\frac{(1-x^{2})\,dx}{(1-x^{3})\sqrt[3]{(1+x^{3})}}. \]
\ednote{The factor is printed $2^{\frac{5}{3}}$; the substitution gives $2^{-\frac{5}{3}}$, as in $\frac{1}{3}Z=2^{-\frac{5}{3}}(M-N)$ below, an apparent misprint.}
I remark next that by making
\[ \begin{aligned} M&=\int\frac{-dx}{(1-x^{3})\sqrt[3]{(1+x^{3})}}\\ N&=\int\frac{-x^{2}\,dx}{(1-x^{3})\sqrt[3]{(1+x^{3})}}, \end{aligned} \]
we shall have $\frac{1}{3}Z=2^{-\frac{5}{3}}(M-N)$. It is therefore only a matter of obtaining the integrals $M$ and $N$. For this,

Let, 1°. $\sqrt[3]{(1+x^{3})}=\frac{x\sqrt[3]{2}}{t}$; we shall have $M=2^{-\frac{1}{3}}\displaystyle\int\frac{-dt}{1-t^{3}}$.

Let, 2°. $\sqrt[3]{(1+x^{3})}=\frac{\sqrt[3]{2}}{u}$; we shall have $N=2^{-\frac{1}{3}}\displaystyle\int\frac{-du}{1-u^{3}}$.
These two integrals are therefore of the same form, and can be expressed by arcs of a circle and logarithms.

If they had to be evaluated separately, they would present a difficulty; for the proposed integral having to be taken from $z=0$, which gives $x=1$ and $t=1$, the integral $M$ would contain an infinite constant. The same would hold for the integral $N$; but as we need only the difference of these integrals, the infinite can be avoided entirely by a very simple means.

We see indeed from the nature of the question that we shall have the value of $Z$ by the single formula
\[ Z=\frac{3}{4}\int\frac{-dv}{1-v^{3}}, \]
provided that we take the integral from $v=t=\frac{x\sqrt[3]{2}}{\sqrt[3]{(1+x^{3})}}$ to $v=u=\frac{\sqrt[3]{2}}{\sqrt[3]{(1+x^{3})}}$. The integration being therefore carried out between these

\origpage{206}
limits both variable, we shall have in general
\[ \begin{aligned} Z={}&\tfrac{1}{8}\log\left(\frac{1+u+u^{2}}{1-2u+u^{2}}\right)+\frac{\sqrt{3}}{4}\operatorname{arc}\operatorname{tang}\frac{2u+1}{\sqrt{3}}\\ &-\tfrac{1}{8}\log\left(\frac{1+t+t^{2}}{1-2t+t^{2}}\right)-\frac{\sqrt{3}}{4}\operatorname{arc}\operatorname{tang}\frac{2t+1}{\sqrt{3}}. \end{aligned} \]
If we wish to have the integral when $z=\infty$, we must suppose $\sqrt[3]{(1+x^{3})}=\omega$, $\omega$ being infinitely small; we shall therefore put $t=-\frac{\sqrt[3]{2}}{\omega}$, and $u=\frac{\sqrt[3]{2}}{\omega}$, which will give $Z=\frac{\pi\sqrt{3}}{4}$; a result in agreement with that which we found by the other method, since in this case we have $\mu=\frac{1}{3}$.

\subsection*{(142).}

Let us now return to the general formula, and suppose that $\nu$ is negative, or, what comes to the same thing, let us consider directly the formula
\[ Z=\int\frac{dz}{(1-\mu z^{2})\sqrt[3]{(1-\nu z^{2})}}, \]
in which we shall suppose $\nu$ positive, $\mu$ being at will positive or negative. If we put $\sqrt[3]{(1-\nu z^{2})}=y$ and $\alpha^{3}=\frac{\nu}{\mu}-1$, we shall have the transformed integral
\[ Z=\frac{3\sqrt{\nu}}{2\mu}\int\frac{-y\,dy}{(\alpha^{3}+y^{3})\sqrt{(1-y^{3})}}, \]
which can be put in the form $Z=\frac{3\sqrt{\nu}}{2\alpha\mu}(P-Q)$, by making
\[ \begin{aligned} P&=\int\frac{-y\,dy}{(\alpha^{2}-\alpha y+y^{2})\sqrt{(1-y^{3})}}\\ Q&=\int\frac{-y^{2}\,dy}{(\alpha^{3}+y^{3})\sqrt{(1-y^{3})}}. \end{aligned} \]
We find immediately the value of the second integral
\[ Q=\frac{2\mu}{3\sqrt{\nu}}\int\frac{dz}{1-\mu z^{2}}=\frac{\sqrt{\mu}}{3\sqrt{\nu}}\log\left(\frac{1+z\sqrt{\mu}}{1-z\sqrt{\mu}}\right); \]
to have that of the first, let $y=1-x^{2}$, and we shall have the new

\origpage{207}
transformed integral
\[ P=\int\frac{2(1-x^{2})\,dx}{(x^{4}-(2-\alpha)x^{2}+\alpha^{2}-\alpha+1)\sqrt{(x^{4}-3x^{2}+3)}}. \]
This integral depends in general on elliptic functions whose parameter is imaginary, and it can be brought back to those whose parameter is real. But in the case of $\alpha=2$, the solution is greatly simplified and becomes entirely independent of elliptic functions; this will be easily verified by treating the formula
\[ P=\int\frac{2(1-x^{2})\,dx}{(x^{4}+3)\sqrt{(x^{4}-3x^{2}+3)}}, \]
like that of the preceding example.

\subsection*{(143).}

We can also, in the same case, arrive directly at the result of the integration without the aid of elliptic functions. And first, since the value of $\nu$ may be taken at will, let us suppose $\nu=3$, which will give $\mu=\frac{1}{3}$, and the proposed formula will become
\[ Z=\int\frac{3\,dz}{(3-z^{2})\sqrt[3]{(1-3z^{2})}}. \]
Let $\sqrt[3]{(1-3z^{2})}=1-\frac{x}{z}$, and we shall have the transformed integral
\[ Z=\frac{3}{2}\int\frac{(1-x^{2})\,dx}{(1+x^{2})(3x^{2}-1)}-\frac{\sqrt{3}}{2}\int\frac{3x^{4}+1}{3x^{4}+2x^{2}-1}\cdot\frac{dx}{\sqrt{(3x^{4}+6x^{2}-1)}}. \]
The first part is rational; the second seems bound to decompose into two elliptic functions of the third species, on account of the two factors of the denominator $3x^{4}+2x^{2}-1$. But if we examine the matter with more attention, we shall find that, on putting $\sqrt{(3x^{4}+6x^{2}-1)}=px$, the second part becomes $\displaystyle-\frac{\sqrt{3}}{2}\int\frac{dp}{p^{2}-4}$, so that the total value of $Z$ is
\[ Z=\frac{3}{2}\int\frac{(1-x^{2})\,dx}{(1+x^{2})(3x^{2}-1)}-\frac{\sqrt{3}}{2}\int\frac{dp}{p^{2}-4}. \]
Thus this integral is determined entirely by arcs of a circle and logarithms.

\origpage{208}

A similar simplification will take place in general for the formula
\[ \int\frac{(\gamma x^{4}-\alpha)\,dx}{(\gamma x^{4}+\varepsilon x^{2}+\alpha)\sqrt{(\alpha+\beta x^{2}+\gamma x^{4})}}; \]
for if we put the radical $\sqrt{(\alpha+\beta x^{2}+\gamma x^{4})}=px$, this formula will become $\displaystyle\int\frac{dp}{p^{2}-\beta+\varepsilon}$; a very remarkable reduction, and analogous to that which serves as the basis of the formula of No.~46.

The integral $\displaystyle\int\frac{dz}{(3\pm z^{2})\sqrt[3]{(1\pm 3z^{2})}}$ which we have just determined in its two different cases, is found mentioned, page 116, in the list of the works of Euler that Fuss has given; but the memoir which concerns it has not yet been published. We have wished to show by so remarkable an example that the theory of elliptic functions leads in a sure manner to the simplest expressions of the integrals related to it.

\subsection*{(144).}

Independently of the two cases of which we have just spoken, there is a third in which the integral $Z$ depends only on arcs of a circle and logarithms. It is that in which we should have
\[ Z=\int\frac{3\,dz}{(3-z^{2})\sqrt[3]{(1+z^{2})}}, \]
an integral which is related to the general formula of No.~139, by putting $\nu=1$, $\mu=-\frac{1}{3}$ and $\alpha^{3}=-4$. One could not then avoid meeting, in the integral $P$, elliptic functions whose parameter is imaginary; but the reduction of these integrals will lead to a result entirely independent of elliptic functions. We shall not enter into the detail of these reductions, and shall confine ourselves to proving by another route that the integral $Z$ can be determined by arcs of a circle and logarithms.

For this purpose, let $1+z^{2}=4y^{3}$, and we shall have the transformed integral
\[ Z=\frac{9}{2\sqrt[3]{4}}\int\frac{y\,dy}{(1-y^{3})\sqrt{(4y^{3}-1)}}: \]
But if we consider the integral $T=\displaystyle\int\frac{t\,dt}{(t^{3}-1)^{\frac{2}{3}}}$, and put $\sqrt[3]{(t^{3}-1)}=\frac{t}{y}-1$, we shall have $t=\frac{2y\sqrt{3}}{\sqrt{3}-\sqrt{(4y^{3}-1)}}$, $\frac{t^{2}\,dt}{(t^{3}-1)^{\frac{2}{3}}}=\frac{dt}{y}-\frac{t\,dy}{yy}$;

\origpage{209}
therefore
\[ T=\frac{3}{2}\int\frac{y\,dy}{1-y^{3}}+\frac{3\sqrt{3}}{2}\int\frac{y\,dy}{(1-y^{3})\sqrt{(4y^{3}-1)}}. \]
Thus we shall know the integral $Z$ if $T$ is known; now it suffices to put $t^{3}=\frac{1}{1-u^{3}}$, and we shall have $T=\displaystyle\int\frac{du}{1-u^{3}}$; therefore $Z$ will be determined by the equation
\[ \frac{\sqrt[3]{4}}{\sqrt{3}}Z=\int\frac{du}{1-u^{3}}-\frac{3}{2}\int\frac{y\,dy}{1-y^{3}}, \]
and consequently $Z$ can be expressed by arcs of a circle and logarithms.

\section*{On the principal cases in which one can evaluate, by elliptic functions, the integral $\left(\frac{p}{q}\right)=\displaystyle\int\frac{x^{p-1}\,dx}{\sqrt[n]{(1-x^{n})^{n-q}}}$, taken from $x=0$ to $x=1$.}

\subsection*{(145).}

These integrals, which we shall treat at great length in the second part, can be determined in several cases by complete elliptic functions of the first species. These cases are those in which the exponent $n$ is equal to one of the numbers 3, 4, 6, 8 and 12. We shall examine them successively.

\section*{First case, $n=3$.}

The only transcendent to be determined is
\[ \left(\frac{1}{1}\right)=2^{\frac{1}{3}}\int\frac{dx}{\sqrt{(1-x^{3})}}, \]
and its value, already found (art.\ 39), is
\[ \left(\frac{1}{1}\right)=2^{\frac{4}{3}}3^{-\frac{3}{4}}F^{1}(\sin 15^{\circ}). \]

\section*{Second case, $n=4$.}

\subsection*{(146).}

It suffices again in this case to determine the transcendent
\[ \left(\frac{1}{1}\right)=2^{\frac{1}{2}}\int\frac{dx}{\sqrt{(1-x^{4})}}=\int\frac{dz}{\sqrt{(1+z^{4})}}, \]

\origpage{210}
the integral with respect to $z$ having to be taken from $z=0$ to $z=\infty$; now one or the other of the two forms leads to the same result
\[ \left(\frac{1}{1}\right)=F^{1}(\sin 45^{\circ}). \]

\section*{Third case, $n=6$.}

\subsection*{(147).}

All the integrals represented by $\left(\frac{p}{q}\right)$ can be determined in this case by the single transcendent
\[ \left(\frac{1}{1}\right)=2^{\frac{2}{3}}\int\frac{dx}{\sqrt{(1-x^{6})}}=2^{\frac{2}{3}}\cos\frac{\pi}{6}\int\frac{dz}{\sqrt{(1+z^{6})}}. \]
The first form will give, on putting $\frac{1}{x^{2}}=1+z^{2}$, the result $\left(\frac{1}{1}\right)=2^{\frac{2}{3}}3^{-\frac{1}{4}}F^{1}(\sin 15^{\circ})$. The second would give, on putting $\frac{1}{z^{2}}=y^{2}-1$, the result $\left(\frac{1}{1}\right)=2^{\frac{2}{3}}3^{-\frac{3}{4}}F^{1}(\sin 75^{\circ})$, and from these two values we deduce $F^{1}(\sin 75^{\circ})=\sqrt{3}E^{1}(\sin 15^{\circ})$\ednote{Printed $E^{1}(\sin 15^{\circ})$; the two values just found give $F^{1}(\sin 75^{\circ})=\sqrt{3}F^{1}(\sin 15^{\circ})$, an apparent misprint for $F^{1}$.}, which is a known property of these two functions.

\section*{Fourth case, $n=8$.}

\subsection*{(148).}

In this case we must evaluate the two transcendents
\[ \begin{aligned} \left(\frac{1}{1}\right)&=2^{\frac{3}{4}}\int\frac{dx}{\sqrt{(1-x^{8})}}=2^{\frac{3}{4}}\cos\frac{\pi}{8}\int\frac{dz}{\sqrt{(1+z^{8})}}\\ \left(\frac{2}{2}\right)&=2^{\frac{1}{2}}\int\frac{x\,dx}{\sqrt{(1-x^{8})}}=2^{\frac{1}{2}}\cos\frac{\pi}{4}\int\frac{z\,dz}{\sqrt{(1+z^{8})}}. \end{aligned} \]
The second, on putting $x^{2}$ in place of $x$, becomes $2^{-\frac{1}{2}}\displaystyle\int\frac{dx}{\sqrt{(1-x^{4})}}$, and its value is $\left(\frac{2}{2}\right)=\frac{1}{2}F^{1}(\sin 45^{\circ})$.

To evaluate the integral $\left(\frac{1}{1}\right)$, we shall choose the second form, in which the differential $\frac{dz}{\sqrt{(1+z^{8})}}$ is to be integrated between the limits $z=0$, $z=\infty$. Let us suppose this integral found from $z=0$ to $z=1$, and call $Q$ the part comprised from $z=1$ to $z=\infty$;

\origpage{211}
if we put $z=\frac{1}{u}$, we shall have the transformed integral $Q=\displaystyle\int\frac{-u^{2}\,du}{\sqrt{(1+u^{8})}}$, which must be taken from $u=1$ to $u=0$; on changing its sign, the integral must be taken from $u=0$ to $u=1$; and as one may then put $z$ in place of $u$, it is clear that the integral sought $\displaystyle\int\frac{dz}{\sqrt{(1+z^{8})}}$, taken from $z=0$ to $z=\infty$, will be equal to the integral $T=\displaystyle\int\frac{dz(1+z^{2})}{\sqrt{(1+z^{8})}}$, taken from $z=0$ to $z=1$.

Let $1+z^{4}=pz^{2}$, we shall have $\sqrt{(1+z^{8})}=z^{2}\sqrt{(p^{2}-2)}$, and consequently $T=\displaystyle\int\frac{dz(1+z^{2})}{z^{2}\sqrt{(p^{2}-2)}}$; but the equation $1+z^{4}=pz^{2}$ gives $1-z^{2}=z\sqrt{(p-2)}$, $\frac{1}{z}-z=\sqrt{(p-2)}$, $\frac{dz}{zz}+dz=-\frac{\frac{1}{2}dp}{\sqrt{(p-2)}}$: therefore we shall have
\[ T=\int\frac{\tfrac{1}{2}dp}{\sqrt{(p-2)}\cdot\sqrt{(p^{2}-2)}}, \]
this integral being taken from $p=2$ to $p=\infty$.

Let $p=2+q^{2}$, we shall have the new transformed integral
\[ T=\int\frac{dq}{\sqrt{[(q^{2}+2+\sqrt{2})(q^{2}+2-\sqrt{2})]}}, \]
which must be taken from $q=0$ to $q=\infty$. Let $m^{2}=2+\sqrt{2}$, $q=m\cot\varphi$ and $c^{2}=2\sqrt{2}-2$, we shall have the indefinite integral $T=\displaystyle\frac{1}{m}\int\frac{d\varphi}{\sqrt{(1-c^{2}\sin^{2}\varphi)}}=\frac{1}{m}F(c,\varphi)$, and the complete integral $T^{1}=\frac{1}{m}F^{1}(c)$; therefore, finally, the value of the transcendent sought $\left(\frac{1}{1}\right)$ will be
\[ \left(\frac{1}{1}\right)=2^{-\frac{1}{4}}F^{1}(c). \]
We have seen besides (No.~64) that for this value of the modulus $c$, to which corresponds the complementary modulus $b=\sqrt{2}-1=\operatorname{tang}\frac{\pi}{8}$, we have $F^{1}(c)=\sqrt{2}F^{1}(b)$. Thus the transcendent $\left(\frac{1}{1}\right)$ is expressed equally by $2^{-\frac{1}{4}}F^{1}(c)$, or by $2^{\frac{1}{4}}F^{1}(b)=2^{\frac{1}{4}}F^{1}\left(\operatorname{tang}\frac{\pi}{8}\right)$.

\origpage{212}

\section*{Fifth case, $n=12$.}

\subsection*{(149).}

All the integrals comprised in this case will be determined if we know the three transcendents
\[ \begin{aligned} \left(\frac{1}{1}\right)&=2^{\frac{5}{6}}\int\frac{dx}{\sqrt{(1-x^{12})}}=2^{\frac{5}{6}}\cos\frac{\pi}{12}\int\frac{dz}{\sqrt{(1+z^{12})}}\\ \left(\frac{2}{2}\right)&=2^{\frac{4}{6}}\int\frac{x\,dx}{\sqrt{(1-x^{12})}}=2^{\frac{4}{6}}\cos\frac{2\pi}{12}\int\frac{z\,dz}{\sqrt{(1+z^{12})}}\\ \left(\frac{3}{3}\right)&=2^{\frac{3}{6}}\int\frac{x^{2}\,dx}{\sqrt{(1-x^{12})}}=2^{\frac{3}{6}}\cos\frac{3\pi}{12}\int\frac{z^{2}\,dz}{\sqrt{(1+z^{12})}}, \end{aligned} \]
the integrals in $x$ being taken from $x=0$ to $x=1$, and the integrals in $z$ being taken from $z=0$ to $z=\infty$.

The second and the third are brought back immediately to the formulas already integrated. Indeed, if in place of $x^{2}$ we put $x$, the second formula will give
\[ \left(\frac{2}{2}\right)=2^{-\frac{1}{3}}\int\frac{dx}{\sqrt{(1-x^{6})}}=2^{-\frac{1}{3}}3^{-\frac{1}{4}}F^{1}(\sin 15^{\circ}); \]
and if in place of $x^{3}$ we put $x$, the third will give
\[ \left(\frac{3}{3}\right)=\frac{\sqrt{2}}{3}\int\frac{dx}{\sqrt{(1-x^{4})}}=\tfrac{1}{3}F^{1}(\sin 45^{\circ}). \]

Let us come now to the first integral, which appears to present rather great difficulties.

Choosing the first form, we have to find the integral $X=\displaystyle\int\frac{dx}{\sqrt{(1-x^{12})}}$ from $x=0$ to $x=1$. For this I put $x^{2}-\frac{1}{x^{2}}=q$\ednote{Printed thus; the lines that follow, and the equation $\frac{1}{x^{2}}-x^{2}=q$ further on, require $\frac{1}{x^{2}}-x^{2}=q$, an apparent misprint.}, whence I derive successively $\frac{1}{x^{4}}+x^{4}=q^{2}+2$. $\frac{1}{x^{6}}-x^{6}=q^{3}+3q$, $1-x^{12}=x^{6}(q^{3}+3q)$, $\frac{dx}{\sqrt{(1-x^{12})}}=\frac{dx}{x^{3}\sqrt{(q^{3}+3q)}}$; it remains to find the value of $\frac{dx}{x^{3}}$; now from the equation $\frac{1}{x^{2}}-x^{2}=q$, we deduce $\frac{1}{x^{2}}=\frac{1}{2}q+\sqrt{(\frac{1}{4}q^{2}+1)}$; therefore $-\frac{2\,dx}{x^{3}}=\frac{1}{2}dq+\frac{1}{2}\cdot\frac{q\,dq}{\sqrt{(q^{2}+4)}}$, therefore finally the transformed integral will be
\[ X=\tfrac{1}{4}\int\frac{dq}{\sqrt{(q^{3}+3q)}}+\tfrac{1}{4}\int\frac{dq\sqrt{q}}{\sqrt{(q^{2}+3)}\cdot\sqrt{(q^{2}+4)}}, \]

\origpage{213}
these two new integrals having to be taken from $q=0$ to $q=\infty$. If in the first we put $q=p^{2}$, we shall have $\displaystyle\int\frac{dq}{\sqrt{(q^{3}+3q)}}=\int\frac{2\,dp}{\sqrt{(p^{4}+3)}}$; this integral, which must likewise be taken from $p=0$ to $p=\infty$, will be found by the formula of No.~38, and we shall have as result $\displaystyle\int\frac{dq}{\sqrt{(q^{3}+3q)}}=\frac{2}{\sqrt[4]{3}}F^{1}(\sin 45^{\circ})$.

\subsection*{(150).}

There remains therefore only to find the integral
\[ Q=\int\frac{dq\sqrt{q}}{\sqrt{[(q^{2}+3)(q^{2}+4)]}}=\int\frac{dq\sqrt{q}}{\sqrt{(q^{4}+7q^{2}+12)}}; \]
this formula falls under case VII of art.\ 137; we shall therefore put $m^{4}=12$ and $q^{2}+m^{2}=qz$, which will give first $q^{4}+7q^{2}+12=q^{2}(z^{2}-2m^{2}+7)$, and $Q=\displaystyle\int\frac{q^{-\frac{1}{2}}\,dq}{\sqrt{(z^{2}-2m^{2}+7)}}$. But from the equation $q^{2}+m^{2}=qz$, we deduce successively
\[ \begin{aligned} q+m&=q^{\frac{1}{2}}\sqrt{(z+2m)}\\ q-m&=\pm q^{\frac{1}{2}}\sqrt{(z-2m)}\\ 2q^{\frac{1}{2}}&=\sqrt{(z+2m)}\pm\sqrt{(z-2m)}\\ q^{-\frac{1}{2}}\,dq&=\frac{\tfrac{1}{2}dz}{\sqrt{(z+2m)}}\pm\frac{\tfrac{1}{2}dz}{\sqrt{(z-2m)}}. \end{aligned} \]
Therefore we shall have the transformed integral
\[ Q=\int\frac{\tfrac{1}{2}dz}{\sqrt{(z+2m)}\cdot\sqrt{(z^{2}-2m^{2}+7)}}\pm\int\frac{\tfrac{1}{2}dz}{\sqrt{(z-2m)}\cdot\sqrt{(z^{2}-2m^{2}+7)}}; \]
the double sign $\pm$ relates to the two values of $q$ which correspond to one and the same value of $z$. The quantity $z$ has for its \emph{maximum}, $z=2m$, which takes place when $q=m$; and the correspondence of these two variables is such that, from $q=0$ to $q=m$, the value of $z$ decreases from $z=\infty$ to $z=2m$; and from $q=m$ to $q=\infty$, the value of $z$ increases from $z=2m$ to $z=\infty$.

From this we see that from $q=0$ to $q=m$, the integral $Q$ will be represented by the difference of the integrals
\[ \int\frac{\tfrac{1}{2}dz}{\sqrt{(z-2m)}\sqrt{(z^{2}-2m^{2}+7)}}-\int\frac{\tfrac{1}{2}dz}{\sqrt{(z+2m)}\sqrt{(z^{2}-2m^{2}+7)}}; \]

\origpage{214}
and the sum of these integrals will represent the integral $Q$ taken from $q=m$ to $q=\infty$. Therefore, on adding these two parts, the total integral $Q$ will be expressed by the single integral
\[ \int\frac{dz}{\sqrt{(z-2m)}\sqrt{(z^{2}-2m^{2}+7)}}. \]
Let $z=2m+p^{2}$, and we shall have finally the transformed integral
\[ Q=\int\frac{2\,dp}{\sqrt{(p^{4}+4mp^{2}+2m^{2}+7)}}, \]
which must be integrated from $p=0$ to $p=\infty$. Comparing this integral with the formula of No.~38, and putting $c=\frac{\sqrt{2}-\sqrt[4]{3}}{1+\sqrt{3}}$, we shall have as result $Q=4\sin 15^{\circ}\cdot F^{1}(c)$; therefore
\[ X=\frac{1}{2\sqrt[4]{3}}F^{1}(\sin 45^{\circ})+\sin 15^{\circ}\cdot F^{1}(c); \]
and finally the transcendent sought
\[ \left(\frac{1}{1}\right)=2^{\frac{5}{6}}X=2^{-\frac{1}{6}}3^{-\frac{1}{4}}F^{1}\sin(45^{\circ})+2^{\frac{5}{6}}\sin 15^{\circ}F^{1}(c). \]
Thus the transcendents necessary to make known all the values of the integral $\left(\frac{p}{q}\right)$ in the case of $n=12$, are determined by the three complete functions of the first species $F^{1}(\sin 45^{\circ})$, $F^{1}(\sin 15^{\circ})$, $F^{1}(c)$.

\subsection*{(151).}

Let us now seek another value of the transcendent $\left(\frac{1}{1}\right)$, by means of the integral $Z=\displaystyle\int\frac{dz}{\sqrt{(1+z^{12})}}$, taken from $z=0$ to $z=\infty$. I observe that if we take this integral first from $z=0$ to $z=1$, and then, to have the other part from $z=1$ to $z=\infty$, we put $\frac{1}{z}$ in place of $z$, we shall have, by the union of the two parts,
\[ Z=\int\frac{dz(1+z^{4})}{\sqrt{(1+z^{12})}}, \]
an integral which must be taken from $z=0$ to $z=1$.

\origpage{215}

Let $1+z^{4}=pz^{2}$, we shall have $1+z^{12}=z^{6}(p^{3}-3p)$, therefore $\frac{dz(1+z^{4})}{\sqrt{(1+z^{12})}}=\frac{dz}{z}\cdot\frac{p}{\sqrt{(p^{3}-3p)}}$; but from the equation $1+z^{4}=pz^{2}$, we get $\frac{dz}{z}=\frac{1}{2}\cdot\frac{dp}{\sqrt{(p^{2}-4)}}$; therefore we have the transformed integral
\[ Z=\int\frac{\tfrac{1}{2}dp\sqrt{p}}{\sqrt{(p^{4}-7p^{2}+12)}}, \]
which must be integrated from $p=2$ to $p=\infty$.

This formula is similar to that which we treated in the preceding article; we shall therefore put likewise $m^{4}=12$, $p^{2}+m^{2}=p\nu$, which will give
\[ \begin{aligned} p+m&=p^{\frac{1}{2}}\sqrt{(\nu+2m)}\\ p+m&=p^{\frac{1}{2}}\sqrt{(\nu-2m)}, \end{aligned} \]
\ednote{The second line is printed with $p+m$; from $p^{2}+m^{2}=p\nu$ it should read $p-m$, an apparent misprint.}
and it must be remarked that there is here no ambiguity, because we always have $p>m$. We shall have next
\[ \begin{aligned} p^{4}-7p^{2}+12&=p^{2}(\nu^{2}-2m^{2}-7)\\ 2p^{\frac{1}{2}}&=\sqrt{(\nu+2m)}+\sqrt{(\nu-2m)}\\ p^{-\frac{1}{2}}dp&=\frac{\tfrac{1}{2}d\nu}{\sqrt{(\nu+2m)}}+\frac{\tfrac{1}{2}d\nu}{\sqrt{(\nu-2m)}}; \end{aligned} \]
therefore we shall have the transformed integral
\[ Z=\tfrac{1}{4}\int\frac{d\nu}{\sqrt{(\nu+2m)}\sqrt{(\nu^{2}-2m^{2}-7)}}+\tfrac{1}{4}\int\frac{d\nu}{\sqrt{(\nu-2m)}\sqrt{(\nu^{2}-2m^{2}+7)}}, \]
\ednote{The second radical is printed with $+7$; from $p^{4}-7p^{2}+12=p^{2}(\nu^{2}-2m^{2}-7)$ it should read $\nu^{2}-2m^{2}-7$, an apparent misprint.}
of which the two parts must be integrated from $\nu=\sqrt{(2m^{2}+7)}$ to $\nu=\infty$. Let $\sqrt{(2m^{2}+7)}=2+\sqrt{3}=\mu$, and $\nu+2m=x^{2}$, the first part will become
\[ \int\frac{\tfrac{1}{2}dx}{\sqrt{[(x^{2}-2m+\mu)(x^{2}-2m-\mu)]}}. \]
Let again $x^{2}=\frac{2m+\mu}{\cos^{2}\varphi}$, and $c^{2}=\frac{\mu-2m}{2\mu}$, or $c=\frac{\sqrt{2}-\sqrt[4]{3}}{1+\sqrt{3}}$, and the transformed integral will be $\frac{1}{2\sqrt{2\mu}}\displaystyle\int\frac{d\varphi}{\sqrt{(1-c^{2}\sin^{2}\varphi)}}$, whose value within the required limits $=\frac{1}{2\sqrt{2\mu}}F^{1}(c)$.

\origpage{216}

We shall find similarly that, on putting $b=\frac{\sqrt{2}+\sqrt[4]{3}}{1+\sqrt{3}}$, the second integral has for its value $\frac{1}{2\sqrt{2\mu}}F^{1}(b)$; and the modulus of this one has been designated by $b$, because in fact, according to the values found, we have $b^{2}+c^{2}=1$, so that these moduli are complements of one another. This being so, on adding the two parts, we shall have the integral sought
\[ Z=\frac{1}{2\sqrt{2\mu}}[F^{1}(b)+F^{1}(c)]=\frac{\sqrt{3}-1}{4}[F^{1}(b)+F^{1}(c)], \]
which will give this second value of the transcendent $\left(\frac{1}{1}\right)$,
\[ \left(\frac{1}{1}\right)=2^{-\frac{5}{3}}[F^{1}(b)+F^{1}(c)]. \]
Equating the two values found, we have a first relation between the three functions $F^{1}(\sin 45^{\circ})$, $F^{1}(b)$, $F^{1}(c)$. But we can obtain a second, which will serve to establish the ratios of the functions $F^{1}(b)$, $F^{1}(c)$, both between themselves and with the function $F^{1}(\sin 45^{\circ})$.

\subsection*{(152).}

To arrive at this result, we must seek directly the integral
\[ M=\int\frac{x^{3}dx}{\sqrt[4]{(1-x^{12})}}, \]
which, being taken from $x=0$ to $x=1$, will be the value of the transcendent $\left(\frac{4}{9}\right)$. Let first $x^{4}=1-y^{4}$, we shall have the transformed integral
\[ M=\int\frac{y^{2}dy}{\sqrt[4]{(y^{8}-3y^{4}+3)}}, \]
which must be integrated from $y=0$ to $y=1$. If we then put $\nu=\sqrt[4]{3}$, $y^{4}+\nu^{2}=y^{2}z$, which will give $y^{2}=\frac{1}{2}z-\sqrt{(\frac{1}{4}z^{2}-\nu^{2})}$, $y^{8}-3y^{4}+3=y^{4}(z^{2}-2\nu^{2}-3)$, the transformed integral in $z$ will be
\[ M=-\tfrac{1}{4}\int\frac{dz}{\sqrt[4]{(z^{2}-2\nu^{2}-3)}}+\tfrac{1}{4}\int\frac{z\,dz}{\sqrt{(z^{2}-4\nu^{2})}\cdot\sqrt[4]{(z^{2}-2\nu^{2}-3)}}. \]
Let $z^{2}-2\nu^{2}-3=u^{4}$, and we shall have for final transformed integral
\[ M=\int\frac{\tfrac{1}{2}u^{2}du}{\sqrt{(u^{4}-2\nu^{2}+3)}}-\int\frac{\tfrac{1}{2}u^{2}du}{\sqrt{(u^{4}+2\nu^{2}+3)}}, \]

\origpage{217}
from which we see that the integral sought depends only on elliptic functions.

By virtue of the limits fixed, the last two integrals must be taken from $u=1$ to $u=\infty$. But as each of these integrals would become infinite at the last limit, in order to avoid this inconvenience we must put
\[ \begin{aligned} P&=\int\left(\frac{u^{2}du}{\sqrt{(u^{4}-2\nu^{2}+3)}}-du\right)\\ Q&=\int\left(du-\frac{u^{2}du}{\sqrt{(u^{4}+2\nu^{2}+3)}}\right). \end{aligned} \]
By this means, the two integrals $P$ and $Q$ will always be finite quantities, and after having found their complete values $P^{1}$, $Q^{1}$, we shall deduce from them $M$ or $\left(\frac{4}{9}\right)=\frac{1}{2}(P^{1}+Q^{1})$.

\subsection*{(153).}

To have the integral $P$, let $\beta^{2}=2\sqrt{3}-3$, and $u^{2}-\sqrt{(u^{4}-\beta^{2})}=p$, or $u^{2}=\frac{1}{2}p+\frac{\beta^{2}}{2p}$, we shall have the transformed integral $P=\displaystyle\int\frac{dp\sqrt{(\frac{1}{2}p)}}{\sqrt{(p^{2}+\beta^{2})}}$. Let then $p=\beta\operatorname{tang}^{2}\frac{1}{2}\varphi$ and $c^{2}=\frac{1}{2}$, we shall have
\[ P=\sqrt{(\tfrac{1}{2}\beta)}\cdot\int\frac{d\varphi\operatorname{tang}^{2}\tfrac{1}{2}\varphi}{\sqrt{(1-c^{2}\sin^{2}\varphi)}}; \]
this integral is found by the formulas of art.\ 138, and we have
\[ P=\sqrt{(\tfrac{1}{2}\beta)}\cdot[2\Delta\operatorname{tang}\tfrac{1}{2}\varphi+F(\varphi)-2E(\varphi)]. \]
We must now take this integral from $\varphi=0$, which corresponds to $u=\infty$, to the value of $\varphi$ which corresponds to $u=1$; at this last limit we have $\operatorname{tang}\frac{1}{2}\varphi=\frac{1-\sqrt{(1-\beta^{2})}}{\beta}$, or $\cos\varphi=\sqrt{\left(\frac{1-\beta}{1+\beta}\right)}$. Now I observe that the angle $\varphi$ thus determined is that which, for the modulus $c=\sqrt{\frac{1}{2}}$, gives $F(\varphi)=\frac{2}{3}F^{1}$; for we found (art.\ 24) that the equation $F(\psi)=\frac{1}{3}F^{1}$ is satisfied by the values $\cos\psi=\sqrt{\beta}$, $\sin\psi=\sqrt{(1-\beta)}$. If we then suppose $F(\varphi)+F(\psi)=F^{1}$, or $\operatorname{tang}\varphi\operatorname{tang}\psi=\frac{1}{b}=\sqrt{2}$, we shall deduce from it $\operatorname{tang}\varphi=\sqrt{\left(\frac{2\beta}{1-\beta}\right)}$ or $\cos\varphi=\sqrt{\left(\frac{1-\beta}{1+\beta}\right)}$, which is the value of $\varphi$ in question. Besides, on putting $\alpha^{2}=2\sqrt{3}+3$, we shall have more simply $\cos\varphi=\frac{\alpha-1}{\alpha+1}$, $\sin\varphi=\frac{2\sqrt{\alpha}}{\alpha+1}$, $\Delta(\varphi)=\frac{1}{2}(\alpha-1)$, $\operatorname{tang}\frac{1}{2}\varphi=\frac{1}{\sqrt{\alpha}}$, $\sin\psi=\frac{2}{\alpha+1}$. But

\origpage{218}
since the value of $\varphi$ is that which gives $F(\varphi)=\frac{2}{3}F^{1}$, we shall have at the same time, by the formulas of No.~32, $3E(\varphi)=2E^{1}+c^{2}\sin\varphi\sin\psi\times(2-\sin\psi)$, which gives
\[ E(\varphi)=\tfrac{2}{3}E^{1}+\frac{4}{3}\cdot\frac{\alpha\sqrt{\alpha}}{(\alpha+1)^{3}}. \]
Substituting all these values in that of $P$, there will result
\[ P^{1}=\frac{\sqrt{2\beta}}{3}(F^{1}-2E^{1})+\sqrt{\left(\frac{\beta}{2\alpha}\right)}\cdot\left(\alpha-1-\frac{8}{3}\cdot\frac{\alpha^{2}}{(\alpha+1)^{3}}\right). \]

\subsection*{(154).}

Let us come now to the integral $Q$. If we put $\sqrt{(u^{4}+\alpha^{2})}=u^{2}+q$, or $u^{2}=\frac{\alpha^{2}-q^{2}}{2q}$, we shall have first $Q=\displaystyle\int\frac{dq\sqrt{(\frac{1}{2}q)}}{\sqrt{(\alpha^{2}-q^{2})}}$. Let then $q=\alpha\cos^{2}\omega$, and we shall have $Q=\sqrt{\alpha}\displaystyle\int\frac{d\omega\cos^{2}\omega}{\sqrt{(1-\frac{1}{2}\sin^{2}\omega)}}$, whence we derive
\[ Q=\sqrt{\alpha}\cdot[2E(\omega)-F(\omega)]+\text{const}, \]
the modulus of these functions being always $c=\sqrt{\frac{1}{2}}$.

This integral must be taken between the two values of $\omega$ which give $u=1$ and $u=\infty$; at the second limit we have $\omega=\frac{1}{2}\pi$, and at the first $\cos^{2}\omega=\frac{\sqrt{(1+\alpha^{2})}-1}{\alpha}=\frac{\sqrt{3}}{\alpha}$, or $\cos^{2}\omega=\beta$; for from the two equations $\alpha^{2}=2\sqrt{3}+3$, $\beta^{2}=2\sqrt{3}-3$, we deduce $\alpha^{2}\beta^{2}=3$, or $\alpha\beta=\sqrt{3}$. But since we have $\cos\omega=\sqrt{\beta}$, the angle $\omega$ is the same as that which was designated in the preceding article by $\psi$, and we have consequently for the complete value of $Q$,
\[ Q^{1}=\sqrt{\alpha}\cdot[2E^{1}-F^{1}-2E(\psi)+F(\psi)]. \]
Now we have $F(\psi)=\frac{1}{3}F^{1}$, $E(\psi)=\frac{1}{3}E^{1}+\frac{1}{6}\sin\psi\sin\varphi(1+\sin\psi)$; therefore
\[ Q^{1}=\frac{2\sqrt{\alpha}}{3}(2E^{1}-F^{1})-\frac{4}{3}\cdot\frac{\alpha(\alpha+3)}{(\alpha+1)^{3}}. \]
If we now add the values of $P^{1}$ and $Q^{1}$, and have regard to the reductions furnished by the equations $\sqrt{(\frac{1}{2}\beta)}=\frac{2\sqrt{\alpha}}{\alpha^{2}-1}$, $\alpha^{4}=6\alpha^{2}+3$, we shall find that the algebraic parts destroy one another, and that we have
\[ P^{1}+Q^{1}=\tfrac{2}{3}(2E^{1}-F^{1})(\sqrt{\alpha}-\sqrt{\tfrac{1}{2}\beta}). \]

\origpage{219}
But the factor $\sqrt{\alpha}-\sqrt{\frac{1}{2}\beta}$ reduces to $\frac{3}{\nu\sqrt{2\alpha}}$; therefore finally the value of the transcendent sought $\left(\frac{4}{9}\right)=\frac{1}{2}(P^{1}+Q^{1})$ will be
\[ \left(\frac{4}{9}\right)=\frac{1}{\nu\sqrt{(2\alpha)}}(2E^{1}-F^{1}), \]
where we have $\nu=\sqrt[4]{3}$, and $\alpha=\nu\sqrt{(2+\sqrt{3})}=2\nu\cos 15^{\circ}$.

This result can still be simplified; for, according to the known property of the functions $E^{1}$, $F^{1}$, whose modulus is $\sin 45^{\circ}$, we have $\frac{\pi}{2}=F^{1}(2E^{1}-F^{1})$; therefore the transcendent $\left(\frac{4}{9}\right)$ is expressed by the single function of the first species $F^{1}$, by means of this very simple formula
\[ \left(\frac{4}{9}\right)=\frac{1}{2\nu\sqrt{2\alpha}}\cdot\frac{\pi}{F^{1}(\sin 45^{\circ})}. \]

\subsection*{(155).}

Let, for brevity, $F^{1}(\sin 45^{\circ})=B$, $F^{1}(\sin 15^{\circ})=C$; we have already found $\left(\frac{2}{2}\right)=\frac{2^{-\frac{1}{3}}}{\nu}C$, $\left(\frac{3}{3}\right)=\frac{1}{3}B$, and the preceding result gives $\left(\frac{4}{9}\right)=\frac{1}{2\omega\sqrt{2\alpha}}\cdot\frac{\pi}{B}$\ednote{Printed $2\omega\sqrt{2\alpha}$; the formula just above has $2\nu\sqrt{2\alpha}$, so $\omega$ is an apparent misprint for $\nu$.}; we shall show that with these data alone we can determine the transcendent $\left(\frac{1}{1}\right)$.

Indeed it results from the known formulas${}^{(1)}$ that we have $\left(\frac{4}{9}\right)\cdot\left(\frac{8}{3}\right)=\frac{\pi\sin\omega}{3\sqrt{6}}$, and $\left(\frac{8}{3}\right)=2^{\frac{1}{3}}\sin\omega\sqrt{\left(\frac{M_{1}B}{3\cos 2\omega}\right)}$, $\omega$ being $\frac{\pi}{12}$ or $15^{\circ}$, and $M_{1}$ representing the transcendent $\left(\frac{1}{1}\right)$; from this we shall derive
\[ M_{1}=\left(\frac{1}{1}\right)=\frac{2^{\frac{1}{3}}\alpha B}{3}=\frac{2^{\frac{4}{3}}\nu\cos 15^{\circ}}{3}\cdot B. \]
Thus the transcendent $\left(\frac{1}{1}\right)$ depends only on the function $B$. But the other values found of this transcendent will furnish us with new results.

We found by the first method (art.\ 150),
\[ \left(\frac{1}{1}\right)=\frac{2^{-\frac{1}{6}}}{\nu}B+2^{\frac{5}{6}}\sin 15^{\circ}\cdot F^{1}(c). \]

\textbf{(1)} See articles 9 and 18 of the second Part.

\origpage{220}

Therefore, on comparing these two values, there will result
\[ F^{1}(c)=\frac{\alpha}{3}B=\frac{2\nu\cos 15^{\circ}}{3}F^{1}(\sin 45^{\circ}). \]
If we compare the same value with that which resulted from the second method (art.\ 151), we shall have
\[ \frac{2^{\frac{4}{3}}\nu\cos 15^{\circ}}{3}B=2^{-\frac{5}{3}}[F^{1}(b)+F^{1}(c)]; \]
therefore $F^{1}(b)+F^{1}(c)=\frac{8\nu\cos 15^{\circ}}{3}B=4F^{1}(c)$, or
\[ F^{1}(b)=3F^{1}(c). \]
Thus the equality which had been found by approximation (No.~84) is confirmed by a rigorous demonstration, a demonstration which might have been regarded as a very difficult problem of analysis.

The modulus $c$ and its complement $b$, which give rise to this singular reduction, are such that if we put $c=\sin\theta$, $b=\cos\theta$, we shall have $\sin 2\theta=\operatorname{tang}^{2}15^{\circ}$, and $\operatorname{tang}(45^{\circ}+\theta)=\sqrt{\left(\frac{2}{\sqrt{3}}\right)}$, or $\operatorname{tang}(45^{\circ}-\theta)=\sqrt{\cos 30^{\circ}}$.

We see moreover that the functions $F^{1}(b)$, $F^{1}(c)$ are both expressed by the function $F^{1}(\sin 45^{\circ})$, so that we have
\[ \left\{\begin{aligned} F^{1}(c)&=\frac{2\cos 15^{\circ}}{\sqrt[4]{27}}F^{1}(\sin 45^{\circ})\\ F^{1}(b)&=2\sqrt[4]{3}\cdot\cos 15^{\circ}\cdot F^{1}(\sin 45^{\circ}). \end{aligned}\right. \]
These ratios are the more remarkable in that the functions $F^{1}(b)$, $F^{1}(c)$ are up to the present the only examples of functions of the first species which have a known ratio to another function of the same species $F^{1}(\sin 45^{\circ})$, without the two functions compared being comprised in one and the same series or scale formed according to a given modulus, and without their moduli being complements of one another.

END OF THE FIRST PART.

\origpage{221}

\section*{SECOND PART.}

\section*{ON THE EULERIAN INTEGRALS.}

Although the name of Euler is attached to almost all the important theories of the Integral Calculus, I have nevertheless thought that I might be permitted to give more especially the name of \emph{Eulerian integrals} to two kinds of transcendents whose properties have been the subject of several fine Memoirs of Euler, and form the most complete theory that is known up to the present on definite integrals.

The first is the integral $\displaystyle\int\frac{x^{p-1}dx}{\sqrt[n]{(1-x^{n})^{n-q}}}$ which is supposed taken between the limits $x=0$, $x=1$. We shall represent it, like Euler, by the abbreviated character $\left(\frac{p}{q}\right)$.

The second is the integral $\displaystyle\int dx\left(\log\frac{1}{x}\right)^{a-1}$, taken likewise between the limits $x=0$, $x=1$. Euler represents this integral by the symbol $\left[\frac{p}{q}\right]$, supposing that $a$ is equal to the rational fraction $\frac{p}{q}$; we shall represent it more generally by $\Gamma(a)$, and we shall regard $\Gamma(a)$ as a continuous function of $a$.

Having as my object to bring together in this work all that the theory of transcendents, and above all that of definite integrals, offers of the most remarkable, I had to include in it the theory of the Eulerian integrals. This is why I give it here, with some additions, such as I have already published it in the \emph{Mémoires de l'Institut}, for the year 1810.

\origpage{222}

\section*{On the Eulerian integrals of the first species.}

\subsection*{(1).}

The expression $\left(\frac{p}{q}\right)$ by which we designate the integral $\displaystyle\int\frac{x^{p-1}dx}{\sqrt[n]{(1-x^{n})^{n-q}}}$, taken from $x=0$ to $x=1$, is in general a function of the two exponents $p$ and $q$, and of the exponent $n$, which are regarded as whole numbers. But we suppose $n$ constant, and our aim is to compare with one another the various values of $\left(\frac{p}{q}\right)$ that correspond to one and the same value of $n$, in order to reduce to the least possible number the transcendents that this expression contains.

We shall observe first that the two exponents $p$ and $q$ can be interchanged. Indeed, if one makes $x^{n}=1-y^{n}$, one will have
\[ \frac{x^{p-1}dx}{\sqrt[n]{(1-x^{n})^{n-q}}}=-\frac{y^{q-1}dy}{\sqrt[n]{(1-y^{n})^{n-p}}}, \]
and the transformed expression in $y$ must be integrated from $y=1$ to $y=0$. By changing its sign, it must be integrated from $y=0$ to $y=1$; and since one may then put $x$ in place of $y$, one will have
\[ \int\frac{x^{p-1}dx}{\sqrt[n]{(1-x^{n})^{n-q}}}=\int\frac{x^{q-1}dx}{\sqrt[n]{(1-x^{n})^{n-p}}}, \]
or, according to our notation,
\[ \left(\frac{p}{q}\right)=\left(\frac{q}{p}\right),\qquad(a) \]
which is the property stated.

\subsection*{(2).}

It must now be shown that if in the formula $\left(\frac{p}{q}\right)$ one of the two numbers $p$ and $q$ is greater than $n$, the formula is easily brought back to the case where $p$ and $q$ are both comprised within the limits 1 and $n$. For this let $Z=x^{k}(1-x^{n})^{r}$; one will have the differential
\[ dZ=kx^{k-1}dx(1-x^{n})^{r-1}-(k+rn)x^{k+n-1}dx(1-x^{n})^{r-1}, \]

\origpage{223}
whence one derives, by integrating,
\[ Z=k\int x^{k-1}dx(1-x^{n})^{r-1}-(k+rn)\int x^{k+n-1}dx(1-x^{n})^{r-1}. \]
Therefore, if one takes the integrals between the limits $x=0$, $x=1$, and at the same time supposes $k>0$ and $r>0$, so that $Z$ vanishes at both limits, one will have
\[ \int x^{k+n-1}dx(1-x^{n})^{r-1}=\frac{k}{k+rn}\int x^{k-1}dx(1-x^{n})^{r-1}; \]
whence it follows that, by making $k+n=p$ and $r=\frac{q}{n}$, there will result
\[ \left(\frac{p}{q}\right)=\frac{p-n}{p+q-n}\cdot\left(\frac{p-n}{q}\right)\qquad(b). \]
By means of this formula, every function $\left(\frac{p}{q}\right)$ in which one has $p>n$ will be reduced successively to the functions $\left(\frac{p-n}{q}\right)$, $\left(\frac{p-2n}{q}\right)$, etc., as far as a term $\left(\frac{p-in}{q}\right)$ or $\left(\frac{p'}{q}\right)$, in which $p'$ will be the remainder of the division of $p$ by $n$.

Having arrived at this term, if on its part $q$ is greater than $n$, the function $\left(\frac{p'}{q}\right)$, which is the same as $\left(\frac{q}{p'}\right)$, will be reduced successively to the functions $\left(\frac{q-n}{p'}\right)$, $\left(\frac{q-2n}{p'}\right)$, etc.\ as far as a term $\left(\frac{q'}{p'}\right)$ or $\left(\frac{p'}{q'}\right)$, in which $q'$ will be the remainder of the division of $q$ by $n$.

From this one sees that one can always suppose the function $\left(\frac{p}{q}\right)$ reduced to a form in which $p$ and $q$ are both comprised in the sequence $1, 2, 3\ldots n$.

\subsection*{(3).}

This being laid down, there are two principal cases in which one can find immediately the value of $\left(\frac{p}{q}\right)$; namely, when $n$ is equal to one of the two numbers $p$ and $q$, or when it is equal to their sum.

Let, 1\textsuperscript{st}. $q=n$; one will have immediately $\left(\frac{p}{n}\right)=\displaystyle\int x^{p-1}dx=\frac{x^{p}}{p}+C$, an integral which, being taken from $x=0$ to $x=1$, reduces to $\frac{1}{p}$; whence results
\[ \left(\frac{p}{n}\right)=\frac{1}{p}\qquad(c). \]

\origpage{224}

Let, 2\textsuperscript{nd}. $p+q=n$, or $p=a$, $q=n-a$; one will have $\left(\frac{p}{q}\right)=\displaystyle\int x^{a-1}dx(1-x^{n})^{-\frac{a}{n}}$; making $1-x^{n}=x^{n}z^{n}$, one will have the rational transformed expression
\[ \left(\frac{p}{q}\right)\text{ or }\left(\frac{a}{n-a}\right)=\int\frac{z^{n-a-1}dz}{1+z^{n}}, \]
an integral which must be taken from $z=0$ to $z=\infty$; now by the known formulas (\emph{Eul.}, \emph{Calc.\ intég.}, tom.~I, pag.~252), this integral $=\frac{\pi}{n\sin\frac{a\pi}{n}}$. Therefore, if one makes $\frac{\pi}{n}=\omega$, one will have
\[ \left(\frac{a}{n-a}\right)=\frac{\omega}{\sin a\omega}.\qquad(d) \]
Except for these two general cases, all the quantities designated by $\left(\frac{p}{q}\right)$ are transcendents more or less composite, according to the value of $n$, and are in no way susceptible of an exact evaluation. But for each value of $n$, one can express all these transcendents by means of a small number of them, and this is the object of the following researches.

\subsection*{(4).}

Let us observe first that, by putting $p+n$ in place of $p$, the equation (b) gives
\[ \left(\frac{p}{q}\right)=\frac{p+q}{p}\cdot\left(\frac{p+n}{q}\right). \]
One would have similarly
\[ \begin{aligned} \left(\frac{p+n}{q}\right)&=\frac{p+q+n}{p+n}\cdot\left(\frac{p+2n}{q}\right)\\ \left(\frac{p+2n}{q}\right)&=\frac{p+q+2n}{p+2n}\cdot\left(\frac{p+3n}{q}\right)\\ &\text{etc.} \end{aligned} \]
Therefore in general, $i$ being any positive whole number whatever, one will have
\[ \left(\frac{p}{q}\right)=\frac{p+q}{p}\cdot\frac{p+q+n}{p+n}\cdot\frac{p+q+2n}{p+2n}\ldots\frac{p+q+in}{p+in}\cdot\left(\frac{p+\overline{i+1}\cdot n}{q}\right); \]

\origpage{225}
putting $p+r$ in place of $p$, one will have similarly
\[ \left(\frac{p+r}{q}\right)=\frac{p+q+r}{p+r}\cdot\frac{p+q+r+n}{p+r+n}\cdot\frac{p+q+r+2n}{p+r+2n}\ldots\frac{p+q+r+in}{p+r+in}\cdot\left(\frac{p+r+\overline{i+1}\cdot n}{q}\right). \]
Let us divide these two formulas by one another, and make, for brevity,
\[ \begin{aligned} M^{0}&=\frac{(p+q)(p+r)}{p\cdot(p+q+r)}\\ M'&=\frac{(p+q+n)(p+r+n)}{(p+n)(p+q+r+n)}\\ M''&=\frac{(p+q+2n)(p+r+2n)}{(p+2n)(p+q+r+2n)}\\ &\text{etc.}, \end{aligned} \]
we shall have
\[ \frac{\left(\frac{p}{q}\right)}{\left(\frac{p+r}{q}\right)}=M^{0}M'M''\ldots M^{(i)}\frac{\left(\frac{p+\overline{i+1}\cdot n}{q}\right)}{\left(\frac{p+r+\overline{i+1}\cdot n}{q}\right)}. \]
Let us suppose $r$ positive and $<n$; then it is clear that the quantity $\left(\frac{p+r+\overline{i+1}\cdot n}{q}\right)$ will be smaller than $\left(\frac{p+\overline{i+1}\cdot n}{q}\right)$ and greater than $\left(\frac{p+\overline{i+2}\cdot n}{q}\right)$. For if one considers various formulas $\left(\frac{p}{q}\right)$ which differ only by the exponent $p$; since these formulas each designate the integral, taken from $x=0$ to $x=1$, of a differential $\frac{x^{p-1}dx}{X}$ whose denominator is the same for all, it is evident that the integral will be the smaller the greater $p$ is.

But the formula (b) gives
\[ \left(\frac{p+\overline{i+1}\cdot n}{q}\right)=\frac{p+q+\overline{i+1}\cdot n}{p+\overline{i+1}\cdot n}\cdot\left(\frac{p+\overline{i+2}\cdot n}{q}\right). \]
Therefore we have on the one hand the ratio
\[ \frac{\left(\frac{p+\overline{i+1}\cdot n}{q}\right)}{\left(\frac{p+r+\overline{i+1}\cdot n}{q}\right)}>1, \]

\origpage{226}
and on the other this same ratio $<\frac{p+q+\overline{i+1}\cdot n}{p+\overline{i+1}\cdot n}$. If therefore one wishes this ratio to be comprised between the limits 1 and $1+\frac{1}{k}$, one must take $i+1>\frac{qk-p}{n}$, and then one will have
\[ \frac{\left(\frac{p}{q}\right)}{\left(\frac{p+r}{q}\right)}=M^{0}M'M''\ldots M^{(i)}\left(1+\frac{1}{k'}\right), \]
$k'$ being greater than $k$.

One knows by this equation how closely one approaches the ratio of $\left(\frac{p}{q}\right)$ to $\left(\frac{p+r}{q}\right)$ by prolonging the product $M^{0}M'M''\ldots$ as far as a term $M^{(i)}$; and it is clear that by continuing this product to infinity, one will have the true value of this ratio, which will be
\[ \frac{\left(\frac{p}{q}\right)}{\left(\frac{p+r}{q}\right)}=M^{0}M'M''M''',\ \text{etc.} \]
Now if in this equation one interchanges the letters $q$ and $r$, the quantities $M^{0}$, $M'$, $M''$, etc.\ will remain the same, so that one will again have
\[ \frac{\left(\frac{p}{r}\right)}{\left(\frac{p+q}{r}\right)}=M^{0}M'M''M''',\ \text{etc.} \]
Therefore, by comparing these equations, one obtains the general formula
\[ \left(\frac{p}{q}\right)\cdot\left(\frac{p+q}{r}\right)=\left(\frac{p}{r}\right)\cdot\left(\frac{p+r}{q}\right)\ldots\ldots(e). \]

This formula, the discovery of which belongs to Euler, is a kind of equation of finite differences, which contains almost the whole theory of the transcendents $\left(\frac{p}{q}\right)$. And first we are going to deduce from it the general expression of the quantities $\left(\frac{p}{q}\right)$.

\origpage{227}

\subsection*{(5).}

The formulas (c) and (d) give the exact values of the function $\left(\frac{p}{q}\right)$ whenever one of the two numbers $p$ and $q$, or their sum, is equal to $n$. Let us suppose now that we also know all the values of $\left(\frac{p}{q}\right)$ when $p+q=n-1$, and let us designate in general by $A_{a}$ the function $\left(\frac{n-a-1}{a}\right)$, so that one has
\[ \left(\frac{n-a-1}{a}\right)=A_{a}\ldots\ldots(f). \]
One will therefore have successively $\left(\frac{n-2}{1}\right)=A_{1}$, $\left(\frac{n-3}{2}\right)=A_{2}$; $\left(\frac{n-4}{3}\right)=A_{3}$, etc.; and because $\left(\frac{q}{p}\right)=\left(\frac{p}{q}\right)$, one will also have $\left(\frac{1}{n-2}\right)=A_{1}$, $\left(\frac{2}{n-3}\right)=A_{2}$, etc.; therefore in general
\[ A_{(n-1-a)}=A_{a}\ldots\ldots(g). \]
Whence one sees that the number of the auxiliaries $A_{1}$, $A_{2}$, $A_{3}$, etc., is always reduced to $\frac{n-2}{2}$ or $\frac{n-1}{2}$, according as $n$ is even or odd.

For example, if $n=7$, there will be three auxiliaries $A_{1}=\left(\frac{5}{1}\right)$, $A_{2}=\left(\frac{4}{2}\right)$, $A_{3}=\left(\frac{3}{3}\right)$, since one would have $A_{4}=\left(\frac{2}{4}\right)=A_{2}$, and $A_{5}=\left(\frac{1}{5}\right)=A_{1}$.

If $n=8$, there will again be only three auxiliaries $A_{1}$, $A_{2}$, $A_{3}$; for one would have, by virtue of the equation (g), $A_{4}=A_{3}$, $A_{5}=A_{2}$, $A_{6}=A_{1}$.

This being laid down, by means of the equations (c), (d), (e) and of the auxiliaries given by the equation (f), we shall be able to find the general expression of $\left(\frac{p}{q}\right)$ in two general cases: 1\textsuperscript{st}. when $p+q$ is $<n$; 2\textsuperscript{nd}. when $p+q$ is $>n$. Here is how one arrives at them.

\subsection*{(6).}

The equation (e) gives immediately $\left(\frac{n-a-1}{a}\right)\cdot\left(\frac{n-1}{1}\right)=\left(\frac{n-a-1}{1}\right)\cdot\left(\frac{n-a}{a}\right)$; substituting in this the known values

\origpage{228}
by the equations (d) and (g), one will have
\[ \left(\frac{n-a-1}{1}\right)=\frac{A_{a}\sin a\omega}{\sin\omega}\ldots\ldots(h). \]
Thus one knows every function $\left(\frac{p}{q}\right)$ in which one of the two numbers $p$ and $q$ is equal to unity; if for greater simplicity one puts $a$ in place of $n-a-1$, one will have for the same object the formula
\[ \left(\frac{a}{1}\right)=\frac{A_{a}\sin(a+1)\omega}{\sin\omega}\ldots\ldots(i). \]
The same equation (e) gives $\left(\frac{n-a-2}{a}\right)\cdot\left(\frac{n-2}{1}\right)=\left(\frac{n-a-2}{1}\right)\cdot\left(\frac{n-a-1}{a}\right)$, and substituting the known values, there results:
\[ \left(\frac{n-a-2}{a}\right)=\frac{A_{a}A_{a+1}}{A_{1}}\cdot\frac{\sin(a+1)\omega}{\sin\omega}. \]
Thus one has the value of every function $\left(\frac{p}{q}\right)$ in which $p+q=n-2$.

From the equation (e) one further deduces immediately
\[ \left(\frac{n-a-3}{a}\right)\cdot\left(\frac{n-3}{1}\right)=\left(\frac{n-a-3}{1}\right)\cdot\left(\frac{n-a-2}{a}\right); \]
which gives
\[ \left(\frac{n-a-3}{a}\right)=\frac{A_{a}A_{a+1}A_{a+2}}{A_{1}A_{2}}\cdot\frac{\sin(a+1)\omega\sin(a+2)\omega}{\sin\omega\sin 2\omega}. \]
Thus one knows the value of every function $\left(\frac{p}{q}\right)$ in which $p+q=n-3$.

In general the equation (e) gives
\[ \left(\frac{n-a-k}{a}\right)\cdot\left(\frac{n-k}{1}\right)=\left(\frac{n-a-k}{1}\right)\cdot\left(\frac{n-a-k+1}{a}\right), \]
and by putting in the values drawn from the equation (h), there results
\[ \left(\frac{n-a-k}{a}\right)=\frac{A_{a+k-1}\sin(a+k-1)\omega}{A_{k-1}\sin(k-1)\omega}\cdot\left(\frac{n-a-k+1}{a}\right). \]
Therefore one will have in general
\[ \left(\frac{n-a-k}{a}\right)=\frac{A_{a}A_{a+1}\ldots A_{a+k-1}}{A_{1}A_{2}\ldots A_{k-1}}\cdot\frac{\sin(a+1)\omega\sin(a+2)\omega\ldots\sin(a+k-1)\omega}{\sin\omega\sin 2\omega\ldots\sin(k-1)\omega}\ldots(k). \]

\origpage{229}
This is the value of every function $\left(\frac{p}{q}\right)$ in which $p+q$ is less than $n$.

\subsection*{(7).}

To have the general expression of $\left(\frac{p}{q}\right)$ when $p+q$ is greater than $n$, let us observe that the equation (e) also gives
\[ \left(\frac{n-a+k}{a}\right)\cdot\left(\frac{n+k}{1}\right)=\left(\frac{n-a+k}{1}\right)\cdot\left(\frac{n-a+k+1}{a}\right). \]
Now by the equation (b) one has
\[ \left(\frac{n+k}{1}\right)=\frac{k}{k+1}\cdot\left(\frac{k}{1}\right); \]
by the equation (i) one has
\[ \left(\frac{k}{1}\right)=\frac{A_{k}\sin(k+1)\omega}{\sin\omega}, \]
and by the equation (h) one has
\[ \left(\frac{n-a+k}{1}\right)=\frac{A_{a-k-1}\sin(a-k-1)\omega}{\sin\omega}; \]
substituting all these values, there will result
\[ \left(\frac{n-a+k+1}{a}\right)=\frac{k}{k+1}\cdot\frac{A_{k}\sin(k+1)}{A_{a-k-1}\sin(a-k-1)\omega}\cdot\left(\frac{n-a+k}{a}\right)\ldots(l). \]
\ednote{The numerator is printed $\sin(k+1)$ without $\omega$; $\sin(k+1)\omega$ is meant, as in equation (i).}
Everything is therefore reduced to finding the value of $\left(\frac{n-a+1}{a}\right)$.

Now the equation (e) gives
\[ \left(\frac{n-a}{a}\right)\cdot\left(\frac{n}{1}\right)=\left(\frac{n-a}{1}\right)\cdot\left(\frac{n-a+1}{a}\right), \]
substituting the known values, there will result
\[ \left(\frac{n-a+1}{a}\right)=\frac{1}{A_{a-1}}\cdot\frac{\omega\sin\omega}{\sin a\omega\sin(a-1)\omega}\ldots\ldots(m). \]
From this and from the equation (l) one deduces successively
\[ \begin{aligned} \left(\frac{n-a+2}{a}\right)&=\frac{1}{2}\cdot\frac{A_{1}}{A_{a-1}A_{a-2}}\cdot\frac{\omega\sin\omega\sin 2\omega}{\sin a\omega\sin(a-1)\omega\sin(a-2)\omega},\\ \left(\frac{n-a+3}{a}\right)&=\frac{1}{3}\cdot\frac{A_{1}A_{2}}{A_{a-1}A_{a-2}A_{a-3}}\cdot\frac{\omega\sin\omega\sin 2\omega\sin 3\omega}{\sin a\omega\sin(a-1)\omega\sin(a-2)\omega\sin(a-3)\omega}, \end{aligned} \]

\origpage{230}
and in general,
\[ \left(\frac{n-a+k}{a}\right)=\frac{1}{k}\cdot\frac{A_{1}A_{2}\ldots A_{k-1}}{A_{a-1}A_{a-2}\ldots A_{a-k}}\cdot\frac{\omega\sin\omega\sin 2\omega\ldots\sin k\omega}{\sin a\omega\sin(a+1)\omega\ldots\sin(a-k)\omega}\ldots(n). \]
\ednote{The denominator is printed $\sin(a+1)\omega$; $\sin(a-1)\omega$ is meant, as in the two preceding formulas.}
This formula gives the value of the function $\left(\frac{p}{q}\right)$ when $p+q$ exceeds $n$; thus, by uniting it with the formula (k), one has generally the expression of every function $\left(\frac{p}{q}\right)$, supposing known only the similar functions for which $p+q=n-1$, and it has already been remarked that the number of these auxiliaries is $\frac{n+2}{2}$\ednote{Printed $\frac{n+2}{2}$; $\frac{n-2}{2}$ is meant, as stated in article (5).} or $\frac{n-1}{2}$, according as $n$ is even or odd.

\subsection*{(8).}

The formula (n) could be regarded as a continuation of the formula (k); and one would thus have only one and the same formula for all the values of $p$ and $q$; but then one would need new auxiliaries $A_{0}$, $A_{-1}$, $A_{-2}$, etc., and one would have to fix their values thus: $A_{0}=\frac{\omega}{\sin 0\omega}$, or rather, $i$ being infinitely small, $A_{i}=\frac{\omega}{\sin i\omega}=\frac{1}{i}$, $A_{-1}=-A_{1}$, $A_{-2}=-A_{2}$, etc. It is to avoid the embarrassment of these substitutions, especially in the case of $A_{0}$, that we have given the two formulas separately.

It is rather astonishing that the general expression of the functions $\left(\frac{p}{q}\right)$ should have escaped Euler; one sees however that he had occupied himself especially with this research, by the passage of volume V of the \emph{Nova acta Petropol.}, pag.\ 125, where he says: \emph{Neque tamen hinc adhuc elucet quânam lege omnes determinationes progrediantur, quandoquidem valores certarum formularum continuò magis evadunt complicati.}

We shall remark moreover that the formulas (k) and (n), which contain the general expression in question, can be regarded as the complete integral of the equation of finite differences (e); so that one can draw from this equation no consequence which is not contained in the formulas (k) and (n). This would be easy to demonstrate by the methods that one follows in this kind of analysis, and which, for the most part, have been indicated

\origpage{231}
by Lagrange in the \emph{Mémoires de l'Académie de Berlin}, an.\ 1775.

\subsection*{(9).}

Here now are some particular formulas which deserve to be cited. From the equation (e) one deduces the two following,
\[ \begin{gathered} \left(\frac{p}{q}\right)\cdot\left(\frac{p+q}{n-p-q}\right)=\left(\frac{p}{n-p-q}\right)\cdot\left(\frac{n-q}{q}\right),\\ \left(\frac{p}{n-p-q}\right)\cdot\left(\frac{n-q}{n-p}\right)=\left(\frac{p}{n-p}\right)\cdot\left(\frac{n}{n-p-q}\right); \end{gathered} \]
multiplying these two equations together, and putting into the product the values known by the formulas (c) and (d), one will have
\[ \left(\frac{p}{q}\right)\cdot\left(\frac{n-p}{n-q}\right)=\frac{\omega\sin(p+q)\omega}{(n-p-q)\sin p\omega\sin q\omega}\ldots\ldots(p); \]
whence it follows that the value of $\left(\frac{n-p}{n-q}\right)$ is deduced immediately from that of $\left(\frac{p}{q}\right)$, which is in some sort its complement. One has in particular,
\[ \left(\frac{a}{a}\right)\cdot\left(\frac{n-a}{n-a}\right)=\frac{2\omega\cot a\omega}{n-2a}\ldots\ldots(q). \]
Thus, knowing the values of $\left(\frac{a}{a}\right)$ when $a$ does not exceed $\frac{1}{2}n$, one deduces from them the values of $\left(\frac{a}{a}\right)$ when $a$ is greater than $\frac{1}{2}n$.

\subsection*{(10).}

To examine more particularly the functions of the form $\left(\frac{a}{a}\right)$, let us take up again the primitive value of these functions, which is
\[ \left(\frac{a}{a}\right)=\int\frac{x^{a-1}\,dx}{\sqrt[n]{(1-x^{n})^{n-a}}}; \]
if one makes
\[ 1-x^{n}=\frac{z^{n}}{4x^{n}},\quad\text{or}\quad x^{n}=\tfrac{1}{2}\pm\tfrac{1}{2}\sqrt{(1-z^{n})}, \]
the transformed expression will be
\[ \left(\frac{a}{a}\right)=\mp 2^{-\frac{2a}{n}}\int\frac{z^{a-1}\,dz}{\sqrt{(1-z^{n})}}. \]

\origpage{232}
As to the limits of this new integral, it must be observed that the values $x^{n}=0$, $x^{n}=\frac{1}{2}$, $x^{n}=1$ give respectively $z=0$, $z=1$, $z=0$. Whence one sees that the integral in $z$ must be taken twice, from $z=0$ to $z=1$. And since then nothing prevents putting $x$ in place of $z$, one will have
\[ \left(\frac{a}{a}\right)=2^{1-\frac{2a}{n}}\int\frac{x^{a-1}\,dx}{\sqrt{(1-x^{n})}}\ldots\ldots(r); \]
this integral is thus reduced to the simplest form of which it is susceptible, since the radical is now only of the second degree.

\subsection*{(11).}

If in this formula one puts $n-a$ in place of $a$, one will have
\[ \left(\frac{n-a}{n-a}\right)=2^{-1+\frac{2a}{n}}\int\frac{x^{n-a-1}\,dx}{\sqrt{(1-x^{n})}}; \]
from this and from the equation (q) results this remarkable formula
\[ \int\frac{x^{a-1}\,dx}{\sqrt{(1-x^{n})}}\cdot\int\frac{x^{n-a-1}\,dx}{\sqrt{(1-x^{n})}}=\frac{2\omega\cot a\omega}{n-2a}\ldots\ldots(s). \]

\subsection*{(12).}

Since the functions $\left(\frac{a}{a}\right)$ are the simplest among the functions $\left(\frac{p}{q}\right)$ not comprised in the formulas (c) and (l), it would seem convenient to substitute them for the auxiliaries designated by $A_{a}$, in order to express by their means all the functions $\left(\frac{p}{q}\right)$.

With this view, let us designate in general the function $\left(\frac{a}{a}\right)$ by $M_{a}$; as one can suppose $a<\frac{1}{2}n$, one will have by the formula (k),
\[ M_{a}=\frac{A_{a}A_{a+1}\ldots A_{n-a-1}}{A_{1}A_{2}\ldots A_{n-2a-1}}\cdot\frac{\sin(a+1)\omega\sin(a+2)\omega\ldots\sin(n-a-1)\omega}{\sin\omega\sin 2\omega\ldots\sin(n-2a-1)\omega}, \]
a value which, by means of the equations $A_{n-1-k}=A_{k}$\ednote{The first letter of this subscript is badly printed; it is read as $n$, following equation (g).}, $\sin(n-k)\omega=\sin k\omega$, is reduced to this form
\[ M_{a}=\frac{A_{a}A_{a+1}\ldots A_{2a-1}}{A_{1}A_{2}\ldots A_{a-1}}\cdot\frac{\sin(a+1)\omega\sin(a+2)\omega\ldots\sin 2a\omega}{\sin\omega\sin 2\omega\ldots\sin a\omega}\ldots\ldots(t); \]

\origpage{233}
whence one draws successively
\[ \begin{aligned} M_{1}&=A_{1}\frac{\sin 2\omega}{\sin\omega}\\ M_{2}&=\frac{A_{2}A_{3}}{A_{1}}\cdot\frac{\sin 3\omega\sin 4\omega}{\sin\omega\sin 2\omega}\\ M_{3}&=\frac{A_{3}A_{4}A_{5}}{A_{1}A_{2}}\cdot\frac{\sin 4\omega\sin 5\omega\sin 6\omega}{\sin\omega\sin 2\omega\sin 3\omega}\\ &\text{etc.} \end{aligned} \]
These equations, which one can also put under the form
\[ \left.\begin{aligned} A_{1}&=M_{1}\cdot\frac{\sin\omega}{\sin 2\omega}\\ \frac{A_{2}A_{3}}{A_{1}A_{1}}&=\frac{M_{2}}{M_{1}}\cdot\frac{\sin^{2}2\omega}{\sin 3\omega\sin 4\omega}\\ \frac{A_{4}A_{5}}{A_{2}A_{2}}&=\frac{M_{3}}{M_{2}}\cdot\frac{\sin^{2}3\omega}{\sin 5\omega\sin 6\omega}\\ &\text{etc.} \end{aligned}\right\}\qquad(u), \]
will serve to determine the auxiliaries $A_{1}$, $A_{2}$, $A_{3}$, etc.\ by means of an equal number of quantities $M_{1}$, $M_{2}$, $M_{3}$, etc., taken in the suitable order. One will therefore be able to express by these latter quantities all the functions $\left(\frac{p}{q}\right)$ which correspond to one and the same value of $n$.

But it must be observed that these substitutions can be carried out only for particular values of $n$, and that thus by the employment of the auxiliaries $M_{a}$, one cannot arrive at formulas as general as the formulas (k) and (n) are.

\subsection*{(13).}

Let us now consider the formula
\[ \left(\frac{n-2a}{a}\right)=\int\frac{x^{n-2a-1}\,dx}{\sqrt[n]{(1-x^{n})^{n-a}}}: \]
if one makes $x^{-n}=\frac{1}{2}+\frac{1}{2}\sqrt{(1-z^{n})}$\ednote{Printed $1-z^{n}$ under the radical; the transformée (v) requires $x^{-n}=\frac{1}{2}+\frac{1}{2}\sqrt{(1+z^{n})}$.}, one will have the transformed expression
\[ \left(\frac{n-2a}{a}\right)=2^{-\frac{2a}{n}}\int\frac{z^{a-1}\,dz}{\sqrt{(1+z^{n})}}\qquad(v), \]
in which the integral must be taken from $z=0$ to $z=\infty$, and which moreover supposes $a<\frac{1}{2}n$. But by virtue of the equations

\origpage{234}
(e) and (p), one has
\[ \begin{gathered} \left(\frac{n-2a}{a}\right)\cdot\left(\frac{2a}{n-a}\right)=\left(\frac{\omega}{a\sin 2a\omega}\right),\\ \left(\frac{a}{a}\right)\cdot\left(\frac{2a}{n-a}\right)=\left(\frac{a}{n-a}\right)\cdot\left(\frac{n}{a}\right)=\frac{\omega}{a\sin a\omega}; \end{gathered} \]
\ednote{In the first line the value $\frac{\omega}{a\sin 2a\omega}$ is printed inside the parentheses of the transcendent symbol; the plain quotient is meant, as in the second line.}
therefore $\left(\frac{a}{a}\right)=2\cos a\omega\cdot\left(\frac{n-2a}{a}\right)$, or, by substituting the value given by the equation (v),
\[ \left(\frac{a}{a}\right)=2^{1-\frac{2a}{n}}\cos a\omega\cdot\int\frac{z^{a-1}\,dz}{\sqrt{(1+z^{n})}}\qquad(x). \]
This formula holds only when $a$ is $<\frac{1}{2}n$; if $a$ is $>\frac{1}{2}n$, one will begin by taking the value of $\left(\frac{n-a}{n-a}\right)$, which will be
\[ \left(\frac{n-a}{n-a}\right)=2^{\frac{2a}{n}-1}\cos(n-a)\omega\cdot\int\frac{z^{n-a-1}\,dz}{\sqrt{(1+z^{n})}}, \]
and one will deduce from it that of $\left(\frac{a}{a}\right)$ by means of the equation (q). One will thus have, $a$ being $>\frac{1}{2}n$:
\[ \left(\frac{a}{a}\right)=\frac{2^{2-\frac{2a}{n}}}{2a-n}\cdot\frac{\omega}{\sin a\omega}:\int\frac{z^{n-a-1}\,dz}{\sqrt{(1+z^{n})}}\qquad(y). \]

\subsection*{(14).}

If one now compares the equations (r), (x) and (y), one will draw from them the formulas
\[ \left.\begin{aligned} \int\frac{x^{a-1}\,dx}{\sqrt{(1-x^{n})}}&=\cos a\omega\cdot\int\frac{z^{a-1}\,dz}{\sqrt{(1+z^{n})}}\\ \int\frac{x^{a-1}\,dx}{\sqrt{(1-x^{n})}}&=\frac{2\omega}{(2a-n)\sin a\omega}:\int\frac{z^{n-a-1}\,dz}{\sqrt{(1+z^{n})}} \end{aligned}\right\}\qquad(z), \]
the first holding when $a$ is $<\frac{1}{2}n$, and the second when $a$ is $>\frac{1}{2}n$.

When $n$ is odd, if one makes in the first equation $x=1-y^{2}$ and $z=y^{2}-1$, the integral formulas comprised in the two members will be reduced both to the form
\[ \int\frac{(1-y^{2})^{a-1}\,dy}{\sqrt{\left(n-\frac{n\cdot n-1}{1\cdot 2}y^{2}+\frac{n\cdot n-1\cdot n-2}{1\cdot 2\cdot 3}y^{4}-\text{etc.}\right)}}: \]

\origpage{235}
one sees therefore that the part of this last integral taken from $y=0$ to $y=1$, and the part taken from $y=1$ to $y=\infty$, are to each other $::\cos\frac{a\pi}{n}:1$.

\subsection*{(15).}

Let us consider anew the formula
\[ \left(\frac{p}{q}\right)=\int\frac{x^{p-1}\,dx}{\sqrt[n]{(1-x^{n})^{n-q}}}, \]
if one makes
\[ 1-x^{n}=\frac{x^{2n}}{4z^{n}}, \]
one obtains first
\[ \left(\frac{p}{q}\right)=2^{2-\frac{2q}{n}}\int x^{p+2q-2n-1}\,dx\cdot z^{n-q}. \]
Let $p=2a$ and $q=n-a$; one will have
\[ \left(\frac{2a}{n-a}\right)=2^{\frac{2a}{n}}\int z^{a}\cdot\frac{dx}{x}, \]
and by completing the substitutions, there will result
\[ \int\left(z^{a-1}\,dz-\frac{z^{a+\frac{1}{2}n-1}\,dz}{\sqrt{(1+z^{n})}}\right)=2^{1-\frac{2a}{n}}\cdot\left(\frac{2a}{n-a}\right); \]
this being laid down, two cases must be distinguished, according as $a$ is $<\frac{1}{2}n$ or $>\frac{1}{2}n$.

Let, 1\textsuperscript{st}. $a<\frac{1}{2}n$; one will have by the formulas of no.~13,
\[ \left(\frac{2a}{n-a}\right)=\frac{\omega}{a\sin 2a\omega}:\left(\frac{n-2a}{a}\right)=2^{\frac{2a}{n}}\cdot\frac{\omega}{a\sin 2a\omega}:\int\frac{z^{a-1}\,dz}{\sqrt{(1+z^{n})}}, \]
therefore
\[ \int\frac{z^{a-1}\,dz}{\sqrt{(1+z^{n})}}\cdot\int\left(z^{a-1}\,dz-\frac{z^{a+\frac{1}{2}n-1}\,dz}{\sqrt{(1+z^{n})}}\right)=\frac{2\omega}{a\sin 2a\omega}\qquad(a'). \]
Let, 2\textsuperscript{nd}. $a>\frac{1}{2}n$; one will have by the equation (b),
\[ \left(\frac{2a}{n-a}\right)=\frac{2a-n}{a}\cdot\left(\frac{2a-n}{n-a}\right); \]
but by making $a=n-c$, one has
\[ \left(\frac{2a-n}{n-a}\right)=\left(\frac{n-2c}{c}\right), \]

\origpage{236}
and by the equation (v), one has
\[ \frac{n-2c}{c}=2^{-\frac{2c}{n}}\int\frac{z^{c-1}\,dz}{\sqrt{(1+z^{n})}}. \]
\ednote{The left side is printed as a plain fraction without the parentheses of the transcendent symbol; $\left(\frac{n-2c}{c}\right)$ is meant, as in the preceding line.}
Therefore in place of the equation (a'), one will have
\[ \int\left(z^{a-1}\,dz-\frac{z^{a+\frac{1}{2}n-1}\,dz}{\sqrt{(1+z^{n})}}\right)=\frac{2a-n}{2a}\int\frac{z^{n-a-1}\,dz}{\sqrt{(1+z^{n})}}\qquad(b'). \]
Moreover this last equation is verified immediately by means of the function $P=z^{a-\frac{1}{2}n}\sqrt{(1+z^{n})}-z^{a}$, which vanishes at both limits, when $z=0$ and when $z=\infty$; for if one takes the differential of this function, and then integrates it, one will find
\[ \int\left(z^{a-1}\,dz-\frac{z^{a+\frac{1}{2}n-1}\,dz}{\sqrt{(1+z^{n})}}\right)=\frac{2a-n}{2a}\cdot\int\frac{z^{a-\frac{n}{2}-1}\,dz}{\sqrt{(1+z^{n})}}; \]
a formula which does not differ from the preceding, because, by putting $\frac{1}{z}$ in place of $z$, the integral $\displaystyle\int\frac{z^{a-\frac{1}{2}n-1}\,dz}{\sqrt{(1+z^{n})}}$ is changed into $\displaystyle\int\frac{z^{n-a-1}\,dz}{\sqrt{(1+z^{n})}}$, the limits being always $z=0$, $z=\infty$.

\subsection*{(16).}

Let us consider finally, on the supposition of $n$ even, the formula
\[ \left(\frac{\frac{1}{2}n-a}{a}\right)=\int\frac{x^{\frac{1}{2}n-a-1}\,dx}{\sqrt[n]{(1-x^{n})^{n-a}}}; \]
if one makes $x^{-n}=1+z^{n}$, one will have by the substitution,
\[ \left(\frac{\frac{1}{2}n-a}{a}\right)=\int\frac{z^{a-1}\,dz}{\sqrt{(1+z^{n})}},\qquad(c) \]
\ednote{The label is printed (c) without a prime; the text below cites this equation as (c').}
and the integral of the second member must be taken between the limits $z=0$, $z=\infty$.

Now, by the combination of the equations (x) and (c'), one obtains
\[ \left(\frac{a}{a}\right)=2^{1-\frac{2a}{n}}\cos a\omega\left(\frac{\frac{1}{2}n-a}{a}\right), \]
and consequently also,
\[ \left(\frac{\frac{1}{2}n-a}{\frac{1}{2}n-a}\right)=2^{\frac{2a}{n}}\sin a\omega\left(\frac{a}{\frac{1}{2}n-a}\right). \]

\origpage{237}
From these two one will conclude
\[ \left(\frac{a}{a}\right)=2^{1-\frac{4a}{n}}\operatorname{cotang}a\omega\cdot\left(\frac{\frac{1}{2}n-a}{\frac{1}{2}n-a}\right)\qquad(d'), \]
whence one sees that, $n$ being even, it suffices to have the values of $\left(\frac{a}{a}\right)$ for all the cases where $a$ does not exceed $\frac{1}{4}n$, and that thus the number of the auxiliaries necessary to determine all the functions $\left(\frac{p}{q}\right)$ which correspond to one and the same value of $n$ is reduced to $\frac{n}{4}$ or $\frac{n-2}{4}$, according as $n$ is of the form $4m$ or $4m+2$.

\subsection*{(17).}

With the aid of the equation (d'), one will find relations between the auxiliaries $A_{1}$, $A_{2}$, $A_{3}$, etc.\ which will reduce their number as has just been said.

For this let us take up again the formula (t),
\[ \left(\frac{a}{a}\right)=\frac{A_{a}A_{a+1}\ldots A_{2a-1}}{A_{1}A_{2}\ldots\ldots A_{a-1}}\cdot\frac{\sin(a+1)\omega\sin(a+2)\omega\ldots\sin 2a\omega}{\sin\omega\sin 2\omega\ldots\sin a\omega}, \]
it gives, on making $n=2m$,
\[ \left(\frac{m-a}{m-a}\right)=\frac{A_{m-a}\cdot A_{m-a+1}\ldots A_{2m-2a-1}}{A_{1}A_{2}\ldots A_{m-a-1}}\cdot\frac{\sin(m-a+1)\omega\sin(m-a+2)\omega\ldots\sin(2m-2a)\omega}{\sin\omega\sin 2\omega\ldots\sin(m-a)\omega}; \]
substituting these values in the equation (d'), and making the reductions on the hypothesis $a<\frac{1}{2}m$, one will have generally
\[ \frac{A_{a}A_{a+1}\ldots A_{2a-1}}{A_{m-1}A_{m-2}\ldots A_{m-a}}=2^{-\frac{a}{m}}\cdot\frac{\sin(m-1)\omega\sin(m-2)\omega\ldots\sin(m-a+1)\omega}{\sin a\omega\sin(a+1)\omega\ldots\sin(2a-1)\omega}\qquad(e'). \]
From this result, on making successively $a=1$, $2$, $3$, etc., particular equations which can be put under this form
\[ \left.\begin{aligned} A_{m-1}&=A_{1}\cdot 2^{\frac{1}{m}}\sin\omega\\ A_{m-2}&=\frac{A_{2}A_{3}}{A_{1}}\cdot 2^{1+\frac{1}{m}}\sin 3\omega\\ A_{m-3}&=\frac{A_{4}A_{5}}{A_{2}}\cdot 2^{1+\frac{1}{m}}\sin 5\omega\\ A_{m-4}&=\frac{A_{6}A_{7}}{A_{3}}\cdot 2^{1+\frac{1}{m}}\sin 7\omega\\ &\text{etc.} \end{aligned}\right\}\qquad(f'). \]

\origpage{238}
They must be continued until their number is $\frac{n-4}{4}$ or $\frac{n-2}{4}$.

\subsection*{(18).}

For example, when $n=12$, there are five auxiliaries $A_{1}=\left(\frac{10}{1}\right)$, $A_{2}=\left(\frac{9}{2}\right)$, $A_{3}=\left(\frac{8}{3}\right)$, $A_{4}=\left(\frac{7}{4}\right)$, $A_{5}=\left(\frac{6}{5}\right)$, among which one has these two equations
\[ \begin{aligned} A_{5}&=A_{1}\cdot 2^{\frac{1}{6}}\sin\omega\\ A_{4}&=\frac{A_{2}A_{3}}{A_{1}}\cdot 2^{\frac{7}{6}}\sin 3\omega. \end{aligned} \]
So that the number of necessary auxiliaries is reduced to three, for which one may take $A_{1}$, $A_{2}$, $A_{3}$.

If one preferred to take as auxiliaries three of the quantities $M_{a}$, one would have to have recourse to the equations (u), which give
\[ \begin{aligned} A_{1}&=M_{1}\cdot\frac{\sin\omega}{\sin 2\omega}\\ \frac{A_{2}A_{3}}{A_{1}A_{1}}&=\frac{M_{2}}{M_{1}}\cdot\frac{\sin^{2}.2\omega}{\sin 3\omega\sin 4\omega}\\ \frac{A_{4}A_{5}}{A_{2}A_{2}}&=\frac{M_{3}}{M_{2}}\cdot\frac{\sin^{2}3\omega}{\sin 5\omega\sin 6\omega}. \end{aligned} \]
By means of these equations and the preceding ones, observing moreover that one has $\omega=\frac{\pi}{12}$, one finds
\[ \begin{aligned} A_{1}&=2\sin\omega\cdot M_{1}\\ A_{2}&=2^{\frac{1}{6}}\sin\omega\cdot M_{2}\sqrt{\left(\frac{M_{1}}{M_{3}\cos 2\omega}\right)}\\ A_{3}&=2^{\frac{1}{3}}\sin\omega\cdot M_{3}\sqrt{\left(\frac{M_{1}}{M_{3}\cos 2\omega}\right)}\\ A_{4}&=2^{\frac{1}{6}}\sin\omega\cdot M_{2}\cdot\frac{1}{\cos 2\omega}\\ A_{5}&=2^{\frac{7}{6}}\sin^{2}\omega\cdot M_{1}. \end{aligned} \]

\subsection*{(19).}

Let us apply here the results that we found in the first part, no.~155, and suppose known the two complete functions of the first species $F^{1}(\sin 45^{\circ})=B$, $F^{1}(\sin 15^{\circ})=C$. If one makes, as in the article cited, $\nu=\sqrt[4]{3}$, one will have

\origpage{239}
\[ \begin{aligned} M_{1}&=\frac{2^{\frac{4}{3}}\nu\cos\omega}{3}B\\ M_{2}&=\frac{2^{-\frac{1}{3}}}{\nu}C\\ M_{3}&=\tfrac{1}{3}B. \end{aligned} \]
By means of these values, we shall find
\[ \begin{aligned} A_{1}&=\frac{2^{\frac{1}{3}}\nu}{3}\cdot B\\ A_{2}&=\sqrt{\left(\frac{\nu\sin\omega}{3}\right)}\cdot C\\ A_{3}&=\frac{\sqrt{(2\nu\sin\omega)}}{3\nu}\cdot B\\ A_{4}&=\frac{2^{\frac{5}{6}}\nu\sin\omega}{3}\cdot C\\ A_{5}&=\frac{2^{\frac{1}{2}}\nu\sin\omega}{3}\cdot B. \end{aligned} \]
Thus from the two data $B$ and $C$ alone, we shall be able to determine all the transcendents designated by $\left(\frac{p}{q}\right)$, in the case of $n=12$. These transcendents, excluding those that are determinable by arcs of a circle, are 60 in number, all different from one another. It is seen moreover that they will be determined rationally in terms of $B$ and $C$ and of the number $\pi$, since they are expressed rationally by the auxiliaries $A_{1}$, $A_{2}$, etc.

If one wishes, for example, to have the value of the transcendent $\left(\frac{5}{9}\right)$, one will first seek by equation (n) its value in terms of $A$, which is
\[ \left(\frac{5}{9}\right)=\tfrac{1}{2}\cdot\frac{A_{1}}{A_{3}A_{4}}\cdot\frac{\sin\omega\sin 2\omega}{\sin 5\omega\sin 4\omega\sin 3\omega}; \]
\ednote{The numerator is printed $\sin\omega\sin 2\omega$; formula (n) gives $\omega\sin\omega\sin 2\omega$, and only with the factor $\omega$ does the equation hold numerically (and lead to the final value below).}
making the substitutions next and observing that $\omega=\frac{\pi}{12}$, there will result
\[ \left(\frac{5}{9}\right)=\frac{\pi}{2C}\sqrt{\left(\frac{\nu\sin\omega}{6}\right)}. \]

\origpage{240}

\section*{Formulas for evaluating the integrals $\left(\frac{p}{q}\right)$ by approximation.}

\subsection*{(20).}

Having reduced to the least possible number the transcendents $\left(\frac{p}{q}\right)$, it remains only to show how one can find by approximation, and in an easy manner, the value of each of these quantities.

For this purpose, let us consider first the formula
\[ \left(\frac{a}{a}\right)=2^{1-\frac{2a}{n}}\int\frac{x^{a-1}\,dx}{\sqrt{(1-x^{n})}}, \]
and let $\sqrt{(1-x^{n})}=1-y$, or $x^{n}=2y-y^{2}$; we shall have as transformed,
\[ \left(\frac{a}{a}\right)=2^{1-\frac{a}{n}}\int\frac{dy}{n}\left(y-\frac{y^{2}}{2}\right)^{\frac{a}{n}-1}; \]
this differential being developed and integrated from $y=0$ to $y=1$, we obtain
\[ \left(\frac{a}{a}\right)=2^{1-\frac{a}{n}}\left\{\begin{aligned} &1+\frac{n-a}{2n}\cdot\frac{a}{n+a}+\frac{n-a\cdot 2n-a}{2n\cdot 4n}\cdot\frac{a}{2n+a}\\ &\qquad+\frac{n-a\cdot 2n-a\cdot 3n-a}{2n\cdot 4n\cdot 6n}\cdot\frac{a}{3n+a}\\ &+\text{etc.} \end{aligned}\right\}:a\qquad(g'), \]
a formula of which each term is less than half of the preceding.

\subsection*{(21).}

In general, if one wishes to have the approximate value of the quantity $\left(\frac{p}{q}\right)=\displaystyle\int\frac{x^{p-1}\,dx}{\sqrt[n]{(1-x^{n})^{n-q}}}$, this integral must be divided into two parts, the one from $x^{n}=0$ to $x^{n}=\frac{1}{2}$, the other from $x^{n}=\frac{1}{2}$ to $x^{n}=1$.

The first part being named $P$, we find by the ordinary developments,
\[ P=2^{-\frac{p}{n}}\left(\frac{1}{p}+\frac{n-q}{2n}\cdot\frac{1}{n+p}+\frac{n-q\cdot 2n-q}{2n\cdot 4n}\cdot\frac{1}{2n+p}+\text{etc.}\right). \]

\origpage{241}
To have the second part, which we shall name $Q$, we must put $x^{n}=1-y^{n}$, and then we shall have
\[ \int\frac{x^{p-1}\,dx}{\sqrt[n]{(1-x^{n})^{n-q}}}=-\int\frac{y^{q-1}\,dy}{\sqrt[n]{(1-x^{n})^{n-q}}}, \]
\ednote{The second denominator is printed $\sqrt[n]{(1-x^{n})^{n-q}}$; the substitution gives $\sqrt[n]{(1-y^{n})^{n-p}}$, as the series for $Q$ below requires. An apparent misprint.}
and the transformed in $y$ must be integrated from $y^{n}=\frac{1}{2}$ to $y^{n}=0$. If its sign is changed, it must be integrated from $y^{n}=0$ to $y^{n}=\frac{1}{2}$; we shall therefore have
\[ Q=2^{-\frac{q}{n}}\left(\frac{1}{q}+\frac{n-p}{2n}\cdot\frac{1}{n+q}+\frac{n-p\cdot 2n-p}{2n\cdot 4n}\cdot\frac{1}{2n+q}+\text{etc.}\right). \]
It remains only to unite these two parts, and we obtain
\[ \left(\frac{p}{q}\right)=\left\{\begin{aligned} &2^{-\frac{p}{n}}\left(\frac{1}{p}+\frac{n-q}{2n}\cdot\frac{1}{n+p}+\frac{n-q\cdot 2n-q}{2n\cdot 4n}\cdot\frac{1}{2n+p}+\frac{n-q\cdot 2n-q\cdot 3n-q}{2n\cdot 4n\cdot 6n}\cdot\frac{1}{3n+p}+\text{etc.}\right)\\ &+2^{-\frac{q}{n}}\left(\frac{1}{q}+\frac{n-p}{2n}\cdot\frac{1}{n+q}+\frac{n-p\cdot 2n-p}{2n\cdot 4n}\cdot\frac{1}{2n+q}+\frac{n-p\cdot 2n-p\cdot 3n-p}{2n\cdot 4n\cdot 6n}\cdot\frac{1}{3n+q}+\text{etc.}\right) \end{aligned}\right\}\qquad(h'). \]
The two series comprised in this formula are always convergent, since each term is less than half of the preceding: we thus obviate the inconvenience that the ordinary method would present, if one wished to integrate the value of $\left(\frac{p}{q}\right)$ all at once from $x=0$ to $x=1$, which would give the very slowly convergent sequence
\[ \left(\frac{p}{q}\right)=\frac{1}{p}+\frac{n-q}{n}\cdot\frac{1}{n+p}+\frac{n-q\cdot 2n-q}{n\cdot 2n}\cdot\frac{1}{2n+p}+\text{etc.}\qquad(i'). \]
Moreover, when one supposes $p=q=a$, formula (h') reduces precisely to formula (g') found by another route.

\subsection*{(22).}

It will not be useless to seek the value of the function $\left(\frac{p}{q}\right)$ in the case where $n$ is very large in relation to the numbers $p$ and $q$. For this, let $p=\alpha n$, $q=\beta n$; we may consider $\alpha$ and $\beta$ as very small quantities of the first order, and the series $P$ and $Q$ must be developed to the suitable degree; we have first,

\origpage{242}
putting $\frac{1}{2}=x$;
\[ P=\frac{2^{-\alpha}}{n\alpha}\left\{\begin{aligned} &1+\frac{1-\beta}{1}\cdot\frac{\alpha x}{1+\alpha}+1-\beta\cdot 1-\tfrac{1}{2}\beta\cdot\frac{\alpha x^{2}}{2+\alpha}\\ &\qquad+1-\beta\cdot 1-\tfrac{1}{2}\beta\cdot 1-\tfrac{1}{3}\beta\cdot\frac{\alpha x^{3}}{3+\alpha}\\ &+\text{etc.} \end{aligned}\right\}. \]

Developing the series enclosed in the parenthesis, up to quantities of the second order inclusive, and putting for brevity,
\[ \begin{aligned} A&=x+\frac{x^{2}}{2}+\frac{x^{3}}{3}+\frac{x^{4}}{4}+\text{etc.}\\ B&=x+\left(1+\tfrac{1}{2}\right)\frac{x^{2}}{2}+\left(1+\tfrac{1}{2}+\tfrac{1}{3}\right)\frac{x^{3}}{3}+\text{etc.}\\ C&=x+\frac{x^{2}}{2^{2}}+\frac{x^{3}}{3^{2}}+\frac{x^{4}}{4^{2}}+\text{etc.} \end{aligned} \]
we shall have
\[ P=\frac{2^{-\alpha}}{n\alpha}\left[1+(A-B\beta)\alpha-C\alpha^{2}\right]. \]
But we also have $2^{-\alpha}=1-\alpha l2+\frac{1}{2}\alpha^{2}(l2)^{2}$, etc.; and moreover $A=-\log(1-x)=-\log\left(1-\frac{1}{2}\right)=\log 2$, so that we may suppose $2^{-\alpha}=1-A\alpha+\frac{1}{2}A^{2}\alpha^{2}$; thus, carrying out the developments, we shall have
\[ P=\frac{1}{n\alpha}\left[1-B\beta\alpha-\left(C+\tfrac{1}{2}A^{2}\right)\alpha^{2}\right]. \]
We shall have similarly
\[ Q=\frac{1}{n\beta}\left[1-B\alpha\beta-\left(C+\tfrac{1}{2}A^{2}\right)\beta^{2}\right]; \]
therefore by the sum of these quantities we find
\[ \left(\frac{p}{q}\right)=\frac{1}{n}\left[\frac{1}{\alpha}+\frac{1}{\beta}-\left(B+C+\tfrac{1}{2}A^{2}\right)(\alpha+\beta)\right], \]
or
\[ \left(\frac{p}{q}\right)=\frac{p+q}{pq}\left[1-\left(B+C+\tfrac{1}{2}A^{2}\right)\frac{pq}{n^{2}}\right]. \]

\subsection*{(23).}

It remains to find the value of $B+C$, and for this we shall consider for a moment $B$ and $C$ as functions of a

\origpage{243}
variable $x$; we shall have first
\[ \begin{aligned} B-C&=\frac{x^{2}}{2}+\left(1+\tfrac{1}{2}\right)\frac{x^{3}}{3}+\left(1+\tfrac{1}{2}+\tfrac{1}{3}\right)\frac{x^{4}}{4}+\text{etc.}\\ \frac{d(B-C)}{dx}&=x+\left(1+\tfrac{1}{2}\right)x^{2}+\left(1+\tfrac{1}{2}+\tfrac{1}{3}\right)x^{3}+\text{etc.}\\ &=\frac{x}{1-x}+\frac{\frac{1}{2}x^{2}}{1-x}+\frac{\frac{1}{3}x^{3}}{1-x}+\text{etc.}=-\frac{\log(1-x)}{1-x}, \end{aligned} \]
therefore $d(B-C)=-\frac{dx}{1-x}\log(1-x)$, and on integrating we have
\[ B-C=\tfrac{1}{2}\log^{2}(1-x), \]
without a constant, because $B$ and $C$ vanish at the same time as $x$.

We have next
\[ \frac{x\,dC}{dx}=x+\frac{x^{2}}{2}+\frac{x^{3}}{3}+\text{etc.}=-\log(1-x); \]
therefore
\[ dC=-\frac{dx}{x}\log(1-x), \]
and on integrating
\[ C=-\log x\log(1-x)-\int\frac{dx}{1-x}\log x. \]
Let $x=1-y$; we shall have
\[ \int\frac{dx}{1-x}\log x=\int\frac{-dy}{y}\log(1-y)=\text{const}+y+\frac{y^{2}}{2^{2}}+\frac{y^{3}}{3^{2}}+\text{etc.} \]
If therefore $C$ is a function of $x$ designated by $\psi(x)$, we shall have
\[ \int\frac{dx}{1-x}\log x=\text{const.}+\psi(y)=\text{const}+\psi(1-x); \]
therefore
\[ \psi(x)+\psi(1-x)=\text{const}-\log x\log(1-x). \]
If we put $x=0$, we find the constant $=\psi(1)=1+\frac{1}{2^{2}}+\frac{1}{3^{2}}+\text{etc.}$, a quantity of which the value is known to be $\frac{\pi^{2}}{6}$, so that we shall have
\[ \psi(x)+\psi(1-x)=\frac{\pi^{2}}{6}-\log x\log(1-x)\qquad(k'). \]

\origpage{244}
This formula shows that, the sum of the sequence
\[ \psi(x)=x+\frac{x^{2}}{2^{2}}+\frac{x^{3}}{3^{2}}+\frac{x^{4}}{4^{2}}+\text{etc.}, \]
being known for every value of $x$ from $x=0$ to $x=\frac{1}{2}$, we shall know the sum of the same sequence for every value of $x$ from $x=\frac{1}{2}$ to $x=1$.

\subsection*{(24).}

In the particular case where one puts $x=\frac{1}{2}$, equation (k') gives
\[ \psi\left(\frac{1}{2}\right)=\frac{\pi^{2}}{12}-\tfrac{1}{2}(\log 2)^{2}; \]
therefore in this same case, we shall have
\[ C=\frac{\pi^{2}}{12}-\tfrac{1}{2}A^{2}; \]
but we have found $B-C=\frac{1}{2}\log^{2}(1-x)=\frac{1}{2}\log^{2}2=\frac{1}{2}A^{2}$; therefore
\[ B=C+\tfrac{1}{2}A^{2}=\frac{\pi^{2}}{12}, \]
and finally
\[ \left(\frac{p}{q}\right)=\frac{p+q}{pq}\left(1-\frac{\pi^{2}}{6}\cdot\frac{pq}{n^{2}}\right)\qquad(l'). \]
This is the limit toward which the function $\left(\frac{p}{q}\right)$ continually tends, as $n$ increases more and more, $p$ and $q$ remaining the same.

It is at the same time an approximate value of $\left(\frac{p}{q}\right)$, when $p$ and $q$ are small in relation to $n$. Let, for example, $p=1$, $q=1$, $n=12$; we shall have very nearly $\left(\frac{1}{1}\right)=2\left(1-\frac{\pi^{2}}{864}\right)$.

\section*{Remark on some particular cases, where the series designated by $\psi(x)$ can be summed, and two others of the same species.}

\subsection*{(25).}

It has been seen in the preceding researches that the series represented by $\psi(x)$ is summable both when $x=1$ and when

\origpage{245}
$x=\frac{1}{2}$. Landen, in his \emph{Mathématical Memoirs}, pag.~112 ff., has indicated two other cases where this same series is summable: we shall show how one arrives at these results.

Having put $\psi(x)=x+\frac{x^{2}}{2^{2}}+\frac{x^{3}}{3^{2}}+\frac{x^{4}}{4^{2}}+\text{etc.}$, we deduce from it
\[ d\psi(x)=dx\left(1+\frac{x}{2}+\frac{x^{2}}{3}+\text{etc.}\right)=-\frac{dx}{x}\log(1-x) \]
and $\psi x)$\ednote{Printed without the opening parenthesis; $\psi(x)$ is meant.} $=\displaystyle\int-\frac{dx}{x}\log(1-x)$. Let $\log x=p$ and $\log(1-x)=q$; we should have at once
\[ \psi(x)=\int-q\,dp,\qquad\psi(1-x)=\int-p\,dq; \]
therefore $\psi(x)+\psi(1-x)=C-pq=\psi(1)-\log x\log(1-x)$.
This is the equation (k') which we have already found; but several others can be found.

Indeed we should have similarly $\psi(-x)=\displaystyle\int-\frac{dx}{x}\log(1+x)$; putting $\frac{x}{1-x}$ in place of $x$, we shall have
\[ \psi\left(-\frac{x}{1-x}\right)=\int\left(\frac{dx}{x}+\frac{dx}{1-x}\right)\log(1-x); \]
therefore
\[ \psi(x)+\psi\left(\frac{-x}{1-x}\right)=\int\frac{dx}{1-x}\log(1-x)=-\tfrac{1}{2}\log^{2}(1-x). \]
No constant is added, because the two members vanish when $x=0$.

Putting again $\frac{x}{1+x}$ in place of $x$, we shall have
\[ \psi\left(\frac{x}{x+1}\right)+\psi(-x)=-\tfrac{1}{2}\log^{2}(1+x). \]
This equation gives an always finite value for $\psi(-x)$, however great $x$ may be, since $\psi(-x)$ depends only on $\psi\left(\frac{x}{x+1}\right)$, where $\frac{x}{x+1}$ is smaller than unity. But as in the case of $x>1$ the series designated by $\psi(x)$ will be divergent, one does not see what its utility can be; let us therefore content ourselves with considering

\origpage{246}
the values of $\psi(-x)$ when $x$ is ${}<1$. Now we have directly
\[ \psi(x)+\psi(-x)=\tfrac{1}{2}\left(x^{2}+\frac{x^{4}}{2^{2}}+\frac{x^{6}}{3^{2}}+\text{etc.}\right)=\tfrac{1}{2}\psi(x^{2}), \]
the preceding equation can therefore be put in the form
\[ \psi\left(\frac{x}{x+1}\right)+\tfrac{1}{2}\psi(x^{2})-\psi(x)=-\tfrac{1}{2}\log^{2}(1+x). \]
Let us give to $x$ a particular value, such that $x^{2}=\frac{x}{x+1}$, or $x^{2}+x=1$; if we put $\beta=-\frac{1}{2}+\frac{1}{2}\sqrt{5}$, we shall have $x=\beta$, $\frac{x}{x+1}=\beta^{2}=1-\beta$. Therefore
\[ \tfrac{3}{2}\psi(1-\beta)-\psi(\beta)=-\tfrac{1}{2}\log^{2}(1+\beta)=-\tfrac{1}{2}\log^{2}\beta. \]
But we have besides, by equation (k'),
\[ \psi(1-\beta)+\psi(\beta)=\frac{\pi^{2}}{6}-2\log^{2}\beta; \]
therefore
\[ \begin{aligned} \psi(\beta)&=\frac{\pi^{2}}{10}-\log^{2}\beta\\ \psi(1-\beta)&=\frac{\pi^{2}}{15}-\log^{2}\beta. \end{aligned} \]
Thus we know the transcendent $\psi(x)$, not only when $x=1$ and $x=\frac{1}{2}$, but also when $x=\beta=-\frac{1}{2}+\frac{1}{2}\sqrt{5}$, and when $x=1-\beta=\frac{3}{2}-\frac{1}{2}\sqrt{5}$.

\subsection*{(26).}

Let us now consider the function
\[ \varphi(x)=x+\frac{x^{3}}{3^{2}}+\frac{x^{5}}{5^{2}}+\text{etc.} \]
We shall have first $\varphi(x)=\psi(x)-\frac{1}{4}\psi(x^{2})$, which is verified immediately by the series expressions of $\psi(x)$ and $\psi(x^{2})$. But we know the values of $\psi(x)$ and $\psi(x^{2})$ when $x=\beta$, since then $x^{2}=1-\beta$; therefore we shall have for this same case,
\[ \varphi(\beta)=\frac{\pi^{2}}{12}-\tfrac{3}{4}\log^{2}\beta. \]
\ednote{Under $\pi^{2}$ only a digit $1$, crossed by a reader's pen stroke, is legible in the scan; read as $12$, the value that $\psi(\beta)-\frac{1}{4}\psi(1-\beta)$ gives.}
To have other values of the function $\varphi$, I observe that we have

\origpage{247}
\[ d\varphi(x)=dx\left(1+\frac{x^{2}}{3}+\frac{x^{4}}{5}+\text{etc.}\right)=\frac{dx}{2x}\log\left(\frac{1+x}{1-x}\right), \]
whence
\[ \varphi(x)=\int\frac{dx}{2x}\log\left(\frac{1+x}{1-x}\right). \]
Let us put $\frac{1-x}{1+x}$ in place of $x$; we shall have similarly
\[ \varphi\left(\frac{1-x}{1+x}\right)=\int\frac{dx}{1-x^{2}}\log x. \]
Adding these two equations and carrying out the indicated integration, we shall have
\[ \varphi(x)+\varphi\left(\frac{1-x}{1+x}\right)=\varphi(1)+\tfrac{1}{2}\log x\log\left(\frac{1+x}{1-x}\right). \]
The constant $\varphi(1)$ has been added, so that the two members may be equal when $x=0$. As for the value of $\varphi(1)$, it is deduced from the equation $\varphi(x)=\psi(x)-\frac{1}{4}\psi(x^{2})$, which gives $\varphi(1)=\frac{3}{4}\psi(1)=\frac{\pi^{2}}{8}$; so that we shall have
\[ \varphi(x)+\varphi\left(\frac{1-x}{1+x}\right)=\frac{\pi^{2}}{8}+\tfrac{1}{2}\log x\log\left(\frac{1+x}{1-x}\right). \]
Let us now give to $x$ a particular value such that $\frac{1-x}{1+x}=x$; there will result $x=\sqrt{2}-1=\gamma$; and from this value, we shall have
\[ \varphi(\gamma)=\frac{\pi^{2}}{8}-\tfrac{1}{2}\log^{2}\gamma. \]
\ednote{As printed. The preceding equation with $x=\gamma$ gives $2\varphi(\gamma)=\frac{\pi^{2}}{8}-\frac{1}{2}\log^{2}\gamma$, that is $\varphi(\gamma)=\frac{\pi^{2}}{16}-\frac{1}{4}\log^{2}\gamma$; the halving appears to have been omitted.}
We therefore know four different values of $\varphi(x)$, as we know four of $\psi(x)$; for the known value $\varphi(\beta)$ makes known another, $\varphi\left(\frac{1-\beta}{1+\beta}\right)$, or $\varphi(2\beta-1)=\varphi(\sqrt{5}-2)$. Thus the values of $x$ for which $\varphi(x)$ is known are $x=1$, $x=\beta=-\frac{1}{2}+\frac{1}{2}\sqrt{5}$, $x=2\beta-1=\sqrt{5}-2$, $x=\gamma=\sqrt{2}-1$.

\subsection*{(27).}

To push these researches still further, let us consider the function
\[ \Lambda(x)=x+\frac{x^{2}}{2^{3}}+\frac{x^{3}}{3^{3}}+\frac{x^{4}}{4^{3}}+\text{etc.}, \]
which in general is a very composite transcendent, since it is a

\origpage{248}
rather difficult problem to determine by a convergent series the value of $\Lambda(1)$. We shall have first by differentiation
\[ d\Lambda(x)=dx\left(1+\frac{x}{2^{2}}+\frac{x^{2}}{3^{2}}+\text{etc.}\right)=\frac{dx}{x}\psi(x); \]
which gives $\Lambda(x)=\displaystyle\int\frac{dx}{x}\psi(x)$, or
\[ \Lambda(x)=\log x\cdot\psi(x)+\int\frac{dx}{x}\log x\log(1-x). \]
Putting as above $\log x=p$, $\log(1-x)=q$, we shall have
\[ \Lambda(x)=\log x\cdot\psi(x)+\int pq\,dp, \]
whence one deduces, putting $1-x$ in place of $x$,
\[ \Lambda(1-x)=\log(1-x)\psi(1-x)+\int pq\,dq. \]
We should have similarly $\Lambda(-x)=\displaystyle\int\frac{dx}{x}\psi(-x)$, or
\[ \Lambda(-x)=\log x\cdot\psi(-x)+\int\frac{dx}{x}\log x\log(1+x). \]
Let us put in this equation $\frac{x}{1-x}$ in place of $x$; there will result
\[ \Lambda\left(-\frac{x}{1-x}\right)=\log\frac{x}{1-x}\cdot\psi\left(-\frac{x}{1-x}\right)+\int q(p-q)(dp-dq). \]
Adding these three equations, and carrying out the indicated integrations, we shall have the general formula
\[ \left.\begin{aligned} &\Lambda(x)+\Lambda(1-x)\\ &+\Lambda\left(-\frac{x}{1-x}\right) \end{aligned}\right\}=\left\{\begin{aligned} &\Lambda(1)+\log x\cdot\psi(x)+\log(1-x)\psi(1-x)\\ &+\log\frac{x}{1-x}\psi\left(-\frac{x}{1-x}\right)+\log x\log^{2}(1-x)-\tfrac{1}{3}\log^{3}(1-x). \end{aligned}\right. \]
Let $x=\frac{1}{2}$; the first member will become $2\Lambda\left(\frac{1}{2}\right)+\Lambda(-1)$; and as we have in general $\Lambda(x)+\Lambda(-x)=\frac{1}{4}\Lambda(x^{2})$, which gives $\Lambda(-1)=-\frac{3}{4}\Lambda(1)$, we shall find
\[ 2\Lambda\left(\tfrac{1}{2}\right)-\tfrac{3}{4}\Lambda(1)=\Lambda(1)+2\log\tfrac{1}{2}\cdot\psi\left(\tfrac{1}{2}\right)+\tfrac{2}{3}\log^{3}\left(\tfrac{1}{2}\right). \]
Substituting the known value of $\psi\left(\frac{1}{2}\right)$, there will result
\[ \tfrac{7}{8}\Lambda(1)=\Lambda\left(\tfrac{1}{2}\right)+\frac{\pi^{2}}{12}\log(2)-\tfrac{1}{6}(\log 2)^{3}; \]

\origpage{249}
whence it is seen that $\Lambda\left(\frac{1}{2}\right)$ can be determined by means of $\Lambda(1)$; and conversely $\Lambda(1)$ can be determined by means of $\Lambda\left(\frac{1}{2}\right)$, which is a first manner of expressing the transcendent $\Lambda(1)$ by a convergent series.

In the second place, if we put $x=1-\beta=\beta^{2}$, $\beta$ having the same value as above, the general formula will give
\[ \begin{aligned} \Lambda(\beta)+\Lambda(\beta^{2})+\Lambda(-\beta)&=\log\beta\left[2\psi(\beta^{2})+\psi(\beta)+\psi(-\beta)\right]\\ &\quad+\Lambda(1)+\tfrac{5}{3}\log^{3}\beta. \end{aligned} \]
But we have $\psi(\beta)+\psi(-\beta)=\frac{1}{2}\psi(\beta^{2})$ and $\Lambda(\beta)+\Lambda(-\beta)=\frac{1}{4}\Lambda(+\beta^{2})$; therefore
\[ \tfrac{5}{4}\Lambda(\beta^{2})=\tfrac{5}{2}\log\beta\,\psi(\beta^{2})+\Lambda(1)+\tfrac{5}{3}\log^{3}\beta; \]
now we have found $\psi(\beta^{2})=\frac{\pi^{2}}{15}-\log^{2}\beta$; therefore finally
\[ \Lambda(1)=\tfrac{5}{4}\Lambda(\beta^{2})-\frac{\pi^{2}}{6}\log\beta+\tfrac{5}{6}\log^{3}\beta. \]
Thus the quantity $\Lambda(1)$, which represents the series $1+\frac{1}{2^{3}}+\frac{1}{3^{3}}+\frac{1}{4^{3}}+\text{etc.}$, can be determined by means of $\Lambda(\beta^{2})$, that is to say by the convergent series $\beta^{2}+\frac{\beta^{4}}{2^{3}}+\frac{\beta^{6}}{3^{3}}+\frac{\beta^{8}}{4^{3}}+\text{etc.}$, in which $\beta^{2}=\frac{1}{2}(3-\sqrt{5})$, a quantity smaller than 0,4. These formulas were given without demonstration in the cited work of Landen, pag.~118; and up to the present no one has gone further in the theory of these sorts of transcendents.

\section*{Consideration of the integral formulas $\displaystyle\int\frac{x^{p-1}\,dx\log\frac{1}{x}}{\sqrt[n]{(1-x^{n})^{n-q}}}$, $\displaystyle\int\frac{x^{p-1}\,dx\log^{2}\frac{1}{x}}{\sqrt[n]{(1-x^{n})^{n-q}}}$, etc. A very remarkable theorem on the first of these formulas.}

\subsection*{(28).}

If we designate by $Z$ the function $\left(\frac{p}{q}\right)$ or the integral $\displaystyle\int\frac{x^{p-1}\,dx}{\sqrt[n]{(1-x^{n})^{n-q}}}$, taken from $x=0$ to $x=1$, the successive differentials of $Z$, taken by varying $p$ alone, will give the

\origpage{250}
new formulas
\[ \begin{aligned} \frac{dZ}{dp}&=-\int\frac{x^{p-1}\,dx\log\frac{1}{x}}{\sqrt[n]{(1-x^{n})^{n-q}}}\\ \frac{ddZ}{dp^{2}}&=\int\frac{x^{p-1}\,dx\log^{2}\frac{1}{x}}{\sqrt[n]{(1-x^{n})^{n-q}}}\\ \frac{d^{3}Z}{dp^{3}}&=-\int\frac{x^{p-1}\,dx\log^{3}\frac{1}{x}}{\sqrt[n]{(1-x^{n})^{n-q}}}\\ &\text{etc.}, \end{aligned} \]
these integrals being again taken from $x=0$ to $x=1$.

As we have generally by formula (i'),
\[ Z=\frac{1}{p}+\frac{n-q}{n}\cdot\frac{1}{n+p}+\frac{n-q\cdot 2n-q}{n\cdot 2n}\cdot\frac{1}{2n+p}+\frac{n-q\cdot 2n-q\cdot 3n-q}{n\cdot 2n\cdot 3n}\cdot\frac{1}{3n+p}+\text{etc.} \]
we deduce from it by successive differentiations,
\[ \left.\begin{aligned} \int\frac{x^{p-1}\,dx\log\frac{1}{x}}{\sqrt[n]{(1-x^{n})^{n-q}}}&=\frac{1}{p^{2}}+\frac{n-q}{n}\cdot\frac{1}{(n+p)^{2}}+\frac{n-q\cdot 2n-q}{n\cdot 2n}\cdot\frac{1}{(2n+p)^{2}}+\text{etc.}\\ \frac{1}{2}\int\frac{x^{p-1}\,dx\log^{2}\frac{1}{x}}{\sqrt[n]{(1-x^{n})^{n-q}}}&=\frac{1}{p^{3}}+\frac{n-q}{n}\cdot\frac{1}{(n+p)^{3}}+\frac{n-q\cdot 2n-q}{n\cdot 2n}\cdot\frac{1}{(2n+p)^{3}}+\text{etc.}\\ \frac{1}{2\cdot 3}\int\frac{x^{p-1}\,dx\log^{3}\frac{1}{x}}{\sqrt[n]{(1-x^{n})^{n-q}}}&=\frac{1}{p^{4}}+\frac{n-q}{n}\cdot\frac{1}{(n+p)^{4}}+\frac{n-q\cdot 2n-q}{n\cdot 2n}\cdot\frac{1}{(2n+p)^{4}}+\text{etc.}\\ &\text{etc.} \end{aligned}\right\}\qquad(m'). \]
But these series have the inconvenience of not being sufficiently convergent.

\subsection*{(29).}

To have these same values expressed in more convergent series, we should have to start from formula (h'), and differentiate it with respect to $p$, as many times as is necessary.

Differentiating it once, we shall have

\origpage{251}
\[ \begin{aligned} \int\frac{x^{p-1}\,dx\log\frac{1}{x}}{\sqrt[n]{(1-x^{n})^{n-q}}}={}&\frac{2^{-\frac{p}{n}}\log 2}{n}\left[\frac{1}{p}+\frac{n-q}{2n}\cdot\frac{1}{n+p}+\frac{n-q\cdot 2n-q}{2n\cdot 4n}\cdot\frac{1}{2n+p}+\text{etc.}\right]\\ &+2^{-\frac{p}{n}}\left[\frac{1}{p^{2}}+\frac{n-q}{2n}\cdot\frac{1}{(n+p)^{2}}+\frac{n-q\cdot 2n-q}{2n\cdot 4n}\cdot\frac{1}{(2n+p)^{2}}+\text{etc}\right]\\ &+2^{-\frac{q}{n}}\left\{\begin{aligned} &\frac{1}{2n}\cdot\frac{1}{n+q}+\frac{n-p\cdot 2n-p}{2n\cdot 4n(2n+q)}\left(\frac{1}{n-p}+\frac{1}{2n-p}\right)\\ &+\frac{n-p\cdot 2n-p\cdot 3n-p}{2n\cdot 4n\cdot 6n\cdot(3n+q)}\left(\frac{1}{n-p}+\frac{1}{2n-p}+\frac{1}{3n-p}\right)\\ &+\text{etc.} \end{aligned}\right\} \end{aligned} \]
These series are convergent, since each term is less than half of the preceding term; but their form is complicated, and it would become more so in the differential of the second order or of a higher order. This is why it is fitting to have recourse to other means if one wishes to evaluate easily the integrals in question.

\subsection*{(30).}

Let us designate by $\varphi\left(\frac{p}{q}\right)$ the integral $\displaystyle\int\frac{x^{p-1}\,dx\log\frac{1}{x}}{\sqrt[n]{(1-x^{n})^{n-q}}}$, and by $\psi\left(\frac{p}{q}\right)$ the ratio of this integral to the function $\left(\frac{p}{q}\right)$, already represented by $Z$, so that we have
\[ \psi\left(\frac{p}{q}\right)=\frac{\varphi\left(\frac{p}{q}\right)}{Z}; \]
according to what was said (art.~28), we shall have
\[ \begin{aligned} \varphi\left(\frac{p}{q}\right)&=-\frac{dZ}{dp}\\ \text{and}\quad\psi\left(\frac{p}{q}\right)&=-\frac{dZ}{dp}\cdot\frac{1}{Z}=-\frac{d(\log Z)}{dp}. \end{aligned} \]
But since we also have $Z=\left(\frac{q}{p}\right)$, it results\ednote{The word ``en'' is printed twice in the French; an apparent misprint.}
\[ \frac{dZ}{dq}=-\varphi\left(\frac{q}{p}\right), \]
where it must be observed that $\varphi\left(\frac{q}{p}\right)$ is not the same thing as $\varphi\left(\frac{p}{q}\right)$.

\origpage{252}

The complete differential of $Z$ or of $\left(\frac{p}{q}\right)$ may therefore be expressed thus:
\[ d\left(\frac{p}{q}\right)=-dp\,\varphi\left(\frac{p}{q}\right)-dq\,\varphi\left(\frac{q}{p}\right)\qquad(n'). \]

\subsection*{(31).}

This being laid down, if we take the general equation
\[ \left(\frac{p}{q}\right)\cdot\left(\frac{p+q}{r}\right)=\left(\frac{p}{r}\right)\cdot\left(\frac{p+r}{q}\right), \]
which can be put in the form
\[ \log\left(\frac{p}{q}\right)+\log\left(\frac{p+q}{r}\right)=\log\left(\frac{p}{r}\right)+\log\left(\frac{p+r}{q}\right), \]
and if we differentiate it with respect to $p$, we shall have
\[ \frac{\varphi\left(\frac{p}{q}\right)}{\left(\frac{p}{q}\right)}+\frac{\varphi\left(\frac{p+q}{r}\right)}{\left(\frac{p+q}{r}\right)}=\frac{\varphi\left(\frac{p}{r}\right)}{\left(\frac{p}{r}\right)}+\frac{\varphi\left(\frac{p+r}{q}\right)}{\left(\frac{p+r}{q}\right)}, \]
or according to the denominations established,
\[ \psi\left(\frac{p}{q}\right)+\psi\left(\frac{p+q}{r}\right)=\psi\left(\frac{p}{r}\right)+\psi\left(\frac{p+r}{q}\right)\qquad(p'). \]
The same equation, differentiated with respect to $q$, gives
\[ \psi\left(\frac{q}{p}\right)+\psi\left(\frac{p+q}{r}\right)=\psi\left(\frac{q}{p+r}\right)\qquad(q'). \]
The differential with respect to $r$ would give a result of the same form as the preceding, and which would consequently be included in it. It can moreover be shown that the equation of four terms (p') is included in equation (q'); for from the latter we deduce, by the permutation of the letters $p$ and $q$,
\[ \psi\left(\frac{p}{q}\right)+\psi\left(\frac{p+q}{r}\right)=\psi\left(\frac{p}{q+r}\right), \]
and from this last we conclude, by the exchange of the letters $q$ and $r$,
\[ \psi\left(\frac{p}{r}\right)+\psi\left(\frac{p+r}{q}\right)=\psi\left(\frac{p}{q+r}\right), \]

\origpage{253}
therefore
\[ \psi\left(\frac{p}{q}\right)+\psi\left(\frac{p+q}{r}\right)=\psi\left(\frac{p}{r}\right)+\psi\left(\frac{p+r}{q}\right), \]
which is equation (p').

\subsection*{(32).}

From this it is seen that it suffices to consider the equation of three terms (q'), and it is from this source that we shall draw all the relations that exist between the various quantities $\psi\left(\frac{p}{q}\right)$ which correspond to one same value of $n$.

Among the quantities $\left(\frac{p}{q}\right)$, one distinguishes those of the form $\left(\frac{a}{n}\right)$, of which the value is $\left(\frac{a}{n}\right)=\frac{1}{a}$; it follows that $\varphi\left(\frac{a}{n}\right)=\frac{1}{a^{2}}$, and consequently
\[ \psi\left(\frac{a}{n}\right)=\frac{1}{a}\qquad(r'), \]
a first formula which will serve for the reduction of the others.

Among the same quantities, one finds in the second place the remarkable formula
\[ \left(\frac{a}{n-a}\right)=\frac{\omega}{\sin a\omega}, \]
whence one deduces, taking the differential of each member with respect to $a$,
\[ -\varphi\left(\frac{a}{n-a}\right)+\varphi\left(\frac{n-a}{a}\right)=-\frac{\omega^{2}\cos a\omega}{\sin^{2}a\omega}. \]
This equation, divided by $\left(\frac{a}{n-a}\right)$, gives
\[ \psi\left(\frac{a}{n-a}\right)-\psi\left(\frac{n-a}{a}\right)=\omega\cot a\omega. \]

The function $\psi\left(\frac{a}{n-a}\right)$ is remarkable in the order of the functions $\psi$, as the quantity $\left(\frac{a}{n-a}\right)$ is among those of its species. We shall therefore make use of this function to express all the others, and we shall put, for brevity,
\[ \psi\left(\frac{a}{n-a}\right)=B_{a}.\qquad(s') \]

\origpage{254}
This being laid down, the equation which we have just found will be expressed thus:
\[ B_{a}-B_{n-a}=\omega\cot a\omega.\qquad(t') \]
It shows that the value of $B_{a}$ is concluded from that of $B_{n-a}$, and conversely; whence it follows that among the quantities $B_{1}$, $B_{2}$, $B_{3}$, etc., it suffices to know the first ones up to that whose rank is $\frac{n}{2}$ or $\frac{n-1}{2}$ inclusive.

\subsection*{(33).}

From equation (q'), we deduce generally
\[ \psi\left(\frac{a}{b}\right)-\psi\left(\frac{b}{a}\right)=\psi\left(\frac{a}{b-a}\right),\qquad(u') \]
therefore in particular,
\[ \psi\left(\frac{a}{n}\right)-\psi\left(\frac{n}{a}\right)=\psi\left(\frac{a}{n-a}\right)=B_{a}, \]
and consequently
\[ \psi\left(\frac{n}{a}\right)=\frac{a}{1}-B_{a}.\qquad(v'). \]
\ednote{Printed $\frac{a}{1}$ instead of $\frac{1}{a}$; an apparent misprint, since $\psi\left(\frac{a}{n}\right)=\frac{1}{a}$ by (r'), and (v') is later quoted as $\psi\left(\frac{n}{r}\right)=\frac{1}{r}-B_{r}$.}
Since we have $\left(\frac{p+n}{q}\right)=\frac{p}{p+q}\cdot\left(\frac{p}{q}\right)$, this equation, differentiated logarithmically with respect to $p$, gives
\[ \psi\left(\frac{q}{p+n}\right)=\psi\left(\frac{p}{q}\right)-\frac{1}{p}+\frac{1}{p+q};\qquad(x') \]
\ednote{The left side is printed $\psi\left(\frac{q}{p+n}\right)$, as in (y'); differentiating $\left(\frac{p+n}{q}\right)$ logarithmically with respect to $p$ gives $\psi\left(\frac{p+n}{q}\right)$. An apparent misprint.}
and the same equation, differentiated with respect to $q$, will give
\[ \psi\left(\frac{q}{p+n}\right)=\psi\left(\frac{q}{p}\right)+\frac{1}{p+q}.\qquad(y') \]
These two equations will serve, as need arises, to transform every function $\psi\left(\frac{a}{b}\right)$ into a similar function, in which $a$ and $b$ would be smaller than $n$: we deduce from them, for example,
\[ \psi\left(\frac{n}{n+c}\right)=\psi\left(\frac{n}{c}\right)+\frac{1}{n+c}. \]
But we have $B_{c}=\frac{1}{c}-\psi\left(\frac{n}{c}\right)$ and $B_{n+c}=\frac{1}{n+c}-\psi\left(\frac{n}{n+c}\right)$; therefore
\[ B_{n+c}=B_{c}-\frac{1}{c},\qquad(z') \]

\origpage{255}
a formula which will serve to find the value of $B_{a}$, if $a$ is greater than $n$. Besides, it is easy to see that we have $B_{n}=0$ and $B_{0}=\infty$.

Let us return to equation (q') and put $p+q+r=n$, so that we have at once $\psi\left(\frac{q}{p+r}\right)=B_{q}$ and $\psi\left(\frac{p+q}{r}\right)=B_{n+q}=B_{n-r}$\ednote{Printed $B_{n+q}$; $B_{p+q}$ is expected, since $\frac{p+q}{r}=\frac{p+q}{n-(p+q)}$. An apparent misprint.}, we shall have
\[ \psi\left(\frac{q}{n-r-q}\right)=B_{q}-B_{n-r},\qquad(a'') \]
this is the value of every function $\psi\left(\frac{a}{b}\right)$, in which we have $a+b<n$.

If in the same equation (q') we put $p=n-q$, we shall have
\[ \psi\left(\frac{q}{n-q+r}\right)=\psi\left(\frac{q}{n-q}\right)+\psi\left(\frac{n}{r}\right). \]
But by equation (v') we have $\psi\left(\frac{n}{r}\right)=\frac{1}{r}-B_{r}$; therefore
\[ \psi\left(\frac{q}{n+r-q}\right)=B_{q}-B_{r}+\frac{1}{r};\qquad(b'') \]
this is the value of every function $\psi\left(\frac{a}{b}\right)$, in which we have $a+b>n$.

\subsection*{(34).}

We shall therefore be able to determine generally the value of every function $\psi\left(\frac{a}{b}\right)$, if we determine that of the auxiliary $B_{a}$, or that of the function $\psi\left(\frac{n}{a}\right)$, since we have $B_{a}=\frac{1}{a}-\psi\left(\frac{n}{a}\right)$.

Now we have $\left(\frac{n}{a}\right)=\frac{1}{a}$, and $\varphi\left(\frac{n}{a}\right)=\displaystyle\int\frac{x^{n-1}\,dx\log\frac{1}{x}}{\sqrt[n]{(1-x^{n})^{n-a}}}$; therefore everything reduces to finding the integral
\[ \psi\left(\frac{n}{a}\right)=a\int\frac{x^{n-1}\,dx\log\frac{1}{x}}{\sqrt[n]{(1-x^{n})^{n-a}}}, \]
taken in the ordinary way from $x=0$ to $x=1$.

For this, let $x^{n}=1-y^{n}$; the preceding integral will have for

\origpage{256}
transformed
\[ \begin{aligned} \frac{a}{n}\int-y^{a-1}\,dy\log(1-y^{n})&=-\frac{y^{a}}{n}\log(1-y^{n})-\int\frac{y^{a+n-1}\,dy}{1-y^{n}}\\ &=\frac{y^{a}}{a}-\frac{y^{a}}{n}\log(1-y^{n})-\int\frac{y^{a-1}\,dy}{1-y^{n}}. \end{aligned} \]
On the other hand, $y^{a}\log(1-y^{n})$ can be put in the form $(y^{a}-1)\log(1-y^{n})-\displaystyle\int\frac{ny^{n-1}\,dy}{1-y^{n}}$; we shall therefore have
\[ \psi\left(\frac{n}{a}\right)=\frac{y^{a}}{a}+\frac{1-y^{a}}{n}\log(1-y^{n})-\int\left(\frac{y^{a-1}-y^{n-1}}{1-y^{n}}\right)dy. \]
The part outside the sign vanishes on putting $y=0$, and it reduces to $\frac{1}{a}$ on putting $y=1$; therefore
\[ \psi\left(\frac{n}{a}\right)=\frac{1}{a}-\int\left(\frac{y^{a-1}-y^{n-1}}{1-y^{n}}\right)dy; \]
therefore in general,
\[ B_{a}=\int\left(\frac{y^{a-1}-y^{n-1}}{1-y^{n}}\right)dy,\qquad(c'') \]
this integral being taken from $y=0$ to $y=1$.

\subsection*{(35).}

This being laid down, always putting $\omega=\frac{\pi}{n}$, we find by the known formulas for the integration of rational fractions:
\[ \begin{aligned} B_{a}={}&\frac{1}{n}\log n+\frac{n-2}{n}\,\omega\sin 2a\omega-\frac{2}{n}\cos 2a\omega\log(2\sin\omega)\\ &+\frac{n-4}{n}\,\omega\sin 4a\omega-\frac{2}{n}\cos 4a\omega\log(2\sin 2\omega)\\ &+\frac{n-6}{n}\,\omega\sin 6a\omega-\frac{2}{n}\cos 6a\omega\log(2\sin 3\omega)\\ &+\text{etc.} \end{aligned} \]
This series must be continued to the term
\[ \frac{1}{n}\,\omega\sin(n-1)a\omega-\frac{2}{n}\cos(n-1)a\omega\log\left(2\sin\frac{n-1}{2}\,\omega\right), \]
if $n$ is odd; but if $n$ is even, only half of the last term must be taken, which will be $-\frac{1}{n}\cos a\pi\log 2$.

\origpage{257}

If we then recall the formulas
\[ \begin{aligned} \sin x+\sin 2x+\sin 3x\ldots+\sin mx&=\frac{\sin x+\sin mx-\sin(m+1)x}{2(1-\cos x)},\\ \sin x+2\sin 2x+3\sin 3x\ldots+m\sin mx&=\frac{(m+1)\sin mx-m\sin(m+1)x}{2(1-\cos x)}, \end{aligned} \]
we shall find that the series
\[ \frac{n-2}{n}\,\omega\sin 2a\omega+\frac{n-4}{n}\,\omega\sin 4a\omega+\frac{n-6}{n}\,\omega\sin 6a\omega+\text{etc.}, \]
continued to a number of terms $\frac{n-1}{2}$ or $\frac{n-2}{2}$, has for its sum $\frac{1}{2}\omega\cot a\omega$. Therefore if $n$ is odd, we shall have
\[ \left.\begin{aligned} B_{a}={}&\frac{1}{n}\log n+\tfrac{1}{2}\omega\cot a\omega-\frac{2}{n}\cos 2a\omega\log(2\sin\omega)\\ &-\frac{2}{n}\cos 4a\omega\log(2\sin 2\omega)\\ &-\frac{2}{n}\cos 6a\omega\log(2\sin 3\omega)\\ &\vdots\\ &-\frac{2}{n}\cos(n-1)a\omega\log\left(2\sin\frac{n-1}{2}\,\omega\right),\\ &\text{and if } n \text{ is even, we shall have}\\ B_{a}={}&\frac{1}{n}\log n+\tfrac{1}{2}\omega\cot a\omega-\frac{1}{n}\cos a\pi\log 2\\ &-\frac{2}{n}\cos 2a\omega\log(2\sin\omega)\\ &-\frac{2}{n}\cos 4a\omega\log(2\sin 2\omega)\\ &\vdots\\ &-\frac{2}{n}\cos(n-2)a\omega\log\left(2\sin\frac{n-2}{n}\,\omega\right). \end{aligned}\right\}\qquad(d'') \]
\ednote{The last row is printed $\log\left(2\sin\frac{n-2}{n}\,\omega\right)$; the general term $\log(2\sin k\omega)$ with $k=\frac{n-2}{2}$ requires $\frac{n-2}{2}$. An apparent misprint.}

\subsection*{(36).}

In the particular case where $a=\frac{1}{2}n$, $n$ being even, we find directly by the formula (c''),
\[ B_{a}=\int\frac{y^{a-1}\,dy}{1+y^{a}}=\frac{1}{a}\log(1+y^{a})=\frac{2}{n}\log 2. \]
In order that this value agree with that which the formula (d'') gives in the same

\origpage{258}
case, we must have
\[ \begin{aligned} &\log(2\sin\omega)-\log(2\sin 2\omega)+\log(2\sin 3\omega)\ldots\\ &\quad\pm\log\left[2\sin\left(\frac{n}{2}-1\right)\omega\right]=\left(1+\tfrac{1}{2}\cos a\pi\right)\log 2-\tfrac{1}{2}\log n; \end{aligned} \]
or, what comes to the same thing, passing to the numbers, and supposing either $n=4i$ or $n=4i+2$,
\[ \frac{\sin\omega\sin 3\omega\sin 5\omega\ldots\sin(2i-1)\omega}{\sin 2\omega\sin 4\omega\sin 6\omega\ldots\sin(2i\omega)}=\sqrt{\left(\frac{2}{n}\right)}. \]
This formula is easy to verify by means of the value of $\sin nz$ given by Euler in his \emph{Introd. in anal.}, pag.~204; for by making successively in this value, $z$ infinitely small, and $nz=\frac{1}{2}\pi$, we derive from it
\[ \begin{gathered} \sin\omega\sin 2\omega\sin 3\omega\ldots\sin\frac{n}{2}\,\omega=2^{\frac{1-n}{2}}\sqrt{n},\\ \sin\omega\sin 3\omega\sin 5\omega\ldots\sin(2i-1)\omega=2^{\frac{2-n}{4}}; \end{gathered} \]
$2i$ being $\frac{n}{2}$ or $\frac{n-2}{2}$, according as $n$ is of the form $4i$ or $4i+2$, and from this the preceding equation results.

\subsection*{(37).}

Having the general expression of $B_{a}$, we deduce from it that of $\psi\left(\frac{p}{q}\right)$, by means of one or the other of the formulas
\[ \left.\begin{aligned} \psi\left(\frac{p}{q}\right)&=B_{p}-B_{p+q}\\ \psi\left(\frac{p}{q}\right)&=B_{p}-B_{p+q-n}+\frac{1}{p+q-n}, \end{aligned}\right\}\qquad(e'') \]
the first holding when $p+q$ is $<n$, and the second when $p+q$ is $>n$. In case of equality, we have simply $\psi\left(\frac{p}{q}\right)=B_{p}$.

These values, which are deduced from the formulas (a'') and (b''), can also be put under a single general form, which is
\[ \psi\left(\frac{p}{q}\right)=\int\left(\frac{y^{p-1}-y^{p+q-1}}{1-y^{n}}\right)dy,\qquad(f'') \]
which proves immediately that the function $\psi\left(\frac{p}{q}\right)$ is always

\origpage{259}
determinable by arcs of a circle and logarithms. We thus arrive at this theorem, very remarkable for its elegance and its generality.

\emph{If we take the three integrals}
\[ \int\frac{x^{p-1}\,dx\log\frac{1}{x}}{\sqrt[n]{(1-x^{n})^{n-q}}},\qquad\int\frac{x^{p-1}\,dx}{\sqrt[n]{(1-x^{n})^{n-q}}},\qquad\int\frac{x^{p-1}-x^{p+q-1}}{1-x^{n}}\,dx, \]
\emph{between the limits} $x=0$, $x=1$, \emph{the first will be equal to the product of the other two.}

\subsection*{(38).}

We find in volume IV of Euler's \emph{Calcul intégral}, pag.~166, a demonstration of the same theorem which, on account of its great simplicity, ought to have been substituted for the preceding one, had there not been some advantage in considering one and the same object from different points of view. Here is this demonstration${}^{(1)}$.

By the formula of art.~4, we have
\[ Z=\frac{p+q}{p}\cdot\frac{p+q+n}{p+n}\cdot\frac{p+q+2n}{p+2n}\cdot\text{etc.} \]
Differentiating both members logarithmically with respect to $p$, there will result
\[ \frac{dZ}{Z}=\left(\frac{dp}{p+q}-\frac{dp}{p}\right)+\left(\frac{dp}{p+q+n}-\frac{dp}{p+n}\right)+\left(\frac{dp}{p+q+2n}-\frac{dp}{p+2n}\right)+\text{etc.} \]
Let $\Phi$ or $\Phi(\nu)$ be a function of $\nu$, such that we have
\[ \Phi=\frac{\nu^{p}}{p}-\frac{\nu^{p+q}}{p+q}+\frac{\nu^{p+n}}{p+n}-\frac{\nu^{p+q+n}}{p+q+n}+\frac{\nu^{p+2n}}{p+2n}-\frac{\nu^{p+q+2n}}{p+q+2n}+\text{etc.} \]
The function $\Phi(\nu)$ becoming $\Phi(1)$ when $\nu=1$, we shall have
\[ \frac{dZ}{Z}=-dp\,\Phi(1). \]

\textbf{(1)} Euler had at first given only a particular case of this theorem, which is found in volume I of his \emph{Opusc. anal.}, pag.~183; he afterwards arrived at the general theorem, of which he gave a demonstration that we report here; but this circumstance had escaped me when I published my \emph{Mémoire sur les intégrales définies}, in which I cited only the particular case of the \emph{Opusc. anal.}

\origpage{260}
But $\frac{dZ}{dp}=-\varphi\left(\frac{p}{q}\right)$, therefore $\frac{1}{Z}\varphi\left(\frac{p}{q}\right)$ or $\psi\left(\frac{p}{q}\right)=\Phi(1)$. Thus, to have the value of the function $\psi\left(\frac{p}{q}\right)$, it is only necessary to find that of $\Phi(1)$.

Now, differentiating the value of $\Phi$ with respect to $\nu$, we have
\[ d\Phi=d\nu\left\{\begin{aligned} &\nu^{p-1}+\nu^{n+p-1}+\nu^{2n+p-1}+\text{etc.}\\ &-\nu^{p+q-1}-\nu^{n+p+q-1}-\nu^{2n+p+q-1}-\text{etc.} \end{aligned}\right\} \]
or, summing these series,
\[ d\Phi=d\nu\left(\frac{\nu^{p-1}-\nu^{p+q-1}}{1-\nu^{n}}\right); \]
therefore
\[ \Phi=\int\left(\frac{\nu^{p-1}-\nu^{p+q-1}}{1-\nu^{n}}\right)d\nu. \]
This integral must be taken in such a way that it vanishes when $\nu=0$; we shall then make $\nu=1$, and we shall have the sought value of $\Phi(1)$. Therefore, since we may put $x$ in place of $\nu$, we shall have in general
\[ \psi\left(\frac{p}{q}\right)=\int\left(\frac{x^{p-1}-x^{p+q-1}}{1-x^{n}}\right)dx, \]
the integral of the second member being taken between the limits $x=0$, $x=1$. From this the theorem of the preceding article results.

\subsection*{(39).}

To extend this theory still further, let us consider the two series of integrals in $Z$ and in $T$, taken from $x=0$ to $x=1$, namely:
\[ \begin{aligned} Z&=\int\frac{x^{p-1}\,dx}{\sqrt[n]{(1-x^{n})^{n-q}}}, &\qquad& T=\int\frac{x^{p-1}-x^{p+q-1}}{1-x^{n}}\,dx\\ Z'&=\int\frac{x^{p-1}\,dx\log\frac{1}{x}}{\sqrt[n]{(1-x^{n})^{n-q}}}, && T'=\int\frac{x^{p-1}-x^{p+q-1}}{1-x^{n}}\,dx\log\frac{1}{x}\\ Z''&=\int\frac{x^{p-1}\,dx\log^{2}\frac{1}{x}}{\sqrt[n]{(1-x^{n})^{n-q}}}, && T''=\int\frac{x^{p-1}-x^{p+q-1}}{1-x^{n}}\,dx\log^{2}\frac{1}{x}\\ &\text{etc.} && \text{etc.} \end{aligned} \]

\origpage{261}
It follows first from the preceding theorem that we have
\[ Z'=ZT. \]
Differentiating this equation with respect to $p$, and observing that we have $\frac{dZ}{dp}=-Z'$, $\frac{dZ'}{dp}=-Z''$, $\frac{dT}{dp}=-T'$, we shall have $Z''=Z'T+ZT'$, or
\[ Z''=Z(T^{2}+T'). \]
This being differentiated anew with respect to $p$, gives
\[ Z'''=Z(T^{3}+3TT'+T''), \]
and so on, the law of these expressions being analogous to that of the successive differentials of the formula $ue^{\displaystyle\int u\,dp}$.

If then we wish to have the values of the quantities $Z'$, $Z''$, $Z'''$, etc., or simply their ratio to the primitive function $Z$, it will be necessary to know the quantities $T$, $T'$, $T''$, etc. in equal number. But as these are rational, and contain lower powers of $\log\frac{1}{x}$, we see that the difficulty is at least diminished.

\subsection*{(40).}

If we integrate the differential $x^{m+\alpha-1}\,dx$, from $x=0$ to $x=1$, we shall have
\[ \int x^{m+\alpha-1}\,dx=\frac{1}{m+\alpha}=\frac{1}{m}-\frac{\alpha}{m^{2}}+\frac{\alpha^{2}}{m^{3}}-\frac{\alpha^{3}}{m^{4}}+\text{etc.} \]
If we put the same differential under the form $x^{m-1}\,dx\cdot x^{\alpha}$, or
\[ x^{m-1}\,dx\left(1+\alpha\log x+\frac{\alpha^{2}}{1\cdot 2}\log^{2}x+\frac{\alpha^{3}}{1\cdot 2\cdot 3}\log^{3}x+\text{etc.}\right), \]
its integral, taken between the same limits, will be
\[ \int x^{m-1}\,dx+\alpha\int x^{m-1}\,dx\log x+\frac{\alpha^{2}}{1\cdot 2}\int x^{m-1}\,dx\log^{2}x+\text{etc.} \]
The identity of these two formulas therefore requires that we have
\[ \begin{aligned} \int x^{m-1}\,dx&=\frac{1}{m}\\ \int x^{m-1}\,dx\log\left(\frac{1}{x}\right)&=\frac{1}{m^{2}}\\ \int x^{m-1}\,dx\log^{2}\left(\frac{1}{x}\right)&=\frac{1\cdot 2}{m^{3}}\\ \int x^{m-1}\,dx\log^{3}\left(\frac{1}{x}\right)&=\frac{1\cdot 2\cdot 3}{m^{4}}, \end{aligned} \]

\origpage{262}
and in general
\[ \int x^{m-1}\,dx\log^{n}\left(\frac{1}{x}\right)=\frac{1\cdot 2\cdot 3\ldots n}{m^{n+1}}.\qquad(g''). \]

\subsection*{(41).}

With the aid of this formula we can express the quantities $T'$, $T''$, etc., by regular series, composed of the reciprocal powers of numbers taken at equal intervals in the sequence of the natural numbers.

Thus, developing first the differential which must be integrated to have $T'$, we have
\[ T'=\int dx\log\frac{1}{x}\left\{\begin{aligned} &x^{p-1}+x^{p+n-1}+x^{p+2n-1}+\text{etc.}\\ &-x^{p+q-1}-x^{p+q+n-1}-x^{p+q+2n-1}-\text{etc.} \end{aligned}\right\}, \]
and carrying out the integration between the limits $x=0$, $x=1$, there results
\[ \begin{aligned} T'={}&\frac{1}{p^{2}}+\frac{1}{(p+n)^{2}}+\frac{1}{(p+2n)^{2}}+\text{etc.}\\ &-\frac{1}{(p+q)^{2}}-\frac{1}{(p+q+n)^{2}}-\frac{1}{(p+q+2n)^{2}}+\text{etc.} \end{aligned} \]
\ednote{The second line closes with $+$ etc. as printed; $-$ etc. would be expected.}
Let us denote in general by $(a,n)^{m}$ the sum of the series
\[ \frac{1}{a^{m}}+\frac{1}{(a+n)^{m}}+\frac{1}{(a+2n)^{m}}+\frac{1}{(a+3n)^{m}}-\text{etc.},\qquad(h'') \]
\ednote{Printed $-$ before etc.; the series has all terms positive, so $+$ etc. is expected. An apparent misprint.}
which will become $(1,n)^{m}$, $(2,n)^{m}$, $(3,n)^{m}$, etc., according as we make $a=1$, $2$, $3$, etc., $n$ being constant: we shall have according to this notation,
\[ \left.\begin{aligned} T'&=(p,n)^{2}-(p+q,n)^{2}\\ \tfrac{1}{2}T''&=(p,n)^{3}-(p+q,n)^{3}\\ \tfrac{1}{2\cdot 3}T'''&=(p,n)^{4}-(p+q,n)^{4}\\ &\text{etc.} \end{aligned}\right\}\qquad(i'') \]

\subsection*{(42).}

It must be observed that, $n$ and $m$ remaining the same, we shall need to consider the values of $(a,n)^{m}$ only from $a=1$ to $a=n$; for it is evident, for example, that $(n+1,n)^{m}$ reduces to $(1,n)^{m}-1$, and that in general we have
\[ (a+n,n)^{m}=(a,n)^{m}-\frac{1}{a^{m}}.\qquad(k'') \]
In consequence of this formula, we should have likewise
\[ (a+2n,n)^{m}=(a,n)^{m}-\frac{1}{a^{m}}-\frac{1}{(a+n)^{m}}, \]

\origpage{263}
and in general every quantity $(a,n)^{m}$, in which $a$ is greater than $n$, will be reduced to a similar quantity in which $a$ does not exceed $n$.

We see further that $(n,n)^{m}$ representing the series
\[ \frac{1}{n^{m}}+\frac{1}{(2n)^{m}}+\frac{1}{(3n)^{m}}+\text{etc.}, \]
this quantity is the same thing as
\[ \frac{1}{n^{m}}\cdot\left(1+\frac{1}{2^{m}}+\frac{1}{3^{m}}+\text{etc.}\right). \]
If therefore we denote by $S_{m}$ the sum of the powers of degree $-m$ of the natural numbers, so that we have
\[ S_{m}=1+\frac{1}{2^{m}}+\frac{1}{3^{m}}+\frac{1}{4^{m}}+\frac{1}{5^{m}}+\text{etc.},\qquad(l'') \]
we shall have $(n,n)^{m}=\frac{1}{n^{m}}S_{m}$. Finally it is evident that we shall also have the equation
\[ S_{m}=(1,n)^{m}+(2,n)^{m}\ldots+(n,n)^{m},\qquad(m'') \]
which, substituting the value of $(n,n)^{m}$, becomes
\[ \left(1-\frac{1}{n^{m}}\right)S_{m}=(1,n)^{m}+(2,n)^{m}+(3,n)^{m}+\ldots+(n-1,n)^{m}. \]

\subsection*{(43).}

Let us consider in particular the case of $n=2$, $p=1$, $q=1$; then we shall have
\[ \begin{aligned} T'&=(1,2)^{2}-(2,2)^{2}=\left(1-\frac{2}{2^{2}}\right)S_{2}=\frac{1}{2}S_{2}\\ T''&=(1,2)^{3}-(2,2)^{3}=\left(1-\frac{2}{2^{3}}\right)S_{3}=\frac{3}{4}S_{3}\\ T'''&=(1,2)^{4}-(2,2)^{4}=\left(1-\frac{2}{2^{4}}\right)S_{4}=\frac{7}{8}S_{4}.\\ &\text{etc.} \end{aligned} \]
\ednote{Printed $T''$ and $T'''$ in the second and third rows; by $(i'')$ these values are $\tfrac{1}{2}T''$ and $\tfrac{1}{2\cdot 3}T'''$.}
It is known that the quantities $S_{2}$, $S_{4}$, $S_{6}$, etc. are known as functions of $\pi$, and that we have $S_{2}=\frac{1}{6}\pi^{2}$, $S_{4}=\frac{1}{90}\pi^{4}$, $S_{6}=\frac{1}{945}\pi^{6}$, etc. As regards the quantities $S_{3}$, $S_{5}$, etc., these are particular transcendents which are not connected at all with the other known transcendents; it is nevertheless easy to find their values with a great approximation,

\origpage{264}
by the beautiful methods which Euler has given for this object (\emph{Calc.\ diff.}, pag.\ 451 et seq.).

By means of these various quantities, we shall therefore know the following integrals, deduced from the equations of art.\ 39.
\[ \begin{aligned} \int\frac{dx}{\sqrt{(1-xx)}}&=\frac{\pi}{2}\\ \int\frac{dx\log\frac{1}{x}}{\sqrt{(1-xx)}}&=\frac{\pi}{2}\log 2\\ \int\frac{dx\log^{2}\frac{1}{x}}{\sqrt{(1-xx)}}&=\frac{\pi}{2}\left(\log^{2}2+\frac{\pi^{2}}{12}\right)\\ \int\frac{dx\log^{3}\frac{1}{x}}{\sqrt{(1-xx)}}&=\frac{\pi}{2}\left(\log^{3}2+3\log 2\cdot\frac{\pi^{2}}{12}+\frac{3}{2}S_{3}\right); \end{aligned} \]
and we shall be able to continue this series indefinitely, in which everything is known, except $S_{3}$, $S_{5}$, etc., of which at least the very approximate values are known, as far as $S_{15}$ (\emph{Calc.\ diff.}, pag.\ 456.).

\section*{On the reduction of the transcendents denoted by $(a,n)^{m}$.}

\subsection*{(44).}

Let us now occupy ourselves with reducing to the smallest possible number the quantities $(1,n)^{m}$, $(2,n)^{m}$, $(3,n)^{m}$, etc. which correspond to one and the same value of $n$. For this purpose, let us take up again the equation (t'), and substitute in it the value of $B$ given by the formula (c''); we shall have
\[ \int\left(\frac{y^{a-1}-y^{n-a-1}}{1-y^{n}}\right)dy=\omega\cot a\omega.\qquad(n'') \]
From which we derive, by differentiating successively with respect to $a$,
\[ \left.\begin{aligned} \int\frac{y^{a-1}+y^{n-a-1}}{1-y^{n}}\,dy\log\frac{1}{y}&=\frac{\omega^{2}}{\sin^{2}a\omega}\\ \frac{1}{2}\int\frac{y^{a-1}-y^{n-a-1}}{1-y^{n}}\,dy\log^{2}\frac{1}{y}&=\frac{\omega^{3}}{\sin^{3}a\omega}\cos a\omega\\ \frac{1}{2\cdot 3}\int\frac{y^{a-1}+y^{n-a-1}}{1-y^{n}}\,dy\log^{3}\frac{1}{y}&=\frac{\omega^{4}}{\sin^{4}a\omega}\left(\frac{1}{3}+\frac{2}{3}\cos^{2}a\omega\right)\\ &\text{etc.} \end{aligned}\right\}\qquad(p'') \]
Putting in place of the first members the values which they obtain

\origpage{265}
by the development in series and the application of the formula (g''), we shall have
\[ \left.\begin{aligned} (a,n)^{2}+(n-a,n)^{2}&=\frac{\omega^{2}}{\sin^{2}a\omega}\\ (a,n)^{3}+(n-a,n)^{3}&=\frac{\omega^{3}}{\sin^{3}a\omega}\cos a\omega\\ (a,n)^{4}+(n-a,n)^{4}&=\frac{\omega^{4}}{\sin^{4}a\omega}\left(\frac{1}{3}+\frac{2}{3}\cos^{2}a\omega\right)\\ &\text{etc.} \end{aligned}\right\}\qquad(q'') \]
\ednote{Second row printed $(a,n)^{3}+(n-a,n)^{3}$; the second row of (p'') and art.\ 47 require $(a,n)^{3}-(n-a,n)^{3}$.}
These formulas will serve to establish between the various quantities $(a,n)^{m}$, which correspond to one and the same value of $n$, all the relations that can be obtained by other ways, and which are given under various forms in the works of Euler. This is what we shall develop in two examples.

\subsection*{(45).}

Let $n=6$ and $m=2$; we shall have the general equation
\[ (a,n)^{2}+(6-a,6)^{2}=\frac{\omega^{2}}{\sin^{2}a\omega}, \]
\ednote{Printed $(a,n)^{2}$ in the first term; with $n=6$ this is $(a,6)^{2}$.}
from which we deduce successively
\[ \left.\begin{aligned} (1,6)^{2}+(5,6)^{2}&=\frac{\omega^{2}}{\sin^{2}\omega}\\ (2,6)^{2}+(4,6)^{2}&=\frac{\omega^{2}}{\sin^{2}2\omega}\\ (3,6)^{2}+(3,6)^{2}&=\frac{\omega^{2}}{\sin^{2}3\omega} \end{aligned}\right\}\qquad(r'') \]
which gives, substituting the value $\omega=\frac{\pi}{6}$,
\[ \begin{aligned} (1,6)^{2}+(5,6)^{2}&=\frac{\pi^{2}}{9}\\ (2,6)^{2}+(4,6)^{2}&=\frac{\pi^{2}}{27}\\ (3,6)^{2}&=\frac{\pi^{2}}{72}. \end{aligned} \]
The sum of these equations is
\[ \left(1-\frac{1}{36}\right)S_{2}=\frac{35}{216}\pi^{2}, \]
from which we deduce $S_{2}=\frac{1}{6}\pi^{2}$, which is a known result, but the

\origpage{266}
demonstration which is given of it here is none the less remarkable.

Besides this value of $S_{2}$, we know $(3,6)^{2}=\frac{1}{72}\pi^{2}$, and $(6,6)^{2}=\frac{1}{36}S=\frac{1}{216}\pi^{2}$\ednote{Printed $S$ without its index at the line end; $S_{2}$ is meant.}; but the four other quantities $(1,6)^{2}$, $(2,6)^{2}$, $(4,6)^{2}$, $(5,6)^{2}$ cannot be determined by the preceding equations.

Let us observe, however, that since we have
\[ (2,6)^{2}=\frac{1}{2^{2}}+\frac{1}{8^{2}}+\frac{1}{14^{2}}+\frac{1}{20^{2}}+\text{etc.}, \]
it follows that
\[ (2,6)^{2}=\frac{1}{4}\left(1+\frac{1}{4^{2}}+\frac{1}{7^{2}}+\frac{1}{10^{2}}+\text{etc.}\right)=\frac{1}{4}(1,6)^{2}+\frac{1}{4}(4,6)^{2}; \]
therefore the four quantities in question can be determined by means of one of them; for example, by means of $(1,6)^{2}$, in the following manner
\[ \left.\begin{aligned} (1,6)^{2}&=\zeta\\ (2,6)^{2}&=\frac{1}{135}\pi^{2}+\frac{1}{5}\zeta\\ (4,6)^{2}&=\frac{4}{135}\pi^{2}-\frac{1}{5}\zeta\\ (5,6)^{2}&=\frac{1}{9}\pi^{2}-\zeta \end{aligned}\right\}\qquad(s'') \]

\subsection*{(46).}

As to the absolute value of $\zeta$, it is not exactly determinable by any known formula; but one can find a value of it as approximate as one pleases by the method which Euler has given (\emph{Calc.\ diff.}, pag.\ 451).

For this, having made
\[ s=\frac{1}{a^{2}}+\frac{1}{(a+n)^{2}}+\frac{1}{(a+2n)^{2}}+\ldots+\frac{1}{(a+nx)^{2}}, \]
we find in general
\[ \begin{aligned} s={}&C-\frac{1}{n}\cdot\frac{1}{a+nx}+\frac{1}{2}\cdot\frac{1}{(a+nx)^{2}}-\frac{A'n}{(a+nx)^{3}}+\frac{B'n^{3}}{(a+nx)^{5}}\\ &-\frac{C'n^{5}}{(a+nx)^{7}}+\frac{D'n^{7}}{(a+nx)^{9}}-\text{etc.} \end{aligned} \]
$A'$, $B'$, $C'$, $D'$, etc. being the sequence of the Bernoullian numbers.

It follows from this formula that, the sum of the series continued to infinity being denoted by $(a,n)^{2}$, we have $(a,n)^{2}=C$; therefore conversely we have

\origpage{267}
\[ \begin{aligned} (a,n)^{2}={}&s+\frac{1}{n}\cdot\frac{1}{a+nx}-\frac{1}{2}\cdot\frac{1}{(a+nx)^{2}}+\frac{A'n}{(a+nx)^{3}}\\ &-\frac{B'n^{3}}{(a+nx)^{5}}+\frac{C'n^{5}}{(a+nx)^{7}}-\text{etc.}\qquad(t'') \end{aligned} \]
It is known that the series $A'$, $B'$, $C'$, $D'$ is divergent from the third term onward, and becomes so more than any given geometric progression; whence it follows that the series contained in the formula (t'') will necessarily become divergent after a certain number of terms. But what is very remarkable is that this formula is none the less fit to give the value of $(a,n)^{2}$ with every degree of approximation that can be desired.

For this we must give to $x$ an arbitrary value, the larger the greater the approximation one wishes to obtain (the value $x=10$ suffices to give 18 or 20 exact decimals). By means of this value, we shall begin by taking the actual sum of the series
\[ s=\frac{1}{a^{2}}+\frac{1}{(a+n)^{2}}+\frac{1}{(a+2n)^{2}}\ldots+\frac{1}{(a+xn)^{2}}; \]
then substituting this value of $s$ as well as that of $x$ in the equation (t''), we shall have for the value of $(a,n)^{2}$ a series at first very convergent, but whose convergence will diminish more and more, up to a certain term where it will become divergent, and this divergence would increase more and more to infinity.

By calculating the successive terms, we shall obtain results alternately greater and smaller than the sought value, and we must stop at the terms where the convergence ceases.

These terms will indicate two very close limits, between which the value of $a_{2}$\ednote{Printed $a_{2}$; the quantity sought is $(a,n)^{2}$.} necessarily lies.

If these two limits did not yet give a sufficient approximation, there would remain no other course to take than to begin a new calculation by giving to $x$ a larger value. But ordinarily, a moderately large value of $x$ will give a very great approximation.

We shall give hereafter an example of the calculation of these sorts of series, which may be called \emph{semi-convergent series}.

\origpage{268}

\subsection*{(47).}

Let now $m=3$ and $n=6$; we shall have the general equation
\[ (a,6)^{3}-(6-a,6)^{3}=\frac{\omega^{3}\cos a\omega}{\sin^{3}a\omega}, \]
from which we deduce the two following:
\[ \left.\begin{aligned} (1,6)^{3}-(5,6)^{3}&=\frac{\omega^{3}\cos\omega}{\sin^{3}\omega}=4\omega^{3}\sqrt{3}\\ (2,6)^{3}-(4,6)^{3}&=\frac{\omega^{3}\cos 2\omega}{\sin^{3}2\omega}=\frac{4\omega^{3}}{3\sqrt{3}} \end{aligned}\right\}\qquad(u'') \]
Moreover we have $(6,6)^{3}=\frac{1}{216}S_{3}$, and
\[ (1,6)^{3}+(2,6)^{3}+(3,6)^{3}+(4,6)^{3}+(5,6)^{3}=\frac{215}{216}S_{3}. \]
These equations are insufficient to determine all the unknowns; but the divisors of 6, which are 2 and 3, furnish new ones.

We have in fact
\[ \begin{aligned} (2,6)^{3}&=\frac{1}{8}\left(1+\frac{1}{4^{3}}+\frac{1}{7^{3}}+\frac{1}{10^{3}}+\text{etc.}\right)=\frac{1}{8}[(1,6)^{3}+(4,6)^{3}]\\ (3,6)^{3}&=\frac{1}{27}\left(1+\frac{1}{3^{3}}+\frac{1}{5^{3}}+\frac{1}{7^{3}}+\text{etc.}\right)=\frac{1}{27}[(1,6)^{3}+(3,6)^{3}+(5,6)^{3}]; \end{aligned} \]
from this we see that the five qnantités\ednote{Printed qnantités, a turned letter; read quantités.} $(1,6)^{3}$, $(2,6)^{3}$, $(3,6)^{3}$, $(4,6)^{3}$, $(5,6)^{3}$ will be determined by supposing one of them known. Thus, making $(1,6)^{3}=\xi$, we shall have
\[ \begin{aligned} (1,6)^{3}&=\xi\\ (2,6)^{3}&=\frac{1}{7}\left(\xi-\frac{4\omega^{3}}{3\sqrt{3}}\right)\\ (3,6)^{3}&=\frac{1}{13}\left(\xi-2\omega^{3}\sqrt{3}\right)\\ (4,6)^{3}&=\frac{1}{7}\left(\xi-\frac{32\omega^{3}}{3\sqrt{3}}\right)\\ (5,6)^{3}&=\xi-4\omega^{3}\sqrt{3}. \end{aligned} \]
The sum of these quantities must be equal to $\frac{215}{216}S_{3}$, thus we shall have
\[ \frac{215}{216}S_{3}=\frac{215}{91}\xi-\frac{430}{91}\omega^{3}\sqrt{3}, \]
or
\[ S_{3}=\frac{216}{91}\xi-\frac{432}{91}\omega^{3}\sqrt{3}. \]

\origpage{269}

Conversely we can express all the quantities $(1,6)^{3}$, $(2,6)^{3}$, etc., by means of $S_{3}$ and of $\omega=\frac{\pi}{6}$, and we find
\[ \left.\begin{aligned} (1,6)^{3}&=\frac{91}{216}S_{3}+2\omega^{3}\sqrt{3}\\ (2,6)^{3}&=\frac{13}{216}S_{3}+\frac{2}{9}\omega^{3}\sqrt{3}\\ (3,6)^{3}&=\frac{7}{216}S_{3}\\ (4,6)^{3}&=\frac{13}{216}S_{3}-\frac{2}{9}\omega^{3}\sqrt{3}\\ (5,6)^{3}&=\frac{91}{216}S_{3}-2\omega^{3}\sqrt{3}\\ (6,6)^{3}&=\frac{1}{216}S_{3} \end{aligned}\right\}\qquad(v'') \]

\subsection*{(48).}

We should easily deduce from the equations (q'') the other known properties of the quantities $(a,n)^{m}$, that is to say of the series which result from the decomposition of the general series
\[ S_{m}=1+\frac{1}{2^{m}}+\frac{1}{3^{m}}+\frac{1}{4^{m}}+\frac{1}{5^{m}}+\text{etc.}, \]
by taking the terms three by three, four by four, or in general $n$ by $n$. We should find, for example, that the series
\[ 1-\frac{1}{3^{m}}+\frac{1}{5^{m}}-\frac{1}{7^{m}}+\text{etc.} \]
is summable when $m$ is odd, and that it is not so when $m$ is even.

We should find on the contrary that the series
\[ 1+\frac{1}{3^{m}}+\frac{1}{5^{m}}+\frac{1}{7^{m}}+\text{etc.} \]
is summable when $m$ is even, and that it is not so when $m$ is odd. The word \emph{summable} is understood here not in an absolute sense, but relatively to the methods known up to the present.

These things presenting no difficulty, and having moreover been demonstrated by other ways, we shall not dwell on them further. We shall show only how one can find, by series which are convergent throughout their extent,

\origpage{270}
the values of the quantities $\zeta$ and $S_{3}$ which have remained unknown in arts.\ 45 and 47; a research more difficult than it appears at first glance, because, in following the methods which present themselves naturally, one falls back on the same difficulty that one wished to resolve.

\subsection*{(49).}

To obtain first the value of $S_{3}$, I observe that we may suppose
\[ \left.\begin{aligned} (1,6)^{3}&=\frac{1}{2}\int\frac{dy}{1-y^{6}}\log^{2}y\\ (2,6)^{3}&=\frac{1}{2}\int\frac{y\,dy}{1-y^{6}}\log^{2}y\\ (3,6)^{3}&=\frac{1}{2}\int\frac{y^{2}dy}{1-y^{6}}\log^{2}y\\ (4,6)^{3}&=\frac{1}{2}\int\frac{y^{3}dy}{1-y^{6}}\log^{2}y\\ (5,6)^{3}&=\frac{1}{2}\int\frac{y^{4}dy}{1-y^{6}}\log^{2}y\\ (6,6)^{3}&=\frac{1}{2}\int\frac{y^{5}dy}{1-y^{6}}\log^{2}y \end{aligned}\right\};\qquad(x'') \]
for the second members being integrated from $y=0$ to $y=1$, by means of the formula (g''), these equations become identities.

Now by the equations (u'') we have
\[ \begin{aligned} (1,6)^{3}-(5,6)^{3}&=\frac{1}{2}\int\frac{1-y^{4}}{1-y^{6}}\,dy\log^{2}y=4\omega^{3}\sqrt{3},\\ (2,6)^{3}-(4,6)^{3}&=\frac{1}{2}\int\frac{y-y^{3}}{1-y^{6}}\,dy\log^{2}y=\frac{4\omega^{3}}{3\sqrt{3}}. \end{aligned} \]
These two equations added together, then subtracted from one another, give the two following
\[ \begin{aligned} \frac{1}{2}\int\frac{dy\log^{2}y}{1-y+y^{2}}&=4\omega^{3}\sqrt{3}+\frac{4\omega^{3}}{3\sqrt{3}}=\frac{40\omega^{3}\sqrt{3}}{9},\\ \frac{1}{2}\int\frac{dy\log^{2}y}{1+y+y^{2}}&=4\omega^{3}\sqrt{3}-\frac{4\omega^{3}}{3\sqrt{3}}=\frac{32\omega^{3}\sqrt{3}}{9}, \end{aligned} \]
so that we have between these integrals the very simple ratio
\[ \int\frac{dy\log^{2}y}{1-y+y^{2}}=\frac{5}{4}\int\frac{dy\log^{2}y}{1+y+y^{2}}; \]

\origpage{271}
but this formula is susceptible of being generalized thus:
\[ \int\frac{dy\log^{m}y}{1-y+y^{2}}=\frac{2^{m}+1}{2^{m}}\int\frac{dy\log^{m}y}{1+y+y^{2}}.\qquad(y'') \]
In fact, let
\[ P=\int\frac{dy\log^{m}y}{1-y+y^{2}},\quad\text{and}\quad Q=\int\frac{dy\log^{m}y}{1+y+y^{2}}, \]
we shall have
\[ \frac{1}{2}(P-Q)=\int\frac{y\,dy\log^{m}y}{1+y^{2}+y^{4}}; \]
putting in this last $y^{\frac{1}{2}}$ in place of $y$, there will result
\[ \frac{1}{2}(P-Q)=\left(\frac{1}{2}\right)^{m+1}\int\frac{dy\log^{m}y}{1+y+y^{2}}=\left(\frac{1}{2}\right)^{M+1}Q, \]
\ednote{The second exponent is printed $M+1$ with a capital; $m+1$ is meant.}
and consequently $P=\frac{2^{m}+1}{2^{m}}Q$.

By the combination of the equations (x''), and the substitution of the values given by the equations (v''), we obtain
\[ \begin{aligned} (1,6)^{3}+(4,6)^{3}&=\frac{1}{2}\int\frac{dy\log^{2}y}{1-y^{3}}=\frac{104}{216}S_{3}+\frac{16}{9}+\omega^{3}\sqrt{3},\\ (3,6)^{3}+(6,6)^{3}&=\frac{1}{2}\int\frac{y^{2}dy\log^{2}y}{1-y^{3}}=\frac{8}{216}S_{3}. \end{aligned} \]
\ednote{In the first line the print sets a stray $+$ between $\frac{16}{9}$ and $\omega^{3}\sqrt{3}$; from the values $(v'')$ the term is $\frac{16}{9}\omega^{3}\sqrt{3}$.}
Subtracting the second from the first, we shall have
\[ \frac{1}{2}\int\frac{(1+y)\,dy\log^{2}y}{1+y+y^{2}}=\frac{4}{9}S_{3}+\frac{19}{9}\omega^{3}\sqrt{3}; \]
\ednote{Printed $\frac{19}{9}$; the subtraction gives $\frac{16}{9}$, which the elimination below requires. An apparent misprint.}
besides, we have already found
\[ \frac{1}{2}\int\frac{dy\log^{2}y}{1+y+y^{2}}=\frac{32}{9}\omega^{3}\sqrt{3}; \]
therefore, eliminating $\omega^{3}$, we shall have
\[ \frac{8}{9}S_{3}=\int\frac{(y+\frac{1}{2})\,dy\log^{2}y}{1+y+y^{2}}.\qquad(z'') \]
It is from this formula that we shall draw the value of $S_{3}$.

\subsection*{(50).}

I remark first that we may put
\[ \int\frac{(y+\frac{1}{2})\,dy\log^{2}y}{1+y+y^{2}}=\int\frac{dy}{1+y}\log^{2}y-\frac{1}{2}\int\frac{(1-y)\,dy\log^{2}y}{(1+y)(1+y+y^{2})}. \]

\origpage{272}
The first part
\[ \int\frac{dy}{1+y}\log^{2}y=2\left(1-\frac{1}{2^{3}}+\frac{1}{3^{3}}-\frac{1}{4^{3}}+\text{etc.}\right)=2S_{3}\left(1-\frac{2}{8}\right)=\frac{3}{2}S_{3}, \]
therefore we shall have
\[ \frac{11}{18}S_{3}=\int\frac{\frac{1}{2}(1-y)\,dy\log^{2}y}{(1+y)(1+y+y^{2})}: \]
let $y=\frac{1-z}{1+z}$, and the transformed will be
\[ \frac{11}{18}S_{3}=\int\frac{z\,dz}{3+z^{2}}\log^{2}y, \]
a formula which must always be integrated from $z=0$ to $z=1$.

We see now that, as $z^{2}$ is always smaller than 1, we can reduce $\dfrac{z\,dz}{3+z^{2}}$ to a convergent series, so that the integral will depend only on terms of the form \ldots\ $\displaystyle\int z^{2m+1}dz\log^{2}y$. But to render the series still more convergent, I make $z^{2}=1-u^{2}$, and I have
\[ \frac{11}{18}S_{3}=\int\frac{u\,du}{4-u^{2}}\log^{2}y, \]
an integral which must again be taken from $u=0$ to $u=1$.

This integral, being taken by parts, becomes
\[ -\frac{1}{2}\log\left(1-\frac{u^{2}}{4}\right)\cdot\log^{2}y+\int\frac{dy}{y}\log y\cdot\log\left(1-\frac{u^{2}}{4}\right); \]
the first part vanishes at the two limits of the integral; thus we have simply
\[ \frac{11}{18}S_{3}=\int\frac{dy}{y}\log y\log\left(1-\frac{u^{2}}{4}\right). \]
Developing $\log\left(1-\dfrac{u^{2}}{4}\right)$, and substituting the value $\dfrac{dy}{y}=-\dfrac{2dz}{u^{2}}$, we shall have
\[ \frac{11}{18}S_{3}=\frac{1}{2}\int dz\log\frac{1}{y}\cdot\left(1+\frac{1}{2}\cdot\frac{u^{2}}{4}+\frac{1}{3}\cdot\frac{u^{4}}{4^{2}}+\frac{1}{4}\cdot\frac{u^{6}}{4^{3}}+\text{etc.}\right). \]

\origpage{273}
Let then
\[ \begin{aligned} \frac{1}{2}\int dz\log\frac{1}{y}&=P^{0},\\ \frac{1}{2}\int u^{2}dz\log\frac{1}{y}&=P',\\ \frac{1}{2}\int u^{4}dz\log\frac{1}{y}&=P'',\\ &\text{etc.} \end{aligned} \]
these integrals being taken from $z=0$, to $z=1$, and we shall have
\[ \frac{11}{18}S_{3}=P^{0}+\frac{1}{2}\cdot\frac{P'}{4}+\frac{1}{3}\cdot\frac{P''}{4^{2}}+\frac{1}{4}\cdot\frac{P'''}{4^{3}}+\text{etc.} \]

\subsection*{(51).}

To have now the quantities $P^{0}$, $P'$, $P''$, etc. I observe that we have in general
\[ \int u^{2m}dz\log\frac{1}{y}=\log\frac{1}{y}\int u^{2m}dz+\int\frac{dy}{y}\cdot\int u^{2m}dz. \]
In this formula we shall suppose that the integral $\displaystyle\int u^{2m}dz$ is taken in such a way that it vanishes when $z=1$ or $u=0$; then the part $\displaystyle\log\frac{1}{y}\int u^{2m}dz$ vanishes at the two limits of the integral, and we have simply
\[ \frac{1}{2}\int u^{2m}dz\log\frac{1}{y}=\frac{1}{2}\int\frac{dy}{y}\int u^{2m}dz=\int-\frac{dz}{u^{2}}\int u^{2m}dz, \]
where we see that the logarithms have entirely disappeared, and that we need no longer take account of anything but the relation $u^{2}=1-z^{2}$.

We have moreover
\[ \int u^{2m}dz=\frac{2m}{2m+1}\int u^{2m-2}dz+\frac{u^{2m}z}{2m+1}; \]
multiplying both sides by $\dfrac{dz}{u^{2}}$ and integrating, we shall have
\[ P^{(m)}=\frac{2m}{2m+1}P^{(m-1)}-\int\frac{u^{2m-2}z\,dz}{2m+1}. \]
Now $z\,dz=-u\,du$, and consequently
\[ -\int\frac{u^{2m-2}z\,dz}{2m+1}=\int\frac{u^{2m-1}du}{2m+1}=\frac{u^{2m}}{(2m+1)2m}+C; \]

\origpage{274}
this integral having to be taken from $z=0$, to $z=1$, it reduces to $-\dfrac{1}{2m(2m+1)}$; therefore we shall have
\[ P^{(m)}=\frac{2m}{2m+1}\cdot P^{(m-1)}-\frac{1}{2m(2m+1)}. \]
The first value
\[ P^{0}=-\int\frac{dz}{u^{2}}\int dz=-\int\frac{dz}{u^{2}}(z-1)=\int\frac{dz}{1+z}=L(1+z)=L2; \]
therefore we shall have successively
\[ \begin{aligned} P^{0}&=\log 2,\\ P'&=\frac{2}{3}P^{0}-\frac{1}{2\cdot 3},\\ P''&=\frac{4}{5}P'-\frac{1}{4\cdot 5},\\ P'''&=\frac{6}{7}P''-\frac{1}{6\cdot 7},\\ &\text{etc.} \end{aligned} \]
By means of this very simple law we shall easily calculate the different terms of the decreasing series $P^{0}$, $P'$, $P''$, $P'''$, etc. Then we shall have $S_{3}$ by the formula
\[ S_{3}=\frac{18}{11}\left(P^{0}+\frac{1}{2}\cdot\frac{P'}{4}+\frac{1}{3}\cdot\frac{P''}{4^{2}}+\frac{1}{4}\cdot\frac{P'''}{4^{3}}+\text{etc.}\right); \]
which gives a very rapid approximation, since each term is less than the fourth part of the preceding.

By calculating this series as far as the term $P^{\text{vi}}$ inclusively, we find $S_{3}=1.2020567$. Euler found by the method of which we have spoken, and by pushing the approximation much further,
\[ S_{3}=1.202056903159594281. \]
(See \emph{Calc.\ diff.}, page 453).

\subsection*{(52).}

To find by similar procedures the value of the transcendent $\zeta$, which remained unknown in the equations (s''), I observe that we have

\origpage{275}
\[ \begin{aligned} (1,6)^{2}&=\int\frac{dy}{1-y^{6}}\log\frac{1}{y}=\zeta,\\ (2,6)^{2}&=\int\frac{y\,dy}{1-y^{6}}\log\frac{1}{y}=\tfrac{1}{135}\pi^{2}+\tfrac{1}{5}\zeta,\\ (4,6)^{2}&=\int\frac{y^{3}\,dy}{1-y^{6}}\log\frac{1}{y}=\tfrac{4}{135}\pi^{2}-\tfrac{1}{5}\zeta,\\ (5,6)^{2}&=\int\frac{y^{4}\,dy}{1-y^{6}}\log\frac{1}{y}=\tfrac{1}{9}\pi^{2}-\zeta: \end{aligned} \]
from this I derive
\[ (1,6)^{2}+(2,6)^{2}-(4,6)^{2}-(5,6)^{2}=\int\frac{dy\log\frac{1}{y}}{1-y+y^{2}}=\frac{12}{5}\zeta-\frac{2}{15}\pi^{2}; \]
therefore
\[ \zeta=\tfrac{1}{18}\pi^{2}+\frac{5}{12}\int\frac{dy\log\frac{1}{y}}{1-y+y^{2}}; \]
everything is thus reduced to finding the value of this integral.

If we put successively $y=\frac{1}{2}(1-u)$, $y=\frac{1}{2}(1+u)$, and add the two transformed integrals taken positively, we shall have
\[ \int\frac{dy\log\frac{1}{y}}{1-y+y^{2}}=\int\frac{2\,du}{3+u^{2}}\log\left(\frac{4}{1-u^{2}}\right), \]
a new integral which must again be taken from $u=0$ to $u=1$.

The first part of this integral
\[ \int\frac{2\,du\log 4}{3+u^{2}}=\frac{2}{\sqrt{3}}\log 4\cdot\operatorname{arc}\operatorname{tang}\frac{u}{\sqrt{3}}, \]
and on putting $u=1$, it reduces to
\[ \frac{2\log 4}{\sqrt{3}}\cdot\frac{\pi}{6}\quad\text{or}\quad\frac{2\pi\log 2}{3\sqrt{3}}. \]
The other part
\[ \int-\frac{2\,du}{3+u^{2}}\log(1-u^{2}) \]
being called $T$, we shall have, by the development of the fraction,
\[ T=\int-\frac{2\,du}{3}\left(1-\frac{u^{2}}{3}+\frac{u^{4}}{3^{2}}-\frac{u^{6}}{3^{3}}+\text{etc.}\right)\log(1-u^{2}): \]
now, integrating by parts, we have
\[ \int-u^{2m}\,du\log(1-u^{2})=\frac{1-u^{2m+1}}{2m+1}\log(1-u^{2})+\int\frac{1-u^{2m+1}}{2m+1}\cdot\frac{2u\,du}{1-u^{2}}, \]

\origpage{276}
and since the quantity $(1-u^{2m+1})\log(1-u^{2})$ vanishes at the two limits of the integral, we simply have
\[ \begin{aligned} \int-u^{2m}\,du\log(1-u^{2})&=\frac{1}{2m+1}\int\frac{1-u^{2m+1}}{1-u^{2}}\cdot 2u\,du\\ &=\frac{2}{2m+1}\int\left(\frac{1-u^{2m+2}}{1-u^{2}}\,du-\frac{du}{1+u}\right). \end{aligned} \]
Developing the quantity under the sign and integrating from $u=0$ to $u=1$, we shall have
\[ \int-u^{2m}\,du\log(1-u^{2})=\frac{2}{2m+1}\left(1+\frac{1}{3}+\frac{1}{5}\ldots+\frac{1}{2m+1}-\log 2\right); \]
therefore finally
\[ \begin{aligned} T={}&\frac{4}{3}\cdot(1-\log 2)\\ &-\frac{4}{3^{2}}\cdot\frac{1}{3}\left(1+\frac{1}{3}-\log 2\right)\\ &+\frac{4}{3^{3}}\cdot\frac{1}{5}\left(1+\frac{1}{3}+\frac{1}{5}-\log 2\right)\\ &-\frac{4}{3^{4}}\cdot\frac{1}{7}\left(1+\frac{1}{3}+\frac{1}{5}+\frac{1}{7}-\log 2\right)\\ &+\text{etc.}, \end{aligned} \]
a convergent series, since each term is less than the third part of the preceding one. $T$ being known, we shall have $\zeta$ by the value
\[ \zeta=\frac{\pi^{2}}{18}+\frac{5\pi\log 2}{18\sqrt{3}}+\frac{5}{12}T. \]

\section*{On the Eulerian integrals of the second species.}

\subsection*{(53).}

Let us now consider the integral
\[ \int x^{m-1}\,dx\left(\log\frac{1}{x}\right)^{n}, \]
which we shall always suppose taken between the limits $x=0$, $x=1$. We have first, integrating by parts,
\[ \int x^{m-1}\,dx\left(\log\frac{1}{x}\right)^{n}=\frac{x^{m}}{m}\left(\log\frac{1}{x}\right)^{n}+\frac{n}{m}\int x^{m-1}\,dx\left(\log\frac{1}{x}\right)^{n-1}, \]
and as the part outside the sign vanishes at the two limits,

\origpage{277}
provided that we suppose $n>0$, we shall then have
\[ \int x^{m-1}\,dx\left(\log\frac{1}{x}\right)^{n}=\frac{n}{m}\int x^{m-1}\,dx\left(\log\frac{1}{x}\right)^{n-1}.\qquad(\alpha) \]
We shall therefore have in general, if $n$ is a positive whole number,
\[ \int x^{m-1}\,dx\left(\log\frac{1}{x}\right)^{n}=\frac{n\cdot n-1\cdot n-2\cdot n-3\ldots 1}{m^{n+1}}.\qquad(\beta) \]
When $n$ is not a whole number, the integral $\displaystyle\int x^{m-1}\,dx\left(\log\frac{1}{x}\right)^{n}$ will in general be a transcendent whose properties it is fitting to examine.

And first, by means of the formula $(\alpha)$, this transcendent can always be brought back to the case where the exponent $n$ lies between 0 and $-1$.

Moreover, I observe that, without diminishing anything of the generality of the calculation, one may take $m=1$; for the formula $\displaystyle\int x^{m-1}\,dx\left(\log\frac{1}{x}\right)^{n}$ being proposed, if one puts $x^{m}=z$, this formula will become $\displaystyle\frac{1}{m^{n+1}}\int dz\left(\log\frac{1}{z}\right)^{n}$.

\subsection*{(54).}

This being laid down, it will suffice to consider the integral $\displaystyle\int dx\left(\log\frac{1}{x}\right)^{a-1}$, in which we shall suppose that $a$ is positive and smaller than unity. This quantity being simply a function of $a$, we shall designate it by $\Gamma(a)$, and we shall put
\[ \Gamma(a)=\int dx\left(\log\frac{1}{x}\right)^{a-1}.\qquad(\gamma) \]
The object of the following researches is to evaluate the function $\Gamma(a)$, when $a$ is a given rational fraction, such as $\frac{1}{2}$, $\frac{2}{3}$, $\frac{3}{4}$, etc., and we propose in particular to compare with one another the functions which correspond to values of $a$ of the same denomination, such as $\Gamma\left(\frac{1}{n}\right)$, $\Gamma\left(\frac{2}{n}\right)$, etc. Finally we shall also seek to determine by approximation the transcendent $\Gamma(a)$ for every value of $a$, rational or irrational.

\subsection*{(55).}

Taking the integrals from $x=0$ to $x=1$, and

\origpage{278}
supposing $m>1$, we have this formula of reduction,
\[ \int x^{\alpha-1}\,dx(1-x^{\beta})^{m-1}=\frac{(m-1)\beta}{\alpha+(m-1)\beta}\int x^{\alpha-1}\,dx(1-x^{\beta})^{m-2}; \]
whence it follows that if $m$ is a whole number, we shall have exactly
\[ \int x^{\alpha-1}\,dx(1-x^{\beta})^{m-1}=\frac{1\cdot 2\cdot 3\ldots(m-1)\beta^{m-1}}{\alpha\cdot\alpha+\beta\cdot\alpha+2\beta\ldots\alpha+(m-1)\beta}. \]

Let us designate by $\Phi(\alpha,m)$ the integral $\displaystyle\int x^{\alpha-1}\,dx(1-x^{\beta})^{m-1}$, taken from $x=0$ to $x=1$, and on the hypothesis that $m$ is a positive whole number; we shall then have
\[ \frac{1\cdot 2\cdot 3\ldots(m-1)}{\alpha+\beta\cdot\alpha+2\beta\ldots\alpha+(m-1)\beta}=\frac{\alpha}{\beta^{m-1}}\Phi(\alpha,m). \]
In this equation, let us put successively $\lambda m$ and $(\lambda+1)m$ in place of $m$; we shall have
\[ \begin{gathered} \frac{1\cdot 2\cdot 3\ldots\overline{\lambda m-1}}{\alpha+\beta\cdot\alpha+2\beta\ldots\alpha+\overline{\lambda m-1}\cdot\beta}=\frac{\alpha}{\beta^{\lambda m-1}}\cdot\Phi(\alpha,\lambda m)\\ \frac{1\cdot 2\cdot 3\ldots\overline{\lambda m+m-1}}{\alpha+\beta\cdot\alpha+2\beta\ldots\alpha+\overline{\lambda m+m-1}\cdot\beta}=\frac{\alpha}{\beta^{\lambda m+m-1}}\cdot\Phi(\alpha,\lambda m+m). \end{gathered} \]
Dividing the second by the first, and putting for brevity $\alpha+\lambda m\beta=\alpha'$, there results
\[ \frac{\alpha'\cdot\alpha'+\beta\cdot\alpha'+2\beta\ldots\alpha'+\overline{m-1}\cdot\beta}{\lambda m\cdot\lambda m+1\cdot\lambda m+2\ldots\lambda m+m-1}=\beta^{m}\cdot\frac{\Phi(\alpha,\lambda m)}{\Phi(\alpha,\lambda m+m)}; \]
but putting $\alpha'$ in place of $\alpha$ in the primitive equation, we have
\[ \frac{1\cdot 2\cdot 3\ldots m-1}{\alpha'+\beta\cdot\alpha'+2\beta\ldots\alpha'+\overline{m-1}\cdot\beta}=\frac{\alpha'}{\beta^{m-1}}\Phi(\alpha+\lambda m\beta,m), \]
and multiplying these two equations together, we have as product
\[ \frac{1\cdot 2\cdot 3\ldots m-1}{\lambda m\cdot\lambda m+1\ldots\lambda m+m-1}=\beta\cdot\frac{\Phi(\alpha,\lambda m)\cdot\Phi(\alpha+\lambda m\beta,m)}{\Phi(\alpha,\lambda m+m)}. \]
The first member, by virtue of the primitive equation, can be represented by $\displaystyle\int x^{\lambda m-1}\,dx(1-x)^{m-1}$; and because one may put $x^{n}$ in place of $x$, without changing the limits of the integral, it can be

\origpage{279}
also represented by $\displaystyle n\int x^{\lambda nm-1}\,dx(1-x^{n})^{m-1}$; therefore we have
\[ n\int x^{\lambda mn-1}\,dx(1-x^{n})^{m-1}=\beta\frac{\Phi(\alpha,\lambda m)\cdot\Phi(\alpha+\lambda m\beta,m)}{\Phi(\alpha,\lambda m+m)}.\qquad(\delta) \]
This equation, thus expressed in a finite number of terms, acquires a greater generality, and no longer supposes that $m$ is a whole number.

Indeed, the two members having to reduce to one and the same function of $m$ and of $\lambda m$, which is
\[ \frac{1}{\lambda m}-\frac{m-1}{1}\cdot\frac{1}{\lambda m+1}+\frac{m-1\cdot m-2}{1\cdot 2}\cdot\frac{1}{\lambda m+2}-\text{etc.}, \]
we are free to give to $m$ and $\lambda$ any positive values whatever, and with all the more reason to the quantities $\alpha$, $\beta$, $n$, which disappear from the two members.

\subsection*{(56).}

Let then $\alpha=1$ and $\beta=$ an infinitely small quantity; we shall have
\[ 1-x^{\beta}=\beta\log\frac{1}{x},\quad\text{and}\quad\Phi(\alpha,m)=\beta^{m-1}\int dx\left(l\,\frac{1}{x}\right)^{m-1}, \]
so that we shall have in general
\[ \Phi(\alpha,k)=\beta^{k-1}\cdot\Gamma(k). \]
By means of this formula, the equation $(\delta)$ becomes
\[ n\int x^{\lambda mn-1}\,dx(1-x^{n})^{m-1}=\frac{\Gamma(\lambda m)\cdot\Gamma(m)}{\Gamma(\lambda m+m)}. \]
Let now $m=\frac{q}{n}$, $\lambda m=\frac{p}{n}$; we shall then have
\[ \int x^{p-1}\,dx(1-x^{n})^{\frac{q}{n}-1}=\frac{\Gamma\left(\frac{q}{n}\right)\cdot\Gamma\left(\frac{q}{n}\right)}{n\Gamma\left(\frac{p+q}{n}\right)}. \]
\ednote{The first factor of the numerator is printed $\Gamma\left(\frac{q}{n}\right)$ (confirmed by magnification); the sense requires $\Gamma\left(\frac{p}{n}\right)$, as in equation $(\varepsilon)$ below. A printer's misprint.}
The first member is nothing other than the transcendent designated above by $\left(\frac{p}{q}\right)$; thus we shall have this remarkable equation
\[ \left(\frac{p}{q}\right)=\frac{\Gamma\left(\frac{p}{n}\right)\cdot\Gamma\left(\frac{q}{n}\right)}{n\Gamma\left(\frac{p+q}{n}\right)}.\qquad(\varepsilon) \]

\origpage{280}
Whence one sees that the transcendent $\left(\frac{p}{q}\right)$ would be known, if one knew, for the same value of $n$, the functions of the form $\Gamma\left(\frac{a}{n}\right)$, $a$ being a whole number.

\subsection*{(57).}

It results first from this value of $\left(\frac{p}{q}\right)$, that one may exchange the numbers $p$ and $q$ with each other, and that thus we have $\left(\frac{p}{q}\right)=\left(\frac{q}{p}\right)$, which is one of the principal properties of these functions.

Moreover we derive from this formula
\[ \left(\frac{p}{q}\right)\cdot\left(\frac{p+q}{r}\right)=\frac{\Gamma\left(\frac{p}{n}\right)\cdot\Gamma\left(\frac{q}{n}\right)\Gamma\left(\frac{r}{n}\right)}{n^{2}\Gamma\left(\frac{p+q+r}{n}\right)} \]
In the second member, it is evident that one may make the permutation between any two of the numbers $p$, $q$, $r$, at will, which gives the fundamental theorem
\[ \left(\frac{p}{q}\right)\cdot\left(\frac{p+q}{r}\right)=\left(\frac{p}{r}\right)\cdot\left(\frac{p+r}{q}\right); \]
of which we have thus a new and very simple demonstration.

\subsection*{(58).}

Let us now show how the functions $\Gamma$ are determined by means of the functions $\left(\frac{p}{q}\right)$.

Let us observe first that by means of the equation $(\alpha)$, we have in general
\[ \Gamma(n+1)=n\Gamma(n),\qquad(\zeta) \]
which will allow us to reduce the cases where $n$ is greater than unity to those where it is smaller.

If we have $n=1$, then $\Gamma(n)$ reduces to $\displaystyle\int dx=x=1$; thus we have
\[ \Gamma(1)=1.\qquad(\eta) \]
This being laid down, let us put $q=n-p$ in the equation $(\varepsilon)$; then the value of the first member is known, and we have
\[ \frac{\omega}{\sin p\omega}=\frac{\Gamma\left(\frac{p}{n}\right)\Gamma\left(\frac{n-p}{n}\right)}{n\Gamma(1)}, \]

\origpage{281}
or, in other terms,
\[ \Gamma(a)\Gamma(1-a)=\frac{\pi}{\sin a\pi}.\qquad(\theta) \]
When $a=\frac{1}{2}$, we have $(\Gamma a)^{2}=\pi$; therefore
\[ \Gamma\left(\tfrac{1}{2}\right)=\sqrt{\pi}.\qquad(\iota) \]

The equation $(\theta)$, very remarkable in this theory, shows that it suffices to know the value of $\Gamma(x)$ from $x=\frac{1}{2}$ to $x=1$, and from it one will deduce the other values of this function from $x=\frac{1}{2}$ to $x=0$. Moreover, the value of $\Gamma(x)$ from $x=\frac{1}{2}$ to $x=1$ varies only between the limits $\sqrt{\pi}$ and 1, or $1{,}77245$ and 1, whereas from $x=\frac{1}{2}$ to $x=0$ it varies from $1{,}77245$ to infinity; and in particular when $a$ is infinitely small, the formula $(\theta)$ gives $\Gamma(a)=\frac{1}{a}$.

\subsection*{(59).}

If in the equation $(\varepsilon)$ we put $p=1$, and take successively $q=1$, 2, 3\ldots\ up to $n=1$\ednote{Printed $n=1$; the sense requires $n-1$, the last equation of the series below being that for $\left(\frac{1}{n-1}\right)$. An apparent misprint.}, we shall have a series of equations
\[ \begin{gathered} \left(\frac{1}{1}\right)=\frac{\Gamma\left(\frac{1}{n}\right)\Gamma\left(\frac{1}{n}\right)}{n\Gamma\left(\frac{2}{n}\right)}\\ \left(\frac{1}{2}\right)=\frac{\Gamma\left(\frac{1}{n}\right)\Gamma\left(\frac{2}{n}\right)}{n\Gamma\left(\frac{3}{n}\right)}\\ \left(\frac{1}{3}\right)=\frac{\Gamma\left(\frac{1}{n}\right)\Gamma\left(\frac{3}{n}\right)}{n\Gamma\left(\frac{4}{n}\right)}\\ \vdots\\ \left(\frac{1}{n-1}\right)=\frac{\Gamma\left(\frac{1}{n}\right)\Gamma\left(\frac{n-1}{n}\right)}{n\Gamma\left(\frac{n}{n}\right)}. \end{gathered} \]
Multiplying the first $a-1$ equations together, we shall have the product
\[ \left(\frac{1}{1}\right)\cdot\left(\frac{1}{2}\right)\cdot\left(\frac{1}{3}\right)\cdot\left(\frac{1}{4}\right)\ldots\left(\frac{1}{a-1}\right)=\frac{\left(\Gamma\frac{1}{n}\right)^{a}}{n^{a-1}\Gamma\left(\frac{a}{n}\right)}, \]

\origpage{282}
if we multiply them all, or put $a=n$, the product will give
\[ \Gamma\left(\frac{1}{n}\right)=n\sqrt[n]{\left[\frac{1}{n}\cdot\left(\frac{1}{1}\right)\cdot\left(\frac{1}{2}\right)\cdot\left(\frac{1}{3}\right)\ldots\left(\frac{1}{n-1}\right)\right]}. \]
Let then for brevity
\[ T=\sqrt[n]{\left[\frac{1}{n}\cdot\left(\frac{1}{1}\right)\cdot\left(\frac{1}{2}\right)\cdot\left(\frac{1}{3}\right)\ldots\left(\frac{1}{n-1}\right)\right]}, \]
and we shall have
\[ \Gamma\left(\frac{1}{n}\right)=nT. \]
Knowing $\Gamma\left(\frac{1}{n}\right)$, we shall deduce $\Gamma\left(\frac{a}{n}\right)$ by the equation already found, which gives
\[ \Gamma\left(\frac{a}{n}\right)=\frac{nT^{a}}{\left(\frac{1}{1}\right)\left(\frac{1}{2}\right)\left(\frac{1}{3}\right)\ldots\left(\frac{1}{a-1}\right)}=\frac{\left(\frac{1}{a}\right)\left(\frac{1}{n+a}\right)\ldots\left(\frac{1}{n-1}\right)}{T^{n-a}}.\qquad(\kappa) \]
\ednote{Printed $\left(\frac{1}{n+a}\right)$ as the second factor; the sense requires $\left(\frac{1}{a+1}\right)$. An apparent misprint.}
And as the quantities $\left(\frac{1}{1}\right)$, $\left(\frac{1}{2}\right)$, $\left(\frac{1}{3}\right)$, etc. are assumed known for each value of $n$, nothing more remains to be desired for the determination of the functions $\Gamma\left(\frac{a}{n}\right)$.

\subsection*{(60).}

Conversely, if one knew the value of $\Gamma(x)$ for every rational value of $x$ less than unity, it would be easy to determine the integral $\left(\frac{p}{q}\right)$; for we have in general
\[ \left(\frac{p}{q}\right)=\frac{\Gamma\left(\frac{p}{n}\right)\Gamma\left(\frac{q}{n}\right)}{n\Gamma\left(\frac{p+q}{n}\right)};\qquad(\lambda) \]
and if it happens that $p+q$ is $>n$, we shall change, by the equation $(\zeta)$, this formula into the following
\[ \left(\frac{p}{q}\right)=\frac{\Gamma\left(\frac{p}{n}\right)\Gamma\left(\frac{q}{n}\right)}{(p+q-n)\Gamma\left(\frac{p+q-n}{n}\right)}.\qquad(\lambda') \]
We shall remark that these formulas are fit to give the general expression of $\left(\frac{p}{q}\right)$ by means of the quantities of the same

\origpage{283}
species $\left(\frac{1}{1}\right)$, $\left(\frac{1}{2}\right)$, $\left(\frac{1}{3}\right)$, etc. Indeed, the values of $\Gamma\left(\frac{p}{n}\right)$, $\Gamma\left(\frac{q}{n}\right)$, $\Gamma\left(\frac{p+q}{n}\right)$ being drawn from the formula $(\kappa)$, we shall deduce from them
\[ \left(\frac{p}{q}\right)=\frac{\left(\frac{1}{p}\right)\left(\frac{1}{p+1}\right)\left(\frac{1}{p+2}\right)\ldots\left(\frac{1}{p+q-1}\right)}{\left(\frac{1}{1}\right)\left(\frac{1}{2}\right)\left(\frac{1}{3}\right)\ldots\left(\frac{1}{q-1}\right)},\qquad(\mu) \]
a formula which will serve so long as $p+q$ is smaller than $n$.

If we have $p+q>n$, we must make use of the second formula, which will give, by similar substitutions,
\[ \left(\frac{p}{q}\right)=\frac{1}{p+q-n}\cdot\frac{\left(\frac{1}{p}\right)\left(\frac{1}{p+1}\right)\left(\frac{1}{p+2}\right)\ldots\left(\frac{1}{n-1}\right)}{\left(\frac{1}{p+q-n}\right)\left(\frac{1}{p+q-n+1}\right)\ldots\left(\frac{1}{q-1}\right)},\qquad(\mu') \]
These formulas correspond to the equations (k) and (n) found above, and one would make them coincide entirely by substituting, in place of each quantity $\left(\frac{1}{a}\right)$, its value
\[ \left(\frac{1}{a}\right)=\frac{A_{a}\sin(a+1)\omega}{\sin\omega}; \]
thus the functions $\Gamma$ offer a new, direct and very simple means of determining the general expression of the quantities $\left(\frac{p}{q}\right)$.

\subsection*{(61).}

To return to the quantities $\Gamma(a)$, we have already found the equation
\[ \Gamma(a)\cdot\Gamma(1-a)=\frac{\pi}{\sin a\pi}, \]
by means of which the values of the function from $a=0$ to $a=\frac{1}{2}$ are deduced from the values supposed known from $a=\frac{1}{2}$ to $a=1$.

We shall prove further on that it suffices to know the values of the function in the half of this interval, that is to say only from $a=\frac{3}{4}$ to $a=1$, and from them one will deduce all the other values.

Indeed, if we suppose $a<\frac{1}{2}n$, the equation $(\varepsilon)$ gives at the\ednote{The print breaks off here: p.~284 opens directly with the display, so a word (presumably ``fois'') is omitted. Kept as printed.}

\origpage{284}
\[ \begin{gathered} \left(\frac{a}{a}\right)=\frac{\Gamma\left(\frac{a}{n}\right)\Gamma\left(\frac{a}{n}\right)}{n\Gamma\left(\frac{2a}{n}\right)}\\ \left(\frac{\tfrac{1}{2}n-a}{\tfrac{1}{2}n-a}\right)=\frac{\Gamma\left(\frac{\tfrac{1}{2}n-a}{n}\right)\Gamma\left(\frac{\tfrac{1}{2}n-a}{n}\right)}{n\Gamma\left(\frac{n-2a}{n}\right)}. \end{gathered} \]
Substituting these values in the equation $(d')$, then simply putting $a$ in place of $\frac{a}{n}$, and reducing the functions by the formula $(\theta)$, we shall have
\[ \Gamma(a)=\frac{2^{1-2a}}{\sqrt{\pi}}\cos a\pi\cdot\Gamma(2a)\cdot\Gamma\left(\tfrac{1}{2}-a\right), \]
an equation which supposes $a<\frac{1}{2}$.

This equation, combined with the equation $(\theta)$, will give
\[ \Gamma(1-a)=\frac{2^{2a}}{\sqrt{\pi}}\cos a\pi\,\Gamma(1-2a)\Gamma\left(\tfrac{1}{2}+a\right); \]
finally from this one deduces, putting $a-\frac{1}{2}$ in place of $a$,
\[ \Gamma(a)=\frac{2^{1-2a}\sqrt{\pi}}{\sin a\pi}\cdot\frac{\Gamma\left(\tfrac{3}{2}-a\right)}{\Gamma(2-2a)}\qquad(\nu) \]
We suppose known the values of $\Gamma(a)$ from $a=\frac{3}{4}$ to $a=1$.

This being laid down, 1°. The second member of the equation $(\nu)$ will be known for every value of $a$, from $a=\frac{1}{2}$ to $a=\frac{5}{8}$; therefore we shall know $\Gamma(a)$ in this same interval from $a=\frac{1}{2}$ to $a=\frac{5}{8}$.

2°. By means of this first case, the second member will be known if $2-2a$ lies between $\frac{1}{2}$ and $\frac{5}{8}$; we shall therefore know $\Gamma(a)$ whenever $a$ lies between $\frac{11}{16}$ and $\frac{12}{16}$.

3°. By means of the first two cases, the second member of the equation $(\nu)$ will be known if $2-2a$ lies between $\frac{11}{16}$ and $\frac{12}{16}$; therefore we shall know $\Gamma(a)$ for all the values of $a$ lying from $a=\frac{5}{8}=\frac{20}{32}$ to $a=\frac{21}{32}$.

4°. The second member will still be known if $2-2a$ lies between $\frac{20}{32}$ and $\frac{21}{32}$; therefore $\Gamma(a)$ will be known from $a=\frac{44}{64}=\frac{11}{16}$ to $a=\frac{43}{64}$, and so on.

By these various operations the values of $a$ for which $\Gamma(a)$

\origpage{285}
becomes known approach alternately the limit $\frac{2}{3}$, which they nevertheless reach only at infinity, since $\frac{2}{3}$ cannot be expressed exactly in fractions whose denominator is a power of 2. But one sees that by four operations only, the interval where $\Gamma(a)$ remains to be determined extends no further than from $a=\frac{44}{64}$ to $a=\frac{43}{64}$. A fifth operation would narrow this interval from $\frac{43}{64}$ or $\frac{86}{128}$ to $\frac{85}{128}$, and so on.

The common limit of these sequences is $\frac{2}{3}$ and $\Gamma\left(\frac{2}{3}\right)$ is determined directly by putting $a=\frac{2}{3}$ in the formnle\ednote{Printed ``formnle'' for ``formule''; a turned letter.} $(\nu)$, which gives
\[ \Gamma\left(\frac{2}{3}\right)^{2}=\frac{2^{-\frac{1}{3}}\sqrt{\pi}}{\sin\frac{2}{3}\pi}\Gamma\left(\frac{5}{6}\right). \]

\subsection*{(62).}

In the particular cases where one would seek to determine a value of $\Gamma(a)$, one need not trouble oneself in any way about the distinction of the preceding cases, and the immediate application of the formula $(\nu)$, repeated as many times as is necessary, or until there is no longer any unknown to determine, will always lead to the result sought.

Let it be proposed, for example, to find the value of $\Gamma(0.675)$; we shall have directly
\[ \Gamma(0.675)=\frac{2^{-0.35}\sqrt{\pi}}{\sin 38^{\circ}30'}\cdot\frac{\Gamma(0.825)}{\Gamma(0.65)}. \]
\ednote{Read $\sin 38^{\circ}30'$ at scan resolution; formula $(\nu)$ with $a=0.675$ requires $\sin 121^{\circ}30'=\sin 58^{\circ}30'$. Either a misprint or a 5 misread as 3; not confirmed by magnification.}
In the second member $\Gamma(0.65)$ is unknown; to find it, a second application of the formula will be necessary, and we shall have
\[ \Gamma(0.65)=\frac{2^{-0.30}\sqrt{\pi}}{\sin 63^{\circ}}\cdot\frac{\Gamma(0.85)}{\Gamma(0.70)}; \]
a third application will give
\[ \Gamma(0.70)=\frac{2^{-0.40}\sqrt{\pi}}{\sin 54^{\circ}}\cdot\frac{\Gamma(0.8)}{\Gamma(0.6)}; \]
finally a fourth
\[ \Gamma(0.60)=\frac{2^{-0.20}\sqrt{\pi}}{\sin 72^{\circ}}\cdot\frac{\Gamma(0.9)}{\Gamma(0.8)}, \]
whence, mounting back, one will conclude the value of $\Gamma(0.675)$ expressed in known quantities. This determination is somewhat long in this case, because 0.675 approaches very closely the limit $\frac{2}{3}$.

\origpage{286}

\subsection*{(63).}

Since we have $\Gamma(1)=1$ and $\Gamma(1+a)=a\Gamma(a)$, the function $\Gamma(a)$ is always determined exactly when $a$ is a whole number, and we shall have generally
\[ \Gamma(a)=1\cdot 2\cdot 3\ldots a-1; \]
whence one sees that the quantity $\Gamma(a)$, considered as a continuous function of $a$, is nothing other than the general term which would result from the interpolation of the series 1, 1, $1\cdot 2$, $1\cdot 2\cdot 3$, $1\cdot 2\cdot 3\cdot 4$, etc.

The first two terms of this series correspond to the values $a=1$, $a=2$, and since they are both equal to unity, it follows that in the interval from $a=1$ to $a=2$, the function $\Gamma(a)$ must become a \emph{maximum} or a \emph{minimum}.

One recognizes easily after some trials that it is the \emph{minimum} which takes place, and in this case we find
\[ \begin{aligned} a&=1.4616038,\\ \Gamma(a)&=0.8856033,\\ \log\Gamma(a)&=9.9472392. \end{aligned} \]

Past the value $a=2$, the function $\Gamma(a)$ increases more and more, since we have $\Gamma(2+a)=(1+a)\Gamma(1+a)$. Therefore the value we have just found is the smallest of all those which the function $\Gamma(a)$ can take from $a=0$ to $a=\infty$.

It follows from this that the best way of constructing a table of the values of $\Gamma(a)$ is to calculate it for the interval between $a=1$ and $a=2$. For in this interval the function varies only between the limits 1 and 0.8856033, so that the differences will be very small, and the table very easy to interpolate.

\subsection*{(64).}

After this observation, and in order to facilitate the use of the functions $\Gamma$, we join here a table of the logarithms of the function $\Gamma$, calculated for all the values of $a$, by thousandths, from $a=1.000$ to $a=2.000$.

This table is easy to interpolate for all intermediate values, and by means of the equation $\Gamma(1+a)=a\Gamma(a)$, its use will be extended to all the other cases.

Indeed, if $a$ lies between 0 and 1, we shall determine $\Gamma(a)$ by the equation $\Gamma(a)=\frac{1}{a}\Gamma(1+a)$.

\origpage{287}
If $a$ lies between 2 and 3, $\Gamma(a)$ will be determined by the equation $\Gamma(a)=(a-1)\Gamma(a-1)$. If $a$ lies between 3 and 4, $\Gamma(a)$ will be determined by the equation $\Gamma(a)=(a-1)(a-2)\Gamma(a-2)$, and so on.

\subsection*{(65).}

By means of this table, it will be easy to evaluate in all cases the transcendent $\left(\frac{p}{q}\right)$ which corresponds to a given value of $n$. For this purpose, the equations $(\lambda)$ being put under these two new forms,
\[ \begin{aligned} \left(\frac{p}{q}\right)&=\frac{n}{pq}\cdot\frac{\Gamma\left(1+\frac{p}{n}\right)\cdot\Gamma\left(1+\frac{q}{n}\right)}{\Gamma\left(\frac{p+q}{n}\right)}\\ \left(\frac{p}{q}\right)&=\frac{p+q}{pq}\cdot\frac{\Gamma\left(1+\frac{p}{n}\right)\cdot\Gamma\left(1+\frac{q}{n}\right)}{\Gamma\left(1+\frac{p+q}{n}\right)}, \end{aligned} \]
we shall use the first, if we have $p+q>n$, and the second if we have $p+q<n$.

Let it be proposed, for example, to find the value of the transcendent $Z=\left(\frac{3}{8}\right)$ when $a=10$\ednote{Printed $a=10$ instead of $n=10$; an apparent misprint. A reader has struck the letter and written n in the margin.}; we shall then make use of the first formula, which gives
\[ Z=\frac{\Gamma(1.300)\Gamma(1.800)}{2.4\Gamma(1.100)}; \]
and with the aid of the table, we find immediately
\[ \log Z=9.5635972,\quad\text{and}\quad Z=0.3660979. \]

\subsection*{(66).}

Let us now seek the value of $\log\Gamma(1+k)$, $k$ being supposed very small. This value could be deduced from the interpolation of the first three terms of the table; but it will be more exact to interpolate terms less close together, so that the second differences may become more sensible. Let us consider for this purpose the three terms which correspond to the values $a=1.000$, $a=1.005$, $a=1.010$; we shall deduce from them the first and second differences, as follows:

\origpage{288}
\[ \begin{aligned} a && \log\Gamma && \text{Diff. }1^{\text{st}}. && \text{Diff. }2^{\text{nd}}.\\ 1.000 && 0.0000000 && -12445 && +178\\ 1.005 && 9.9987555 && -12267 &&\\ 1.010 && 9.9975288 && && \end{aligned} \]
Let $\Delta'=-0.0012445$, $\Delta''=0.0000178$; we shall have, neglecting the third differences, $\log\Gamma\left(1+\frac{5x}{1000}\right)=x\Delta'+\frac{x(x-1)}{2}\Delta''$, and putting $x=200k$,
\[ \log\Gamma(1+k)=-k(0.25068)+k^{2}(0.356). \]
This logarithm is a vulgar logarithm; to make it hyperbolic, it must be multiplied by 2.3025, etc., which will give
\[ \log\Gamma(1+k)=-k(0.57721)+k^{2}(0.820). \]
This value is only approximate; let us suppose that the exact value is, still rejecting the terms affected by $k^{3}$ and by higher powers of $k$,
\[ \log\Gamma(1+k)=-Pk+Qk^{2}, \]
we shall deduce from it
\[ \Gamma(1+k)=1-Pk+(Q+\tfrac{1}{2}P^{2})k^{2}: \]
from this results, on changing the sign of $k$,
\[ \Gamma(1-k)=1+Pk+(Q+\tfrac{1}{2}P^{2})k^{2}. \]
But the equation $\Gamma(1+k)=k\Gamma(k)$ gives $\Gamma(k)=\frac{1}{k}\Gamma(1+k)$; therefore
\[ \Gamma(k)=\frac{1}{k}-P+(Q+\tfrac{1}{2}P^{2})k. \]
Multiplying together these last two equations, we shall have
\[ \Gamma(k)\Gamma(1-k)=\frac{1}{k}(1+2Qk^{2}): \]
but on the other hand we know that the first member $=\frac{\pi}{\sin\pi k}=\frac{1}{k}\left(1+\frac{\pi^{2}k^{2}}{6}\right)$; we have therefore exactly
\[ Q=\frac{\pi^{2}}{12}=0.822467. \]

\origpage{289}
We had found $Q=0.820$, but this value could only be approximate.

As to the value of $P$, we found it to be 0.57721; but pushing the approximation further, one would find
\[ P=0.5772156649, \]
as will be shown hereafter.

\subsection*{(67).}

We now know, then, in a very approximate manner the values of $\Gamma(1+k)$ and $\Gamma(k)$, when $k$ is very small: these values are
\[ \begin{aligned} \Gamma(1+k)&=1-Pk+\left(\tfrac{\pi^{2}}{12}+\tfrac{1}{2}P^{2}\right)k^{2},\\ \Gamma(k)&=\frac{1}{k}-P+\left(\tfrac{\pi^{2}}{12}+\tfrac{1}{2}P^{2}\right)k. \end{aligned} \]
By means of the last formula, it is easy to find the value of the transcendent $\left(\frac{p}{q}\right)$ when $p$ and $q$ are very small in relation to $n$. Indeed, if we substitute the values of $\Gamma\left(\frac{p}{n}\right)$, $\Gamma\left(\frac{q}{n}\right)$, $\Gamma\left(\frac{p+q}{n}\right)$, given by the preceding formula, in the equation $(\varepsilon)$, we shall find
\[ \left(\frac{p}{q}\right)=\frac{p+q}{pq}\left(1-\frac{\pi^{2}}{6}\cdot\frac{pq}{n^{2}}\right), \]
which agrees with the equation $(l')$.

\subsection*{(68).}

It remains to show how we constructed the table by means of which one finds so easily, in all cases, the value of the functions $\Gamma$ and $\left(\frac{p}{q}\right)$. The simplest method that can be proposed for this object is that which results from a formula given by Euler in his \emph{Calc.\ diff.}, page 465, and which we are going to relate.

If we call $S$ the sum of the series
\[ \log 1+\log 2+\log 3\ldots+\log k, \]
we shall have
\[ S=k\log k+\tfrac{1}{2}\log(2\pi k)-k+\frac{A'}{1\cdot2k}-\frac{B'}{3\cdot4k^{3}}+\frac{C'}{5\cdot6k^{5}}-\text{etc.} \]
$A'$, $B'$, $C'$, etc. being the Bernoullian numbers. Let then $e$ be the

\origpage{290}
number whose logarithm is 1, and $R$ a number such that we have
\[ \log R=\frac{A'}{1\cdot2k}-\frac{B'}{3\cdot4k^{3}}+\frac{C'}{5\cdot6k^{5}}-\text{etc.} \]
we shall have the product
\[ 1\cdot2\cdot3\ldots k=\left(\frac{k}{e}\right)^{k}\cdot(2\pi k)^{\frac{1}{2}}\cdot R.\qquad(\rho) \]
The first member is the value of $\Gamma(k+1)$, when $k$ is a whole number; and as the second member is a continuous function of $k$, we have generally, whatever $k$ may be,
\[ \Gamma(k+1)=\left(\frac{k}{e}\right)^{k}(2\pi k)^{\frac{1}{2}}R.\qquad(\sigma) \]
Such is the formula by which one can in all cases determine the approximate value of $\Gamma(k+1)$; but it is fitting to make some observations on this subject.

\subsection*{(69).}

The quantity $R$ can be developed according to the powers of $\frac{1}{k}$; for we have $R=1+\log R+\frac{1}{2}\log^{2}R+\frac{1}{2\cdot3}\log^{3}R+\text{etc.}$ Substituting therefore the value of $\log R$, and putting in place of the coefficients $A'$, $B'$, $C'$, etc. their known values $A'=\frac{1}{6}$, $B'=\frac{1}{30}$, $C'=\frac{1}{42}$, $D'=\frac{1}{30}$, etc., we shall have
\[ R=1+\frac{1}{12k}+\frac{1}{2(12k)^{2}}-\frac{139}{30(12k)^{3}}-\frac{571}{120(12k)^{4}}-\text{etc.} \]
In this series, if we call $M$ the part
\[ 1+\frac{1}{2(12k)^{2}}-\frac{571}{120(12k)^{4}}-\text{etc.} \]
where $k$ is raised to even powers, and $N$ the other part, we shall have $M^{2}-N^{2}=1$; so that one may take indifferently $R=M+N$, or $R=\frac{1}{M-N}$. Indeed, as all the powers of $k$ are odd in $\log R$, the change of sign of $k$ will give $\log(M+N)=-\log.(M-N)$, or $\log(M^{2}-N^{2})=0$. Therefore $M^{2}-N^{2}=1$.

\subsection*{(70).}

It is to be remarked that the series
\[ \frac{A'}{1\cdot2k}-\frac{B'}{3\cdot4k^{3}}+\frac{C'}{5\cdot6k^{3}}-\text{etc.} \]
\ednote{Printed $5\cdot6k^{3}$ in the last denominator instead of $5\cdot6k^{5}$; an apparent misprint, cf. the same series at the head of this page.}

\origpage{291}
even supposing $k$ sufficiently large, is convergent only over a certain number of the first terms; for it is known that the Bernoullian numbers, whose expressions are very irregular, increase continually, so that if $T$ and $V$ are two consecutive terms far from the beginning, the one of rank $n$, the other of rank $n+1$, the ratio $\frac{V}{T}$ has for its limit $\frac{n^{2}}{\pi^{2}}$. This sequence, which begins by being convergent for a fairly large number of terms, especially if $k$ is somewhat large, therefore ends by being divergent, and would give a value of $\log R$ all the more faulty, the more terms one took beyond those where it ceases to be convergent.

From this one sees that for a given value of $k$ there is a term which must not be passed in the calculation of the sequence
\[ \frac{A'}{1\cdot2k}-\frac{B'}{3\cdot4k^{3}}+\text{etc.} \]
The term at which one must stop is the one that would be followed by a larger term; the approximation can then go no further; but it will be as extensive as one wishes, by taking $k$ sufficiently large.

The same would hold for the series
\[ R=1+\frac{1}{12k}+\frac{1}{2(12k)^{2}}-\text{etc.}; \]
but this one is not of such easy use as the series
\[ \frac{A'}{1\cdot2k}-\frac{B'}{3\cdot4k^{3}}+\text{etc.} \]
whose law is manifest, and which depends only on the Bernoullian numbers.

One can fix \emph{a priori} the number of terms after which the sequence
\[ \frac{A'}{1\cdot2k}-\frac{B'}{3\cdot4k^{3}}+\text{etc.} \]
ceases to be convergent; for, considering the two consecutive terms
\[ \pm\frac{T^{(n)}}{2n\cdot2n-1\cdot k^{2n-1}}\mp\frac{T^{(n+1)}}{2n+2\cdot2n+1\cdot k^{2n+1}}, \]
and supposing them equal, we shall have
\[ \frac{T^{(n+1)}}{T^{(n)}}=\frac{2n+2\cdot2n+1}{2n\cdot2n-1}k^{2}; \]

\origpage{292}
but the larger $n$ is, the more the first member approaches the limit $\frac{2n+2\cdot2n+1}{4\pi^{2}}$ (\emph{Eul.\ Calc.\ diff.}, page 429). Hence we shall have very nearly $n=\pi k$. Thus, putting $k=5$, we have $n=15$ or 16, that is to say the series ceases to be convergent about the 15\textsuperscript{th} term; if we put $k=10$, the series would cease to be convergent only about the 31\textsuperscript{st} term, and so on.

\subsection*{(71).}

At the same time one can have the measure of the degree of approximation that one can obtain with a given value of $k$. Indeed, if we call $\Omega$ the $n^{\text{th}}$ term of the sequence
\[ \frac{A'}{1\cdot2k^{2}}-\frac{B'}{3\cdot4k^{3}}+\text{etc.}, \]
\ednote{Printed $1\cdot2k^{2}$ in the first denominator instead of $1\cdot2k$; an apparent misprint.}
we shall have
\[ \Omega=\frac{T^{(n)}}{2n(2n-1)k^{2n-1}}, \]
and as we have very nearly
\[ T^{(n)}=\frac{1\cdot2\cdot3\ldots2n}{2^{2n-1}\pi^{2n}}, \]
we may put
\[ \Omega=\frac{1\cdot2\cdot3\ldots2n-2}{\pi(2\pi k)^{2n-1}}. \]
This value, by means of the formula $(\rho)$, becomes
\[ \Omega=\frac{\left(\frac{2n-2}{e}\right)^{2n-\frac{3}{2}}(2\pi e)^{\frac{1}{2}}}{\pi(2\pi k)^{2n-1}}, \]
and putting $n$ in place of $\pi k$, we easily deduce from it
\[ \log\Omega=-2n-\tfrac{1}{2}\log(\pi n). \]
Thus, putting $k=5$ and $n=16$, we shall have $\log\Omega=-33.96$. To this hyperbolic logarithm corresponds the common logarithm $-14.75$, so that we have $\Omega=10^{-14.75}$. Therefore by means of the first 16 terms of the sequence $\frac{A'}{1\cdot2k}-\frac{B'}{3\cdot4k^{3}}+\text{etc.}$, we shall have the value of $\log R$ approximated to about 15 decimal places. If we put $k=10$, we could have 29 or 30 exact decimal places, and so on.

\subsection*{(72).}

This theory is easy to verify, since whenever $k$ is an integer, the value of $\Gamma(k+1)$ is exactly $1\cdot2\cdot3\ldots k$.

\origpage{293}

Let, for example, $k=3$; it follows from the preceding formulas that the series equal to $\log R$ will cease to be convergent after a number of terms $n=k\pi=9$ or 10, and that the number of exact decimal places obtained by these nine or ten terms will be 8 or 9.

Indeed, the true value of $\log R$ is deduced from the equation $6=\left(\frac{3}{e}\right)^{3}(6\pi)^{\frac{1}{2}}R$, which gives
\[ \log R=0.02767\ 79256\ 86. \]
This same value, deduced from the sequence
\[ \log R=\frac{A'}{1\cdot2k}-\frac{B'}{3\cdot4k^{3}}+\frac{C'}{5\cdot6k^{5}}-\text{etc.}, \]
is found by calculating successively the different terms as follows:
\[ \begin{aligned} 1^{\text{st}}\ \text{term} && + && 0.02777\ 77777\ 78\\ 2^{\text{nd}} && - && 0.00010\ 28806\ 58\\ && && 0.02767\ 48971\ 20\\ 3^{\text{rd}} && + && 32660\ 53\\ && && 0.02767\ 81631\ 73\\ 4^{\text{th}} && - && 2721\ 71\\ && && 0.02767\ 78910\ 02\\ 5^{\text{th}} && + && 427\ 65\\ && && 0.02767\ 79337\ 67\\ 6^{\text{th}} && - && 108\ 24\\ && && 0.02767\ 79229\ 43\\ 7^{\text{th}} && + && 40\ 21\\ && && 0.02767\ 79269\ 64\\ 8^{\text{th}} && - && 20\ 59\\ && && 0.02767\ 79249\ 05\\ 9^{\text{th}} && + && 13\ 91\\ && && 0.02767\ 79262\ 96\\ 10^{\text{th}} && - && 11\ 98\\ && && 0.02767\ 79250\ 98\\ 11^{\text{th}} && + && 12\ 81\\ && && 0.02767\ 79263\ 79 \end{aligned} \]

\origpage{294}

We see that, in conformity with the formula, the series ceases to be convergent after the 10\textsuperscript{th} term, and that the value of $\log R$ deduced from it must lie between $0.02767\ 7926296$ and $0.02767\ 7925098$, which gives, by a mean,
\[ \log R=0.02767\ 7925697, \]
a value exact almost to the eleventh decimal. But by continuing the sequence further, one would depart more and more from the true result.

This example sets in full light the manner of drawing all possible advantage, for approximations, from \emph{semi-convergent} sequences, that is to say sequences which are convergent in the first terms, and which afterwards become divergent.

\subsection*{(73).}

By means of the formula $(\sigma)$, one can develop in series the function $\Gamma(k)$ when $k$ is very small. For this, let us observe first that $\Gamma(k+1)=k\Gamma(k)$, and that thus we shall have
\[ \Gamma(k)=c^{-k}k^{k-\frac{1}{2}}(2\pi)^{\frac{1}{2}}R; \]
\ednote{The base of the first factor reads $c$ instead of $e$ cf. formula $(\sigma)$; a misprint, or an $e$ whose bar failed to print.}
whence
\[ \log\Gamma(k)=(k-\tfrac{1}{2})\log k-k+\tfrac{1}{2}l(2\pi)+\frac{A'}{1\cdot2k}-\frac{B'}{3\cdot4k^{3}}+\text{etc.}\ldots(\tau) \]
This formula can serve only for values of $k$ greater than unity; but if one puts $1+k$ in place of $k$, and in place of the first member, which will become $\log\Gamma(1+k)$, one puts its value $\log k+\log\Gamma(k)$, one will draw from it anew
\[ \begin{aligned} \log\Gamma(k)={}&-\log k+(\tfrac{1}{2}+k)\log(1+k)-1-k+\tfrac{1}{2}\log(2\pi)\\ &+\frac{A'}{1\cdot2(1+k)}-\frac{B'}{3\cdot4(1+k)^{3}}+\frac{C'}{5\cdot6(1+k)^{5}}-\text{etc.} \end{aligned} \]
and the development being carried out up to quantities of the order $k^{2}$ exclusively, we shall have
\[ \begin{aligned} \log\Gamma(k)={}&-\log k+\tfrac{1}{2}\log(2\pi)-1+\frac{A'}{1\cdot2}-\frac{B'}{3\cdot4}+\frac{C'}{5\cdot6}-\text{etc.}\\ &-k\left(\tfrac{1}{2}+\frac{A'}{2}-\frac{B'}{4}+\frac{C'}{6}-\text{etc.}\right) \end{aligned} \]
When we put $k=0$, we must have $\log\Gamma(k)=-\log k$, because $k\Gamma(k)=\Gamma(1+k)$, and because, putting $k=0$, the second member

\origpage{295}
$=\Gamma(1)=1$; therefore we must have
\[ \tfrac{1}{2}\log(2\pi)-1+\frac{A'}{1\cdot2}-\frac{B'}{3\cdot4}+\frac{C'}{5\cdot6}-\text{etc.}=0. \]
And indeed, Euler found this equality, page 466 of his \emph{Calc.\ diff.} This being laid down, the preceding value reduces to this,
\[ \log\Gamma(k)=-\log k-k\left(\tfrac{1}{2}+\frac{A'}{2}-\frac{B'}{4}+\frac{C'}{6}-\text{etc.}\right); \]
now in the cited work, page 444, one still finds the equality
\[ C=\tfrac{1}{2}+\frac{A'}{2}-\frac{B'}{4}+\frac{C'}{6}-\text{etc.}, \]
$C$ being a constant whose value, calculated with precision by another route, is
\[ C=0.5772156649015325; \]
therefore finally we shall have, $k$ being very small,
\[ \log\Gamma(k)=-\log k-Ck. \]
We found above, by pushing the approximation further,
\[ \Gamma(k)=\frac{1}{k}\left[1-Pk+\left(\tfrac{1}{2}P^{2}+\frac{\pi^{2}}{12}\right)k^{2}\right]; \]
from this we see that $P=C$, and that thus $P$ is nothing other than the constant $C$, whose value has just been given, approximated to sixteen decimal places.

We now know, therefore, the value of $\Gamma(k)$ very closely when $k$ is very small, and one could approach it still more closely by pushing the development further.

\subsection*{(74).}

But here are considerations which lead more generally and more directly to the same end. Let
\[ \begin{aligned} M&=1+\tfrac{1}{2}+\tfrac{1}{3}\ldots+\frac{1}{x},\\ N&=\frac{1}{x+1}+\frac{1}{x+2}+\frac{1}{x+3}+\text{etc.}, \end{aligned} \]
this last sequence being continued to infinity.

We shall regard the quantity $M$ as a continuous function

\origpage{296}
of $x$, since indeed we have (\emph{Calc.\ diff.}, pag.\ 443),
\[ M=lx+C+\frac{1}{2x}-\frac{A'}{2x^{2}}+\frac{B'}{4x^{4}}-\frac{C'}{6x^{6}}+\text{etc.}; \]
a formula which is fit to give the approximate value of $M$, whatever $x$ may be, provided that one supposes $x$ greater than unity. This same formula would give the value of $M$ when $x$ is smaller than unity; for, representing by $M(x)$ and by $M(x+1)$ similar functions of $x$ and of $x+1$, we evidently have
\[ M(x)=M(x+1)-\frac{1}{x+1}, \]
or again
\[ \begin{aligned} M(x)&=M(x+2)-\frac{1}{x+1}-\frac{1}{x+2},\\ &\text{etc.}; \end{aligned} \]
so that, $x$ being smaller than unity, we shall have the function $M(x)$ or $M$ by means of a similar function in which $x$ is increased by several units, and which will render the preceding sequence convergent enough in its first terms that it can give all the approximation one may desire.

\subsection*{(75).}

This being laid down, if $x$ becomes $x+\omega$ ($\omega$ being smaller than unity), the functions $M$ and $N$, which may be represented by $M(x)$ and $N(x)$, will become $M(x+\omega)$ and $N(x+\omega)$. Now I say that the sum $M(x)+N(x)=M(x+\omega)+N(x+\omega)$, and that the two sums are equal to one and the same constant.

Indeed, if the sequence which has $N(x)$ for its sum, instead of being continued to infinity, were continued only to the term $\frac{1}{x+m}$, $m$ being a very large number, the sum $M(x)+N(x)$ and the sum $M(x+\omega)+N(x+\omega)$ could differ from each other only by a quantity less than $\frac{1}{x+m+1}$, since this difference is the one that holds when $\omega=1$. Now $m$ being very large, the difference $\frac{1}{x+m+1}$ is reckoned null; all the more will it be so when the sequence $N(x)$ is continued to infinity. We shall therefore have
\[ M+N=\text{const}=C'; \]

\origpage{297}
but since the value of $N$ is
\[ \frac{1}{1+x}+\frac{1}{2+x}+\frac{1}{3+x}+\text{etc.}, \]
if we develop each of these fractions on the hypothesis that $x$ is smaller than unity, we shall have
\[ \begin{aligned} N={}&1+\tfrac{1}{2}+\tfrac{1}{3}+\tfrac{1}{4}+\text{etc.}\\ &-x\left(1+\frac{1}{2^{2}}+\frac{1}{3^{2}}+\frac{1}{4^{2}}+\text{etc.}\right)\\ &+x^{2}\left(1+\frac{1}{2^{3}}+\frac{1}{3^{3}}+\frac{1}{4^{3}}+\text{etc.}\right)\\ &-\text{etc.} \end{aligned} \]
The first part $1+\frac{1}{2}+\frac{1}{3}+\frac{1}{4}+\text{etc.}$ is nothing other than the constant $C'$; therefore we shall have
\[ M=S_{2}x-S_{3}x^{2}+S_{4}x^{3}-S_{5}x^{4}+\text{etc.},\qquad(\upsilon) \]
$S_{n}$ representing in general the sum of the reciprocal powers of degree $n$ of the natural numbers.

\subsection*{(76).}

But on the hypothesis $x>1$, we have, from equation $(\tau)$,
\[ \frac{d\log\Gamma(x)}{dx}=\log x-\frac{1}{2x}-\frac{A'}{2x^{2}}+\frac{B'}{4x^{4}}-\text{etc.}, \]
and on the same hypothesis we have
\[ M=\log x+C+\frac{1}{2x}-\frac{A'}{2x^{2}}+\frac{B'}{4x^{4}}-\text{etc.} \]
Therefore
\[ \frac{d\log\Gamma(x)}{dx}=-\frac{1}{x}+M-C,\qquad(\varphi) \]
an equation which must hold whatever $x$ may be, since $M$ may be regarded as a continuous function of $x$.

If now we suppose $x<1$, which will permit the use of the series $(\upsilon)$, we shall have
\[ \frac{d\log\Gamma(x)}{dx}=-\frac{1}{x}-C+S_{2}x-S_{3}x^{2}+S_{4}x^{3}-\text{etc.}; \]
whence results, on integrating,
\[ \log\Gamma(x)=-\log x-Cx+\tfrac{1}{2}S_{2}x^{2}-\tfrac{1}{3}S_{3}x^{3}+\tfrac{1}{4}S_{4}x^{4}-\text{etc.}\qquad(\psi) \]

\origpage{298}
We add no constant, because $x\Gamma(x)$ or $\Gamma(1+x)$ reduces to unity when $x=0$.

\subsection*{(77).}

The formula $(\psi)$, which gives a complete development of $\log\Gamma(x)$, will serve to express the value of any proposed function $\Gamma(z)$, since this function can always be brought back to the case where $z$ is smaller than $\frac{1}{2}$, and even to that where $z$ is smaller than $\frac{1}{4}$.

From the formula $(\psi)$ we deduce immediately the following,
\[ \log\Gamma(1+x)=-Cx+\tfrac{1}{2}S_{2}x^{2}-\tfrac{1}{3}S_{3}x^{3}+\tfrac{1}{4}S_{4}x^{4}-\text{etc.},\qquad(\omega) \]
and changing in this the sign of $x$, we should have
\[ \log\Gamma(1-x)=Cx+\tfrac{1}{2}S_{2}x^{2}+\tfrac{1}{3}S_{3}x^{3}+\tfrac{1}{4}S_{4}x^{4}+\text{etc.},\qquad(\alpha') \]
formulas equally fit to make us find in each proposed case the value of $\Gamma(z)$, since this function can always be brought either to the form $\Gamma(1+x)$ or to the form $\Gamma(1-x)$, $x$ not exceeding $\frac{1}{4}$.

Thus we have regular and always convergent series for calculating the value of a given function $\Gamma(z)$; they suppose only the use of the quantities $S_{2}$, $S_{3}$, $S_{4}$, etc., whose numerical values Euler has given very closely as far as $S_{15}$ (\emph{Calc.\ diff.}, page 456).

We must observe that the formula $(\omega)$ amounts to the one seen in the cited work, page 800; it appears only that Euler did not perceive the use that this series could have in the determination of the functions $\Gamma(x)$ with which he occupied himself in other places in his works.

Let us remark further that from the formula $(\omega)$ one can deduce this value of $C$,
\[ C=-\frac{1}{x}\log\Gamma(1+x)+\frac{1}{2}S_{2}x-\frac{1}{3}S_{3}x^{2}+\frac{1}{4}S_{4}x^{3}-\text{etc.}, \]
whence one sees that it suffices to know a particular value of $\Gamma(1+x)$, and that we shall have the value of $C$ expressed by a series all the more convergent as $x$ is smaller. If we put $x=\frac{1}{2}$, we shall have
\[ \Gamma\left(1+\tfrac{1}{2}\right)=\tfrac{1}{2}\Gamma\tfrac{1}{2}=\tfrac{1}{2}\sqrt{\pi}, \]

\origpage{299}
whence results
\[ C=\log\frac{4}{\pi}+\tfrac{1}{2}S_{2}\cdot\tfrac{1}{2}-\tfrac{1}{3}S_{3}\cdot\tfrac{1}{4}+\tfrac{1}{4}S_{4}\cdot\tfrac{1}{8}-\text{etc.} \]

\subsection*{(78).}

The two equations $\Gamma(1+x)=x\Gamma(x)$, and $\Gamma(x)\Gamma(1-x)=\frac{\pi}{\sin\pi x}$, give
\[ \Gamma(1+x)\Gamma(1-x)=\frac{\pi x}{\sin\pi x};\qquad(\beta') \]
taking the logarithms of the two members, and substituting the values given by the formulas $(\omega)$ and $(\alpha')$, we shall have
\[ \log\left(\frac{\pi x}{\sin\pi x}\right)=S_{2}x^{2}+\tfrac{1}{2}S_{4}x^{4}+\tfrac{1}{3}S_{6}x^{6}+\text{etc.}, \]
a known formula, which by its differential serves to determine the values of the quantities $S_{2}$, $S_{4}$, $S_{6}$, etc.

One can make use of this formula to render still more convergent the series contained in the equations $(\omega)$ and $(\alpha')$: we shall thus have
\[ \left.\begin{aligned} \log\Gamma(1+x)&=\tfrac{1}{2}\log\left(\frac{\pi x}{\sin\pi x}\right)-Cx-\tfrac{1}{3}S_{3}x^{3}-\tfrac{1}{5}S_{5}x^{5}-\text{etc.}\\ \log\Gamma(1-x)&=\tfrac{1}{2}\log\left(\frac{\pi x}{\sin\pi x}\right)+Cx+\tfrac{1}{3}S_{3}x^{3}+\tfrac{1}{5}S_{5}x^{5}+\text{etc.} \end{aligned}\right\}\qquad(\gamma') \]
The latter, putting $x=\frac{1}{2}$, gives
\[ C=L2-\tfrac{1}{3}S_{3}\cdot\tfrac{1}{4}-\tfrac{1}{5}S_{5}\cdot\tfrac{1}{16}-\text{etc.} \]
a value more convergent than that of art.\ 77.

\subsection*{(79).}

Having the developed expression of $\log\Gamma(x)$ by the formula $(\psi)$, one can likewise have that of the logarithm of the function $\left(\frac{p}{q}\right)$; for since we have found
\[ \left(\frac{p}{q}\right)=\frac{\Gamma\left(\frac{p}{n}\right)\Gamma\left(\frac{p}{n}\right)}{n\Gamma\left(\frac{p+q}{n}\right)}, \]
\ednote{The second factor of the numerator is printed $\Gamma\left(\frac{p}{n}\right)$; read $\Gamma\left(\frac{q}{n}\right)$.}
if we take the logarithms of each member, and designate by $Z$ the function $\left(\frac{p}{q}\right)$ which corresponds to a given value of $n$, we

\origpage{300}
shall have
\[ \left.\begin{aligned} \log Z=\log.\left(\frac{p+q}{pq}\right)&-\tfrac{1}{2}S_{2}\cdot\frac{(p+q)^{2}-p^{2}-q^{2}}{n^{2}}\\ &+\tfrac{1}{3}S_{3}\cdot\frac{(p+q)^{3}-p^{3}-q^{3}}{n^{3}}\\ &-\tfrac{1}{4}S_{4}\cdot\frac{(p+q)^{4}-p^{4}-q^{4}}{n^{4}}\\ &+\text{etc.} \end{aligned}\right\}\qquad(\delta') \]
The first two terms of this formula agree with the equation $(l')$; but here we see the general law of the development, which one can continue at will, and which gives a series all the more convergent as $p$ and $q$ are smaller relatively to $n$.

Thus for the function $\left(\frac{a}{a}\right)$ designated above by $M_{a}$, we shall have
\[ \log.M_{a}=\log\left(\frac{2}{a}\right)-\tfrac{1}{2}S_{2}(2^{2}-2)\frac{a^{2}}{n^{2}}+\tfrac{1}{3}S_{3}(2^{3}-2)\frac{a^{3}}{n^{3}}-\text{etc.} \]

\subsection*{(80).}

From these formulas it was easy to calculate the table here annexed of the values of $\log\Gamma(a)$. For this, we put $k=4+a$ in the formula $(\tau)$ (one could equally have taken $k=3+a$, $k=5+a$, etc.). The first member then giving the value of $\log\Gamma(4+a)$, we deduced from it $\log\Gamma(a)$ by the known relation between these quantities; namely:
\[ \Gamma(4+a)=(3+a)(2+a)(1+a)a\Gamma(a). \]
We have therefore had to calculate the formula
\[ \begin{aligned} \log\Gamma(a)={}&\left(k-\tfrac{1}{2}\right)\log k+\tfrac{1}{2}l(2\pi)-mk+\frac{mA'}{1\cdot2k}-\frac{mB'}{3\cdot4k^{3}}+\frac{mC'}{5\cdot6k^{5}}-\text{etc.}\\ &-\log[a(1+a)(2+a)(3+a)], \end{aligned} \]
in which we have introduced the factor $m=0.43429$, etc. in order to reduce everything to the logarithms of the tables.

In this manner we never had need to calculate more than two or three terms of the series $\frac{mA'}{1\cdot2k}-\frac{mB'}{3\cdot4k^{3}}+\frac{mC'}{5\cdot6k^{5}}-\text{etc.}$, to have $\log\Gamma(a)$ approximated to seven decimal places, throughout the whole interval from $a=1$ to $a=2$.

\subsection*{(81).}

We shall remark in closing that the integrals of the form

\origpage{301}
$\displaystyle\int e^{-t^{m}}t^{n}\,dt$, taken from $t=0$ to $t=\infty$, can be brought back to the functions $\Gamma(a)$. Indeed, let first $t^{n+1}=z$, and $\frac{m}{n+1}=\alpha$; the preceding integral will become $\displaystyle\int\frac{m}{n+1}\cdot e^{-z^{\alpha}}\,dz$,\ednote{Printed $\frac{m}{n+1}$; the substitution gives the factor $\frac{1}{n+1}$; an apparent misprint.} this one again having to be taken from $z=0$ to $z=\infty$.

From this one sees that the proposed integral loses nothing of its generality by putting $n=0$, and that thus one may propose simply the integral $\displaystyle\int e^{-t^{m}}\,dt$. If in this one puts $e^{-t^{m}}=x$, or $t=\left(\log\frac{1}{x}\right)^{\frac{1}{m}}$, the transformed integral will be $\displaystyle\frac{1}{m}\int dx\left(\log\frac{1}{x}\right)^{\frac{1}{m}-1}$, and this new integral will have to be taken from $x=0$ to $x=1$; we shall therefore have generally
\[ \int e^{-t^{m}}\,dt=\frac{1}{m}\Gamma\left(\frac{1}{m}\right). \]

Thus the integrals in question offer no new species of transcendents, and relate to the functions $\Gamma$.

It is not useless for the history of science to observe that the integral $\displaystyle\int dx\left(\log\frac{1}{x}\right)^{-\frac{1}{2}}=\sqrt{\pi}$, which one finds in § 28 of Euler's memoir, printed in vol.\ XVI of the \emph{Novi Comm.\ Petrop.}, had been given long before by the same author, in vol.\ V of the old Mémoires of Petersburg, pag.\ 44. It is therefore to this epoch that the discovery of the integral $\displaystyle\int e^{-x^{2}}\,dx=\frac{1}{2}\sqrt{\pi}$ goes back, taken between the limits $x=0$, $x=\infty$, since the simple substitution $e^{-x^{2}}=z$ suffices to bring the integral $\displaystyle\int e^{-x^{2}}\,dx$ to the form $\displaystyle\frac{1}{2}\int dz\left(\log\frac{1}{z}\right)^{-\frac{1}{2}}$.

END OF THE SECOND PART.

\origpage{302}

\section*{TABLE OF THE LOGARITHMS OF THE FUNCTION $\Gamma(a)$.}

\[ \begin{aligned} a && \text{Log }\Gamma &\qquad& a && \text{Log }\Gamma. &\qquad& a. && \text{Log }\Gamma. &\qquad& a. && \text{Log }\Gamma.\\ 1.000 && 0.0000000 && 1.050 && 9.9883379 && 1.100 && 9.9783407 && 1.150 && 9.9699006\\ 1.001 && 9.9997497 && 1.051 && 9.9881220 && 1.101 && 9.9781570 && 1.151 && 9.9697470\\ 1.002 && 9.9995001 && 1.052 && 9.9879068 && 1.102 && 9.9779739 && 1.152 && 9.9695941\\ 1.003 && 9.9992512 && 1.053 && 9.9876922 && 1.103 && 9.9777914 && 1.153 && 9.9694417\\ 1.004 && 9.9990030 && 1.054 && 9.9874783 && 1.104 && 9.9776095 && 1.154 && 9.9692898\\ 1.005 && 9.9987555 && 1.055 && 9.9872651 && 1.105 && 9.9774283 && 1.155 && 9.9691386\\ 1.006 && 9.9985088 && 1.056 && 9.9870525 && 1.106 && 9.9772477 && 1.156 && 9.9689879\\ 1.007 && 9.9982627 && 1.057 && 9.9868406 && 1.107 && 9.9770677 && 1.157 && 9.9688378\\ 1.008 && 9.9980174 && 1.058 && 9.9866293 && 1.108 && 9.9768883 && 1.158 && 9.9686883\\ 1.009 && 9.9977727 && 1.059 && 9.9864187 && 1.109 && 9.9767096 && 1.159 && 9.9685394\\ 1.010 && 9.9975288 && 1.060 && 9.9862088 && 1.110 && 9.9765314 && 1.160 && 9.9683910\\ 1.011 && 9.9972855 && 1.061 && 9.9859996 && 1.111 && 9.9763539 && 1.161 && 9.9682432\\ 1.012 && 9.9970430 && 1.062 && 9.9857910 && 1.112 && 9.9761769 && 1.162 && 9.9680960\\ 1.013 && 9.9968011 && 1.063 && 9.9855830 && 1.113 && 9.9760006 && 1.163 && 9.9679494\\ 1.014 && 9.9965599 && 1.064 && 9.9853758 && 1.114 && 9.9758249 && 1.164 && 9.9678033\\ 1.015 && 9.9963195 && 1.065 && 9.9851691 && 1.115 && 9.9756498 && 1.165 && 9.9676578\\ 1.016 && 9.9960798 && 1.066 && 9.9849631 && 1.116 && 9.9754753 && 1.166 && 9.9675129\\ 1.017 && 9.9958408 && 1.067 && 9.9847577 && 1.117 && 9.9753014 && 1.167 && 9.9673686\\ 1.018 && 9.9956025 && 1.068 && 9.9845530 && 1.118 && 9.9751281 && 1.168 && 9.9672248\\ 1.019 && 9.9953648 && 1.069 && 9.9843489 && 1.119 && 9.9749555 && 1.169 && 9.9670816\\ 1.020 && 9.9951279 && 1.070 && 9.9841455 && 1.120 && 9.9747834 && 1.170 && 9.9669390\\ 1.021 && 9.9948917 && 1.071 && 9.9839427 && 1.121 && 9.9746120 && 1.171 && 9.9667969\\ 1.022 && 9.9946561 && 1.072 && 9.9837406 && 1.122 && 9.9744411 && 1.172 && 9.9666554\\ 1.023 && 9.9944213 && 1.073 && 9.9835391 && 1.123 && 9.9742709 && 1.173 && 9.9665145\\ 1.024 && 9.9941871 && 1.074 && 9.9833383 && 1.124 && 9.9741012 && 1.174 && 9.9663742\\ 1.025 && 9.9939536 && 1.075 && 9.9831381 && 1.125 && 9.9739322 && 1.175 && 9.9662344\\ 1.026 && 9.9937208 && 1.076 && 9.9829386 && 1.126 && 9.9737638 && 1.176 && 9.9660952\\ 1.027 && 9.9934887 && 1.077 && 9.9827398 && 1.127 && 9.9735960 && 1.177 && 9.9659566\\ 1.028 && 9.9932572 && 1.078 && 9.9825415 && 1.128 && 9.9734287 && 1.178 && 9.9658186\\ 1.029 && 9.9930265 && 1.079 && 9.9823440 && 1.129 && 9.9732621 && 1.179 && 9.9656810\\ 1.030 && 9.9927964 && 1.080 && 9.9821470 && 1.130 && 9.9730961 && 1.180 && 9.9655441\\ 1.031 && 9.9925670 && 1.081 && 9.9819507 && 1.131 && 9.9729307 && 1.181 && 9.9654077\\ 1.032 && 9.9923383 && 1.082 && 9.9817550 && 1.132 && 9.9727659 && 1.182 && 9.9652718\\ 1.033 && 9.9921103 && 1.083 && 9.9815599 && 1.133 && 9.9726016 && 1.183 && 9.9651366\\ 1.034 && 9.9918830 && 1.084 && 9.9813655 && 1.134 && 9.9724380 && 1.184 && 9.9650019\\ 1.035 && 9.9916564 && 1.085 && 9.9811717 && 1.135 && 9.9722750 && 1.185 && 9.9648677\\ 1.036 && 9.9914305 && 1.086 && 9.9809786 && 1.136 && 9.9721126 && 1.186 && 9.9647341\\ 1.037 && 9.9912052 && 1.087 && 9.9807860 && 1.137 && 9.9719508 && 1.187 && 9.9646011\\ 1.038 && 9.9909807 && 1.088 && 9.9805942 && 1.138 && 9.9717896 && 1.188 && 9.9644687\\ 1.039 && 9.9907567 && 1.089 && 9.9804029 && 1.139 && 9.9716289 && 1.189 && 9.9643367\\ 1.040 && 9.9905335 && 1.090 && 9.9802123 && 1.140 && 9.9714689 && 1.190 && 9.9642054\\ 1.041 && 9.9903109 && 1.091 && 9.9800223 && 1.141 && 9.9713095 && 1.191 && 9.9640746\\ 1.042 && 9.9900890 && 1.092 && 9.9798330 && 1.142 && 9.9711506 && 1.192 && 9.9639444\\ 1.043 && 9.9898678 && 1.093 && 9.9796443 && 1.143 && 9.9709923 && 1.193 && 9.9638147\\ 1.044 && 9.9896472 && 1.094 && 9.9794562 && 1.144 && 9.9708346 && 1.194 && 9.9636856\\ 1.045 && 9.9894273 && 1.095 && 9.9792687 && 1.145 && 9.9706775 && 1.195 && 9.9635570\\ 1.046 && 9.9892081 && 1.096 && 9.9790819 && 1.146 && 9.9705209 && 1.196 && 9.9634290\\ 1.047 && 9.9889895 && 1.097 && 9.9788956 && 1.147 && 9.9703650 && 1.197 && 9.9633016\\ 1.048 && 9.9887717 && 1.098 && 9.9787100 && 1.148 && 9.9702096 && 1.198 && 9.9631747\\ 1.049 && 9.9885545 && 1.099 && 9.9785251 && 1.149 && 9.9700548 && 1.199 && 9.9630483\\ 1.050 && 9.9883379 && 1.100 && 9.9783407 && 1.150 && 9.9699006 && 1.200 && 9.9629225 \end{aligned} \]

\origpage{303}

\[ \begin{aligned} a && \text{Log }\Gamma &\qquad& a && \text{Log }\Gamma. &\qquad& a. && \text{Log }\Gamma. &\qquad& a. && \text{Log }\Gamma.\\ 1.200 && 9.9629225 && 1.250 && 9.9573211 && 1.300 && 9.9530204 && 1.350 && 9.9499515\\ 1.201 && 9.9627972 && 1.251 && 9.9572226 && 1.301 && 9.9529472 && 1.351 && 9.9499022\\ 1.202 && 9.9626725 && 1.252 && 9.9571246 && 1.302 && 9.9528744 && 1.352 && 9.9498535\\ 1.203 && 9.9625484 && 1.253 && 9.9570271 && 1.303 && 9.9528021 && 1.353 && 9.9498051\\ 1.204 && 9.9624248 && 1.254 && 9.9569302 && 1.304 && 9.9527304 && 1.354 && 9.9497573\\ 1.205 && 9.9623017 && 1.255 && 9.9568337 && 1.305 && 9.9526591 && 1.355 && 9.9497099\\ 1.206 && 9.9621792 && 1.256 && 9.9567377 && 1.306 && 9.9525883 && 1.356 && 9.9496630\\ 1.207 && 9.9620573 && 1.257 && 9.9566423 && 1.307 && 9.9525180 && 1.357 && 9.9496165\\ 1.208 && 9.9619359 && 1.258 && 9.9565474 && 1.308 && 9.9524482 && 1.358 && 9.9495706\\ 1.209 && 9.9618150 && 1.259 && 9.9564530 && 1.309 && 9.9523789 && 1.359 && 9.9495251\\ 1.210 && 9.9616947 && 1.260 && 9.9563591 && 1.310 && 9.9523101 && 1.360 && 9.9494800\\ 1.211 && 9.9615749 && 1.261 && 9.9562657 && 1.311 && 9.9522417 && 1.361 && 9.9494354\\ 1.212 && 9.9614556 && 1.262 && 9.9561729 && 1.312 && 9.9521739 && 1.362 && 9.9493913\\ 1.213 && 9.9613369 && 1.263 && 9.9560805 && 1.313 && 9.9521065 && 1.363 && 9.9493476\\ 1.214 && 9.9612187 && 1.264 && 9.9559887 && 1.314 && 9.9520396 && 1.364 && 9.9493044\\ 1.215 && 9.9611011 && 1.265 && 9.9558974 && 1.315 && 9.9519732 && 1.365 && 9.9492617\\ 1.216 && 9.9609840 && 1.266 && 9.9558066 && 1.316 && 9.9519073 && 1.366 && 9.9492194\\ 1.217 && 9.9608675 && 1.267 && 9.9557163 && 1.317 && 9.9518419 && 1.367 && 9.9491776\\ 1.218 && 9.9607516 && 1.268 && 9.9556266 && 1.318 && 9.9517769 && 1.368 && 9.9491363\\ 1.219 && 9.9606361 && 1.269 && 9.9555373 && 1.319 && 9.9517125 && 1.369 && 9.9490953\\ 1.220 && 9.9605212 && 1.270 && 9.9554486 && 1.320 && 9.9516485 && 1.370 && 9.9490549\\ 1.221 && 9.9604068 && 1.271 && 9.9553604 && 1.321 && 9.9515850 && 1.371 && 9.9490149\\ 1.222 && 9.9602930 && 1.272 && 9.9552727 && 1.322 && 9.9515220 && 1.372 && 9.9489754\\ 1.223 && 9.9601797 && 1.273 && 9.9551855 && 1.323 && 9.9514594 && 1.373 && 9.9489363\\ 1.224 && 9.9599669 && 1.274 && 9.9550988 && 1.324 && 9.9513974 && 1.374 && 9.9488977\\ 1.225 && 9.9599547 && 1.275 && 9.9550126 && 1.325 && 9.9513358 && 1.375 && 9.9488595\\ 1.226 && 9.9598430 && 1.276 && 9.9549269 && 1.326 && 9.9512747 && 1.376 && 9.9488218\\ 1.227 && 9.9597318 && 1.277 && 9.9548417 && 1.327 && 9.9512141 && 1.377 && 9.9487845\\ 1.228 && 9.9596212 && 1.278 && 9.9547570 && 1.328 && 9.9511540 && 1.378 && 9.9487477\\ 1.229 && 9.9595111 && 1.279 && 9.9546728 && 1.329 && 9.9510944 && 1.379 && 9.9487114\\ 1.230 && 9.9594015 && 1.280 && 9.9545891 && 1.330 && 9.9510352 && 1.380 && 9.9486755\\ 1.231 && 9.9592925 && 1.281 && 9.9545059 && 1.331 && 9.9509765 && 1.381 && 9.9486401\\ 1.232 && 9.9591839 && 1.282 && 9.9544232 && 1.332 && 9.9509183 && 1.382 && 9.9486051\\ 1.233 && 9.9590760 && 1.283 && 9.9543410 && 1.333 && 9.9508605 && 1.383 && 9.9485706\\ 1.234 && 9.9589685 && 1.284 && 9.9542594 && 1.334 && 9.9508033 && 1.384 && 9.9485365\\ 1.235 && 9.9588616 && 1.285 && 9.9541782 && 1.335 && 9.9507465 && 1.385 && 9.9485029\\ 1.236 && 9.9587552 && 1.286 && 9.9540976 && 1.336 && 9.9506902 && 1.386 && 9.9484698\\ 1.237 && 9.9586494 && 1.287 && 9.9540174 && 1.337 && 9.9506344 && 1.387 && 9.9484371\\ 1.238 && 9.9585441 && 1.288 && 9.9539377 && 1.338 && 9.9505791 && 1.388 && 9.9484048\\ 1.239 && 9.9584393 && 1.289 && 9.9538586 && 1.339 && 9.9505242 && 1.389 && 9.9483730\\ 1.240 && 9.9583350 && 1.290 && 9.9537799 && 1.340 && 9.9504698 && 1.390 && 9.9483417\\ 1.241 && 9.9582313 && 1.291 && 9.9537017 && 1.341 && 9.9504159 && 1.391 && 9.9483108\\ 1.242 && 9.9581280 && 1.292 && 9.9536240 && 1.342 && 9.9503624 && 1.392 && 9.9482804\\ 1.243 && 9.9580253 && 1.293 && 9.9535468 && 1.343 && 9.9503094 && 1.393 && 9.9482504\\ 1.244 && 9.9579232 && 1.294 && 9.9534701 && 1.344 && 9.9502569 && 1.394 && 9.9482208\\ 1.245 && 9.9578215 && 1.295 && 9.9533939 && 1.345 && 9.9502048 && 1.395 && 9.9481917\\ 1.246 && 9.9577204 && 1.296 && 9.9533182 && 1.346 && 9.9501532 && 1.396 && 9.9481631\\ 1.247 && 9.9576198 && 1.297 && 9.9532430 && 1.347 && 9.9501021 && 1.397 && 9.9481349\\ 1.278 && 9.9575197 && 1.298 && 9.9531683 && 1.348 && 9.9500514 && 1.398 && 9.9481071\\ 1.249 && 9.9574201 && 1.299 && 9.9530941 && 1.349 && 9.9500012 && 1.399 && 9.9480798\\ 1.250 && 9.9573211 && 1.300 && 9.9530204 && 1.350 && 9.9499515 && 1.400 && 9.9480529 \end{aligned} \]

\ednote{In the row $a=1.224$ the table prints $9.9599669$ instead of $9.9600669$, and the argument $1.248$ is printed $1.278$; apparent misprints.}

\origpage{304}

\[ \begin{aligned} a. && \text{Log }\Gamma. &\qquad& a. && \text{Log }\Gamma. &\qquad& a. && \text{Log }\Gamma. &\qquad& a. && \text{Log }\Gamma.\\ 1.400 && 9.9480529 && 1.450 && 9.9472677 && 1.500 && 9.9475449 && 1.550 && 9.9488375\\ 1.401 && 9.9480264 && 1.451 && 9.9472630 && 1.501 && 9.9475610 && 1.551 && 9.9488734\\ 1.402 && 9.9480004 && 1.452 && 9.9472587 && 1.502 && 9.9475775 && 1.552 && 9.9489097\\ 1.403 && 9.9479748 && 1.453 && 9.9472548 && 1.503 && 9.9475943 && 1.553 && 9.9489464\\ 1.404 && 9.9479497 && 1.454 && 9.9472514 && 1.504 && 9.9476116 && 1.554 && 9.9489835\\ 1.405 && 9.9479250 && 1.455 && 9.9472484 && 1.505 && 9.9476293 && 1.555 && 9.9490209\\ 1.406 && 9.9479008 && 1.456 && 9.9472458 && 1.506 && 9.9476474 && 1.556 && 9.9490587\\ 1.407 && 9.9478770 && 1.457 && 9.9472436 && 1.507 && 9.9476659 && 1.557 && 9.9490969\\ 1.408 && 9.9478536 && 1.458 && 9.9472419 && 1.508 && 9.9476847 && 1.558 && 9.9491355\\ 1.409 && 9.9478307 && 1.459 && 9.9472406 && 1.509 && 9.9477040 && 1.559 && 9.9491745\\ 1.410 && 9.9478083 && 1.460 && 9.9472397 && 1.510 && 9.9477237 && 1.560 && 9.9492139\\ 1.411 && 9.9477863 && 1.461 && 9.9472392 && 1.511 && 9.9477438 && 1.561 && 9.9492536\\ 1.412 && 9.9477648 && 1.462 && 9.9472392 && 1.512 && 9.9477643 && 1.562 && 9.9492938\\ 1.413 && 9.9477437 && 1.463 && 9.9472396 && 1.513 && 9.9477851 && 1.563 && 9.9493343\\ 1.414 && 9.9477230 && 1.464 && 9.9472404 && 1.514 && 9.9478064 && 1.564 && 9.9493752\\ 1.415 && 9.9477028 && 1.465 && 9.9472416 && 1.515 && 9.9478281 && 1.565 && 9.9494165\\ 1.416 && 9.9476830 && 1.466 && 9.9472432 && 1.516 && 9.9478502 && 1.566 && 9.9494582\\ 1.417 && 9.9476637 && 1.467 && 9.9472453 && 1.517 && 9.9478726 && 1.567 && 9.9495003\\ 1.418 && 9.9476447 && 1.468 && 9.9472477 && 1.518 && 9.9478955 && 1.568 && 9.9495428\\ 1.419 && 9.9476263 && 1.469 && 9.9472506 && 1.519 && 9.9479189 && 1.569 && 9.9495857\\ 1.420 && 9.9476082 && 1.470 && 9.9472539 && 1.520 && 9.9479425 && 1.570 && 9.9496289\\ 1.421 && 9.9475906 && 1.471 && 9.9472576 && 1.521 && 9.9479666 && 1.571 && 9.9496725\\ 1.422 && 9.9475734 && 1.472 && 9.9472618 && 1.522 && 9.9479911 && 1.572 && 9.9497165\\ 1.423 && 9.9475566 && 1.473 && 9.9472664 && 1.523 && 9.9480160 && 1.573 && 9.9497608\\ 1.424 && 9.9475403 && 1.474 && 9.9472713 && 1.524 && 9.9480413 && 1.574 && 9.9498056\\ 1.425 && 9.9475244 && 1.475 && 9.9472767 && 1.525 && 9.9480670 && 1.575 && 9.9498507\\ 1.426 && 9.9475089 && 1.476 && 9.9472825 && 1.526 && 9.9480931 && 1.576 && 9.9498962\\ 1.427 && 9.9474939 && 1.477 && 9.9472887 && 1.527 && 9.9481196 && 1.577 && 9.9499421\\ 1.428 && 9.9474794 && 1.478 && 9.9472953 && 1.528 && 9.9481465 && 1.578 && 9.9499884\\ 1.429 && 9.9474652 && 1.479 && 9.9473023 && 1.529 && 9.9481737 && 1.579 && 9.9500350\\ 1.430 && 9.9474515 && 1.480 && 9.9473097 && 1.530 && 9.9482014 && 1.580 && 9.9500821\\ 1.431 && 9.9474382 && 1.481 && 9.9473175 && 1.531 && 9.9482295 && 1.581 && 9.9501295\\ 1.432 && 9.9474254 && 1.482 && 9.9473258 && 1.532 && 9.9482579 && 1.582 && 9.9501773\\ 1.433 && 9.9474130 && 1.483 && 9.9473345 && 1.533 && 9.9482868 && 1.583 && 9.9502255\\ 1.434 && 9.9474010 && 1.484 && 9.9473436 && 1.534 && 9.9483160 && 1.584 && 9.9502741\\ 1.435 && 9.9473895 && 1.485 && 9.9473531 && 1.535 && 9.9483457 && 1.585 && 9.9503230\\ 1.436 && 9.9473784 && 1.486 && 9.9473631 && 1.536 && 9.9483757 && 1.586 && 9.9503723\\ 1.437 && 9.9473677 && 1.487 && 9.9473734 && 1.537 && 9.9484062 && 1.587 && 9.9504220\\ 1.438 && 9.9473574 && 1.488 && 9.9473842 && 1.538 && 9.9484370 && 1.588 && 9.9504721\\ 1.439 && 9.9473476 && 1.489 && 9.9473953 && 1.539 && 9.9484683 && 1.589 && 9.9505225\\ 1.440 && 9.9473382 && 1.490 && 9.9474069 && 1.540 && 9.9484999 && 1.590 && 9.9505733\\ 1.441 && 9.9473292 && 1.491 && 9.9474189 && 1.541 && 9.9485319 && 1.591 && 9.9506245\\ 1.442 && 9.9473207 && 1.492 && 9.9474312 && 1.542 && 9.9485643 && 1.592 && 9.9506761\\ 1.443 && 9.9473126 && 1.493 && 9.9474440 && 1.543 && 9.9485971 && 1.593 && 9.9507280\\ 1.444 && 9.9473049 && 1.494 && 9.9474572 && 1.544 && 9.9486302 && 1.594 && 9.9507803\\ 1.445 && 9.9472976 && 1.495 && 9.9474708 && 1.545 && 9.9486638 && 1.595 && 9.9508330\\ 1.446 && 9.9472908 && 1.496 && 9.9474848 && 1.546 && 9.9486978 && 1.596 && 9.9508860\\ 1.447 && 9.9472844 && 1.497 && 9.9474992 && 1.547 && 9.9487321 && 1.597 && 9.9509395\\ 1.448 && 9.9472784 && 1.498 && 9.9475140 && 1.548 && 9.9487669 && 1.598 && 9.9509933\\ 1.449 && 9.9472728 && 1.499 && 9.9475293 && 1.549 && 9.9488020 && 1.599 && 9.9510474\\ 1.450 && 9.9472677 && 1.500 && 9.9475449 && 1.550 && 9.9488375 && 1.600 && 9.9511020 \end{aligned} \]

\origpage{305}

\[ \begin{aligned} a. && \text{Log }\Gamma. &\qquad& a. && \text{Log }\Gamma. &\qquad& a. && \text{Log }\Gamma. &\qquad& a. && \text{Log }\Gamma.\\ 1.600 && 9.9511020 && 1.650 && 9.9542989 && 1.700 && 9.9583911 && 1.750 && 9.9633451\\ 1.601 && 9.9511569 && 1.651 && 9.9543721 && 1.701 && 9.9584819 && 1.751 && 9.9634527\\ 1.602 && 9.9512123 && 1.652 && 9.9544456 && 1.702 && 9.9585730 && 1.752 && 9.9635607\\ 1.603 && 9.9512679 && 1.653 && 9.9545195 && 1.703 && 9.9586644 && 1.753 && 9.9636689\\ 1.604 && 9.9513240 && 1.654 && 9.9545938 && 1.704 && 9.9587562 && 1.754 && 9.9637775\\ 1.605 && 9.9513804 && 1.655 && 9.9546684 && 1.705 && 9.9588483 && 1.755 && 9.9638865\\ 1.606 && 9.9514372 && 1.656 && 9.9547434 && 1.706 && 9.9589408 && 1.756 && 9.9639958\\ 1.607 && 9.9514943 && 1.657 && 9.9548187 && 1.707 && 9.9590335 && 1.757 && 9.9641054\\ 1.608 && 9.9515518 && 1.658 && 9.9548944 && 1.708 && 9.9591267 && 1.758 && 9.9642154\\ 1.609 && 9.9516097 && 1.659 && 9.9549704 && 1.709 && 9.9592202 && 1.759 && 9.9643257\\ 1.610 && 9.9516680 && 1.660 && 9.9550468 && 1.710 && 9.9593140 && 1.760 && 9.9644363\\ 1.611 && 9.9517267 && 1.661 && 9.9551236 && 1.711 && 9.9594082 && 1.761 && 9.9645473\\ 1.612 && 9.9517857 && 1.662 && 9.9552007 && 1.712 && 9.9595027 && 1.762 && 9.9646586\\ 1.613 && 9.9518451 && 1.663 && 9.9552782 && 1.713 && 9.9595975 && 1.763 && 9.9647702\\ 1.614 && 9.9519049 && 1.664 && 9.9553560 && 1.714 && 9.9596928 && 1.764 && 9.9648821\\ 1.615 && 9.9519650 && 1.665 && 9.9554342 && 1.715 && 9.9597883 && 1.765 && 9.9649944\\ 1.616 && 9.9520255 && 1.666 && 9.9555127 && 1.716 && 9.9598842 && 1.766 && 9.9651070\\ 1.617 && 9.9520863 && 1.667 && 9.9565916 && 1.717 && 9.9599805 && 1.767 && 9.9652199\\ 1.618 && 9.9521476 && 1.668 && 9.9556709 && 1.718 && 9.9600770 && 1.768 && 9.9653331\\ 1.619 && 9.9522092 && 1.669 && 9.9557505 && 1.719 && 9.9601739 && 1.769 && 9.9654467\\ 1.620 && 9.9522711 && 1.670 && 9.9558304 && 1.720 && 9.9602712 && 1.770 && 9.9655606\\ 1.621 && 9.9523334 && 1.671 && 9.9559107 && 1.721 && 9.9603688 && 1.771 && 9.9656748\\ 1.622 && 9.9523961 && 1.672 && 9.9559914 && 1.722 && 9.9604667 && 1.772 && 9.9657894\\ 1.623 && 9.9524591 && 1.673 && 9.9560724 && 1.723 && 9.9605650 && 1.773 && 9.9659043\\ 1.624 && 9.9525225 && 1.674 && 9.9561537 && 1.724 && 9.9606636 && 1.774 && 9.9660195\\ 1.625 && 9.9525863 && 1.675 && 9.9562354 && 1.725 && 9.9607625 && 1.775 && 9.9661350\\ 1.626 && 9.9526504 && 1.676 && 9.9563174 && 1.726 && 9.9608618 && 1.776 && 9.9662509\\ 1.627 && 9.9527149 && 1.677 && 9.9563998 && 1.727 && 9.9609614 && 1.777 && 9.9663671\\ 1.628 && 9.9527798 && 1.678 && 9.9564826 && 1.728 && 9.9610614 && 1.778 && 9.9664836\\ 1.629 && 9.9528450 && 1.679 && 9.9565657 && 1.729 && 9.9611617 && 1.779 && 9.9666004\\ 1.630 && 9.9529106 && 1.680 && 9.9566491 && 1.730 && 9.9612623 && 1.780 && 9.9667176\\ 1.631 && 9.9529766 && 1.681 && 9.9567329 && 1.731 && 9.9613632 && 1.781 && 9.9668351\\ 1.632 && 9.9530429 && 1.682 && 9.9568171 && 1.732 && 9.9614645 && 1.782 && 9.9669529\\ 1.633 && 9.9531096 && 1.683 && 9.9569016 && 1.733 && 9.9615661 && 1.783 && 9.9670711\\ 1.634 && 9.9531767 && 1.684 && 9.9569864 && 1.734 && 9.9616681 && 1.784 && 9.9671895\\ 1.635 && 9.9532441 && 1.685 && 9.9570716 && 1.735 && 9.9617704 && 1.785 && 9.9673083\\ 1.636 && 9.9533119 && 1.686 && 9.9571572 && 1.736 && 9.9618731 && 1.786 && 9.9674274\\ 1.637 && 9.9533801 && 1.687 && 9.9572430 && 1.737 && 9.9619760 && 1.787 && 9.9675468\\ 1.638 && 9.9534487 && 1.688 && 9.9573293 && 1.738 && 9.9620794 && 1.788 && 9.9676665\\ 1.639 && 9.9535176 && 1.689 && 9.9574159 && 1.739 && 9.9621830 && 1.789 && 9.9677866\\ 1.640 && 9.9535868 && 1.690 && 9.9575028 && 1.740 && 9.9622870 && 1.790 && 9.9679070\\ 1.641 && 9.9536564 && 1.691 && 9.9575901 && 1.741 && 9.9623913 && 1.791 && 9.9680277\\ 1.642 && 9.9537263 && 1.692 && 9.9576777 && 1.742 && 9.9624960 && 1.792 && 9.9681487\\ 1.643 && 9.9537966 && 1.693 && 9.9577656 && 1.743 && 9.9626009 && 1.793 && 9.9682701\\ 1.644 && 9.9538673 && 1.694 && 9.9578539 && 1.744 && 9.9627063 && 1.794 && 9.9683918\\ 1.645 && 9.9539383 && 1.695 && 9.9579426 && 1.745 && 9.9628119 && 1.795 && 9.9685138\\ 1.646 && 9.9540097 && 1.696 && 9.9580316 && 1.746 && 9.9629179 && 1.796 && 9.9686361\\ 1.647 && 9.9540815 && 1.697 && 9.9581210 && 1.747 && 9.9630242 && 1.797 && 9.9687588\\ 1.648 && 9.9541536 && 1.698 && 9.9582107 && 1.748 && 9.9631308 && 1.798 && 9.9688818\\ 1.649 && 9.9542261 && 1.699 && 9.9583007 && 1.749 && 9.9632378 && 1.799 && 9.9690051\\ 1.650 && 9.9542989 && 1.700 && 9.9583911 && 1.750 && 9.9633451 && 1.800 && 9.9691287 \end{aligned} \]

\ednote{In the row $a=1.667$ the table prints $9.9565916$ instead of $9.9555916$; an apparent misprint.}

\origpage{306}

\[ \begin{aligned} a. && \text{Log }\Gamma. &\qquad& a. && \text{Log }\Gamma. &\qquad& a. && \text{Log }\Gamma. &\qquad& a. && \text{Log }\Gamma.\\ 1.800 && 9.9691287 && 1.850 && 9.9757126 && 1.900 && 9.9830693 && 1.950 && 9.9911732\\ 1.801 && 9.9692526 && 1.851 && 9.9758522 && 1.901 && 9.9832241 && 1.951 && 9.9913427\\ 1.802 && 9.9693769 && 1.852 && 9.9759922 && 1.902 && 9.9833793 && 1.952 && 9.9915125\\ 1.803 && 9.9695015 && 1.853 && 9.9761325 && 1.903 && 9.9835347 && 1.953 && 9.9916826\\ 1.804 && 9.9696264 && 1.854 && 9.9762730 && 1.904 && 9.9836905 && 1.954 && 9.9918530\\ 1.805 && 9.9697516 && 1.855 && 9.9764139 && 1.905 && 9.9838465 && 1.955 && 9.9920237\\ 1.806 && 9.9698771 && 1.856 && 9.9765551 && 1.906 && 9.9840028 && 1.956 && 9.9921947\\ 1.807 && 9.9700030 && 1.857 && 9.9766966 && 1.907 && 9.9841594 && 1.957 && 9.9923659\\ 1.808 && 9.9701291 && 1.858 && 9.9768384 && 1.908 && 9.9843164 && 1.958 && 9.9925374\\ 1.809 && 9.9702556 && 1.859 && 9.9769805 && 1.909 && 9.9844736 && 1.959 && 9.9927093\\ 1.810 && 9.9703824 && 1.860 && 9.9771229 && 1.910 && 9.9846311 && 1.960 && 9.9928814\\ 1.811 && 9.9705095 && 1.861 && 9.9772657 && 1.911 && 9.9847889 && 1.961 && 9.9930538\\ 1.812 && 9.9706369 && 1.862 && 9.9774087 && 1.912 && 9.9849471 && 1.962 && 9.9932266\\ 1.813 && 9.9707647 && 1.863 && 9.9775521 && 1.913 && 9.9851055 && 1.963 && 9.9933996\\ 1.814 && 9.9708927 && 1.864 && 9.9776957 && 1.914 && 9.9852642 && 1.964 && 9.9935728\\ 1.815 && 9.9710211 && 1.865 && 9.9778397 && 1.915 && 9.9854232 && 1.965 && 9.9937464\\ 1.816 && 9.9711498 && 1.866 && 9.9779840 && 1.916 && 9.9855825 && 1.966 && 9.9939202\\ 1.817 && 9.9712788 && 1.867 && 9.9781286 && 1.917 && 9.9857421 && 1.967 && 9.9940943\\ 1.818 && 9.9714081 && 1.868 && 9.9782734 && 1.918 && 9.9859020 && 1.968 && 9.9942687\\ 1.819 && 9.9715377 && 1.869 && 9.9784186 && 1.919 && 9.9860621 && 1.969 && 9.9944434\\ 1.820 && 9.9716677 && 1.870 && 9.9785641 && 1.920 && 9.9862226 && 1.970 && 9.9946184\\ 1.821 && 9.9717980 && 1.871 && 9.9787099 && 1.921 && 9.9863834 && 1.971 && 9.9947937\\ 1.822 && 9.9719286 && 1.872 && 9.9788559 && 1.922 && 9.9865444 && 1.972 && 9.9949692\\ 1.823 && 9.9720595 && 1.873 && 9.9790023 && 1.923 && 9.9867058 && 1.973 && 9.9951451\\ 1.824 && 9.9721907 && 1.874 && 9.9791490 && 1.924 && 9.9868675 && 1.974 && 9.9953212\\ 1.825 && 9.9723223 && 1.875 && 9.9792960 && 1.925 && 9.9870294 && 1.975 && 9.9954976\\ 1.826 && 9.9724542 && 1.876 && 9.9794433 && 1.926 && 9.9871916 && 1.976 && 9.9956743\\ 1.827 && 9.9725864 && 1.877 && 9.9795909 && 1.927 && 9.9873542 && 1.977 && 9.9958513\\ 1.828 && 9.9727189 && 1.878 && 9.9797389 && 1.928 && 9.9875170 && 1.978 && 9.9960286\\ 1.829 && 9.9728517 && 1.879 && 9.9798871 && 1.929 && 9.9876801 && 1.979 && 9.9962062\\ 1.830 && 9.9729848 && 1.880 && 9.9800356 && 1.930 && 9.9878435 && 1.980 && 9.9963840\\ 1.831 && 9.9731182 && 1.881 && 9.9801844 && 1.931 && 9.9880072 && 1.981 && 9.9965621\\ 1.832 && 9.9732520 && 1.882 && 9.9803335 && 1.932 && 9.9881712 && 1.982 && 9.9967405\\ 1.833 && 9.9733860 && 1.883 && 9.9804829 && 1.933 && 9.9883355 && 1.983 && 9.9969192\\ 1.834 && 9.9735204 && 1.884 && 9.9806327 && 1.934 && 9.9885001 && 1.984 && 9.9970982\\ 1.835 && 9.9736551 && 1.885 && 9.9807827 && 1.935 && 9.9886650 && 1.985 && 9.9972774\\ 1.836 && 9.9737901 && 1.886 && 9.9809330 && 1.936 && 9.9888302 && 1.986 && 9.9974569\\ 1.837 && 9.9739254 && 1.887 && 9.9810837 && 1.937 && 9.9889956 && 1.987 && 9.9976367\\ 1.838 && 9.9740610 && 1.888 && 9.9812346 && 1.938 && 9.9891614 && 1.988 && 9.9978168\\ 1.839 && 9.9741970 && 1.889 && 9.9813859 && 1.939 && 9.9893275 && 1.989 && 9.9979972\\ 1.840 && 9.9743332 && 1.890 && 9.9815374 && 1.940 && 9.9894938 && 1.990 && 9.9981779\\ 1.841 && 9.9744697 && 1.891 && 9.9816892 && 1.941 && 9.9896604 && 1.991 && 9.9983589\\ 1.842 && 9.9746066 && 1.892 && 9.9818414 && 1.942 && 9.9898274 && 1.992 && 9.9985401\\ 1.843 && 9.9747437 && 1.893 && 9.9819938 && 1.943 && 9.9899946 && 1.993 && 9.9987216\\ 1.844 && 9.9748812 && 1.894 && 9.9821466 && 1.944 && 9.9901621 && 1.994 && 9.9989034\\ 1.845 && 9.9750190 && 1.895 && 9.9822996 && 1.945 && 9.9903299 && 1.995 && 9.9990855\\ 1.846 && 9.9751571 && 1.896 && 9.9824529 && 1.946 && 9.9904980 && 1.996 && 9.9992678\\ 1.847 && 9.9752955 && 1.897 && 9.9826066 && 1.947 && 9.9906664 && 1.997 && 9.9994504\\ 1.848 && 9.9754342 && 1.898 && 9.9827605 && 1.948 && 9.9908350 && 1.998 && 9.9996333\\ 1.849 && 9.9755733 && 1.899 && 9.9829148 && 1.949 && 9.9910040 && 1.999 && 9.9998165\\ 1.850 && 9.9757126 && 1.900 && 9.9830693 && 1.950 && 9.9911732 && 2.000 && 0.0000000 \end{aligned} \]

\origpage{307}

\section*{THIRD PART.}

\section*{ON QUADRATURES.}

Various methods relating to quadratures are the subject of this third part. These methods naturally complete the theory of transcendents, since if the transcendents are too complicated, or if they cannot be reduced to those given by the Tables, one must necessarily have recourse to quadratures in order to evaluate them in particular cases.

The method that I propose as the most general and the surest consists in expressing the sought integral, taken in infinitely small differences, by an integral in finite differences to which one adds the corrections that analysis indicates and which serve to direct the approximation. In this way one can arrive at results whose exactness extends to whatever order of decimals one wishes.

The same method, considered from another point of view, can serve to construct a curve whose coordinates each depend on a particular quadrature. As an example of this kind I have given the calculation of the trajectory of a projectile in a resisting medium, and by this means I have carried the approximation further than I had done in the piece crowned by the Academy of Berlin in 1782.

One finds in the Mémoires de l'Académie des Sciences, for the years 1778 and 1782, a very ingenious method for obtaining the value of the integral $\displaystyle\int y\,dx$, in the case where the function $y$ vanishes at the two limits of the integral. I have set out this method, with the developments that it requires, in some examples, and principally

\origpage{308}
in its application to the case where the limits of the integral are supposed imaginary.

Finally, I have taken advantage of the liberty that the title of this work gave me to treat of various kinds of definite integrals which belong to the theory of transcendents. Among these are several formulas that Laplace has given in various collections, and which hold a distinguished rank in this theory; the others are, for the most part, taken from the works of Euler, and particularly from the Supplements which make up volume IV of his \emph{Calcul intégral}.

I flatter myself that these various subjects, brought together under a single point of view, may interest the Geometers; they have moreover been chosen in such a way that each of them may offer new formulas or new demonstrations.

\section*{General formula for quadratures.}

\subsection*{(1).}

A problem is considered to have been solved when it is reduced to quadratures, that is, when it no longer depends on anything but one or several integrals of the form $\displaystyle\int y\,dx$, where $y$ is known as a function of $x$.

If the integrals in question are capable of being expressed by convergent series, the solution is as complete as one can desire, having regard to the nature of the question; but most often one has expressions too complicated to integrate by convergent series, and there remains no other course than to seek their values in particular cases only, that is, for a given part of the line of abscissas.

Let it therefore be proposed to find the value of the integral $\displaystyle Z=\int y\,dx$, from $x=a$ to $x=b$, or only from $x=0$ to $x=a$; for one may suppose that the origin of the abscissas is taken at the point where the area begins, so that $Z$ and $x$ vanish at the same time.

In order to obtain more easily the degree of approximation that one may desire, we shall suppose,

1°. That in the given interval from $x=0$ to $x=a$, the

\origpage{309}
function $y$ remains constantly of the same sign, and we shall regard this sign as positive. If it were otherwise, one would seek successively the different parts of the area $\displaystyle\int y\,dx$ situated on different sides of the line of abscissas, and one would subtract the sum of the negative areas from the sum of the positive areas.

2°. That the curvature of the curve, from $x=0$ to $x=a$, undergoes no variation abrupt enough for one of the coefficients $\frac{dy}{dx}$, $\frac{ddy}{dx^{2}}$, $\frac{d^{3}y}{dx^{3}}$ to become infinite. Thus, throughout the whole extent of the portion of the curve that we wish to square, the ordinate corresponding to an abscissa $x+\alpha$, very little different from $x$, can be expressed with sufficient exactness by the series $y+\alpha\frac{dy}{dx}+\frac{\alpha^{2}}{2}\frac{ddy}{dx^{2}}+\frac{\alpha^{3}}{2\cdot3}\cdot\frac{d^{3}y}{dx^{3}}+\text{etc.}$, without there being room for exception at any point. This condition is above all necessary at the two limits of the integral; it is less important at the intermediate points, because the error due to this cause is bounded, and can be corrected fairly easily, as will be seen below.

This being laid down, let $y=F(x)$, and $x=n\omega$, $n$ being a whole number which is the larger the greater the approximation one wishes to obtain; if one denotes by $\Sigma F(x+\frac{1}{2}\omega)$ the sum of the series
\[ F(\tfrac{1}{2}\omega)+F(\tfrac{3}{2}\omega)+F(\tfrac{5}{2}\omega)\ldots+F(x-\tfrac{1}{2}\omega), \]
it is evident that one will have at first, very nearly, $Z=\omega\Sigma F(x+\frac{1}{2}\omega)$. To have a more exact value, let
\[ Z=\omega\Sigma F(x+\tfrac{1}{2}\omega)+\zeta, \]
one will have, by making $x$ vary by $\omega$, and taking the finite differences of each member,
\[ \Delta\zeta=\Delta Z-\omega F(x+\tfrac{1}{2}\omega). \]
Developing the second member according to Taylor's theorem, and observing that $\frac{dZ}{dx}=y=F(x)$, one will have
\[ \begin{aligned} \Delta\zeta={}&\omega F+\frac{\omega^{2}}{2}\cdot\frac{dF}{dx}+\frac{\omega^{3}}{2\cdot3}\cdot\frac{ddF}{dx^{2}}+\frac{\omega^{4}}{2\cdot3\cdot4}\cdot\frac{d^{3}F}{dx^{3}}+\text{etc.}\\ &-\omega F-\omega^{2}\cdot\frac{1}{2}\frac{dF}{dx}-\frac{\omega^{3}}{2}\cdot\frac{1}{4}\cdot\frac{ddF}{dx^{2}}-\frac{\omega^{4}}{2\cdot3}\cdot\frac{1}{8}\cdot\frac{d^{3}F}{dx^{3}}-\text{etc.}, \end{aligned} \]

\origpage{310}
or, on reducing,
\[ \Delta\zeta=\frac{\omega^{3}}{2}\left(\frac{1}{3}-\frac{1}{4}\right)\frac{ddF}{dx^{2}}+\frac{\omega^{4}}{2\cdot3}\left(\frac{1}{4}-\frac{1}{8}\right)\frac{d^{3}F}{dx^{3}}+\frac{\omega^{5}}{2\cdot3\cdot4}\left(\frac{1}{5}-\frac{1}{16}\right)\frac{d^{4}F}{dx^{4}}+\text{etc.} \]
One sees that the first term of the value of $\zeta$ is $\frac{\omega^{2}}{2}\cdot\frac{1}{12}\cdot\frac{dF}{dx}$; so that one may put $\zeta=\frac{\omega^{2}}{24}\cdot\frac{dF}{dx}+\zeta'$; one would then find that the first term of the value of $\zeta'$ is $-\frac{7\omega^{4}}{5760}\cdot\frac{d^{3}F}{dx^{3}}$. But in order to know the law of the series which expresses $\zeta$, we shall leave the coefficients of its different terms undetermined, and we shall suppose
\[ \begin{aligned} Z={}&\omega\Sigma F(x+\tfrac{1}{2}\omega)+\text{const.}\\ &+a'\omega^{2}\frac{dF}{dx}+b'\omega^{3}\frac{ddF}{dx^{2}}+c'\omega^{4}\frac{d^{3}F}{dx^{3}}+f'\omega^{5}\frac{d^{4}F}{dx^{4}}+\text{etc.} \end{aligned} \]

\subsection*{(2).}

It will suffice to determine the coefficients $a'$, $b'$, $c'$, etc. in a particular case; let therefore $F(x)=e^{x}$, one will have $\frac{dF}{dx}=e^{x}$, $\frac{ddF}{dx^{2}}=e^{x}$, etc., $\displaystyle Z=\int e^{x}\,dx=e^{x}-1$, $\Sigma F(x+\frac{1}{2}\omega)=e^{\frac{1}{2}\omega}\Sigma e^{x}=\frac{\omega e^{\frac{1}{2}\omega}}{e^{\omega}-1}(e^{x}-1)$,\ednote{The factor $\omega$ in this numerator belongs to $\omega\Sigma F(x+\frac{1}{2}\omega)$, the quantity substituted next, not to $\Sigma F(x+\frac{1}{2}\omega)$; as printed.} and the substitution of these values will give the following equation, which must be identical:
\[ e^{x}-1=\frac{\omega(e^{x}-1)}{e^{\frac{1}{2}\omega}-e^{-\frac{1}{2}\omega}}+\text{const.}+(a'\omega^{2}+b'\omega^{3}+c'\omega^{4}+\text{etc.})e^{x}. \]
Putting $x=0$, one finds the constant $=-a'\omega^{2}-b'\omega^{3}-c'\omega^{4}-\text{etc.}$; so that, dividing the whole equation by $e^{x}-1$, there will result
\[ \frac{\omega}{e^{\frac{1}{2}\omega}-e^{-\frac{1}{2}\omega}}=1-a'\omega^{2}-b'\omega^{3}-c'\omega^{4}-f'\omega^{5}-\text{etc.} \]

The first member is an even function of $\omega$, since it remains the same when the sign of $\omega$ is changed; therefore in the second member all the coefficients of the odd powers of $\omega$ are zero. To take account of this circumstance, we shall put anew
\[ \frac{\omega}{e^{\frac{1}{2}\omega}-e^{-\frac{1}{2}\omega}}=1-A\omega^{2}+B\omega^{4}-C\omega^{6}+D\omega^{8}-\text{etc.}, \tag{1} \]

\origpage{311}
and the value of $Z$ will in general be
\[ \begin{aligned} Z={}&\omega\Sigma F(x+\tfrac{1}{2}\omega)+A\omega^{2}\frac{dF}{dx}-B\omega^{4}\frac{d^{3}F}{dx^{3}}+C\omega^{6}\frac{d^{5}F}{dx^{5}}-\text{etc.}\\ &+\text{const.} \end{aligned} \]
Let $\frac{dF^{0}}{dx}$, $\frac{ddF^{0}}{dx^{2}}$, etc., denote what the coefficients $\frac{dF}{dx}$, $\frac{ddF}{dx^{2}}$, etc. become when $x=0$; and the constant being determined so that $Z$ and $x$ vanish at the same time, one has
\[ \begin{aligned} Z={}&\omega\Sigma F(x+\tfrac{1}{2}\omega)+A\omega^{2}\left(\frac{dF}{dx}-\frac{dF^{0}}{dx^{0}}\right)\\ &-B\omega^{4}\left(\frac{d^{3}F}{dx^{3}}-\frac{d^{3}F^{0}}{dx^{3}}\right)\\ &+C\omega^{6}\left(\frac{d^{5}F}{dx^{5}}-\frac{d^{5}F^{0}}{dx^{5}}\right)\\ &-\text{etc.}: \end{aligned} \tag{2} \]
\ednote{Printed $dx^{0}$ in the denominator of the second fraction of the first line; $dx$ is meant, as in the other terms. An apparent misprint.}
this is the required integral taken from $x=0$ to $x=n\omega$.

\subsection*{(3).}

If in equation (1) one puts $-\omega^{2}$ in place of $\omega^{2}$, there will result
\[ \frac{\tfrac{1}{2}\omega}{\sin\tfrac{1}{2}\omega}=1+A\omega^{2}+B\omega^{4}+C\omega^{6}+D\omega^{8}+\text{etc.} \]
The first member can be put in the form
\[ 1+\frac{2\omega^{2}}{4\pi^{2}-\omega^{2}}-\frac{2\omega^{2}}{16\pi^{2}-\omega^{2}}+\frac{2\omega^{2}}{36\pi^{2}-\omega^{2}}-\text{etc.}, \]
(\emph{Introd. in An. inf.}, page 142), and by its development one has
\[ \begin{aligned} A&=\frac{2}{4\pi^{2}}\left(1-\frac{1}{2^{2}}+\frac{1}{3^{2}}-\frac{1}{4^{2}}+\text{etc.}\right)=\frac{1}{24}\\ B&=\frac{2}{16\pi^{4}}\left(1-\frac{1}{2^{4}}+\frac{1}{3^{4}}-\frac{1}{4^{4}}+\text{etc.}\right)=\frac{7}{5760}\\ C&=\frac{2}{64\pi^{6}}\left(1-\frac{1}{2^{6}}+\frac{1}{3^{6}}-\frac{1}{4^{6}}+\text{etc.}\right)=\frac{31}{967680}\\ D&=\frac{2}{256\pi^{8}}\left(1-\frac{1}{2^{8}}+\frac{1}{3^{8}}-\frac{1}{4^{8}}+\text{etc.}\right)=\frac{127}{(128)^{2}\cdot9450}\\ &\text{etc.} \end{aligned} \]

\origpage{312}
From which one sees that all the coefficients $A$, $B$, $C$, etc. are positive, and that they decrease according to a ratio which approaches more and more closely to $\frac{1}{4\pi^{2}}$, so that each term will be about the fortieth of the preceding.

The simplest way of calculating the numerical values of these coefficients is to deduce them from the development of the fraction
\[ \frac{1}{1-\frac{1}{2\cdot3}\cdot\frac{\omega^{2}}{4}+\frac{1}{2\cdot3\cdot4\cdot5}\cdot\frac{\omega^{4}}{16}-\text{etc.}}=1+A\omega^{2}+B\omega^{4}+C\omega^{6}+\text{etc.}, \]
from which one draws
\[ \begin{aligned} A&=\frac{1}{2\cdot3}\cdot\frac{1}{4}\\ B&=\frac{1}{2\cdot3}\cdot\frac{1}{4}A-\frac{1}{2\cdot3\cdot4\cdot5}\cdot\frac{1}{16}\\ C&=\frac{1}{2\cdot3}\cdot\frac{1}{4}B-\frac{1}{2\cdot3\cdot4\cdot5}\cdot\frac{1}{16}A+\frac{1}{2\cdot3\cdot4\cdot5\cdot6\cdot7}\cdot\frac{1}{64}\\ &\text{etc.} \end{aligned} \]
If one calls $S_{n}$ the sum of the series $1+\frac{1}{2^{n}}+\frac{1}{3^{n}}+\frac{1}{4^{n}}+\text{etc.}$, one will have again
\[ \begin{aligned} A&=\frac{1}{2\pi^{2}}\left(1-\frac{1}{2}\right)S_{2}\\ B&=\frac{1}{2^{3}\pi^{4}}\left(1-\frac{1}{2^{3}}\right)S_{4}\\ C&=\frac{1}{2^{5}\pi^{6}}\left(1-\frac{1}{2^{5}}\right)S_{6}\\ &\text{etc.} \end{aligned} \]
And in general, $N$ being the $n^{\text{th}}$ term of the series $A$, $B$, $C$, etc., one will have
\[ N=\frac{1}{2^{2n-1}\pi^{2n}}\left(1-\frac{1}{2^{2n-1}}\right)S_{2n}. \]

We now know the law of the terms of equation (2); but in order that this formula may be employed with success in the determination of the area $Z$, one must take $\omega$ small enough that the term $B\omega^{4}\left(\frac{d^{3}F}{dx^{3}}-\frac{d^{3}F^{0}}{dx^{3}}\right)$ and the following ones may be neglected;

\origpage{313}
and this is always easy to do, unless $\frac{d^{3}F}{dx^{3}}$ is infinite at one or the other limit of the integral.

Thus, this case excepted, and a fortiori that in which $\frac{dF}{dx}$ would become infinite at one of these limits, one will determine the area $Z$ by the formula
\[ Z=\omega\Sigma F(x+\tfrac{1}{2}\omega)+\frac{\omega^{2}}{24}\left(\frac{dF}{dx}-\frac{dF^{0}}{dx}\right), \]
the second part being the correction of the first.

\subsection*{(4).}

If the coefficient $\frac{dF}{dx}$ were infinite at one of the limits of the integral, that is, if at one of these points the ordinate were tangent to the curve, one would have to seek by another procedure the area contained between this ordinate and an ordinate not far removed from it; the rest of the area would be subject to no difficulty.

Suppose, for example, that $\frac{dy}{dx}$ is infinite at the second limit of the integral, and let $y=b$ be the ordinate which is tangent to the curve at this point. If, on putting $x=a-\alpha$, one has $y=c$, then the area contained between the two ordinates $c$, $b$, will be very nearly $\frac{\alpha}{3}(2c+b)$. This results from the supposition that the arc of the curve in question can be likened to an arc of a parabola whose axis is parallel to the line of abscissas.

To the quantity $\frac{\alpha}{3}(2c+b)$ one must therefore join the area contained from $x=0$ to $x=a-\alpha$.

A more intimate knowledge of the nature of the curve may lead to a closer value of the area contained from $x=a-\alpha$ to $x=a$; but the preceding determination will suffice in almost all cases, and the procedure would be the same if $\frac{dF}{dx}$ were infinite at the first limit.

Likewise if $\frac{dF}{dx}$ is not infinite at one of the limits of the integral, but $\frac{d^{3}F}{dx^{3}}$ is, or has a very large value; then it will be proper to calculate in a particular manner the area of the portion

\origpage{314}
of the curve in which this coefficient has too large a value for formula (2) to be applied to it with safety.

\subsection*{(5).}

Let us now inquire what error can result from our formula when, between the limits of the integral, there is a point where the value of $x$ renders infinite either the coefficient $\frac{dy}{dx}$, or one of the two following, $\frac{ddy}{dx^{2}}$, $\frac{d^{3}y}{dx^{3}}$; I observe that if $x=\alpha$ and $y=\beta$ are the coordinates of this singular point, one must in general have
\[ y=\beta+A(x-\alpha)^{\mu}+B(x-\alpha)^{\mu+\nu}+C(x-\alpha)^{\mu+2\nu}+\text{etc.}; \]
$\mu$ being a positive fractional number, and $\nu$ a number likewise positive, but which may be supposed equal to unity, except in extraordinary cases. I observe moreover that, since the curve is supposed to form a continuous branch from $x=0$ to $x=a$, $a$ being $>\alpha$, $\mu$ must be represented by a fraction $\frac{p}{q}$ whose denominator will always be odd, so that the supposition of $x-\alpha$ negative may not render $y-\beta$ imaginary.

Suppose now that the interval $\omega$, which is an aliquot part of $a$, is one also of $\alpha$, which is always possible by cutting off a small part of the proposed area, and afterwards adding back the part cut off; the finite integral $\Sigma\omega F(x+\frac{1}{2}\omega)$ will contain the two terms $\omega F(\alpha-\frac{1}{2}\omega)+\omega F(\alpha+\frac{1}{2}\omega)$ which will correspond to the part of the area contained from $x=\alpha-\omega$ to $x=\alpha+\omega$, and formula (2) will suppose that this part of the area $\zeta$ is represented by the formula
\[ \zeta=\omega F(\alpha-\tfrac{1}{2}\omega)+\omega F(\alpha+\tfrac{1}{2}\omega)+\frac{\omega^{2}}{24}\left(\frac{dF''}{dx}-\frac{dF'}{dx}\right), \]
$\frac{dF'}{dx}$ and $\frac{dF''}{dx}$ being the values of $\frac{dy}{dx}$ which correspond to the abscissas $\alpha-\omega$ and $\alpha+\omega$. Now the supposed value of $y$ gives
\[ \begin{aligned} F(\alpha+\tfrac{1}{2}\omega)&=\beta+A(\tfrac{1}{2}\omega)^{\mu}+B(\tfrac{1}{2}\omega)^{\mu+\nu}+\text{etc.}\\ F(\alpha-\tfrac{1}{2}\omega)&=\beta+A(-\tfrac{1}{2}\omega)^{\mu}+B(-\tfrac{1}{2}\omega)^{\mu+\nu}+\text{etc.}\\ \frac{dF''}{dx}&=A\mu\omega^{\mu-1}+B(\mu+\nu)\omega^{\mu+\nu-1}+\text{etc.}\\ \frac{dF'}{dx}&=A\mu(-\omega)^{\mu-1}+B(\mu+\nu)(-\omega)^{\mu+\nu-1}+\text{etc.} \end{aligned} \]

\origpage{315}
Therefore
\[ \begin{aligned} \zeta={}&2\beta\omega+A\omega^{\mu+1}\left(\frac{1}{2^{\mu}}+\frac{\mu}{24}\right)[1+(-1)^{\mu}]\\ &+B\omega^{\mu+\nu+1}\left(\frac{1}{2^{\mu+\nu}}+\frac{\mu+\nu}{24}\right)[1+(-1)^{\mu+\nu}]+\text{etc.} \end{aligned} \]
But the integral $\displaystyle\int y\,dx$, taken from $x=\alpha-\omega$ to $x=\alpha+\omega$, has for its exact value
\[ 2\beta\omega+A\omega^{\mu+1}\left[\frac{1+(-1)^{\mu}}{\mu+1}\right]+B\omega^{\mu+\nu+1}\left[\frac{1+(-1)^{\mu+\nu}}{\mu+\nu+1}\right]+\text{etc.} \]
Thus the correction that must be applied to $\zeta$ or to $Z$ in order to have the true value of the area sought, will be
\[ \begin{aligned} \delta Z={}&A\omega^{\mu+1}\left(\frac{1}{\mu+1}-\frac{1}{2^{\mu}}-\frac{\mu}{24}\right)[1+(-1)^{\mu}]\\ &+B\omega^{\mu+\nu+1}\left(\frac{1}{\mu+\nu+1}-\frac{1}{2^{\mu+\nu}}-\frac{(\mu+\nu)}{24}\right)[1+(-1)^{\mu+\nu}]. \end{aligned} \]
In this formula two cases must be distinguished, according as $\mu$ is of the form $\frac{2k}{2i+1}$ or $\frac{2k+1}{2i+1}$; supposing besides $\nu=1$, which can suffer only very few exceptions, one will have in the first case,
\[ \delta Z=2A\omega^{\mu+1}\left(\frac{1}{\mu+1}-\frac{1}{2^{\mu}}-\frac{\mu}{24}\right), \]
and in the second case,
\[ \delta Z=2B\omega^{\mu+2}\left[\frac{1}{\mu+2}-\frac{1}{2^{\mu+1}}-\left(\frac{\mu+1}{24}\right)\right]. \]
Let $F(\alpha+\frac{1}{2}\omega)-2F(\alpha)+F(\alpha-\frac{1}{2}\omega)=M$; the quantity $M$ will be known from the terms which compose $\Sigma\omega F(x+\frac{1}{2}\omega)$, and by means of this quantity one will have, in the first case,
\[ \delta Z=M\omega\left(\frac{2^{\mu}}{\mu+1}-1-\frac{\mu\cdot2^{\mu}}{24}\right), \]
and in the second
\[ \delta Z=M\omega\left(\frac{2^{\mu+1}}{\mu+2}-1-\frac{(\mu+1)2^{\mu+1}}{24}\right). \]
Thus the correction to be applied to the sum found will be known,

\origpage{316}
% The print numbers this page "163" (digits transposed); 316 by sequence.
provided that one knows the exponent $\mu$ which serves to characterize the singular point that renders infinite one or several of the coefficients $\frac{dy}{dx}$, $\frac{d^{2}y}{dx^{2}}$, $\frac{d^{3}y}{dx^{3}}$. One sees at the same time that the correction will be much smaller when $\mu$ is of the form $\frac{2k+1}{2i+1}$ than when $\mu$ is of the form $\frac{2k}{2i+1}$; thus the case where $y$ becomes a \emph{maximum} or a \emph{minimum} when $x=\alpha$ is that which requires the greatest correction. The other case brings to the formula only a very slight correction, and often a negligible one.

\subsection*{(6).}

When it has been recognized that the portion of the curve which one wishes to square has a singular point determined by the equation
\[ y-\beta=A(x-\alpha)^{\mu}+B(x-a)^{\mu+\nu}+\text{etc.}; \]
\ednote{Printed $(x-a)$ in the second term; $(x-\alpha)$ is meant, as in the first term and in art. (5). An apparent misprint.}
one will be able, by a very simple transformation, to forestall the inconvenience which arises from the infinite value of the coefficients $\frac{dy}{dx}$, $\frac{ddy}{dx^{2}}$, $\frac{d^{3}y}{dx^{3}}$. It suffices for this to put $x=\alpha+(z-f)^{m}$, $m$ being the odd denominator of the fraction equal to $\mu$; then $y$ will become a known function of $z$, and if one takes $f=\sqrt[m]{\alpha}$, one will have $\displaystyle\int y\,dx=\displaystyle\int m(z-f)^{m-1}y\,dz$. This new integral must be taken from $z=0$ to $z=\sqrt[m]{\alpha}+\sqrt[m]{(a-\alpha)}$; and throughout this whole interval the new curve which is to be squared, and which has for ordinate the function $m(z-f)^{m-1}y$, will have no singular point; so that the application of formula (2) will be subject to no difficulty.

\subsection*{(7).}

Most of the methods that are ordinarily given for finding the area of a curve by approximation are less exact or less convenient in practice than the one we have just set out. I regard as one of the most defective, above all, the one which supposes that the ordinate of the curve is represented throughout its whole extent by the formula $y=a+bx+cx^{2}+ex^{3}+\text{etc.}$, or by an equivalent formula; for from the fact that one curve passes through a great number of points of a given curve, it does not follow that the two curves are very close to one another; it may happen on the contrary that the two areas, despite all the common points, are as different from each other as one pleases.

\origpage{317}
The only practical method which is at once simple and exact for quadratures consists in dividing the curve by equidistant ordinates, in such a way that the arcs contained between every second ordinate are likened to arcs of a parabola. Let us denote by $c$, $c'$, $c''\ldots c^{(2n)}$ the series of ordinates corresponding to the abscissas $0$, $\frac{1}{2}\omega$, $\omega$, $\frac{3}{2}\omega\ldots n\omega$; the parabolic area from $x=0$ to $x=\omega$ will be expressed by $\frac{\omega}{6}(c+4c'+c'')$; likewise the parabolic area from $x=\omega$ to $x=2\omega$ will be $\frac{\omega}{6}(c''+4c'''+c^{\text{iv}})$; and so on. Adding all these parts, one has for the entire area $S$, contained from $x=0$ to $x=n\omega$, the formula
\[ S=\frac{\omega}{6}(c+4c'+2c''+4c'''+2c^{\text{iv}}\ldots+4c^{(2n-1)}+c^{(2n)}), \tag{3} \]
or, what comes to the same thing, according to the preceding denominations,
\[ S=\frac{\omega}{6}[F(0)+4F(\tfrac{1}{2}\omega)+2F(\omega)+4F(\tfrac{3}{2}\omega)\ldots+4F(x-\tfrac{1}{2}\omega)+F(x)]. \]
This formula can be written thus:
\[ \begin{aligned} S={}&\frac{\omega}{6}[F(0)-F(x)]\\ &+\frac{2\omega}{3}[F(\tfrac{1}{2}\omega)+F(\tfrac{3}{2}\omega)+F(\tfrac{5}{2}\omega)\ldots+F(x-\tfrac{1}{2}\omega)]\\ &+\frac{\omega}{3}[F(\omega)+F(2\omega)+F(3\omega)\ldots+F(x-\omega)+F(x)]. \end{aligned} \]
But according to formula (2), the area contained from $x=0$ to $x=n\omega$ being called $Z$, one will have
\[ \begin{aligned} Z={}&\omega[F(\tfrac{1}{2}\omega)+F(\tfrac{3}{2}\omega)+F(\tfrac{5}{2}\omega)\ldots F(x-\tfrac{1}{2}\omega)]\\ &+A\omega^{2}P-B\omega^{4}Q+C\omega^{6}R-\text{etc.}, \end{aligned} \]
$P$, $Q$, $R$, etc. depending on the value of the differential coefficients $\frac{dF}{dx}$, $\frac{d^{3}F}{dx^{3}}$, $\frac{d^{5}F}{dx^{5}}$, etc. at the limits of the integral.

If one calls $Z'$ the area contained from $x=\frac{1}{2}\omega$ to $x=(n+\frac{1}{2})\omega$, the expression of this area will be similarly
\[ \begin{aligned} Z'={}&\omega[F(\omega)+F(2\omega)+F(3\omega)\ldots+F(n\omega)]\\ &+A\omega^{2}P'-B\omega^{4}Q'+C\omega^{6}R'-\text{etc.}, \end{aligned} \]

\origpage{318}
denoting by $P'$, $Q'$, $R'$, etc. what $P$, $Q$, $R$, etc. become when, instead of $x$, one puts $x+\frac{1}{2}\omega$. But if one denotes by $\psi$ the increment of the area $Z$ when $x$ becomes $x+\frac{1}{2}\omega$, and by $\psi^{0}$ what $\psi$ becomes when $x=0$, we shall have $Z'=Z+\psi-\psi^{0}$.

Now since $Z=\displaystyle\int dx\,F(x)$, we deduce from it
\[ \psi=\frac{\omega}{2}F-\frac{\omega^{2}}{2\cdot4}\frac{dF}{dx}+\frac{\omega^{3}}{2\cdot3\cdot8}\cdot\frac{ddF}{dx^{2}}-\text{etc.}, \]
and consequently,
\[ \psi^{0}=\frac{\omega}{2}F^{0}-\frac{\omega^{2}}{2\cdot4}\cdot\frac{dF^{0}}{dx}+\frac{\omega^{3}}{2\cdot3\cdot8}\cdot\frac{ddF^{0}}{dx^{2}}-\text{etc.} \]
From these equations we obtain,
\[ \begin{aligned} \omega[F(\tfrac{1}{2}\omega)+F(\tfrac{3}{2}\omega)\ldots+F(x-\tfrac{1}{2}\omega)]&=Z-A\omega^{2}P+B\omega^{4}Q-\text{etc.}\\ \omega[F(\omega)+F(2\omega)\ldots+F(x)]&=Z+\psi-\psi^{0}-A\omega^{2}P'+B\omega^{4}Q'-\text{etc.} \end{aligned} \]
Substituting these values in that of $S$, there will result
\[ S=\frac{\omega}{6}(F^{0}-F)+Z+\tfrac{1}{3}(\psi-\psi^{0})-A\omega^{2}\left(\frac{2P+P'}{3}\right)+B\omega^{4}\left(\frac{2Q+Q'}{3}\right)-\text{etc.}, \]
or, substituting the value of $\psi-\psi^{0}$,
\[ \begin{aligned} S=Z&-\frac{\omega^{2}}{2\cdot4}\cdot\frac{1}{3}\left(\frac{dF}{dx}-\frac{dF^{0}}{dx}\right)+\frac{\omega^{2}}{2\cdot3\cdot8}\cdot\frac{1}{3}\left(\frac{ddF}{dx^{2}}-\frac{ddF^{0}}{dx^{2}}\right)\\ &-\frac{\omega^{4}}{2\cdot3\cdot4\cdot16}\cdot\frac{1}{3}\left(\frac{d^{3}F}{dx^{3}}-\frac{d^{3}F^{0}}{dx^{3}}\right)+\text{etc.}\\ &-A\omega^{2}\left(\frac{2P+P'}{3}\right)+B\omega^{4}\left(\frac{2Q+Q'}{3}\right)-\text{etc.} \end{aligned} \]
\ednote{Printed $\omega^{2}$ over $2\cdot3\cdot8$ in the second term of the first line; $\omega^{3}$ is meant, as in the value of $\psi$ above. An apparent misprint.}
But we have $P=\frac{dF}{dx}-\frac{dF^{0}}{dx}$; if in this quantity one puts $x+\frac{1}{2}\omega$ in place of $x$, $\frac{dF}{dx}$ will become $\frac{dF}{dx}+\frac{\omega}{2}\cdot\frac{ddF}{dx^{2}}+\frac{1}{2}\cdot\frac{\omega^{2}}{4}\cdot\frac{d^{3}F}{dx^{3}}+\text{etc.}$; therefore
\[ P'=\frac{dF}{dx}-\frac{dF^{0}}{dx}+\frac{\omega}{2}\left(\frac{ddF}{dx^{2}}-\frac{ddF^{0}}{dx^{2}}\right)+\frac{1}{2}\cdot\frac{\omega^{2}}{4}\left(\frac{d^{3}F}{dx^{3}}-\frac{d^{3}F^{0}}{dx^{3}}\right)+\text{etc.}; \]
likewise, having $Q=\frac{d^{3}F}{dx^{3}}-\frac{d^{3}F^{0}}{dx^{3}}$, we shall deduce from it
\[ Q'=\frac{d^{3}F}{dx^{3}}-\frac{d^{3}F^{0}}{dx^{3}}+\frac{1}{2}\omega\left(\frac{d^{4}F}{dx^{4}}-\frac{d^{4}F^{0}}{dx^{4}}\right)+\text{etc.} \]

\origpage{319}
This being laid down, making the substitutions and rejecting the powers of $\omega$ beyond the fourth, we shall have
\[ S=Z+\frac{\omega^{4}}{2880}\left(\frac{d^{3}F}{dx^{3}}-\frac{d^{3}F^{0}}{dx^{3}}\right); \]
whence one sees how small the difference is between the exact value of the area $Z$, and its value $S$ approximated by a sum of parabolic areas.

Thus in practice one may, if one wishes, keep to formula (3), which gives the value of $S$, and add to it the small correction furnished by the preceding formula in order to have a more exact value of the area $Z$, which will give
\[ Z=S-\frac{\omega^{4}}{2880}\left(\frac{d^{3}F}{dx^{3}}-\frac{d^{3}F^{0}}{dx^{3}}\right). \]
Moreover formula (3) supposes the calculation of twice as many terms as formula (2); thus on this ground alone, the use of formula (2) is preferable.

\subsection*{(8).}

We have explained sufficiently in what precedes how one can find a proposed integral $\displaystyle\int y\,dx$ between two given limits $x=0$, $x=a$; and we have given the means of pushing the approximation to whatever degree one may wish to fix. But if the integral had to extend from $x=0$ to $x=\infty$, one might fear that the method would be impracticable, since the finite integral $\Sigma\omega F(x+\frac{1}{2}\omega)$ would then contain an infinity of terms. The answer to this objection is that if the integral sought is to be a finite quantity, the only case in which there is occasion to calculate it, the ordinate $F(x)$ must decrease with great rapidity when $x$ becomes somewhat large, so that the terms which compose $\Sigma\omega F(x+\frac{1}{2}\omega)$ will soon become very small and entirely negligible.

It is easy, moreover, to remedy this inconvenience by a transformation. Indeed if one makes, for example, $x=\frac{z}{1-z}$, the integral $\displaystyle\int y\,dx$ will be transformed into another $\displaystyle\int P\,dz$, in which $P$ will be a known function of $z$ and which will have to be integrated from $z=0$ to $z=1$.

One could also make $x=k\operatorname{tang}\varphi$, $k$ being at will; and the

\origpage{320}
transform in $\varphi$, which would be of the form $\displaystyle\int Q\,d\varphi$, would have to be integrated from $\varphi=0$ to $\varphi=\frac{1}{2}\pi$.

In general it is by the nature of the question that one will judge which substitution it is fitting to employ, in preference to any other, in order to transform the proposed integral into another which is comprised between finite limits.

\section*{Construction of the curve in which the arc $s$ is given as a function of the quantity $\frac{dy}{dx}$:}

\subsection*{(9).}

Let $s$ be the arc of a curve, $\theta$ the angle which the tangent at the extremity of the arc makes with the line of abscissas, so that $\operatorname{tang}\theta=\frac{dy}{dx}$; if the equation of the curve is not given, but one has simply the expression of the arc $s$ as a function of the angle $\theta$, namely $s=F(\theta)$, and if from this expression one cannot deduce the values of the coordinates $x$ and $y$, because they depend on integrals that are too complicated, the problem is to find at least by approximation the values of $x$ and $y$ which correspond to a given value of the angle $\theta$.

Since $s$ is a known function of $\theta$, one may suppose $ds=Q\,d\theta$, and then the values of $x$ and $y$ depend immediately on quadratures, since one has $x=\displaystyle\int Q\,d\theta\cos\theta$, $y=\displaystyle\int Q\,d\theta\sin\theta$. One could therefore make use of the method given in the preceding chapter to calculate the values of $x$ and $y$ which correspond to a given value of $\theta$. But the curve will be constructed more easily by a particular method which we are going to set out.

Suppose that from a given point where $\theta=\alpha$, to any other point where $\theta$ has a given value, one wishes to know the values of $x$ and $y$; one will divide the interval $\theta-\alpha$ into a certain number of equal parts, so much the smaller as one wishes to push the approximation further. Let $\omega$ be one of these parts, and $n$ their number, so that $\theta=\alpha+n\omega$.

By means of the given equation between $s$ and $\theta$, one will calculate successively the values of $s$ which answer to the angles $\alpha$, $\alpha+\omega$, $\alpha+2\omega\ldots\alpha+n\omega$, and one will take their consecutive differences.

\origpage{321}
Let $\Delta s$ be one of these differences, supposed to correspond to the angle $\theta$, so that $s=F(\theta)$, $s+\Delta s=F(\theta+\omega)$; we shall have for first approximate values of $x$ and $y$, the formulas
\[ \begin{aligned} x&=\Sigma\Delta s\cos(\theta+\tfrac{1}{2}\omega)\\ y&=\Sigma\Delta s\sin(\theta+\tfrac{1}{2}\omega), \end{aligned} \]
where one sees that each $\Delta s$ is multiplied by the cosine or the sine of the angle $\theta+\frac{1}{2}\omega$, mean between the angle $\theta$ which answers to the arc $s$, and the angle $\theta+\omega$ which answers to the arc $s+\Delta s$.

\subsection*{(10).}

We must now see what corrections must be applied to these first approximate values. For this purpose let us seek in general the value of $\xi$ from the equation
\[ x=\Sigma\Delta s\cos(\theta+\tfrac{1}{2}\omega)+\xi, \]
we shall have
\[ \Delta\xi=\Delta x-\Delta s\cos(\theta+\tfrac{1}{2}\omega); \]
Now regarding $x$ and $s$ as functions of $\theta$, we have
\[ \begin{aligned} \Delta s&=\omega\frac{ds}{d\theta}+\frac{\omega^{2}}{2}\cdot\frac{dds}{d\theta^{2}}+\frac{\omega^{3}}{2\cdot3}\cdot\frac{d^{3}s}{d\theta^{3}}+\text{etc.}\\ \Delta x&=\omega\frac{dx}{d\theta}+\frac{\omega^{2}}{2}\cdot\frac{ddx}{d\theta^{2}}+\frac{\omega^{3}}{2\cdot3}\cdot\frac{d^{3}x}{d\theta^{3}}+\text{etc.} \end{aligned} \]
We have at the same time
\[ \cos(\theta+\tfrac{1}{2}\omega)=\cos\theta-\frac{\omega}{2}\sin\theta-\frac{1}{1\cdot2}\cdot\frac{\omega^{2}}{4}\cos\theta+\frac{1}{1\cdot2\cdot3}\cdot\frac{\omega^{3}}{8}\sin\theta+\text{etc.} \]
But from $dx=ds\cos\theta$, we deduce
\[ \begin{aligned} \frac{dx}{d\theta}&=\frac{ds}{d\theta}\cos\theta\\ \frac{ddx}{d\theta^{2}}&=\frac{dds}{d\theta^{2}}\cos\theta-\frac{ds}{d\theta}\sin\theta\\ \frac{d^{3}x}{d\theta^{3}}&=\frac{d^{3}s}{d\theta^{3}}\cos\theta-2\frac{dds}{d\theta^{2}}\sin\theta-\frac{ds}{d\theta}\cos\theta\\ \frac{d^{4}x}{d\theta^{4}}&=\frac{d^{4}s}{d\theta^{4}}\cos\theta-3\frac{d^{3}s}{d\theta^{3}}\sin\theta-3\frac{d^{2}s}{d\theta^{2}}\cos\theta+\frac{ds}{d\theta}\sin\theta.\\ &\text{etc.} \end{aligned} \]

\origpage{322}
Thus we have on the one hand,
\[ \begin{aligned} \Delta x={}&\frac{ds}{d\theta}\left(\omega\cos\theta-\frac{\omega^{2}}{2}\sin\theta-\frac{\omega^{3}}{2\cdot3}\cos\theta+\frac{\omega^{4}}{2\cdot3\cdot4}\sin\theta+\text{etc.}\right)\\ &+\frac{dds}{d\theta^{2}}\left(\frac{\omega^{2}}{2}\cos\theta-2\frac{\omega^{3}}{2\cdot3}\sin\theta-3\frac{\omega^{4}}{2\cdot3\cdot4}\cos\theta+\text{etc.}\right)\\ &+\frac{d^{3}s}{d\theta^{3}}\left(\frac{\omega^{3}}{2\cdot3}\cos\theta-3\frac{\omega^{4}}{2\cdot3\cdot4}\sin\theta-\text{etc.}\right)\\ &+\frac{d^{4}s}{d\theta^{4}}\left(\frac{\omega^{4}}{2\cdot3\cdot4}\cos\theta-\text{etc.}\right)+\text{etc.}; \end{aligned} \]
and on the other hand,
\[ \begin{aligned} \Delta s\cos(\theta+\tfrac{1}{2}\omega)={}&\frac{ds}{d\theta}\left(\omega\cos\theta-\frac{\omega^{2}}{2}\sin\theta-\frac{\omega^{3}}{2\cdot4}\cos\theta+\frac{\omega^{4}}{2\cdot4\cdot6}\sin\theta+\text{etc.}\right)\\ &+\frac{dds}{d\theta^{2}}\left(\frac{\omega^{2}}{2}\cos\theta-\frac{\omega^{3}}{4}\sin\theta-\frac{\omega^{4}}{2\cdot4}\cdot\frac{1}{2}\cos\theta+\text{etc.}\right)\\ &+\frac{d^{3}s}{d\theta^{3}}\left(\frac{\omega^{3}}{2\cdot3}\cos\theta-\frac{\omega^{4}}{2\cdot3}\cdot\frac{1}{2}\sin\theta-\text{etc.}\right)\\ &+\frac{d^{4}s}{d\theta^{4}}\left(\frac{\omega^{4}}{2\cdot3\cdot4}\cos\theta-\text{etc.}\right)+\text{etc.} \end{aligned} \]
The difference of these two quantities will give, stopping at the $\omega^{4}$,
\[ \begin{aligned} \Delta\xi={}&-\frac{ds}{d\theta}\left(\frac{\omega^{3}}{2\cdot3\cdot4}\cos\theta-\frac{\omega^{4}}{2\cdot4\cdot6}\sin\theta\right)\\ &-\frac{dds}{d\theta^{2}}\left(\frac{\omega^{3}}{3\cdot4}\sin\theta+\frac{\omega^{4}}{16}\cos\theta\right)\\ &-\frac{d^{3}s}{d\theta^{3}}\left(\frac{\omega^{4}}{2\cdot3\cdot4}\sin\theta\right). \end{aligned} \]
Let $\xi=P\omega^{2}+Q\omega^{3}+R\omega^{4}$, we shall have
\[ \Delta\xi=\omega^{3}\frac{dP}{d\theta}+\frac{\omega^{4}}{2}\cdot\frac{ddP}{d\theta^{2}}+\omega^{4}\frac{dQ}{d\theta}. \]
Thus the coefficients $P$ and $Q$ must be determined by the equations
\[ \begin{aligned} \frac{dP}{d\theta}&=-\frac{ds}{d\theta}\cdot\frac{\cos\theta}{2\cdot3\cdot4}-\frac{dds}{d\theta^{2}}\cdot\frac{\sin\theta}{3\cdot4}\\ \frac{1}{2}\frac{ddP}{d\theta^{2}}+\frac{dQ}{d\theta}&=\frac{ds}{d\theta}\cdot\frac{\sin\theta}{48}-\frac{dds}{d\theta^{2}}\cdot\frac{\cos\theta}{16}-\frac{d^{3}s}{d\theta^{3}}\cdot\frac{\sin\theta}{24}. \end{aligned} \]
The first gives on integrating,
\[ P=-\frac{ds}{d\theta}\cdot\frac{\sin\theta}{2\cdot3\cdot4}+\frac{1}{2\cdot3\cdot4}\int ds\cos\theta, \]
\ednote{The first term is printed with the denominator $2\cdot3\cdot4$ here and in the next formula; integrating the preceding equation gives $3\cdot4$, which the corrected value of $x$ below (coefficient $\frac{\omega^{2}}{12}$) presupposes. An apparent misprint.}

\origpage{323}
or
\[ P=-\frac{ds}{d\theta}\cdot\frac{\sin\theta}{2\cdot3\cdot4}+\frac{x}{2\cdot3\cdot4}+\text{constant}. \]
The other gives $\frac{dQ}{d\theta}=0$; we have therefore for first corrected value
\[ x=\Sigma\Delta s\cos(\theta+\tfrac{1}{2}\omega)-\frac{\omega^{2}}{12}\left(\frac{ds}{d\theta}\sin\theta-\tfrac{1}{2}x+\text{const.}\right); \]
and since we have found $Q=0$, it is clear that this result is exact, to quantities of the order $\omega^{4}$. Pushing the approximation further, we should find
\[ \begin{aligned} x={}&\Sigma\Delta s\cos(\theta+\tfrac{1}{2}\omega)-\frac{\omega^{2}}{2}\left(\frac{ds}{d\theta}\sin\theta-\tfrac{1}{2}x+\text{const.}\right)\\ &+\frac{\omega^{4}}{720}\left(\frac{d^{3}s}{d\theta^{3}}\sin\theta+3\frac{d^{2}s}{d\theta^{2}}\cos\theta-\frac{1}{2}\cdot\frac{ds}{d\theta}\sin\theta-\frac{3}{8}x+\text{const.}\right); \end{aligned} \]
\ednote{The first parenthesis is printed with the coefficient $\frac{\omega^{2}}{2}$, where the preceding formula has $\frac{\omega^{2}}{12}$; an apparent misprint.}
and this new result is exact to quantities of the order $\omega^{6}$; for the series of the second member contains no odd power of $\omega$, as will be demonstrated below.

On examining this value of $x$ more closely, one sees that it can be put in the form
\[ \begin{aligned} \frac{x\sin\frac{1}{2}\omega}{\frac{1}{2}\omega}={}&\Sigma\Delta s\cos(\theta+\tfrac{1}{2}\omega)-\frac{\omega^{2}}{12}\left(\frac{ds}{d\theta}\sin\theta+\text{const.}\right)\\ &+\frac{\omega^{4}}{720}\left(\frac{d^{3}s}{d\theta^{3}}\sin\theta+3\frac{dds}{d\theta^{2}}\cos\theta-\frac{1}{2}\frac{ds}{d\theta}\sin\theta+\text{const.}\right); \end{aligned} \]
whence one draws
\[ \begin{aligned} x={}&\frac{\frac{1}{2}\omega}{\sin\frac{1}{2}\omega}\Sigma\Delta s\cos(\theta+\tfrac{1}{2}\omega)-\frac{\omega^{2}}{12}\left(\frac{ds}{d\theta}\sin\theta+\text{const.}\right)\\ &+\frac{\omega^{4}}{720}\left(\frac{d^{3}s}{d\theta^{3}}\sin\theta+3\frac{dds}{d\theta^{2}}\cos\theta-3\frac{ds}{d\theta}\sin\theta+\text{const.}\right)\\ &-\text{etc.} \end{aligned} \]
The constants enclosed in the parentheses are such that they must cause the quantities joined with them to vanish, when one supposes in these $\theta=\alpha$.

The form of this expression is fairly evident, but it is important to determine the law of its coefficients, in order to be able to continue it indefinitely. For this purpose, we shall take a particular form

\origpage{324}
of $s$, such as may indicate, as simply as possible, the different terms of this development.

\subsection*{(11).}

Let then $s=e^{m\theta}-e^{m\alpha}$, so that $s=0$ when $\theta=\alpha$; the coordinates $x$ and $y$ will be determined by the differential equations $dx=ds\cos\theta=md\theta e^{m\theta}\cos\theta$, $dy=ds\sin\theta=md\theta e^{m\theta}\sin\theta$, whose integrals are
\[ \begin{aligned} (m^{2}+1)x&=m^{2}e^{m\theta}\cos\theta+me^{m\theta}\sin\theta+\text{const.}\\ (m^{2}+1)y&=m^{2}e^{m\theta}\sin\theta-me^{m\theta}\cos\theta+\text{const.} \end{aligned} \]
In this same case, we shall have $\Delta s=e^{m\theta}(e^{m\omega}-1)$, and
\[ \begin{aligned} \Sigma\Delta s\cos(\theta+\tfrac{1}{2}\omega)&=(e^{m\omega}-1)\Sigma e^{m\theta}\cos(\theta+\tfrac{1}{2}\omega)\\ &=\frac{(e^{m\omega}-1)^{2}\cos\frac{1}{2}\omega}{Re^{m\omega}}e^{m\theta}\cos\theta+\frac{(e^{2m\omega}-1)\sin\frac{1}{2}\omega}{Re^{m\omega}}e^{m\theta}\sin\theta+\text{const.}, \end{aligned} \]
a formula in which, for brevity, we have made $R=e^{m\omega}+e^{-m\omega}-2\cos\omega$. This value of $\Sigma$ can be put in the form
\[ \cos\tfrac{1}{2}\omega\cdot e^{m\theta}\cos\theta-\frac{2\sin\omega\sin\frac{1}{2}\omega}{R}e^{m\theta}\cos\theta+\frac{e^{m\omega}-e^{-m\omega}}{R}\sin\tfrac{1}{2}\omega e^{m\theta}\sin\theta+\text{const.} \]
Let then
\[ x=\frac{\frac{1}{2}\omega}{\sin\frac{1}{2}\omega}\Sigma\Delta s\cos(\theta+\tfrac{1}{2}\omega)-\xi', \]
and we shall have, making the substitutions,
\[ \begin{aligned} \xi'={}&\left(\tfrac{1}{2}\omega\cot\tfrac{1}{2}\omega-\frac{m^{2}}{m^{2}+1}-\frac{\omega\sin\omega}{R}\right)e^{m\theta}\cos\theta\\ &+\left(\frac{\frac{1}{2}\omega(e^{m\omega}-e^{-m\omega})}{R}-\frac{m}{m^{2}+1}\right)e^{m\theta}\sin\theta+\text{const.} \end{aligned} \]
We shall remark first that this value of $\xi'$ does not change sign when $\omega$ changes it, and that thus the development of $\xi'$ contains only even powers of $\omega$, as we have already announced.

\subsection*{(12).}

We must now develop the value of $\xi'$ according to the powers of $\omega$; and for that I observe that we have
\[ \frac{1}{m^{2}+1}-\frac{\omega\sin\omega}{R}=\frac{\frac{1}{3\cdot4}\omega^{2}+\frac{m^{2}-2}{3\cdot4\cdot5\cdot6}\omega^{4}+\frac{m^{4}-2m^{2}+3}{3\cdot4\ldots8}\omega^{6}+\frac{m^{6}-2m^{4}+3m^{2}-4}{3\cdot4\ldots10}\omega^{8}+\text{etc.}}{1+\frac{m^{2}-1}{3\cdot4}\omega^{2}+\frac{m^{4}-m^{2}+1}{3\cdot4\cdot5\cdot6}\omega^{4}+\frac{m^{6}-m^{4}+m^{2}-1}{3\cdot4\ldots8}\omega^{6}+\text{etc.}} \]

\origpage{325}
This quantity being a function of $m$, I denote it by $\Psi(m)$, and I remark that on making $m=0$, we shall have $1-\frac{1}{2}\omega\cot\frac{1}{2}\omega=\Psi(0)$. Suppose that on carrying out the development we have
\[ \begin{aligned} \Psi(m)&=A^{m}\omega^{2}+B^{m}\omega^{4}+C^{m}\omega^{6}+\text{etc.}\\ \Psi(0)&=A^{0}\omega^{2}+B^{0}\omega^{4}+C^{0}\omega^{6}+\text{etc.} \end{aligned} \]
As the first term $A^{m}=\frac{1}{3\cdot4}=A^{0}$, we shall have
\[ \tfrac{1}{2}\omega\cot\tfrac{1}{2}\omega-1+\frac{1}{m^{2}+1}-\frac{\omega\sin\omega}{R}=(B^{m}-B^{0})\omega^{4}+(C^{m}-C^{0})\omega^{6}+\text{etc.}; \]
this is the coefficient of $e^{m\theta}\cos\theta$ in the value of $\xi'$: we shall have similarly
\[ \frac{\frac{1}{2}\omega(e^{m\omega}-e^{-m\omega})}{mR}-\frac{1}{m^{2}+1}=\frac{\frac{\omega^{2}}{3\cdot4}+\frac{2m^{2}-1}{3\cdot4\cdot5\cdot6}\omega^{4}+\frac{3m^{4}-2m^{2}+1}{3\cdot4\cdot5\ldots8}\omega^{6}+\text{etc.}}{1+\frac{m^{2}-1}{3\cdot4}\omega^{2}+\frac{m^{4}-m^{2}+1}{3\cdot4\cdot5\cdot6}\omega^{4}+\text{etc.}}. \]
Let us call this quantity $\Phi(m)$, and suppose that its development gives
\[ \Phi(m)=\alpha^{m}\omega^{2}+\beta^{m}\omega^{4}+\gamma^{m}\omega^{6}+\delta^{m}\omega^{8}+\text{etc.}, \]
we shall have finally
\[ \begin{aligned} \xi'={}&[(B^{m}-B^{0})\omega^{4}+(C^{m}-C^{0})\omega^{6}+(D^{m}-D^{0})\omega^{8}+\text{etc.}]e^{m\theta}\cos\theta\\ &+(\alpha^{m}\omega^{2}+\beta^{m}\omega^{4}+\gamma^{m}\omega^{6}+\text{etc.})me^{m\theta}\sin\theta+\text{const.} \end{aligned} \]
These series being developed to whatever power of $\omega$ one wishes, and the coefficients $B^{m}$, $C^{m}$, $D^{m}$, $\alpha^{m}$, $\beta^{m}$, $\gamma^{m}$, etc. being expressed as functions of $m$, one will make the substitutions $me^{m\theta}=\frac{ds}{d\theta}$, $m^{2}e^{m\theta}=\frac{dds}{d\theta^{2}}$, $m^{3}e^{m\theta}=\frac{d^{3}s}{d\theta^{3}}$, etc., and one will have the value of $\xi'$, whatever the function of $\theta$ equal to $s$ may be.

\subsection*{(13).}

If for example one wishes to develop the value of $\xi'$ as far as the $\omega^{6}$ inclusively, one will first develop the functions $\Psi(m)$ and $\Phi(m)$, stopping at the $\omega^{6}$; which will give
\[ \begin{aligned} \Psi(m)&=\frac{1}{12}\omega^{2}-\frac{(3m^{2}-1)}{720}\omega^{4}+\frac{5m^{4}-10m^{2}+1}{1440\cdot21}\omega^{6}\\ \Phi(m)&=\frac{1}{12}\omega^{2}-\frac{(m^{2}-3)}{720}\omega^{4}+\frac{m^{4}-10m^{2}+5}{1440\cdot21}\omega^{6}; \end{aligned} \]

\origpage{326}
These two values would be deduced from one another by putting $\frac{1}{m}$ in place of $m$, suppressing the denominators and changing the signs of the terms taken alternately; this is also what one would deduce from the form of the fractions whose development gives the functions $\Psi(m)$ and $\Phi(m)$.

Substituting these values in the formula
\[ \xi'=[\Psi(m)-\Psi(0)]e^{m\theta}\cos\theta+m\Phi(m)e^{m\theta}\sin\theta+\text{const.}, \]
we shall have in the case supposed,
\[ \begin{aligned} \xi'={}&-\frac{3\omega^{4}}{720}m^{2}e^{m\theta}\cos\theta+\frac{5m^{4}-10m^{2}}{30240}\omega^{6}e^{m\theta}\cos\theta\\ &+\frac{\omega^{2}}{12}me^{m\theta}\sin\theta-\frac{\omega^{4}}{720}(m^{3}-3m)e^{m\theta}\sin\theta\\ &+\frac{\omega^{6}}{30240}(m^{5}-10m^{3}+5m)e^{m\theta}\sin\theta+\text{const.}; \end{aligned} \]
and consequently the general value of $\xi'$ is
\[ \begin{aligned} \xi'={}&\frac{\omega^{2}}{12}\cdot\frac{ds}{d\theta}\sin\theta-\frac{\omega^{4}}{720}\left(\frac{d^{3}s}{d\theta^{3}}\sin\theta+3\frac{dds}{d\theta^{2}}\cos\theta-3\frac{ds}{d\theta}\sin\theta\right)\\ &+\frac{\omega^{6}}{30240}\left(\frac{d^{5}s}{d\theta^{5}}\sin\theta+5\frac{d^{4}s}{d\theta^{4}}\cos\theta-10\frac{d^{3}s}{d\theta^{3}}\sin\theta-10\frac{d^{2}s}{d\theta^{2}}\cos\theta+5\frac{ds}{d\theta}\sin\theta\right)\\ &+\text{const.} \end{aligned} \]
This formula is such that one sees at first glance the law followed by the differential factors; as for the constant coefficients $\frac{1}{12}$, $\frac{1}{720}$, $\frac{1}{30240}$, etc., they are nothing other than those which are deduced from the development of the quantity $1-\frac{1}{2}\omega\cot\frac{1}{2}\omega$; so that if we make
\[ 1-\tfrac{1}{2}\omega\cot\tfrac{1}{2}\omega=A^{0}\omega^{2}+B^{0}\omega^{4}+C^{0}\omega^{6}+D^{0}\omega^{8}+\text{etc.}, \]
we shall have generally
\[ \begin{aligned} \xi'={}&\text{const}+A^{0}\omega^{2}\left(\frac{d(s\sin\theta)}{d\theta}-s\cos\theta\right)\\ &-B^{0}\omega^{4}\left(\frac{d^{3}(s\sin\theta)}{d\theta^{3}}+s\cos\theta\right)\\ &+C^{0}\omega^{6}\left(\frac{d^{5}(s\sin\theta)}{d\theta^{5}}-s\cos\theta\right)\\ &-\text{etc.} \end{aligned} \]

\origpage{327}
The constant is such that the value of $\xi'$ must vanish at the beginning of the integral where $s=0$ and $\theta=\alpha$; it will therefore be necessary, from each differential coefficient reduced to finite quantities, to subtract what this term becomes when $s=0$ and $\theta=\alpha$.

\subsection*{(14).}

The state of simplicity to which we have reduced the value of $\xi'$ makes one presume that it is possible to arrive at this formula by a more direct and less laborious route. But without stopping at this inquiry, we shall content ourselves with verifying the formula found, in its whole extent, by means of a value of $s$ which will permit us to find generally, and in a very simple manner, the difference of any order whatever of $s\sin\theta$.

We shall choose for this object the value $s=\sin a\theta$, which gives $s\sin\theta=\frac{1}{2}\cos(a-1)\theta-\frac{1}{2}\cos(a+1)\theta$. We shall then have by successive differentiations,
\[ \begin{aligned} \frac{d(s\sin\theta)}{d\theta}&=-\tfrac{1}{2}(a-1)\sin(a-1)\theta+\tfrac{1}{2}(a+1)\sin(a+1)\theta\\ \frac{d^{3}(s\sin\theta)}{d\theta^{3}}&=\tfrac{1}{2}(a-1)^{3}\sin(a-1)\theta-\tfrac{1}{2}(a+1)^{3}\sin(a+1)\theta\\ \frac{d^{5}(s\sin\theta)}{d\theta^{5}}&=-\tfrac{1}{2}(a-1)^{5}\sin(a-1)\theta+\tfrac{1}{2}(a+1)^{5}\sin(a+1)\theta\\ &\text{etc.} \end{aligned} \]
Substituting these values in that of $\xi'$, we shall have
\[ \begin{aligned} \xi'={}&-(A^{0}\omega^{2}+B^{0}\omega^{4}(a-1)^{2}+C^{0}\omega^{6}(a-1)^{4}+\text{etc.})\frac{a-1}{2}\sin(a-1)\theta\\ &+(A^{0}\omega^{2}+B^{0}\omega^{4}(a+1)^{2}+C^{0}\omega^{6}(a+1)^{4}+\text{etc.})\frac{a+1}{2}\sin(a+1)\theta\\ &-(A^{0}\omega^{2}+B^{0}\omega^{4}+C^{0}\omega^{6}+\text{etc.})s\cos\theta. \end{aligned} \]
This being laid down, since we have in general $A^{0}z^{2}+B^{0}z^{4}+C^{0}z^{6}+\text{etc.}=1-\frac{1}{2}z\cot\frac{1}{2}z$, the series comprised in the value of $\xi'$ are easy to sum, and there results from it
\[ \begin{aligned} \xi'={}&\left[1-\tfrac{1}{2}(a+1)\omega\cot(a+1)\frac{\omega}{2}\right]\frac{\sin(a+1)\theta}{2(a+1)}\\ &-\left[1-\tfrac{1}{2}(a-1)\omega\cot(a-1)\frac{\omega}{2}\right]\frac{\sin(a-1)\theta}{2(a-1)}\\ &-(1-\tfrac{1}{2}\omega\cot\tfrac{1}{2}\omega)[\tfrac{1}{2}\sin(a+1)\theta+\tfrac{1}{2}\sin(a-1)\theta]. \end{aligned} \]

\origpage{328}
I add no constant, because I suppose $\theta=0$ at the beginning of the integral.

But from the value $s=\sin a\theta$ we deduce $ds=ad\theta\cos a\theta$, $dx=ad\theta\cos a\theta\cos\theta=\frac{1}{2}ad\theta\cos(a+1)\theta+\frac{1}{2}ad\theta\cos(a-1)\theta$, and consequently
\[ x=\frac{\frac{1}{2}a}{a+1}\sin(a+1)\theta+\frac{\frac{1}{2}a}{a-1}\sin(a-1)\theta. \]
Now we have made
\[ x=\frac{\frac{1}{2}\omega}{\sin\frac{1}{2}\omega}\Sigma\Delta s\cos(\theta+\tfrac{1}{2}\omega)-\xi'; \]
thus there remains nothing more to verify than the equation
\[ \Sigma\Delta s\cos(\theta+\tfrac{1}{2}\omega)=\frac{\frac{1}{2}\sin\frac{a\omega}{2}}{\sin(a+1)\frac{\omega}{2}}\sin(a+1)\theta+\frac{\frac{1}{2}\sin\frac{a\omega}{2}}{\sin(a-1)\frac{\omega}{2}}\sin(a-1)\theta. \]
Taking the differences on both sides, supposing that $\theta$ becomes $\theta+\omega$, we find that the equation is entirely identical.

\subsection*{(15).}

The general formula found for the value of $\xi'$ is thereby established in a certain manner; for if there were in the general series a single term which did not conform to the law observed, this same term would be found again in the application to the case of $s=\sin a\theta$, since there are not two terms which are affected at once by the same power of $\omega$ and the same power of $a$, and since moreover there exists no term in the general formula which has not its correspondent in the formula proper to the particular case.

We shall have then generally
\[ x=\frac{\frac{1}{2}\omega}{\sin\frac{1}{2}\omega}\Sigma\Delta s\cos(\theta+\tfrac{1}{2}\omega)-X(\theta)+X(\alpha), \]
$X$ or $X(\theta)$ being a function of $\theta$ represented by the formula
\[ \begin{aligned} X(\theta)={}&A^{0}\omega^{2}\left(\frac{d(s\sin\theta)}{d\theta}-s\cos\theta\right)\\ &-B^{0}\omega^{4}\left(\frac{d^{3}(s\sin\theta)}{d\theta^{3}}+s\cos\theta\right)\\ &+C^{0}\omega^{6}\left(\frac{d^{5}(s\sin\theta)}{d\theta^{5}}-s\cos\theta\right)\\ &-\text{etc.}, \end{aligned} \]
and $X(\alpha)$ representing a similar function of $\alpha$.

\origpage{329}

If one changes $\theta$ into $90^{\circ}-\theta$ and $\omega$ into $-\omega$, the value of $x$ will become that of $y$; we shall therefore have
\[ y=\frac{\frac{1}{2}\omega}{\sin\frac{1}{2}\omega}\Sigma\Delta s\sin(\theta+\tfrac{1}{2}\omega)+Y(\theta)-Y(\alpha), \]
$Y(\theta)$ being a function of $\theta$ given by the equation
\[ \begin{aligned} Y(\theta)={}&A^{0}\omega^{2}\left(\frac{d(s\cos\theta)}{d\theta}+s\sin\theta\right)\\ &-B^{0}\omega^{4}\left(\frac{d^{3}(s\cos\theta)}{d\theta^{3}}-s\sin\theta\right)\\ &+C^{0}\omega^{6}\left(\frac{d^{5}(s\cos\theta)}{d\theta^{5}}+s\sin\theta\right)\\ &-\text{etc.}, \end{aligned} \]
and $Y(\alpha)$ a similar function of $\alpha$.

\subsection*{(16).}

We shall remark that the coefficients $A^{0}$, $B^{0}$, $C^{0}$, etc. have very simple relations to the coefficients $A$, $B$, $C$, etc. of which we made use in the preceding chapter (art.~3). Indeed, if in the formula
\[ 1-\tfrac{1}{2}\omega\cot\tfrac{1}{2}\omega=A^{0}\omega^{2}+B^{0}\omega^{4}+C^{0}\omega^{6}+\text{etc.}, \]
one puts $2\omega$ in place of $\omega$, we shall have
\[ 1-\omega\cot\omega=A^{0}\cdot2^{2}\omega^{2}+B^{0}\cdot2^{4}\omega^{4}+C^{0}\cdot2^{6}\omega^{6}+\text{etc.} \]
But we have $\frac{1}{\sin\omega}=\cot\frac{1}{2}\omega-\cot\omega$; therefore
\[ \frac{\omega}{\sin\omega}=1+(2^{2}-2)A^{0}\omega^{2}+(2^{4}-2)B^{0}\omega^{4}+(2^{6}-2)C^{0}\omega^{6}+\text{etc.}, \]
and it results from it that
\[ \frac{\frac{1}{2}\omega}{\sin\frac{1}{2}\omega}=1+\left(1-\frac{1}{2}\right)A^{0}\omega^{2}+\left(1-\frac{1}{2^{3}}\right)B^{0}\omega^{4}+\left(1-\frac{1}{2^{5}}\right)C^{0}\omega^{6}+\text{etc.} \]
But in the article cited, we made $\frac{\frac{1}{2}\omega}{\sin\frac{1}{2}\omega}=1+A\omega^{2}+B\omega^{4}+C\omega^{6}+\text{etc.}$; therefore

\origpage{330}
\[ \begin{aligned} A&=\left(1-\frac{1}{2}\right)A^{0} &&\text{or}\qquad & A^{0}&=\frac{2A}{2-1}=\frac{1}{12}\\ B&=\left(1-\frac{1}{2^{3}}\right)B^{0} &&& B^{0}&=\frac{2^{3}B}{2^{3}-1}=\frac{1}{720}\\ C&=\left(1-\frac{1}{2^{5}}\right)C^{0} &&& C^{0}&=\frac{2^{5}C}{2^{5}-1}=\frac{1}{30240}\\ D&=\left(1-\frac{1}{2^{7}}\right)D^{0} &&& D^{0}&=\frac{2^{7}D}{2^{7}-1}=\frac{1}{1209600}\\ &\text{etc.} &&& &\text{etc.} \end{aligned} \]
The series $A^{0}$, $B^{0}$, $C^{0}$, $D^{0}$, etc. has this in common with the series $A$, $B$, $C$, $D$, etc., that the ratio of two consecutive terms converges toward the limit $\frac{1}{4\pi^{2}}$; but this ratio is still smaller in the first terms, since we have $\frac{B^{0}}{A^{0}}=\frac{1}{60}$, $\frac{C^{0}}{B^{0}}=\frac{1}{42}$, $\frac{D^{0}}{C^{0}}=\frac{1}{40}$, and one sees that from the fourth term on, the ratio is almost equal to its limit.

In the application of these formulas, it will be fitting to take $\omega$ small enough that the first term or at most the first two of the corrections $X(\theta)$, $Y(\theta)$ suffice for the degree of approximation one has in view; thus everything depends on the magnitude of the coefficients $\frac{ds}{d\theta}$, $\frac{dds}{d\theta^{2}}$, $\frac{d^{3}s}{d\theta^{3}}$, which must be calculated beforehand for the first and the last point of the arc of the curve whose coordinates one wishes to know. When the curve has very little curvature, these coefficients are very large; then a greater arc of the curve $\Delta s$ will be needed to answer to the same difference $\Delta\theta$, and it is not astonishing that the corrections to be made to the sum of the finite differences, in order to have the sum of the infinitely small differences, should be more considerable. In this case it will therefore be necessary to take $\omega$ smaller than if the curvature were more sensible. Here, moreover, is an example in which are united all the difficulties that may be encountered in calculations of this sort, with the means of overcoming them.

\section*{Application of the preceding method to the calculation of the trajectory of a projectile.}

\subsection*{(17).}

We shall take for example the curve described by a projectile in a medium of constant density, and whose resistance

\origpage{331}
is proportional to the square of the velocity. The equation of this curve cannot be obtained in finite terms except in the case of a very small resistance or of a very small angle of projection; but one can supply this by an equation between the arc $s$ traversed since the beginning of the motion and the angle $\theta$ which the tangent to the curve makes with the horizon. This equation is
\[ e^{\frac{s}{k}}-1=\frac{h\cos^{2}\alpha}{k}[f(\alpha)-f(\theta)], \]
$f(\theta)$ denoting the function $\frac{\sin\theta}{\cos^{2}\theta}+\log\operatorname{tang}(45^{\circ}+\frac{1}{2}\theta)$, and $f(\alpha)$ a similar function of the angle $\alpha$, initial value of $\theta$.

Let us suppose the angle of projection $\alpha=45^{\circ}$, and the velocity of projection such that $\frac{h}{k}=10$, we shall have, taking $k$ for the unit,
\[ e^{s}=1+5f(\alpha)-5f(\theta). \]
From this equation, the problem is to find the values of $x$ and $y$ at the summit of the curve; one will next propose to calculate the amplitude of the descending branch.

It seems at first that it suffices to calculate the successive values of $\Delta x$ and $\Delta y$ by making $\theta$ vary by $5^{\circ}$, from $\theta=45^{\circ}$ to $\theta=0^{\circ}$. But in the part of the curve comprised from $\theta=45^{\circ}$ to $\theta=40^{\circ}$, the value of $\Delta s$ would be found very large, because the curvature is very small in this part, and this is what results from the values of $\frac{ds}{d\theta}$, $\frac{dds}{d\theta^{2}}$, etc., which are very large when $s=0$ or $\theta=45^{\circ}$.

Indeed from the given equation we draw
\[ \begin{aligned} \frac{ds}{d\theta}&=-\frac{10e^{-s}}{\cos^{3}\theta}\\ \frac{dds}{d\theta^{2}}&=3\operatorname{tang}\theta\cdot\frac{ds}{d\theta}-\frac{ds^{2}}{d\theta^{2}}\\ \frac{d^{3}s}{d\theta^{3}}&=\left(3\operatorname{tang}\theta-2\frac{ds}{d\theta}\right)\frac{dds}{d\theta^{2}}+\frac{3}{\cos^{2}\theta}\cdot\frac{ds}{d\theta}. \end{aligned} \]
If one makes in these values $\theta=45^{\circ}$ and $s=0$, we shall have

\origpage{332}
\[ \begin{aligned} \log\left(-\frac{ds}{d\theta}\right)&=1.451545\\ \log\left(-\frac{dds}{d\theta^{2}}\right)&=2.946871\\ \log\left(-\frac{d^{3}s}{d\theta^{3}}\right)&=4.723285. \end{aligned} \]
These quantities increase fairly rapidly, and if $\theta$ were made to vary by $5^{\circ}$, the corrections would turn out too considerable. This is why it is fitting to make it vary only by $1^{\circ}$, and then the second term of the correction $X(\theta)-X(\alpha)$ will affect only the seventh order of decimals.

\subsection*{(18).}

Let therefore $\omega=-1^{\circ}=-\frac{\pi}{180}$, and let us first propose to determine the portion of the curve comprised from $\theta=45^{\circ}$ to $\theta=40^{\circ}$. Here is the calculation of the successive values of $\Delta s$, $\Delta s\cos(\theta+\frac{1}{2}\omega)$, $\Delta s\sin(\theta+\frac{1}{2}\omega)$.
\[ \begin{aligned} &\theta. && s. && \Delta s. && \Delta s\cos(\theta+\tfrac{1}{2}\omega). && \Delta s\sin(\theta+\tfrac{1}{2}\omega).\\ &45^{\circ} && 0.000000 && 0.392782 && 0.280152 && 0.275305\\ &44 && 0.392782 && 0.269115 && 0.195209 && 0.185247\\ &43 && 0.661897 && 0.202695 && 0.149442 && 0.136938\\ &42 && 0.864592 && 0.161284 && 0.120795 && 0.106870\\ &41 && 1.025876 && 0.133064 && 0.101183 && 0.086418\\ &40 && 1.158940 && \text{Sums}\ldots && 0.846781 && 0.790778 \end{aligned} \]
These sums must be multiplied by the factor $\frac{\frac{1}{2}\omega}{\sin\frac{1}{2}\omega}$; now we have $\frac{\frac{1}{2}\omega}{\sin\frac{1}{2}\omega}=1+\frac{\omega^{2}}{24}+\frac{7\omega^{4}}{5760}$, and putting $\Omega=\frac{\omega^{2}}{24}+\frac{7\omega^{4}}{5760}$, then $\omega=-\frac{\pi}{180}$, we find $\log\Omega=5.103547$; which gives
\[ \begin{aligned} &\text{For the sums}\ldots\ldots && 0.846781\ldots\ldots && 0.790778\\ &\text{the corrections}\ldots\ldots && +\ 0.000011 && +\ 0.000010\\ & && 0.846792 && 0.790788. \end{aligned} \]
To obtain the corrections denoted by $X(\alpha)-X(\theta)$ and $Y(\alpha)-Y(\theta)$, we must first calculate the coefficients $\frac{ds}{d\theta}$, $\frac{dds}{d\theta^{2}}$, $\frac{d^{3}s}{d\theta^{3}}$ for the value

\origpage{333}
$\theta=40^{\circ}$. By the formulas above we shall find
\[ \begin{aligned} L\left(-\frac{ds}{d\theta}\right)&=0.843916\\ L\left(-\frac{dds}{d\theta^{2}}\right)&=1.821562\\ L\left(-\frac{d^{3}s}{d\theta^{3}}\right)&=3.052452. \end{aligned} \]
Moreover we have $L\left(\frac{\omega^{2}}{12}\right)=5.404573$ and $L\left(\frac{\omega^{4}}{720}\right)=0.110176$; from this results
\[ \begin{aligned} X(\theta)&=-0.0001139 &\qquad Y(\theta)&=-0.0001357\\ X(\alpha)&=-0.0005026 & Y(\alpha)&=-0.0005031\\ X(\theta)-X(\alpha)&=\phantom{-}0.0003887 & Y(\theta)-Y(\alpha)&=\phantom{-}0.0003674. \end{aligned} \]
Applying these corrections to the sums found,
\[ \begin{aligned} &\text{Sum of the }\Delta x\ldots && 0.846792 &\qquad& \text{of the }\Delta y\ldots && 0.790788\\ &\text{Correction}\ldots\ldots && -\ 0.000389 && \ldots\ldots && +\ 0.000367\\ &x= && 0.846403 && y= && 0.791155 \end{aligned} \]
we have the values of $x$ and of $y$ which correspond to the point where $\theta=40^{\circ}$.

\subsection*{(19).}

We can now set out from this point to make the angles $\theta$ vary by larger intervals: taking these intervals of $5^{\circ}$, we shall form the following table:
\[ \begin{aligned} &\theta. && s. && \Delta s. && \Delta s\cos(\theta+\tfrac{1}{2}\omega). && \Delta s\sin(\theta+\tfrac{1}{2}\omega).\\ &40^{\circ} && 1.158940 && 0.438379 && 0.347789 && 0.266868\\ &35 && 1.597319 && 0.258678 && 0.218167 && 0.138988\\ &30 && 1.855997 && 0.178839 && 0.158632 && 0.082579\\ &25 && 2.034836 && 0.135291 && 0.124993 && 0.051774\\ &20 && 2.170127 && 0.108854 && 0.103816 && 0.032733\\ &15 && 2.278981 && 0.091783 && 0.089607 && 0.019865\\ &10 && 2.370764 && 0.080408 && 0.079720 && 0.010495\\ &5 && 2.451172 && 0.072792 && 0.072723 && 0.003175\\ &0 && 2.523964 && \text{Sums}\ldots\ldots && 1.195447 && 0.606477 \end{aligned} \]
When $\omega=-5^{\circ}=-\frac{\pi}{36}$, we have $\log\Omega=6.501579$, so on account

\origpage{334}
of the factor $\frac{\frac{1}{2}\omega}{\sin\frac{1}{2}\omega}$, we shall have the following corrections:
\[ \begin{aligned} & 1.195447 &\qquad& 0.606477\\ +\ & 0.0003794 && 0.0001925\\ & 1.1958264 && 0.6066695. \end{aligned} \]
We have next, putting $\theta=0$,
\[ \begin{aligned} L\left(-\frac{ds}{d\theta}\right)&=9.903857 &\qquad L\left(\frac{\omega^{2}}{12}\right)&=6.802513\\ L\left(-\frac{dds}{d\theta^{2}}\right)&=9.807714 & L\left(\frac{\omega^{4}}{720}\right)&=2.906056\\ L\left(-\frac{d^{3}s}{d\theta^{3}}\right)&=0.053577; \end{aligned} \]
whence results
\[ \begin{aligned} X(0^{\circ})&=+0.0000002 &\qquad Y(0^{\circ})&=-0.0005085\\ X(40^{\circ})&=-0.0027781 & Y(40^{\circ})&=-0.0033358\\ X(0^{\circ})-X(40^{\circ})&=+0.0027783 & Y(0^{\circ})-Y(40^{\circ})&=+0.0028273\\ &\phantom{=+}1.1958264 & &\phantom{=+}0.6066695\\ x&=\phantom{+}1.193048 & y&=\phantom{+}0.609497 \end{aligned} \]
These are the corrected values of $x$ and $y$, from $\theta=40^{\circ}$ to $\theta=0$; if we add to them the values of the same coordinates, from $\theta=45^{\circ}$ to $\theta=40^{\circ}$, we shall have the coordinates which correspond to the vertex of the curve
\[ x=2.039451,\quad y=1.400652. \]

\subsection*{(20).}

We must now calculate the descending branch, and for this purpose we must change the sign of $\theta$ in the equation of the arcs, which will give
\[ e^{s}=1+5f(\alpha)+sf(\theta). \]
\ednote{Printed $sf(\theta)$ with an italic $s$; read $5f(\theta)$, as the constant $2.495587$ below requires.}

When $\theta$ becomes $\theta+\omega$, the arc $s$ will become $s+\Delta s$, and we shall have $\Delta s$ by the equation
\[ \Delta s=\log\left(\frac{2.495587+f(\theta+\omega)}{2.495587+f(\theta)}\right). \]
Putting therefore successively $\theta=0^{\circ}$, $\theta=5^{\circ}$, $\theta=10^{\circ}$, etc., and $\omega=5^{\circ}$, we shall form the following table:

\origpage{335}
\[ \begin{aligned} &\theta. && \Delta s. && \Delta s\cos(\theta+\tfrac{1}{2}\omega). && \Delta s\sin(\theta+\tfrac{1}{2}\omega).\\ &0^{\circ} && 0.067849 && 0.067785 && 0.002960\\ &5 && 0.064967 && 0.064411 && 0.008480\\ &10 && 0.063804 && 0.062291 && 0.013810\\ &15 && 0.064212 && 0.061240 && 0.019309\\ &20 && 0.066195 && 0.061156 && 0.025332\\ &25 && 0.069902 && 0.062003 && 0.032277\\ &30 && 0.075647 && 0.063800 && 0.040645\\ &35 && 0.083965 && 0.066614 && 0.051114\\ &40 && 0.095712 && 0.070566 && 0.064662\\ &45 && 0.112505 && 0.075833 && 0.082757\\ &50 && 0.135758 && 0.082644 && 0.107704\\ &55 && 0.169891 && 0.091282 && 0.143284\\ &60 && 0.221005 && 0.102049 && 0.196034\\ &65 && 0.300999 && 0.115187 && 0.278087\\ &70 && 0.434744 && 0.130730 && 0.414623\\ & && \text{Sums}\ldots && 1.177591 && 1.481078\\ & && \text{Product by }\Omega\ldots && 0.000374 && 0.000470\\ & && && 1.177965 && 1.481548 \end{aligned} \]
\ednote{Printed $0.112505$ in the $\Delta s$ column for $\theta=45^{\circ}$; the two products in the same row, and the formula for $\Delta s$, give $0.112246$, an apparent misprint. The column sums are unaffected.}

To correct these sums, we must have the values of the coefficients $\frac{ds}{d\theta}$, $\frac{dds}{d\theta^{2}}$, $\frac{d^{3}s}{d\theta^{3}}$ when $\theta=75^{\circ}$; these coefficients then become positive, and we find their logarithms as follows:
\[ \begin{aligned} \log\frac{ds}{d\theta}&=0.784599\\ \log\frac{dds}{d\theta^{2}}&=1.492717\\ \log\frac{d^{3}s}{d\theta^{3}}&=2.521532 \end{aligned} \]
By means of these values, the calculation of the corrections gives

\origpage{336}
\[ \begin{aligned} X(75^{\circ})&=0.0037066 &\qquad Y(75^{\circ})&=0.0010010\\ -X(0^{\circ})&=0.0000002 & Y(0^{\circ})&=0.0005085\\ X(75^{\circ})-X(0^{\circ})&=0.0037068 & Y(75^{\circ})-Y(0^{\circ})&=0.0004925\\ &\phantom{=}1.177965 & &\phantom{=}1.481548\\ x&=1.174258 & y&=1.482040. \end{aligned} \]
These are the values of $x$ and $y$ comprised from the vertex to the point where $\theta=75^{\circ}$.

But since the height $y$ found is greater than the height of the ascending branch $1.400652$, their difference $0.081388$ will require that a proportional diminution be made on the value of $x$, in order to have the true amplitude of the descending branch. And since the difference $0.414623$ corresponds to $5^{\circ}$ of difference in the angle $\theta$, we shall find that $0.081388$ corresponds to a difference of $58'54''$; so that the angle of fall must be about $74^{\circ}1'6''$. Taking the mean between this angle and $75^{\circ}$, we shall have $74^{\circ}30'33''$, and the quantity by which $x$ must be diminished will be $0.081388\cot 74^{\circ}30'33''=0.022557$, whence we shall conclude
\[ \begin{aligned} & && 1.174258\\ & && 0.022557\\ &\text{Ampl. of the desc. branch}\ldots\ldots && 1.151701\\ &\text{Ampl. of the asc. branch}\ldots\ldots && 2.039451\\ &\text{Total ampl.}\ldots\ldots\ldots\ldots && 3.191152; \end{aligned} \]
these results must be exact at least to the fifth decimal.

\subsection*{(21).}

Having followed the course of the trajectory to the point where $\theta=75^{\circ}$, it will not be useless to determine the position of the vertical asymptote, that is to say, to seek the value of $x$ when $\theta=90^{\circ}$.

For this purpose we shall count the $x$ from the point where $\theta=75^{\circ}$, and putting $\operatorname{tang}\theta=p$, we shall have to integrate the equation
\[ \tfrac{1}{2}\,dx=\frac{dp}{A+p\sqrt{(1+pp)}+\log[p+\sqrt{(1+pp)}]}, \]
in which $A=\frac{1}{5}+f(\alpha)=2.495587$, and the integral must extend from $p=\operatorname{tang}75^{\circ}$ to $p=\infty$.

\origpage{337}

The means we shall employ for this object consists in taking successively, for the denominator of the second member, a quantity always greater, or always smaller than this denominator, and we shall conclude that the true value of $x$ is comprised between the two which will result from each hypothesis.

Let $D=A+p\sqrt{(1+pp)}+\log[p+\sqrt{(1+pp)}]$ and $m=\operatorname{tang}75^{\circ}$; since $p$ is always comprised between $\operatorname{tang}75^{\circ}$ and $\operatorname{tang}90^{\circ}$, we shall always have $D>A+\log[m+\sqrt{(1+mm)}]+m\sqrt{(1+mm)}-m^{2}+p^{2}$; therefore if we make $c^{2}=A+f(75^{\circ})-\operatorname{tang}^{2}75^{\circ}=5.014525$, we shall have
\[ \tfrac{1}{2}\,dx<\frac{dp}{c^{2}+p^{2}}. \]
This gives, on integrating, $\frac{1}{2}cx<\operatorname{arc}\operatorname{tang}\frac{p}{c}-\operatorname{arc}\operatorname{tang}\frac{m}{c}$, and putting $p=\infty$, we shall have $\frac{1}{2}cx<\frac{\pi}{2}-\operatorname{arc}\operatorname{tang}\frac{m}{c}$, or $\operatorname{tang}\frac{1}{2}cx<\frac{c}{m}$; whence results $x<0.48268$.

To have the other limit of $x$, I put $n=\frac{1}{2m}\log[m+\sqrt{(1+mm)}]$, which gives $\log n=9.434002$; I observe next that we shall have throughout the whole extent of the integral, $p\sqrt{(1+pp)}<p^{2}+\frac{1}{2}$, $\log[p+\sqrt{(1+pp)}]<2np$; therefore we have constantly
\[ \tfrac{1}{2}\,dx>\frac{dp}{A+\frac{1}{2}+2np+p^{2}}. \]
Let $A+\frac{1}{2}-n^{2}=g^{2}$, or $\log g=0.232825$, we shall have
\[ \tfrac{1}{2}\,dx>\frac{dp}{g^{2}+(p+n)^{2}}; \]
whence results, on integrating, $\frac{1}{2}gx>\operatorname{arc}\operatorname{tang}\frac{p+n}{g}-\operatorname{arc}\operatorname{tang}\frac{m+n}{g}$; Putting $p=\infty$, we shall have $\frac{1}{2}gx>\frac{\pi}{2}-\operatorname{arc}\operatorname{tang}\frac{m+n}{g}>\operatorname{arc}\operatorname{tang}\frac{g}{m+n}$, or $x>0.47212$.

We see therefore that the value of $x$ which, starting from the point where $\theta=75^{\circ}$, corresponds to the vertical asymptote, is certainly comprised between two fairly close limits, namely, between $0.47212$ and $0.48268$. By a mean, we find $x=0.4774$, and this value cannot be in error by more than $0.0053$.

\subsection*{(22).}

It remains to know the position of the other asymptote on the side

\origpage{338}
of the negative $x$. Now it follows from the equation $e^{s}=1+5f(\alpha)-5f(\theta)$; that, making $s$ infinite negative, we shall have $f(\theta)=\frac{1}{5}+f(\alpha)=2.495587$; from this value, we find $\theta=46^{\circ}55'23''4$; this is the angle which the asymptote in question makes with the line of abscissas.

Calling this angle $\beta$, if we consider $s$, $x$ and $y$ together as positive for all the points of the branch which descends toward the asymptote; if next we put as above, $\operatorname{tang}\theta=p$, $\operatorname{tang}\beta=b$, $p\sqrt{(1+pp)}+\log[p+\sqrt{(1+pp)}]=\psi(p)$, we shall have, to determine the coordinates $x$ and $y$, the equations
\[ \tfrac{1}{2}\,dx=\frac{dp}{\psi(b)-\psi(p)},\quad \tfrac{1}{2}\,dy=\frac{p\,dp}{\psi(b)-\psi(p)}. \]
Let us put $x-\frac{y}{b}=z$, and the value of $z$ calculated for the point where $p=b$, will be the distance from the origin of abscissas to the point where the asymptote meets the line of abscissas. We shall therefore have to integrate the following formula, from $p=1$ to $p=b$,
\[ \tfrac{1}{2}b\,dz=\tfrac{1}{2}(b\,dx-dy)=\frac{(b-p)\,dp}{\psi(b)-\psi(p)}. \]
As $p$ differs little from $b$ in the interval over which we must extend the integral, we may put $p=b-u$, and we shall have the very convergent series
\[ \psi(p)=\psi(b)-\psi'(b)u+\psi''(b)\cdot\frac{u^{2}}{2}-\psi'''(b)\cdot\frac{u^{3}}{2\cdot3}+\text{etc.}, \]
where we have $\psi'(b)=2\sqrt{(1+bb)}=\frac{2}{\cos\beta}$, $\psi''(b)=\frac{2b}{\sqrt{(1+bb)}}=2\sin\beta$, etc., and we shall have to integrate the equation
\[ \tfrac{1}{2}b\,dz=\frac{-du}{\psi'(b)-\frac{u}{2}\psi''(b)+\frac{u^{2}}{2\cdot3}\psi'''(b)-\text{etc.}}. \]
To have this integral by approximation, I put
\[ \tfrac{1}{2}b\psi'(b)\,dz=-\frac{du}{1-Au}, \]
and I observe that $A$ will always be comprised between two values $A'$ and $A''$, the first of which corresponds to the beginning of the integral, when $u=b-1$, and the second to the end of the integral, when $u=0$. The latter is found exactly by the development

\origpage{339}
indicated, which gives $A''=\frac{1}{2}\cdot\frac{\psi''(b)}{\psi'(b)}=\frac{1}{2}\sin\beta\cos\beta$; the other value $A'$ will be found by the equation
\[ 1-A'(b-1)=\frac{\psi(b)-\psi(1)}{(b-1)\psi'(b)}=\frac{0.2}{(b-1)\psi'(b)}. \]
Thus we shall have $\log A'=9.3921695$ and $\log A''=9.3969607$.

But regarding $A$ as constant, the integral of the preceding equation is $\frac{1}{2}b\psi'(b)\cdot z=\frac{1}{A}\log(1-Au)+\text{const.}$ Taking therefore the integral between the limits $u=b-1$, $u=0$, we shall have
\[ z=-\frac{\log[1-A(b-1)]}{\frac{1}{2}bA\psi'(b)}. \]
Substituting successively in place of $A$ the two values $A'$ and $A''$, we shall find $z=0.0447606$ and $z=0.0447650$. The slight difference there is between these two values proves how exact this determination is, and by a mean taken between them, we shall have still more exactly $z=0.0447628$. This is the sought abscissa of the point where the asymptote in question meets the line of abscissas.

\section*{On the indefinite integral $\displaystyle\int dx\left(\log\frac{1}{x}\right)^{a-1}$, taken from $x=0$.}

\subsection*{(23).}

We have shown in the second part how the integral $\displaystyle\int dx\left(l\frac{1}{x}\right)^{a-1}$ is found when $x=1$, and we have denoted it in this case by the function $\Gamma(a)$. But it may be useful to determine this integral for any value whatever of $x$: we shall represent it generally by $\Gamma(a,x)$, so that $\Gamma(a,1)$ will be the same thing as $\Gamma(a)$.

Let us observe first that integration by parts gives the formula
\[ \Gamma(a,x)=x\left(\log\frac{1}{x}\right)^{a-1}+(a-1)\Gamma(a-1,x), \tag{1} \]
whence it follows that the integral $\Gamma(a,x)$ can always be brought back to a similar integral, in which $a$ is comprised between $1$ and $2$, as we have done for the definite integral $\Gamma(a)$.

Let $\log\frac{1}{x}=z$, we shall have the transformed expression
\[ \Gamma(a,x)=\int -z^{a-1}\,dz\,e^{-z}=\int -z^{a-1}\,dz\left(1-z+\frac{z^{2}}{2}-\frac{z^{3}}{2\cdot3}+\text{etc.}\right); \]

\origpage{340}
whence we obtain on integrating,
\[ \Gamma(a,x)=\Gamma(a)-z^{a}\left(\frac{1}{a}-\frac{z}{a+1}+\frac{1}{2}\cdot\frac{z^{2}}{a+2}-\frac{1}{2\cdot3}\cdot\frac{z^{3}}{a+2}+\text{etc.}\right). \tag{2} \]
\ednote{Printed $\frac{z^{3}}{a+2}$ in the fourth term; the integration requires $\frac{z^{3}}{a+3}$, an apparent misprint.}
This formula gives the value of $\Gamma(a,x)$ by a sequence which may be divergent in the first terms, but which always ends by being convergent.

This sequence is convergent from the very first terms, if $x$ is not smaller than $\frac{1}{e}$; it can nevertheless be employed with success for much smaller values of $x$, such as $x=\frac{1}{10}$, or even $x=\frac{1}{20}$; but then it will be necessary to calculate a fairly large number of terms of the series; for example, if $x=\frac{1}{10}$, twelve or thirteen terms of the series must be calculated to have the value of the integral approximated to the fifth or sixth decimal. In this way we find $\Gamma\left(\frac{1}{2},\frac{1}{10}\right)=0.056497$.

As $x$ becomes smaller, the sequence must be prolonged farther to obtain an equal degree of approximation; so that it is fitting to have recourse to another means to evaluate the integral with precision, when $x$ is very small, such as $\frac{1}{100}$, $\frac{1}{1000}$, etc.

\subsection*{(24).}

The best method for calculating the integral $\Gamma(a,x)$ when $x$ is very small, is to deduce it from formula (1) which gives, by successive transformations,
\[ \Gamma(a,x)=x[z^{a-1}+(a-1)z^{a-2}+(a-1)(a-2)z^{a-3}+\text{etc.}] \tag{3} \]
This series, being continued sufficiently, will have its terms become alternately positive and negative, so that we shall have values alternately greater and smaller than the integral sought. But as the sequence, which is at first convergent, necessarily becomes divergent after a certain number of terms, we must stop at the point where the divergence begins, and we shall thus have only a limited approximation.

If we call $P^{n}$ the term of rank $n$ in formula (3), and $P^{n+1}$ the following term, we shall have $P^{n+1}=\frac{a-n}{z}P^{n}$; thus the divergence will begin

\origpage{341}
when we have $n=$ or $>a+z$, and then the sequence must not be prolonged farther.

\subsection*{(25).}

Let, for example, $a=\frac{1}{2}$, $x=\frac{1}{100}$, we shall have $z=\log 100=4.60517$, and $a+z=5.1$. Thus the sequence must not be prolonged beyond the first 5 or 6 terms. Indeed, we have, putting $z=\log 10$\ednote{Printed $\log 10$; the value used, $z=4.60517$, is $\log 100$ as in the line above, an apparent misprint.},
\[ \Gamma\left(\frac{1}{2},\frac{1}{100}\right)=\frac{1}{100}\left(z^{-\frac{1}{2}}-\frac{1}{2}z^{-\frac{3}{2}}+\frac{1\cdot3}{2\cdot2}z^{-\frac{5}{2}}-\frac{1\cdot3\cdot5}{2\cdot2\cdot2}z^{-\frac{7}{2}}+\text{etc.}\right), \]
and by the calculation of successive terms, we find
\[ \begin{aligned} z^{-\frac{1}{2}}&=\phantom{+\ }0.465991\\ -\frac{1}{2}z^{-\frac{3}{2}}&=-\ 0.050594\\ &\phantom{=+\ }0.415397\\ +\frac{1\cdot3}{2\cdot2}z^{-\frac{5}{2}}&=+\ 0.016479\\ &\phantom{=+\ }0.431876\\ -\frac{1\cdot3\cdot5}{2\cdot2\cdot2}z^{-\frac{7}{2}}&=-\ 0.008946\\ &\phantom{=+\ }0.422930\\ +\frac{1\cdot3\cdot5\cdot7}{2\cdot2\cdot2\cdot2}z^{-\frac{9}{2}}&=+\ 0.006799\\ &\phantom{=+\ }0.429729\\ -\frac{1\cdot3\cdot5\cdot7\cdot9}{2\cdot2\cdot2\cdot2\cdot2}z^{-\frac{11}{2}}&=-\ 0.006644\\ &\phantom{=+\ }0.423085\\ +\frac{1\cdot3\ldots11}{2\cdot2\ldots2}z^{-\frac{13}{2}}&=+\ 0.007935 \end{aligned} \]
We see that it is useless to go farther than the sixth term, since the sequence becomes divergent from this term on.

It results from this calculation that the sum sought is greater than $0.423085$, and smaller than $0.429729$. By a mean, we find the sum of the sequence $=0.4264$, and we see that we cannot count on more than four exact decimals in this evaluation; there results from it, for the value of the integral sought, $\Gamma(a,x)=0.004264$.

\origpage{342}
\subsection*{(26).}

If, still taking $a=\frac{1}{2}$, we had put $x=\frac{1}{1000}$, we should have seen first that the sequence need be prolonged only to the ninth term, and by the actual calculation of the terms, we should have found that the integral sought is comprised between $x\times0.357708$ and $x\times0.357173$. The mean is $x\times0.35744$; whence results the integral $\Gamma(a,x)=0.00035744$.

In general formula (3) will give so much the more precision as $x$ is smaller; but this precision will be obtained only by the calculation of a greater number of terms, since in order to draw from the sequence all the approximation it can offer, it must be prolonged until the number of its terms is $a+z$ or $a+\log\frac{1}{x}$.

If we applied formula (3) to the case of $x=\frac{1}{10}$ already solved by formula (2), we should find that the sequence ceases to be convergent at the fourth term. The first two terms give $\Gamma>x\times0.5159$, and the first three give $\Gamma<x\times0.6091$; there results from it by a mean, $\Gamma=x\times0.5625=0.05625$, whereas the true value is $0.056497$. It is fitting therefore to prefer formula (2) when we wish to have at least four exact significant figures, and the value of $x$ is not smaller than $0.01$.

\subsection*{(27).}

Up to now we have tacitly supposed that $x$ does not exceed unity; but one may also ask for the integral $\displaystyle\int dx\left(l\frac{1}{x}\right)^{a-1}$ for a value of $x$ greater than unity. The part comprised from $x=0$ to $x=1$ is known and represented by $\Gamma(a)$; thus everything reduces to finding the integral from $x=1$ to any value whatever of $x>1$.

Let us remark first that in this whole interval, $l\left(\frac{1}{x}\right)$ being negative, we must put the integral in the form $(-1)^{a-1}\displaystyle\int dx(lx)^{a-1}$. I set aside the factor $(-1)^{a-1}$ which may be real or imaginary, according to the various values of $a$, and I consider simply the integral $\displaystyle\int dx(lx)^{a-1}$ which I denote by $\psi(a,x)$, and which is supposed to vanish when $x=1$.

If we put $lx=u$, we shall have $x=e^{u}$, and the integral in question

\origpage{343}
will become $\displaystyle\int u^{a-1}\,du\cdot e^{u}=\displaystyle\int u^{a-1}\,du\left(1+u+\frac{u^{2}}{2}+\text{etc.}\right)$; whence we obtain
\[ \psi(a,x)=\frac{u^{a}}{a}+\frac{u^{a+1}}{a+1}+\frac{1}{2}\cdot\frac{u^{a+2}}{a+2}+\frac{1}{2\cdot3}\cdot\frac{u^{a+3}}{a+3}+\text{etc.}, \tag{4} \]
a formula in which there is no constant to add, because it vanishes when $x=1$.

The sequence comprised in this integral may be divergent in the first terms; but it will always end by being convergent. Thus we shall draw from it in all cases a value as approximate as one wishes of the true integral, and it is evident that this value will become infinite if we put $x=\infty$.

\subsection*{(28).}

We can still put $\psi(a,x)$ in a more convenient form. Let $\displaystyle\int u^{a-1}\,du\cdot e^{u}=e^{u}P$, we shall have $u^{a-1}=\frac{dP}{du}+P$; let next $P=\frac{u^{a}}{a}+Au^{a+1}+Bu^{a+2}+\text{etc.}$, we shall find $A=-\frac{1}{a\cdot a+1}$, $B=\frac{1}{a\cdot a+1\cdot a+2}$, etc.; therefore
\[ \psi(a,x)=x\left(\frac{u^{a}}{a}-\frac{u^{a+1}}{a\cdot a+1}+\frac{u^{a+2}}{a\cdot a+1\cdot a+2}-\text{etc.}\right), \tag{5} \]
a series which will have, like the preceding one, the property of always becoming convergent, but which will have besides the advantage of giving values alternately greater and smaller than the integral sought; so that the degree of approximation will be known in each case by the difference of two consecutive terms.

\section*{On the integral $\displaystyle\int y\,dx$ taken between two limits which make the function $y$ vanish.}

\subsection*{(29).}

If the function $y$ vanishes at the two limits of the integral; if we suppose at the same time that it does not change sign from one limit to the other, and that in this interval it is susceptible of only a single \emph{maximum}, we may make use of the following method to determine the integral $Z=\displaystyle\int y\,dx$~${}^{(1)}$.

Let $m$ be the value of $x$ which makes $y$ a \emph{maximum}, and let this \emph{maximum} be

\textbf{(1)} Mém.\ de l'Acad.\ des Sc., ann.\ 1778 et 1782.

\origpage{344}
$=M$, we shall put in general
\[ y=Me^{-t^{2}}, \]
$t$ being a new variable which will extend from $t=-\infty$, to $t=+\infty$, these values corresponding to the two limits of the integral where we have $y=0$.

From the supposed value of $y$, we shall find an expression of $x$ as a function of $t$, which will be of the form
\[ x=m+At+Bt^{2}+Ct^{3}+Dt^{4}+\text{etc.}, \]
and we shall have
\[ \int y\,dx=M\int e^{-t^{2}}\,dt(A+2Bt+3Ct^{2}+4Dt^{3}+\text{etc.}). \]
In this integral two parts must be distinguished; one from $t=-\infty$ to $t=0$; the other from $t=0$ to $t=\infty$. If we change the sign of the first, it will have to be taken from $t=0$ to $t=\infty$, and will be represented by the formula
\[ M\int e^{-t^{2}}\,dt(A-2Bt+3Ct^{2}-4Dt^{3}+\text{etc.}); \]
the second, which will have to be taken between the same limits, will be
\[ M\int e^{-t^{2}}\,dt(A+2Bt+3Ct^{2}+4Dt^{3}+\text{etc.}). \]
Adding these two parts, we have the total integral
\[ Z=2M\int e^{-t^{2}}\,dt(A+3Ct^{2}+5Et^{4}+\text{etc.}). \]
Now putting $e^{-t^{2}}=z$, the integral $\displaystyle\int t^{2n}e^{-t^{2}}\,dt$ will have for its transformed expression $\frac{1}{2}\displaystyle\int dz\left(l\frac{1}{z}\right)^{\frac{2n-1}{2}}$, and this, having to be taken from $z=0$ to $z=1$, will be represented by $\frac{1}{2}\Gamma\left(\frac{2n+1}{2}\right)$. Thus we shall have generally $\displaystyle\int e^{-t^{2}}t^{2n}\,dt=\frac{1}{2}\Gamma\left(\frac{2n+1}{2}\right)$, whence results in particular
\[ \int e^{-t^{2}}\,dt=\tfrac{1}{2}\sqrt{\pi},\quad \int t^{2}e^{-t^{2}}\,dt=\tfrac{1}{2}\sqrt{\pi}\cdot\tfrac{1}{2},\quad \int t^{4}e^{-t^{2}}\,dt=\tfrac{1}{2}\sqrt{\pi}\cdot\frac{1\cdot3}{2\cdot2},\ \text{etc.} \]
Therefore, finally, the integral sought,
\[ Z=M\sqrt{\pi}\cdot\left(A+\frac{3}{2}C+\frac{3\cdot5}{2,2}E+\frac{3\cdot5\cdot7}{2\cdot2\cdot2}G+\text{etc.}\right); \tag{1} \]
\ednote{The first denominator is printed $2,2$ with a comma, where the numerators and the next denominator have product dots; kept as printed. The separators in $2\cdot2\cdot2$ are faint and could also be commas.}

\origpage{345}
and considering $x$ as a known function of $t$, then putting $t=0$ in the differential coefficients, we shall have
\[ Z=M\sqrt{\pi}\left(\frac{dx}{dt}+\frac{1}{2^{2}}\cdot\frac{d^{3}x}{dt^{3}}+\frac{1}{1\cdot2}\cdot\frac{1}{2^{4}}\cdot\frac{d^{5}x}{dt^{5}}+\frac{1}{1\cdot2\cdot3}\cdot\frac{1}{2^{6}}\cdot\frac{d^{7}x}{dt^{7}}+\text{etc.}\right). \]

\subsection*{(30).}

If in the equation $y=Me^{-t^{2}}$, or $\log M-\log y=t^{2}$, we substitute the value of $y$ in $x$, we shall have an equation between $x$ and $t$, from which we must deduce the values of the coefficients $\frac{dx}{dt}$, $\frac{d^{3}x}{dt^{3}}$, etc. which enter into the preceding formula, and in which we shall then put $t=0$ or $x=m$. It is in particular examples that it is fitting to determine these coefficients; nevertheless they may also be deduced generally from the coefficients $\frac{ddy}{dx^{2}}$, $\frac{d^{3}y}{dx^{3}}$, etc. supposed known at the point of the \emph{maximum}.

Indeed the equation $\log M-\log y=t^{2}$, being differentiated, gives $\frac{dy}{dx}\cdot\frac{dx}{dt}+2ty=0$. If we put in it $x=m$ or $t=0$, we draw no result from it, because then $\frac{dy}{dx}=0$; but if we differentiate a second time, we shall have
\[ \frac{ddy}{dx^{2}}\cdot\frac{dx^{2}}{dt^{2}}+\frac{dy}{dx}\cdot\frac{ddx}{dt^{2}}+2y+2t\frac{dy}{dx}\cdot\frac{dx}{dt}=0; \]
putting in this $t=0$, we shall have $\frac{ddy}{dx^{2}}\cdot\frac{dx^{2}}{dt^{2}}+2y=0$; whence results $\frac{dx}{dt}$, or
\[ A=\sqrt{\left(\frac{2y}{-\frac{ddy}{dx^{2}}}\right)}. \]
If therefore the series (1) is convergent enough that it can be reduced to its first term, we shall have the integral sought
\[ Z=y\sqrt{\left(\frac{2\pi y}{-\frac{ddy}{dx^{2}}}\right)}, \]
$y$ and $\frac{ddy}{dx^{2}}$ being the values of these quantities at the point of the \emph{maximum} where $x=m$.

By differentiating further the equation from which we have drawn

\origpage{346}
the value of $\frac{dx}{dt}$, we should find the values of the other coefficients $\frac{ddx}{dt^{2}}$, $\frac{d^{3}x}{dt^{3}}$, etc., expressed as functions of the quantities $y$, $\frac{ddy}{dx^{2}}$, $\frac{d^{3}y}{dx^{3}}$, etc. at the point of the \emph{maximum}. But these values become complicated as the calculation is pushed farther, and it is better, as we have said, to determine these coefficients in particular cases.

If one reflects ever so little on the spirit of the preceding method, one will see that it must give a result so much the more approximate as the function $y$ decreases more promptly on departing from the \emph{maximum}; and this is what will happen if the factors of which $y$ is composed are powers of a very high order. This circumstance permits us to have regard only to a small part of the integral, from $x=m-\alpha$ to $x=m+\alpha$; it would also contribute to simplifying the use of the other methods; but these would not give so elegant a result, because there would always remain something vague in the determination of $\alpha$. For the rest, a juster idea of this method will be gained from the following examples.

\section*{EXAMPLE I.}

\subsection*{(31).}

Let it be proposed to find the integral $Z=\displaystyle\int x^{\alpha}e^{-x}\,dx$, from $x=0$ to $x=\infty$, $\alpha$ being a positive number.

In this case, the function $x^{\alpha}e^{-x}$ vanishes at the two limits of the integral; its value in this interval is always positive, and it is susceptible of only a single \emph{maximum}, which takes place when $x=\alpha$; thus the preceding method can be applied to this integral.

Let therefore $y=x^{\alpha}e^{-x}=\left(\frac{\alpha}{e}\right)^{\alpha}e^{-t^{2}}$, there will result
\[ \alpha\log\frac{x}{\alpha}-(x-\alpha)=-t^{2}, \]
and supposing $x=\alpha+u$, we shall have by the development of this equation,
\[ \frac{u^{2}}{2\alpha^{2}}-\frac{u^{3}}{3\alpha^{3}}+\frac{u^{4}}{4\alpha^{4}}-\text{etc.}=\frac{t^{2}}{\alpha}. \]
Let, as above, $u=At+Bt^{2}+Ct^{3}+Dt^{4}+\text{etc.}$, we shall find
\[ A=\sqrt{2\alpha},\quad B=\frac{2}{3},\quad C=\frac{1}{9A},\quad D=-\frac{4}{135A^{2}},\quad E=\frac{1}{270A^{3}},\ \text{etc.} \]

\origpage{347}
Substituting these values and that of $M=\left(\frac{\alpha}{e}\right)^{\alpha}$ in formula (1), we shall have the integral sought
\[ Z=\left(\frac{\alpha}{e}\right)^{\alpha}\sqrt{(2\alpha\pi)}\left[1+\frac{1}{12\alpha}+\frac{1}{288\alpha^{2}}+\text{etc.}\right]\ldots\ldots \tag{2} \]

If in the proposed integral $Z=\displaystyle\int x^{\alpha}e^{-x}\,dx$, we put $e^{-x}=z$, we shall have the transformed expression $Z=\displaystyle\int dz\left(l\frac{1}{z}\right)^{\alpha}$ which will have to be integrated from $z=0$ to $z=1$. We have therefore $Z=\Gamma(\alpha+1)$; and indeed, the value we have just found for $Z$ agrees with formula $(\sigma)$ of No.~67, second part.

\subsection*{(32).}

This formula resolves very well the case where $\alpha$ is a large number, since it then gives a very convergent sequence (at least in the first terms); but it resolves only imperfectly the case where $\alpha$ is comprised between 0 and 1, and it gives an imaginary result when $\alpha$ is negative and smaller than unity, although the value of $\Gamma(\alpha+1)$ is then real.

It is easy to remedy the inconvenience presented by the last two cases; for integration by parts gives
\[ \int x^{\alpha}e^{-x}\,dx=\frac{1}{\alpha+1}x^{\alpha+1}e^{-x}+\frac{1}{\alpha+1}\int x^{\alpha+1}e^{-x}\,dx; \]
and as the term outside the sign vanishes at the two limits of the integral, we have simply
\[ \int x^{\alpha}e^{-x}\,dx=\frac{1}{\alpha+1}\int x^{\alpha+1}e^{-x}\,dx; \]
whence it follows that, the integral $\displaystyle\int x^{\alpha}e^{-x}\,dx$ being proposed, we may transform it into another in which the exponent of $x$ will be as great as one wishes; and then we shall determine this one in a very approximate manner by formula (2).

It is easy to see why in this example the result of the formula is so much the more exact as $\alpha$ is greater; it is that the factor $x^{\alpha}$ decreases so much the more rapidly in the neighbourhood of the \emph{maximum} as $\alpha$ is greater. There is no longer a \emph{maximum} when $\alpha$ is negative; this is why the result of the formula is entirely faulty in this case, although the true value of the integral is real, as long as $1+\alpha$ is positive.

\origpage{348}

\section*{EXAMPLE II.}

\subsection*{(33).}

Let it be proposed to find the integral $Z=\displaystyle\int x^{\alpha}(1-x)^{\beta}\,dx$, between the limits $x=0$, $x=1$, the numbers $\alpha$ and $\beta$ being positive.

Having put $y=x^{\alpha}(1-x)^{\beta}$, one finds that $y$ is a \emph{maximum} when $x=\frac{\alpha}{\alpha+\beta}=m$. Let us therefore put in general
\[ x^{\alpha}(1-x)^{\beta}=m^{\alpha}(1-m)^{\alpha}e^{-t^{2}}; \]
\ednote{Printed $(1-m)^{\alpha}$ instead of $(1-m)^{\beta}$; an apparent misprint.}
If we take the logarithms of both sides, and put $x=m+u$, we shall have
\[ \alpha\log\left(1+\frac{u}{m}\right)+\beta\log\left(1-\frac{u}{1-m}\right)=-t^{2}, \]
or
\[ \frac{u^{2}}{2}\left(\frac{\alpha}{m^{2}}+\frac{\beta}{(1-m)^{2}}\right)-\frac{u^{3}}{3}\left(\frac{\alpha}{m^{3}}-\frac{\beta}{(1-m)^{3}}\right)+\frac{u^{4}}{4}\left(\frac{\alpha}{m^{4}}+\frac{\beta}{(1-m)^{4}}\right)-\text{etc.}=t^{2}; \]
Let, as above, $u=At+Bt^{2}+Ct^{3}+\text{etc.}$, we shall find
\[ A=\frac{1}{\alpha+\beta}\sqrt{\left(\frac{2\alpha\beta}{\alpha+\beta}\right)},\quad B=\frac{2}{3}\cdot\frac{\beta-\alpha}{(\alpha+\beta)^{2}},\quad C=\frac{\alpha^{2}-11\alpha\beta+\beta^{2}}{18\alpha\beta(\alpha+\beta)}A,\ \text{etc.}; \]
whence results the integral sought
\[ Z=\frac{\alpha^{\alpha}\beta^{\beta}\sqrt{(2\pi\alpha\beta)}}{(\alpha+\beta)^{\alpha+\beta+\frac{3}{2}}}\left(1+\frac{\alpha^{2}-11\alpha\beta+\beta^{2}}{12\alpha\beta(\alpha+\beta)}+\text{etc.}\right). \]

If the exponents $\alpha$ and $\beta$ are both large numbers, this series will be very convergent; but if only one of the two is a large number, the series will be but little convergent, so that one would have to calculate many of its terms to obtain only a mediocre approximation: now these terms are difficult to calculate by the method we have just followed.

\subsection*{(34).}

Let us examine in particular the case where $\alpha$ and $\beta$ are both large numbers; then the series will reduce sensibly to its first term, and we shall have
\[ Z=\frac{\alpha^{\alpha+\frac{1}{2}}\beta^{\beta+\frac{1}{2}}\sqrt{2\pi}}{(\alpha+\beta)^{\alpha+\beta+\frac{3}{2}}}. \]

\origpage{349}
But the exact integral, when $\alpha$ and $\beta$ are integers, is
\[ Z^{1}=\frac{1\cdot2\cdot3\ldots\alpha}{\beta+1\cdot\beta+2\ldots\beta+\alpha+1}=\frac{1\cdot2\cdot3\ldots\alpha\cdot1\cdot2\cdot3\ldots\beta}{1\cdot2\cdot3\ldots\alpha+\beta+1}; \]
This integral can be expressed in the same case by the formula
\[ Z^{1}=\frac{\Gamma(\alpha+1)\cdot\Gamma(\beta+1)}{(\alpha+\beta+1)\,\Gamma(\alpha+\beta+1)}; \]
and then it holds for all values of $\alpha$ and $\beta$. If we then substitute the values of the functions $\Gamma$ according to the formula
\[ \Gamma(z+1)=\left(\frac{z}{e}\right)^{z+\frac{1}{2}}\sqrt{(2e\pi)}\cdot\Phi(z), \]
$\Phi(z)$ being put for the function $1+\frac{1}{12z}+\frac{1}{288z^{2}}+\text{etc.}$, we find the integral sought
\[ Z^{1}=\frac{\alpha^{\alpha+\frac{1}{2}}\beta^{\beta+\frac{1}{2}}\sqrt{2\pi}}{(\alpha+\beta+1)(\alpha+\beta)^{\alpha+\beta+\frac{1}{2}}}\cdot\frac{\Phi(\alpha)\cdot\Phi(\beta)}{\Phi(\alpha+\beta)}. \]
But by developing the quantity $\frac{\Phi(\alpha)\,\Phi(\beta)}{\Phi(\alpha+\beta)}$, or, what comes to the same, $\Phi(\alpha)\cdot\Phi(\beta)\cdot\Phi(-\alpha-\beta)$, we have for the first three terms of the development
\[ 1+\frac{\alpha^{2}+\alpha\beta+\beta^{2}}{12\alpha\beta(\alpha+\beta)}+\frac{(\alpha^{2}+\alpha\beta+\beta^{2})^{2}}{288\alpha^{2}\beta^{2}(\alpha+\beta)^{2}}. \]
Therefore the value of $Z^{1}$, developed as far as the third term of the series, will be
\[ Z^{1}=\frac{\alpha^{\alpha+\frac{1}{2}}\beta^{\beta+\frac{1}{2}}\sqrt{2\pi}}{(\alpha+\beta)^{\alpha+\beta+\frac{3}{2}}}\left(1+\frac{\alpha^{2}-11\alpha\beta+\beta^{2}}{12\alpha\beta(\alpha+\beta)}+\frac{(\alpha^{2}-11\alpha\beta+\beta^{2})^{2}+144\alpha^{2}\beta^{2}}{288\alpha^{2}\beta^{2}(\alpha+\beta)^{2}}\right), \]
a formula which agrees entirely with the value found for $Z$, which was calculated only as far as the second term.

We now see \emph{a priori} why the value found for $Z$ does not give \uncertain{a} convergent series when only one of the exponents $\alpha$ and $\beta$ \uncertain{is} a large number; it is that the series in question, being equal to the development of the function $\frac{\Phi(\alpha)\,\Phi(\beta)}{\Phi(\alpha+\beta)}\left(\frac{\alpha+\beta}{\alpha+\beta+1}\right)$, is indeed convergent with respect to the function $\Phi(\alpha+\beta)$, and to one of the functions $\Phi(\alpha)$, $\Phi(\beta)$, but is not so with respect to

\origpage{350}
the other; and it follows that the total series is no more convergent than that of the two series $\Phi(\alpha)$, $\Phi(\beta)$ which corresponds to the lesser of the two numbers $\alpha$ and $\beta$.

\subsection*{(35).}

The hypothesis of the preceding art.\ still holding, if we further put $\alpha=\beta$, we shall have the very close value
\[ Z=\left(\tfrac{1}{2}\right)^{2\alpha+1}\sqrt{\left(\frac{\pi}{\alpha}\right)}\cdot\left(1-\frac{3}{8\alpha}+\frac{25}{144\alpha^{2}}-\text{etc.}\right) \]
\ednote{Printed $\frac{25}{144\alpha^{2}}$; the expansion (art.~34 with $\alpha=\beta$) gives $\frac{25}{128\alpha^{2}}$, an apparent misprint.}
The exact value is
\[ Z^{1}=\frac{1\cdot2\ldots\alpha}{\alpha+1\cdot\alpha+2\ldots2\alpha+1}. \]
Thus, when $\alpha$ is very large, we must have
\[ \frac{1\cdot2\cdot3\ldots\alpha}{\alpha+1\cdot\alpha+2\ldots2\alpha+1}=\left(\tfrac{1}{2}\right)^{2\alpha+1}\sqrt{\left(\frac{\pi}{\alpha}\right)}. \]
Here is how this formula may be verified.

Let $\psi(m)=\frac{1\cdot2\cdot3\ldots m}{m+1\cdot m+2\ldots2m+1}$, we shall have
\[ \psi(m+1)=\psi(m)\cdot\frac{(m+1)^{2}}{(2m+2)(2m+3)}=\tfrac{1}{4}\psi(m)\cdot\frac{2m+2}{2m+3}; \]
from this results
\[ \begin{aligned} &\psi(0)=1,\\ &\psi(1)=\tfrac{1}{4}\cdot\tfrac{2}{3},\\ &\psi(2)=\left(\tfrac{1}{4}\right)^{2}\cdot\frac{2\cdot4}{3\cdot5},\\ &\psi(3)=\left(\tfrac{1}{4}\right)^{3}\cdot\frac{2\cdot4\cdot6}{3\cdot5\cdot7},\\ &\text{etc.} \end{aligned} \]
But by the expression of \emph{Wallis}, we have $\frac{\pi}{2}=\frac{2\cdot2}{1\cdot3}\cdot\frac{4\cdot4}{3\cdot5}\cdot\frac{6\cdot6}{5\cdot7}\cdot\frac{8\cdot8}{7\cdot9}$, etc. that is to say that, $m$ being a very large number, we may suppose
\[ \frac{\pi}{2}=\frac{2^{2}\cdot4^{2}\cdot6^{2}\ldots(2m)^{2}}{3^{2}\cdot5^{2}\cdot7^{2}\ldots(2m+1)^{2}}(2m+1); \]
there results $\sqrt{\left(\frac{\pi}{4m+2}\right)}=\frac{2\cdot4\cdot6\ldots2m}{3\cdot5\cdot7\ldots2m+1}$; therefore $\psi(m)=\left(\frac{1}{4}\right)^{m}\sqrt{\left(\frac{\pi}{4m+2}\right)}$, or, since $m$ is very large, $\psi(m)=\left(\frac{1}{2}\right)^{2m+1}\sqrt{\left(\frac{\pi}{m}\right)}$, which agrees with the preceding formula.

\origpage{351}

Our formula is therefore verified in general, and one may be assured that it will give a close approximation when the exponents $\alpha$ and $\beta$ are both large numbers; but if only one of the two is a large number, the series will have only the degree of convergence which the development of $\Phi(\alpha)$ allows, $\alpha$ being the smaller of the two exponents. For the rest, as the value of $Z$ is given generally by the functions $\Gamma$, the observations we have just made relate rather to the general method than to the particular example, of which one will always have the solution as closely approximated as one wishes by the properties of the functions $\Gamma$.

\subsection*{(36).}

We shall remark further that in example II is comprised the general determination of the definite integrals which we have designated above by the formula $\left(\frac{p}{q}\right)=\displaystyle\int x^{p-1}\,dx\,(1-x^{n})^{\frac{q-n}{n}}$. Indeed, if in place of $x^{n}$ one puts $x$, one will have $\left(\frac{p}{q}\right)=\frac{1}{n}\displaystyle\int x^{\frac{p}{n}-1}(1-x)^{\frac{q}{n}-1}\,dx$; let again $\frac{p}{n}=\alpha$, and $\frac{q}{n}=\beta$; one will have $\left(\frac{p}{q}\right)=\frac{1}{n}\displaystyle\int x^{\alpha-1}(1-x)^{\beta-1}\,dx$. But provided that $\alpha$ and $\beta$ are positive, we have between the limits $x=0$, $x=1$,
\[ \int x^{\alpha-1}(1-x)^{\beta-1}\,dx=\frac{(\alpha+\beta)(\alpha+\beta+1)}{\alpha\beta}\int x^{\alpha}(1-x)^{\beta}\,dx; \]
and the integral of the second member $=\frac{\Gamma(\alpha+1)\Gamma(\beta+1)}{(\alpha+\beta+1)\Gamma(\alpha+\beta+1)}$; therefore
\[ \int x^{\alpha-1}(1-x)^{\beta-1}\,dx=\frac{(\alpha+\beta)\Gamma(\alpha+1)\Gamma(\beta+1)}{\alpha\beta\Gamma(\alpha+\beta+1)}=\frac{\Gamma(\alpha)\Gamma(\beta)}{\Gamma(\alpha+\beta)}, \]
and consequently
\[ \left(\frac{p}{q}\right)=\frac{\Gamma(\alpha)\Gamma(\beta)}{n\Gamma(\alpha+\beta)}=\frac{\Gamma\left(\frac{p}{n}\right)\Gamma\left(\frac{q}{n}\right)}{n\Gamma\left(\frac{p+q}{n}\right)}, \]
which is the formula of no.\ 56, second part; we thus have a very simple demonstration of the formula which serves to express the integrals $\left(\frac{p}{q}\right)$ by means of the functions $\Gamma$.

\origpage{352}

\section*{On the integral $Z=\displaystyle\int x^{-\alpha}e^{-x}\,dx$, taken between the imaginary limits which render $x^{-\alpha}e^{-x}$ null.}

\subsection*{(37).}

The quantity $x^{-\alpha}e^{-x}$, in which $\alpha$ is supposed positive, becomes a \emph{minimum} when $x=-\alpha$; but this \emph{minimum} is changed into a \emph{maximum}, if instead of giving real values to $x$, one supposes
\[ x=-\alpha+z\sqrt{-1}, \]
and gives to $z$ any values whatever, from $-\infty$ to $+\infty$. This supposition is admissible analytically, and the consequences deduced from it deserve to be remarked.

According to this value of $x$, the ordinate $x^{-\alpha}e^{-x}$ will be null at the two limits of the integral, namely, when $z=-\infty$ and when $z=+\infty$. One can therefore apply to the integral $Z$ the method given in the preceding chapter.

Before all, I observe that $\alpha$ may be supposed greater than unity, and even as great as one wishes; for we have $d(x^{-m}e^{-x})=-mx^{-m-1}e^{-x}\,dx-x^{-m}e^{-x}\,dx$; and consequently $x^{-m}e^{-x}=-m\displaystyle\int x^{-m-1}e^{-x}\,dx-\displaystyle\int x^{-m}e^{-x}\,dx$; now $m$ being positive, the quantity $x^{-m}e^{-x}$ vanishes at the two limits of the integral; we have therefore
\[ \int x^{-m}e^{-x}\,dx=-m\int x^{-m-1}e^{-x}\,dx; \]
whence we see that if $\alpha$ were less than unity, one could transform the proposed formula into another, in which $\alpha$ would be greater by one unit, and so on. The proposed formula can therefore be prepared in such a way that $\alpha$ is large enough for the series which result from the integration to be convergent.

\subsection*{(38).}

This being laid down, if we remark that the calculation necessary to obtain the value of the integral $\displaystyle\int x^{-\alpha}e^{-x}\,dx$ in the case of imaginary limits, is absolutely the same as that which we made to obtain the integral $\displaystyle\int x^{\alpha}e^{-x}\,dx$ in the case of real limits, we shall see that it suffices to change the sign of $\alpha$ in the integral already found, in order to have the one we seek. Thus since we have found
\[ \int x^{\alpha}e^{-x}\,dx=\left(\frac{\alpha}{e}\right)^{\alpha+\frac{1}{2}}\sqrt{(2e\pi)}\cdot\Phi(\alpha), \]

\origpage{353}
$\Phi(\alpha)$ designating the series $1+\frac{1}{12\alpha}+\frac{1}{288\alpha^{2}}+\text{etc.}$, we shall deduce from it, by simply changing the sign of $\alpha$,
\[ \int x^{-\alpha}e^{-x}\,dx=\left(\frac{-\alpha}{e}\right)^{-\alpha+\frac{1}{2}}\sqrt{(2e\pi)}\cdot\Phi(-\alpha). \]
Multiplying these two equations by one another, and observing that by the property of the function $\Phi$, we have $\Phi(\alpha)\times\Phi(-\alpha)=1$, there will come
\[ \int x^{-\alpha}e^{-x}\,dx\times\int x^{\alpha}e^{-x}\,dx=2\alpha\pi(-1)^{-\alpha+\frac{1}{2}}, \]
an equation by which one will deduce generally the integral $\displaystyle\int x^{-\alpha}e^{-x}\,dx$, whose limits are imaginary, from the integral $\displaystyle\int x^{\alpha}e^{-x}\,dx$, whose limits are real.

This latter integral is represented by $\Gamma(\alpha+1)$, and we have $\Gamma(\alpha+1)=\alpha\Gamma(\alpha)$; therefore if we designate by $\Gamma'(\alpha)$ the integral $\displaystyle\int x^{-\alpha}e^{-x}\,dx$ whose limits are imaginary, we shall have
\[ \Gamma'(\alpha)=\frac{2\pi(-1)^{\frac{1}{2}-\alpha}}{\Gamma(\alpha)}. \tag{1} \]

\subsection*{(39).}

In the particular case where $\alpha$ is an integer, we have $\Gamma(\alpha)=1\cdot2\cdot3\ldots(\alpha-1)$; we shall therefore have in this same case,
\[ \Gamma'(\alpha)=\frac{2\pi(-1)^{\alpha}\sqrt{-1}}{1\cdot2\cdot3\ldots\alpha-1}; \]
which gives successively $\Gamma'(1)=-2\pi\sqrt{-1}$, $\Gamma'(2)=2\pi\sqrt{-1}$,
\[ \Gamma'(3)=-\frac{2\pi\sqrt{-1}}{1\cdot2},\quad\Gamma'(4)=\frac{2\pi\sqrt{-1}}{1\cdot2\cdot3},\quad\Gamma'(5)=-\frac{2\pi\sqrt{-1}}{1\cdot2\cdot3\cdot4},\ \text{etc.} \]

If we put $\alpha=\frac{1}{2}$ (for the formula found, being independent of the series, is no longer subject to the conditions which concern the convergence of those series, and it supposes only $\alpha$ positive), we shall have $\Gamma'\left(\frac{1}{2}\right)=\frac{2\pi}{\Gamma\left(\frac{1}{2}\right)}=2\sqrt{\pi}$; whence we see that $\Gamma'\left(\frac{1}{2}\right)$ is real and double $\Gamma\left(\frac{1}{2}\right)$.

It is remarkable that $\Gamma'\left(\frac{1}{2}\right)$ and $\Gamma\left(\frac{1}{2}\right)$ both represent the integral $\displaystyle\int x^{-\frac{1}{2}}e^{-x}\,dx$; but the first is taken between the imaginary limits which render $x^{-\frac{1}{2}}e^{-x}$ null, and the second is taken between

\origpage{354}
the real limits $x=0$, $x=\infty$; it is therefore not surprising that they are unequal.

In the same way one would distinguish $\Gamma'\left(\frac{1}{3}\right)$, which represents $\displaystyle\int x^{-\frac{1}{3}}e^{-x}\,dx$, taken between imaginary limits, from $\Gamma\left(\frac{2}{3}\right)$, which represents the integral $\displaystyle\int x^{-\frac{1}{3}}e^{-x}\,dx$ taken between real limits. To compare them with one another, we shall apply the general formula which gives
\[ \Gamma'\left(\tfrac{1}{3}\right)=\frac{2\pi\sqrt{-1}}{\Gamma\left(\frac{1}{3}\right)}(-1)^{-\frac{1}{3}}; \]
then, as we have $\Gamma\left(\frac{1}{3}\right)\Gamma\left(\frac{2}{3}\right)=\frac{\pi}{\sin\frac{1}{3}\pi}$, there results
\[ \frac{\Gamma'\left(\frac{1}{3}\right)}{\Gamma\left(\frac{2}{3}\right)}=2\sin\frac{\pi}{3}\cdot(-1)^{\frac{1}{3}}\sqrt{-1}=2\sqrt{-1}\sin\tfrac{1}{3}\pi\left(\cos\frac{2k+1}{3}\pi+\sqrt{-1}\sin\frac{2k+1}{3}\pi\right), \]
which gives three values for the ratio sought of $\Gamma'\left(\frac{1}{3}\right)$ to $\Gamma\left(\frac{2}{3}\right)$. These three values correspond to those which $x^{\frac{1}{3}}$ can take in the formula $\displaystyle\int x^{-\frac{1}{3}}e^{-x}\,dx$, which are $x^{\frac{1}{3}}$, $x^{\frac{1}{3}}\left(-\frac{1}{2}+\frac{1}{2}\sqrt{-3}\right)$, $x^{\frac{1}{3}}\left(-\frac{1}{2}-\frac{1}{2}\sqrt{-3}\right)$; if we confine ourselves to the first, we shall have simply
\[ F'\left(\tfrac{1}{3}\right)=-\sqrt{(-3)}\cdot\Gamma\left(\tfrac{2}{3}\right). \]
\ednote{Printed $F'$ (a clear capital F) where $\Gamma'$ is meant; a misprint.}

\subsection*{(40).}

Let us now resume the formula $Z=\displaystyle\int x^{-\alpha}e^{-x}\,dx$; since we have put $x=-\alpha+z\sqrt{-1}$, and the integral must be taken between the limits $z=-\infty$, $z=+\infty$, it follows that if for brevity we put $\left(-\frac{\alpha}{e}\right)^{-\alpha}=M$, we shall have
\[ Z=M\int\left[\left(1-\frac{z}{\alpha}\sqrt{-1}\right)^{-\alpha}e^{-z\sqrt{-1}}+\left(1+\frac{z}{\alpha}\sqrt{-1}\right)^{-\alpha}e^{z\sqrt{-1}}\right]dz\sqrt{-1}, \]
this integral being taken from $z=0$ to $z=\infty$.

Let $z=\alpha\operatorname{tang}\varphi$, we shall have the transformed expression
\[ Z=2M\alpha\sqrt{-1}\cdot\int d\varphi\cos^{\alpha-2}\varphi\cos(\alpha\operatorname{tang}\varphi-\alpha\varphi), \]
a new integral which will have to be taken from $\varphi=0$ to $\varphi=\frac{1}{2}\pi$.

In applications, one may suppose $\alpha>2$; thus the integral $\displaystyle\int dx\cos^{\alpha-2}\varphi\cos(\alpha\operatorname{tang}\varphi-\alpha\varphi)$\ednote{Printed $dx$ instead of $d\varphi$; an apparent misprint.} will present only ordinary difficulties, when one wishes to evaluate it by the method of quadratures.

\origpage{355}
We shall therefore know, with whatever degree of approximation one wishes, the integral $Z$ or the quantity which we have designated by $\Gamma'(\alpha)$.

Conversely, as the value of the integral $\Gamma'(\alpha)$ is known and represented by the expression $\frac{2\pi(-1)^{\frac{1}{2}-\alpha}}{\Gamma(\alpha)}$, if we put
\[ \int d\varphi\cos^{\alpha-2}\varphi\cos(\alpha\operatorname{tang}\varphi-\alpha\varphi)=Q(\alpha), \]
we shall have
\[ \frac{2\pi(-1)^{\frac{1}{2}-\alpha}}{\Gamma(\alpha)}=2\alpha\sqrt{-1}\cdot\left(-\frac{\alpha}{e}\right)^{-\alpha}Q(\alpha), \]
which gives the integral
\[ Q(\alpha)=\frac{\pi}{\Gamma(\alpha+1)}\cdot\left(\frac{\alpha}{e}\right)^{\alpha}. \tag{2} \]

\subsection*{(41).}

This formula, being independent of the series, must hold whatever $\alpha$ may be, provided it be positive; we shall therefore know $Q(\alpha)$ in all the cases where $\Gamma(\alpha+1)$ is known; these cases are those where $2\alpha$ is an integer.

Thus, putting successively $\alpha=1$, 2, 3, 4, etc., we shall have
\[ \begin{aligned} Q(1)&=\frac{\pi}{e}, &\qquad Q(2)&=\frac{\pi}{2}\cdot\left(\frac{2}{e}\right)^{2}=\frac{2\pi}{e^{2}},\\ Q(3)&=\frac{\pi}{2\cdot3}\cdot\left(\frac{3}{e}\right)^{3}=\frac{9\pi}{2e^{3}}, & Q(4)&=\frac{\pi}{2\cdot3\cdot4}\left(\frac{4}{e}\right)^{4}=\frac{32\pi}{3e^{4}},\ \text{etc.} \end{aligned} \]
Likewise, putting $\alpha=\frac{1}{2}$, $\frac{3}{2}$, $\frac{5}{2}$, etc., we shall have
\[ \begin{aligned} Q\left(\frac{1}{2}\right)&=\sqrt{\left(\frac{2\pi}{e}\right)}, &\qquad Q\left(\frac{3}{2}\right)&=\frac{\sqrt{2\pi}}{1\cdot3}\left(\frac{3}{e}\right)^{\frac{3}{2}},\\ Q\left(\frac{5}{2}\right)&=\frac{\sqrt{2\pi}}{1\cdot3\cdot5}\left(\frac{5}{e}\right)^{\frac{5}{2}}, & Q\left(\frac{7}{2}\right)&=\frac{\sqrt{2\pi}}{1\cdot3\cdot5\cdot7}\left(\frac{7}{e}\right)^{\frac{7}{2}},\ \text{etc.} \end{aligned} \]
In general, if in the value of $Q(\alpha)$ we substitute the known value of $\Gamma(\alpha+1)$ developed in series, we shall have
\[ Q(\alpha)=\left(\frac{\pi}{2\alpha}\right)^{\frac{1}{2}}\left(1-\frac{1}{12\alpha}+\frac{1}{288\alpha^{2}}-\text{etc.}\right). \tag{3} \]
But this formula supposes $\alpha>1$, and it will give a result the more exact the greater $\alpha$ is. When therefore one can

\origpage{356}
neglect $\frac{1}{12\alpha}$ with respect to unity, we shall have $Q(\alpha)=\sqrt{\left(\frac{\pi}{2\alpha}\right)}$, or
\[ \int d\varphi\cos^{\alpha-2}\varphi\cos(\alpha\operatorname{tang}\varphi-\alpha\varphi)=\sqrt{\left(\frac{\pi}{2\alpha}\right)}; \]
a theorem which it would perhaps be very difficult to demonstrate by another route.

The particular theorems concerning $Q(1)$, $Q(2)$, etc. can be presented under another form. Let $\operatorname{tang}\varphi=z$, and suppose that the following integrals are taken from $z=0$ to $z=\infty$, we shall have
\[ \begin{aligned} Q(1)&=\int\frac{dz\,(\cos z+z\sin z)}{1+zz}=\frac{\pi}{e}\\ Q(2)&=\int\frac{dz\,[(1-z^{2})\cos 2z+2z\sin 2z]}{(1+zz)^{2}}=\frac{2\pi}{e^{2}}\\ &\text{etc.} \end{aligned} \tag{4} \]
As these results are deduced from a very thorny analysis, one must be curious to verify them, at least in a particular case. Let, for example, $\alpha=2$, the formula gives $Q(2)=\frac{2\pi}{e^{2}}=0.850337$; we must therefore see whether this value is that of the integral $\displaystyle\int d\varphi\cos(2\operatorname{tang}\varphi-2\varphi)$, taken from $\varphi=0$ to $\varphi=\frac{1}{2}\pi$.

If we put $y=\cos(2\operatorname{tang}\varphi-2\varphi)$, the integral $\displaystyle\int y\,d\varphi$ will be easy to find by approximation, from $\varphi=0$ to the first point where $y=0$. This first point occurs when $\operatorname{tang}\varphi-\varphi=\frac{1}{4}\pi$, and then we find $\varphi=61^{\circ}46'40''$. The following points where $y=0$ are those where $\operatorname{tang}\varphi-\varphi$ has the successive values $\frac{3}{4}\pi$, $\frac{5}{4}\pi$, $\frac{7}{4}\pi$, etc., and their number is visibly infinite. We see therefore that from $\varphi=61^{\circ}46'40''$ to $\varphi=90^{\circ}$, the area of the curve is composed of an infinity of parts alternately positive and negative. These parts are difficult to calculate with a certain degree of precision; but they soon escape by their smallness, and we find that the total result in effect approaches very closely the number given by the preceding formula.

For the rest, we shall see in the preceding chapter that one can, by a direct integration, verify the values of $Q(1)$, $Q(2)$, etc. which will finish dispelling all doubts on the exactness of the preceding formulas. We shall at the same time have occasion to consider

\origpage{357}
a new series of integrals which can be determined generally by the numbers $e$ and $\pi$.

\section*{On the integral $Z=\displaystyle\int\frac{dx\cos ax}{1+xx}$, and others similar, taken from $x=0$ to $x=\infty$.}

\subsection*{(42).}

Let us suppose that the integral is taken from $x=0$ to $x=\frac{2k\pi}{a}$, $k$ being an integer, we shall have, by differentiating with respect to $a$, and having regard to the variability of the second limit,
\[ \frac{dZ}{da}=-\int\frac{x\,dx\sin ax}{1+xx}-\frac{2k\pi}{a^{2}+4k^{2}\pi^{2}}. \]
Differentiating a second time with respect to $a$, we shall have
\[ \frac{ddZ}{da^{2}}=-\int\frac{x^{2}\,dx\cos ax}{1+xx}+\frac{4ka\pi}{(a^{2}+4k^{2}\pi^{2})^{2}}; \]
whence results
\[ Z-\frac{ddZ}{da^{2}}=\int dx\cos ax-\frac{4ka\pi}{(a^{2}+4k^{2}\pi^{2})^{2}}; \]
But we have $\displaystyle\int dx\cos ax=\frac{1}{a}\sin ax$, and this integral vanishes at the supposed limit; therefore we have simply
\[ Z-\frac{ddZ}{da^{2}}=-\frac{4ka\pi}{(a^{2}+4k^{2}\pi^{2})^{2}}. \]
Let us now suppose that $k$ is a very large number with respect to $a$, so that the quantities of the order $\frac{a}{k}$ may be neglected; \emph{a fortiori} those of the order $\frac{a}{k^{3}}$ may be neglected; thus the preceding equation reduces to this
\[ Z-\frac{ddZ}{da^{2}}=0, \]
and there results $Z=Ae^{a}+Be^{-a}$, $A$ and $B$ being two arbitrary constants.

In the case of $a=0$, the integral $Z=\displaystyle\int\frac{dx}{1+xx}=\operatorname{arc}\operatorname{tang}x$, and

\origpage{358}
putting $x=\frac{2k\pi}{a}$, we have $Z=\frac{1}{2}\pi-\frac{a}{2k\pi}$, or simply $Z=\frac{1}{2}\pi$; therefore $A+B=\frac{1}{2}\pi$.

I now suppose that one gives to $a$ values greater and greater; since $\cos ax$ is always smaller than unity, we shall always have $Z<\displaystyle\int\frac{dx}{1+xx}$ or $Z<\frac{1}{2}\pi$. But if $A$ were not zero, the value $Ae^{a}+Be^{-a}$, when $a$ has become a very large number, would reduce to $Ae^{a}$, a quantity infinitely greater than $\frac{1}{2}\pi$; therefore we have generally $A=0$; therefore $B=\frac{1}{2}\pi$, and finally the integral sought
\[ \int\frac{dx\cos ax}{1+xx}=\tfrac{1}{2}\pi e^{-a}. \]
If in this formula we put $\frac{x}{m}$ in place of $x$, and $am$ in place of $a$, we shall have more generally
\[ \int\frac{dx\cos ax}{m^{2}+x^{2}}=\frac{\pi}{2m}e^{-am}; \tag{1} \]
and this formula, being differentiated with respect to $a$, gives a second one no less remarkable, namely,
\[ \int\frac{x\,dx\sin ax}{m^{2}+x^{2}}=\frac{\pi}{2}e^{-am}. \tag{2} \]

\subsection*{(43).}

If we differentiate formula (1) with respect to $m$, and repeat the differentiations, we shall have successively
\[ \begin{aligned} \int\frac{dx\cos ax}{(m^{2}+x^{2})^{2}}&=\frac{\pi e^{-am}}{2^{2}}\left(\frac{a}{m^{2}}+\frac{1}{m^{3}}\right)\\ \int\frac{dx\cos ax}{(m^{2}+x^{2})^{3}}&=\frac{1}{2}\cdot\frac{\pi e^{-am}}{2^{3}}\left(\frac{a^{2}}{m^{3}}+\frac{3a}{m^{4}}+\frac{3}{m^{5}}\right)\\ \int\frac{dx\cos ax}{(m^{2}+x^{2})^{4}}&=\frac{1}{2\cdot3}\cdot\frac{\pi e^{-am}}{2^{4}}\left(\frac{a^{3}}{m^{4}}+\frac{6a^{2}}{m^{5}}+\frac{15a}{m^{6}}+\frac{15}{m^{7}}\right).\\ &\text{etc.} \end{aligned} \]
The law of these expressions is easy to find, and if we put in general,
\[ \int\frac{dx\cos ax}{(m^{2}+x^{2})^{k}}=\frac{A^{k}}{1\cdot2\cdot3\ldots k-1}\cdot\frac{\pi e^{-am}}{2^{k}}, \tag{3} \]
the coefficient $A^{k}$ will have for its value

\origpage{359}
\[ \begin{aligned} A^{k}&=\frac{a^{k-1}}{m^{k}}+\frac{k\cdot k-1}{2}\cdot\frac{a^{k-2}}{m^{k+1}}+\frac{k+1\cdot k\cdot k-1\cdot k-2}{2\cdot4}\cdot\frac{a^{k-3}}{m^{k+2}}\\ &\quad+\frac{k+2\cdot k+1\cdot k\cdot k-1\cdot k-2\cdot k-3}{2\cdot4\cdot6}\cdot\frac{a^{k-4}}{m^{k+3}}+\text{etc.} \end{aligned} \]
Likewise, if we differentiate successively with respect to $m$ the equation (2), we shall deduce from it this series of formulas
\[ \begin{aligned} \int\frac{x\,dx\sin ax}{(m^{2}+x^{2})^{2}}&=\frac{\pi e^{-am}}{2^{2}}\cdot\frac{a}{m}\\ \int\frac{x\,dx\sin ax}{(m^{2}+x^{2})^{3}}&=\frac{1}{1\cdot2}\cdot\frac{\pi e^{-am}}{2^{3}}\left(\frac{a^{2}}{m^{2}}+\frac{a}{m^{3}}\right)\\ \int\frac{x\,dx\sin ax}{(m^{2}+x^{2})^{4}}&=\frac{1}{1\cdot2\cdot3}\cdot\frac{\pi e^{-am}}{2^{4}}\left(\frac{a^{3}}{m^{3}}+\frac{3a^{2}}{m^{4}}+\frac{3a}{m^{5}}\right)\\ &\text{etc.}, \end{aligned} \]
and in general,
\[ \int\frac{x\,dx\sin ax}{(m^{2}+x^{2})^{k}}=\frac{A^{k-1}}{1\cdot2\cdot3\ldots k-1}\cdot\frac{\pi a}{2^{k}}e^{-am}, \tag{4} \]
$A^{k-1}$ being determined according to the same law as $A^{k}$.

\subsection*{(44).}

If in equation (1) we put $m=n(\cos\theta+\sqrt{-1}\sin\theta)$, we shall deduce from it the two following formulas:
\[ \begin{aligned} \int\frac{x^{2}dx\cos ax}{x^{4}+2n^{2}x^{2}\cos2\theta+n^{4}}&=\frac{\pi}{2n\sin2\theta}e^{-an\cos\theta}\sin(\theta-an\sin\theta)\\ \int\frac{dx\cos ax}{x^{4}+2n^{2}x^{2}\cos2\theta+n^{4}}&=\frac{\pi}{2n^{3}\sin2\theta}e^{-an\cos\theta}\sin(\theta+an\sin\theta). \end{aligned} \tag{5} \]
The same substitution being made in formula (2), there will result these two other formulas:
\[ \begin{aligned} \int\frac{x^{3}dx\sin ax}{x^{4}+2n^{2}x^{2}\cos2\theta+n^{4}}&=\frac{\pi}{2\sin2\theta}e^{-an\cos\theta}\sin(2\theta-an\sin\theta)\\ \int\frac{x\,dx\sin ax}{x^{4}+2n^{2}x^{2}\cos2\theta+n^{4}}&=\frac{\pi}{2n^{2}\sin2\theta}e^{-an\cos\theta}\sin(an\sin\theta). \end{aligned} \tag{6} \]
One would find similarly, by means of formulas (3) and (4), the general values of the integrals
\[ \begin{aligned} &\int\frac{dx\cos ax}{(x^{2}+n^{2}\cos2\theta+n^{2}\sin2\theta\cdot\sqrt{-1})^{k}}\\ &\int\frac{x\,dx\sin ax}{(x^{2}+n^{2}\cos2\theta+n^{2}\sin2\theta\cdot\sqrt{-1})^{k}}. \end{aligned} \tag{7} \]

\origpage{360}

\subsection*{(45).}

Let $M$, $N$, $P$ be rational and entire functions of $x^{2}$, if we suppose that the highest exponent of $x$ in $P$ is greater than in $M$ and $N$, and that moreover $P$ has no factor of the form $x^{2}-m^{2}$, that is to say that there is no real value of $x$ which renders $P$ equal to zero; then it is evident that one can generally find, by means of the preceding formulas, the integral
\[ T=\int\left(\frac{M\cos ax+Nx\sin ax}{P}\right)dx, \tag{8} \]
taken from $x=0$ to $x=\infty$. The operation to be done for this is entirely similar to that which is practised for the integration of rational fractions.

It is now easy to apply these formulas to the integrals designated (art 41) by $Q(1)$, $Q(2)$, etc., and we shall find that the results agree perfectly with those which we have deduced by a very different and much less direct method. In general one can verify the value of $Q(\alpha)$, $\alpha$ being any integer whatever; for, putting $\operatorname{tang}\varphi=x$, this integral can always be reduced to a sum of terms of the form $\displaystyle\int\frac{A\,dx\cos ax+Bx\,dx\sin ax}{(1+xx)^{k}}$, and will thus be comprised in formula (8).

\subsection*{(46).}

If we multiply by $da$ both members of the equation $\displaystyle\int\frac{dx\cos ax}{1+xx}=\tfrac{1}{2}\pi e^{-a}$, and then take the integral with respect to $a$, from $a=0$, we shall have
\[ \int\frac{dx\sin ax}{x(1+x^{2})}=\tfrac{1}{2}\pi(1-e^{-a}). \]
Adding this equation to the equation $\displaystyle\int\frac{x\,dx\sin ax}{1+xx}=\tfrac{1}{2}\pi e^{-a}$, we shall have this remarkable result
\[ \int\frac{dx}{x}\sin ax=\tfrac{1}{2}\pi. \tag{9} \]
It appears at first astonishing that this integral should be independent of $a$; but if we put $\frac{x}{a}$ in place of $x$, which does not change the limits $0$ and $\infty$ between which the integral must be comprised; we find $\displaystyle\int\frac{dx}{x}\sin ax=\int\frac{dx}{x}\sin x$; whence it follows that in effect the integral must be independent of $a$.

\origpage{361}

We have therefore also $\displaystyle\int\frac{dx}{x}\sin ax\cos bx=\frac{\pi}{2}$ or zero, according as $a$ is greater or less than $b$; for in the first case we have $\displaystyle\int\frac{dx}{x}\sin ax\cos bx=\tfrac{1}{2}\int\frac{dx}{x}\sin(a+b)x+\tfrac{1}{2}\int\frac{dx}{x}\sin(a-b)x=\tfrac{1}{4}\pi+\tfrac{1}{4}\pi=\tfrac{1}{2}\pi$, and in the second we have $\displaystyle\int\frac{dx}{x}\sin ax\cos bx=\tfrac{1}{2}\int\frac{dx}{x}\sin(a+b)x-\tfrac{1}{2}\int\frac{dx}{x}\sin(b-a)x=\tfrac{1}{4}\pi-\tfrac{1}{4}\pi=0$. As to the integral $\displaystyle\int\frac{dx}{x}\cos ax$, it is evident that it is infinite.

\subsection*{(47).}

Formula (1) and the principal consequences which result from it, are due to Laplace, who published them in the Bulletin de la Société philomatique (April 1811). The demonstration which this illustrious author gives of it is different from that which we have reported; and as it is founded on a principle which may have other applications, we believe that one will be glad to find it here.

Let us consider the double integral
\[ Z=\int2y\,dy\,e^{-y^{2}(1+x^{2})}dx\cos ax, \]
and suppose that it must be taken from $x=0$ to $x=\infty$, and from $y=0$ to $y=\infty$. If we integrate first with respect to $y$, we shall have
\[ Z=\int\frac{dx\cos ax}{1+xx}, \]
this is the integral which it is a question of finding.

If we then return to the double integral, and wish to carry out the integration with respect to $x$, we shall have to seek the value of $\displaystyle\int e^{-x^{2}y^{2}}dx\cos ax$. Now by a formula which will be demonstrated in the following article, we have $\displaystyle\int e^{-x^{2}}dx\cos ax=\tfrac{1}{2}\sqrt{\pi}\cdot e^{-\frac{a^{2}}{4}}$; therefore, putting $\frac{a}{y}$ in place of $a$, and $xy$ in place of $x$, we shall have
\[ \int e^{-x^{2}y^{2}}dx\cos ax=\frac{\sqrt{\pi}}{2y}e^{-\frac{a^{2}}{4y^{2}}}. \]
It remains therefore to integrate the formula
\[ Z=\sqrt{\pi}\int dy\,e^{-y^{2}-\frac{a^{2}}{4y^{2}}}. \]
But we shall demonstrate hereafter (art.\ 50) that between the limits $z=0$ and

\origpage{362}
$z=\infty$, we have $\displaystyle\int z^{-\frac{1}{2}}dz\,e^{-\left(\frac{1+z^{2}}{2nz}\right)}=e^{-\frac{1}{n}}\sqrt{(2n\pi)}$; therefore, putting $a=\frac{1}{n}$, and $y^{2}=\frac{z}{2n}$, we find $Z=\frac{1}{2}\pi e^{-a}$; therefore finally
\[ \int\frac{dx\cos ax}{1+xx}=\frac{\pi}{2}e^{-a}; \]
this is the principal formula from which the others are easy to deduce.

\section*{On the integral $\displaystyle Z=\int e^{-x^{2}}dx\cos ax$, taken from $x=0$ to $x=\infty$.}

\subsection*{(48).}

If we differentiate this formula with respect to $a$, we shall have
\[ \frac{dZ}{da}=-\int e^{-x^{2}}x\,dx\sin ax=\tfrac{1}{2}e^{-x^{2}}\sin ax-\frac{a}{2}\int e^{-x^{2}}dx\cos ax. \]
The part outside the sign vanishes at the two limits of the integral; thus we have $\frac{dZ}{da}=-\frac{a}{2}Z$, and on integrating $Z=Ae^{-\frac{a^{2}}{4}}$.

To determine the constant $A$, let $a=0$, then $\displaystyle Z=\int e^{-x^{2}}dx=\tfrac{1}{2}\sqrt{\pi}$; therefore $Z$ or
\[ \int e^{-x^{2}}dx\cos ax=\frac{\sqrt{\pi}}{2}e^{-\frac{a^{2}}{4}}. \tag{1} \]
Laplace, who gave this formula in the Mémoires de l'Institut, ann.\ 1809, pag.\ 367, demonstrates it in the following manner.

If one substitutes in place of $\cos ax$ its value developed in series, one will have
\[ Z=\int e^{-x^{2}}dx\left(1-\frac{a^{2}x^{2}}{2}+\frac{a^{4}x^{4}}{2\cdot3\cdot4}-\text{etc.}\right) \]
Now in general $\displaystyle\int x^{2n}e^{-x^{2}}dx=\tfrac{1}{2}\Gamma\left(\frac{2n+1}{2}\right)$; therefore we have
\[ Z=\frac{\sqrt{\pi}}{2}\left(1-\frac{a^{2}}{4}+\frac{1}{2}\cdot\frac{a^{4}}{16}-\frac{1}{2\cdot3}\cdot\frac{a^{6}}{64}+\text{etc.}\right), \]
or
\[ Z=\frac{\sqrt{\pi}}{2}e^{-\frac{a^{2}}{4}}. \]

\subsection*{(49).}

If we take the successive differentials of equation (1) with respect to $a$, we shall deduce from it this series of integrals,

\origpage{363}
\[ \begin{aligned} \int e^{-x^{2}}x\,dx\sin ax&=\frac{\sqrt{\pi}}{4}e^{-\frac{a^{2}}{4}}a\\ \int e^{-x^{2}}x^{2}dx\cos ax&=\frac{\sqrt{\pi}}{4}e^{-\frac{a^{2}}{4}}\left(1-\frac{a^{2}}{2}\right)\\ \int e^{-x^{2}}x^{3}dx\sin ax&=\frac{\sqrt{\pi}}{8}e^{-\frac{a^{2}}{4}}\left(3a-\frac{a^{3}}{2}\right)\\ \int e^{-x^{2}}x^{4}dx\cos ax&=\frac{\sqrt{\pi}}{8}e^{-\frac{a^{2}}{4}}\left(3-3a^{2}+\frac{a^{4}}{4}\right)\\ \int e^{-x^{2}}x^{5}dx\sin ax&=\frac{\sqrt{\pi}}{16}e^{-\frac{a^{2}}{4}}\left(15a-5a^{3}+\frac{a^{5}}{4}\right)\\ &\text{etc.} \end{aligned} \tag{2} \]
from which it is seen that one can find in general the integral
\[ \int(M\cos ax+Nx\sin ax)e^{-x^{2}}dx, \]
$M$ and $N$ being rational and integral functions of $x^{2}$.

If it were proposed to find between the same limits the integral $\displaystyle T=\int e^{-x^{2}}dx\sin ax$, we would have first $\displaystyle\frac{dT}{da}=\int e^{-x^{2}}x\,dx\cos ax=\tfrac{1}{2}-\tfrac{1}{2}e^{-x^{2}}\cos ax-\tfrac{1}{2}a\int e^{-x^{2}}dx\sin ax$. Making $x=\infty$ in the the\ednote{The word ``la'' is printed twice across the line break; an apparent misprint.} part outside the sign, there would result $\frac{dT}{da}=\frac{1}{2}-\frac{1}{2}aT$, whence $\displaystyle T=\tfrac{1}{2}e^{-\frac{a^{2}}{4}}\int e^{\frac{a^{2}}{4}}da$. This integral, taken from $a=0$, is simpler than the proposed one; but its expression can be found only by this convergent series
\[ T=\frac{1}{2}\left(a-\frac{a^{3}}{2\cdot3}+\frac{a^{5}}{3\cdot4\cdot5}-\frac{a^{7}}{4\cdot5\cdot6\cdot7}+\text{etc.}\right), \]
which does not appear capable of being reduced to the known transcendents.

$T$ being supposed known, one will deduce from it by successive differentiations,
\[ \begin{aligned} \int e^{-x^{2}}dx\sin ax&=T\\ \int e^{-x^{2}}x\,dx\cos ax&=\frac{dT}{da}=\frac{1}{2}-\frac{1}{2}aT\\ \int e^{-x^{2}}x^{2}dx\sin ax&=-\frac{ddT}{da^{2}}=\frac{1}{4}a+\frac{1}{2}T\left(1-\frac{a^{2}}{2}\right)\\ \int e^{-x^{2}}x^{3}dx\cos ax&=-\frac{d^{3}T}{da^{3}}=\frac{1}{2}-\frac{a^{2}}{8}-\frac{1}{4}T\left(3a-\frac{a^{3}}{2}\right)\\ \int e^{-x^{2}}x^{4}dx\sin ax&=\frac{d^{4}T}{da^{4}}=\frac{5a}{8}-\frac{a^{3}}{16}+\frac{1}{4}T\left(3-3a^{2}+\frac{a^{4}}{4}\right)\\ &\text{etc.}; \end{aligned} \]

\origpage{364}
and since $T$ is expressed by a very convergent series, it is seen that it will be easy to find in all cases the closely approximate values of these various integrals.

Moreover, by combining these formulas, one could find some rather remarkable results, such as the following two:
\[ \begin{aligned} \int e^{-x^{2}}dx\left(\tfrac{1}{2}a\sin ax+x\cos ax\right)&=\tfrac{1}{2}\\ \int e^{-x^{2}}dx\left(1-\tfrac{1}{2}a^{2}-2x^{2}\right)\sin ax&=-\tfrac{1}{2}a. \end{aligned} \]

\section*{On the integral $\displaystyle Z(k)=\int x^{\frac{2k-1}{2}}dx\,e^{-\left(\frac{1+xx}{2nx}\right)}$, taken from $x=0$ to $x=\infty$.}

\subsection*{(50).}

Let us consider first the integral
\[ Z(0)=\int x^{-\frac{1}{2}}dx\cdot e^{-\left(\frac{1+zz}{2nx}\right)}, \]
\ednote{Printed $1+zz$ in the numerator of the exponent instead of $1+xx$; an apparent misprint.} and let us divide it into two parts, one from $x=0$ to $x=1$; the other from $x=1$ to $x=\infty$. To have the first part, let $x=\frac{1}{z}$; we shall have the transformed integral $\displaystyle\int-z^{\frac{3}{2}}dz\,e^{-\left(\frac{1+2z}{2nz}\right)}$\ednote{Printed $z^{\frac{3}{2}}$ instead of $z^{-\frac{3}{2}}$, and $1+2z$ in the numerator of the exponent where $1+zz$ is meant (the printed Errata, p.\ 387, corrects it to $1+zz$); apparent misprints.}, which will have to be integrated from $z=\infty$ to $z=1$; if its sign is changed, it will have to be integrated from $z=1$ to $z=\infty$; thus, uniting the two parts, it is seen that everything reduces to finding the integral
\[ \int\left(x^{-\frac{1}{2}}+x^{-\frac{3}{2}}\right)dx\cdot e^{-\left(\frac{1+xx}{2nx}\right)} \]
between the limits $x=1$ and $x=\infty$.

Let $1+x^{2}=2xz$; we shall have $x-1=x^{\frac{1}{2}}\sqrt{(2z-2)}$, and consequently
\[ \left(x^{-\frac{1}{2}}+x^{-\frac{3}{2}}\right)dx=\frac{dz\sqrt{2}}{\sqrt{(z-1)}}. \]
Hence the integral in question will have for its transformed integral
\[ Z(0)=\sqrt{2}\int\frac{dz}{\sqrt{(z-1)}}e^{-\frac{z}{n}}, \]
this being to be taken from $z=1$ to $z=\infty$.

Let $z=1+y^{2}$, and we shall have a new transformed integral
\[ Z(0)=2^{\frac{3}{2}}e^{-\frac{1}{n}}\int dy\,e^{-\frac{y^{2}}{n}}, \]

\origpage{365}
which will have to be integrated from $y=0$ to $y=\infty$. But by making, according to our custom, $e^{-\frac{y^{2}}{n}}=u$, we shall have
\[ Z(0)=(2n)^{\frac{1}{2}}e^{-\frac{1}{n}}\int du\left(l\frac{1}{u}\right)^{-\frac{1}{2}}; \]
and as this integral must be taken from $u=0$ to $u=1$, we shall have finally
\[ Z(0)=e^{-\frac{1}{n}}\sqrt{(2n\pi)}. \tag{1} \]

\subsection*{(51).}

Let us consider in the second place the formula
\[ Z(1)=\int x^{\frac{1}{2}}dx\,e^{-\left(\frac{1+xx}{2nx}\right)}; \]
if it is divided into two parts, as in the preceding article, we shall have
\[ Z(1)=\int\left(x^{\frac{1}{2}}+x^{-\frac{5}{2}}\right)dx\,e^{-\left(\frac{1+xx}{2nx}\right)}, \]
which will have to be integrated from $x=1$ to $x=\infty$. Now, making $1+x^{2}=2xz$, then $z=1+y^{2}$, we shall have successively
\[ \begin{aligned} x^{\frac{3}{2}}&=(2z-1)\sqrt{\left(\frac{z+1}{2}\right)}+(2z+1)\sqrt{\left(\frac{z-1}{2}\right)}\\ x^{-\frac{3}{2}}&=(2z-1)\sqrt{\left(\frac{z+1}{2}\right)}-(2z+1)\sqrt{\left(\frac{z-1}{2}\right)}\\ x^{\frac{3}{2}}-x^{-\frac{3}{2}}&=(2z+1)\sqrt{(2z-2)}=(2y^{3}+3y)\sqrt{2}\\ \left(x^{\frac{1}{2}}+x^{-\frac{5}{2}}\right)dx&=2^{\frac{3}{2}}(2y^{2}+1)\,dy. \end{aligned} \]
Hence the transformed integral in $y$ will be
\[ Z(1)=2^{\frac{3}{2}}e^{-\frac{1}{n}}\int(1+2y^{2})e^{-\frac{y^{2}}{n}}dy. \]
Let again $e^{-\frac{y^{2}}{n}}=u$, and we shall have the new transformed integral
\[ Z(1)=(2n)^{\frac{1}{2}}e^{-\frac{1}{n}}\int du\left[\left(l\frac{1}{u}\right)^{-\frac{1}{2}}+2n\left(l\frac{1}{u}\right)^{\frac{1}{2}}\right]; \]
whence results, on integrating,
\[ Z(1)=(2n)^{\frac{1}{2}}e^{-\frac{1}{n}}\left[\Gamma(\tfrac{1}{2})+2n\Gamma(\tfrac{3}{2})\right]. \]
Substituting the known values $\Gamma(\frac{1}{2})=\sqrt{\pi}$, $\Gamma(\frac{3}{2})=\frac{1}{2}\sqrt{\pi}$, we

\origpage{366}
shall have finally
\[ Z(1)=(1+n)e^{-\frac{1}{n}}\sqrt{(2n\pi)}=(1+n)Z(0). \tag{2} \]

\subsection*{(52).}

One could calculate in the same manner the values of $Z(2)$, $Z(3)$, etc. But the law of progression of these quantities is easy to find; indeed, if one differentiates the quantity\ldots\ $T=x^{\frac{2k+1}{2}}e^{-\left(\frac{1+xx}{2nx}\right)}$, we shall have
\[ dT=\frac{2k+1}{2}x^{\frac{2k-1}{2}}dx\cdot e^{-\left(\frac{1+xx}{2nx}\right)}+x^{\frac{2k+1}{2}}\left(\frac{1}{2nx^{2}}-\frac{1}{2n}\right)dx\cdot e^{-\left(\frac{1+xx}{2nx}\right)}. \]
Integrating and observing that $T$ vanishes at the two limits of the integral, we find
\[ 0=\frac{2k+1}{2}Z(k)+\frac{1}{2n}Z(k-1)-\frac{1}{2n}Z(k+1); \]
whence results
\[ Z(k+1)=(2k+1)nZ(k)+Z(k-1). \tag{3} \]
Thus we shall have successively
\[ \begin{aligned} Z(2)&=3nZ(1)+Z(0)=Z(0)(1+3n+3n^{2})\\ Z(3)&=5nZ(2)+Z(1)=Z(0)(1+6n+15n^{2}+15n^{3})\\ &\text{etc.} \end{aligned} \]
The general expression of these quantities is
\[ \begin{aligned} Z(k)=Z(0)\cdot\Big[1&+\frac{k+1\cdot k}{2}n+\frac{k+2\cdot k+1\cdot k\cdot k-1}{2\cdot4}n^{2}\\ &+\frac{k+3\cdot k+2\cdot k+1\cdot k\cdot k-1\cdot k-2}{2\cdot4\cdot6}n^{3}+\text{etc.}\Big]; \end{aligned} \tag{4} \]
which can easily be verified by substitution in equation (3). The series which we meet with here therefore follows the same law as that which was represented above by $A^{k}$ (art.\ 43).

This formula also holds when $k$ is negative, and it gives immediately $Z(-k-1)=Z(k)$; this is also what results from the direct consideration of the integrals. Indeed we have generally
\[ Z(k)=\int\left(x^{\frac{2k-1}{2}}+x^{-\left(\frac{2k+3}{2}\right)}\right)dx\,e^{-\left(\frac{1+xx}{2nx}\right)}, \]
this integral being taken from $x=1$ to $x=\infty$.

\origpage{367}

We would have similarly
\[ Z(-k-1)=\int\left(x^{-\left(\frac{2k+3}{2}\right)}+x^{\frac{2k-1}{2}}\right)dx\,e^{-\left(\frac{1+xx}{2nx}\right)}; \]
therefore $Z(-k-1)=Z(k)$.

\subsection*{(53).}

Euler, in volume IV of the \emph{Supplémens au Calcul intégral}, pag.\ 415, makes mention of the integrals
\[ A=\frac{1}{2n}\int x^{-\frac{5}{2}}dx\,e^{-\left(\frac{1+xx}{2nx}\right)},\qquad B=\frac{1}{2n}\int x^{-\frac{3}{2}}dx\,e^{-\left(\frac{1+xx}{2nx}\right)}, \]
which seem to him incapable of being brought back to the known methods. This difficulty is resolved by the preceding formulas, which give
\[ \begin{aligned} A&=\frac{1}{2n}Z(-2)=\frac{1}{2n}Z(1)=(1+n)e^{-\frac{1}{n}}\sqrt{\left(\frac{\pi}{2n}\right)}\\ B&=\frac{1}{2n}Z(-1)=\frac{1}{2n}Z(0)=e^{-\frac{1}{n}}\sqrt{\left(\frac{\pi}{2n}\right)}, \end{aligned} \]
whence $\frac{A}{B}=1+n$. The same author adds that if the values of these two integrals cannot be found separately, at least their ratio is known, namely,
\[ \frac{A}{B}=\frac{1+e^{\frac{2}{n}}}{1-e^{\frac{2}{n}}}. \]
But there is evidently an error in the calculations which led to this result, since it would give a negative value of $\frac{A}{B}$, whereas $A$ and $B$ are positive.

\section*{On the integrals $\displaystyle\int x^{a-1}e^{-mx}dx\cos nx$, $\displaystyle\int x^{a-1}e^{-mx}dx\sin nx$, taken between the limits $x=0$, $x=\infty$.}

\subsection*{(54).}

We shall suppose that $a$ and $m$ are positive, a condition necessary in order that the integrals in question be finite quantities. To find their values, let us consider first the formula
\[ Z=\int x^{a-1}dx\,e^{-(m-n\sqrt{-1})x}; \]

\origpage{368}
if we make $e^{-(m-n\sqrt{-1})x}=z$, we shall have the transformed integral
\[ Z=\int\frac{dz\left(l\frac{1}{z}\right)^{a-1}}{(m-n\sqrt{-1})^{a}}, \]
which will have to be integrated from $z=0$ to $z=1$. Let therefore $\frac{n}{m}=\operatorname{tang}\theta$ and $\sqrt{(m^{2}+n^{2})}=r$, we shall have
\[ Z=\frac{\cos a\theta+\sqrt{-1}\sin a\theta}{r^{a}}\Gamma(a). \]
It is now sufficient to substitute in the proposed formulas the values of $\cos nx$ and $\sin nx$ in imaginary exponentials, and one will conclude immediately
\[ \begin{aligned} \int x^{a-1}e^{-mx}dx\cos nx&=\frac{\cos a\theta}{r^{a}}\Gamma(a)\\ \int x^{a-1}e^{-mx}dx\sin nx&=\frac{\sin a\theta}{r^{a}}\Gamma(a). \end{aligned} \tag{1} \]
By means of these formulas, one will find the values of the integrals $\displaystyle\int x^{a-1}e^{-mx}dx\sin^{k}nx$, $\displaystyle\int x^{a-1}e^{-mx}dx\cos^{k}nx$, and in general that of the integral $\displaystyle\int Tx^{a-1}e^{-mx}dx$, $T$ being a function of sines and cosines which can be developed into a finite series of linear sines and cosines of the form $A\sin(\alpha x+\beta)+$ etc.

If the number $a$ is an integer, the integrals (1) can be found by the ordinary processes of integration, and we shall moreover have $\Gamma(a)=1\cdot2\cdot3\ldots a-1$, which will permit these formulas to be verified.

\subsection*{(55).}

If one makes $m=0$ or $\theta=\frac{1}{2}\pi$, the formulas (1) will reduce to the following two:
\[ \begin{aligned} \int x^{a-1}dx\cos nx&=\frac{\cos\frac{1}{2}a\pi}{n^{a}}\Gamma(a)\\ \int x^{a-1}dx\sin nx&=\frac{\sin\frac{1}{2}a\pi}{n^{a}}\Gamma(a); \end{aligned} \tag{2} \]
and in particular when $a=\frac{1}{2}$, we shall have
\[ \begin{aligned} \int\frac{dx\cos nx}{\sqrt{x}}&=\sqrt{\left(\frac{\pi}{2n}\right)}\\ \int\frac{dx\sin nx}{\sqrt{x}}&=\sqrt{\left(\frac{\pi}{2n}\right)}. \end{aligned} \tag{3} \]
If in the formulas (1) one supposes $a$ infinitely small, which

\origpage{369}
gives $\Gamma(a)=\frac{1}{a}$, we shall have
\[ \begin{aligned} \int\frac{dx}{x}e^{-mx}\cos nx&=\infty,\\ \int\frac{dx}{x}e^{-mx}\sin nx&=\theta=\operatorname{arc}\operatorname{tang}\frac{n}{m}. \end{aligned} \tag{4} \]

If in this last one makes $m=0$, one recovers the formula of no.\ 46.

The integrals with which we have just been occupied are found in the fourth volume of Euler's Calcul intégral, pag.\ 337 et suiv.\ They have also been treated by Laplace, in volume VIII of the Journal de l'École Polytechnique, pag.\ 244 et suiv.

\subsection*{(56).}

If one differentiates the formulas (1) with respect to $a$, one will deduce from them
\[ \begin{aligned} \int x^{a-1}e^{-mx}dx\cos nx\log x&=\frac{\cos a\theta}{r^{a}}\cdot\frac{d\Gamma(a)}{da}-\left(\frac{\theta\sin a\theta+\cos a\theta\log r}{r^{a}}\right)\Gamma(a)\\ \int x^{a-1}e^{-mx}dx\sin nx\log x&=\frac{\sin a\theta}{r^{a}}\cdot\frac{d\Gamma(a)}{da}+\left(\frac{\theta\cos a\theta-\sin a\theta\log r}{r^{a}}\right)\Gamma(a). \end{aligned} \tag{5} \]
The new transcendent $\frac{d\Gamma(a)}{da}$ which enters into these formulas, can be found approximately from the tables; its value is also had in convergent series by the formulas of no.\ 76, IInd part.

If from these last two equations one eliminates $\frac{d\Gamma(a)}{da}$, one draws from them this remarkable result:
\[ \int x^{a-1}e^{-mx}dx\sin(a\theta-nx)\log\frac{1}{x}=\frac{\theta}{r^{a}}\Gamma(a), \tag{6} \]
and in particular when $a=1$,
\[ \int e^{-mx}dx\sin(\theta-nx)\log\frac{1}{x}=\frac{\theta}{r}; \tag{7} \]
a formula in which we have $\operatorname{tang}\theta=\frac{n}{m}$ and $r=\sqrt{(m^{2}+n^{2})}$.

\origpage{370}

\section*{On the integral $\displaystyle\int\frac{(x^{n}-1)\,dx}{\log x}$ and others like it, taken from $x=0$ to $x=1$.}

\subsection*{(57).}

The integral $\displaystyle\int\frac{dx}{\log x}$ taken between these limits, is infinite; the integral $\displaystyle\int\frac{x^{n}dx}{\log x}$ reduces to the same form $\displaystyle\int\frac{du}{\log u}$, by making $x^{n+1}=u$; it is therefore also infinite; but the difference of these infinites is a finite quantity which is to be determined.

For this, let $\displaystyle P=\int\frac{(x^{n}-1)\,dx}{lx}$; if this equation is differentiated with respect to $n$, we shall have
\[ \frac{dP}{dn}=\int x^{n}dx=\frac{1}{n+1}; \]
therefore $dP=\frac{dn}{n+1}$, and consequently $P=\log(n+1)$, without a constant, because $P$ must vanish when $n=0$. We have therefore between the given limits,
\[ \int(x^{n}-1)\frac{dx}{lx}=\log(n+1). \tag{1} \]
From this it is easy to find, between the same limits, the integral $\displaystyle\int\frac{(x^{n}-1)x^{m}dx}{lx}$. Let for this purpose $x^{m+1}=u$, and $\frac{n}{m+1}=\alpha$, we shall have
\[ \int\frac{(x^{n}-1)x^{m}dx}{lx}=\int\frac{(u^{\alpha}-1)\,du}{lu}=\log(\alpha+1)=\log\left(\frac{m+n+1}{m+1}\right); \]
therefore
\[ \int\frac{(x^{n}-1)x^{m}dx}{lx}=\log\left(\frac{m+n+1}{m+1}\right). \tag{2} \]
One could arrive immediately at this result by observing that $(x^{n}-1)x^{m}=x^{m+n}-1-(x^{m}-1)$; which gives
\[ \int\frac{(x^{n}-1)x^{m}dx}{lx}=\int(x^{m+n}-1)\frac{dx}{lx}-\int(x^{m}-1)\frac{dx}{lx}=\log\left(\frac{m+n+1}{m+1}\right). \]
In general, if we have a polynomial
\[ X=Ax^{n}+Bx^{n-1}+Cx^{n-2}+\text{etc.}, \]
which reduces to zero when $x=1$, it may be put in the form
\[ X=A(x^{n}-1)+B(x^{n-1}-1)+C(x^{n-2}-1)+\text{etc.}; \]

\origpage{371}
and then it is clear that we shall have
\[ \int\frac{X\,dx}{lx}=A\log(n+1)+B\log(n)+C\log(n-1)+\text{etc.} \]

\subsection*{(58).}

Let us suppose now that one wishes to have the integral
\[ Q=\int\frac{(x^{n}-1)^{2}dx}{(lx)^{2}}; \]
differentiating with respect to $n$, we shall have
\[ \frac{dQ}{dn}=\int\frac{2(x^{n}-1)x^{n}dx}{lx}=2\log(2n+1)-2\log(n+1); \]
whence results, on integrating,
\[ Q=(2n+1)\log(2n+1)-2(n+1)\log(n+1). \tag{3} \]
Let the integral $\displaystyle\int\frac{(x^{n}-1)^{2}x^{m}dx}{(lx)^{2}}$ be proposed more generally; we shall make $x^{m+1}=u$, $\frac{n}{m+1}=\alpha$, and the transformed integral will be $\displaystyle(m+1)\int\frac{(u^{\alpha}-1)^{2}du}{lu}$.\ednote{Printed $lu$ in the denominator; the substitution gives $(lu)^{2}$, as the value that follows requires. An apparent misprint.} This one, according to the preceding formula, has for its value,
\[ (m+1)(2\alpha+1)\log(2\alpha+1)-2(m+1)(\alpha+1)\log(\alpha+1). \]
Therefore, restoring the value of $\alpha$, we shall have
\[ \begin{aligned} \int\frac{(x^{n}-1)^{2}x^{m}dx}{(lx)^{2}}={}&(2n+m+1)\log(2n+m+1)-2(n+m+1)\log(n+m+1)\\ &+(m+1)\log(m+1), \end{aligned} \]
an integral which could be represented more simply by the formula
\[ \int\left(\frac{x^{n}-1}{lx}\right)^{2}x^{m}dx=\frac{1}{2}\Delta^{3}[(m+1)^{2}\log(m+1)], \tag{4} \]
\ednote{The right side of (4) is printed as read here, the same form as in formula (5); by the general formulae of art.\ 60 it should be $\Delta^{2}[(m+1)\log(m+1)]$. An apparent misprint. A reader has struck the $\frac{1}{2}$ and the exponents in pencil, so the printed exponents under the pencil are read with some doubt.}
supposing that the difference $\Delta m$ is equal to $n$.

\subsection*{(59).}

We shall find similarly
\[ \int\left(\frac{x^{n}-1}{lx}\right)^{3}x^{m}dx=\frac{1}{2}\Delta^{3}[(m+1)^{2}\log(m+1)]. \]
Indeed, if we make the first member $=R$, we shall have
\[ \frac{dR}{dn}=3\int\left(\frac{x^{n}-1}{lx}\right)^{2}x^{m+n}\,dx, \]
or
\[ \begin{aligned} \frac{dR}{dn}={}&3(3n+m+1)\log(3n+m+1)-6(2n+m+1)\log(2n+m+1)\\ &+3(n+m+1)\log(n+m+1). \end{aligned} \]

\origpage{372}
Multiplying by $dn$, and integrating the second member by the formula $\displaystyle\int x\,dx\log x=\frac{x^{2}}{2}(\log x-\tfrac{1}{2})$; determining then the constant in such a way that the integral shall be zero when $n=0$, we shall have
\[ \begin{aligned} R={}&\tfrac{1}{2}(3n+m+1)^{2}\log(3n+m+1)-\tfrac{3}{2}(2n+m+1)^{2}\log(2n+m+1)\\ &+\tfrac{3}{2}(n+m+1)^{2}\log(n+m+1)-\tfrac{1}{2}(m+1)^{2}\log(m+1), \end{aligned} \]
a quantity which can be put in the form $\frac{1}{2}\Delta^{3}[(m+1)^{2}\log(m+1)]$; we shall therefore have
\[ \int\left(\frac{x^{n}-1}{lx}\right)^{3}x^{m}dx=\frac{1}{2}\Delta^{3}[(m+1)^{2}\log(m+1)]. \tag{5} \]

\subsection*{(60).}

These formulas can be presented in the simplest and most general manner, as follows:
\[ \begin{aligned} \int\left(\frac{x^{n}-1}{lx}\right)x^{m-1}dx&=\Delta(\log m)\\ \int\left(\frac{x^{n}-1}{lx}\right)^{2}x^{m-1}dx&=\Delta^{2}(m\log m)\\ \int\left(\frac{x^{n}-1}{lx}\right)^{3}x^{m-1}dx&=\frac{1}{2}\Delta^{3}(m^{2}\log m)\\ \int\left(\frac{x^{n}-1}{lx}\right)^{4}x^{m-1}dx&=\frac{1}{2\cdot3}\Delta^{4}(m^{3}\log m)\\ \text{etc.}\qquad& \end{aligned} \]
the differences indicated by $\Delta$ being taken on the supposition that $\Delta m=n$.

Euler has considered these integrals in only a few cases; see volume IV of his \emph{Calcul intégral}, pag.~271 et suiv.

\section*{On the integrals $\displaystyle\int\frac{d\varphi\cos\lambda\varphi}{\Delta^{n}}$ and $\displaystyle\int\Delta^{n}d\varphi\cos\lambda\varphi$, taken from $\varphi=0$ to $\varphi=\pi$, supposing the numbers $\lambda$ and $n$ integers, and $\Delta=1+a^{2}-2a\cos\varphi$.}

\subsection*{(61).}

From the theory of recurring series, it is easy to see that we have the equation
\[ \frac{1-a\cos\varphi}{1-2a\cos\varphi+a^{2}}=1+a\cos\varphi+a^{2}\cos2\varphi+a^{3}\cos3\varphi+\text{etc.} \]
Multiplying each member by 2 and subtracting unity from both sides, there will result
\[ \frac{1-a^{2}}{1-2a\cos\varphi+a^{2}}=1+2a\cos\varphi+2a^{2}\cos2\varphi+2a^{3}\cos3\varphi+\text{etc.} \]

\origpage{373}
therefore
\[ \frac{d\varphi\cos\lambda\varphi}{\Delta}=\frac{d\varphi\cos\lambda\varphi}{1-a^{2}}(1+2a\cos\varphi+2a^{2}\cos2\varphi+2a^{3}\cos3\varphi+\text{etc.}). \]
Let us call in general $P^{n}$ the integral $\displaystyle\int\frac{d\varphi\cos\lambda\varphi}{\Delta^{n}}$, taken between the limits $\varphi=0$, $\varphi=\pi$; to have the value of $P^{1}$, it will be necessary to integrate the second member of the preceding equation. Now it is evident that for this purpose, one can reduce the series $1+2a\cos\varphi+2a^{2}\cos2\varphi+\text{etc.}$, to the single term $2a^{\lambda}\cos\lambda\varphi$; for $k$ being different from $\lambda$, we have\ldots\ $\displaystyle\int d\varphi\cos\lambda\varphi\cos k\varphi=\frac{1}{2}\int d\varphi\cos(\lambda-k)\varphi+\frac{1}{2}\int d\varphi\cos(\lambda+k)\varphi=\frac{\sin(\lambda-k)\varphi}{2(\lambda-k)}+\frac{\sin(\lambda+k)\varphi}{2(\lambda+k)}$, an integral which vanishes between the given limits, since $\lambda$ and $k$ are integer numbers different from one another. We shall therefore have simply
\[ P^{1}=\int\frac{d\varphi\cos\lambda\varphi}{1-a^{2}}\cdot2a^{\lambda}\cos\lambda\varphi=\frac{a^{\lambda}}{1-a^{2}}\int d\varphi(1+\cos2\lambda\varphi). \]
But the integral $\displaystyle\int d\varphi(1+\cos2\lambda\varphi)$, taken between the limits $\varphi=0$, $\varphi=\pi$, reduces to $\pi$; therefore finally
\[ P^{1}=\frac{\pi a^{\lambda}}{1-a^{2}}. \tag{1} \]
We have tacitly supposed that $a$ was smaller than unity; if it were greater, we would make $a=\frac{1}{\alpha}$, and then $\Delta$ would become $\frac{1}{\alpha^{2}}(1+\alpha^{2}-2\alpha\cos\varphi)$, so that the integral $\displaystyle\int\frac{d\varphi\cos\lambda\varphi}{\Delta^{n}}$ would always be brought back to a similar integral in which we would have $a<1$.

\subsection*{(62).}

Knowing thus the value of $P^{1}$, one can deduce from it by successive differentiations, the values of $P^{2}$, $P^{3}$, etc. Indeed if one differentiates with respect to $a$ the equation $\displaystyle P^{n}=\int\frac{d\varphi\cos\lambda\varphi}{\Delta^{n}}$, one draws from it
\[ \frac{dP^{n}}{da}=\int\frac{n\,d\varphi\cos\lambda\varphi}{\Delta^{n+1}}\left(\frac{1-a^{2}-\Delta}{a}\right), \]
or
\[ \frac{dP^{n}}{da}=\frac{n}{a}(1-a^{2})P^{n+1}-\frac{n}{a}P^{n}; \]
therefore
\[ P^{n+1}=\frac{1}{1-a^{2}}\cdot\left(P^{n}+\frac{n}{a}\cdot\frac{dP^{n}}{da}\right)=\frac{d(a^{n}P^{n})}{(1-a^{2})\,d(a^{n})}. \tag{2} \]
\ednote{Printed $\frac{n}{a}\cdot\frac{dP^{n}}{da}$; the preceding equation and the last member of (2) require $\frac{a}{n}\cdot\frac{dP^{n}}{da}$. An apparent misprint, corrected in pencil by a reader.}
By means of this equation and the known value of $P^{1}$, we shall find successively

\origpage{374}
\[ \begin{aligned} P^{2}&=\frac{\pi a^{\lambda}}{(1-a^{2})^{3}}\{\lambda+1-a^{2}(\lambda-1)\}\\ P^{3}&=\frac{1}{1\cdot2}\cdot\frac{\pi a^{\lambda}}{(1-a^{2})^{5}}\{(\lambda+1)(\lambda+2)-2a^{2}(\lambda+2)(\lambda-2)+a^{4}(\lambda-2)(\lambda-1)\}\\ P^{4}&=\frac{1}{1\cdot2\cdot3}\cdot\frac{\pi a^{\lambda}}{(1-a^{2})^{7}}\{(\lambda+1)(\lambda+2)(\lambda+3)-3a^{2}(\lambda+2)(\lambda+3)(\lambda-3)\\ &\qquad+3a^{4}(\lambda+3)(\lambda-3)(\lambda-2)-a^{6}(\lambda-3)(\lambda-2)(\lambda-1)\}\\ P^{5}&=\frac{1}{1\cdot2\cdot3\cdot4}\cdot\frac{\pi a^{\lambda}}{(1-a^{2})^{9}}\{(\lambda+1)(\lambda+2)(\lambda+3)(\lambda+4)-4a^{2}(\lambda+2)(\lambda+3)(\lambda+4)(\lambda-4)\\ &\qquad+6a^{4}(\lambda+3)(\lambda+4)(\lambda-4)(\lambda-3)-4a^{6}(\lambda+4)(\lambda-4)(\lambda-3)(\lambda-2)\\ &\qquad+a^{8}(\lambda-4)(\lambda-3)(\lambda-2)(\lambda-1)\}\\ &\text{etc.} \end{aligned} \]
The law of these expressions is easy to grasp, and in general if we make
\[ \left.\begin{aligned} P^{n}&=\frac{\pi a^{\lambda}}{(1-a^{2})^{2n-1}}\cdot\frac{\lambda+1\cdot\lambda+2\cdot\lambda+3\ldots\lambda+n-1}{1\cdot2\cdot3\ldots\ldots n-1}A^{n},\\ &\text{the coefficient }A^{n}\text{ will have for its value}\\ A^{n}&=1+\frac{n-1}{1}a^{2}\cdot\frac{n-\lambda-1}{\lambda+1}+\frac{n-1\cdot n-2}{1\cdot2}a^{4}\cdot\frac{n-\lambda-1\cdot n-\lambda-2}{\lambda+1\cdot\lambda+2}\\ &\qquad+\frac{n-1\cdot n-2\cdot n-3}{1\cdot2\cdot3}a^{6}\cdot\frac{n-\lambda-1\cdot n-\lambda-2\cdot n-\lambda-3}{\lambda+1\cdot\lambda+2\cdot\lambda+3}+\text{etc.} \end{aligned}\right\} \tag{3} \]

To make sure of the exactness of this formula, it suffices to substitute the general value of $P^{n}$ in the second member of equation (2); we shall find, after the reductions, a value of $P^{n+1}$ which will be nothing other than what $P^{n}$ becomes on putting $n+1$ in place of $n$. Any other mode of verification would be less easy than that which we indicate.

The simplest case is that where $n=\lambda+1$; then $A^{n}$ reduces to its first term 1, and we have
\[ P^{\lambda+1}=\int\frac{d\varphi\cos\lambda\varphi}{\Delta^{\lambda+1}}=\frac{\pi a^{\lambda}}{(1-a^{2})^{2\lambda+1}}\cdot\frac{\lambda+1\cdot\lambda+2\cdot\lambda+3\ldots2\lambda}{1\cdot2\cdot3\ldots\lambda}. \tag{4} \]
Another case which deserves to be noted, is that of $\lambda=0$; then we have
\[ \begin{aligned} \int\frac{d\varphi}{\Delta^{n}}=\frac{\pi}{(1-a^{2})^{2n-1}}\Big[&1+\left(\frac{n-1}{1}\right)^{2}a^{2}+\left(\frac{n-1\cdot n-2}{1\cdot2}\right)^{2}a^{4}\\ &+\left(\frac{n-1\cdot n-2\cdot n-3}{1\cdot2\cdot3}\right)^{2}a^{6}+\text{etc.}\Big], \end{aligned} \tag{5} \]
a formula whose coefficients are the squares of the coefficients of the binomial raised to the power $n-1$. We would have, for example,

\origpage{375}
\[ \int\frac{d\varphi}{\Delta^{5}}=\frac{\pi}{(1-a^{2})^{9}}[1+4^{2}a^{2}+6^{2}a^{4}+4^{2}a^{6}+a^{8}]. \]

\subsection*{(63).}

Let us come now to the integral $\displaystyle\int\Delta^{n}d\varphi\cos\lambda\varphi$ which we shall designate by $Q^{n}$. Without passing through the particular cases, one can determine all at once the general value of $Q^{n}$. Indeed, if the polynomial $1-2a\cos\varphi+a^{2}$ is decomposed into its two imaginary factors, we shall have
\[ \Delta=(1-ae^{\varphi\sqrt{-1}})(1-ae^{-\varphi\sqrt{-1}}). \]
If next we suppose
\[ \begin{aligned} (1-ae^{\varphi\sqrt{-1}})^{n}={}&1+K_{1}e^{\varphi\sqrt{-1}}+K_{2}e^{2\varphi\sqrt{-1}}\ldots+K_{\lambda}e^{\lambda\varphi\sqrt{-1}}\\ &+K_{\lambda+1}e^{(\lambda+1)\varphi\sqrt{-1}}+\text{etc.} \end{aligned} \]
we shall have similarly
\[ \begin{aligned} (1-ae^{-\varphi\sqrt{-1}})^{n}={}&1+K_{1}e^{-\varphi\sqrt{-1}}+K_{2}e^{-2\varphi\sqrt{-1}}\ldots+K_{\lambda}e^{-\lambda\varphi\sqrt{-1}}\\ &+K_{\lambda+1}e^{-(\lambda+1)\varphi\sqrt{-1}}+\text{etc.} \end{aligned} \]
Multiplying these two equations by one another, in order to have the value of $\Delta^{n}$, and observing that it suffices to retain in the second member the terms affected by $e^{\lambda\varphi\sqrt{-1}}$ and by $e^{-\lambda\varphi\sqrt{-1}}$, we shall have the product sought
\[ \Delta^{n}=(K_{\lambda}+K_{\lambda+1}K_{1}+K_{\lambda+2}K_{2}+K_{\lambda+3}K_{3}+\text{etc.})(e^{\lambda\varphi\sqrt{-1}}+e^{-\lambda\varphi\sqrt{-1}}). \]
But the factor $e^{\lambda\varphi\sqrt{-1}}+e^{-\lambda\varphi\sqrt{-1}}$ reduces to $2\cos\lambda\varphi$, and we have between the given limits $\displaystyle\int d\varphi\cdot2\cos^{2}\lambda\varphi=\pi$; therefore the general value of $Q^{n}$ is
\[ Q^{n}=\pi[K_{\lambda}+K_{\lambda+1}K_{1}+K_{\lambda+2}K_{2}+K_{\lambda+3}K_{3}+\text{etc.}]. \]
Now we have by the binomial formula,
\[ \begin{aligned} K_{\lambda}&=\frac{n\cdot n-1\cdot n-2\ldots n+1-\lambda}{1\cdot2\cdot3\ldots\lambda}(-a)^{\lambda}\\ K_{\lambda+1}&=K_{\lambda}\cdot\frac{n-\lambda}{\lambda+1}(-a) &\qquad K_{1}&=\frac{n}{1}(-a)\\ K_{\lambda+2}&=K_{\lambda}\cdot\frac{n-\lambda\cdot n-\lambda-1}{\lambda+1\cdot\lambda+2}(-a)^{2} & K_{2}&=\frac{n\cdot n-1}{1\cdot2}(-a)^{2}\\ K_{\lambda+3}&=K_{\lambda}\cdot\frac{n-\lambda\cdot n-\lambda-1\cdot n-\lambda-2}{\lambda+1\cdot\lambda+2\cdot\lambda+3}(-a)^{3} & K_{3}&=\frac{n\cdot n-1\cdot n-2}{1\cdot2\cdot3}(-a)^{3}\\ &\text{etc.} && \text{etc.} \end{aligned} \]
Therefore, making $Q^{n}$ or
\[ \int\Delta^{n}d\varphi\cos\lambda\varphi=\pi(-a)^{\lambda}\cdot\frac{n\cdot n-1\cdot n-2\ldots n+1-\lambda}{1\cdot2\cdot3\ldots\lambda}B^{n}, \]

\origpage{376}
we shall have for the value of the coefficient $B^{n}$,
\[ \begin{aligned} B^{n}=1&+\frac{n}{1}a^{2}\cdot\frac{n-\lambda}{\lambda+1}+\frac{n\cdot n-1}{1\cdot2}a^{4}\cdot\frac{n-\lambda\cdot n-\lambda-1}{\lambda+1\cdot\lambda+2}\\ &+\frac{n\cdot n-1\cdot n-2}{1\cdot2\cdot3}a^{6}\cdot\frac{n-\lambda\cdot n-\lambda-1\cdot n-\lambda-2}{\lambda+1\cdot\lambda+2\cdot\lambda+3}+\text{etc.} \end{aligned} \tag{6} \]

\subsection*{(64).}

Let us remark that in the application of this formula one must suppose $n=$ or $>\lambda$. Indeed, if $n$ were $<\lambda$, the development of $\Delta^{n}$ could not give any term similar to $\cos\lambda\varphi$, so the integral $\displaystyle\int\Delta^{n}d\varphi\cos\lambda\varphi$, taken between the given limits, would be zero.

The smallest value that can be given to $n$ is therefore $n=\lambda$; then we have $\displaystyle\int\Delta^{\lambda}d\varphi\cos\lambda\varphi=\pi(-a)^{\lambda}$, which is verified immediately.

The case of $\lambda=0$ deserves also to be noted; then we have
\[ \int\Delta^{n}d\varphi=\pi\left[1+n^{2}a^{2}+\left(\frac{n\cdot n-1}{1\cdot2}\right)^{2}a^{4}+\left(\frac{n\cdot n-1\cdot n-2}{1\cdot2\cdot3}\right)^{2}a^{6}+\text{etc.}\right]. \tag{7} \]
Comparing this formula with formula (5), one draws from it
\[ \int\Delta^{n}d\varphi=(1-a^{2})^{2n+1}\int\frac{d\varphi}{\Delta^{n+1}}, \tag{8} \]
which establishes a remarkable relation between these two integrals, a relation which would hold even if $n$ were not an integer number, and this one will easily verify for the case of $n=\frac{1}{2}$, by the formulas of the first part.

In general if one observes that the value of the coefficient $B^{n}$ is nothing other than that of the coefficient $A^{n}$, in which one would have put $n+1$ instead of $n$, one will conclude that the two integrals designated by $P^{n+1}$ and $Q^{n}$, have between them a very simple relation; we have in fact, from equations (3) and (6), this very remarkable formula
\[ \int\Delta^{n}d\varphi\cos\lambda\varphi=\frac{n\cdot n-1\cdot n-2\ldots n-\lambda+1}{n+1\cdot n+2\cdot n+3\ldots n+\lambda}(-1)^{\lambda}(1-a^{2})^{2n+1}\cdot\int\frac{d\varphi\cos\lambda\varphi}{\Delta^{n+1}} \]
and equation (8) is only a particular case of this one.

All these formulas are due to Euler; but the demonstrations of this illustrious Author are, if I am not mistaken, much less simple than those which I have just set forth. See volume IV of his \emph{Calcul intégral}, pag.~194 et suiv.

END OF THE THIRD PART.

\origpage{377}
% No page number is printed on this first page of the Table (only the signature "48"); 377 by sequence (leaf 384 prints 378).

\section*{TABLE OF CONTENTS.}

\section*{FIRST PART.}

\section*{ON ELLIPTIC FUNCTIONS.}

§ I. \emph{General idea of the different sorts of transcendents contained in the integral formula $\displaystyle\int\frac{P\,dx}{R}$}, pag.~4

It is demonstrated that, $P$ being a rational function of $x$, and $R$ a radical of the form $\sqrt{(\alpha+\beta x+\gamma x^{2}+\delta x^{3}+\varepsilon x^{4})}$, the integral $\displaystyle\int\frac{P\,dx}{R}$ contains only two sorts of transcendents, one of the form $\displaystyle\int(A+Bx+Cx^{2})\frac{dx}{R}$, the other of the form $\displaystyle\int\frac{D}{1+nx}\cdot\frac{dx}{R}$, $A$, $B$, $C$, $D$, $n$ being constant coefficients.

§ II. \emph{Manner of making the odd powers of the variable under the radical disappear}, 7

§ III. \emph{Reduction of the differential $\frac{P\,dx}{R}$ to the form $\frac{Q\,d\varphi}{\sqrt{(1-c^{2}\sin^{2}\varphi)}}$}, 9

The principal advantages of this transformation are, 1°. that the radical $R$, which in its generality contains four arbitrary coefficients, is replaced by the radical $\Delta=\sqrt{(1-c^{2}\sin^{2}\varphi)}$, where there is but one; 2°. that the variable $x$, which might be subjected to certain limits, is replaced by the amplitude $\varphi$ which increases indefinitely; 3°. that the integral, in its new form, can be determined for every value of $\varphi$, if it is known only from $\varphi=0$ to $\varphi=\frac{1}{2}\pi$.

§ IV. \emph{Development of the formula $\displaystyle\int\frac{Q\,d\varphi}{\Delta}$}, 12

It is demonstrated that the formula $\displaystyle\int\frac{Q\,d\varphi}{\Delta}$, where $Q$ is an even rational function of $\sin\varphi$, contains only two sorts of transcendents; one of the form

\origpage{378}
$\displaystyle\int(A+B\sin^{2}\varphi)\frac{d\varphi}{\Delta}$, the other of the form $\displaystyle\int\frac{N}{1+n\sin^{2}\varphi}\cdot\frac{d\varphi}{\Delta}$, $N$ and $n$ having any values whatever, real or imaginary.

§ V. \emph{Definition of the elliptic functions, and their division into three species}, page 14

The simplest of the elliptic functions is represented by $\displaystyle\int\frac{d\varphi}{\Delta}$, it constitutes the first species; the arc of an ellipse designated by $\displaystyle\int\Delta\,d\varphi$ is the second species, and the integral $\displaystyle\int\frac{d\varphi}{(1+n\sin^{2}\varphi)\Delta}$ is the third.

§ VI. \emph{Comparison of elliptic functions of the first species}, 20

There is given, after Euler, the complete algebraic integral of the equation $\frac{d\varphi}{\Delta(\varphi)}+\frac{d\psi}{\Delta(\psi)}=0$, this integral corresponds to the transcendental equation $F(\varphi)+F(\psi)=F(\mu)$; whence it follows that all the problems relating to the addition, subtraction, multiplication and division of the functions $F$ can be solved algebraically, as is done with regard to arcs of a circle.

§ VII. \emph{Application to the lemniscate}, 37

The arcs of this curve represent the function $F(c,\varphi)$ in the case where the modulus $c=\sin45^{\circ}$.

§ VIII. \emph{Application to the motion of the simple pendulum}, 40

The time of motion through any arc whatever, whether the pendulum makes only oscillations, or turns always in the same direction around the point of suspension, is in general an elliptic function of the first species, and enjoys all the properties of this function.

§ IX. \emph{Comparison of elliptic functions of the second species}, 41

The same algebraic equation, which gives $F(\varphi)+F(\psi)-F(\mu)=0$, applies to the functions of the second species, and gives $E(\varphi)+E(\psi)-E(\mu)=$ an algebraic quantity $c^{2}\sin\varphi\sin\psi\sin\mu$. Thus all the comparisons of arcs of an ellipse are established on the same basis as those of the functions $F$.

§ X. \emph{Comparison of arcs of a hyperbola}, 52

§ XI. \emph{Particular development of the formula $\displaystyle Z=\int\frac{(f+gx^{2})\,dx}{\sqrt{(\alpha^{2}+2\alpha\beta x^{2}\cos\theta+\beta^{2}x^{4})}}$}, 54

The application of this formula to various definite integrals, offers different results

\origpage{379}
remarkable results on the relations that exist between the complete functions $F^{1}(c)$, $F^{1}(b)$, $E^{1}(c)$, $E^{1}(b)$, in the case where $c=\sin15^{\circ}$.

§ XII. \emph{Theorem on the complete functions of the first and second species, whose moduli are complements of one another}, page 61

§ XIII. \emph{Differential equations that express the mutual connection of the functions $E$ and $F$}, 62

It is remarked that a table of the functions $E$ would serve to determine the functions $F$, and conversely.

§ XIV. \emph{Development of the functions $F^{1}$ and $E^{1}$ in series}, 65

These functions each depend on a linear differential equation of the second order: they are developed in series in two ways, one according to the powers of the modulus $c$, the other according to the powers of the complementary modulus $b$.

§ XV. \emph{On the changes that may be made to the parameter in elliptic functions of the third species}, 68

Formula for transforming a function of the third species whose parameter is $n$, into another whose parameter is $\frac{c^{2}}{n}$, 69

In elliptic functions of the third species, three forms of the parameter are distinguished; namely, $n=\cot^{2}\theta$, $n=-1+b^{2}\sin^{2}\theta$, $n=-c^{2}\sin^{2}\theta$, 72

Formula for transforming into one another the elliptic functions whose parameters $n$ and $n'$ satisfy the equation $(1+n)(1+n')=b^{2}$, 73

By means of this formula any function whose parameter is of the form $n=\cot^{2}\theta$ can be reduced to a function whose parameter is of the form $n=-1+b^{2}\sin^{2}\theta$, and conversely; but those that relate to the third form $n=-c^{2}\sin^{2}\theta$ are reducible only to functions whose parameter is of the same form, 74

§ XVI. \emph{Comparison of elliptic functions of the third species}, 75

The same algebraic equation, which gives $F(\varphi)+F(\psi)-F(\mu)=0$ and $E(\varphi)+E(\psi)-E(\mu)=c^{2}\sin\varphi\sin\psi\sin\mu$, applies to functions of the third species, and gives $\Pi(\varphi)+\Pi(\psi)-\Pi(\mu)=$ a quantity determinable by arcs of a circle or by logarithms. Thus the functions of the third species are compared like those of the first species, apart from the differences that the nature of these functions entails.

\origpage{380}

§ XVII. \emph{Formation of an infinite sequence of elliptic functions of the first species, connected with one another by constant ratios}, p.~81

The modulus of one of these functions serves to determine all the others; there results a sequence or scale of moduli which extends to infinity in both directions, from zero to unity. Likewise, the amplitude of one of these functions serves to determine all the others, which increase to infinity in one direction, and decrease in the other, up to a determinate limit.

§ XVIII. \emph{Application of the same law to elliptic functions of the second species}, 85

The first result of this application is that the function of the first species $F(c,\varphi)$ can be expressed generally by two functions of the second species $E(c,\varphi)$, $E(c',\varphi')$; whence it follows that any arc of a hyperbola can be determined by two arcs of an ellipse, which is the theorem of Landen.

One can form, according to the same law, an infinite sequence of ellipses such that, supposing the arcs of two of these ellipses to be known, the arcs of all the others can be found.

§ XIX. \emph{Method of approximation applied to elliptic functions of the first species}, 89

Two formulas serve this purpose, one for the case where the modulus $c$ is not too near unity, the other for the case where the difference $1-c$ is extremely small.

The same formulas serve to solve the inverse problem, that is, to determine the amplitude $\varphi$, when the function $F(c,\varphi)$ is known.

§. XX. \emph{Particular properties of the functions $F(c)$, $F(b)$, whose moduli are complements of one another}, 97

These properties lead to singularly close approximations of the number $\pi$ in logarithmic quantities.

§ XXI. \emph{Method of approximation applied to elliptic functions of the second species}, 102

Two formulas are again given on this subject, one for the case where $c$ is not too near unity, the other for the case where the difference $1-c$ is extremely small.

These formulas are the simplest that analysis can offer for the rectification of the ellipse and that of the hyperbola.

§ XXII. \emph{Remarkable formulas for determining the complete functions}

\origpage{381}
\emph{$F^{1}(c)$, $E^{1}(c)$, when $c$ is little removed from one of its limits}, page 113

A table has been added of the logarithms of the functions $F^{1}(c)$, $E^{1}(c)$, calculated for all the angles of the modulus, from degree to degree, from zero up to $90^{\circ}$, 118

§ XXIII. \emph{Method of approximation applied to elliptic functions of the third species}, 119

After having determined the law followed by the parameters and the coefficients of the successive transforms, the general expression of the function is given, for all the cases where the modulus is not too near unity.

The case where the modulus differs very little from unity is then discussed at great length, and the general solution of it is given by a convergent series.

§ XXIV. \emph{On the most general cases in which the reduction of elliptic functions of the third species can be effected}, 134

It is demonstrated generally, 1°. that every complete function of the third species can be expressed by functions of the first and second species.

2°. That the same holds for all the incomplete functions, infinite in number, which can be measured by the complete function.

3°. That the function of the third species $\Pi(n,c,\varphi)$ can be reduced indefinitely to the first species, in an infinity of cases for which one has a general symptom.

§ XXV. \emph{General reduction of elliptic functions whose parameter is imaginary}, 147

It is demonstrated generally that every elliptic function of the third species whose parameter is imaginary can be reduced to two functions of the same species whose parameters are real, one being of the form $-1+b^{2}\sin^{2}\theta$, the other of the form $-c^{2}\sin^{2}\theta$; which fully confirms the division that has been made of elliptic functions into three species.

The case where $n=c(\cos\theta+\sqrt{-1}\sin\theta)$ is discussed in particular, and that where $1+n=b(\cos\lambda+\sqrt{-1}\sin\lambda)$.

§ XXVI. \emph{On a general symptom for recognizing whether two given functions of the third species, which differ only in the parameters, can be reduced one to the other}, 159

§ XXVII. \emph{Example of a particular transformation of elliptic functions of the first and second species}, 161

This example relates to the functions whose moduli are $\sin75^{\circ}$ and $\sin15^{\circ}$, and are consequently complements of one another.

\origpage{382}

§ XXVIII. \emph{On the series that give, without transformations, the approximate values of elliptic functions}, page 169

These series proceed according to the sines of the multiple arcs of the amplitude, and their coefficients can be determined generally by the formulas relating to the complete functions of the first and second species.

§ XXIX. \emph{Surface of the oblique cone}, 173

It is demonstrated that the total surface is always determinable by arcs of an ellipse, whether the base be a circle or an ellipse.

§ XXX. \emph{Construction of the shortest line on the surface of the spheroid}, 178

The latitude of any point of this curve being given, the length of an arc of the curve is found by that of an equal arc of an ellipse; as for the longitude, it depends on an elliptic function of the third species, and one of the first species: it can also be constructed by the development of an oblique cone with circular base. The points where the curve meets the equator and those where it touches the parallels between which it is contained are determined by arcs of an ellipse alone.

§ XXXI. \emph{Determination of the area of the ellipsoid}, 182

The approximate value of this area is first given by a regular and convergent series.

The exact value of the area in elliptic functions is then determined, by the consideration of the lines of greatest and least curvature traced on the surface of the ellipsoid, and it is found that the total area can be expressed by arcs of an ellipse.

§ XXXII. \emph{On some integral formulas that can be brought back to elliptic functions}, 194

§ XXXIII. \emph{On the integral $\displaystyle V=\int\frac{(x^{2}+m)\,dx}{\sqrt{(\alpha+\beta x+\gamma x^{2}+\varepsilon x^{4})}}$}, 200

It is proved that this integral depends only on elliptic functions of the first and second species.

§ XXXIV. \emph{On the integral $\displaystyle Z=\int\frac{dz}{(1+\mu z^{2})\sqrt[3]{(1+\nu z^{2})}}$}, 201

The case of $\nu=3$ and $\mu=\frac{1}{3}$ is examined successively, that of $\nu=-3$ and $\mu=-\frac{1}{3}$, and finally that of $\nu=1$ and $\mu=-\frac{1}{8}$,\ednote{The denominator of this fraction is heavily inked; read as $8$.} all three capable of being resolved by arcs of a circle and logarithms.

§ XXXV. \emph{On the principal cases where one can evaluate, by elliptic}

\origpage{383}
\emph{functions, the integral $\left(\frac{p}{q}\right)=\displaystyle\int x^{p-1}\,dx(1-x^{n})^{\frac{q}{n}-1}$, taken from $x=0$ to $x=1$}, page 209

These cases are those where $n$ is equal to one of the numbers 3, 4, 6, 8 and 12. In the case of $n=12$, one arrives at very remarkable results on the comparison of the complete functions $F^{1}(c)$, $F^{1}(b)$, both among themselves and with the function $F^{1}(\sin 45^{\circ})$, the angle $\theta$ of the modulus $c$ being such that $\sin 2\theta=\operatorname{tang}^{2}15^{\circ}$.

\section*{SECOND PART.}

\section*{ON EULERIAN INTEGRALS.}

§ I. \emph{On Eulerian integrals of the first species}, 222

General formula which comprises nearly the whole theory of these functions, and which may be regarded as a sort of equation in partial finite differences, 227

By means of the auxiliaries denoted by $A_{a}$, the general expression of the functions $\left(\frac{p}{q}\right)$ is given in two different forms, according as $p+q$ is $>$ or $<n$, 230

It is shown that if $n$ is even, the number of auxiliaries can be reduced by half by the formula $(d')$, 237

§ II. \emph{Formula for evaluating the integrals $\left(\frac{p}{q}\right)$ by approximation}, 240

The approximate value of this integral is sought in particular when $p$ and $q$ are very small in relation to $n$.

§ III. \emph{Remark on some particular cases where one can sum the series denoted by $\psi(x)$, and two others of the same species}, 244

Several remarkable theorems are given on these series which are due to \emph{Landen}, and the demonstrations of some of them, which the author suppressed, are added.

§ IV. \emph{Consideration of the integral formulas $\displaystyle\int\frac{x^{p-1}\,dx}{\sqrt[n]{(1-x^{n})^{n-q}}}\log\frac{1}{x}$, $\displaystyle\int\frac{x^{p-1}\,dx}{\sqrt[n]{(1-x^{n})^{n-q}}}\log^{2}\frac{1}{x}$, etc. Very remarkable theorem on the first of these formulas}, 249

This theorem serves to determine exactly the first of these integrals, by

\origpage{384}
supposing the integral $\left(\frac{p}{q}\right)$ known; as for the following formulas, their determination depends in general on the transcendent $(a,n)^{m}$, that is, on the sum of the terms taken from $n$ to $n$, in the series $1+\frac{1}{2^{m}}+\frac{1}{3^{m}}+\frac{1}{4^{m}}+\text{etc.}$

§. V. \emph{On the reduction of the transcendents denoted by $(a,n)^{m}$}, p.~265

The case of $m=2$, $n=6$ is treated in particular, and that of $m=3$, $n=6$; they depend on two transcendents whose expression is sought in convergent series.

§ VI. \emph{On Eulerian integrals of the second species}, 276

Every integral of the first species $\left(\frac{p}{q}\right)$ can be expressed very simply by the functions $\Gamma$; conversely every function $\Gamma(a)$ in which $a$ is a rational number can be expressed by the integrals $\left(\frac{p}{q}\right)$, 283

Expression of the function $\Gamma(k)$ when $k$ is very small, 289

It is proved that the function $\Gamma(a)$ will be known for every value of $a$, if only this function is known from $a=\frac{3}{4}$ to $a=1$, \emph{ibid.}

Formula for calculating directly the value of each function $\Gamma$, 290

This formula is among the \emph{semi-convergent} series; the point at which one must stop in the calculation of this series is sought \emph{a priori}, and the degree of approximation it can give in each case is fixed, 292

General considerations from which are deduced convergent and regular series, to express the values of $\log\Gamma(1+x)$, $\log\Gamma(1-x)$, etc., 298

Application of these series to the general expression of the logarithm of the function $\left(\frac{p}{q}\right)$, 300

The integral $\displaystyle\int t^{n}e^{-t^{m}}\,dt$, taken from $t=0$ to $t=\infty$, is brought back immediately to the functions $\Gamma$, 301

Table of the logarithms of the function $\Gamma(a)$ for every value of $a$, from thousandth to thousandth, from $a=1.000$ to $a=2.000$, 302

\section*{THIRD PART.}

\section*{ON QUADRATURES.}

§ I. \emph{General formula for quadratures}, 308

This formula is composed of an integral in finite differences and a series of corrections whose law is known, 311

\origpage{385}
Means of obviating the inconvenience the formula presents when any of the coefficients $\frac{dy}{dx}$, $\frac{ddy}{dx^{2}}$, $\frac{d^{3}y}{dx^{3}}$ becomes infinite in the part of the curve that one wishes to square, page 313

The result of the general formula comes very near to that obtained by considering the area of the proposed curve as a sum of parabolic areas, each determined by three consecutive and equidistant ordinates, 319

§ II. \emph{Construction of the curve in which the arc $s$ is given as a function of the quantity $\frac{dy}{dx}$}, 320

The expression of each of the coordinates is found generally by an integral in finite differences, joined to a series of corrections whose law is known.

§ III. \emph{Application of the preceding method to the calculation of the trajectory of a projectile}, 330

The details of the calculation are given for finding, in a particular example, the height of the throw, the amplitude of the ascending branch and that of the descending branch. The position of the vertical asymptote and that of the inclined asymptote are then determined.

§ IV. \emph{On the indefinite integral $\displaystyle\int dx\left(\log\frac{1}{x}\right)^{a-1}$, taken from $x=0$}, 339

Two formulas are given for evaluating this integral by approximation, one for all cases, the other for the case only where $x$ is very small. A formula is also given for finding the integral when $x$ is greater than unity.

§ V. \emph{On the integral $\displaystyle\int y\,dx$ taken between two limits which make the function $y$ vanish}, 343

Application to the integral $\displaystyle\int x^{\alpha}e^{-x}\,dx$, taken from $x=0$ to $x=\infty$, $\alpha$ being a positive number. There results the same formula that was found in the second part for the value of $\Gamma(\alpha+1)$, 346

Application to the integral $\displaystyle\int x^{\alpha}(1-x)^{\beta}\,dx$, taken between the limits $x=0$, $x=1$, $\alpha$ and $\beta$ being positive, 348

This integral comprises the general determination of the functions $\left(\frac{p}{q}\right)$, 351

§ VI. \emph{On the integral $\displaystyle\int x^{-\alpha}e^{-x}\,dx$, taken between the imaginary limits which make $x^{-\alpha}e^{-x}$ vanish}, 352\ednote{In this copy the last digit is overwritten by a reader's ink; read as $352$, since the section begins on p.~352 (reviewer's reading from context).}

The value of this integral denoted by $\Gamma'(a)$ is found generally by supposing the function $\Gamma(a)$ known. 353

\origpage{386}
From it is deduced the general expression of the integral $Q(\alpha)=\displaystyle\int d\varphi\cos^{\alpha-2}\varphi\ldots\ \cos(\alpha\operatorname{tang}\varphi-\alpha\varphi)$, taken from $\varphi=0$ to $\varphi=\frac{1}{2}\pi$. It follows that this integral can be determined exactly as a function of $\pi$ and $e$, whenever $2\alpha$ is a whole number, page 355

§ VII. \emph{On the integral $Z=\displaystyle\int\frac{dx\cos ax}{1+xx}$ and others like it, taken from $x=0$ to $x=\infty$.} 357

From the principal formula are deduced two series of integrals which can all be determined as functions of the numbers $e$ and $\pi$. In general the exact value of the integral $\displaystyle\int\left(\frac{M\cos ax+Nx\sin ax}{P}\right)dx$ can be found, from $x=0$ to $x=\infty$, if $M$, $N$, $P$ are rational and integral functions of $x^{2}$, $P$ being the highest, and if moreover $P$ has no factor of the form $x^{2}-m^{2}$, 360

§ VIII. \emph{On the integral $Z=\displaystyle\int e^{-x^{2}}\,dx\cos ax$, taken from $x=0$ to $x=\infty$.} 362

Other integrals that are deduced from the same formula, 363

§ IX. \emph{On the integral $Z(k)=\displaystyle\int x^{\frac{2k-1}{2}}\,dx\,e^{-\left(\frac{1+xx}{2nx}\right)}$, taken from $x=0$ to $x=\infty$}, 364

§ X. \emph{On the integrals $\displaystyle\int x^{a-1}e^{-mx}\,dx\cos nx$, $\displaystyle\int x^{a-1}e^{-mx}\,dx\sin nx$, taken between the limits $x=0$, $x=\infty$}, 367

§ XI. \emph{On the integral $\displaystyle\int\frac{(x^{n}-1)\,dx}{\log x}$ and others like it, taken from $x=0$ to $x=1$}, 370

Very simple law followed by the integrals $\displaystyle\int\left(\frac{x^{n}-1}{\log x}\right)^{k}x^{m-1}\,dx$, by giving to $k$ the successive values 1; 2, 3, etc.

§ XII. \emph{On the integrals $\displaystyle\int\frac{d\varphi\cos\lambda\varphi}{\Delta^{n}}$ and $\displaystyle\int\Delta^{n}\,d\varphi\cos\lambda\varphi$, taken from $\varphi=0$ to $\varphi=\pi$, supposing the numbers $\lambda$ and $n$ whole, and $\Delta=1+a^{2}-2a\cos\varphi$}, 372

END OF THE TABLE.

\origpage{387}
% No page number is printed on the Errata page; 387 by sequence (the Table ends on printed p. 386, leaf 392).

\section*{ERRATA.}

\[
\begin{aligned}
&\text{Pages.} && \text{Lines.} && \text{Faults.} && \text{Corrections.}\\
&26 && 2 && \cot\tfrac{1}{2}\beta\ldots\ldots && \cos\tfrac{1}{2}\beta.\\
&38 && 18 && \text{retranchés}\ldots\ldots && \text{retranchés, multipliés.}\\
&61 && \text{title} && \textit{functions}\ldots\ldots && \textit{complete functions.}\\
&62 && \text{end of the page} && F-c\frac{dE}{dc}\ldots\ldots && E-c\frac{dE}{dc}.\\
& && && b^{2}\left(E+c\frac{dF}{dc}\right)\ldots && b^{2}\left(F+c\frac{dF}{dc}\right).\\
&66 && 2 && -\frac{1^{2}}{2^{2}}c^{2}\ldots\ldots && +\frac{1^{2}}{2^{2}}c^{2}.\\
&\textit{Ibid.} && 3 && -\frac{1^{2}}{2^{2}\cdot4^{2}}3c^{2}\ldots\ldots && -\frac{1^{2}}{2^{2}\cdot4^{2}}3c^{4}.\\
&85 && 15 && (2+c)E(c',\varphi')\ldots && (2+2c)E(c',\varphi').\\
&104 && \text{next to last} && \tfrac{1}{2}B''\sin\varphi''\ldots\ldots && \tfrac{1}{2}B''\sin\varphi'.\\
&156 && 15 && \text{art.\ 58}\ldots\ldots && \text{art.\ 60.}\\
&179 && 15 && CA=A,\ CB=B. && CE=A,\ CP=B.\\
&194 && \text{13 and 15} && am\beta\ldots am'\beta'\ldots && amb\ldots am'b'.\\
&287 && 14 && a=10\ldots\ldots && n=10.\\
&364 && 14 && \frac{1+2z}{2nz}\ldots\ldots && \frac{1+zz}{2nz}.
\end{aligned}
\]


\end{document}
